% !TEX program = xelatex
\documentclass[11pt,openany,fontset=fandol]{ctexbook}
\usepackage[paperwidth=176mm,paperheight=250mm,inner=21mm,outer=18mm,top=20mm,bottom=20mm,headheight=14pt,headsep=6mm,footskip=10mm]{geometry}
\usepackage{amsmath,amssymb,amsthm,bm,mathtools,booktabs,array,enumitem,float}
\usepackage{tikz,fancyhdr,hyperref}
\usetikzlibrary{arrows.meta,calc,patterns,angles,quotes,decorations.pathmorphing,decorations.markings}
\hypersetup{hidelinks,pdftitle={动力学综合：手写讲义修订版},pdfauthor={}}
\setlength{\parindent}{2em}\setlength{\parskip}{3pt}
\setlength{\emergencystretch}{2em}\linespread{1.18}\raggedbottom
\setcounter{tocdepth}{1}
\ctexset{chapter={name={第,讲},number=\chinese{chapter},format=\LARGE\bfseries\centering,beforeskip=16pt,afterskip=22pt},section={format=\Large\bfseries},subsection={format=\large\bfseries}}
\pagestyle{fancy}\fancyhf{}\fancyhead[LE,RO]{\small\thepage}
\fancyhead[LO]{\small 动力学综合}\fancyhead[RE]{\small\nouppercase{\leftmark}}
\renewcommand{\headrulewidth}{0pt}
\fancypagestyle{plain}{\fancyhf{}\fancyfoot[C]{\small\thepage}\renewcommand{\headrulewidth}{0pt}}
\theoremstyle{definition}\newtheorem{example}{例}[chapter]\newtheorem{exercise}{习题}[chapter]
\newdimen\LectureSpaceDim
\newcommand{\LectureNeedspace}[1]{\par\penalty-100\begingroup\LectureSpaceDim=\dimexpr\pagegoal-\pagetotal\relax\ifdim\LectureSpaceDim<#1\relax\newpage\fi\endgroup}
\newcommand{\SolutionStart}{\par\medskip\ifdim\pagetotal>0pt\ifdim\dimexpr\pagegoal-\pagetotal\relax<3\baselineskip\newpage\fi\fi\noindent\textbf{解}\quad\ignorespaces}
\newenvironment{solution}{\SolutionStart}{\par\medskip}
\newcommand{\dd}{\mathop{}\!\mathrm d}\newcommand{\vct}[1]{\bm{#1}}
\tikzset{axis/.style={->,>=Stealth,line width=.4pt},curve/.style={line width=.7pt},guide/.style={dashed,line width=.35pt},every node/.style={font=\small},spring/.style={decorate,decoration={coil,aspect=.45,segment length=4pt,amplitude=3pt}}}

\newcommand{\ee}{\mathrm e}
\tikzset{flow/.style={curve,postaction={decorate},decoration={markings,mark=at position .22 with {\arrow{Stealth}},mark=at position .68 with {\arrow{Stealth}}}}}
\providecommand{\ExpansionDiagram}[3]{\begin{figure}[H]\centering\begin{tikzpicture}#3\end{tikzpicture}\caption{#2}\label{#1}\end{figure}}

% 分册与十编合订本共用的定位页码；局部章、图、题号保持原有对应。
\newcommand{\CoursePartNumber}{3}
\geometry{inner=19.5mm,outer=19.5mm,headheight=16pt}
\AtBeginDocument{%
  \fancyhead{}\fancyhead[C]{\small 动力学综合}%
  \fancyfoot[C]{\small\CoursePartNumber-\thepage}%
  \fancypagestyle{plain}{\fancyhf{}\fancyfoot[C]{\small\CoursePartNumber-\thepage}\renewcommand{\headrulewidth}{0pt}}%
}

\begin{document}\frontmatter
\begin{titlepage}
\centering\vspace*{30mm}
{\Large 物理竞赛二轮讲义\quad 第三编\par}\vspace{8mm}

{\Huge\bfseries 动力学综合\par}
\vspace{12mm}{\Large 物理竞赛二轮讲义\par}
\vspace{10mm}{\large 第一讲至第八讲\par}
\vfill{\large 2026\par}
\end{titlepage}

\begingroup\small\setlength{\parskip}{0pt}\linespread{1}\selectfont\tableofcontents\endgroup\let\cleardoublepage\clearpage\mainmatter
\chapter{能量解法问题}


能量方程给出速度与位置的关系。对单自由度系统，这个关系还可以用来求加速度；对多自由度系统，则须配合动量积分或其他运动方程。计算动能时，质心平动与绕质心转动须分别计入，速度关系则由约束确定。纯滚动中的静摩擦是否做功，要看接触面的运动。

本讲采用地面惯性系。轻绳不可伸长，轻杆视为刚性；滑轮是否有转动惯量、接触面是否允许滑动，均在题设中另行说明。图中的力矢量只表示作用点和方向，长度不代表力的大小。



\section{能量积分、坐标加速度与约束反力}

可动支承和连杆系统中的坐标加速度，往往不同于质点的实际加速度。各质点的位置与速度由约束确定；能量方程给出运动后，反力仍由受力方程求出。



\subsection{由能量积分求运动方程}
\subsubsection{坐标、速度和惯性系数}
设固定理想约束下的系统只有一个自由度，取坐标 $q$。质点的位置为 $\vct r_i(q)$，速度为
\[
\vct v_i=\frac{\dd\vct r_i}{\dd q}\dot q.
\]
若含刚体，还须加上绕质心的转动动能。将全部动能合并后，写成
\begin{equation}\label{sup1:eq:d1-A}
T=\frac12 A(q)\dot q^2,\qquad
A(q)=\sum_i m_i\left|\frac{\dd\vct r_i}{\dd q}\right|^2
+\sum_j I_j\left(\frac{\dd\varphi_j}{\dd q}\right)^2.
\end{equation}
$A(q)$ 是相对于坐标 $q$ 的惯性系数。若 $q$ 为长度，它具有质量量纲；若 $q$ 为角度，它具有转动惯量量纲。换用另一坐标后，$A$ 的表达式随之改变；它由各物体的速度关联决定。

设保守力势能为 $U(q)$，非保守广义力为 $Q$，则功率方程为
\[
\frac{\dd}{\dd t}(T+U)=Q\dot q.
\]
将动能和势能分别对时间求导，
\[
\frac{\dd T}{\dd t}=\frac12A'(q)\dot q^3+A(q)\dot q\ddot q,
\qquad \frac{\dd U}{\dd t}=U'(q)\dot q.
\]
代入功率方程，在 $\dot q\ne0$ 的区间约去 $\dot q$，得
\begin{equation}\label{sup1:eq:d1-one}
A(q)\ddot q+\frac12 A'(q)\dot q^2+U'(q)=Q.
\end{equation}
其中 $\tfrac12A'\dot q^2$ 来自构形变化，常出现在连杆机构和曲线轨道中。即使外力只依赖位置，坐标加速度也可能依赖速度。

式\eqref{sup1:eq:d1-one}也可以由拉格朗日方程直接推出。令 $L=A(q)\dot q^2/2-U(q)$，则
\[
\begin{aligned}
\frac{\partial L}{\partial\dot q}&=A\dot q,
&\frac{\dd}{\dd t}\frac{\partial L}{\partial\dot q}&=A\ddot q+A'\dot q^2,\\
\frac{\partial L}{\partial q}&=\frac12A'\dot q^2-U',
&\frac{\dd}{\dd t}\frac{\partial L}{\partial\dot q}-\frac{\partial L}{\partial q}
&=A\ddot q+\frac12A'\dot q^2+U'.
\end{aligned}
\]
这一推导未除以速度，因此在普通构形的瞬时静止点仍成立。只凭某一瞬间的 $\dot E=0$ 不能约去零速度；应使用邻近运动的连续极限，或直接代入运动方程。若 $A=0$ 或约束本身发生改变，则须重新分析。

\subsubsection{速度平方与转折点}
机械能守恒时
\begin{equation}\label{sup1:eq:d1-speed}
\dot q^2=F(q)=\frac{2[E-U(q)]}{A(q)}.
\end{equation}
在速度不为零的一段运动中，
\[
2\dot q\ddot q=F'(q)\dot q,
\qquad\boxed{\ddot q=\frac12F'(q)}.
\]
其中按商法则，
\[
F'(q)=\frac{-2U'(q)A(q)-2[E-U(q)]A'(q)}{A(q)^2}
=-\frac{2U'}A-\frac{A'}A\dot q^2.
\]
所得加速度与式\eqref{sup1:eq:d1-one}在 $Q=0$ 时一致。这一写法使往返运动的两个速度分支得到同一个加速度。运动方向仍由初始条件决定：$\dot q=\pm\sqrt{F(q)}$，符号在普通转折点改变。

\begin{figure}[H]\centering
\begin{tikzpicture}
\draw[axis](-0.1,0)--(4.3,0)node[right]{$q$};\draw[axis](0,-0.15)--(0,2.5)node[above]{$U$};\draw[curve,domain=0:3.9,samples=140]plot(\x,{.45*(\x-1)^2*(4-\x)+.2});\draw[guide](0,1.15)--(4.05,1.15)node[right]{$E_1$};\draw[guide](0,2)--(4.05,2)node[right]{$E_2$};\node at(1.5,-.75){允许区间由 $U\le E$ 决定};\node at(3.1,2.65){势垒};
\fill(0.249719483,1.15)circle(1.2pt);\draw[guide](0.249719483,0)--(0.249719483,1.15);\node[below]at(0.249719483,0){$q_-$};\fill(2.037053995,1.15)circle(1.2pt);\draw[guide](2.037053995,0)--(2.037053995,1.15);\node[below]at(2.037053995,0){$q_+$};\fill(3,2)circle(1.2pt);\draw[guide](3,0)--(3,2);\node[below]at(3,0){$q_b$};\end{tikzpicture}
\caption{$E_1$ 在势阱内对应转折点 $q_-,q_+$；$E_2$ 等于势垒顶 $q_b$ 处的势能。}\label{sup1:fig:d1-energy}
\end{figure}

若 $U(q_t)=E$ 且 $U'(q_t)\ne0$，则转折点加速度为 $-U'(q_t)/A(q_t)$，粒子在有限时间内到达并反向。若 $U'(q_t)=0$ 且该点是光滑势垒顶点，临界轨道通常需无限长时间接近它。因而“能量等于势垒高度”与“在有限时间越过势垒”须加以区分。

需要时间时，再作积分
\[
t-t_0=\int_{q_0}^{q}\frac{\dd\xi}{\pm\sqrt{F(\xi)}}.
\]
积分上下限和速度分支应随运动方向选取，不能将含转折点的整段运动强行使用同一个正号。


\subsection{坐标加速度与实际受力}
\label{sec:complete1}
单自由度系统中，位置矢量的二阶导数为
\[
 \vct a_i=\vct r_i'(q)\ddot q+\vct r_i''(q)\dot q^2.
\]
第一项随坐标加速度变化，第二项来自速度方向或速度比例随构形的变化。例\ref{ex:hw1-3}中，即使五个质点只受重力和杆的约束力，加速度仍含端点速率的平方。这些速度项由几何关系确定；仅对速度大小求导会漏掉方向变化。

固定理想约束下，允许的虚位移为$\delta\vct r_i=\vct r_i'\delta q$。将各质点的牛顿方程投影到这一允许方向，并相加，约束力的总虚功为零。惯性项满足
\[
 \sum_i m_i\vct a_i\cdot\vct r_i'
 =\left(\sum_i m_i|\vct r_i'|^2\right)\ddot q
 +\frac12\frac{\dd}{\dd q}\left(\sum_i m_i|\vct r_i'|^2\right)\dot q^2.
\]
这里使用了平方模的求导关系
\[
\frac{\dd}{\dd q}|\vct r_i'|^2
=\vct r_i''\cdot\vct r_i'+\vct r_i'\cdot\vct r_i''
=2\vct r_i'\cdot\vct r_i''.
\]
加入刚体的转动惯性项，即得到$A\ddot q+A'\dot q^2/2+U'=Q$。这个推导没有除以$\dot q$，所以普通转向点与释放瞬间也可使用。在这些点使用能量积分时，须沿保持同一约束的连续运动取极限。

非保守广义力由实际作用点求出。若第$j$个外力作用于$P_j$，另有平面力偶$\tau_j$，则
\[
 Q=\sum_j\vct F_j\cdot\frac{\partial\vct r_{P_j}}{\partial q}
   +\sum_j\tau_j\frac{\partial\varphi_j}{\partial q}.
\]
因而$Q\dot q$等于这些力和力偶的实际功率。力作用点与质心不重合时，用质心速度计算其功率会漏掉转动部分。例\ref{ex:hw1-7}的动摩擦就须同时计入平动与自转的功率。

\subsubsection{由运动方程恢复反力}
求得$q,\dot q,\ddot q$后，先按上式恢复各质心加速度，再对需要的子系统列平动和转动方程。对滚轮，$I_C\dot\omega$给出轮地摩擦力矩；对车身与车轮的整体，沿斜面的合力给出共同加速度。将二者联立，才能判断所需静摩擦是否超过$\mu_sN$。

能量式所得的接触解须满足反力与摩擦条件。运动首次达到允许边界时，保留该处的位置和速度，再按改变后的约束列方程。例\ref{ex:hw1-8}中的最高位置还要保留外筒的反冲动能，不能把小筒相对外筒的转向误当成全系统静止。



\LectureNeedspace{4\baselineskip}


\subsection{例题与模型变式}

\begin{example}[固定端轻绳约束下的可动斜面]\label{ex:hw1-1}

质量为$M$、倾角为$\alpha$的斜面放在光滑水平地面上，质量为$m$的物块在斜面上无摩擦滑动。轻绳由物块沿斜面绕过固定在斜面顶端的轻滑轮，再水平连接到右侧固定墙上，如图\ref{fig:hw1-1-q}。斜面和绳足够长，求斜面的加速度。

\end{example}


\begin{figure}[H]\centering
\begin{tikzpicture}[>=Stealth,line cap=round,line join=round]
\draw[curve](0,0)--(4,0)--(3,1.8)--cycle;\node at(2.7,.55){$M$};\draw(.65,0)arc(0:30.963756:.65)node[midway,right]{$\alpha$};\begin{scope}[shift={(1.5,.9)},rotate=30.963756]\draw[fill=white](-.25,0)rectangle(.25,.4);\end{scope}\node[above left]at(1.3971,1.0715){$m$};\draw(1.3971,1.0715)--(2.8971,1.9715)--(5,1.9715);\draw(2.8971,1.9715)circle(.08);\draw[thick](5,1.6)--(5,2.35);\fill[pattern=north east lines](5,1.6)rectangle(5.12,2.35);
\end{tikzpicture}
\caption{固定墙、轻绳与可动斜面}\label{fig:hw1-1-q}
\end{figure}

\begin{solution}
取斜面向右位移为$x$，物块相对斜面向下位移为$s$。图\ref{fig:hw1-1-s}中水平绳段缩短$x$，斜绳段必须伸长同样的量，因此$s=x+\text{常量}$。

\begin{figure}[H]\centering
\begin{tikzpicture}[>=Stealth,line cap=round,line join=round]
\draw[curve](0,0)--(4,0)--(3,1.8)--cycle;\node at(2.7,.55){$M$};\draw(.65,0)arc(0:30.963756:.65)node[midway,right]{$\alpha$};\begin{scope}[shift={(1.5,.9)},rotate=30.963756]\draw[fill=white](-.25,0)rectangle(.25,.4);\end{scope}\node[below right]at(1.58,1.12){$m$};\draw(1.3971,1.0715)--(2.8971,1.9715)--(5,1.9715);\draw(2.8971,1.9715)circle(.08);\draw[thick](5,1.6)--(5,2.35);\fill[pattern=north east lines](5,1.6)rectangle(5.12,2.35);\draw[axis](2,-.4)--(3,-.4)node[above]{$x$};\draw[axis](1.2,1.4)--(.5,.98)node[above left]{$s$};
\end{tikzpicture}
\caption{绳长不变给出两段位移的关系}\label{fig:hw1-1-s}
\end{figure}

物块的地面系速度为
\[
 \vct v_m=\dot x(1-\cos\alpha)\vct e_x-\dot x\sin\alpha\vct e_y,
 \quad v_m^2=2(1-\cos\alpha)\dot x^2.
\]
物块下降$x\sin\alpha$，故
\[
 T=\frac12[M+2m(1-\cos\alpha)]\dot x^2,
 \qquad V=-mgx\sin\alpha+\text{常量}.
\]
惯性系数为常数，广义动量为 $\partial T/\partial\dot x=[M+2m(1-\cos\alpha)]\dot x$，而 $\partial T/\partial x=0$、$\partial V/\partial x=-mg\sin\alpha$。对广义动量求时间导数，运动方程为
\[
[M+2m(1-\cos\alpha)]\ddot x-mg\sin\alpha=0,
\]
所以
\[
 \boxed{a=\ddot x=\frac{mg\sin\alpha}{M+2m(1-\cos\alpha)}}.
\]
为确认轻绳保持绷紧，再沿斜面写物块的运动方程。物块在地面系中的加速度沿斜面向下的分量为$a(1-\cos\alpha)$，故绳张力满足
\[
 mg\sin\alpha-T=m a(1-\cos\alpha).
\]
将能量法求得的$a$代入，得$T=mg\sin\alpha[M+m(1-\cos\alpha)]/[M+2m(1-\cos\alpha)]>0$，与轻绳保持绷紧相容。

加速度与初速度无关。墙通过轻绳向系统施加水平力，因此系统水平动量不守恒。若$M\gg m$，分母近于$M$，所得加速度也相应减小。


物块贴着斜面运动，还要求支持力非负。取斜面外法向为$\vct n=(-\sin\alpha,\cos\alpha)$，由上面的地面系加速度可得$\vct a_m\cdot\vct n=-a\sin\alpha$。图\ref{fig:review1-contact}中张力沿斜面，不参加法向方程，故
\[
 N-mg\cos\alpha=-ma\sin\alpha,
 \qquad
 N=mg\frac{M\cos\alpha-m(1-\cos\alpha)^2}
 {M+2m(1-\cos\alpha)}.
\]
因此，所求贴面运动要求
\[
 \boxed{M\cos\alpha\ge m(1-\cos\alpha)^2}.
\]
等号为法向力恰好降至零的临界状态；若不满足该式，斜面不能向内拉住物块，须从脱离时刻起另列绳约束下的运动方程。例如$\alpha=60^\circ$时，保持接触要求$M\ge m/2$。
\begin{figure}[H]\centering
\begin{tikzpicture}[>=Stealth,line cap=round,line join=round]
\coordinate(O)at(0,0);\draw[guide](-1.7,-.98)--(1.7,.98);\fill(O)circle(2pt)node[below right]{$m$};\draw[axis](O)--(1.25,.72)node[right]{$T$};\draw[axis](O)--(-.72,1.25)node[above]{$N$};\draw[axis](O)--(0,-1.6)node[below]{$mg$};\draw[axis](2.6,.2)--(3.6,.2)node[right]{$a$};\node[align=left]at(3.15,1.15){外法向：$\vct n$};
\end{tikzpicture}
\caption{物块的受力与斜面加速度；外法向上的加速度分量为负}\label{fig:review1-contact}
\end{figure}


\end{solution}



\begin{example}[光滑夹具中的六圆柱堆]\label{ex:hw1-2}

六个均匀实心圆柱质量均为$m$，半径可不同，在质量为$M$的光滑夹具中保持图\ref{fig:hw1-2-q}所示的接触构形，沿倾角为$\theta$的固定斜面运动。夹具足够重，各接触不脱离；圆柱之间及圆柱与斜面间均不滑动。求沿斜面向下的加速度。图仅示意堆放位置，圆柱半径取满足这些接触约束的值。

\end{example}


\begin{figure}[H]\centering
\begin{tikzpicture}[>=Stealth,line cap=round,line join=round]
\begin{scope}[rotate=-20]\draw[thick](-3,0)--(3,0);\draw[curve](-2.5,0)--(-2.5,2.200220688)--(2.5,2.200220688)--(2.5,0);\node at(2.05,1.8){$M$};\foreach \x/\y/\r/\i in {-1.8/.7/.7/1,0/.3/.3/2,1.8/.7/.7/3,-.571236611/.929435250/.55/4,.571236611/.929435250/.55/5,0/1.750220688/.45/6}{\draw(\x,\y)circle(\r);\node at(\x,\y){$\i$};}\end{scope}
\draw[guide](2.81908,-1.02606)--(1.3,-1.02606);
\draw(1.61908,-1.02606)arc(180:160:1.2);\node at(1.35,-.8){$\theta$};
\end{tikzpicture}
\caption{六圆柱与光滑夹具的接触构形}\label{fig:hw1-2-q}
\end{figure}

\begin{solution}
所有圆柱的圆心随夹具以速度$\dot x$平移，$x$沿斜面向下为正。以逆时针为角速度正向，两圆柱接触点的平动速度相同，转动速度相等要求$r_i\omega_i+r_j\omega_j=0$。与斜面接触的底层圆柱满足$r_i\omega_i=-\dot x$；沿接触链传递，便有$r_i|\omega_i|=|\dot x|$，相邻圆柱转向相反。图\ref{fig:hw1-2-s}按$\dot x>0$画出一个圆柱的平动、转动。

\begin{figure}[H]\centering
\begin{tikzpicture}[>=Stealth,line cap=round,line join=round]
\draw[thick](-1,0)--(1,0);
\draw(0,.6)circle(.6);\draw(0,1.8)circle(.6);
\fill(0,.6)circle(1pt);\node[font=\scriptsize]at(-.05,.38){$C_i$};\fill(0,1.8)circle(1pt)node[below]{$C_j$};
\draw[axis](0,.6)--(1.15,.6)node[right]{$\dot x$};
\draw[axis](0,1.8)--(1.15,1.8)node[right]{$\dot x$};
\draw[axis](-.39,.825)arc(150:-90:.45);\node[left]at(-.6,.25){$\omega_i$};
\draw[axis](.39,1.575)arc(-30:210:.45);\node[left]at(-.6,2.15){$\omega_j$};
\fill(0,1.2)circle(1.4pt)node[right]{$P$};
\node[align=left]at(3.6,1.9){$r_i\omega_i+r_j\omega_j=0$\\[5pt]$r_i|\omega_i|=r_j|\omega_j|=|\dot x|$};
\node[align=left]at(3.6,.5){$T_{i,\text{平}}=m\dot x^2/2$\\[5pt]$T_{i,\text{转}}=m\dot x^2/4$};
\end{tikzpicture}
\caption{相邻圆柱在接触点的速度相同，转向相反}\label{fig:hw1-2-s}
\end{figure}

图\ref{fig:review1-graph}把每个圆柱缩为一个顶点，只保留实际接触的圆柱对。底层$1,2,3$与斜面无滑动，其$r_i\omega_i$均为$-\dot x$；从这些顶点出发，沿每一条接触边换一次符号，得到$4,5$为$+\dot x$、$6$为$-\dot x$。图中的闭合接触链有偶数条边，沿不同路径传递约束会得到相同转向。

\begin{figure}[H]\centering
\begin{tikzpicture}[>=Stealth,line cap=round,line join=round]
\foreach \n/\x/\y in {1/0/0,2/2/0,3/4/0,4/1/1.25,5/3/1.25,6/2/2.5}{\coordinate(P\n)at(\x,\y);}
\foreach \a/\b in {1/4,2/4,2/5,3/5,4/6,5/6}{\draw(P\a)--(P\b);}
\foreach \n in {1,2,3,6}{\node[circle,draw,fill=gray!15,minimum size=6mm]at(P\n){$\n$};}
\foreach \n in {4,5}{\node[circle,draw,fill=white,minimum size=6mm]at(P\n){$\n$};}
\node[right,align=left]at(4.5,1.8){灰点：$r_i\omega_i=-\dot x$\\白点：$r_i\omega_i=+\dot x$};
\node[below]at(2,-.45){连线只表示实际接触；$4,5$之间没有接触边};
\end{tikzpicture}
\caption{圆柱接触图与各圆柱的转向}\label{fig:review1-graph}
\end{figure}


每个圆柱的动能为
\[
 \frac12m\dot x^2+\frac12\left(\frac12mr_i^2\right)
 \frac{\dot x^2}{r_i^2}=\frac34m\dot x^2.
\]
半径虽在动能中消去，无滑动条件仍依赖接触构形。以下计算以图中接触关系彼此相容为前提。全系统有
\[
 T=\frac12(M+9m)\dot x^2,\qquad
 V=-(M+6m)gx\sin\theta+\text{常量}.
\]
广义动量为 $(M+9m)\dot x$，对时间求导并与重力广义力 $(M+6m)g\sin\theta$ 相等，得到
\[
(M+9m)\ddot x=(M+6m)g\sin\theta,
\qquad
 \boxed{a=\frac{M+6m}{M+9m}g\sin\theta}.
\]
夹具质量很大时$a\to g\sin\theta$，圆柱转动对总惯性的影响相对减小。接触是否保持、摩擦是否足以维持无滑动，仍须由各接触点的受力另行判断。


\end{solution}



\begin{example}[五质点双杆的速度项]\label{ex:hw1-3}

两根长度均为$l$的轻杆铰接成图\ref{fig:hw1-3-q}的对称构形，两个底端、两个杆中点和顶端各固定一个质量为$m$的质点，底端在光滑水平地面上运动。某瞬间两底端分别以速率$v$向外运动，杆与竖直方向夹角$\theta=30^\circ$。求五个质点的加速度，并说明加速度为何与$v$有关。

\end{example}


\begin{figure}[H]\centering
\begin{tikzpicture}[>=Stealth,line cap=round,line join=round]
\coordinate(A)at(-1.2,0);\coordinate(B)at(1.2,0);\coordinate(C)at(0,2.08);\draw[curve](A)--(C)--(B);\draw[thick](-2,0)--(2,0);\foreach \p in {A,B,C}{\fill(\p)circle(2pt)node[above right]{$\p$};}\fill(-.6,1.04)circle(2pt)node[left]{$D$};\fill(.6,1.04)circle(2pt)node[right]{$E$};\draw[guide](C)--(0,0);\draw(0,1.5)arc(-90:-120:.58)node[midway,below]{$\theta$};
\end{tikzpicture}
\caption{五个等质量质点组成的双杆装置}\label{fig:hw1-3-q}
\end{figure}

\begin{solution}
以对称轴为$y$轴，置$A=(-x,0)$、$B=(x,0)$、$C=(0,y)$。图\ref{fig:hw1-3-s}中$x=l\sin\theta$、$y=l\cos\theta$，两中点坐标分别为$(-x/2,y/2)$、$(x/2,y/2)$。

\begin{figure}[H]\centering
\begin{tikzpicture}[>=Stealth,line cap=round,line join=round]
\coordinate(A)at(-1.2,0);\coordinate(B)at(1.2,0);\coordinate(C)at(0,2.08);\draw[curve](A)--(C)--(B);\draw[thick](-2,0)--(2,0);\foreach \p in {A,B,C}{\fill(\p)circle(2pt)node[above right]{$\p$};}\fill(-.6,1.04)circle(2pt)node[left]{$D$};\fill(.6,1.04)circle(2pt)node[right]{$E$};\draw[guide](C)--(0,0);\draw(0,1.5)arc(-90:-120:.58)node[midway,below]{$\theta$};\draw[axis](1.2,-.3)--(2,-.3)node[below]{$v$};\draw[axis](-1.2,-.3)--(-2,-.3)node[below]{$v$};\node at(0,-.85){$x=l\sin\theta,\quad y=l\cos\theta$};
\draw[axis](-1.2,0)--(-1.2,.7)node[left]{$N$};\draw[axis](1.2,0)--(1.2,.7)node[right]{$N$};
\end{tikzpicture}
\caption{底端位移与杆倾角的几何关系}\label{fig:hw1-3-s}
\end{figure}

五质点的动能、势能为
\[
 T=\frac54m\dot x^2+\frac34m\dot y^2
 =\frac{ml^2}{4}(5\cos^2\theta+3\sin^2\theta)\dot\theta^2,
 \quad V=2mgl\cos\theta.
\]
将动能写成$T=B(\theta)\dot\theta^2/2$，有
\[
\begin{aligned}
B&=\frac{ml^2}{2}(5\cos^2\theta+3\sin^2\theta),\\
B'&=\frac{ml^2}{2}(-10\sin\theta\cos\theta+6\sin\theta\cos\theta)
=-2ml^2\sin\theta\cos\theta,\\
V'&=-2mgl\sin\theta.
\end{aligned}
\]
代入$B\ddot\theta+B'\dot\theta^2/2+V'=0$，整理得
\[
 \ddot\theta=\frac{4(g/l)\sin\theta+
 2\sin\theta\cos\theta\dot\theta^2}
 {5\cos^2\theta+3\sin^2\theta}.
\]
给定瞬间$\dot\theta=v/(l\cos30^\circ)$，所以
\[
 \ddot\theta=\frac{4g}{9l}+\frac{4\sqrt3v^2}{27l^2}.
\]
再对底端和顶点的几何关系各求两次时间导数：
\[
\ddot x=l\cos\theta\,\ddot\theta-l\sin\theta\,\dot\theta^2,
\qquad
\ddot y=-l\sin\theta\,\ddot\theta-l\cos\theta\,\dot\theta^2.
\]
角速度平方项来自杆方向随时间的转动。代入$\theta=30^\circ$以及刚求得的$\dot\theta,\ddot\theta$，得到
\begin{align*}
 a_x:=\ddot x&=\frac{2\sqrt3}{9}g-\frac{4v^2}{9l},\\
 a_y:=\ddot y&=-\frac29g-\frac{20\sqrt3v^2}{27l}.
\end{align*}
因此
\begin{align*}
 \vct a_A&=(-a_x,0),&\vct a_B&=(a_x,0),&\vct a_C&=(0,a_y),\\
 \vct a_D&=(-a_x/2,a_y/2),&\vct a_E&=(a_x/2,a_y/2).
\end{align*}

还须核对两个底端是否仍接地。由对称性，两端支持力均为$N$，五质点竖直加速度的质量加权和为$2m a_y$，故
\[
 2N-5mg=2ma_y,\qquad
 N=\frac{41}{18}mg-\frac{20\sqrt3}{27}\frac{mv^2}{l}.
\]
因此本题接触模型要求$v^2\le123gl/(40\sqrt3)$。若超过该值，地面无法提供负支持力，不能继续沿用底端高度恒为零的解。

速度平方项来自曲线几何的二阶导数。即使重力保持不变，改变瞬时速率也会改变杆所需的法向加速度和约束力；当$v$足够大时，底端虽仍向外运动，其水平加速度可以向内。


\end{solution}



\begin{example}[固定端折杆的支承力]\label{ex:hw1-4}

两根长度为$2l$的轻杆在$C$处铰接，$A$端固定在墙角，$B$端在光滑地面上滑动。$A$端及$AC$中点各有质量$m$的小球，$C$、$CB$中点、$B$各有质量$M$的小球，如图\ref{fig:hw1-4-q}。系统由接近$90^\circ$的角度$\theta_0<\pi/2$从静止开始下落。求$AC$与水平夹角为$\theta$时$B$端支持力，并取$\theta_0\to\pi/2$的极限。

\end{example}


\begin{figure}[H]\centering
\begin{tikzpicture}[>=Stealth,line cap=round,line join=round]
\coordinate(A)at(0,0);\coordinate(C)at(2,1.5);\coordinate(B)at(4,0);\draw[curve](A)--(C)--(B);\draw[thick](0,2.2)--(0,0)--(4.8,0);\foreach \x/\y/\s in {1/.75/m,3/.75/M}{\fill(\x,\y)circle(2pt)node[above]{$\s$};}\node[below]at(A){$A$};\fill(C)circle(2pt)node[above left]{$C$};\node[above right]at(C){$M$};\node[below]at(B){$B$};\draw(.7,0)arc(0:37:.7)node[midway,right]{$\theta$};\fill(A)circle(2pt)node[above left]{$m$};\fill(B)circle(2pt);\node at(4.45,.27){$M$};
\end{tikzpicture}
\caption{固定端折杆及五个质点的质量分布}\label{fig:hw1-4-q}
\end{figure}

\begin{solution}
图\ref{fig:hw1-4-s}中令 $s=\sin\theta$、$c=\cos\theta$，$C=(2lc,2ls)$、$B=(4lc,0)$。左、右中点分别为 $(lc,ls)$ 和 $(3lc,ls)$，固定在 $A$ 的质点没有动能。左中点、$C$、右中点和 $B$ 的坐标逐项求导，速度依次为
\[
l\dot\theta(-s,c),\quad 2l\dot\theta(-s,c),\quad
l\dot\theta(-3s,c),\quad l\dot\theta(-4s,0).
\]
相应的速度平方为
\[
l^2\dot\theta^2,\quad4l^2\dot\theta^2,\quad
l^2(9s^2+c^2)\dot\theta^2,\quad16l^2s^2\dot\theta^2.
\]
分别乘以$m/2,M/2,M/2,M/2$后相加，得到
\[
 T=\frac12l^2K\dot\theta^2,\quad
 V=(m+3M)gls,\quad K=m+5M+24Ms^2.
\]

\begin{figure}[H]\centering
\begin{tikzpicture}[>=Stealth,line cap=round,line join=round]
\coordinate(A)at(0,0);\coordinate(C)at(2,1.5);\coordinate(B)at(4,0);\draw[curve](A)--(C)--(B);\draw[thick](0,2.2)--(0,0)--(4.8,0);\foreach \x/\y/\s in {1/.75/m,3/.75/M}{\fill(\x,\y)circle(2pt)node[above]{$\s$};}\node[below]at(A){$A$};\fill(C)circle(2pt)node[above left]{$C$};\node[above right]at(C){$M$};\node[below]at(B){$B$};\draw(.7,0)arc(0:37:.7)node[midway,right]{$\theta$};\fill(A)circle(2pt)node[above left]{$m$};\fill(B)circle(2pt);\node at(4.45,.27){$M$};\draw[axis](4,0)--(4,1.05)node[right]{$N$};\draw[axis](-.45,.15)arc(30:290:.32);\node at(-.8,-.6){$L_A$};
\end{tikzpicture}
\caption{绕固定点取矩求底端支持力}\label{fig:hw1-4-s}
\end{figure}

记$s_0=\sin\theta_0$，初态能量为$(m+3M)gls_0$。因此
\[
\dot\theta^2=\frac{2(m+3M)g(s_0-s)}{lK}.
\]
为得到也适用于释放瞬间的角加速度，使用单自由度方程。由$K'=48Msc$，
\[
l^2K\ddot\theta+24Ml^2sc\dot\theta^2+(m+3M)glc=0.
\]
代入速度平方，得
\[
\ddot\theta=-\frac{(m+3M)gc[K+48Ms(s_0-s)]}{lK^2}.
\]
支持力不直接出现在能量中。对整个系统绕固定点$A$取矩，$A$端作用力的力矩为零。由各球坐标与速度逐项计算，有
\[
 L_A=(m+7M)l^2\dot\theta,\qquad
 \dot L_A=4lcN-(m+9M)glc.
\]
上述角动量的四个活动质点贡献依次为
\[
 ml^2\dot\theta,\quad4Ml^2\dot\theta,\quad
 3Ml^2\dot\theta,\quad0.
\]
最后一个质点$B$始终沿通过$A$的水平线运动，故其绕$A$的角动量为零。动能中的$K$随构形变化，而角动量的系数$m+7M$保持不变；两者不能互相代替。

当$c\ne0$时，代入角加速度得到
\[
 \boxed{N=\frac g4\left[(m+9M)-
 \frac{(m+7M)(m+3M)[K+48Ms(s_0-s)]}{K^2}\right]}.
\]
令$\theta_0\to\pi/2$，即$s_0\to1$，便得到由近竖直位置释放的极限公式。竖直构形处的支持力也应由此连续运动取极限；在原取矩式中直接置$c=0$，已无法再约去该因子。若恰从$\theta_0=\pi/2$且角速度严格为零开始，理想运动方程还保留静止解，不能据此断言系统自行下落。若计算得$N<0$，说明地面接触约束已经失效。上述下落解只适用于$N\ge0$的阶段。


\end{solution}



\begin{example}[抛物线轨道中的变惯性系数]\label{sup1:problem-example-2}\label{sup1:example:d1-e2}
质点 $m$ 套在固定的光滑细轨道上，轨道位于竖直平面内，方程为 $y=x^2/(2a)$，$a>0$。质点从 $x=x_0>0$ 处静止释放。求任意位置的 $\dot x^2$、$\ddot x$ 以及轨道对质点的法向力大小。
\end{example}

\begin{figure}[H]\centering
\begin{tikzpicture}
\draw[axis](-2.1,0)--(2.6,0)node[right]{$x$};\draw[axis](0,-0.3)--(0,2.8)node[above]{$y$};\draw[curve,domain=-1.9:1.9,samples=100]plot(\x,{.5*(\x)^2});\fill(1.00000,0.50000)circle(1.3pt)node[right]{$P$};
\node at(-1.15,2.25){$y=x^2/(2a)$};\fill(1.70000,1.44500)circle(1.3pt)node[right]{$P_0$};
\draw[guide](1.7,1.445)--(1.7,0)node[below]{$x_0$};
\end{tikzpicture}
\caption{以横坐标描述抛物线轨道，释放点为 P$_0$。}\label{sup1:fig:d1-e2-q}
\end{figure}
\par\medskip
\begin{solution}
由图\ref{sup1:fig:d1-e2-q}的轨道方程，$\dot y=(x/a)\dot x$，所以
\[
T=\frac m2\left(1+\frac{x^2}{a^2}\right)\dot x^2,
\quad U=\frac{mgx^2}{2a},\quad E=\frac{mgx_0^2}{2a}.
\]
于是
\begin{equation}\label{sup1:eq:d1-parabola}
\boxed{\dot x^2=\frac{g(x_0^2-x^2)}{a(1+x^2/a^2)}}.
\end{equation}
在初次下降阶段取负的速度分支。这里$M(x)=m(1+x^2/a^2)$，$M'(x)=2mx/a^2$、$U'(x)=mgx/a$。把这些量代入式\eqref{sup1:eq:d1-one}，即
\[
M(x)\ddot x+\frac12M'(x)\dot x^2+U'(x)=0
\]
并解出加速度，得到
\[
\boxed{\ddot x=-\frac{gx/a+(x/a^2)\dot x^2}{1+x^2/a^2}}.
\]
这个过程不需要在释放点除以$\dot x=0$。将速度平方进一步代回，还可化为只含位置的表达式：
\[
\ddot x=-\frac{ga x(a^2+x_0^2)}{(a^2+x^2)^2}.
\]
在右半轨道 $x>0$ 时，水平加速度向左；通过最低点后 $x<0$，它反向向右。

分子中的速度平方项来自 $M\prime(x)$。以 $x$ 描述曲线运动时，即使支持力不做功，它仍会影响地面水平方向的加速度。

\begin{figure}[H]\centering
\begin{tikzpicture}
\draw[axis](-2.1,0)--(2.6,0)node[right]{$x$};\draw[axis](0,-0.3)--(0,2.8)node[above]{$y$};\draw[curve,domain=-1.9:1.9,samples=100]plot(\x,{.5*(\x)^2});\fill(1.00000,0.50000)circle(1.3pt)node[right]{$P$};
\node at(-1.15,2.25){$y=x^2/(2a)$};\draw[axis](1.00000,0.50000)--(0.30000,1.20000)node[above right]{$N$};
\draw[axis](1.00000,0.50000)--(1.00000,-0.50000)node[below]{$mg$};
\draw[axis](1.00000,0.50000)--(0.30000,-0.20000)node[below left]{$\vct v$};
\draw[guide](.45,-.05)--(1.8,1.3);\node[right]at(1.95,1.4){切线};
\end{tikzpicture}
\caption{P 点的速度沿切线，法向力指向轨道凹侧。}\label{sup1:fig:d1-e2-s}
\end{figure}

图\ref{sup1:fig:d1-e2-s}中 P 点的切线斜率为 $p=y'=x/a$，二阶导数为 $y''=1/a$。指向凹侧的单位法向量为
\[
\vct n=\frac{(-p,1)}{\sqrt{1+p^2}}.
\]
曲率和实际速率分别为
\[
\kappa_c=\frac{1/a}{(1+p^2)^{3/2}},\qquad
v^2=(1+p^2)\dot x^2.
\]
沿法向投影，$N-mg/\sqrt{1+p^2}=m\kappa_c v^2$，故
\[
\boxed{N=\frac{m(g+\dot x^2/a)}{\sqrt{1+x^2/a^2}}}.
\]
质点将在$-x_0$与$x_0$之间往复运动，法向力始终指向凹侧。
法向力也可整理为
\[
N=\frac{mga(a^2+x_0^2)}{(a^2+x^2)^{3/2}}.
\]
在两个转向点 $x=\pm x_0$，速率为零，但重力仍有法向分量；到达最低点时，向心加速度使支持力增大。两处的数值分别为
\[
N(\pm x_0)=\frac{mg}{\sqrt{1+x_0^2/a^2}},\qquad
N(0)=mg\left(1+\frac{x_0^2}{a^2}\right).
\]
这与图\ref{sup1:fig:d1-e2-s}中法向加速度指向凹侧相符，且全过程 $N>0$。

能量式只给出 $\dot x^2$。初次下降取 $\dot x<0$，到达 $-x_0$ 后应改取 $\dot x>0$，不能把负根用于整个往复过程。在两个转向点，运动方程仍给出指向轨道最低处的加速度，质点由此沿原轨道返回。

\end{solution}



\section{纯滚动系统的动能与接触条件}

带轮小车的平动与车轮转动由接触条件联系。只有确认各处均无滑动，才能把转动动能折合到同一个惯性系数中；所得运动还须满足静摩擦上限。



\subsection{纯滚动系统的动能与约束力}
\subsubsection{固定接触面上的静摩擦}
刚体在固定表面上纯滚动时，接触点的瞬时速度为零，静摩擦力的瞬时功率也是零。质心所受摩擦力可以不为零，质心平动动能与绕质心的转动动能之间仍会发生能量分配。

对半径 $r$、质量 $m$、绕轴转动惯量 $I=\kappa mr^2$ 的滚动物体，若质心沿固定直线运动且 $v=r\omega$，则
\[
T=\frac12mv^2+\frac12I\omega^2
=\frac12m(1+\kappa)v^2.
\]
接触面本身运动时，先在地面系写出质心速度，再由接触点的相对速度确定自转角速度。此时动能系数须重新计算。

\subsubsection{理想约束与实际接触条件}
能量法消去静摩擦力后，仍须用转动方程恢复所需摩擦力，并与 $|f_s|\le\mu_sN$ 比较。支持力还应满足 $N\ge0$。这两项限制决定所求轨道能否由实际接触实现。


\subsection{例题与模型变式}

\begin{example}[固定斜面与可动斜面上的带轮小车]\label{ex:hw1-5}

倾角为$\theta$的斜面上有一小车，车身质量为$M$，两个车轮均为均匀实心圆柱，每轮质量$m$、半径$r$，轮轴无摩擦，轮与斜面间纯滚动。(1)斜面固定，求小车下坡加速度；(2)斜面质量为$M^\prime$，可在光滑水平地面上平移，求小车相对斜面的加速度和斜面加速度。

\end{example}


\begin{figure}[H]\centering
\begin{tikzpicture}[>=Stealth,line cap=round,line join=round]
\draw[curve](0,0)--(4,0)--(0,1.6)--cycle;\node at(1,.35){$M^\prime$};\begin{scope}[shift={(1.6,.96)},rotate=-21.801409]\draw[fill=white](-.7,.4)rectangle(.7,.95);\node at(0,.68){$M$};\draw(-.45,.2)circle(.2);\draw(.45,.2)circle(.2);\node[right]at(.8,.4){$m,r$};\end{scope}\draw[thick](-.5,-.05)--(4.8,-.05);
\draw(3.45,0)arc(180:158.198591:.55);\node at(2.95,.25){$\theta$};
\end{tikzpicture}
\caption{带轮小车与可平移的斜面}\label{fig:hw1-5-q}
\end{figure}

\begin{solution}
令 $q$ 为小车相对斜面的下坡位移，$X$ 为斜面向右的位移，$B=M+2m$ 为小车总质量。图\ref{fig:hw1-5-s}中的速度合成为 $\vct v_C=(\dot X+\dot q\cos\theta,-\dot q\sin\theta)$，轮的绝对角速度大小为 $|\dot q|/r$。

\begin{figure}[H]\centering
\begin{tikzpicture}[>=Stealth,line cap=round,line join=round]
\draw[curve](0,0)--(4,0)--(0,1.6)--cycle;\node at(1,.35){$M^\prime$};\begin{scope}[shift={(1.6,.96)},rotate=-21.801409]\draw[fill=white](-.7,.4)rectangle(.7,.95);\node at(0,.68){$M$};\draw(-.45,.2)circle(.2);\draw(.45,.2)circle(.2);\node[right]at(.8,.4){$m,r$};\end{scope}\draw[thick](-.5,-.05)--(4.8,-.05);\draw[axis](1.5,2)--(2.5,1.6)node[above]{$\dot q$};\draw[axis](1,-.55)--(2,-.55)node[above]{$\dot X$};
\draw(3.45,0)arc(180:158.198591:.55);\node at(2.95,.25){$\theta$};
\end{tikzpicture}
\caption{小车相对斜面的速度及斜面速度}\label{fig:hw1-5-s}
\end{figure}

固定斜面时，车轮平动已计入$B$，再加两个车轮的转动动能：
\[
 T=\frac12(M+3m)\dot q^2,\quad V=-Bgq\sin\theta,
 \quad\boxed{\ddot q=\frac{M+2m}{M+3m}g\sin\theta}.
\]
可动斜面时
\[
 \mathcal L=\tfrac12(M^\prime+B)\dot X^2+B\cos\theta\dot X\dot q
 +\tfrac12(B+m)\dot q^2+Bgq\sin\theta.
\]
分别对两个广义速度求偏导，
\[
\frac{\partial\mathcal L}{\partial\dot q}=(B+m)\dot q+B\cos\theta\dot X,
\qquad
\frac{\partial\mathcal L}{\partial\dot X}=(M^\prime+B)\dot X+B\cos\theta\dot q.
\]
再用 $\partial\mathcal L/\partial q=Bg\sin\theta$、$\partial\mathcal L/\partial X=0$，列式为
\[
 (B+m)\ddot q+B\cos\theta\ddot X=Bg\sin\theta,
 \quad(M^\prime+B)\ddot X+B\cos\theta\ddot q=0.
\]
置$D=(B+m)(M^\prime+B)-B^2\cos^2\theta>0$，解得
\[
 \boxed{\ddot q=\frac{B(M^\prime+B)g\sin\theta}{D},\qquad
 \ddot X=-\frac{B^2g\sin\theta\cos\theta}{D}}.
\]
第二个方程也可由水平动量守恒得到：地面只给斜面竖直支持力，因此$P_x=(M^\prime+B)\dot X+B\cos\theta\dot q$为常量。消去$\ddot X$后，小车沿$q$的惯性系数为
\[
 B+m-\frac{B^2\cos^2\theta}{M^\prime+B}=\frac{D}{M^\prime+B}>0.
\]
它小于斜面固定时的$B+m$，所以可动斜面会增大小车相对下坡的加速度。

若所问为小车在地面系中的加速度，则应再组成$(\ddot X+\ddot q\cos\theta,-\ddot q\sin\theta)$。让$M^\prime\to\infty$，第二种模型回到第一种。


\end{solution}



\begin{example}[两组车轮的转动惯性]\label{ex:hw1-6}

将上一题的斜面改为一辆大车：大车车身质量$M^\prime$，小车车身质量$M$，两车各有两个相同的实心圆柱轮，每轮质量$m$、半径$r$。大车轮在水平地面纯滚动，小车轮在大车的倾角$\theta$斜面上纯滚动。求小车相对斜面的下坡加速度$a_1$与大车向右的加速度$a_2$。

\end{example}


\begin{figure}[H]\centering
\begin{tikzpicture}[>=Stealth,line cap=round,line join=round]
\draw[curve](0,0)--(4,0)--(0,1.6)--cycle;\node at(1,.35){$M^\prime$};\begin{scope}[shift={(1.6,.96)},rotate=-21.801409]\draw[fill=white](-.7,.4)rectangle(.7,.95);\node at(0,.68){$M$};\draw(-.45,.2)circle(.2);\draw(.45,.2)circle(.2);\node[right]at(.8,.4){$m,r$};\end{scope}\draw(.6,-.2)circle(.2);\draw(3.2,-.2)circle(.2);\draw[thick](-.5,-.4)--(4.8,-.4);
\draw(3.45,0)arc(180:158.198591:.55);\node at(2.95,.25){$\theta$};
\end{tikzpicture}
\caption{大小两车的车身与四个车轮}\label{fig:hw1-6-q}
\end{figure}

\begin{solution}
仍用图\ref{fig:hw1-6-s}中的$q,X$。小车总平动质量$B=M+2m$；大车自身总平动质量为$M^\prime+2m$。四个轮的转动动能必须按各自承托面计算。

\begin{figure}[H]\centering
\begin{tikzpicture}[>=Stealth,line cap=round,line join=round]
\draw[curve](0,0)--(4,0)--(0,1.6)--cycle;\node at(1,.35){$M^\prime$};\begin{scope}[shift={(1.6,.96)},rotate=-21.801409]\draw[fill=white](-.7,.4)rectangle(.7,.95);\node at(0,.68){$M$};\draw(-.45,.2)circle(.2);\draw(.45,.2)circle(.2);\node[right]at(.8,.4){$m,r$};\end{scope}\draw(.6,-.2)circle(.2);\draw(3.2,-.2)circle(.2);\draw[thick](-.5,-.4)--(4.8,-.4);\draw[axis](1.5,2)--(2.5,1.6)node[above]{$a_1$};\draw[axis](1,-.8)--(2,-.8)node[above]{$a_2$};\node at(2,-1.4){$C=M^\prime+M+5m$};
\draw(3.45,0)arc(180:158.198591:.55);\node at(2.95,.25){$\theta$};
\end{tikzpicture}
\caption{两组滚动约束；$a_1,a_2$箭头表示所取正向，实际$a_2<0$}\label{fig:hw1-6-s}
\end{figure}

小车两轮的转动项为$m\dot q^2/2$，大车两轮的转动项为$m\dot X^2/2$。因此
\[
 T=\frac12(M^\prime+M+5m)\dot X^2
 +B\cos\theta\dot X\dot q+\frac12(M+3m)\dot q^2,
 \quad V=-Bgq\sin\theta.
\]
令 $C=M^\prime+M+5m$、$A=M+3m$。两个广义动量分别为 $\partial T/\partial\dot q=A\dot q+B\cos\theta\dot X$、$\partial T/\partial\dot X=C\dot X+B\cos\theta\dot q$。对时间求导，并用 $V_q=-Bg\sin\theta$、$V_X=0$，得到
\[
 A\ddot q+B\cos\theta\ddot X=Bg\sin\theta,
 \qquad C\ddot X+B\cos\theta\ddot q=0.
\]
消去$\ddot X$即得
\[
 \boxed{a_1=\frac{BCg\sin\theta}{AC-B^2\cos^2\theta},\quad
 a_2=-\frac{B^2g\sin\theta\cos\theta}{AC-B^2\cos^2\theta}}.
\]
与上一题相比，$C$多出大车两轮的平动质量$2m$及转动惯性折合质量$m$。循环坐标对应的守恒量为$C\dot X+B\cos\theta\dot q$；实际水平动量的$\dot X$系数只有$M^\prime+M+4m$，二者相差$m\dot X$，这正是地面静摩擦传递水平动量的结果。


图\ref{fig:hw1-6-s}中 $a_2$ 的箭头表示选取的正方向。对于 $0<\theta<\pi/2$，$AC-B^2\cos^2\theta>0$，因而 $a_1>0$、$a_2<0$：小车沿斜面向右下方运动，大车实际向左加速。这里的大车轮角速度满足 $r\omega_{\text{大}}=-\dot X$，小车轮相对斜面的滚动角速度大小为 $|\dot q|/r$；把后一式误写成小车地面速度除以 $r$，会使动能交叉项和最终加速度同时出错。

\end{solution}



\begin{example}[斜面纯滚动与最小静摩擦系数]\label{sup1:problem-example-4}\label{sup1:example:d1-e4}
质量为 $m$、半径为 $r$、转动惯量为 $I=\kappa mr^2$ 的轴对称刚体，从固定斜面上静止释放。斜面倾角为 $\alpha$。先假定刚体始终纯滚动，求质心加速度、下降竖直高度 $h$ 后的速率，以及该假定所需的最小静摩擦系数。
\end{example}

\begin{figure}[H]\centering
\begin{tikzpicture}
\begin{scope}[rotate=-24]\draw[curve](-1.6,0)--(3.4,0);\draw[fill=white,curve](0,.7)circle(.7);\fill(0.00000,0.70000)circle(1.3pt)node[above left]{$C$};
\fill(0.00000,0.00000)circle(1.3pt)node[below left]{$P$};
\draw(0,.7)--(0.520201378,1.168391424)node[midway,above]{$r$};\draw[axis](1.20000,0.70000)--(2.10000,0.70000)node[above]{$s$};
\node at(1.8,1.4){$I=\kappa mr^2$};\end{scope}\draw[guide](1.3060545559848429,-1.3829045864577207)--(3.306054555984843,-1.3829045864577207);\draw(2.456054555984843,-1.3829045864577207)arc(180:156:.65);\node at(2.126054555984843,-1.1529045864577208){$\alpha$};
\end{tikzpicture}
\caption{刚体沿固定斜面滚动，P 为瞬时接触点。}\label{sup1:fig:d1-e4-q}
\end{figure}
\par\medskip
\begin{solution}
本题先用纯滚动约束计算运动，再由受力方程核对约束能否实现。绕质心的转动惯量为 $\kappa mr^2$；纯滚动使平动速率与自转角速度的大小满足 $v=r|\omega|$，所以自转动能可以改写为 $\kappa mv^2/2$。下降势能应同时分配给平动和自转。

取沿斜面向下的位移 $s$，纯滚动关系为 $\dot\varphi=\dot s/r$，其中 $\varphi$ 以沿下行滚动方向为正。由图\ref{sup1:fig:d1-e4-q}，质心下降高度为 $s\sin\alpha$，故
\[
\frac12m(1+\kappa)\dot s^2-mgs\sin\alpha=0.
\]
于是
\[
\boxed{v^2=\frac{2gh}{1+\kappa}}.
\]

\begin{figure}[H]\centering
\begin{tikzpicture}
\begin{scope}[rotate=-24]\draw[curve](-1.6,0)--(3.4,0);\draw[fill=white,curve](0,.7)circle(.7);\fill(0.00000,0.70000)circle(1.3pt)node[above left]{$C$};
\fill(0.00000,0.00000)circle(1.3pt)node[below left]{$P$};
\draw(0,.7)--(0.520201378,1.168391424)node[midway,above]{$r$};\draw[axis](0.00000,0.70000)--(1.00000,0.70000)node[above]{$a$};
\draw[axis](0.00000,0.00000)--(-1.05000,0.00000)node[above]{$f_s$};
\draw[axis](0.00000,0.00000)--(0.00000,1.85000)node[above left]{$N$};
\end{scope}\draw[guide](1.3060545559848429,-1.3829045864577207)--(3.306054555984843,-1.3829045864577207);\draw(2.456054555984843,-1.3829045864577207)arc(180:156:.65);\node at(2.126054555984843,-1.1529045864577208){$\alpha$};\draw[axis](0.2847157,0.6394818)--(0.2847157,-0.8105182)node[below right]{$mg$};

\end{tikzpicture}
\caption{支持力作用于接触点，静摩擦沿上坡方向；重力作用于质心。}\label{sup1:fig:d1-e4-s}
\end{figure}

图\ref{sup1:fig:d1-e4-s}中，重力和法向力对质心均不产生绕轴力矩，使刚体自转加快的力矩由沿斜面向上的静摩擦提供。分别列平动与绕质心的转动方程，
\[
mg\sin\alpha-f_s=ma,\qquad f_s r=I\frac a r.
\]
第二式给$f_s=\kappa ma$，代入第一式即有$(1+\kappa)ma=mg\sin\alpha$。因此
\[
\boxed{a=\frac{g\sin\alpha}{1+\kappa}},\qquad
N=mg\cos\alpha.
\]
上述消元没有约去速度，适用于释放瞬间。所需静摩擦及其上限条件为
\[
f_s=\frac{\kappa}{1+\kappa}mg\sin\alpha,
\qquad
\boxed{\mu_s\ge\frac{\kappa}{1+\kappa}\tan\alpha}.
\]
均匀实心圆柱 $\kappa=1/2$，薄圆环 $\kappa=1$。同样的竖直落差下，转动惯量系数较大的物体质心速率较小，更多能量分配给自转。纯滚动接触不产热。摩擦在平动与转动之间传递能量，整体能量收支中不再另扣 $f_s s$。
相应能量分配为
\[
T_{\text{平动}}=\frac{mgh}{1+\kappa},\qquad
T_{\text{转动}}=\frac{\kappa mgh}{1+\kappa}.
\]
图\ref{sup1:fig:d1-e4-s}的摩擦对质心平动贡献功率 $-f_s v$，对自转贡献功率 $f_s r|\omega|=f_s v$，二者相消。

当 $\kappa=1/2$，$a=2g\sin\alpha/3$、$v^2=4gh/3$、$\mu_{s,\min}=\tan\alpha/3$；当 $\kappa=1$，三者分别为 $g\sin\alpha/2$、$gh$、$\tan\alpha/2$。这些数值均须先满足各自的静摩擦条件，才是实际运动解。

\end{solution}



\section{多自由度、循环坐标与整体反冲}

内部运动会带动可动支承反冲。对称性对应的守恒量允许消去部分速度，但内外两物体的接触速度关系仍须保留。



\subsection{多自由度系统与循环坐标}
\subsubsection{交叉动能与第二个方程}
当斜面也能移动时，小车轮心速度是斜面平移速度与相对斜面速度的矢量和。两者一般不正交，动能中会出现 $\dot X\dot s$ 项。两个自由度只有一个能量积分，通常不足以求解；还须使用对应的运动方程或守恒量。

对固定理想约束，拉格朗日函数为 $L=T-U$，
\begin{equation}\label{sup1:eq:d1-lagrange}
\frac{\dd}{\dd t}\frac{\partial L}{\partial\dot q_j}
-\frac{\partial L}{\partial q_j}=Q_j.
\end{equation}
若 $L$ 不显含某个坐标 $X$ 且 $Q_X=0$，则 $P=\partial L/\partial\dot X$ 守恒。它是与所选坐标共轭的广义动量。若该坐标只表示所有质点的同向平移，它等于该方向的总平动动量；若改变该坐标还带动车轮转动，则须计入相应的转动项。

\subsubsection{消去循环速度}
常见的动能形式为
\[
T=\frac12D\dot X^2+C(q)\dot X\dot q+\frac12B(q)\dot q^2,
\qquad U=U(q),
\]
其中 $D$ 为正常量。守恒量与配方结果为
\begin{align*}
P&=D\dot X+C(q)\dot q,\\
T&=\frac{P^2}{2D}+\frac12\left[B(q)-\frac{C(q)^2}{D}\right]\dot q^2.
\end{align*}
令 $B_{\text{eff}}=B-C^2/D$。无非保守力做功时，能量积分化为
\begin{equation}\label{sup1:eq:d1-reduced}
\frac12B_{\text{eff}}(q)\dot q^2+U(q)
=E-\frac{P^2}{2D}.
\end{equation}
$B_{\text{eff}}>0$ 是动能正定性的要求。在这一常 $D$ 情形，内坐标仍满足变惯性系数方程。求出 $q(t)$ 后，再由 $\dot X=[P-C(q)\dot q]/D$ 恢复整体运动。

直接把循环速度代回原拉格朗日函数，然后当作新的拉格朗日函数变分，一般会得到错误的方程。若要在拉格朗日层面消元，应先作劳斯变换 $\mathcal R=L-P\dot X$，再代入循环速度：
\[
\begin{aligned}
\mathcal R&=\frac12B\dot q^2-U-\frac{(P-C\dot q)^2}{2D}\\
&=\frac12B_{\rm eff}\dot q^2-U-\frac{P^2}{2D}+\frac{PC}{D}\dot q.
\end{aligned}
\]
最后一项是 $\dd[\int^q PC(\xi)/D\,\dd\xi]/\dd t$，不改变内坐标方程；常量项也不参与变分。因此约化后的方程确为
\[
B_{\rm eff}\ddot q+\frac12B_{\rm eff}'\dot q^2+U'=Q_q,
\qquad B_{\rm eff}'=B'-\frac{2CC'}D,
\]
其中保留的非保守广义力为 $Q_q$，且消元仍要求 $Q_X=0$。本讲用式\eqref{sup1:eq:d1-reduced}求速度，并与原来的两条方程对应。


\subsection{内滚圆筒与整体反冲}
圆筒在另一圆筒内部滚动时，圆心的公转与自身转动由接触条件联系。若外筒也能沿地面滚动，还会出现第三个角速度；各角速度应先统一正方向，再从接触点速度相等建立关系。直接套用 $v=r\omega$ 容易漏掉外筒的运动。


\subsection{例题与模型变式}

\begin{example}[小车上的有限振幅摆]\label{sup1:problem-example-8}\label{sup1:example:d1-e8}
质量为 $M$ 的小车可在光滑水平轨道上平移，车轮惯量忽略。车上固定铰点 O，通过长为 $l$ 的刚性轻杆连接质点 $m$。取摆角 $\theta$ 从竖直向下方向量起。系统总水平动量为零。在 $\theta=0$ 处给予系统一组内力冲量，使摆角速度为 $\omega_0>0$。求 $\dot\theta$ 与 $\theta$ 的关系，并给出小振动角频率。
\end{example}

\begin{figure}[H]\centering
\begin{tikzpicture}
\draw[curve](-1.8,0)--(3.3,0);\draw[fill=white,curve](-.8,.35)rectangle(.8,.95);\node at(0,.65){$M$};\draw(-.55,.175)circle(.175);\draw(.55,.175)circle(.175);\draw[curve](-.68,.95)--(-.68,3.9)--(0,3.9);\fill(0.0000000,3.9000000)circle(1.3pt)node[above right]{$O$};
\draw[line width=.8pt](0,3.9)--(1.469463131,1.877457514);\fill(1.4694631,1.8774575)circle(1.3pt)node[right]{$m$};
\draw[guide](0,3.9)--(0,1.2);\draw(0,3.28)arc(-90:-54:.62);\node at(.27,3.0){$\theta$};\node at(0.9847315653655915,2.9887287570313155){$l$};\draw[guide](0,1.4)arc(-90:-54:2.5);\draw[fill=white](0,1.4)circle(1.4pt);\node[right]at(.12,1.22){最低位置};\draw[axis](-1.8500000,1.2500000)--(-1.0500000,1.2500000)node[above]{$X$};
\node[right]at(2.3,3.1){刚性轻杆};
\end{tikzpicture}
\caption{摆球的最低位置高于车身，支架位于摆动平面的后方。}\label{sup1:fig:d1-e8-q}
\end{figure}
\par\medskip
\begin{solution}
摆球地面坐标为 $x=X+l\sin\theta$、$y=y_O-l\cos\theta$。由图\ref{sup1:fig:d1-e8-s}的速度分解，
\begin{align*}
T&=\frac12(M+m)\dot X^2+ml\cos\theta\dot X\dot\theta
+\frac12ml^2\dot\theta^2,\\
U&=mgl(1-\cos\theta).
\end{align*}

\begin{figure}[H]\centering
\begin{tikzpicture}
\draw[curve](-1.8,0)--(3.3,0);\draw[fill=white,curve](-.8,.35)rectangle(.8,.95);\node at(0,.65){$M$};\draw(-.55,.175)circle(.175);\draw(.55,.175)circle(.175);\draw[curve](-.68,.95)--(-.68,3.9)--(0,3.9);\fill(0.0000000,3.9000000)circle(1.3pt)node[above right]{$O$};
\draw[line width=.8pt](0,3.9)--(1.469463131,1.877457514);\fill(1.4694631,1.8774575)circle(1.3pt)node[right]{$m$};
\draw[guide](0,3.9)--(0,1.2);\draw(0,3.28)arc(-90:-54:.62);\node at(.27,3.0){$\theta$};\node at(0.9847315653655915,2.9887287570313155){$l$};\draw[guide](0,1.4)arc(-90:-54:2.5);\draw[fill=white](0,1.4)circle(1.4pt);\node[right]at(.12,1.22){最低位置};\draw[axis](-1.8500000,1.2500000)--(-1.0500000,1.2500000)node[above]{$X$};
\fill(5,2.3)circle(1.3pt)node[below]{$m$};\node at(5,3.9){摆球速度分解};\draw[axis](5.0000000,2.3000000)--(4.0000000,2.3000000)node[above left]{$\dot X$};
\draw[axis](5.0000000,2.3000000)--(5.8899187,2.9465638)node[above right]{$l\dot\theta$};

\end{tikzpicture}
\caption{右侧在同一摆球处画出小车平移速度与相对转动速度。}\label{sup1:fig:d1-e8-s}
\end{figure}

水平动量为零，故
\[
\dot X=-\frac{ml\cos\theta}{M+m}\dot\theta.
\]
代回动能后得到
\[
B_{\text{eff}}(\theta)=ml^2\left(1-\frac{m}{M+m}\cos^2\theta\right).
\]
最低点刚受内力冲量以后，小车也已有速度：
\[
\dot X_0=-\frac{ml}{M+m}\omega_0,
\qquad v_{m,x}(0)=\dot X_0+l\omega_0=\frac{Ml}{M+m}\omega_0.
\]
这两个速度使水平动量相消。将小车和摆球的动能相加，才得到下面的初始能量。

初态能量为 $E=\tfrac12ml^2M\omega_0^2/(M+m)$，因此
\begin{equation}\label{sup1:eq:d1-cartpend}
\boxed{\dot\theta^2=
\frac{\dfrac{M}{M+m}\omega_0^2-\dfrac{2g}{l}(1-\cos\theta)}
{1-\dfrac{m}{M+m}\cos^2\theta}}.
\end{equation}
转折角满足分子为零。到达转折角时 $\dot\theta=0$，由零动量关系可知小车也瞬时静止。

在 $\theta$ 很小时，$B_{\text{eff}}(0)=ml^2M/(M+m)$，$U\simeq\tfrac12mgl\theta^2$。保留一阶项，
\[
\ddot\theta+\frac{(M+m)g}{Ml}\theta=0,
\quad\boxed{\Omega=\sqrt{\frac{(M+m)g}{Ml}}}.
\]
减小 M 会增强小车反冲，并提高小振动频率。若需有限振幅的周期，则设振幅为 $\theta_m<\pi$，由能量积分写成
\[
\boxed{\mathcal T=4\int_0^{\theta_m}
\sqrt{\frac{B_{\text{eff}}(\theta)}{2[U(\theta_m)-U(\theta)]}}\,\dd\theta}.
\]
分母在转折点为零，但普通转折点的平方根奇性可积。振幅趋近 $\pi$ 时转折点接近势垒顶，周期发散。
令 $\rho=m/(M+m)$。在任意普通位置，完整的约化角运动方程为
\[
(1-\rho\cos^2\theta)\ddot\theta
+\rho\sin\theta\cos\theta\dot\theta^2+\frac gl\sin\theta=0.
\]
这里 $B_{\rm eff}'=ml^2[-2\rho\cos\theta(-\sin\theta)]=2ml^2\rho\sin\theta\cos\theta$，势能导数为 $U'=mgl\sin\theta$。代入 $B_{\rm eff}\ddot\theta+B_{\rm eff}'\dot\theta^2/2+U'=0$，便得到上式的速度平方项；它在小角度线性化时才成为高阶项而被舍去。

若初始能量满足 $0<E<2mgl$，摆角在两侧转向，最大正角由
\[
\cos\theta_m=1-\frac{Ml\omega_0^2}{2(M+m)g}
\]
确定。位置关系也能由零动量方程直接积分：
\[
X-X_0=-\frac{ml}{M+m}\int_0^\theta\cos u\dd u
=-\frac{ml}{M+m}[\sin u]_0^\theta
=-\frac{ml}{M+m}\sin\theta.
\]
图\ref{sup1:fig:d1-e8-s}中摆球相对小车向右偏移时，小车向左反冲；到 $\theta=\theta_m$ 时二者同时瞬时静止。若 $E\ge2mgl$，便不能使用普通有限振幅的转向周期积分，临界与越顶情形须按习题\ref{sup1:exercise:d1-p8}区分。

小振幅周期可由所给积分取极限。此时 $B_{\rm eff}\simeq ml^2M/(M+m)$、$U(\theta_m)-U(\theta)\simeq mgl(\theta_m^2-\theta^2)/2$，积分化为
\[
\begin{aligned}
\mathcal T&\longrightarrow4\sqrt{\frac{Ml}{(M+m)g}}
\int_0^{\theta_m}\frac{\dd\theta}{\sqrt{\theta_m^2-\theta^2}}\\
&=4\sqrt{\frac{Ml}{(M+m)g}}\left[\arcsin\frac\theta{\theta_m}\right]_0^{\theta_m}
=2\pi\sqrt{\frac{Ml}{(M+m)g}},
\end{aligned}
\]
与线性化频率相符。

\end{solution}



\begin{example}[大圆筒内的小圆筒]\label{ex:hw1-8}

薄壁大圆筒质量$M$、半径$R$，内有薄壁小圆筒质量$m$、半径$r<R$，两筒接触处不滑动。初始小筒在最低处，给予其质心沿水平方向的速度$v_0$；初始大筒静止，初速度满足接触处纯滚动条件。$v_0$足够小，小筒始终保持接触。(1)大筒固定；(2)大筒在地面上纯滚动。分别求小筒质心能上升的最大高度。

\end{example}


\begin{figure}[H]\centering
\begin{tikzpicture}[>=Stealth,line cap=round,line join=round]
\draw(0,0)circle(1.8);\coordinate(C)at(0,0);\coordinate(P)at(.8660254,-.5);\draw(P)circle(.8);\draw[guide](0,0)--(0,-1.8);\draw(0,0)--(P);\node[above]at(C){$C$};\fill(P)circle(1pt)node[above right]{$P$};\node at(-1.2,.5){$M,R$};\node at(.85,-.77){$m,r$};\draw[thick](-2.5,-1.8)--(2.5,-1.8);
\end{tikzpicture}
\caption{内部相切的两个薄壁圆筒}\label{fig:hw1-8-q}
\end{figure}

\begin{solution}
令$a=R-r$，$\theta$从向下铅直线量起，$X$为大筒圆心的水平位移。图\ref{fig:hw1-8-s}标出小筒质心高度增量$h=a(1-\cos\theta)$。

\begin{figure}[H]\centering
\begin{tikzpicture}[>=Stealth,line cap=round,line join=round]
\begin{scope}
\draw(0,0)circle(1.8);\coordinate(P)at(.8660254,-.5);\draw(P)circle(.8);
\draw[guide](0,0)--(0,-1.8);\draw(0,0)--(P);
\node[above]at(0,0){$C$};\fill(P)circle(1pt)node[above right]{$P$};
\node at(-1.2,.5){$M,R$};\node at(.85,-.77){$m,r$};
\draw[thick](-2.2,-1.8)--(2.4,-1.8);\node[above]at(.43,-.25){$a$};
\draw(0,-.45)arc(-90:-30:.45)node[midway,below]{$\theta$};
\draw[guide](.8660254,-.5)--(2.5,-.5);\draw[guide](0,-1)--(2.5,-1);
\draw[<->](2.3,-1)--(2.3,-.5)node[midway,right]{$h$};
\end{scope}
\begin{scope}[shift={(5.6,0)}]
\draw(0,0)circle(1.4);\draw(.69282,-.4)circle(.6);\draw[thick](-1.9,-1.4)--(2,-1.4);
\fill(0,0)circle(1pt)node[above left]{$C$};\fill(.69282,-.4)circle(1pt)node[left]{$P$};
\draw[axis](0,0)--(1.05,0)node[above]{$\dot X$};
\draw[axis](.69282,-.4)--(1.75,-.4)node[below]{$\dot X$};
\node at(0,1.8){转向瞬间：$\dot\theta=0$};
\node[align=center]at(.1,-2.05){$\Omega=-\dot X/R$\\[4pt]$\omega=-\dot X/r$};
\end{scope}
\end{tikzpicture}
\caption{上升高度与转向瞬间仍保留的平动、自转}\label{fig:hw1-8-s}
\end{figure}

大筒固定时，小筒质心速率为 $a|\dot\theta|$，自身角速度大小为 $a|\dot\theta|/r$。薄壁筒 $I=mr^2$，因此 $T=ma^2\dot\theta^2$。初态 $a\dot\theta=v_0$，动能为 $mv_0^2$；转向点两项动能都为零，能量关系为 $mv_0^2=mgh$，所以
\[
 \boxed{h_{\text{固定}}=\frac{v_0^2}{g}}.
\]
大筒滚动时，$\Omega=-\dot X/R$。由接触点速度相等，小筒的有向角速度为 $\omega=-(\dot X+a\dot\theta)/r$。相加两筒平动、转动动能，得到
\[
 T=(M+m)\dot X^2+ma(1+\cos\theta)\dot X\dot\theta
 +ma^2\dot\theta^2.
\]
$X$为循环坐标。初态$\dot X=0$、$a\dot\theta=v_0$，故
\[
 2(M+m)\dot X+ma(1+\cos\theta)\dot\theta=2mv_0.
\]
在相对运动转向点$\dot\theta=0$，却仍有$\dot X=mv_0/(M+m)$；不能把全部动能置零。代入能量方程：
\[
 mv_0^2=\frac{m^2v_0^2}{M+m}+mgh,
 \qquad\boxed{h_{\text{滚动}}=\frac{M}{M+m}\frac{v_0^2}{g}}.
\]
在转向构形中，两个圆心的速度都为$\dot X\vct e_x$，但两筒仍有$\Omega=-\dot X/R$、$\omega=-\dot X/r$。图\ref{fig:hw1-8-s}右侧画出这一状态。此时的$2(M+m)\dot X$是循环坐标的广义动量；两筒的实际水平动量只有$(M+m)\dot X$，不能把地面静摩擦的水平作用删去。

转向时，两筒的共同平移和约束所要求的自转仍有动能。$M\to\infty$时，大筒的运动趋于消失，第二式回到固定大筒的结果。


\end{solution}



\section{滑动摩擦与滚动状态的转换}

滑动摩擦的耗散取决于接触点的相对速度。滑动与纯滚动采用不同的方程，两阶段由速度连续及接触条件连接。



\subsection{滑动摩擦与滚动过渡}
\subsubsection{摩擦功率取决于相对接触速度}
设刚体与接触面的接触点速度分别为 $\vct v_P$、$\vct v_S$，接触面施于刚体的摩擦力为 $\vct f$。这一对摩擦力对两个物体的总功率为
\[
\mathcal P_f=\vct f\cdot(\vct v_P-\vct v_S).
\]
纯滚动时接触点速度相同，总功率为零；若接触面也在运动，单个刚体所受静摩擦的功率则可以不为零。出现相对滑动时，库仑动摩擦模型给出
\begin{equation}\label{sup1:eq:d1-dissipation}
\mathcal P_f=-\mu_kN|u|,\qquad
Q_{\text{热}}=\int\mu_kN|u|\,\dd t,
\end{equation}
其中 $u$ 为切向相对滑动速度。热量对应相对滑移路程，不对应质心的总路程。

\subsubsection{接触状态的分段处理}
先根据相对滑动速度决定摩擦方向，分别写平动和转动方程。若随后 $u$ 到达零，再检验维持 $u=0$ 所需的静摩擦是否可实现。若可以，转入纯滚动方程；若不能，须根据新的相对滑动趋势继续列式。动摩擦力不能在 $u=0$ 后仍机械地保持原方向。


\subsection{例题与模型变式}

\begin{example}[圆柱边滚边滑]\label{ex:hw1-7}

质量$m$、半径$r$的均匀实心圆柱在倾角$\varphi$的固定斜面上运动，动摩擦因数为$\mu$。求发生滑动时圆柱的质心加速度和角加速度，并判断滑动何时可能结束。

\end{example}


\begin{figure}[H]\centering
\begin{tikzpicture}[>=Stealth,line cap=round,line join=round]
\begin{scope}[rotate=-20]\draw[thick](-1,0)--(3.5,0);\draw(1,.65)circle(.65);\fill(1,.65)circle(1pt)node[above]{$C$};\node at(2.8,1){$m,r$};\end{scope}
\draw[axis](1.16201,.26887)--(2.1017,-.07315)node[above]{$v$};
\draw[axis](1.38701,.65858)arc(60:-180:.45);\node[left]at(.55,.68){$\omega$};
\draw[guide](3.28892,-1.19707)--(1.4,-1.19707);\draw(1.98892,-1.19707)arc(180:160:1.3);\node at(1.68,-.94){$\varphi$};
\end{tikzpicture}
\caption{在固定斜面上滑动的实心圆柱}\label{fig:hw1-7-q}
\end{figure}

\begin{solution}
取下坡为$x$正向，正转向使$r\omega$也沿下坡。图\ref{fig:hw1-7-s}中接触点相对斜面的切向速度为$w=v-r\omega$，摩擦方向由$w$决定。

\begin{figure}[H]\centering
\begin{tikzpicture}[>=Stealth,line cap=round,line join=round]
\begin{scope}[rotate=-20]\draw[thick](-1,0)--(3.5,0);\draw(1,.65)circle(.65);\fill(1,.65)circle(1pt)node[left]{$C$};\node at(2.8,1){$m,r$};\end{scope}\begin{scope}[rotate=-20]\draw[axis](1,0)--(.05,0)node[above left]{$f$};\draw[axis](1,0)--(1,1.65)node[above]{$N$};\end{scope}\node at(1.5,-1.65){$w=v-r\omega$};\draw[axis](1.16201,.26887)--(1.16201,-.85)node[left]{$mg$};\node at(3,-1.3){$w>0$};
\end{tikzpicture}
\caption{接触点向下滑动时的摩擦方向}\label{fig:hw1-7-s}
\end{figure}

令$\sigma=\operatorname{sgn}w$。滑动时$N=mg\cos\varphi$，摩擦的下坡分量为$-\sigma\mu mg\cos\varphi$，故
\[
 \boxed{a=g\sin\varphi-\sigma\mu g\cos\varphi,\qquad
 \beta=\frac{2\sigma\mu g\cos\varphi}{r}}.
\]
于是
\[
\dot w=\dot v-r\dot\omega
=g\sin\varphi-\sigma\mu g\cos\varphi-2\sigma\mu g\cos\varphi
=g\sin\varphi-3\sigma\mu g\cos\varphi.
\]
在同一滑动支上 $\sigma$ 固定，故 $\dot w$ 为常数，$w=w_0+\int_0^t\dot w\dd s=w_0+\dot w\,t$。只有 $t_*=-w_0/\dot w>0$ 时才会到达 $w=0$。随后若能维持纯滚动，所需静摩擦为 $mg\sin\varphi/3$，还须满足 $\mu_s\ge\tan\varphi/3$；否则应按新的滑动方向继续计算。

例如$w_0>0$时，只有$3\mu\cos\varphi>\sin\varphi$才能消除这一方向的滑动，此时
\[
 t_* =\frac{w_0}{g(3\mu\cos\varphi-\sin\varphi)}.
\]
$w_0<0$时，$\dot w=g(\sin\varphi+3\mu\cos\varphi)>0$，可在有限时间到达零点。到达零点后，用纯滚动条件重新求静摩擦，并检验它是否超过静摩擦极限。图中的$f$画的是$w>0$这一支的实际方向。

取机械能 $E=mv^2/2+I\omega^2/2-mgx\sin\varphi$。逐项求导，并用平动、转动方程代换，
\[
\begin{aligned}
\dot E&=mv\dot v+I\omega\dot\omega-mgv\sin\varphi\\
&=-\sigma\mu mg\cos\varphi\,v
+\sigma\mu mgr\cos\varphi\,\omega\\
&=-\mu mg\cos\varphi|v-r\omega|.
\end{aligned}
\]
耗散率涉及接触点相对速度；若仅用质心位移计算摩擦损失，会漏掉摩擦对转动动能的贡献。


\end{solution}



\begin{example}[圆柱从滑动转为纯滚动]\label{sup1:problem-example-9}\label{sup1:example:d1-e9}
质量为 $m$、半径为 $r$、转动惯量 $I=\kappa mr^2$ 的圆柱形刚体，在倾角为 $\alpha$ 的固定斜面上向下运动。初始质心速率为 $v_0>0$，自转角速度为零。动摩擦系数为 $\mu_k$，静摩擦系数为 $\mu_s$，并有
\[
\mu_s\ge\mu_k>\frac{\kappa}{1+\kappa}\tan\alpha.
\]
求开始纯滚动的时刻、该时刻的速率和此前产生的热量。
\end{example}

\begin{figure}[H]\centering
\begin{tikzpicture}
\begin{scope}[rotate=-24]\draw[curve](-1.6,0)--(3.4,0);\draw[fill=white,curve](0,.7)circle(.7);\fill(0.00000,0.70000)circle(1.3pt)node[above left]{$C$};
\fill(0.00000,0.00000)circle(1.3pt)node[below left]{$P$};
\draw(0,.7)--(0.520201378,1.168391424)node[midway,above]{$r$};\draw[axis](1.20000,0.70000)--(2.10000,0.70000)node[above]{$v_0$};
\node at(1.8,1.4){$I=\kappa mr^2$};\draw[axis](0.0000000,0.0000000)--(-1.0500000,0.0000000)node[above]{$f_k$};
\end{scope}\draw[guide](1.3060545559848429,-1.3829045864577207)--(3.306054555984843,-1.3829045864577207);\draw(2.456054555984843,-1.3829045864577207)arc(180:156:.65);\node at(2.126054555984843,-1.1529045864577208){$\alpha$};\node at(-1.1,2.1){$v_0>0,\quad\omega_0=0$};
\end{tikzpicture}
\caption{初始相对滑动沿下坡方向，动摩擦力 $f_k$ 沿上坡方向。}\label{sup1:fig:d1-e9-q}
\end{figure}
\par\medskip
\begin{solution}
取向下平动和顺下坡滚动的自转为正，接触滑动速度为 $u=v-r\omega$。初始 $u=v_0>0$，所以动摩擦沿上坡。平动和转动方程为
\[
\dot v=g(\sin\alpha-\mu_k\cos\alpha)=a_k,
\qquad
\dot\omega=\frac{\mu_kg\cos\alpha}{\kappa r}.
\]
因此
\[
\dot u=-K,\qquad
K=g\left[\left(1+\frac1\kappa\right)\mu_k\cos\alpha-\sin\alpha\right]>0.
\]
图\ref{sup1:fig:d1-e9-s}右图中的 $u(t)$ 为一直线，首次到零的时刻为
\[
\boxed{t_*=\frac{v_0}{K}}.
\]

\begin{figure}[H]\centering
\begin{tikzpicture}
\begin{scope}[rotate=-24]\draw[curve](-1.6,0)--(3.4,0);\draw[fill=white,curve](0,.7)circle(.7);\fill(0.00000,0.70000)circle(1.3pt)node[above left]{$C$};
\fill(0.00000,0.00000)circle(1.3pt)node[below left]{$P$};
\draw(0,.7)--(0.520201378,1.168391424)node[midway,above]{$r$};\draw[axis](1.20000,0.70000)--(2.10000,0.70000)node[above]{$v_0$};
\node at(1.8,1.4){$I=\kappa mr^2$};\draw[axis](0.0000000,0.0000000)--(-1.0500000,0.0000000)node[above]{$f_k$};
\end{scope}\draw[guide](1.3060545559848429,-1.3829045864577207)--(3.306054555984843,-1.3829045864577207);\draw(2.456054555984843,-1.3829045864577207)arc(180:156:.65);\node at(2.126054555984843,-1.1529045864577208){$\alpha$};\node at(-1.1,2.1){$v_0>0,\quad\omega_0=0$};\begin{scope}[shift={(5.4,0.2)}]
\draw[axis](-0.15,0)--(3.4,0)node[right]{$t$};\draw[axis](0,-0.2)--(0,1.8)node[above]{$u$};\draw[curve](0,1.4)--(2.8,0);\draw[guide](0,1.4)--(-.15,1.4)node[left]{$v_0$};\node[below]at(2.8,0){$t_*$};\draw[guide](2.8,0)--(3.25,0);\node at(1.5,-.6){$Q_{\text{热}}\propto$ 面积};
\end{scope}

\end{tikzpicture}
\caption{滑动阶段的相对速度线性减小；曲线下的面积为相对滑移路程。}\label{sup1:fig:d1-e9-s}
\end{figure}

在 $0\le t\le t_*$ 内，速度的完整表达式为
\[
\begin{aligned}
v&=v_0+\int_0^t a_k\dd s=v_0+a_kt,\\
r\omega&=\frac{\mu_kg\cos\alpha}{\kappa}\int_0^t\dd s
=\frac{\mu_kg\cos\alpha}{\kappa}t,\\
u&=v_0-Kt\ge0.
\end{aligned}
\]
这验证了整个滑动阶段的摩擦始终沿上坡方向。即使 $\mu_k>\tan\alpha$ 使 $a_k<0$，下式给出的 $v_*$ 仍为正，质心不会在达到纯滚动以前先倒向上坡。

在 $t=t_*$，接触滑动速度首次为零，即 $v_*=r\omega_*$。代入任一速度式得
\[
\boxed{v_*=v_0+a_kt_*
=\frac{\mu_kg\cos\alpha}{\kappa K}v_0}.
\]
维持纯滚动所需的摩擦系数正是例\ref{sup1:example:d1-e4}所得下限，题设保证静摩擦足够，因此 $t>t_*$ 后的加速度为 $g\sin\alpha/(1+\kappa)$。

由式\eqref{sup1:eq:d1-dissipation}，热量等于相对滑移路程乘以摩擦力。图\ref{sup1:fig:d1-e9-s}右图的三角形面积也可直接积分为
\[
\int_0^{t_*}(v_0-Kt)\dd t
=\left[v_0t-\frac K2t^2\right]_0^{t_*}
=v_0t_*-\frac12Kt_*^2=\frac12v_0t_*.
\]
其中用到了 $Kt_*=v_0$，故
\[
\boxed{Q_{\text{热}}=\frac12\mu_kmg\cos\alpha\,v_0t_*
=\frac{\mu_kmg\cos\alpha}{2K}v_0^2}.
\]
质心在这段时间沿斜面走过 $s_*=v_0t_*+\tfrac12a_kt_*^2$，但相对滑移路程只有 $v_0t_*/2$。完整能量关系为
\[
\frac12mv_0^2+mgs_*\sin\alpha
=\frac12m(1+\kappa)v_*^2+Q_{\text{热}}.
\]
若 $\mu_k\le\kappa\tan\alpha/(1+\kappa)$，则 $K\le0$，在本题初态下 $u$ 不会降到零，上述有限过渡时刻也就不存在。
令 $a_r=g\sin\alpha/(1+\kappa)$。过渡后的运动应另写一段：
\[
\begin{aligned}
v(t)&=v_*+a_r(t-t_*),\qquad \omega(t)=v(t)/r,\\
s(t)&=s_*+v_*(t-t_*)+\frac12a_r(t-t_*)^2,
\qquad t\ge t_*.
\end{aligned}
\]
此时静摩擦取所需值 $\kappa mg\sin\alpha/(1+\kappa)$，不再固定为 $\mu_kmg\cos\alpha$；接触不继续滑动，故此前产生的热量保持不变。

图\ref{sup1:fig:d1-e9-s}的面积对应 $\int_0^{t_*}u\,\dd t$。若改用质心路程来算热量，会漏掉摩擦同时增加自转动能的部分。这一分段解成立需要两项条件：滑动阶段的 K 为正，保证 u 能降到零；达到零后静摩擦足够，保证新约束能持续。题设已分别给出了这两项保证。

\end{solution}



\section{习题与解答}

\begin{exercise}[螺旋轨道上的下滑]\label{sup1:problem-exercise-1}\label{sup1:exercise:d1-p1}
光滑螺旋细管固定在空间，中心线为 $x=R\cos\varphi$、$y=R\sin\varphi$、$z=z_0-b\varphi$，其中 $R,b>0$。质点 $m$ 从 $\varphi=0$ 处静止释放。求 $\ddot\varphi$、运动一整圈所需时间，以及该时刻的实际速率。
\end{exercise}

\begin{figure}[H]\centering
\begin{tikzpicture}
\draw[curve]plot coordinates{(1.200000,3.700000)(1.198355,3.712031)(1.193426,3.723990)(1.185226,3.735806)(1.173777,3.747407)(1.159111,3.758724)(1.141268,3.769685)(1.120297,3.780224)(1.096255,3.790271)(1.069208,3.799761)(1.039230,3.808628)(1.006405,3.816811)(0.970820,3.824247)(0.932575,3.830877)(0.891774,3.836645)(0.848528,3.841496)(0.802957,3.845378)(0.755184,3.848241)(0.705342,3.850039)(0.653567,3.850729)(0.600000,3.850269)(0.544789,3.848623)(0.488084,3.845755)(0.430042,3.841635)(0.370820,3.836236)(0.310583,3.829534)(0.249494,3.821507)(0.187721,3.812141)(0.125434,3.801420)(0.062803,3.789337)(0.000000,3.775885)(-0.062803,3.761063)(-0.125434,3.744872)(-0.187721,3.727318)(-0.249494,3.708410)(-0.310583,3.688162)(-0.370820,3.666590)(-0.430042,3.643715)(-0.488084,3.619560)(-0.544789,3.594154)(-0.600000,3.567526)(-0.653567,3.539711)(-0.705342,3.510747)(-0.755184,3.480675)(-0.802957,3.449537)(-0.848528,3.417381)(-0.891774,3.384256)(-0.932575,3.350213)(-0.970820,3.315309)(-1.006405,3.279598)(-1.039230,3.243142)(-1.069208,3.206000)(-1.096255,3.168236)(-1.120297,3.129914)(-1.141268,3.091101)(-1.159111,3.051865)(-1.173777,3.012274)(-1.185226,2.972399)(-1.193426,2.932309)(-1.198355,2.892075)(-1.200000,2.851770)(-1.198355,2.811465)(-1.193426,2.771231)(-1.185226,2.731141)(-1.173777,2.691265)(-1.159111,2.651675)(-1.141268,2.612438)(-1.120297,2.573626)(-1.096255,2.535304)(-1.069208,2.497540)(-1.039230,2.460398)(-1.006405,2.423942)(-0.970820,2.388231)(-0.932575,2.353327)(-0.891774,2.319284)(-0.848528,2.286159)(-0.802957,2.254003)(-0.755184,2.222865)(-0.705342,2.192792)(-0.653567,2.163829)(-0.600000,2.136014)(-0.544789,2.109386)(-0.488084,2.083980)(-0.430042,2.059825)(-0.370820,2.036950)(-0.310583,2.015378)(-0.249494,1.995130)(-0.187721,1.976222)(-0.125434,1.958668)(-0.062803,1.942477)(-0.000000,1.927655)(0.062803,1.914203)(0.125434,1.902120)(0.187721,1.891399)(0.249494,1.882033)(0.310583,1.874006)(0.370820,1.867304)(0.430042,1.861905)(0.488084,1.857785)(0.544789,1.854917)(0.600000,1.853271)(0.653567,1.852811)(0.705342,1.853500)(0.755184,1.855299)(0.802957,1.858162)(0.848528,1.862044)(0.891774,1.866895)(0.932575,1.872663)(0.970820,1.879293)(1.006405,1.886729)(1.039230,1.894912)(1.069208,1.903779)(1.096255,1.913269)(1.120297,1.923316)(1.141268,1.933854)(1.159111,1.944816)(1.173777,1.956133)(1.185226,1.967734)(1.193426,1.979550)(1.198355,1.991509)(1.200000,2.003540)(1.198355,2.015571)(1.193426,2.027530)(1.185226,2.039346)(1.173777,2.050947)(1.159111,2.062264)(1.141268,2.073225)(1.120297,2.083764)(1.096255,2.093811)(1.069208,2.103301)(1.039230,2.112168)(1.006405,2.120351)(0.970820,2.127787)(0.932575,2.134417)(0.891774,2.140185)(0.848528,2.145036)(0.802957,2.148918)(0.755184,2.151781)(0.705342,2.153579)(0.653567,2.154269)(0.600000,2.153809)(0.544789,2.152163)(0.488084,2.149295)(0.430042,2.145175)(0.370820,2.139776)(0.310583,2.133074)(0.249494,2.125047)(0.187721,2.115681)(0.125434,2.104960)(0.062803,2.092877)(0.000000,2.079425)(-0.062803,2.064603)(-0.125434,2.048412)(-0.187721,2.030858)(-0.249494,2.011950)(-0.310583,1.991702)(-0.370820,1.970130)(-0.430042,1.947255)(-0.488084,1.923100)(-0.544789,1.897694)(-0.600000,1.871066)(-0.653567,1.843251)(-0.705342,1.814287)(-0.755184,1.784215)(-0.802957,1.753077)(-0.848528,1.720921)(-0.891774,1.687796)(-0.932575,1.653753)(-0.970820,1.618849)(-1.006405,1.583138)(-1.039230,1.546682)(-1.069208,1.509540)(-1.096255,1.471776)(-1.120297,1.433454)(-1.141268,1.394641)(-1.159111,1.355405)(-1.173777,1.315814)(-1.185226,1.275939)(-1.193426,1.235849)(-1.198355,1.195615)(-1.200000,1.155310)(-1.198355,1.115005)(-1.193426,1.074771)(-1.185226,1.034681)(-1.173777,0.994805)(-1.159111,0.955215)(-1.141268,0.915978)(-1.120297,0.877166)(-1.096255,0.838844)(-1.069208,0.801080)(-1.039230,0.763938)(-1.006405,0.727482)(-0.970820,0.691771)(-0.932575,0.656867)(-0.891774,0.622824)(-0.848528,0.589699)(-0.802957,0.557543)(-0.755184,0.526405)(-0.705342,0.496332)(-0.653567,0.467368)(-0.600000,0.439554)(-0.544789,0.412926)(-0.488084,0.387520)(-0.430042,0.363365)(-0.370820,0.340490)(-0.310583,0.318918)(-0.249494,0.298670)(-0.187721,0.279762)(-0.125434,0.262208)(-0.062803,0.246017)(-0.000000,0.231195)(0.062803,0.217743)(0.125434,0.205660)(0.187721,0.194939)(0.249494,0.185572)(0.310583,0.177546)(0.370820,0.170844)(0.430042,0.165445)(0.488084,0.161325)(0.544789,0.158457)(0.600000,0.156811)(0.653567,0.156351)(0.705342,0.157040)(0.755184,0.158839)(0.802957,0.161702)(0.848528,0.165584)(0.891774,0.170435)(0.932575,0.176203)(0.970820,0.182833)(1.006405,0.190269)(1.039230,0.198452)(1.069208,0.207319)(1.096255,0.216809)(1.120297,0.226856)(1.141268,0.237394)(1.159111,0.248356)(1.173777,0.259673)(1.185226,0.271274)(1.193426,0.283090)(1.198355,0.295049)(1.200000,0.307080)};
\draw[axis](2.1,.15)--(2.1,4.25)node[above]{$z$};\draw[guide](0,.1)--(0,4.15);\draw(0,3.7)--(1.2,3.7)node[midway,below]{$R$};\fill(-1.14127,2.61244)circle(1.3pt)node[above left]{$P$};
\node[right]at(2.75,2.9){$z=z_0-b\varphi$};\node[right]at(2.75,1.9){$b>0$};
\end{tikzpicture}
\caption{螺旋轴竖直；图作斜投影，$\varphi$ 增大时高度降低。}\label{sup1:fig:d1-p1-q}
\end{figure}
\par\medskip
\begin{solution}
图\ref{sup1:fig:d1-p1-s}中的竖直下降与沿螺旋的运动由同一个 $\varphi$ 决定，坐标求导给出 $\vct v=\dot\varphi(-R\sin\varphi,R\cos\varphi,-b)$。平方相加，
\[
v^2=R^2\dot\varphi^2+b^2\dot\varphi^2,
\quad T=\frac12m(R^2+b^2)\dot\varphi^2,
\quad U=-mgb\varphi.
\]

\begin{figure}[H]\centering
\begin{tikzpicture}
\draw[curve]plot coordinates{(1.200000,3.700000)(1.198355,3.712031)(1.193426,3.723990)(1.185226,3.735806)(1.173777,3.747407)(1.159111,3.758724)(1.141268,3.769685)(1.120297,3.780224)(1.096255,3.790271)(1.069208,3.799761)(1.039230,3.808628)(1.006405,3.816811)(0.970820,3.824247)(0.932575,3.830877)(0.891774,3.836645)(0.848528,3.841496)(0.802957,3.845378)(0.755184,3.848241)(0.705342,3.850039)(0.653567,3.850729)(0.600000,3.850269)(0.544789,3.848623)(0.488084,3.845755)(0.430042,3.841635)(0.370820,3.836236)(0.310583,3.829534)(0.249494,3.821507)(0.187721,3.812141)(0.125434,3.801420)(0.062803,3.789337)(0.000000,3.775885)(-0.062803,3.761063)(-0.125434,3.744872)(-0.187721,3.727318)(-0.249494,3.708410)(-0.310583,3.688162)(-0.370820,3.666590)(-0.430042,3.643715)(-0.488084,3.619560)(-0.544789,3.594154)(-0.600000,3.567526)(-0.653567,3.539711)(-0.705342,3.510747)(-0.755184,3.480675)(-0.802957,3.449537)(-0.848528,3.417381)(-0.891774,3.384256)(-0.932575,3.350213)(-0.970820,3.315309)(-1.006405,3.279598)(-1.039230,3.243142)(-1.069208,3.206000)(-1.096255,3.168236)(-1.120297,3.129914)(-1.141268,3.091101)(-1.159111,3.051865)(-1.173777,3.012274)(-1.185226,2.972399)(-1.193426,2.932309)(-1.198355,2.892075)(-1.200000,2.851770)(-1.198355,2.811465)(-1.193426,2.771231)(-1.185226,2.731141)(-1.173777,2.691265)(-1.159111,2.651675)(-1.141268,2.612438)(-1.120297,2.573626)(-1.096255,2.535304)(-1.069208,2.497540)(-1.039230,2.460398)(-1.006405,2.423942)(-0.970820,2.388231)(-0.932575,2.353327)(-0.891774,2.319284)(-0.848528,2.286159)(-0.802957,2.254003)(-0.755184,2.222865)(-0.705342,2.192792)(-0.653567,2.163829)(-0.600000,2.136014)(-0.544789,2.109386)(-0.488084,2.083980)(-0.430042,2.059825)(-0.370820,2.036950)(-0.310583,2.015378)(-0.249494,1.995130)(-0.187721,1.976222)(-0.125434,1.958668)(-0.062803,1.942477)(-0.000000,1.927655)(0.062803,1.914203)(0.125434,1.902120)(0.187721,1.891399)(0.249494,1.882033)(0.310583,1.874006)(0.370820,1.867304)(0.430042,1.861905)(0.488084,1.857785)(0.544789,1.854917)(0.600000,1.853271)(0.653567,1.852811)(0.705342,1.853500)(0.755184,1.855299)(0.802957,1.858162)(0.848528,1.862044)(0.891774,1.866895)(0.932575,1.872663)(0.970820,1.879293)(1.006405,1.886729)(1.039230,1.894912)(1.069208,1.903779)(1.096255,1.913269)(1.120297,1.923316)(1.141268,1.933854)(1.159111,1.944816)(1.173777,1.956133)(1.185226,1.967734)(1.193426,1.979550)(1.198355,1.991509)(1.200000,2.003540)(1.198355,2.015571)(1.193426,2.027530)(1.185226,2.039346)(1.173777,2.050947)(1.159111,2.062264)(1.141268,2.073225)(1.120297,2.083764)(1.096255,2.093811)(1.069208,2.103301)(1.039230,2.112168)(1.006405,2.120351)(0.970820,2.127787)(0.932575,2.134417)(0.891774,2.140185)(0.848528,2.145036)(0.802957,2.148918)(0.755184,2.151781)(0.705342,2.153579)(0.653567,2.154269)(0.600000,2.153809)(0.544789,2.152163)(0.488084,2.149295)(0.430042,2.145175)(0.370820,2.139776)(0.310583,2.133074)(0.249494,2.125047)(0.187721,2.115681)(0.125434,2.104960)(0.062803,2.092877)(0.000000,2.079425)(-0.062803,2.064603)(-0.125434,2.048412)(-0.187721,2.030858)(-0.249494,2.011950)(-0.310583,1.991702)(-0.370820,1.970130)(-0.430042,1.947255)(-0.488084,1.923100)(-0.544789,1.897694)(-0.600000,1.871066)(-0.653567,1.843251)(-0.705342,1.814287)(-0.755184,1.784215)(-0.802957,1.753077)(-0.848528,1.720921)(-0.891774,1.687796)(-0.932575,1.653753)(-0.970820,1.618849)(-1.006405,1.583138)(-1.039230,1.546682)(-1.069208,1.509540)(-1.096255,1.471776)(-1.120297,1.433454)(-1.141268,1.394641)(-1.159111,1.355405)(-1.173777,1.315814)(-1.185226,1.275939)(-1.193426,1.235849)(-1.198355,1.195615)(-1.200000,1.155310)(-1.198355,1.115005)(-1.193426,1.074771)(-1.185226,1.034681)(-1.173777,0.994805)(-1.159111,0.955215)(-1.141268,0.915978)(-1.120297,0.877166)(-1.096255,0.838844)(-1.069208,0.801080)(-1.039230,0.763938)(-1.006405,0.727482)(-0.970820,0.691771)(-0.932575,0.656867)(-0.891774,0.622824)(-0.848528,0.589699)(-0.802957,0.557543)(-0.755184,0.526405)(-0.705342,0.496332)(-0.653567,0.467368)(-0.600000,0.439554)(-0.544789,0.412926)(-0.488084,0.387520)(-0.430042,0.363365)(-0.370820,0.340490)(-0.310583,0.318918)(-0.249494,0.298670)(-0.187721,0.279762)(-0.125434,0.262208)(-0.062803,0.246017)(-0.000000,0.231195)(0.062803,0.217743)(0.125434,0.205660)(0.187721,0.194939)(0.249494,0.185572)(0.310583,0.177546)(0.370820,0.170844)(0.430042,0.165445)(0.488084,0.161325)(0.544789,0.158457)(0.600000,0.156811)(0.653567,0.156351)(0.705342,0.157040)(0.755184,0.158839)(0.802957,0.161702)(0.848528,0.165584)(0.891774,0.170435)(0.932575,0.176203)(0.970820,0.182833)(1.006405,0.190269)(1.039230,0.198452)(1.069208,0.207319)(1.096255,0.216809)(1.120297,0.226856)(1.141268,0.237394)(1.159111,0.248356)(1.173777,0.259673)(1.185226,0.271274)(1.193426,0.283090)(1.198355,0.295049)(1.200000,0.307080)};
\draw[axis](2.1,.15)--(2.1,4.25)node[above]{$z$};\draw[guide](0,.1)--(0,4.15);\draw(0,3.7)--(1.2,3.7)node[midway,below]{$R$};\fill(-1.14127,2.61244)circle(1.3pt)node[above left]{$P$};
\draw[axis](-1.14127,2.61244)--(-1.14127,1.61244)node[below]{$mg$};
\node[right]at(2.75,3.25){$v^2=(R^2+b^2)\dot\varphi^2$};\node[right]at(2.75,2.2){$U=-mgb\varphi$};
\end{tikzpicture}
\caption{重力势能只由螺旋的竖直下降量决定；动能含水平绕转与竖直运动两项。}\label{sup1:fig:d1-p1-s}
\end{figure}

广义动量为 $m(R^2+b^2)\dot\varphi$，其时间导数等于广义重力 $mgb$。令 $A_\varphi=gb/(R^2+b^2)$，从零初角速度积分，$\dot\varphi=\int_0^t A_\varphi\dd s=A_\varphi t$；再从零初角度积分，$\varphi=A_\varphi\int_0^t s\dd s=A_\varphi[s^2/2]_0^t$。因此
\[
\boxed{\ddot\varphi=\frac{gb}{R^2+b^2}},\qquad
\varphi=\frac{gb}{2(R^2+b^2)}t^2.
\]
一整圈为 $\varphi=2\pi$，所以
\[
\boxed{t_1=\sqrt{\frac{4\pi(R^2+b^2)}{gb}}},\qquad
\boxed{v_1=\sqrt{4\pi gb}}.
\]
后一结果也可由下降高度 $2\pi b$ 直接得到。改变 $R$ 会改变运动时间，但不改变同一竖直落差所对应的实际速率。
图\ref{sup1:fig:d1-p1-s}中，参数角增加 $\dd\varphi$ 时，水平绕转路程为 $R\dd\varphi$，竖直下降量为 $b\dd\varphi$。实际弧长微元因而为
\[
\dd s_h=\sqrt{R^2+b^2}\,\dd\varphi,
\qquad s_h=\sqrt{R^2+b^2}\,\varphi.
\]
等效直线斜坡的下降高度与弧长之比为 $b/\sqrt{R^2+b^2}$，所以切向加速度为 $a_t=gb/\sqrt{R^2+b^2}$。再由 $a_t=\sqrt{R^2+b^2}\ddot\varphi$，恢复前面的角加速度。

到一整圈时，弧长为 $2\pi\sqrt{R^2+b^2}$。由 $s_h=a_tt_1^2/2$ 或 $v_1=a_tt_1$ 都能得到相同结果。轨道反力还须提供水平绕转所需的向心加速度，因此实际加速度还含法向分量，但轨道法向力对速度正交，不进入能量损失。

\end{solution}


\begin{exercise}[铰接细杆的角速度与铰链力]\label{sup1:problem-exercise-2}\label{sup1:exercise:d1-p2}
长为 $l$、质量为 $m$ 的均匀细杆，一端铰接在固定点 O，可在竖直平面内转动。细杆从水平位置静止释放。取图\ref{sup1:fig:d1-p2-q}的 $\theta$ 为杆与竖直向下方向夹角。求下降过程中 $\dot\theta^2$、$\ddot\theta$，以及杆到达最低处时铰链的作用力。
\end{exercise}

\begin{figure}[H]\centering
\begin{tikzpicture}
\draw[curve](-.7,.15)--(.7,.15);\draw(0,.15)--(0,0);\draw[guide](0,0)--(0,-3.2);\draw[line width=2pt](0,0)--(1.8038095242655348,-2.1415581243965676);\fill(0.00000,0.00000)circle(1.3pt)node[above left]{$O$};
\fill(0.90190,-1.07078)circle(1.3pt)node[left]{$G$};
\draw(0,-.65)arc(-90:-49.893:.65);\node at(.33,-.84){$\theta$};\node[right]at(2.5,0){均匀细杆 $m,l$};\node[right]at(2.5,-.75){从水平静止释放};
\end{tikzpicture}
\caption{杆下降时 $\theta$ 从 $\pi/2$ 减小到零。}\label{sup1:fig:d1-p2-q}
\end{figure}
\par\medskip
\begin{solution}
绕 O 的转动惯量为 $I_O=ml^2/3$，质心 G 距 O 为 $l/2$。以初始水平位置为势能零点，
\[
\frac12I_O\dot\theta^2-\frac12mgl\cos\theta=0.
\]
因此
\[
\boxed{\dot\theta^2=\frac{3g}{l}\cos\theta},\qquad
\boxed{\ddot\theta=-\frac{3g}{2l}\sin\theta}.
\]

\begin{figure}[H]\centering
\begin{tikzpicture}
\draw[curve](-.7,.15)--(.7,.15);\draw(0,.15)--(0,0);\draw[guide](0,0)--(0,-3.2);\draw[line width=2pt](0,0)--(1.8038095242655348,-2.1415581243965676);\fill(0.00000,0.00000)circle(1.3pt)node[above left]{$O$};
\fill(0.90190,-1.07078)circle(1.3pt)node[left]{$G$};
\draw(0,-.65)arc(-90:-49.893:.65);\node at(.33,-.84){$\theta$};\draw[axis](0.90190,-1.07078)--(0.90190,-2.02078)node[below]{$mg$};
\draw[axis](0.00000,0.00000)--(1.00000,0.00000)node[above]{$H_x$};
\draw[axis](0.00000,0.00000)--(0.00000,1.00000)node[right]{$H_y$};

\end{tikzpicture}
\caption{恢复铰链力时对整杆使用质心加速度；图示一般位置的作用点和分力正方向。}\label{sup1:fig:d1-p2-s}
\end{figure}

下降时 $\theta$ 由 $\pi/2$ 减小到 0，故实际取 $\dot\theta=-\sqrt{3g\cos\theta/l}$。角加速度可直接由重力矩得到
\[
I_O\ddot\theta=-\frac{mgl}{2}\sin\theta,
\]
这一列式也适用于释放瞬间，避免对能量求导后在 $\dot\theta=0$ 处约去速度。

取 O 为原点、水平向右和竖直向上为正，质心坐标为
\[
x_G=\frac l2\sin\theta,\qquad y_G=-\frac l2\cos\theta.
\]
先按乘积法则保留两类项，
\[
a_{Gx}=\frac l2(\cos\theta\ddot\theta-\sin\theta\dot\theta^2),
\qquad
a_{Gy}=\frac l2(\sin\theta\ddot\theta+\cos\theta\dot\theta^2).
\]
再代入角速度、角加速度，得到
\[
a_{Gx}=-\frac{9g}{4}\sin\theta\cos\theta,
\qquad a_{Gy}=\frac{3g}{4}(3\cos^2\theta-1).
\]
因此一般位置的铰链力为
\[
H_x=-\frac94mg\sin\theta\cos\theta,
\qquad H_y=\frac{mg}{4}(1+9\cos^2\theta).
\]
图\ref{sup1:fig:d1-p2-s}中的两个箭头只定义正分量方向；例如 $H_x<0$ 表示实际向左。

图\ref{sup1:fig:d1-p2-s}中铰链力作用于 O，重力作用于 G。到最低点时 $\dot\theta^2=3g/l$，$\ddot\theta=0$。质心只有指向 O 的法向加速度，大小为 $(l/2)\dot\theta^2=3g/2$。于是
\[
H_x=0,\qquad H_y-mg=m\frac{3g}{2},
\quad\boxed{\vct H=(0,\tfrac52mg)}.
\]
绕 O 的能量方程消去了铰链力，恢复该力时应使用质心的平动方程，不能把整根杆视为集中在末端的质点。
在水平释放点，$\dot\theta=0$、$a_{Gy}=-3g/4$，于是 $H_y=mg/4$；到最低点，切向加速度为零而向心加速度最大，铰链竖直力增大到 $5mg/2$。这两个位置的铰链力不同，是因为质心加速度发生了变化。最低点的角加速度虽为零，质心仍有向心加速度，故不能作静力平衡。

\end{solution}


\begin{exercise}[有转动惯量的滑轮]\label{sup1:problem-exercise-6}\label{sup1:exercise:d1-p6}
轻绳跨过轴固定的滑轮，绳两端悬挂 $m_1,m_2$，其中 $m_2>m_1$。滑轮半径为 $r$，转动惯量为 $I$，轴承无摩擦，绳与滑轮之间不滑动。系统从静止释放，求加速度及绳两侧张力。
\end{exercise}

\begin{figure}[H]\centering
\begin{tikzpicture}
\draw[curve](-1.5,3.3)--(1.6,3.3);\draw(0,3.3)--(0,2.5);\draw[curve,fill=white](0,2.5)circle(.5);\draw[curve](-.5,.65)--(-.5,2.5)arc(180:0:.5)--(.5,.05);\draw[fill=white](-0.825,0.09999999999999998)rectangle(-0.175,0.65);\node at(-0.5,0.375){$m_1$};\draw[fill=white](0.175,-0.5)rectangle(0.825,0.05000000000000002);\node at(0.5,-0.225){$m_2$};\draw(0,2.5)--(.35,2.85)node[midway,left]{$r$};\node at(1.5,2.7){$I$};\draw[axis](1.80000,1.00000)--(1.80000,0.00000)node[right]{$q$};

\end{tikzpicture}
\caption{两端位移大小相等，滑轮自转由绳速决定。}\label{sup1:fig:d1-p6-q}
\end{figure}
\par\medskip
\begin{solution}
取 $m_2$ 向下位移为 $q$，绳速为 $\dot q$，滑轮角速度大小为 $\dot q/r$。于是
\[
T=\frac12\left(m_1+m_2+\frac I{r^2}\right)\dot q^2,
\quad U=-(m_2-m_1)gq.
\]
能量方程给出
\[
\boxed{a=\frac{(m_2-m_1)g}{m_1+m_2+I/r^2}}.
\]

\begin{figure}[H]\centering
\begin{tikzpicture}
\draw[curve](-1.5,3.3)--(1.6,3.3);\draw(0,3.3)--(0,2.5);\draw[curve,fill=white](0,2.5)circle(.5);\draw[curve](-.5,.65)--(-.5,2.5)arc(180:0:.5)--(.5,.05);\draw[fill=white](-0.825,0.09999999999999998)rectangle(-0.175,0.65);\node at(-0.5,0.375){$m_1$};\draw[fill=white](0.175,-0.5)rectangle(0.825,0.05000000000000002);\node at(0.5,-0.225){$m_2$};\draw(0,2.5)--(.35,2.85)node[midway,left]{$r$};\node at(1.5,2.7){$I$};\draw[axis](-0.50000,0.65000)--(-0.50000,1.60000)node[left]{$T_1$};
\draw[axis](0.50000,0.05000)--(0.50000,1.00000)node[right]{$T_2$};
\draw[axis](1.80000,0.50000)--(1.80000,-0.40000)node[right]{$a$};

\end{tikzpicture}
\caption{由重物平动恢复两个张力，二者的力矩差驱动滑轮转动。}\label{sup1:fig:d1-p6-s}
\end{figure}

图\ref{sup1:fig:d1-p6-s}中轻绳两侧张力不相等，差值提供滑轮角加速度。对两个重物分别列式，
\[
\boxed{T_1=m_1(g+a)},\qquad
\boxed{T_2=m_2(g-a)}.
\]
二者确实满足 $(T_2-T_1)r=I a/r$。轻绳内部任一不跨越滑轮接触区的自由段张力相同，但轻绳绕过有转动惯量的滑轮后，两自由端的张力可以不同。
若从受力直接消元，三个方程为
\[
T_1-m_1g=m_1a,\qquad m_2g-T_2=m_2a,
\qquad T_2-T_1=\frac I{r^2}a.
\]
前两式相加再用第三式，得到 $[(m_1+m_2)+I/r^2]a=(m_2-m_1)g$。由此看出，$I/r^2$ 只增加运动惯性，并不增加悬挂物的重力驱动。

两个张力的显式形式为
\[
T_1=\frac{m_1(2m_2+I/r^2)}{m_1+m_2+I/r^2}g,
\qquad
T_2=\frac{m_2(2m_1+I/r^2)}{m_1+m_2+I/r^2}g.
\]
因此，$I>0$ 时 $T_2>T_1>0$；$I=0$ 时两张力相等，恢复无质量滑轮。$I$ 很大时，加速度趋零，两张力分别接近各自重物的重量，但仍有张力差提供很小的角加速度。图\ref{sup1:fig:d1-p6-s}中滑轮与绳的接触可传递切向力，理想无摩擦轴承并不意味着这一接触区也没有切向作用。

\end{solution}


\begin{exercise}[水平地面上的滑滚转换]\label{sup1:problem-exercise-9}\label{sup1:exercise:d1-p9}
刚体质量为 $m$、半径为 $r$、转动惯量为 $\kappa mr^2$，在粗糙水平地面上运动。初始质心速度 $v_0\ge0$ 向右，初始自转角速度 $\omega_0\ge0$ 按图示顺时针方向为正。动摩擦系数为 $\mu_k>0$。求最终纯滚动速度、转换所需时间和产热。
\end{exercise}

\begin{figure}[H]\centering
\begin{tikzpicture}
\draw[curve](-1.7,0)--(3,0);\draw[fill=white,curve](0,.7)circle(.7);\fill(0.00000,0.70000)circle(1.3pt)node[above left]{$C$};
\fill(0.00000,0.00000)circle(1.3pt)node[below left]{$P$};
\draw[axis](0.00000,0.70000)--(1.30000,0.70000)node[above]{$v$};
\draw[axis](-0.845723359,1.007818129)arc(160:20:0.9);\node at(0,1.95){$\omega$};\node at(2.3,1.8){$I=\kappa mr^2$};
\end{tikzpicture}
\caption{向右平动与顺时针自转取正；初始接触滑动可向左也可向右。}\label{sup1:fig:d1-p9-q}
\end{figure}
\par\medskip
\begin{solution}
设 $u_0=v_0-r\omega_0$。若 $u_0=0$，初态已是纯滚动，转换时间和热量均为零。以下取 $u_0\ne0$。图\ref{sup1:fig:d1-p9-s}画的是 $u>0$ 的情形；$u<0$ 时摩擦方向反向。

\begin{figure}[H]\centering
\begin{tikzpicture}
\draw[curve](-1.7,0)--(3,0);\draw[fill=white,curve](0,.7)circle(.7);\fill(0.00000,0.70000)circle(1.3pt)node[above left]{$C$};
\fill(0.00000,0.00000)circle(1.3pt)node[below left]{$P$};
\draw[axis](0.00000,0.70000)--(1.30000,0.70000)node[above]{$v$};
\draw[axis](-0.845723359,1.007818129)arc(160:20:0.9);\node at(0,1.95){$\omega$};\draw[axis](0.00000,0.00000)--(-1.05000,0.00000)node[below]{$f$};
\node[right]at(2.2,1.75){$u=v-r\omega$};\node[right]at(2.2,.85){图取 $u>0$};
\end{tikzpicture}
\caption{摩擦方向由接触滑动速度决定，顺时针自转对应接触点相对轮心向左的速度。}\label{sup1:fig:d1-p9-s}
\end{figure}

以 $\sigma=\operatorname{sgn}u_0$ 表示摩擦方向，滑动阶段有
\[
\dot v=-\mu_kg\sigma,\qquad
\dot\omega=\frac{\mu_kg}{\kappa r}\sigma.
\]
两式相加时，$\dd(v+\kappa r\omega)/\dd t=-\mu_kg\sigma+\mu_kg\sigma=0$，故该组合为常量。相减则有 $\dot u=-(1+1/\kappa)\mu_kg\sigma$，用初值 $u(0)=u_0$ 积分，得到
\[
u=u_0-\left(1+\frac1\kappa\right)\mu_kg\sigma t.
\]
首次达到 $u=0$ 的时间及速率为
\[
\boxed{t_*=\frac{\kappa|v_0-r\omega_0|}{(1+\kappa)\mu_kg}},\qquad
\boxed{v_f=\frac{v_0+\kappa r\omega_0}{1+\kappa}}.
\]
以后无需水平摩擦即可保持匀速纯滚动。在滑动阶段 $|u(t)|=|u_0|(1-t/t_*)$，产热为
\[
Q_{\text{热}}=\mu_kmg|u_0|\int_0^{t_*}\left(1-\frac t{t_*}\right)\dd t
=\mu_kmg|u_0|\left[t-\frac{t^2}{2t_*}\right]_0^{t_*}
=\frac12\mu_kmg|u_0|t_*.
\]
这也等于相对滑动速度图下的三角形面积乘以摩擦力。代入 $t_*$，得到
\[
\boxed{Q_{\text{热}}=\frac12m\frac{\kappa}{1+\kappa}
(v_0-r\omega_0)^2}.
\]
这一结果也等于初始总动能与 $\tfrac12m(1+\kappa)v_f^2$ 的差。$\mu_k$ 改变转换时间，却不改变给定初态的最终速度和总产热。
用摩擦冲量也能独立求出最终速度。记滑动阶段摩擦的有向冲量为 J，则
\[
m(v_f-v_0)=J,\qquad I(\omega_f-\omega_0)=-rJ.
\]
代入 $r\omega_f=v_f$，可得 $J=-m\kappa(v_0-r\omega_0)/(1+\kappa)$。由 $|J|=\mu_kmg t_*$ 恢复同一转换时间。自转过快时 $u_0<0$，摩擦向右，质心加速而自转减慢；这类初态下，摩擦将部分自转能量转入质心平动。

动能还可精确分解为
\[
T=\frac{m}{2(1+\kappa)}(v+\kappa r\omega)^2
+\frac{m\kappa}{2(1+\kappa)}(v-r\omega)^2.
\]
第一项由滑动阶段的守恒组合决定，第二项在 $u\to0$ 时全部耗散，因此直接给出前述总热量。图\ref{sup1:fig:d1-p9-speeds}画出两种相反滑动初态，v 与 $r\omega$ 在 $t=t_*$ 相遇后共同保持常量。
\begin{figure}[H]\centering
\begin{tikzpicture}[scale=.95]
\draw[axis](-.1,0)--(3.6,0)node[right]{$t/t_*$};\draw[axis](0,-.1)--(0,1.9)node[above]{速度$/U_0$};
\draw[curve](0,1.5)--(2.4,1)--(3.3,1);\draw[guide](0,0)--(2.4,1)--(3.3,1);
\node[left]at(0,1.5){$1$};\node[left]at(0,1){$2/3$};\node[below]at(2.4,0){$1$};\draw[guide](2.4,0)--(2.4,1);
\node at(.75,1.6){$v$};\node at(1.3,.35){$r\omega$};
\begin{scope}[shift={(5.5,0)}]
\draw[axis](-.1,0)--(3.6,0)node[right]{$t/t_*$};\draw[axis](0,-.1)--(0,1.9)node[above]{速度$/U_0$};
\draw[guide](0,1.5)--(2.4,.5)--(3.3,.5);\draw[curve](0,0)--(2.4,.5)--(3.3,.5);
\node[left]at(0,1.5){$1$};\node[left]at(0,.5){$1/3$};\node[below]at(2.4,0){$1$};\draw[guide](2.4,0)--(2.4,.5);
\node at(.6,1.6){$r\omega$};\node at(1.3,.12){$v$};
\end{scope}
\end{tikzpicture}
\caption{取 $\kappa=1/2$、$U_0>0$。左图初态为 $v_0=U_0,\omega_0=0$；右图为 $v_0=0,r\omega_0=U_0$。实线表示平动速率，虚线表示轮缘自转速率。}\label{sup1:fig:d1-p9-speeds}
\end{figure}

\end{solution}


\begin{exercise}[摆杆越过最高点的条件（选做）]\label{sup1:problem-exercise-8}\label{sup1:exercise:d1-p8}
沿用例\ref{sup1:example:d1-e8}的小车与刚性轻杆摆。系统总水平动量为零，摆从最低点 $\theta=0$ 以角速度 $\omega_0>0$ 开始运动。求在有限时间内越过最高点所需的 $\omega_0$ 条件，并说明临界等号对应的运动。
\end{exercise}

\begin{figure}[H]\centering
\begin{tikzpicture}
\draw[curve](-1.8,0)--(3.3,0);\draw[fill=white,curve](-.8,.35)rectangle(.8,.95);\node at(0,.65){$M$};\draw(-.55,.175)circle(.175);\draw(.55,.175)circle(.175);\draw[curve](-.68,.95)--(-.68,3.9)--(0,3.9);\fill(0.0000000,3.9000000)circle(1.3pt)node[above right]{$O$};
\draw[line width=.8pt](0,3.9)--(1.469463131,1.877457514);\fill(1.4694631,1.8774575)circle(1.3pt)node[right]{$m$};
\draw[guide](0,3.9)--(0,1.2);\draw(0,3.28)arc(-90:-54:.62);\node at(.27,3.0){$\theta$};\node at(0.9847315653655915,2.9887287570313155){$l$};\draw[guide](0,1.4)arc(-90:-54:2.5);\draw[fill=white](0,1.4)circle(1.4pt);\node[right]at(.12,1.22){最低位置};\draw[axis](-1.8500000,1.2500000)--(-1.0500000,1.2500000)node[above]{$X$};
\node[right]at(2.3,3.1){刚性轻杆};
\end{tikzpicture}
\caption{图示摆杆上升时的位置。支架位于摆动平面后方，摆球的整周轨迹高于车身。}\label{sup1:fig:d1-p8-q}
\end{figure}
\par\medskip
\begin{solution}
消去循环速度后的惯性系数为
\[
B_{\text{eff}}(\theta)=ml^2\left(1-\frac{m}{M+m}\cos^2\theta\right)>0.
\]
初态可用能量为 $E=ml^2M\omega_0^2/[2(M+m)]$。图\ref{sup1:fig:d1-p8-s}中势能在 $\theta=\pi$ 处达到 $2mgl$。有限时间越过该势垒须有 $E>2mgl$，即
\[
\boxed{\omega_0^2>\frac{4(M+m)g}{Ml}}.
\]

\begin{figure}[H]\centering
\begin{tikzpicture}
\draw[axis](-0.1,0)--(3.6,0)node[right]{$\theta$};\draw[axis](0,-0.2)--(0,2.8)node[above]{$U/(mgl)$};\draw[curve]plot coordinates{(0.000000,0.000000)(0.031416,0.000493)(0.062832,0.001973)(0.094248,0.004438)(0.125664,0.007885)(0.157080,0.012312)(0.188496,0.017713)(0.219911,0.024083)(0.251327,0.031417)(0.282743,0.039706)(0.314159,0.048943)(0.345575,0.059119)(0.376991,0.070224)(0.408407,0.082245)(0.439823,0.095173)(0.471239,0.108993)(0.502655,0.123693)(0.534071,0.139258)(0.565487,0.155672)(0.596903,0.172919)(0.628319,0.190983)(0.659734,0.209845)(0.691150,0.229487)(0.722566,0.249889)(0.753982,0.271031)(0.785398,0.292893)(0.816814,0.315453)(0.848230,0.338688)(0.879646,0.362576)(0.911062,0.387093)(0.942478,0.412215)(0.973894,0.437917)(1.005310,0.464173)(1.036726,0.490959)(1.068142,0.518246)(1.099557,0.546010)(1.130973,0.574221)(1.162389,0.602852)(1.193805,0.631875)(1.225221,0.661262)(1.256637,0.690983)(1.288053,0.721009)(1.319469,0.751310)(1.350885,0.781857)(1.382301,0.812619)(1.413717,0.843566)(1.445133,0.874667)(1.476549,0.905892)(1.507964,0.937209)(1.539380,0.968589)(1.570796,1.000000)(1.602212,1.031411)(1.633628,1.062791)(1.665044,1.094108)(1.696460,1.125333)(1.727876,1.156434)(1.759292,1.187381)(1.790708,1.218143)(1.822124,1.248690)(1.853540,1.278991)(1.884956,1.309017)(1.916372,1.338738)(1.947787,1.368125)(1.979203,1.397148)(2.010619,1.425779)(2.042035,1.453990)(2.073451,1.481754)(2.104867,1.509041)(2.136283,1.535827)(2.167699,1.562083)(2.199115,1.587785)(2.230531,1.612907)(2.261947,1.637424)(2.293363,1.661312)(2.324779,1.684547)(2.356194,1.707107)(2.387610,1.728969)(2.419026,1.750111)(2.450442,1.770513)(2.481858,1.790155)(2.513274,1.809017)(2.544690,1.827081)(2.576106,1.844328)(2.607522,1.860742)(2.638938,1.876307)(2.670354,1.891007)(2.701770,1.904827)(2.733186,1.917755)(2.764602,1.929776)(2.796017,1.940881)(2.827433,1.951057)(2.858849,1.960294)(2.890265,1.968583)(2.921681,1.975917)(2.953097,1.982287)(2.984513,1.987688)(3.015929,1.992115)(3.047345,1.995562)(3.078761,1.998027)(3.110177,1.999507)(3.141593,2.000000)};
\draw[guide](0,2)--(3.35,2)node[right]{$E_c/(mgl)=2$};\draw[guide](0,2.5)--(3.35,2.5)node[right]{$E>E_c$};\draw[guide](3.14159,0)--(3.14159,2);\node[below]at(3.14159,0){$\pi$};
\end{tikzpicture}
\caption{临界能量恰过势垒顶，超临界能量曲线在全程保留正的动能。}\label{sup1:fig:d1-p8-s}
\end{figure}

等号时，在最高点附近令 $\varepsilon=\pi-\theta$，有
\[
E-U\simeq\frac12mgl\varepsilon^2,\qquad
B_{\text{eff}}(\pi)=\frac{ml^2M}{M+m}.
\]
故 $\dot\theta\simeq\Omega\varepsilon$，其中 $\Omega=\sqrt{(M+m)g/(Ml)}$。以接近顶点的某时刻 $t_1$ 为起点，分离变量并积分，得
\[
t-t_1\simeq\frac1\Omega\int_{\varepsilon(t)}^{\varepsilon(t_1)}\frac{\dd\xi}{\xi}
=\frac1\Omega[\ln\xi]_{\varepsilon(t)}^{\varepsilon(t_1)}
=\frac1\Omega\ln\frac{\varepsilon(t_1)}{\varepsilon(t)}.
\]
当 $\varepsilon(t)\to0$ 时，时间呈对数发散。系统只渐近接近最高点，不会在有限时间越过。因为连接物是刚性杆，途中不必另加绳张力非负的限制；若改为软绳，能量超过势垒并不一定足以保证绳始终绷紧。
超临界时，将 $\theta=\pi$ 代入能量积分，可直接得到过顶角速度
\[
\dot\theta_{\rm top}^2=
\omega_0^2-\frac{4(M+m)g}{Ml}>0.
\]
由于全过程 $B_{\rm eff}>0$，而 $U\le2mgl<E$，角速度沿上升过程始终不为零。故确实可以连续越过最高点，并且行程时间在这一有限角度区间内为有限值。

将临界情况下的时间关系反解为偏角，得到
\[
\varepsilon(t)\simeq\varepsilon(t_1)\ee^{-\Omega(t-t_1)}.
\]
摆球越来越接近顶点，却在任意有限时刻保留一个非零偏角。若初始能量更低，转向角已在 $\pi$ 以前出现，此时不能把能量式继续延伸到负的角速度平方区域。

图\ref{sup1:fig:d1-p8-s}中的势能曲线在顶端是极大值，过顶条件的严格不等号由这一势垒附近的时间积分决定。本题的刚性轻杆不另设拉压强度限制。

\end{solution}


\begin{exercise}[牵连斜面的接触余量]\label{deep1:problem-exercise-3}\label{deep1:exercise:d1-p3}
沿用例\ref{ex:hw1-1}的装置，取 $M=m$、$\alpha=60^\circ$，各接触面光滑。求斜面加速度、绳张力和物块所受支持力。再求该角度下允许维持接触的最小斜面质量。
\end{exercise}

\begin{figure}[H]\centering
\begin{tikzpicture}
\draw[curve](-.4,.1)--(5.5,.1);\draw[curve](5.1,.1)--(5.1,3.35);\draw[curve](2.070132856636795,0.1)--(3.5192820323027556,2.61)--(3.5192820323027556,0.1)--cycle;\begin{scope}[shift={(2.3942820323027556,0.6614428414850132)},rotate=60]\draw[fill=white](-.36,0)rectangle(.36,.48);\node at(0,.24){$m$};\end{scope}\draw[curve](2.366435935,1.093211987)--(3.31143593539449,2.73)arc(150.0:90:.16)--(5.1,2.81);\draw[fill=white](3.45,2.65)circle(.16);\draw(3.45,2.65)--(3.5192820323027556,2.61);\fill(3.45000,2.65000)circle(1.3pt)node[above right]{$P$};
\node at(3.0362323070807684,0.9366666666666666){$M$};\draw(2.720132856636795,.1)arc(0:60:.65);\node at(3.0201328566367946,.38){$\alpha$};\draw[axis](3.75000,0.45000)--(4.75000,0.45000)node[above]{$X$};
\draw[axis](1.3936156,1.0082378)--(1.0186156,0.3587187)node[above left]{$s$};

\end{tikzpicture}
\caption{斜面与物块由固定绳端关联，沿用 $s=X$ 的坐标。}\label{deep1:fig:d1-p3-q}
\end{figure}
\par\medskip
\begin{solution}
图\ref{deep1:fig:d1-p3-s}中物块的地面加速度有水平和竖直分量。惯性系数为
\[
M+2m(1-\cos60^\circ)=2m,
\]
故
\[
\boxed{a=\frac{\sqrt3}{4}g}.
\]

\begin{figure}[H]\centering
\begin{tikzpicture}
\draw[curve](-.4,.1)--(5.5,.1);\draw[curve](5.1,.1)--(5.1,3.35);\draw[curve](2.070132856636795,0.1)--(3.5192820323027556,2.61)--(3.5192820323027556,0.1)--cycle;\begin{scope}[shift={(2.3942820323027556,0.6614428414850132)},rotate=60]\draw[fill=white](-.36,0)rectangle(.36,.48);\node at(0,.24){$m$};\end{scope}\draw[curve](2.366435935,1.093211987)--(3.31143593539449,2.73)arc(150.0:90:.16)--(5.1,2.81);\draw[fill=white](3.45,2.65)circle(.16);\draw(3.45,2.65)--(3.5192820323027556,2.61);\fill(3.45000,2.65000)circle(1.3pt)node[above right]{$P$};
\node at(3.0362323070807684,0.9366666666666666){$M$};\draw(2.720132856636795,.1)arc(0:60:.65);\node at(3.0201328566367946,.38){$\alpha$};\draw[axis](3.75000,0.45000)--(4.75000,0.45000)node[above]{$X$};
\draw[axis](1.3936156,1.0082378)--(1.0186156,0.3587187)node[above left]{$s$};
\begin{scope}[shift={(7.1,1.2)}]
\fill(0.00000,0.00000)circle(1.3pt)node[below left]{$C$};
\draw[axis](0.00000,0.00000)--(-0.69282,0.40000)node[above left]{$N$};
\draw[axis](0.00000,0.00000)--(0.42500,0.73612)node[above right]{$T$};
\draw[axis](0.00000,0.00000)--(0.00000,-1.00000)node[below]{$mg$};
\node at(0,1.6){物块受力};
\end{scope}

\end{tikzpicture}
\caption{加速度求出后，分别沿斜面与法线方向恢复张力和支持力。}\label{deep1:fig:d1-p3-s}
\end{figure}

沿图中法向与斜面方向投影，代入例\ref{ex:hw1-1}所得受力式，得到
\[
\boxed{N=mg\cos60^\circ-ma\sin60^\circ=\frac18mg},
\]
\[
\boxed{T=mg\sin60^\circ-ma(1-\cos60^\circ)
=\frac{3\sqrt3}{8}mg}.
\]
两者均为正，当前运动能够保持绳绷紧和接触。令 $N=0$，接触边界为
\[
M\cos60^\circ=m(1-\cos60^\circ)^2,
\quad\boxed{M_{\min}=\frac m2}.
\]
绳长关系确定运动关联，单侧接触条件则给出这一质量下限。
物块的地面加速度由坐标关系给出：
\[
\vct a_C=a(1-\cos60^\circ,-\sin60^\circ)
=\left(\frac{\sqrt3}{8}g,-\frac38g\right).
\]
其外法线分量为 $-a\sin60^\circ=-3g/8$，所以法向方程为 $N-mg/2=-3mg/8$，只剩 $N=mg/8$。若将相对斜面的法向加速度为零，误当成地面系中也为零，就会错误地得到 $mg/2$。

固定该角度，以 $\lambda=M/m$ 表示质量比，三个结果可统一写成
\[
a=\frac{\sqrt3g}{2(\lambda+1)},\qquad
N=\frac{(2\lambda-1)mg}{4(\lambda+1)},\qquad
T=\frac{\sqrt3(2\lambda+1)mg}{4(\lambda+1)}.
\]
张力对所有 $\lambda>0$ 均为正，支持力则仅在 $\lambda\ge1/2$ 时允许。当前质量比 $\lambda=1$ 大于临界值 $1/2$，故支持力为正；减轻斜面会增大牵连加速度、减小法向接触余量。$\lambda=1/2$ 给出 N 为零的临界约束解，$\lambda<1/2$ 则已要求负支持力，不能再沿用贴面运动公式。

\end{solution}


\begin{exercise}[圆弧轨道上的滚动]\label{deep1:problem-exercise-5}\label{deep1:exercise:d1-p5}
轴对称刚体的质量为 $m$、半径为 $r$，转动惯量为 $\kappa mr^2$。刚体在固定圆弧轨道内侧纯滚动，质心 C 的轨迹半径为 $\rho$。从 $\theta=\theta_0$ 处静止释放，其中 $0<\theta_0<\pi/2$，$\theta$ 从竖直向下方向量起。求下降时的质心速率、$\ddot\theta$ 和法向力 $N$。
\end{exercise}

\begin{figure}[H]\centering
\begin{tikzpicture}
\draw[curve](2.15,0)arc(0:-180:2.15);\draw[guide](-1.9,0)arc(180:360:1.9);\draw(0,0)--(-1.343502884,-1.343502884)node[midway,above left]{$\rho$};\draw[guide](0,.2)--(0,-2.4);\fill(0.0000000,0.0000000)circle(1.3pt)node[above left]{$O$};
\draw[guide](1.645448267,-0.950000000)circle(0.25);\draw[guide](0,0)--(1.645448267,-0.950000000);\fill(1.6454483,-0.9500000)circle(1.3pt)node[right=10pt]{$C_0$};
\draw[fill=white,curve](1.089795229,-1.556388884)circle(0.25);\fill(1.0897952,-1.5563889)circle(1.3pt)node[left=10pt,yshift=-2pt,fill=white,inner sep=1pt]{$C$};
\draw(0,0)--(1.089795229,-1.556388884);\draw(0,-.62)arc(-90:-55:.62);\node at(.2,-.8){$\theta$};\draw(0,-1.1)arc(-90:-30:1.1);\node at(.55,-1.2){$\theta_0$};\fill(1.2331893,-1.7611769)circle(1.3pt)node[below right]{$P$};

\end{tikzpicture}
\caption{实线圆与虚线圆分别表示当前状态和释放状态，圆心为 $C$、$C_0$。两角均从 $O$ 处的竖直向下方向量起。}\label{deep1:fig:d1-p5-q}
\end{figure}
\par\medskip
\begin{solution}
接触点相对固定轨道静止，故 $|\omega|=\rho|\dot\theta|/r$。由图\ref{deep1:fig:d1-p5-s}中 C$_0$ 与 C 的高度差，能量关系为
\[
\frac12m(1+\kappa)\rho^2\dot\theta^2
=mg\rho(\cos\theta-\cos\theta_0).
\]
由 $v^2=\rho^2\dot\theta^2$，先得
\[
\boxed{v^2=\frac{2g\rho(\cos\theta-\cos\theta_0)}{1+\kappa}}.
\]
求角加速度时，取 $U=mg\rho(1-\cos\theta)$，有
\[
\frac{\partial T}{\partial\dot\theta}=m(1+\kappa)\rho^2\dot\theta,
\qquad \frac{\partial U}{\partial\theta}=mg\rho\sin\theta.
\]
将第一式对时间求导，代入拉格朗日方程，
\[
m(1+\kappa)\rho^2\ddot\theta+mg\rho\sin\theta=0,
\qquad\boxed{\ddot\theta=-\frac{g\sin\theta}{(1+\kappa)\rho}}.
\]
这一推导不需要除以 $\dot\theta$，在静止释放点同样成立。

\begin{figure}[H]\centering
\begin{tikzpicture}
\draw[curve](2.15,0)arc(0:-180:2.15);\draw[guide](-1.9,0)arc(180:360:1.9);\draw(0,0)--(-1.343502884,-1.343502884)node[midway,above left]{$\rho$};\draw[guide](0,.2)--(0,-2.4);\fill(0.0000000,0.0000000)circle(1.3pt)node[above left]{$O$};
\draw[guide](1.645448267,-0.950000000)circle(0.25);\draw[guide](0,0)--(1.645448267,-0.950000000);\fill(1.6454483,-0.9500000)circle(1.3pt)node[right=10pt]{$C_0$};
\draw[fill=white,curve](1.089795229,-1.556388884)circle(0.25);\fill(1.0897952,-1.5563889)circle(1.3pt)node[left=10pt,yshift=-2pt,fill=white,inner sep=1pt]{$C$};
\draw(0,0)--(1.089795229,-1.556388884);\draw(0,-.62)arc(-90:-55:.62);\node at(.2,-.8){$\theta$};\draw(0,-1.1)arc(-90:-30:1.1);\node at(.55,-1.2){$\theta_0$};\fill(1.2331893,-1.7611769)circle(1.3pt)node[below right]{$P$};
\draw[axis](1.2331893,-1.7611769)--(0.7169705,-1.0239401)node[right=5pt]{$N$};
\draw[axis](1.0897952,-1.5563889)--(1.0897952,-2.4563889)node[below]{$mg$};
\draw[axis](1.2331893,-1.7611769)--++({.8*cos(35)},{.8*sin(35)})node[below right,xshift=2pt]{$f_s$};
\node[right]at(2.9,-1.65){$v^2=\rho^2\dot\theta^2$};
\end{tikzpicture}
\caption{支持力 $N$ 由接触点 $P$ 沿半径指向 $O$，重力作用于质心 $C$；静摩擦 $f_s$ 在 $P$ 沿角度增大的切向作用。}\label{deep1:fig:d1-p5-s}
\end{figure}

图\ref{deep1:fig:d1-p5-s}中法向力指向 O。沿 C 到 O 方向列式，
\[
N-mg\cos\theta=m\frac{v^2}{\rho},
\quad\boxed{N=mg\left[\cos\theta+
\frac{2(\cos\theta-\cos\theta_0)}{1+\kappa}\right]}.
\]
在下降区间 $0\le\theta\le\theta_0$，支持力为正。上述角加速度只适用于保持纯滚动的阶段，不能用于因静摩擦不足而发生滑动的轨道。
沿图\ref{deep1:fig:d1-p5-s}中角度增大的切向取正，重力切向分量为 $-mg\sin\theta$。由角加速度可得质心切向加速度 $a_t=\rho\ddot\theta=-g\sin\theta/(1+\kappa)$。于是摩擦的有向分量为
\[
f_s-mg\sin\theta=m\rho\ddot\theta,
\qquad f_s=\frac{\kappa}{1+\kappa}mg\sin\theta>0.
\]
在初次下降中 $\dot\theta<0$，这个方向沿轨道向上，所产生的力矩恰好维持相应自转。轨道曲率进入质心法向加速度时使用的是 $\rho$，自转约束使用的滚动物体半径则是 r，两者不可互换。

如需核对纯滚动条件，
\[
\frac{|f_s|}{N}=\frac{\kappa\sin\theta}
{(\kappa+3)\cos\theta-2\cos\theta_0}.
\]
记分母为 $D(\theta)=(\kappa+3)\cos\theta-2\cos\theta_0$。在 $0\le\theta\le\theta_0$ 上，$D\ge(\kappa+1)\cos\theta_0>0$，按商法则求导，
\[
\begin{aligned}
\frac{\dd}{\dd\theta}\frac{|f_s|}{N}
&=\frac{\kappa\cos\theta\,D(\theta)+\kappa(\kappa+3)\sin^2\theta}{D(\theta)^2}\\
&=\frac{\kappa[(\kappa+3)-2\cos\theta_0\cos\theta]}{D(\theta)^2}>0.
\end{aligned}
\]
所以最大比值出现在释放点，所需 $\mu_s\ge\kappa\tan\theta_0/(1+\kappa)$。这一条件满足后，下降过程中所需摩擦比值反而减小，且 N 保持为正。

\end{solution}


\begin{exercise}[带内部阻尼的可动斜面]\label{deep1:problem-exercise-7}\label{deep1:exercise:d1-p7}
沿用例\ref{ex:hw1-5}的自由斜面与带轮小车，另在小车和斜面之间安装阻尼器。本题将斜面质量记为 $M_w$、倾角记为 $\alpha$；车身质量为 $M$，两个均匀实心轮的质量各为 $m$、半径各为 $r$，每个轮的转动惯量为 $\kappa mr^2$，其中 $\kappa=1/2$。按图\ref{deep1:fig:d1-p7-q}，取沿坡向下的相对位移为 $s$，斜面向右的位移为 $X$。阻尼器对小车的力沿斜面，大小为 $b|\dot s|$，方向与相对速度相反；对斜面的反作用力与其等大反向。$b>0$。系统从静止释放，滚动始终保持。求 $\dot s(t)$、$\dot X(t)$ 和单位时间产热，并给出长期速度极限。
\end{exercise}

\begin{figure}[H]\centering
\begin{tikzpicture}
\draw[curve](-.7,0)--(6.6,0);\draw[curve](0,0)--(0,3)--(6,0.20215405107000795)--(6,0)--cycle;\begin{scope}[shift={(2.15,1.997438534966753)},rotate=-25]\draw[curve,fill=white](-.9,.7)rectangle(1.1,1.5);\node at(.1,1.1){$M$};\draw[curve,fill=white](-0.6,.35)circle(.35);\draw(-0.6,.35)--(-0.6,.7);\fill(-0.6,.35)circle(1pt);\draw[curve,fill=white](0.8,.35)circle(.35);\draw(0.8,.35)--(0.8,.7);\fill(0.8,.35)circle(1pt);\node at(2,1.2){$m,r$};\draw[axis](1.5,.4)--(2.5,.4)node[right]{$s$};\end{scope}\begin{scope}[shift={(2.15,1.997438534966753)},rotate=-25]\draw(-2.2,0)--(-2.2,1.1)--(-1.8,1.1);\draw(-1.8,.94)rectangle(-1.45,1.26);\draw(-1.6,.98)--(-1.6,1.22);\draw(-1.6,1.1)--(-.9,1.1);\node[above]at(-1.7,1.4){$b$};\end{scope}\draw[guide](4.7,0.20215405107000795)--(6.2,0.20215405107000795);\draw(5.3,0.20215405107000795)arc(180:155:.7);\node at(5.03,0.5621540510700079){$\alpha$};\node at(2.5,.7){$M_w$};\draw[axis](5.40000,1.30000)--(6.40000,1.30000)node[above]{$X$};

\end{tikzpicture}
\caption{阻尼器阻碍小车相对斜面的运动，其作用力和反作用力均属于系统内部。}\label{deep1:fig:d1-p7-q}
\end{figure}
\par\medskip
\begin{solution}
阻尼器的两力对 $X$ 的广义力相消，对 $s$ 的广义力为 $-b\dot s$，所以总水平动量仍为零。沿用题设的质量和倾角符号，记
\[
A=M+2m,\ B=A+2\kappa m,\ D=M_w+A,
\quad B_{\text{eff}}=B-\frac{A^2\cos^2\alpha}{D}.
\]
其中 $A$ 是小车及两轮的总质量，$B$ 还计入轮子的转动惯性。系统动能与水平动量方程为
\[
T=\frac12D\dot X^2+A\cos\alpha\,\dot X\dot s
+\frac12B\dot s^2,\qquad
D\dot X+A\cos\alpha\,\dot s=0.
\]
沿 $s$ 的运动方程为
\[
A\cos\alpha\,\ddot X+B\ddot s=Ag\sin\alpha-b\dot s.
\]
将水平动量方程对时间求导，得到
$\ddot X=-A\cos\alpha\,\ddot s/D$。代回上式，合并 $\ddot s$ 的系数，得
\[
B_{\text{eff}}\ddot s+b\dot s=Ag\sin\alpha.
\]

\begin{figure}[H]\centering
\begin{tikzpicture}
\draw[axis](-0.1,0)--(4.5,0)node[right]{$t/\tau$};\draw[axis](0,-0.2)--(0,2.2)node[above]{$\dot s/V$};\draw[curve]plot coordinates{(0.000000,0.000000)(0.040000,0.062737)(0.080000,0.123014)(0.120000,0.180927)(0.160000,0.236570)(0.200000,0.290031)(0.240000,0.341395)(0.280000,0.390746)(0.320000,0.438162)(0.360000,0.483718)(0.400000,0.527488)(0.440000,0.569542)(0.480000,0.609947)(0.520000,0.648767)(0.560000,0.686065)(0.600000,0.721901)(0.640000,0.756332)(0.680000,0.789413)(0.720000,0.821196)(0.760000,0.851734)(0.800000,0.881074)(0.840000,0.909263)(0.880000,0.936347)(0.920000,0.962370)(0.960000,0.987371)(1.000000,1.011393)(1.040000,1.034473)(1.080000,1.056647)(1.120000,1.077952)(1.160000,1.098422)(1.200000,1.118089)(1.240000,1.136985)(1.280000,1.155140)(1.320000,1.172584)(1.360000,1.189343)(1.400000,1.205445)(1.440000,1.220916)(1.480000,1.235780)(1.520000,1.250061)(1.560000,1.263782)(1.600000,1.276966)(1.640000,1.289632)(1.680000,1.301802)(1.720000,1.313494)(1.760000,1.324728)(1.800000,1.335522)(1.840000,1.345892)(1.880000,1.355856)(1.920000,1.365429)(1.960000,1.374627)(2.000000,1.383464)(2.040000,1.391954)(2.080000,1.400112)(2.120000,1.407949)(2.160000,1.415480)(2.200000,1.422715)(2.240000,1.429666)(2.280000,1.436345)(2.320000,1.442762)(2.360000,1.448928)(2.400000,1.454851)(2.440000,1.460543)(2.480000,1.466011)(2.520000,1.471265)(2.560000,1.476312)(2.600000,1.481162)(2.640000,1.485822)(2.680000,1.490299)(2.720000,1.494600)(2.760000,1.498733)(2.800000,1.502704)(2.840000,1.506519)(2.880000,1.510184)(2.920000,1.513706)(2.960000,1.517090)(3.000000,1.520341)(3.040000,1.523464)(3.080000,1.526465)(3.120000,1.529349)(3.160000,1.532119)(3.200000,1.534780)(3.240000,1.537338)(3.280000,1.539795)(3.320000,1.542155)(3.360000,1.544424)(3.400000,1.546603)(3.440000,1.548697)(3.480000,1.550708)(3.520000,1.552641)(3.560000,1.554498)(3.600000,1.556282)(3.640000,1.557996)(3.680000,1.559643)(3.720000,1.561226)(3.760000,1.562746)(3.800000,1.564207)(3.840000,1.565610)(3.880000,1.566959)(3.920000,1.568254)(3.960000,1.569499)(4.000000,1.570695)};
\draw[guide](0,1.6)--(4.1,1.6);\node[left]at(0,1.6){$1$};\node at(2.6,.3){$\dot Q=b\dot s^2$};
\end{tikzpicture}
\caption{相对速度按指数规律接近终端速度；阻尼只消耗相对运动的机械能。}\label{deep1:fig:d1-p7-s}
\end{figure}

取 $\tau=B_{\text{eff}}/b$、$V=Ag\sin\alpha/b$，令相对速度为 $w=\dot s$，方程变为一阶线性方程
\[
\dot w+\frac w\tau=\frac V\tau,\qquad
\frac{\dd}{\dd t}(\ee^{t/\tau}w)=\frac V\tau\ee^{t/\tau}.
\]
对第二式从 0 积分至 $t$，使用 $w(0)=0$，
\[
\ee^{t/\tau}w(t)=\frac V\tau\int_0^t\ee^{u/\tau}\dd u
=\frac V\tau[\tau\ee^{u/\tau}]_0^t
=V(\ee^{t/\tau}-1).
\]
这里 $\tau$ 表示速度接近极限的时间尺度，$V$ 则由重力沿坡驱动与阻尼力平衡决定。

满足初始静止条件的解为
\[
\boxed{\dot s=V(1-\ee^{-t/\tau})},\qquad
\boxed{\dot X=-\frac{A\cos\alpha}{D}V(1-\ee^{-t/\tau})}.
\]
图\ref{deep1:fig:d1-p7-s}中相对速度趋于 $V$，相对位移为
\[
s=V\int_0^t(1-\ee^{-u/\tau})\dd u
=V[u+\tau\ee^{-u/\tau}]_0^t
=V\big[t-\tau(1-\ee^{-t/\tau})\big].
\]
阻尼器功率为两力总功率的相反数，故
\[
\boxed{\dot Q=b\dot s^2=bV^2(1-\ee^{-t/\tau})^2}.
\]
其长期极限为 $\dot s\to V$、$\dot X\to-A\cos\alpha V/D$、$\dot Q\to bV^2$。极限状态中小车和斜面均匀运动，重力势能的下降率恰好等于产热率，转动动能和平动动能均不再增长。
阻尼器两力的总功率为
\[
\mathcal P_b=-b\dot s\,\vct e_s\cdot(\vct v_C-\vct v_w)
=-b\dot s^2.
\]
共同的平移速度相消，因而产热取决于相对速度。把瞬时产热率展开后积分，
\[
\begin{aligned}
Q(t)&=bV^2\int_0^t(1-2\ee^{-u/\tau}+\ee^{-2u/\tau})\dd u\\
&=bV^2\left[u+2\tau\ee^{-u/\tau}-\frac\tau2\ee^{-2u/\tau}\right]_0^t.
\end{aligned}
\]
代入上下限，得到
\[
Q(t)=bV^2\left[t-2\tau(1-\ee^{-t/\tau})
+\frac\tau2(1-\ee^{-2t/\tau})\right].
\]
它满足 $Q(0)=0$、$\dot Q\ge0$。从静止开始的整体能量收支为
\[
Ags\sin\alpha=\frac12B_{\rm eff}\dot s^2+Q(t).
\]
左侧是下降势能，右侧已同时计入斜面、小车和轮子自转的动能。在短时间 $t\ll\tau$，$\dot s\simeq(Ag\sin\alpha/B_{\rm eff})t$，恢复无阻尼释放时的初始加速度；达到长期状态后，势能下降率才全部转化为产热率。

\end{solution}


\begin{exercise}[摩擦不足时会不会自行滚起来]\label{deep1:problem-exercise-10}\label{deep1:exercise:d1-p10}
刚体参数为 $m,r,I=\kappa mr^2$，从倾角为 $\alpha$ 的固定粗糙斜面上静止释放。静摩擦系数和动摩擦系数满足
\[
0<\mu_k\le\mu_s<\frac{\kappa}{1+\kappa}\tan\alpha.
\]
求释放后的质心加速度、角加速度和接触滑动速度 $u(t)$，并判断是否会在以后自动转为纯滚动。
\end{exercise}

\begin{figure}[H]\centering
\begin{tikzpicture}
\begin{scope}[rotate=-24]\draw[curve](-1.6,0)--(3.4,0);\draw[fill=white,curve](0,.7)circle(.7);\fill(0.00000,0.70000)circle(1.3pt)node[above left]{$C$};
\fill(0.00000,0.00000)circle(1.3pt)node[below left]{$P$};
\draw(0,.7)--(0.520201378,1.168391424)node[midway,above]{$r$};\draw[axis](1.20000,0.70000)--(2.10000,0.70000)node[above]{$s$};
\node at(1.8,1.4){$I=\kappa mr^2$};\end{scope}\draw[guide](1.3060545559848429,-1.3829045864577207)--(3.306054555984843,-1.3829045864577207);\draw(2.456054555984843,-1.3829045864577207)arc(180:156:.65);\node at(2.126054555984843,-1.1529045864577208){$\alpha$};
\end{tikzpicture}
\caption{从静止释放不代表能保持纯滚动，所需静摩擦首先受到上限限制。}\label{deep1:fig:d1-p10-q}
\end{figure}
\par\medskip
\begin{solution}
若假定纯滚动，所需摩擦系数为 $\kappa\tan\alpha/(1+\kappa)$，超过题设静摩擦上限，所以接触立即向下滑动。图\ref{deep1:fig:d1-p10-s}中取沿斜面向下与下行自转为正，动摩擦沿上坡，
\[
\boxed{a=g(\sin\alpha-\mu_k\cos\alpha)},\qquad
\boxed{\beta=\frac{\mu_kg\cos\alpha}{\kappa r}}.
\]

\begin{figure}[H]\centering
\begin{tikzpicture}
\begin{scope}[shift={(0,0)}]
\begin{scope}[rotate=-24]\draw[curve](-1.6,0)--(3.4,0);\draw[fill=white,curve](0,.7)circle(.7);\fill(0.00000,0.70000)circle(1.3pt)node[above left]{$C$};
\fill(0.00000,0.00000)circle(1.3pt)node[below left]{$P$};
\draw(0,.7)--(0.520201378,1.168391424)node[midway,above]{$r$};\draw[axis](0.00000,0.70000)--(1.00000,0.70000)node[above]{$a$};
\draw[axis](0.00000,0.00000)--(-1.05000,0.00000)node[above]{$f_k$};
\draw[axis](0.00000,0.00000)--(0.00000,1.85000)node[above left]{$N$};
\end{scope}\draw[guide](1.3060545559848429,-1.3829045864577207)--(3.306054555984843,-1.3829045864577207);\draw(2.456054555984843,-1.3829045864577207)arc(180:156:.65);\node at(2.126054555984843,-1.1529045864577208){$\alpha$};\draw[axis](0.2847157,0.6394818)--(0.2847157,-0.8105182)node[below right]{$mg$};

\end{scope}
\begin{scope}[shift={(5.4,0.2)}]
\draw[axis](-0.1,0)--(2.8,0)node[right]{$t$};\draw[axis](0,-0.1)--(0,2.2)node[above]{$u$};\draw[curve](0,0)--(2.3,1.8);\node at(1.2,2.2){$\dot u>0$};
\end{scope}

\end{tikzpicture}
\caption{动摩擦力 $f_k$ 沿上坡方向，右图中的相对滑动速度 $u$ 随时间增大。}\label{deep1:fig:d1-p10-s}
\end{figure}

加速度均为常量，使用静止初值分别积分，
\[
v(t)=\int_0^t a\,\dd\xi=[a\xi]_0^t=at,
\qquad \omega(t)=\int_0^t\beta\,\dd\xi=[\beta\xi]_0^t=\beta t.
\]
接触滑动速度为
\[
\boxed{u(t)=v-r\omega=
g\left[\sin\alpha-\left(1+\frac1\kappa\right)\mu_k\cos\alpha\right]t}.
\]
题设保证方括号为正，故滑动速度随时间增大，系统不会自行转入纯滚动。本题中自转虽在加快，质心平动增加得更快，接触处的相对滑动因而继续增大。

记上述正的滑动加速度为 $\gamma=\dot u=a-r\beta$。若运动持续至时刻 $t$，累计相对滑移为
\[
\int_0^t u(\xi)\dd\xi=\int_0^t\gamma\xi\dd\xi
=\gamma[\xi^2/2]_0^t=\frac12\gamma t^2.
\]
动摩擦力大小为 $\mu_kmg\cos\alpha$，乘以相对滑移，得到产热
\[
Q_{\text{热}}=\frac12\mu_kmg^2\cos\alpha
\left[\sin\alpha-\left(1+\frac1\kappa\right)\mu_k\cos\alpha\right]t^2.
\]
这一表达式为正，与重力势能损失扣除平动、自转动能后的结果相符。
从静止释放时 u 也为零，但静摩擦已不足以维持 $u=0$。应先假定向下滑动，按上坡动摩擦求解，再检验该方向是否自洽。题设给出
\[
\left(1+\frac1\kappa\right)\mu_k\cos\alpha
\le\left(1+\frac1\kappa\right)\mu_s\cos\alpha<\sin\alpha,
\]
所以求出的 $\dot u>0$，确实立即进入并始终保持向下滑动。

对任意 $t>0$，平动、自转的增长比为常量：
\[
\frac{r\omega}{v}
=\frac{\mu_k\cos\alpha}{\kappa(\sin\alpha-\mu_k\cos\alpha)}<1.
\]
故两种速度的差值始终为正，并按图\ref{deep1:fig:d1-p10-s}的直线增大。在当前坡度、摩擦条件及静止初态下，不会自行转为纯滚动；这一结论不适用于另加外力或改变初态的情况。


\end{solution}


\begin{exercise}[固定外筒的支持力与摩擦上限]\label{deep1:problem-exercise-11}\label{deep1:exercise:d1-p11}
薄壁大圆筒固定，小薄壁圆筒在其内侧纯滚动。记 $a=R-r$。初始小筒在最低处，其圆心速度为 $v_0=\sqrt{ga/2}$。求最高位置的摆角、内接触法向力 $N(\theta)$，以及全程保持纯滚动所需的最小静摩擦系数。
\end{exercise}

\begin{figure}[H]\centering
\begin{tikzpicture}
\draw[curve](0,0)circle(2.2);\draw[fill=white,curve](1.2639733311463137,-1.06059955598279)circle(0.55);\fill(0.00000,0.00000)circle(1.3pt)node[above left]{$O$};
\fill(1.26397,-1.06060)circle(1.3pt)node[below left]{$C$};
\draw[guide](0,0)--(0,-2.4000000000000004);\draw(0,0)--(1.2639733311463137,-1.06059955598279);\draw(0,-.65)arc(-90:-40.0:.65);\node at(.3,-.86){$\theta$};\draw(0,0)--(-1.6852977748617515,1.414132741310387)node[midway,above]{$R$};\draw(1.2639733311463137,-1.06059955598279)--(1.8139733311463138,-1.06059955598279)node[midway,above]{$r$};\fill(1.68530,-1.41413)circle(1.3pt)node[below right]{$P$};
\draw[guide](0,-1.6500000000000001)circle(0.55);\draw[axis](0.00000,-1.65000)--(1.20000,-1.65000)node[below]{$v_0$};
\node at(-1.7,-1.15){初态};
\end{tikzpicture}
\caption{外筒固定，圆心 C 沿半径 $a$ 的圆运动；接触点为 P。}\label{deep1:fig:d1-p11-q}
\end{figure}
\par\medskip
\begin{solution}
由例\ref{ex:hw1-8}固定外筒的能量关系，
\[
v_C^2=v_0^2-ga(1-\cos\theta)
=ga(\cos\theta-\tfrac12).
\]
故最高位置满足 $\cos\theta_m=1/2$，即
\[
\boxed{\theta_m=\pi/3}.
\]

\begin{figure}[H]\centering
\begin{tikzpicture}
\draw[curve](0,0)circle(2.2);\draw[fill=white,curve](1.4289419162443238,-0.8250000000000003)circle(0.55);\fill(0.00000,0.00000)circle(1.3pt)node[above left]{$O$};
\fill(1.42894,-0.82500)circle(1.3pt)node[below left]{$C$};
\draw[guide](0,0)--(0,-2.4000000000000004);\draw(0,0)--(1.4289419162443238,-0.8250000000000003);\draw(0,-.65)arc(-90:-30.000000000000007:.65);\node at(.3,-.86){$\theta$};\draw(0,0)--(-1.6852977748617515,1.414132741310387)node[midway,above]{$R$};\draw(1.4289419162443238,-0.8250000000000003)--(1.9789419162443238,-0.8250000000000003)node[midway,above]{$r$};\fill(1.90526,-1.10000)circle(1.3pt)node[below right]{$P$};
\draw[guide](0,-1.6500000000000001)circle(0.55);\draw[axis](0.00000,-1.65000)--(1.20000,-1.65000)node[below]{$v_0$};
\node at(-1.7,-1.15){初态};\draw[axis](1.90526,-1.10000)--(0.69282,-0.40000)node[above right]{$N$};
\draw[axis](1.90526,-1.10000)--++(.4,.69282)node[above right]{$f_s$};
\begin{scope}[shift={(5.4,-1.2)}]
\draw[axis](-0.1,0)--(2.7,0)node[right]{$\theta$};\draw[axis](0,-0.1)--(0,2.1)node[above]{$N/(mg)$};\draw[curve]plot coordinates{(0.000000,1.500000)(0.020944,1.499890)(0.041888,1.499561)(0.062832,1.499013)(0.083776,1.498246)(0.104720,1.497259)(0.125664,1.496053)(0.146608,1.494629)(0.167552,1.492986)(0.188496,1.491124)(0.209440,1.489044)(0.230383,1.486746)(0.251327,1.484229)(0.272271,1.481496)(0.293215,1.478545)(0.314159,1.475377)(0.335103,1.471992)(0.356047,1.468391)(0.376991,1.464575)(0.397935,1.460542)(0.418879,1.456295)(0.439823,1.451834)(0.460767,1.447158)(0.481711,1.442269)(0.502655,1.437166)(0.523599,1.431852)(0.544543,1.426325)(0.565487,1.420587)(0.586431,1.414639)(0.607375,1.408481)(0.628319,1.402113)(0.649262,1.395537)(0.670206,1.388753)(0.691150,1.381762)(0.712094,1.374564)(0.733038,1.367161)(0.753982,1.359553)(0.774926,1.351741)(0.795870,1.343726)(0.816814,1.335509)(0.837758,1.327091)(0.858702,1.318472)(0.879646,1.309654)(0.900590,1.300638)(0.921534,1.291424)(0.942478,1.282013)(0.963422,1.272407)(0.984366,1.262607)(1.005310,1.252613)(1.026254,1.242428)(1.047198,1.232051)(1.068142,1.221484)(1.089085,1.210729)(1.110029,1.199785)(1.130973,1.188656)(1.151917,1.177341)(1.172861,1.165842)(1.193805,1.154161)(1.214749,1.142298)(1.235693,1.130256)(1.256637,1.118034)(1.277581,1.105635)(1.298525,1.093060)(1.319469,1.080310)(1.340413,1.067387)(1.361357,1.054292)(1.382301,1.041026)(1.403245,1.027592)(1.424189,1.013990)(1.445133,1.000222)(1.466077,0.986290)(1.487021,0.972194)(1.507964,0.957937)(1.528908,0.943520)(1.549852,0.928945)(1.570796,0.914214)(1.591740,0.899327)(1.612684,0.884286)(1.633628,0.869094)(1.654572,0.853752)(1.675516,0.838261)(1.696460,0.822624)(1.717404,0.806841)(1.738348,0.790915)(1.759292,0.774848)(1.780236,0.758641)(1.801180,0.742296)(1.822124,0.725814)(1.843068,0.709198)(1.864012,0.692450)(1.884956,0.675571)(1.905900,0.658562)(1.926843,0.641427)(1.947787,0.624167)(1.968731,0.606783)(1.989675,0.589278)(2.010619,0.571654)(2.031563,0.553912)(2.052507,0.536054)(2.073451,0.518083)(2.094395,0.500000)};
\draw[guide](2.0944,0)--(2.0944,.5)--(0,.5);\node[left]at(0,.5){$1/2$};\node[left]at(0,1.5){$3/2$};\node[below]at(2.0944,0){$\pi/3$};\node at(1.6,1.9){$N>0$};
\end{scope}

\end{tikzpicture}
\caption{左图画出最高位置的法向力和静摩擦，后者沿角度增大的切向；右图为法向力随角度的变化。}\label{deep1:fig:d1-p11-s}
\end{figure}

图\ref{deep1:fig:d1-p11-s}中的法向力沿 C 到 O 方向。法向方程为
\[
N-mg\cos\theta=m\frac{v_C^2}{a},
\quad\boxed{N=mg(2\cos\theta-\tfrac12)}.
\]
在 $0\le\theta\le\pi/3$ 内，最小法向力为最高处的 $mg/2$，因此不会失去接触。初始平动与自转动能各为 $mv_0^2/2$，总能量为 $mv_0^2=mga/2$。最高处的速度虽为零，接触压力仍为正，小筒将在轨道上反向运动。

在习题\ref{deep1:exercise:d1-p5}已推导的滚动方程中取 $\kappa=1$、$\rho=a$，得 $\ddot\theta=-g\sin\theta/(2a)$。绕小筒圆心的转动方程则给出切向静摩擦大小
\[
|f_s|=\frac12mg\sin\theta.
\]
于是
\[
\frac{|f_s|}{N}=\frac{\sin\theta}{4\cos\theta-1}.
\]
按商法则求导，
\[
\frac{\dd}{\dd\theta}\frac{\sin\theta}{4\cos\theta-1}
=\frac{\cos\theta(4\cos\theta-1)+4\sin^2\theta}{(4\cos\theta-1)^2}
=\frac{4-\cos\theta}{(4\cos\theta-1)^2}>0.
\]
最大值出现在最高处，故
\[
\boxed{\mu_{s,\min}=\frac{\sqrt3}{2}}.
\]
小筒在最高处瞬时静止，但重力切向分量不为零，所以维持随后的滚动反向仍需非零静摩擦。
摩擦方向可用切向方程明确确定。沿 $\theta$ 增大方向取正，
\[
f_s-mg\sin\theta=ma\ddot\theta=-\frac12mg\sin\theta,
\qquad f_s=\frac12mg\sin\theta.
\]
所以初次上升时摩擦沿切向向上，虽与圆心速度同向，却使自转减慢；到转向点，摩擦方向仍不立即改变。摩擦方向取决于滚动约束所需的角加速度。

图\ref{deep1:fig:d1-p11-friction}按 $\sin\theta/(4\cos\theta-1)$ 绘出所需比值。端点 $\theta=\pi/3$ 处，$f_s=\sqrt3mg/4$、$N=mg/2$，两者的比正是 $\sqrt3/2$。最大摩擦需求恰发生在转向点。
\begin{figure}[H]\centering
\begin{tikzpicture}[scale=.95]
\draw[axis](-.1,0)--(3.6,0)node[right]{$\theta$};\draw[axis](0,-.1)--(0,1.9)node[above]{$|f_s|/N$};
\draw[curve,domain=0:60,samples=110]plot({3*\x/60},{1.6*sin(\x)/(4*cos(\x)-1)});
\fill(3,{.8*sqrt(3)})circle(1.4pt);\draw[guide](0,{.8*sqrt(3)})--(3,{.8*sqrt(3)})--(3,0);
\node[left]at(0,{.8*sqrt(3)}){$\sqrt3/2$};\node[below]at(1.5,0){$\pi/6$};\node[below]at(3,0){$\pi/3$};
\node[anchor=west]at(4.1,1.1){$\mu_s\ge\max |f_s|/N$};
\end{tikzpicture}
\caption{全程纯滚动所需的摩擦比值。最高处为本题的控制位置，曲线只取保持接触的上升分支。}\label{deep1:fig:d1-p11-friction}
\end{figure}

\end{solution}


\begin{exercise}[任意转动惯量的内滚系统（选做）]\label{deep1:problem-exercise-12}\label{deep1:exercise:d1-p12}
将例\ref{ex:hw1-8}的两个圆筒推广为轴对称滚动物体。外物体质量 $M$、外接触半径 $R$、轴向转动惯量 $I_o=\kappa_oMR^2$；内物体质量 $m$、接触半径 $r$、转动惯量 $I_i=\kappa_imr^2$。外物体的内侧圆轨道半径也为 $R$，可用不妨碍内滚运动的轴向配重改变其转动惯量。外物体在水平地面纯滚动。初始外物体静止，内物体位于最低处，圆心向右速率为 $v_0$。设各接触在上升阶段均保持，求内物体圆心的最大上升高度。
\end{exercise}

\begin{figure}[H]\centering
\begin{tikzpicture}
\draw[curve](0,0)circle(2.2);\draw[fill=white,curve](1.4289419162443238,-0.8250000000000003)circle(0.55);\fill(0.00000,0.00000)circle(1.3pt)node[above left]{$O$};
\fill(1.42894,-0.82500)circle(1.3pt)node[below left]{$C$};
\draw[guide](0,0)--(0,-2.4000000000000004);\draw(0,0)--(1.4289419162443238,-0.8250000000000003);\draw(0,-.65)arc(-90:-30.000000000000007:.65);\node at(.3,-.86){$\theta$};\draw(0,0)--(-1.6852977748617515,1.414132741310387)node[midway,above]{$R$};\draw(1.4289419162443238,-0.8250000000000003)--(1.9789419162443238,-0.8250000000000003)node[midway,above]{$r$};\fill(1.90526,-1.10000)circle(1.3pt)node[below right]{$P$};
\draw[curve](-2.65,-2.2)--(3.1,-2.2);\draw[axis](-1.70000,-2.75000)--(-0.60000,-2.75000)node[above]{$X$};
\draw[guide](0,-1.6500000000000001)circle(0.55);\draw[axis](0.00000,-1.65000)--(1.20000,-1.65000)node[below]{$v_0$};
\node at(-1.7,-1.15){初态};
\end{tikzpicture}
\caption{内外接触半径决定速度关联，质量分布由两个惯量系数表示。}\label{deep1:fig:d1-p12-q}
\end{figure}
\par\medskip
\begin{solution}
仍取 $a=R-r$。图\ref{deep1:fig:d1-p12-s}中两个自转角速度的约束与质量分布无关，
\[
\Omega=-\frac{\dot X}{R},\qquad
\omega=-\frac{\dot X+a\dot\theta}{r}.
\]
因此动能为
\[
T=\frac12D\dot X^2+ma(\cos\theta+\kappa_i)\dot X\dot\theta
+\frac12ma^2(1+\kappa_i)\dot\theta^2,
\]
其中 $D=(1+\kappa_o)M+(1+\kappa_i)m$。势能为 $mga(1-\cos\theta)$。

\begin{figure}[H]\centering
\begin{tikzpicture}
\draw[curve](0,0)circle(2.2);\draw[fill=white,curve](1.4289419162443238,-0.8250000000000003)circle(0.55);\fill(0.00000,0.00000)circle(1.3pt)node[above left]{$O$};
\fill(1.42894,-0.82500)circle(1.3pt)node[below left]{$C$};
\draw[guide](0,0)--(0,-2.4000000000000004);\draw(0,0)--(1.4289419162443238,-0.8250000000000003);\draw(0,-.65)arc(-90:-30.000000000000007:.65);\node at(.3,-.86){$\theta$};\draw(0,0)--(-1.6852977748617515,1.414132741310387)node[midway,above]{$R$};\draw(1.4289419162443238,-0.8250000000000003)--(1.9789419162443238,-0.8250000000000003)node[midway,above]{$r$};\fill(1.90526,-1.10000)circle(1.3pt)node[below right]{$P$};
\draw[curve](-2.65,-2.2)--(3.1,-2.2);\draw[axis](-1.70000,-2.75000)--(-0.60000,-2.75000)node[above]{$X$};
\draw[axis](1.90526,-1.10000)--(0.69282,-0.40000)node[above right]{$N$};
\node at(0,2.65){相对转折：$\dot\theta=0$};\draw[axis](0.00000,0.00000)--(1.00000,0.00000)node[above]{$V_*$};
\draw[axis](1.42894,-0.82500)--(2.42894,-0.82500)node[below right]{$V_*$};
\begin{scope}[shift={(5.7,0)}]
\node at(0,1.25){$E=\tfrac12m(1+\kappa_i)v_0^2$};\draw[axis](0,.8)--(0,.15);\node at(0,-.25){$\tfrac12D V_*^2+mgh$};
\end{scope}

\end{tikzpicture}
\caption{若内物体在圆轨道最高点之前发生相对转向，此时两圆心速度相同，能量分为共同运动的动能与重力势能。}\label{deep1:fig:d1-p12-s}
\end{figure}

循环动量及初态能量分别为
\[
P=D\dot X+ma(\cos\theta+\kappa_i)\dot\theta
=m(1+\kappa_i)v_0,
\]
\[
E=\frac12m(1+\kappa_i)v_0^2.
\]
几何高度为 $h=a(1-\cos\theta)$，最多为 $2a$。因此必须先判断内物体在上升途中转向，还是具有越过轨道最高点的能量。令
\[
C(\theta)=ma(\cos\theta+\kappa_i),\qquad
B=ma^2(1+\kappa_i).
\]
动能配方为
\[
T=\frac D2\left(\dot X+\frac CD\dot\theta\right)^2
+\frac12\left(B-\frac{C^2}{D}\right)\dot\theta^2.
\]
第一个平方项等于常量 $P^2/(2D)$，剩余的相对运动能量为
\[
E_{\rm rel}=E-\frac{P^2}{2D}
=\frac{m(1+\kappa_i)(1+\kappa_o)M}{2D}v_0^2.
\]
记 $K(\theta)=B-C^2/D$，约化能量式为
\[
\frac12K(\theta)\dot\theta^2+mga(1-\cos\theta)=E_{\rm rel}.
\]
因为 $|\cos\theta+\kappa_i|\le1+\kappa_i$ 且 $D>m(1+\kappa_i)$，有 $K(\theta)>0$。故能否到达某一位置，只需比较 $E_{\rm rel}$ 与该处的势能。

当 $0<E_{\rm rel}<2mga$ 时，在 $0<\theta_*<\pi$ 处有 $\dot\theta=0$，因此
\[
\cos\theta_*=1-\frac{E_{\rm rel}}{mga},\qquad
\boxed{h_{\max}=\frac{E_{\rm rel}}{mg}
=\frac{(1+\kappa_i)(1+\kappa_o)M}{2D}\frac{v_0^2}{g}}.
\]
这正是图\ref{deep1:fig:d1-p12-s}所示的转向情况。此时两圆心速度均为 $V_*=P/D$，但两物体仍在自转：$\Omega_*=-V_*/R$、$\omega_*=-V_*/r$。平动和自转动能之和为 $DV_*^2/2$，不能将最高处的全部动能置零。

当 $E_{\rm rel}=2mga$ 时，内物体渐近接近 $\theta=\pi$。令 $\delta=\pi-\theta$，在最高点附近有
\[
E_{\rm rel}-mga(1-\cos\theta)
=\frac12mga\delta^2+O(\delta^4),\qquad
\dot\delta\simeq-\sqrt{\frac{mga}{K(\pi)}}\,\delta.
\]
记 $\lambda=\sqrt{mga/K(\pi)}$，分离变量，从接近最高点的时刻 $t_1$ 积分至 $t$，
\[
\int_{\delta(t_1)}^{\delta(t)}\frac{\dd\xi}{\xi}
=[\ln\xi]_{\delta(t_1)}^{\delta(t)}
\simeq-\lambda(t-t_1).
\]
所以 $\delta(t)\simeq\delta(t_1)\ee^{-\lambda(t-t_1)}$，在有限时间内不等于零；高度的上确界为 $2a$。

当 $E_{\rm rel}>2mga$ 时，内物体越过轨道最高点。在该点
\[
h_{\max}=2a,\qquad
\dot\theta_{\rm top}^{\,2}
=\frac{2(E_{\rm rel}-2mga)}{K(\pi)}>0.
\]
最高位置由 $\dd h/\dd\theta=a\sin\theta=0$ 决定，此时相对角速度并不为零。三种情况可合写为
\[
\boxed{\sup h(t)=\min\!\left\{
\frac{(1+\kappa_i)(1+\kappa_o)M}{2D}\frac{v_0^2}{g},\ 2(R-r)\right\}}.
\]
除临界能量外，式中的上确界都能实际达到；$v_0=0$ 时系统保持初始状态。

\begin{figure}[H]\centering
\begin{tikzpicture}[x=2.4cm,y=1.8cm]
\draw[axis](0,0)--(2.35,0)node[right]{$\eta$};
\draw[axis](0,0)--(0,1.4)node[above]{$\sup h/(2a)$};
\draw[curve](0,0)--(1,1)--(2.1,1);
\draw[guide](1,0)--(1,1)--(0,1);
\node[below]at(1,0){$1$};\node[left]at(0,1){$1$};
\node[above,rotate=37]at(.48,.55){上升后转向};
\node[above]at(1.7,1){越过最高点};
\fill(1,1)circle(1.2pt);
\draw[guide](1,1)--(1.15,.55);
\node[below right,align=left]at(1.05,.55){临界：仅渐近到达};
\end{tikzpicture}
\caption{以 $\eta=E_{\rm rel}/(2mga)$ 为横坐标。转向高度随能量增加，达到轨道几何上限后保持为 $2a$。}\label{deep1:fig:d1-p12-height}
\end{figure}

图\ref{deep1:fig:d1-p12-height}中的折点给出两类运动的分界。取 $\kappa_o=\kappa_i=1$，在 $E_{\rm rel}<2mga$ 的分支恢复两个薄壁圆筒的转向高度；令 $M\to\infty$，该分支趋于 $h=(1+\kappa_i)v_0^2/(2g)$。半径在转向高度公式中消去，但仍通过 $2(R-r)$ 限制几何高度。上述各分支均以接触不脱离、各处维持纯滚动为前提；若实际摩擦或支持力不足，应在约束先失效的位置转入新的运动模型。

\end{solution}


\chapter{约束临界问题}


求约束临界状态，既要算出物体怎样运动，也要算出维持该运动所需的力。轨道只能推小球，轻绳只能拉质点，静摩擦受接触压力限制；接触合力还必须落在支承区域内。达到哪一条边界，决定了原运动能否继续。杆内强度问题则要进一步求出各截面传递的力与弯矩。

除另有说明外，以地面为惯性参考系，轻杆、轻环的质量忽略不计，固定光滑接触不耗散能量。图中力箭头标出作用点与方向，长度通常不代表力的比例；结果曲线按相应函数绘制。



\section{单侧约束与临界分离}

绳只能拉，接触面只能沿允许方向施力。按现有约束求出的反力达到零时，还要判断后续运动趋势，才能确定是否松绳或分离。



\subsection{单侧约束与临界分离}
\subsubsection{约束力的允许方向}
刚性光滑轨道对轨道外侧的小球只能施加向外的支持力，轻绳只能沿绳方向拉质点。计算时先保留原有几何约束，求出维持它所需的力，再与允许范围比较。

以接触间隙 $h$ 向允许脱离的一侧为正，支持力大小为 $N$，单侧接触满足
\[
h\ge0,\qquad N\ge0,\qquad Nh=0.
\]
在保持接触的一段运动中 $h=0$；分离后 $N=0$，轨迹由新的运动方程决定。若还存在静摩擦，则必须同时满足 $|f|\le\mu N$。闭合套环、铰链和可承压的刚杆通常属于双侧约束，反力可以改变方向，不能套用单侧接触的非负条件。

对轻绳，记绳长为 $L$、两端距离为 $r$，相应条件为
\[
r\le L,\qquad T\ge0,\qquad T(L-r)=0.
\]
绳松弛后只受实际外力；再次绷紧时若沿绳有非零的向外相对速度，还会出现约束冲量，不能直接延续原来的能量守恒式。

\subsubsection{零反力是否意味着约束改变}
解出 $N=0$ 或 $T=0$ 后，还要沿实际运动方向判断反力的符号：此前须处于允许范围，之后若要求负反力，原约束才会失效。反力若触零后回升，约束可能继续保持。判断时还应比较速度转向、摩擦达到极限和接触合力到达支承边缘的先后。
图\ref{sup2:fig:d2-zero-test}画出两种反力零点。
\begin{figure}[H]\centering
\begin{tikzpicture}[scale=.95]
\draw[axis](-.1,0)--(3.3,0)node[right]{$t$};\draw[axis](0,-1.1)--(0,1.5)node[above]{$N$};
\draw[curve](0,1)--(1.5,0);\draw[guide](1.5,0)--(3,-1);\node[above right]at(1.5,0){$t_s$};
\begin{scope}[shift={(5,0)}]
\draw[axis](-.1,0)--(3.3,0)node[right]{$t$};\draw[axis](0,-.2)--(0,1.5)node[above]{$N$};
\draw[curve,domain=0:3,samples=80]plot(\x,{.5*(\x-1.5)^2});\node[below]at(1.5,0){$t_s$};
\end{scope}
\end{tikzpicture}
\caption{原约束延拓所需反力的两种零点。左图会进入不允许的负反力区；右图只触零后回升，须结合运动判断。}\label{sup2:fig:d2-zero-test}
\end{figure}

若切向运动由能量给出 $v^2(s)$，普通运动点满足 $a_t=\tfrac12\dd(v^2)/\dd s$。法向运动则满足
\[
\sum F_n=m\frac{v^2}{\rho}.
\]
这里主法线指向轨迹的曲率中心，支持力未必沿主法线正向。圆外侧运动中，重力的向心分量与向外支持力相减；圆内侧运动中，支持力可以直接提供向心力。

下文所说的从竖直位置或凸轨道顶端静止倒下，均指从极小非零偏角释放，再取偏角趋零的极限。严格处在光滑不稳定平衡点且初速度为零时，理想运动方程的解仍是静止。


\subsection{例题与模型变式}

\begin{example}[圆弧与端点质量杆的分离]\label{ex:hw2-1}

(1)质量$m$的小球由半径$r$的光滑圆弧最高点附近从静止沿外侧下滑，求离开圆弧的位置。(2)长$2l$的轻杆两端各有质量$m$的小球，分别接触光滑竖直墙和光滑水平地面，杆从近乎竖直状态由静止释放，求杆上端离墙的位置。图\ref{fig:hw2-1-q}中两题的角度均从竖直方向量起。

\end{example}


\begin{figure}[H]\centering
\begin{tikzpicture}[>=Stealth,line cap=round,line join=round]
\begin{scope}[shift={(0,0)}]\draw[curve,domain=0:90,samples=50]plot({2*sin(\x)},{2*cos(\x)});\draw[guide](0,0)--(0,2);\coordinate(P)at(1.41421356,1.41421356);\draw(0,0)--(P);\fill(P)circle(2pt)node[right]{$P$};\node[below]at(0,0){$O$};\node[left]at(0,2){$A$};\draw(0,.6)arc(90:45:.6)node[midway,above]{$\theta$};\end{scope}\begin{scope}[shift={(4,0)}]\draw[thick](0,2.7)--(0,0)--(3,0);\draw[curve](0,2.25)--(1.7,0);\fill(0,2.25)circle(2pt)node[left]{$A$};\fill(1.7,0)circle(2pt)node[below]{$B$};\node[right]at(.9,1.3){$2l$};\draw[guide](1.7,0)--(1.7,2.5);\draw(1.7,.55)arc(90:127.073:.55)node[midway,above left]{$\theta$};\end{scope}
\end{tikzpicture}
\caption{圆弧外侧下滑与端点质量杆下落}\label{fig:hw2-1-q}
\end{figure}

\begin{solution}
圆弧上由能量守恒有$v^2=2gr(1-\cos\theta)$。图\ref{fig:hw2-1-s}左图向心方向的方程为$mg\cos\theta-N=mv^2/r$，故
\[
 N=mg(3\cos\theta-2),\qquad
 \boxed{\cos\theta_c=2/3}.
\]
离开前$N>0$，继续沿圆弧运动会要求$N<0$，所以该零点确为分离点。

\begin{figure}[H]\centering
\begin{tikzpicture}[>=Stealth,line cap=round,line join=round]
\begin{scope}[shift={(0,0)}]\draw[curve,domain=0:90,samples=50]plot({2*sin(\x)},{2*cos(\x)});\draw[guide](0,0)--(0,2);\coordinate(P)at(1.41421356,1.41421356);\draw(0,0)--(P);\fill(P)circle(2pt)node[right]{$P$};\node[below]at(0,0){$O$};\node[left]at(0,2){$A$};\draw(0,.6)arc(90:45:.6)node[midway,above]{$\theta$};\draw[axis](P)--(.75,.75);\draw[axis](P)--(1.838478,1.838478)node[right]{$N$};\node[right]at(1.2,.55){$mg\cos\theta$};\end{scope}\begin{scope}[shift={(4,0)}]\draw[thick](0,2.7)--(0,0)--(3,0);\draw[curve](0,2.25)--(1.7,0);\fill(0,2.25)circle(2pt)node[left]{$A$};\fill(1.7,0)circle(2pt)node[below]{$B$};\node[right]at(.9,1.3){$2l$};\draw[guide](1.7,0)--(1.7,2.5);\draw(1.7,.55)arc(90:127.073:.55)node[midway,above left]{$\theta$};\draw[axis](0,2.25)--(.9,2.25)node[above]{$N_w$};\draw[axis](1.7,0)--(1.7,.85)node[right]{$N_B$};\end{scope}
\end{tikzpicture}
\caption{径向重力分量与墙的水平反力}\label{fig:hw2-1-s}
\end{figure}

杆题中$x_B=2l\sin\theta$、$y_A=2l\cos\theta$，两端速度分别为 $2l\cos\theta\,\dot\theta$ 和 $-2l\sin\theta\,\dot\theta$，故两端动能和为 $2ml^2\dot\theta^2$。势能的下降量为 $2mgl(1-\cos\theta)$，绕角度坐标的运动方程为 $4ml^2\ddot\theta=2mgl\sin\theta$。因此
\[
 \dot\theta^2=\frac gl(1-\cos\theta),\quad
 \ddot\theta=\frac g{2l}\sin\theta.
\]
对两端坐标继续求导，
\[
\ddot x_B=2l(\cos\theta\,\ddot\theta-\sin\theta\,\dot\theta^2),
\qquad
\ddot y_A=-2l(\sin\theta\,\ddot\theta+\cos\theta\,\dot\theta^2).
\]
系统水平外力只有墙的反力$N_w$；上端水平速度为零，故
\[
 N_w=m\ddot x_B
 =mg\sin\theta(3\cos\theta-2).
\]
地面对下端的支持力可由竖直动量求得：$N_B-2mg=m\ddot y_A$，于是
\[
 N_B=mg(1-2c+3c^2)>0,\qquad c=\cos\theta.
\]
在墙反力首次由正变负的位置，地面接触仍可保持。因此先失效的是墙面约束。

杆上端离墙时同样有$\boxed{\cos\theta_c=2/3}$。这里判定分离的是墙的水平反力；圆弧题采用的向心力方程不适用于该杆系。


\end{solution}



\begin{example}[抛物线与椭圆外轨道]\label{ex:hw2-2}

质量$m$的质点分别在光滑轨道$y=-x^2/b$（$b>0$）及$x^2/a^2+y^2/b^2=1$的上部外侧，从最高点附近由静止下滑。求是否分离；若分离，确定分离点。图\ref{fig:hw2-2-q}中椭圆半轴为$a,b>0$。

\end{example}


\begin{figure}[H]\centering
\begin{tikzpicture}[>=Stealth,line cap=round,line join=round]
\begin{scope}[shift={(0,0)}]\draw[axis](0,0)--(2.5,0)node[right]{$x$};\draw[axis](0,0)--(0,1.2)node[above]{$y$};\draw[curve,domain=-1.7:1.7,samples=60]plot(\x,{-(\x)^2/2});\node[above left]at(0,0){$A$};\node at(0,-2){$y=-x^2/b$};\end{scope}\begin{scope}[shift={(6,0)}]\draw(0,0)ellipse(2 and 1.35);\coordinate(P)at(1.3,1.026);\draw[guide](0,0)--(P);\fill(P)circle(2pt)node[right]{$P$};\draw[axis](-2.3,0)--(2.4,0)node[right]{$x$};\draw[axis](0,-1.6)--(0,1.7)node[above]{$y$};\node[left]at(0,1.35){$A$};\end{scope}
\end{tikzpicture}
\caption{抛物线外轨道与椭圆外轨道}\label{fig:hw2-2-q}
\end{figure}

\begin{solution}
在抛物线上，$y'=-2x/b$、$y''=-2/b$，因此曲率为
\[
\frac1{\rho_c}=\frac{|y''|}{(1+y'^2)^{3/2}}
=\frac{2/b}{(1+4x^2/b^2)^{3/2}}.
\]
能量守恒给出 $mv^2/2=mgx^2/b$。以向下的主法向列式，得到图\ref{fig:hw2-2-s}左侧的支持力
\[
 \boxed{N=\frac{mg}{(1+4x^2/b^2)^{3/2}}>0}.
\]
因此不存在有限位置的分离。严格静止在不稳定的最高点会保持静止，题意采用从其邻域开始运动的极限。

\begin{figure}[H]\centering
\begin{tikzpicture}[>=Stealth,line cap=round,line join=round]
\begin{scope}[shift={(0,0)}]\draw[axis](0,0)--(2.5,0)node[right]{$x$};\draw[axis](0,0)--(0,1.6)node[above]{$N$};\draw[curve,domain=0:2.3,samples=60]plot(\x,{1.4/(1+\x*\x)^(1.5)});\node at(1,-.5){$N>0$};\end{scope}\begin{scope}[shift={(6,0)}]\draw(0,0)ellipse(2 and 1.35);\coordinate(P)at(1.3,1.026);\draw[axis](P)--(1.65,1.632)node[right]{$N$};\draw[guide](0,0)--(P);\fill(P)circle(2pt)node[right]{$P$};\draw[axis](-2.3,0)--(2.4,0)node[right]{$x$};\draw[axis](0,-1.6)--(0,1.7)node[above]{$y$};\node[left]at(0,1.35){$A$};\node at(0,-2){$D=a^2c^2+b^2(1-c^2)$};\end{scope}
\end{tikzpicture}
\caption{抛物线支持力保持正值，椭圆须另求分离点}\label{fig:hw2-2-s}
\end{figure}

抛物线单位主法向的竖直分量大小为$1/\sqrt{1+4x^2/b^2}$，因此
\[
 N=\frac{mg}{\sqrt{1+4x^2/b^2}}-\frac{mv^2}{\rho_c}.
\]
代入 $v^2=2gx^2/b$ 并通分，
\[
N=\frac{mg(1+4x^2/b^2)-4mgx^2/b^2}{(1+4x^2/b^2)^{3/2}},
\]
分子中的 $x^2$ 项抵消，支持力始终为正。椭圆的参数 $\theta$ 采用下文的坐标定义，一般不等于位置矢量与坐标轴的极角。

椭圆取 $x=a\sin\theta$、$y=b\cos\theta$，记 $c=\cos\theta$、$D=a^2c^2+b^2(1-c^2)$。参数的一阶与二阶导数为
\[
x'=a\cos\theta,\quad y'=-b\sin\theta,\quad
x''=-a\sin\theta,\quad y''=-b\cos\theta.
\]
外单位法线为 $(b\sin\theta,a\cos\theta)/\sqrt D$；曲率分子 $|x'y''-y'x''|=ab(\cos^2\theta+\sin^2\theta)=ab$。能量给出 $v^2=2gb(1-c)$，故法向方程为
\[
N=\frac{mgac}{\sqrt D}-\frac{2mgb(1-c)ab}{D^{3/2}}.
\]
用 $D=b^2+(a^2-b^2)c^2$ 合并分子，得到
\[
 N=\frac{mga}{D^{3/2}}\big[(a^2-b^2)c^3+3b^2c-2b^2\big].
\]
所以分离点由
\[
 \boxed{(a^2-b^2)c_*^3+3b^2c_*-2b^2=0,\quad
 (x_*,y_*)=(a\sqrt{1-c_*^2},bc_*)}
\]
确定。该多项式在$c=0$时为负、$c=1$时为$a^2>0$，导数$3[b^2+(a^2-b^2)c^2]>0$，故$(0,1)$内恰有一根。$a=b$时退化为$c_*=2/3$。


\end{solution}



\begin{example}[轻绳竖直圆运动的三种结局]\label{sup2:problem-example-3}\label{sup2:example:d2-e3}

长为 $R$ 的轻绳一端固定，另一端系质量为 $m$ 的质点。质点在最低点 B 获得水平速率 $u$，以后在竖直面内运动。角度从向下半径起算，见图\ref{sup2:fig:d2-e3-q}。按 $k=u^2/(gR)$ 分类讨论转向、松绳与完整周转。若把绳换成可拉可压的轻杆，越过顶点的条件有何改变？

\end{example}
\begin{figure}[H]\centering
\begin{tikzpicture}[scale=.95]
\coordinate(O)at(0,1.8);\coordinate(B)at(0,.2);\coordinate(P)at({1.6*sin(60)},{1.8-1.6*cos(60)});
\draw[guide](O)circle(1.6);\draw[thick](O)--(P);\draw[guide](O)--(B);
\fill(O)circle(1.3pt)node[above left]{$O$};\fill(B)circle(1.3pt)node[below]{$B$};\fill(P)circle(1.7pt)node[right]{$P$};
\draw(0,1.25)arc(-90:-30:.55);\node at(.4,1){$\theta$};
\draw[axis](0,.2)--(.9,.2)node[right]{$u$};\node at(2.8,2.5){$k=u^2/(gR)$};
\node at(2.8,1.5){$OP=R$};
\end{tikzpicture}
\caption{最低点发射的绳摆。角度沿右半圆向上增加，虚线表示绳保持绷紧时的圆轨道。}\label{sup2:fig:d2-e3-q}
\end{figure}
\par\medskip

\begin{solution}
在绳仍绷紧时，质点上升高度为 $R(1-\cos\theta)$。能量与沿绳向内的牛顿方程给出
\[
v^2=gR(k-2+2\cos\theta),\qquad
T-mg\cos\theta=\frac{mv^2}{R},
\]
因此
\[
T=mg(k-2+3\cos\theta).
\]
必须同时满足 $v^2\ge0$、$T\ge0$。前者的零点表示可能转向，后者的零点表示可能松绳；应比较沿上升过程先遇到哪一个。图\ref{sup2:fig:d2-e3-s}只画各条绷紧运动实际能够到达的部分。

\begin{figure}[H]\centering
\begin{tikzpicture}[scale=.95]
\draw[axis](-.15,0)--(3.7,0)node[right]{$\theta$};\draw[axis](0,-.2)--(0,2.65)node[above]{$T/(mg)$};
\draw[curve,domain=0:60,samples=60]plot({\x/55},{.35*(-1+3*cos(\x))});
\draw[curve,domain=0:109.47122,samples=70]plot({\x/55},{.35*(1+3*cos(\x))});
\draw[curve,domain=0:180,samples=100]plot({\x/55},{.35*(3+3*cos(\x))});
\node[left]at(0,.7){$k=1$};\node[left]at(0,1.4){$k=3$};\node[left]at(0,2.1){$k=5$};
\fill({60/55},.175)circle(1.4pt);\draw[guide]({60/55},.175)--(.6,-.25);\node at(.55,-.6){转向};
\fill({109.47122/55},0)circle(1.4pt)node[below]{松绳};
\fill({180/55},0)circle(1.4pt)node[below]{$\pi$};
\node at(3,1.4){$k=5$ 在顶点触零};
\end{tikzpicture}
\caption{张力曲线与运动终点。速度为零和张力为零须按各自的物理意义处理。}\label{sup2:fig:d2-e3-s}
\end{figure}

当 $0<k<2$，转向点满足 $\cos\theta_t=1-k/2>0$。在该点
\[
T_t=mg\cos\theta_t>0.
\]
所以质点先在水平半径以下转向，绳保持绷紧。$k=0$ 则始终停在最低点。

当 $k=2$，质点到水平半径处同时有 $v=0,T=0$。随后质点在重力作用下沿圆弧返回，张力 $T$ 重新为正。若假设自由落体，短时间后到 O 的距离反而超过 R，不符合绳长限制，故这一瞬时零张力并不导致松绳。

当 $2<k<5$，先到达
\[
\cos\theta_s=\frac{2-k}{3},\qquad
v_s^2=\frac{k-2}{3}gR>0.
\]
这一点位于上半圆，继续上行会要求负张力，因而从此进入松绳抛体运动。到达速度零点所需的余弦为 $(2-k)/2$，比上述余弦更小，证实松绳先于转向发生。

当 $k\ge5$，整周中 T 在顶点最小，$T_{\rm top}=mg(k-5)$。$k=5$ 时顶点速度仍为 $\sqrt{gR}>0$，T 只触及零后回升，可以完成周转；$k>5$ 时整周张力为正。因此完整绷紧周转的条件为 $u^2\ge5gR$。

轻杆允许轴向力为负，只有能量条件限制能否上到顶端。杆摆在顶点的速率平方为 $gR(k-4)$。$k>4$ 时以非零速度越过；$k=4$ 时，令 $\delta=\pi-\theta$，有 $v\sim\sqrt{gR}\,\delta$，因而 $\dot\delta\sim-\sqrt{g/R}\,\delta$。在顶点附近，
\[
t-t_1\sim\sqrt{\frac Rg}\int_{\delta(t)}^{\delta(t_1)}\frac{\dd\xi}{\xi}
=\sqrt{\frac Rg}[\ln\xi]_{\delta(t)}^{\delta(t_1)}
=\sqrt{\frac Rg}\ln\frac{\delta(t_1)}{\delta(t)}.
\]
当 $\delta\to0$ 时，时间发散。故轻杆真正越过顶点要求 $u^2>4gR$，等号只是渐近倒立位置。

\end{solution}



\section{刚体支承、滑动与内部受力}

刚体可能先滑动或失去支承，也可能先在内部达到强度极限。运动由整体受力确定，内力由截取部分求出，各种临界条件应沿同一运动过程比较。



\subsection{刚体支承与内部受力}
\subsubsection{同时存在多个约束时的事件顺序}
杆沿墙和地面运动时，墙面支持力、地面支持力与摩擦条件都须检查。某一反力先达到边界后，原来的全部几何约束就不再同时成立。分离瞬间通常保留位置和速度连续，再用剩余约束重建动能、动量方程或牛顿方程。

约束力还可以由某个动量分量的变化率求得。例如只有墙面力提供水平外力时，墙面支持力等于总水平动量的导数。动量达到极值是反力为零的另一种表达；极值对应哪个方向的动量，必须由外力方向确定。

\subsubsection{截取刚体的一部分}
刚体内部也要传递力和力偶。截取杆的一段后，切面上一般同时出现轴向力、剪力与弯矩。对子段质心应用平动方程，再对该质心列转动方程，可以避免把加速中的切口误当作固定转轴。

下文的内力符号均由图中箭头定义。强度限制若针对力的大小，应对整个运动范围内的绝对值求最大值；仅求正向峰值会漏掉反向载荷。支承铰链可承受两个方向的力，与地面对杆端的单侧支持不同。


\subsection{例题与模型变式}

\begin{example}[粗糙地面上的刚性棒]\label{ex:hw2-3}

刚性棒总质量$m$，质心到下端$O$的距离为$l$，关于下端的转动惯量为$I_O=\lambda ml^2$。棒从竖直附近由静止释放，底端与地面间摩擦足够大。求何时离地。角度$\theta$从向下倾倒前的竖直位置量起。

\end{example}


\begin{figure}[H]\centering
\begin{tikzpicture}[>=Stealth,line cap=round,line join=round]
\draw[thick](-.7,0)--(2.5,0);\draw[curve](0,0)--(1.4,2.4);\fill(.7,1.2)circle(2pt)node[left]{$C$};\node[below]at(0,0){$O$};\draw[guide](0,0)--(0,2.7);\node at(.2,1){$\theta$};
\end{tikzpicture}
\caption{粗糙地面上以底端为支点的棒}\label{fig:hw2-3-q}
\end{figure}

\begin{solution}
在离地前下端不动，图\ref{fig:hw2-3-s}中质心$C=(l\sin\theta,l\cos\theta)$。平行轴定理要求$\lambda\ge1$。由能量及绕$O$的力矩方程，
\[
 \dot\theta^2=\frac{2g}{\lambda l}(1-\cos\theta),\quad
 \ddot\theta=\frac g{\lambda l}\sin\theta.
\]

\begin{figure}[H]\centering
\begin{tikzpicture}[>=Stealth,line cap=round,line join=round]
\draw[curve](0,0)--(1.4,2.4);\fill(.7,1.2)circle(2pt)node[left]{$C$};\draw[axis](0,0)--(0,.8)node[left]{$N$};\draw[axis](.7,1.2)--(.7,.4)node[right]{$mg$};\draw[axis](0,0)--(.8,0)node[below]{$f_x$};\node at(3,1.4){$I_O=I_C+ml^2$};
\end{tikzpicture}
\caption{水平反力按向右为正，其实际方向由解出的符号决定}\label{fig:hw2-3-s}
\end{figure}

由$y_C=l\cos\theta$逐次求导，
\[
\dot y_C=-l\sin\theta\,\dot\theta,\qquad
\ddot y_C=-l\bigl(\cos\theta\,\dot\theta^2+\sin\theta\,\ddot\theta\bigr)
=-\frac g\lambda(1+2c-3c^2),\quad c=\cos\theta.
\]
竖直方向合力为$N-mg=m\ddot y_C$，代入上式得
\[
 N=\frac{mg}{\lambda}[\lambda-1+3c^2-2c],\quad c=\cos\theta.
\]
令方括号为 $p(c)=\lambda-1+3c^2-2c$，有 $p'=6c-2$、$p''=6>0$。所以 $p$ 在 $c=1/3$ 处最小，最小值为 $\lambda-4/3$。因此
\[
 \boxed{1\le\lambda<\tfrac43:\quad
 \cos\theta_c=\frac{1+\sqrt{4-3\lambda}}3}.
\]
若$\lambda>4/3$，支持力始终为正。若$\lambda=4/3$，$N$在临界角处触零后仍为非负，约束运动可以继续，因而这里不发生分离。

若改用质心惯量$I_C=\kappa ml^2$表示，则$\lambda=\kappa+1$，相应分界为$\kappa=1/3$。这一参数平移来自平行轴定理，比较分离条件时应先统一取轴。

上述判据以底端不滑动为约束假设。实际地面若有有限静摩擦因数$\mu_s$，还须满足
\[
 |f|\le\mu_sN,\qquad
 f=m\ddot x_C=ml(\cos\theta\,\ddot\theta-\sin\theta\,\dot\theta^2)
 =\frac{mg}{\lambda}\sin\theta(3\cos\theta-2).
\]
接近分离根时，$N\to0$，而$f$通常仍不为零，因此有限静摩擦会先达到极限。滑动发生后，实际离地角须结合摩擦条件重新计算。这里的离地角适用于题设采用的理想不滑动约束。


\end{solution}



\begin{example}[倒杆各截面的强度]\label{ex:hw2-4}

均匀刚性棒质量$m$、长度$L$，下端无摩擦铰接于固定点$O$，由近乎竖直的状态静止释放。各截面承受的轴向力、横向剪力及弯矩的绝对值上限分别为$N_0,f_0,M_0$。求保证棒倒地前不破坏所需的最小上限。

\end{example}


\begin{figure}[H]\centering
\begin{tikzpicture}[>=Stealth,line cap=round,line join=round]
\draw[curve](0,0)--(3.1,1.8);\fill(0,0)circle(2pt)node[below]{$O$};\draw[guide](0,0)--(0,2.4);\draw[guide](1.4,.81)--(1.8,.12);\node[above]at(1.4,.81){$s$};\node[right]at(3.1,1.8){$L$};\node at(.4,1){$\theta$};
\end{tikzpicture}
\caption{倒杆上距固定端为$s$的截面}\label{fig:hw2-4-q}
\end{figure}

\begin{solution}
取截面距$O$为$s$，置$z=s/L\in[0,1]$。图\ref{fig:hw2-4-s}单独取截面以外的上段为研究对象，其每一微元距$O$为$r$。均匀杆有
\[
 \omega^2=\frac{3g}{L}(1-\cos\theta),\qquad
 \alpha=\frac{3g}{2L}\sin\theta.
\]

\begin{figure}[H]\centering
\begin{tikzpicture}[>=Stealth,line cap=round,line join=round]
\begin{scope}
\draw[curve](0,0)--(2.6,1.5);\fill(0,0)circle(2pt)node[below]{$O$};
\draw[guide](0,0)--(0,2);\draw[guide](1.15,.6635)--(1.4,.23);
\node[above]at(1.15,.6635){$s$};\node[right]at(2.6,1.5){$L$};\node at(.35,.9){$\theta$};
\end{scope}
\begin{scope}[shift={(4.5,.2)}]
\draw[curve](0,0)--(2.6,1.5);\fill(0,0)circle(1pt)node[below right]{$s$};
\node[right]at(2.6,1.5){$L$};
\draw[axis](0,0)--(.78,.45)node[above]{$N$};\draw[axis](0,0)--(.45,-.78)node[right]{$f$};
\draw[axis](0,.5)arc(90:-170:.5);\node[left]at(-.55,.25){$M$};
\draw[axis](1.55,.894)--(1.55,-.1)node[right]{$g\dd m$};
\node at(1.2,2.05){截面以外的一段};\node at(1.3,-1.2){$a_r=-r\omega^2,\quad a_t=r\alpha$};
\end{scope}
\end{tikzpicture}
\caption{隔离截面上段，分别列轴向、切向及力矩方程}\label{fig:hw2-4-s}
\end{figure}

截面以外一段的质量为$m(1-z)$。重力沿杆向内的分量为$-(m/L)g\cos\theta\dd r$，沿倒下切向的分量为$(m/L)g\sin\theta\dd r$。相应的微元平衡式为
\[
 N=\int_s^L\frac mL(g\cos\theta-r\omega^2)\dd r,
 \qquad f=\int_s^L\frac mL(r\alpha-g\sin\theta)\dd r.
\]
轴力正负对应截面对上段沿杆向外或向内作用，强度条件取其绝对值。

设截面对上段的轴力沿杆向外为正，剪力沿倒下的切向为正。把两个积分分别算出：
\[
\begin{aligned}
N&=\frac mL\left[g\cos\theta\,r-\frac{\omega^2r^2}{2}\right]_s^L
=mg\cos\theta(1-z)-\frac{mL\omega^2}{2}(1-z^2),\\
f&=\frac mL\left[\frac{\alpha r^2}{2}-g\sin\theta\,r\right]_s^L
=\frac{mL\alpha}{2}(1-z^2)-mg\sin\theta(1-z).
\end{aligned}
\]
弯矩取与剪力约定相容的符号，对上段各微元的切向惯性项与重力取矩，
\[
\begin{aligned}
M&=\frac mL\int_s^L(r-s)(r\alpha-g\sin\theta)\,\dd r\\
&=\frac mL\left[\alpha\left(\frac{r^3}{3}-\frac{sr^2}{2}\right)
-g\sin\theta\left(\frac{r^2}{2}-sr\right)\right]_s^L\\
&=\frac mL\left[\alpha\left(\frac{L^3}{3}-\frac{sL^2}{2}+\frac{s^3}{6}\right)
-\frac{g\sin\theta}{2}(L-s)^2\right].
\end{aligned}
\]
代入已求出的$\omega^2,\alpha$并用$s=zL$整理，得到
\begin{align*}
 N(s,\theta)&=\frac{mg}{2}(1-z)[(5+3z)\cos\theta-3(1+z)],\\
 f(s,\theta)&=\frac{mg}{4}\sin\theta(1-z)(3z-1),\\
 M(s,\theta)&=\frac{mgL}{4}\sin\theta\,z(1-z)^2.
\end{align*}
轴力随 $\cos\theta$ 线性变化，绝对值最大值在竖直或水平极限取到。竖直时 $N=mg(1-z)$，最大绝对值为 $mg$；水平时 $N=-3mg(1-z^2)/2$，最大绝对值为 $3mg/2$。剪力的角度因子 $\sin\theta$ 在水平极限为 1，再令
\[
p(z)=(1-z)(3z-1)=-3z^2+4z-1,\qquad p'(z)=4-6z.
\]
比较 $p(0)=-1$、$p(1)=0$、$p(2/3)=1/3$，得剪力绝对值最大为 $mg/4$。弯矩同样在水平时最大，其位置因子满足
\[
\frac{\dd}{\dd z}[z(1-z)^2]
=(1-z)^2-2z(1-z)=(1-z)(1-3z).
\]
端点值均为零，内点 $z=1/3$ 处取最大值 $4/27$。因此
\[
 \boxed{N_{0,\min}=\frac32mg,\quad f_{0,\min}=\frac14mg,
 \quad M_{0,\min}=\frac{mgL}{27}}.
\]
这些上限针对杆落地前的连续运动。落地碰撞中的冲击载荷须另作计算。


\end{solution}



\begin{example}[固定杆端假设与先滑动的边界]\label{deep2:problem-example-6}\label{deep2:example:d2-e6}

刚体质量为 m，质心 C 到底端 O 的距离为 $\ell$，绕 O 的转动惯量为 $I_O=\lambda m\ell^2$，$\lambda\ge1$。刚体从竖直附近静止向右倒下，底端接触粗糙水平地面，见图\ref{deep2:fig:d2-e6-q}。先假定底端不动，求支持力与摩擦力，再讨论 N=0 的候选条件。对均匀细杆，若 $\mu=3\sqrt3$，求首次滑动角。

\end{example}
\begin{figure}[H]\centering
\begin{tikzpicture}[scale=.95]
\coordinate(O)at(0,0);\coordinate(C)at({2*sin(35)},{2*cos(35)});
\draw[thick](-.6,0)--(3.5,0);\draw[line width=2pt](O)--({3.3*sin(35)},{3.3*cos(35)});
\draw[guide](O)--(0,3);\draw(0,.7)arc(90:55:.7);\node at(.35,.9){$\theta$};
\fill(O)circle(1.5pt)node[below]{$O$};\fill(C)circle(1.5pt)node[left]{$C$};\node[anchor=west]at(2.6,3.1){$OC=\ell$};
\node[anchor=west]at(2.6,2){$I_O=\lambda m\ell^2$};
\end{tikzpicture}
\caption{底端尚未滑动时的刚体构形。惯量参数相对于底端定义。}\label{deep2:fig:d2-e6-q}
\end{figure}
\par\medskip

\begin{solution}
在底端保持不动的阶段，重力势能的减少等于绕 O 的转动动能。令 $c=\cos\theta$、$s=\sin\theta$，得到
\[
\dot\theta^2=\frac{2g}{\lambda\ell}(1-c),\qquad
\ddot\theta=\frac g{\lambda\ell}s.
\]
质心坐标为 $(\ell s,\ell c)$。对坐标求两次导数，
\[
a_{Cx}=\ell(c\ddot\theta-s\dot\theta^2)
=\frac g\lambda s(3c-2),
\]
\[
a_{Cy}=-\ell(s\ddot\theta+c\dot\theta^2)
=-\frac g\lambda(1+2c-3c^2).
\]
图\ref{deep2:fig:d2-e6-s}规定 f 向右为正、N 向上为正。由质心平动方程，
\[
f=\frac{mg}{\lambda}s(3c-2),\qquad
N=\frac{mg}{\lambda}(\lambda-1-2c+3c^2).
\]

\begin{figure}[H]\centering
\begin{tikzpicture}[scale=.95]
\coordinate(O)at(0,0);\coordinate(C)at({2*sin(35)},{2*cos(35)});
\draw[thick](-.6,0)--(3.5,0);\draw[line width=2pt](O)--({3.3*sin(35)},{3.3*cos(35)});
\draw[guide](O)--(0,3);\draw(0,.7)arc(90:55:.7);\node at(.35,.9){$\theta$};
\fill(O)circle(1.5pt)node[below]{$O$};\fill(C)circle(1.5pt)node[left]{$C$};\node[anchor=west]at(2.6,3.1){$OC=\ell$};
\node[anchor=west]at(2.6,2){$I_O=\lambda m\ell^2$};

\draw[axis](O)--(0,1.15)node[left]{$N$};\draw[axis](O)--(.95,0)node[below]{$f$};
\draw[axis](C)--++(0,-1)node[right]{$mg$};
\node at(4,1){$|f|\le\mu N$};
\end{tikzpicture}
\caption{固定端阶段的实际外力。摩擦方向由 f 的符号确定，图上正向只用于列式。}\label{deep2:fig:d2-e6-s}
\end{figure}

在 $0\le c\le1$ 内，N 在 $c=1/3$ 处最小，$N_{\min}=mg[1-4/(3\lambda)]$。若 $\lambda>4/3$，固定端模型的 N 始终为正。若 $\lambda=4/3$，只在 $c=1/3$ 触零。若 $1\le\lambda<4/3$，从 c=1 开始首先遇到的零点为
\[
c_N=\frac{1+\sqrt{4-3\lambda}}3.
\]
这只是“底端不动”模型的候选边界。对 $1<\lambda\le4/3$，在该候选点附近 f 不为零而 N 趋于零，$|f|/N$ 发散。任何有限静摩擦因数都不能将底端一直固定到那里，实际应先滑动。

对均匀杆，$\ell=L/2$、$I_O=mL^2/3$，所以 $\lambda=4/3$，
\[
N=\frac{mg}{4}(3c-1)^2,\qquad
\frac{|f|}{N}=\frac{3\sqrt{1-c^2}|3c-2|}{(3c-1)^2}.
\]
在 $2/3\le c\le1$，分母至少为 1、分子不超过 3，故比值小于给定的 $3\sqrt3$，这一阶段静摩擦尚未达到极限。

在 $1/3<c<2/3$，$f$ 已向左。比值为 $R(c)=3\sqrt{1-c^2}(2-3c)/(3c-1)^2$，取对数，
\[
\ln R=\ln3+\tfrac12\ln(1-c^2)+\ln(2-3c)-2\ln(3c-1).
\]
逐项使用链式法则，得到
\[
\frac{R'}R=-\frac c{1-c^2}-\frac3{2-3c}-\frac6{3c-1}<0.
\]
真实运动中 c 减小，R 因而严格增大。在 c=1/2 时
\[
N=\frac{mg}{16},\qquad |f|=\frac{3\sqrt3}{16}mg=\mu N.
\]
所以首次滑动角是 $\theta=60^\circ$，早于候选 N=0 的角度 $\arccos(1/3)$。滑动后 O 的位置也成为未知量，不能继续用固定端的能量与反力公式。

\end{solution}



\section{约束极值、虚功与二阶变化}

平衡要求势能在允许方向上的一阶变化为零，稳定性则由更高阶变化决定。几何约束可用独立坐标消去，也可用乘数保留；两种方法都须检查可行边界和扰动方向。



\subsection{条件极值与约束平衡}
\subsubsection{可行方向与拉格朗日乘数}
约束 $G(q_1,\ldots,q_n)=0$ 限制允许位移满足 $\nabla G\cdot\delta\vct q=0$。在约束流形的正则点，若 f 对全部可行方向的一阶变化均为零，其梯度必须与约束梯度平行，即
\[
\nabla f=\lambda\nabla G.
\]
若有若干个梯度线性无关的约束，则 $\nabla f$ 是这些梯度的线性组合。将约束一起代入，才能确定候选点。

乘数方程给出驻值点。确定全局极值时，还须考察可行集的范围、边界及退化构形。推导中若除以某个坐标或三角函数，必须把其为零的情形另列出来。正则性失效的点也不能仅凭乘数方程排除。

对保守静力系统，固定理想约束不对可行虚位移做功，平衡条件为 $\delta U=0$。若所选坐标独立，应令各偏导为零；若坐标仍受固定几何关系约束，则应先消元或引入乘数。稳定性由沿可行方向的二阶变化决定。接触为单侧时，位于边界的平衡还须保留相应反力的非负条件。

\subsubsection{变分方程的初步}
有限维条件极值中改变的是若干坐标；若要改变整条曲线 $y(x)$，可研究泛函
\[
J[y]=\int_{x_0}^{x_1}\mathcal L(x,y,y')\dd x.
\]
令 $y\mapsto y+\epsilon\eta$，且端点固定，$\eta(x_0)=\eta(x_1)=0$。在 $\epsilon=0$ 处对泛函求导，
\[
\delta J=\int_{x_0}^{x_1}(\mathcal L_y\eta+\mathcal L_{y'}\eta')\dd x.
\]
第二项分部积分，
\[
\int_{x_0}^{x_1}\mathcal L_{y'}\eta'\dd x
=[\mathcal L_{y'}\eta]_{x_0}^{x_1}
-\int_{x_0}^{x_1}\frac{\dd\mathcal L_{y'}}{\dd x}\eta\dd x.
\]
边界项因端点变化为零而消失，得到
\[
\delta J=\int_{x_0}^{x_1}
\left(\frac{\partial\mathcal L}{\partial y}
-\frac\dd{\dd x}\frac{\partial\mathcal L}{\partial y'}\right)\eta\dd x.
\]
若对任意允许的 $\eta$ 均为零，曲线内部须满足欧拉—拉格朗日方程
\[
\frac{\partial\mathcal L}{\partial y}
-\frac\dd{\dd x}\frac{\partial\mathcal L}{\partial y'}=0.
\]
当 $\mathcal L$ 不显含 x 时，可得到第一积分 $\mathcal L-y'\mathcal L_{y'}=\text{常量}$。第\ref{sec:complete2-variation}节将展开沿曲线求导的乘积项，并用欧拉方程使它们相消。第一积分仍是驻值的必要条件；是否为极小、能否处理端点奇性，需要结合具体问题。


\subsection{受约束的二阶变化}
\label{sec:complete2}
设驻值点满足$\nabla f=\sum_a\lambda_a\nabla g_a$，各约束写成$g_a=0$。在该点取切向变化$\vct\xi$，满足$\nabla g_a\cdot\vct\xi=0$。为保留约束曲率，取实际落在约束面上的曲线$\vct q(\varepsilon)$，令$\vct q'(0)=\vct\xi$、$\vct q''(0)=\vct\eta$。对$g_a(\vct q(\varepsilon))=0$连续求导，
\[\nabla g_a\cdot\vct\xi=0,\qquad
\vct\xi^{\mathsf T}\nabla^2g_a\vct\xi+\nabla g_a\cdot\vct\eta=0.
\]
目标函数沿该曲线的二阶导数为
\[\frac{\dd^2f}{\dd\varepsilon^2}
=\vct\xi^{\mathsf T}\nabla^2f\vct\xi+\nabla f\cdot\vct\eta.
\]
把驻值条件$\nabla f=\sum_a\lambda_a\nabla g_a$代入第二项，再由约束的二阶关系消去$\vct\eta$，得到
\[
 \delta^2 f=\vct\xi^{\mathsf T}
 \left[\nabla^2 f-\sum_a\lambda_a\nabla^2g_a\right]\vct\xi.
\]
局部极值由允许切空间内的二次型决定。若直接在整个坐标空间中考察$f$的二阶导数，就会混入不满足约束的变化。若二阶变化退化，应继续计算高阶项；若约束含不等式，还要检验边界上的单侧变化。


\subsection{例题与模型变式}

\begin{example}[竖直导杆上的滑环]\label{ex:hw2-5}

长$3l$的轻杆在$A$端固定质量$m$的小球，距$A$为$l$处固定轻环$C$，环套在固定光滑竖直导杆上；另一端$B$接触光滑地面。杆从近乎竖直状态由静止释放，求何时离开地面。图\ref{fig:hw2-5-q}中$AC=l$、$CB=2l$。

\end{example}


\begin{figure}[H]\centering
\begin{tikzpicture}[>=Stealth,line cap=round,line join=round]
\draw[thick](-.5,0)--(3,0);\draw[thick](.676189,.05)--(.676189,2.85);
\coordinate(B)at(0,0);\coordinate(C)at(.676189,1.450092);\coordinate(A)at(1.014284,2.175138);
\draw[curve](B)--(A);\draw(C)circle(.1);\fill(A)circle(2pt)node[right]{$A,m$};
\node[left]at(C){$C$};\node[below]at(B){$B$};\node[left]at(.33,1.15){$2l$};\node[right]at(.88,1.75){$l$};
\draw[guide](C)--(.676189,2.2);\draw(.676189,1.950092)arc(90:65:.5)node[midway,above]{$\theta$};
\end{tikzpicture}
\caption{竖直导杆约束下的滑环与轻杆}\label{fig:hw2-5-q}
\end{figure}

\begin{solution}
取导杆为$y$轴，$x$轴沿地面向右，杆与竖直方向夹角为$\theta$。图\ref{fig:hw2-5-q}中$AC=l$、$CB=2l$，接触地面期间有
\[
 A=(l\sin\theta,3l\cos\theta),\qquad
 B=(-2l\sin\theta,0).
\]
令$c=\cos\theta$、$D=9-8c^2$。由$A$的速度平方，
\[
 T=\frac12ml^2(c^2+9\sin^2\theta)\dot\theta^2
 =\frac12ml^2D\dot\theta^2,\qquad V=3mglc.
\]
以近乎竖直、从静止释放的极限取总能量$3mgl$，于是
\[
 \dot\theta^2=\frac{6g(1-c)}{lD},\qquad
 v_{Ax}^2=l^2c^2\dot\theta^2
 =6gl\,\frac{c^2(1-c)}{D}.
\]
\begin{figure}[H]\centering
\begin{tikzpicture}[>=Stealth,line cap=round,line join=round]
\draw[thick](-.5,0)--(3,0);\draw[thick](.676189,.05)--(.676189,2.85);
\coordinate(B)at(0,0);\coordinate(C)at(.676189,1.450092);\coordinate(A)at(1.014284,2.175138);
\draw[curve](B)--(A);\draw(C)circle(.1);\fill(A)circle(2pt)node[right]{$A,m$};
\node[left]at(C){$C$};\node[below]at(B){$B$};\node[left]at(.33,1.15){$2l$};\node[right]at(.88,1.75){$l$};
\draw[guide](C)--(.676189,2.2);\draw(.676189,1.950092)arc(90:65:.5)node[midway,above]{$\theta$};\draw[axis](B)--(0,.8)node[left]{$N$};\draw[axis](C)--(1.6,1.450092)node[below]{$H$};
\end{tikzpicture}
\caption{地面支持力和导杆水平反力}\label{fig:hw2-5-s}
\end{figure}
图\ref{fig:hw2-5-s}中的$H$向右为正，$N$向上为正。对整个系统列平动方程：
\[
 H=m\ddot x_A,\qquad N-mg=m\ddot y_A.
\]
轻杆、轻环没有惯性，对它们绕$A$取矩。$B$点的竖直支持力到$A$的水平力臂为$AB\sin\theta=3l\sin\theta$；$C$点的水平反力到$A$的竖直力臂为$AC\cos\theta=l\cos\theta$。两力矩方向相反，故
$3l\sin\theta\,N=l\cos\theta\,H$。
因此在$0<\theta<\pi/2$内，$N,H$同时为零。此时$v_{Ax}>0$，且$c$随时间严格减小，由
$2v_{Ax}\ddot x_A=(\dd v_{Ax}^2/\dd c)\dot c$，可用速度平方的驻值求临界位置：
\[
 \frac{\dd}{\dd c}\frac{c^2(1-c)}D
 =\frac{(2c-3c^2)D+16c(c^2-c^3)}{D^2}
 =\frac{c(8c^3-27c+18)}{D^2}.
\]
排除不在该开区间内的$c=0$，得到临界候选条件
\[
 \boxed{8c_*^3-27c_*+18=0,\qquad 0<c_*<1}.
\]
为确定离地前的接触区间，由$\dot c=-\sin\theta\,\dot\theta$计算支持力：
\[
 \dot c^2=\frac{6g}{l}\frac{(1-c)^2(1+c)}D.
\]
在$\dot c\ne0$的运动阶段求导并约去$2\dot c$，得
\[
 \ddot c=\frac{3g}{l}
 \frac{(3c^2-2c-1)D+16c(1-c)^2(1+c)}{D^2}.
\]
将$y_A=3lc$代入竖直动力学方程，合并分子后为
\[
 N=mg+3ml\ddot c
 =-\frac{mgc(8c^3-27c+18)}{D^2}.
\]
多项式$p(c)=8c^3-27c+18$在$(0,1)$内有
$p'(c)=24c^2-27<0$，且$p(0)=18$、$p(1)=-1$，故根唯一。
释放后$c$从接近$1$向下减小，在$c_*<c<1$时$p(c)<0$，所以$N>0$；越过$c_*$，接触解将要求$N<0$，地面不能提供这样的拉力。因而
\[
 \boxed{c_*\simeq0.846200,\qquad\theta_c=\arccos c_*\simeq32.20^\circ}.
\]
严格竖直且没有扰动的构形可静止不动；这里取微小倾角释放再令初始倾角趋零。离地后应撤去$y_B=0$，不能延用本阶段的几何关系。
\end{solution}



\begin{example}[定表面积开口圆筒的容量]\label{ex:hw2-6}

开口圆筒底面半径$r>0$、高度$h>0$，材料总面积固定为$S_0=\pi r^2+2\pi rh$。求最大容积。

\end{example}


\begin{figure}[H]\centering
\begin{tikzpicture}[>=Stealth,line cap=round,line join=round]
\draw(0,0)ellipse(1.1 and .3);\draw(-1.1,0)--(-1.1,2);\draw(1.1,0)--(1.1,2);\draw(0,2)ellipse(1.1 and .3);\draw[<->](1.6,0)--(1.6,2)node[midway,right]{$h$};\draw[axis](0,0)--(1.1,0)node[midway,below]{$r$};
\end{tikzpicture}
\caption{定材料面积的开口圆筒}\label{fig:hw2-6-q}
\end{figure}

\begin{solution}
图\ref{fig:hw2-6-s}的底面与侧面分别贡献$\pi r^2$和$2\pi rh$。先由约束消去$h$：
\[
 h=\frac{S_0-\pi r^2}{2\pi r},\quad
 0<r<\sqrt{S_0/\pi},\quad V=\frac12(S_0r-\pi r^3).
\]

\begin{figure}[H]\centering
\begin{tikzpicture}[>=Stealth,line cap=round,line join=round]
\draw(0,0)ellipse(1.1 and .3);\draw(-1.1,0)--(-1.1,1.1);\draw(1.1,0)--(1.1,1.1);
\draw(0,1.1)ellipse(1.1 and .3);\draw[<->](1.6,0)--(1.6,1.1)node[midway,right]{$h=r$};
\draw[axis](0,0)--(1.1,0)node[midway,below]{$r$};
\node at(4,1){$S_{\text{底}}=S_0/3$};\node at(4,.3){$S_{\text{侧}}=2S_0/3$};
\end{tikzpicture}
\caption{最优圆筒的底面与侧面面积分配}\label{fig:hw2-6-s}
\end{figure}

在允许区间内
\[
 V'=\frac12(S_0-3\pi r^2)=0
 \quad\Longrightarrow\quad r_*^2=\frac{S_0}{3\pi},\quad h_*=r_*.
\]
$V''=-3\pi r<0$，且两个端点极限的容积都为零，因此该驻点给出全局最大值：
\[
 \boxed{V_{\max}=\pi\left(\frac{S_0}{3\pi}\right)^{3/2}}.
\]


也可用乘数法求最优面积分配。记
\[
 \mathcal L=\pi r^2h-\lambda(\pi r^2+2\pi rh-S_0).
\]
对$h,r$分别求偏导，有
\[
 \pi r^2-2\pi\lambda r=0,\qquad
 2\pi rh-2\pi\lambda(r+h)=0.
\]
由$r>0$知$\lambda=r/2$，第二式随即化为$2h=r+h$，故$h=r$。代回面积约束，底面积占$S_0/3$，侧面积占$2S_0/3$。这种比例适用于有底、无盖的圆筒；若另加上盖，约束式改变，最优高径比也随之改变。

\end{solution}



\begin{example}[球面与交圆上的乘积极值]\label{ex:hw2-7}

求$f=xyz$在以下条件下的最大值和最小值：(1)$x^2+y^2+z^2=1$；(2)在前一条件外另有$x+y+z=0$。

\end{example}


\begin{figure}[H]\centering
\begin{tikzpicture}[>=Stealth,line cap=round,line join=round]
\draw(0,0)circle(1.4);\draw[guide](0,0)ellipse(1.4 and .4);\draw[axis](0,0)--(2,0)node[right]{$x$};\draw[axis](0,0)--(0,1.9)node[above]{$z$};\draw[axis](0,0)--(-1,-.8)node[left]{$y$};\node at(3,1){$x^2+y^2+z^2=1$};
\end{tikzpicture}
\caption{球面上的三坐标约束}\label{fig:hw2-7-q}
\end{figure}

\begin{solution}
只有球面约束时，对$x^2,y^2,z^2$应用均值不等式，有$|xyz|\le1/(3\sqrt3)$。等号要求三个坐标的绝对值均为$1/\sqrt3$，选择符号得到
\[
 \boxed{f_{\max}=\frac1{3\sqrt3},\qquad f_{\min}=-\frac1{3\sqrt3}}.
\]


\begin{figure}[H]\centering
\begin{tikzpicture}[>=Stealth,line cap=round,line join=round]
\draw[axis](-2.05,0)--(2.05,0)node[right]{$\xi$};\draw[axis](0,-2.05)--(0,2.05)node[above]{$\eta$};\draw[curve](0,0)circle(1.65);
\foreach \ang/\sgn in {30/+,90/-,150/+,210/-,270/+,330/-}{\fill({1.65*cos(\ang)},{1.65*sin(\ang)})circle(1.8pt);\ifnum\ang=90\node[right,fill=white,inner sep=1pt]at(.12,1.78){$\sgn$};\else\ifnum\ang=270\node[right,fill=white,inner sep=1pt]at(.12,-1.78){$\sgn$};\else\node at({1.96*cos(\ang)},{1.96*sin(\ang)}){$\sgn$};\fi\fi}
\node[right,align=left]at(2.5,.65){$\xi^2+\eta^2=1$\\[4pt]$x+y+z=0$\\[4pt]$+$：极大值点\\$-$：极小值点};
\node[below]at(0,-2.45){交圆所在平面内的正交坐标};
\end{tikzpicture}
\caption{交圆的准确平面图；六个驻点的函数值正负交替}\label{fig:hw2-7-s}
\end{figure}


若改用乘数法处理第一问，$yz=2\lambda x$等方程在$xyz\ne0$时给出$x^2=y^2=z^2$；若有零坐标，函数值为零。两类候选点共同覆盖球面上的驻值情形。

加入平面后，可行集是图\ref{fig:hw2-7-s}中的交圆。两个约束的梯度在交圆上处处线性无关，因此全部极值均满足
\[
 yz=2\lambda x+\mu,\quad zx=2\lambda y+\mu,\quad xy=2\lambda z+\mu.
\]
两两相减，得$(x-y)(z+2\lambda)=0$及其循环式。若三个坐标互不相同，会推出三者都等于$-2\lambda$，矛盾，故至少两坐标相等。令$y=z=t$，则$x=-2t$、$6t^2=1$，产生六个排列候选点，函数值为$\pm2/(6\sqrt6)$。
\[
 \boxed{f_{\max}=\frac1{3\sqrt6},\qquad f_{\min}=-\frac1{3\sqrt6}}.
\]
含零坐标的可行点函数值为零，不超过上述上下界。交圆紧致且候选点已经穷尽，故两界均为全局极值。


为把六个点与图\ref{fig:hw2-7-s}对应，取交圆平面内的单位正交基
\[
 \vct e_1=\frac{(1,-1,0)}{\sqrt2},\qquad
 \vct e_2=\frac{(1,1,-2)}{\sqrt6},\qquad
 (x,y,z)=\xi\vct e_1+\eta\vct e_2.
\]
此时$\xi^2+\eta^2=1$。再置$\xi=\cos\phi,\eta=\sin\phi$，则
\[
 xyz=\frac{\sin\phi(3\cos^2\phi-\sin^2\phi)}{3\sqrt6}
 =\frac{3\sin\phi-4\sin^3\phi}{3\sqrt6}
 =\frac{\sin3\phi}{3\sqrt6}.
\]
对圆周参数求导，$\dd(xyz)/\dd\phi=\cos3\phi/\sqrt6=0$，因此六个驻点恰在 $\phi=\pi/6+j\pi/3$，$j=0,\ldots,5$，沿圆周依次取极大、极小值。它们正是 $(2,-1,-1)/\sqrt6$ 的三个排列及其相反数；参数法给出的六个极值点与乘数法一致。

\end{solution}



\begin{example}[固定悬点与弹簧连接悬点的四杆链]\label{ex:hw2-8}

四根长$l$的轻杆依次铰接，三个中间结点的质量依次为$m,2m,m$，构形左右对称，上、下层杆与水平夹角分别为$\alpha,\beta\in(0,\pi/2)$。(1)两个上悬点固定且相距$3l$；(2)上悬点改成在光滑水平杆上滑动的轻环，两环用原长$4l$、劲度系数$k=mg/l$的弹簧连接。分别求平衡角。

\end{example}


\begin{figure}[H]\centering
\begin{tikzpicture}[>=Stealth,line cap=round,line join=round]
\coordinate(A)at(-2.3999999999999995,0);\coordinate(B)at(2.3999999999999995,0);\coordinate(C)at(-1.373381785857214,-1.2272143424806758);\coordinate(D)at(0,-2.0480812441618497);\coordinate(E)at(1.373381785857214,-1.2272143424806758);\draw[thick](-3,0)--(3,0);\draw[curve](A)--(C)--(D)--(E)--(B);\foreach \p/\m in {C/m,D/2m,E/m}{\fill(\p)circle(2pt)node[below]{$\m$};}\node[above left]at(A){$A$};\node[above right]at(B){$B$};\draw[guide](C)--++(.8,0);\draw(-1.8999999999999995,0)arc(0:-50.08605097462309:.5)node[midway,right]{$\alpha$};\draw(-0.8733817858572139,-1.2272143424806758)arc(0:-30.866656674099772:.5)node[midway,right]{$\beta$};\draw[<->](-2.3999999999999995,.65)--(2.3999999999999995,.65)node[midway,above]{$3l$};
\end{tikzpicture}
\caption{悬点相距三倍杆长的四杆链}\label{fig:hw2-8-q}
\end{figure}

\begin{solution}
图\ref{fig:hw2-8-s}中两侧球下降$l\sin\alpha$，中央球下降$l(\sin\alpha+\sin\beta)$，故重力势能为$U_g=-2mgl(2\sin\alpha+\sin\beta)$。

\begin{figure}[H]\centering
\begin{tikzpicture}[>=Stealth,line cap=round,line join=round]
\coordinate(A)at(-2.150266565775473,0);\coordinate(B)at(2.150266565775473,0);\coordinate(C)at(-1.2725720522998616,-1.337778883453768);\coordinate(D)at(0,-2.307603801521622);\coordinate(E)at(1.2725720522998616,-1.337778883453768);\draw[thick](-3,0)--(3,0);\draw[curve](A)--(C)--(D)--(E)--(B);\foreach \p/\m in {C/m,D/2m,E/m}{\fill(\p)circle(2pt)node[below]{$\m$};}\node[above left]at(A){$A$};\node[above right]at(B){$B$};\draw[guide](C)--++(.8,0);\draw(-1.650266565775473,0)arc(0:-56.73178472991683:.5)node[midway,right]{$\alpha$};\draw(-0.7725720522998616,-1.337778883453768)arc(0:-37.3109614418283:.5)node[midway,right]{$\beta$};\draw(A)circle(.08);\draw(B)circle(.08);\draw[spring](-2.150266565775473,.6)--(2.150266565775473,.6);\draw[thin](A)--(-2.150266565775473,.6);\draw[thin](B)--(2.150266565775473,.6);\node at(0,1){$d=2l(\cos\alpha+\cos\beta)$};
\end{tikzpicture}
\caption{滑环间距随弹簧压缩量改变}\label{fig:hw2-8-s}
\end{figure}

固定悬点时$c_\alpha+c_\beta=3/2$。令
\[
\mathcal L=\frac{U_g}{mgl}+\Lambda(c_\alpha+c_\beta-3/2).
\]
对两角求偏导，得到
\[
-4c_\alpha-\Lambda s_\alpha=0,\qquad
-2c_\beta-\Lambda s_\beta=0.
\]
两式分别除以非零的正弦，消去$\Lambda$，即$4\cot\alpha=2\cot\beta$，故$\tan\alpha=2\tan\beta$。令$t=\tan\beta>0$，平衡由
\[
 \frac1{\sqrt{1+4t^2}}+\frac1{\sqrt{1+t^2}}=\frac32,
 \quad\alpha=\arctan(2t),\ \beta=\arctan t
\]
唯一确定。左端随$t$严格递减，从$t=0$时的2趋向零，故恰有一个正根。数值求得$t\simeq0.597698$，即$\alpha\simeq50.09^\circ$、$\beta\simeq30.87^\circ$。由结点平衡，上杆竖直载荷为$2mg$、下杆竖直载荷为$mg$；两杆水平分力相等，也给出相同的角度关系。

滑环可动时，弹簧长度$d=2l(c_\alpha+c_\beta)$，总势能为
\[
 U=-2mgl(2s_\alpha+s_\beta)+2mgl(c_\alpha+c_\beta-2)^2.
\]
记$C=c_\alpha+c_\beta$。分别对两个独立角求导，得
\[
\frac{U_\alpha}{mgl}=-4c_\alpha+4(2-C)s_\alpha=0,
\qquad
\frac{U_\beta}{mgl}=-2c_\beta+4(2-C)s_\beta=0.
\]
除以各自的正弦并整理，得到
\[
 \cot\alpha=2-c_\alpha-c_\beta,\qquad
 \cot\beta=2(2-c_\alpha-c_\beta).
\]
仍有$\tan\alpha=2\tan\beta$，但固定间距条件已被弹簧方程替代。置$d_*=2-c_\alpha-c_\beta>0$，则
\[
 \boxed{d_*=2-\frac{d_*}{\sqrt{1+d_*^2}}
 -\frac{2d_*}{\sqrt{1+4d_*^2}},\quad
 \alpha=\arctan(1/d_*),\quad\beta=\arctan[1/(2d_*)]}.
\]
右端随$d_*$严格递减，左端严格递增；在$d_*=0$时右端大于左端，而在$d_*=2$时右端小于左端，因此正根唯一。数值为$d_*\simeq0.656083$，于是$\alpha\simeq56.73^\circ$、$\beta\simeq37.31^\circ$。

两个平衡点的稳定性可由势能的二阶变化判定。第一问对间距约束求变分，$-s_\alpha\delta\alpha-s_\beta\delta\beta=0$，故 $\delta\beta=-(s_\alpha/s_\beta)\delta\alpha$。乘数函数的二阶偏导为
\[
\mathcal L_{\alpha\alpha}=4s_\alpha-\Lambda c_\alpha
=4s_\alpha+\frac{4c_\alpha^2}{s_\alpha}=\frac4{s_\alpha},
\qquad
\mathcal L_{\beta\beta}=2s_\beta-\Lambda c_\beta=\frac2{s_\beta},
\]
其中分别代入了平衡式 $\Lambda=-4c_\alpha/s_\alpha=-2c_\beta/s_\beta$；混合偏导为零。因此约束下的二阶变化为
\[
 \frac{\delta^2U_{\rm constrained}}{mgl}
 =\frac4{\sin\alpha}(\delta\alpha)^2+
 \frac2{\sin\beta}(\delta\beta)^2>0.
\]
第二问有 $C_\alpha=-s_\alpha$、$C_\beta=-s_\beta$。在势能的一阶偏导上继续求导，
\[
\begin{aligned}
\frac{U_{\alpha\alpha}}{mgl}&=4s_\alpha+4s_\alpha^2+4(2-C)c_\alpha,\\
\frac{U_{\beta\beta}}{mgl}&=2s_\beta+4s_\beta^2+4(2-C)c_\beta,\\
\frac{U_{\alpha\beta}}{mgl}&=4s_\alpha s_\beta.
\end{aligned}
\]
将平方项和交叉项合并，平衡点的二次型为
\[
 \frac{\delta^2U}{mgl}=
 4\bigl[(2-C)c_\alpha+ s_\alpha\bigr](\delta\alpha)^2
 +4\bigl[(2-C)c_\beta+\tfrac12s_\beta\bigr](\delta\beta)^2
 +4(s_\alpha\delta\alpha+s_\beta\delta\beta)^2.
\]
在题设的向下对称分支上各系数均为正，因此平衡是局部势能极小。

滑环间距由弹簧变形决定，已不再满足固定悬点题中的$d=3l$。


\end{solution}



\section{曲线约束与最速下降}

轨道形状待定时，未知量是一条曲线。运动时间成为曲线的泛函，其第一积分给出候选轨道；端点和可行区域仍决定哪些候选符合题意。



\subsection{最速下降曲线的第一积分}
\label{sec:complete2-variation}
在竖直平面内以向下为$y$正向。质点从$O=(0,0)$静止出发，沿固定光滑轨道到达$B=(x_B,y_B)$，其中$x_B,y_B>0$。任意下降轨道上的速率由能量给出$v=\sqrt{2gy}$，时间泛函为
\[
 t[y]=\int_0^{x_B}\sqrt{\frac{1+y'^2}{2gy}}\,\dd x.
\]
图\ref{fig:hw2-variation}标出端点和坐标。端点固定的驻值曲线在内部满足欧拉--拉格朗日方程。起点速率为零，其邻域用参数曲线处理。
\begin{figure}[H]\centering
\begin{tikzpicture}[>=Stealth]
\draw[axis](0,0)--(4.3,0)node[right]{$x$};
\draw[axis](0,0)--(0,-3.4)node[below]{$y$};
\draw[curve,domain=0:145,samples=90]plot({1.5*(\x*pi/180-sin(\x))},{-1.5*(1-cos(\x))});
\draw[guide](0,0)--(2.936,-2.729);
\fill(0,0)circle(1.4pt)node[above left]{$O$};
\fill(2.936,-2.729)circle(1.4pt)node[right]{$B$};
\end{tikzpicture}
\caption{固定端点之间的下降轨道，虚线为直线连接}\label{fig:hw2-variation}
\end{figure}
对不显含$x$的积分核$\mathcal L$，直接微分得
\[
 \begin{aligned}
 \frac\dd{\dd x}(\mathcal L-y'\mathcal L_{y'})
 &=\mathcal L_y y'+\mathcal L_{y'}y''-y''\mathcal L_{y'}
 -y'\frac\dd{\dd x}\mathcal L_{y'}\\
 &=y'\left(\mathcal L_y-\frac\dd{\dd x}\mathcal L_{y'}\right)=0.
 \end{aligned}
\]
将本题积分核代入，得到
\[
 \frac1{\sqrt{2gy}\sqrt{1+y'^2}}=C,
 \qquad y(1+y'^2)=2a,
\]
其中$a>0$为常数。令$y=a(1-\cos\psi)$，第一下降段满足$y'=\cot(\psi/2)$，于是
\[
 \frac{\dd x}{\dd\psi}
 =\frac{a\sin\psi}{\cot(\psi/2)}=a(1-\cos\psi).
\]
由 $x(0)=0$ 积分，
\[
x=a\int_0^\psi(1-\cos u)\dd u
=a[u-\sin u]_0^\psi=a(\psi-\sin\psi).
\]
图\ref{fig:complete2-cycloid}中的参数曲线为摆线。由端点求$a,\psi_B$：
\[
 \frac{x_B}{y_B}=\frac{\psi_B-\sin\psi_B}{1-\cos\psi_B},\qquad
 a=\frac{y_B}{1-\cos\psi_B}.
\]
选取$0<\psi_B<2\pi$的第一摆线弧时，曲线可以先降至$\psi=\pi$，再上升至终点；不能预先强加全程$y'>0$。沿参数曲线，$\dd s=2a\sin(\psi/2)\dd\psi$、$v=2\sqrt{ga}\sin(\psi/2)$，故
\[
 t_B=\sqrt{a/g}\int_0^{\psi_B}\dd\psi
 =\sqrt{a/g}[\psi]_0^{\psi_B}=\sqrt{a/g}\,\psi_B.
\]
\begin{figure}[H]\centering
\begin{tikzpicture}[>=Stealth,scale=.9]
\draw[axis](0,0)--(5,0)node[right]{$x$};\draw[axis](0,0)--(0,-3)node[below]{$y$};
\draw[curve,domain=0:155,samples=90]plot({1.2*(\x*pi/180-sin(\x))},{-1.2*(1-cos(\x))});
\coordinate(P)at({1.2*(100*pi/180-sin(100))},{-1.2*(1-cos(100))});
\fill(P)circle(1.5pt)node[above right]{$P(\psi)$};
\draw[axis](P)--++(.75,-.63)node[right]{$\vct v$};
\draw[guide](P)--(0,{-1.2*(1-cos(100))});
\node[left]at(0,{-1.2*(1-cos(100))}){$y$};
\node[above left]at(0,0){$O$};
\end{tikzpicture}
\caption{摆线上的切向速度及高度；时间积分中的弧长与速率含相同的正弦因子}\label{fig:complete2-cycloid}
\end{figure}
图\ref{fig:complete2-cycloid}在起点有竖直切线，$y'$在该处发散，而$\dd t/\dd\psi=\sqrt{a/g}$有限。欧拉--拉格朗日方程给出内部驻值条件；最速下降问题采用端点间的第一摆线弧，经过后续尖点的多弧延拓不属于同一最短时间分支。



\LectureNeedspace{4\baselineskip}


\subsection{例题与模型变式}

\begin{example}[下降导轨的变分候选（选做）]\label{sup2:problem-exercise-11}\label{sup2:exercise:d2-h11}

取 y 轴竖直向下，质点从 O=(0,0) 静止沿固定光滑双侧导轨下滑，终点 $B=(\pi a,2a)$，$a>0$。见图\ref{sup2:fig:d2-h11-q}。在曲线内部 $y>0$ 的条件下，写出行程时间泛函，求满足变分驻值方程的一支单调下降轨线，并与直线 OB 的时间比较。此题不要求证明全体允许曲线中的全局最小性。

\end{example}
\begin{figure}[H]\centering
\begin{tikzpicture}[scale=.95]
\draw[axis](-.2,0)--(4.1,0)node[right]{$x$};\draw[axis](0,.2)--(0,-2.8)node[below]{$y$（向下）};
\coordinate(B)at(pi,-2);\draw[guide](0,0)--(B);\fill(0,0)circle(1.4pt)node[above left]{$O$};\fill(B)circle(1.5pt)node[right]{$B=(\pi a,2a)$};
\end{tikzpicture}
\caption{两端点固定。虚线表示作为比较的直线导轨，y 正向朝下。}\label{sup2:fig:d2-h11-q}
\end{figure}
\par\medskip

\begin{solution}
由能量守恒，质点在深度 y 处的速率为 $v=\sqrt{2gy}$，弧长微元为 $\sqrt{1+y'^2}\dd x$。因此时间泛函为
\[
T[y]=\frac1{\sqrt{2g}}\int_0^{\pi a}
\sqrt{\frac{1+y'^2}{y}}\,\dd x.
\]
取 $\mathcal L=\sqrt{(1+y'^2)/y}$，它不显含 x。第一积分给出
\[
\mathcal L-y'\mathcal L_{y'}
=\frac1{\sqrt{y(1+y'^2)}}=\text{常量}.
\]
这里沿用第\ref{sec:complete2-variation}节已逐项证明的第一积分。本题对斜率求偏导，有
\[
\mathcal L_{y'}=\frac{y'}{\sqrt y\sqrt{1+y'^2}},\qquad
\mathcal L-y'\mathcal L_{y'}
=\frac{1+y'^2-y'^2}{\sqrt y\sqrt{1+y'^2}}
=\frac1{\sqrt{y(1+y'^2)}}.
\]
代入后即得到上面的常量关系。

把常量写成 $1/\sqrt{2a}$，得到 $y(1+y'^2)=2a$。对下降支取参数 $\varphi$，使
\[
y=a(1-\cos\varphi),\qquad y'=\cot\frac\varphi2.
\]
由 $\dd y/\dd\varphi=a\sin\varphi$，再除以 $y'$，得到
\[
\frac{\dd x}{\dd\varphi}=a\sin\varphi\tan\frac\varphi2
=a(1-\cos\varphi).
\]
沿用前文已计算的原函数，并取 $x(0)=0$，得 $x=a(\varphi-\sin\varphi)$。
在 $0\le\varphi\le\pi$ 内，这支摆线从 O 到 B，且深度单调增加，见图\ref{sup2:fig:d2-h11-s}。

\begin{figure}[H]\centering
\begin{tikzpicture}[scale=.95]
\draw[axis](-.2,0)--(4.1,0)node[right]{$x$};\draw[axis](0,.2)--(0,-2.8)node[below]{$y$（向下）};
\coordinate(B)at(pi,-2);\draw[guide](0,0)--(B);\fill(0,0)circle(1.4pt)node[above left]{$O$};\fill(B)circle(1.5pt)node[right]{$B=(\pi a,2a)$};

\draw[curve,domain=0:180,samples=140]plot({\x*pi/180-sin(\x)},{cos(\x)-1});
\draw[axis]({60*pi/180-sin(60)},{cos(60)-1})--({75*pi/180-sin(75)},{cos(75)-1});
\node at(1.1,-2.15){摆线候选};\node at(2.6,-.65){直线};
\end{tikzpicture}
\caption{摆线候选与直线的几何比较。摆线起始下降更陡，随后趋于水平。}\label{sup2:fig:d2-h11-s}
\end{figure}

沿摆线有 $\dd s=2a\sin(\varphi/2)\dd\varphi$，而速率为 $2\sqrt{ga}\sin(\varphi/2)$，所以
\[
\dd t=\sqrt{\frac ag}\dd\varphi,
\qquad T_{\rm cycloid}=\sqrt{\frac ag}\int_0^\pi\dd\varphi
=\sqrt{\frac ag}[\varphi]_0^\pi=\pi\sqrt{\frac ag}.
\]

直线长为 $L=a\sqrt{\pi^2+4}$，沿线的重力加速度为 $a_\parallel=2g/\sqrt{\pi^2+4}$。从静止开始，沿线速度为 $v(t)=\int_0^t a_\parallel\dd u=a_\parallel t$，路程为 $s(t)=\int_0^t a_\parallel u\dd u=a_\parallel t^2/2$。令 $s(T_{\rm line})=L$，得
\[
T_{\rm line}=\sqrt{\frac ag}\sqrt{\pi^2+4}
>\pi\sqrt{\frac ag}.
\]
因此，这支摆线比直线所需时间短。两条轨线的比较尚不能证明对所有允许曲线的全局最小性。本题的双侧导轨允许反力改变方向，轨线不另受单侧接触条件限制。

为判断起点附近的时间积分是否收敛，将图\ref{sup2:fig:d2-h11-s}中的参数式在 $\varphi\to0^+$ 处展开：
\[
x\sim\frac a6\varphi^3,\quad y\sim\frac a2\varphi^2,
\quad\dd s\sim a\varphi\dd\varphi,\quad v\sim\sqrt{ga}\,\varphi.
\]
于是 $\dd s/v\sim\sqrt{a/g}\dd\varphi$，虽速度在起点为零，时间积分却不发散。求时间时应比较弧长与速率的共同趋零阶数。

\end{solution}



\section{移动约束的反力与分离条件}
\subsection{把轨道的加速度写入法向方程}
固定轨道上的分离由速度和曲率决定。轨道本身加速时，质点的绝对加速度还包含支架的加速度。设光滑轨道只作水平平移，位置为$X(t)$；相对于支架的曲线写成$y=f(\xi)$，于是
\[
 x=X(t)+\xi,\qquad y=f(\xi),\qquad A(t)=\ddot X(t).
\]
取朝轨道上方的单位法向量及沿$\xi$增大方向的单位切向量：
\[
 \vct n=\frac{(-f',1)}{\sqrt{1+f'^2}},\qquad
 \vct t=\frac{(1,f')}{\sqrt{1+f'^2}}.
\]
将坐标对时间求导两次：
\[
 \vct a=(A+\ddot\xi,\ f''\dot\xi^2+f'\ddot\xi).
\]
若轨道从下方支承质点，接触力为$N\vct n$且$N\geq0$。沿两个方向投影得到
\begin{align}
 (1+f'^2)\ddot\xi+f'f''\dot\xi^2&=-A-gf',\\
 N&=\frac{m(g+f''\dot\xi^2-Af')}{\sqrt{1+f'^2}}.
\end{align}
这里$f''$带有符号。凹面朝上时，曲率项提高反力；凸面朝上时，曲率项降低反力。$N=0$以后还应检验自由运动是否朝允许离开的方向发展。
\ExpansionDiagram{fig:dep-d2-moving}{移动曲线上质点的坐标及法向方向}{
\draw[->](-.6,0)--(5.8,0)node[right]{$x$};\draw[->](0,-1.3)--(0,2.7)node[above]{$y$};
\draw[thick,domain=.2:4.7,samples=60]plot(\x,{1.8-.13*(\x-1.5)^2});
\fill(3,1.5075)circle(2pt)node[above left]{$P$};
\draw[->,thick](3,1.5075)--(3.468,2.7075)node[right]{$N\vct n$};\draw[->](3,1.5075)--(3,.5)node[right]{$mg$};
\draw[dashed](1.5,-.4)--(1.5,1.8);\node[below]at(1.5,-.4){$X$};\draw[<->](1.5,-.65)--(3,-.65)node[midway,below]{$\xi$};
\draw[->](.3,2.3)--(1.3,2.3)node[midway,above]{$A>0$};}
图\ref{fig:dep-d2-moving}中的反力垂直于曲线，却不一定垂直于质点的绝对速度。它对质点的瞬时功率为
\[
 \mathcal P_N=N\vct n\cdot\dot{\vct r}
 =-\frac{Nf'\dot X}{\sqrt{1+f'^2}}.
\]
轨道驱动装置通过反力交换能量，所以实验室系中的机械能通常变化。

\subsection{恒加速支架中的第一积分}
当$A$为常量时，支架参考系中的惯性力$-mA\vct e_x$可由势$mA\xi$表示。相对动能包含曲线的度量因子：
\[
 E_{\rm rel}=\frac m2(1+f'^2)\dot\xi^2+mgf(\xi)+mA\xi.
\]
将上式逐项求导。因为 $\dd(f'^2)/\dd t=2f'f''\dot\xi$，有
\[
\begin{aligned}
\frac{\dd E_{\rm rel}}{\dd t}
&=m(1+f'^2)\dot\xi\ddot\xi+mf'f''\dot\xi^3
+mgf'\dot\xi+mA\dot\xi\\
&=m\dot\xi\big[(1+f'^2)\ddot\xi+f'f''\dot\xi^2+gf'+A\big]=0.
\end{aligned}
\]
括号正好是切向方程，故 $E_{\rm rel}$ 守恒。如果 $A$ 随时间变化，求导时还会多出 $m\dot A\,\xi$，于是 $\dd E_{\rm rel}/\dd t=m\dot A\,\xi$；此时应直接积分运动方程，不能套用同一个常量。

\begin{example}[向左加速的凸抛物线支架]
光滑轨道在支架坐标系中为$y=-\xi^2/(2R)$，支架始终以$g/4$向左加速。质量$m$的小珠放在轨道上侧，从顶点以相对速度$u=\sqrt{gR/2}$向右运动。求第一次分离的位置和分离时的相对速度；说明分离前能否出现转向。
\ExpansionDiagram{fig:dep-d2-parabola-q}{小珠从凸抛物线顶点向右滑动，轨道加速度向左}{
\draw[->](-1,0)--(4.4,0)node[right]{$\xi$};\draw[->](0,-2.2)--(0,1.2)node[above]{$y$};
\draw[thick,domain=-.8:3.5,samples=50]plot(\x,{-.13*\x*\x});
\fill(0,0)circle(2pt)node[above left]{$O$};\draw[->,thick](.1,.25)--(1.2,.25)node[right]{$u$};
\draw[->,thick](2,.9)--(.6,.9)node[midway,above]{$A=-g/4$};
\node[right]at(2.7,-1.8){$y=-\xi^2/(2R)$};}
\end{example}
\begin{solution}
图\ref{fig:dep-d2-parabola-q}中从顶点向右的坐标为$\xi>0$。支架系内的惯性力恒为$mg/4$，方向向右，其势能为$-mg\xi/4$。取顶点的总势能为零，重力势能为$-mg\xi^2/(2R)$，相对速率平方为$(1+\xi^2/R^2)\dot\xi^2$。令$s=\xi/R$，相对能量守恒给出
\[
\frac m2(1+s^2)\dot\xi^2-\frac{mgR}{2}s^2-\frac{mgR}{4}s
=\frac m2u^2,
\]
即
\[
 (1+s^2)\dot\xi^2=u^2+gRs^2+\frac{gR}{2}s,
 \qquad
 \frac{\dot\xi^2}{gR}=\frac{1/2+s/2+s^2}{1+s^2}.
\]
对 $s\geq0$，分子严格为正，小珠分离前不会先停下或转向。由轨道方程逐次求导，
\[
\dot y=-\frac\xi R\dot\xi=-s\dot\xi,
\qquad \ddot y=-\dot s\dot\xi-s\ddot\xi
=-\frac{\dot\xi^2}R-s\ddot\xi.
\]
取指向轨道上方的单位法向量 $\vct n=(s,1)/\sqrt{1+s^2}$，相对加速度的法向投影中两项 $s\ddot\xi$ 相消，只剩 $-\dot\xi^2/[R\sqrt{1+s^2}]$。重力与惯性力的合投影为 $mg(s/4-1)/\sqrt{1+s^2}$，故
\[
N+\frac{mg(s/4-1)}{\sqrt{1+s^2}}
=-\frac{m\dot\xi^2}{R\sqrt{1+s^2}}.
\]
将速度关系代入，得到
\[
 \frac{N\sqrt{1+s^2}}{mg}
 =1-\frac{\dot\xi^2}{gR}-\frac s4
 =\frac{2-3s-s^3}{4(1+s^2)}.
\]
起点$N=mg/2>0$。函数$2-3s-s^3$的导数为$-3-3s^2<0$，所以只可能有一次正根。分离位置满足
\[
 s_*^3+3s_*=2.
\]
令 $s_*=2\sinh u$，利用 $\sinh3u=4\sinh^3u+3\sinh u$，左端为 $2\sinh3u$，所以 $\sinh3u=1$，得到
\[
 s_*=2\sinh\!\left(\frac{\operatorname{arsinh}1}{3}\right)
 \simeq0.59607.
\]
\ExpansionDiagram{fig:dep-d2-parabola-a}{分离点处反力降为零；沿切线离开后转入自由运动}{
\draw[->](-.3,0)--(4.2,0)node[right]{$s$};
\draw[->](0,-.7)--(0,1.6)node[above]{$N/(mg)$};
\draw[thick,domain=0:.85,samples=70]plot({4*\x},{2*(2-3*\x-\x^3)/(4*(1+\x^2)^1.5)});
\draw[dashed](2.3843,-.5)--(2.3843,.6);\fill(2.3843,0)circle(2pt)node[above right]{$s_*$};
\node[left]at(0,1){$1/2$};\node[right]at(3.3,-.5){约束延拓};
\draw[->](2.38,-.1)--(3,-.35);}
由图\ref{fig:dep-d2-parabola-a}可见，将接触解延拓到$s>s_*$会得到负反力，轨道需要拉住小珠，违反单侧支承条件。分离后应改用重力下的抛体方程。在分离瞬间，两条轨迹法向加速度相同；再向右，保持接触所需反力将变成向下，实际小珠留在轨道上方，因而分离有效。

利用$N=0$直接消去一次速度计算，得到
\[
 \dot\xi_*=\sqrt{gR(1-s_*/4)}\simeq0.92249\sqrt{gR}.
\]
相对速度沿图\ref{fig:dep-d2-parabola-q}所示曲线切线，大小为
\[
 v_{{\rm rel},*}=\sqrt{1+s_*^2}\,\dot\xi_*
 \simeq1.07394\sqrt{gR}.
\]
题目未给支架初速度，因此绝对速度应保留为
$\vct v_*=(\dot X_*+\dot\xi_*,\,-s_*\dot\xi_*)$。分离位置和相对速度均不依赖$\dot X_*$。
\end{solution}

\section{习题与解答}

\begin{exercise}[顶端有初速的圆弧分离]\label{sup2:problem-exercise-1}\label{sup2:exercise:d2-h1}

沿用例\ref{ex:hw2-1}的圆弧，从顶端以 $u_0=\sqrt{gR}/2$ 向右发射。求分离角、分离速率和分离后的水平、竖直初速度，并判断分离发生在顶端近静止释放的分离点之前还是之后。题设几何见图\ref{sup2:fig:d2-h1-q}。

\end{exercise}
\begin{figure}[H]\centering
\begin{tikzpicture}[scale=.95]
\draw[axis](-.3,0)--(3.1,0)node[right]{$x$};\draw[axis](0,-.2)--(0,2.7)node[above]{$y$};
\draw[curve](2,0)arc(0:100:2);\coordinate(P)at(1.3382612127177165,1.4862896509547885);
\draw[guide](0,0)--(P);\draw(0,.65)arc(90:48:.65);
\node at(.35,.8){$\theta$};\node[below left]at(0,0){$O$};
\fill(P)circle(1.6pt)node[above right]{$P$};\node at(.95,.55){$R$};
\fill(0,2)circle(1.2pt)node[above left]{$A$};\draw[axis](.45,2.35)--(1.2,2.35)node[midway,above]{$u_0$};
\end{tikzpicture}
\caption{顶端以非零切向速度发射，圆心与角度定义保持不变。}\label{sup2:fig:d2-h1-q}
\end{figure}
\par\medskip

\begin{solution}
这里初态已在顶端运动，不涉及无扰动静止平衡。初始支持力为 $mg-mu_0^2/R=3mg/4>0$，所以先有一段接触运动。能量及反力分别为
\[
v^2=gR\left(\frac94-2\cos\theta\right),\qquad
N=mg\left(3\cos\theta-\frac94\right).
\]
随着角度增大，N 由正减到零，得到
\[
\cos\theta_s=\frac34,\qquad
\sin\theta_s=\frac{\sqrt7}4,\qquad
v_s=\frac{\sqrt3}{2}\sqrt{gR}.
\]
图\ref{sup2:fig:d2-h1-s}中的 S 位于右上方，速度沿圆弧切线向右下。分离位置和速度分量为
\[
(x_s,y_s)=\left(\frac{\sqrt7}{4}R,\frac34R\right),
\]
\[
v_{xs}=\frac{3\sqrt3}{8}\sqrt{gR},\qquad
v_{ys}=-\frac{\sqrt{21}}8\sqrt{gR}.
\]
二者平方之和为 $3gR/4$，与分离速率一致。

\begin{figure}[H]\centering
\begin{tikzpicture}[scale=.95]
\draw[axis](-.2,0)--(3.7,0)node[right]{$x$};\draw[axis](0,-.2)--(0,2.6)node[above]{$y$};
\draw[curve](2,0)arc(0:90:2);\coordinate(S)at({sqrt(7)/2},1.5);
\fill(S)circle(1.5pt)node[above right]{$S$};\draw[guide](0,0)--(S);
\draw[axis](S)--++(.65,{- .65*sqrt(7)/3})node[right]{$\vct v_s$};
\draw[guide,domain=0:.9,samples=60]plot({sqrt(7)/2+3*sqrt(6)*\x/8},{1.5-sqrt(42)*\x/8-.5*\x*\x});
\node at(2.8,2.35){$\cos\theta_s=3/4$};
\end{tikzpicture}
\caption{分离点和切向抛出方向。虚线是分离后的自由轨迹局部段。}\label{sup2:fig:d2-h1-s}
\end{figure}

若取分离时刻为 $\tau=0$，自由运动满足 $\dot v_x=0$、$\dot v_y=-g$。先积分得 $v_x=v_{xs}$、$v_y=v_{ys}-g\tau$，再用分离位置作初值，
\[
x-x_s=\int_0^\tau v_{xs}\dd u=[v_{xs}u]_0^\tau=v_{xs}\tau,
\]
\[
y-y_s=\int_0^\tau(v_{ys}-gu)\dd u
=[v_{ys}u-gu^2/2]_0^\tau=v_{ys}\tau-\frac12g\tau^2.
\]
分离点处径向加速度在两模型中连续，但其后的圆弧约束已不再成立。

近静止释放的临界余弦为 $2/3$，而本题为 $3/4>2/3$。在 $0<\theta<\pi/2$ 内，$\cos\theta$ 随 $\theta$ 增大而减小，因此本题分离角较小，沿圆弧更早分离。增加初速提高了同一位置所需的向心加速度，使向外支持力更早降到零。

\end{solution}


\begin{exercise}[曲率改变的双曲余弦轨道]\label{sup2:problem-exercise-3}\label{sup2:exercise:d2-h3}

光滑凸轨道为 $y=H-a\cosh(x/a)$，$a>0$，小球从顶端附近向右静止释放，采用初始偏角趋零的极限。求首次分离点及速率，见图\ref{sup2:fig:d2-h3-q}。

\end{exercise}
\begin{figure}[H]\centering
\begin{tikzpicture}[scale=.95]
\draw[axis](-.25,0)--(3.2,0)node[right]{$x$};\draw[axis](0,-2.2)--(0,.7)node[above]{$y-H+a$};
\draw[curve,domain=0:1.65,samples=90]plot({1.5*\x},{1.1*(1-cosh(\x))});
\fill(0,0)circle(1.5pt)node[above left]{$A$};\node[anchor=west]at(3,-.6){$y=H-a\cosh(x/a)$};
\end{tikzpicture}
\caption{双曲余弦轨道的右支。图中竖直坐标减去顶端高度，只用于显示形状。}\label{sup2:fig:d2-h3-q}
\end{figure}
\par\medskip

\begin{solution}
记 $q=x/a$。轨道导数为 $y'=-\sinh q$、$y''=-\cosh q/a$，因而
\[
\sqrt{1+y'^2}=\cosh q,\qquad
\rho=\frac{(1+y'^2)^{3/2}}{|y''|}=a\cosh^2q.
\]
小球从顶端下降的高度为 $a(\cosh q-1)$，故 $v^2=2ga(\cosh q-1)$。重力在向曲率中心的法线上投影为 $mg/\cosh q$，支持力方向相反，所以
\[
N=\frac{mg}{\cosh q}-\frac{mv^2}{a\cosh^2q}
=mg\frac{2-\cosh q}{\cosh^2q}.
\]
图\ref{sup2:fig:d2-h3-s}按这一表达式画出反力变化。

\begin{figure}[H]\centering
\begin{tikzpicture}[scale=.95]
\draw[axis](-.2,0)--(3.4,0)node[right]{$q=x/a$};\draw[axis](0,-.7)--(0,1.5)node[above]{$N/(mg)$};
\draw[curve,domain=0:1.6,samples=100]plot({1.8*\x},{(2-cosh(\x))/(cosh(\x)*cosh(\x))});
\fill({1.8*ln(2+sqrt(3))},0)circle(1.3pt)node[below]{$\operatorname{arcosh}2$};
\node[left]at(0,1){$1$};\node[anchor=west]at(1,1.7){$N=mg(2-\cosh q)/\cosh^2q$};
\end{tikzpicture}
\caption{按约束延拓计算的支持力。零点之后的负值仅表示原接触模型失效。}\label{sup2:fig:d2-h3-s}
\end{figure}

在右支 $q>0$，$\cosh q$ 严格增大。N 从 mg 开始，首次在 $\cosh q_s=2$ 处为零；继续运动所需的 N 随即变负，确实发生分离。因此
\[
x_s=a\operatorname{arcosh}2=a\ln(2+\sqrt3),\qquad
y_s=H-2a,\qquad v_s=\sqrt{2ga}.
\]
分离点的斜率为 $-\sinh q_s=-\sqrt3$，切线与水平线向下成 $60^\circ$。若需继续求抛体运动，其初速度为
\[
v_{xs}=\frac{v_s}{2},\qquad
v_{ys}=-\frac{\sqrt3}{2}v_s.
\]
分离条件同时用到了当前曲率半径和已下降的高度；只比较轨道斜率与某个固定角度不能得到这一结果。

\end{solution}


\begin{exercise}[由高度控制的内环通行]\label{sup2:problem-exercise-4}\label{sup2:exercise:d2-h4}

质点从高于内环最低点 B 的高度 H 静止释放，经过光滑引导段从 B 沿切线进入半径 R 的竖直圆环内侧，见图\ref{sup2:fig:d2-h4-q}。假定引导段已保证顺利进入。按 H 分类讨论能否保持接触完成一周，并给出中途分离的角度。

\end{exercise}
\begin{figure}[H]\centering
\begin{tikzpicture}[scale=.95]
\coordinate(O)at(2,1.6);\draw[curve](O)circle(1.6);\draw[thick](-1.6,0)--(2,0);
\draw[curve](-1.4,2.8)..controls(-1.2,.7)and(-.5,0)..(1,0);
\fill(-1.4,2.8)circle(1.5pt)node[left]{$A$};\fill(2,0)circle(1.4pt)node[below]{$B$};\fill(O)circle(1.3pt)node[left]{$O$};
\draw[guide](-1.7,2.8)--(-2.1,2.8);\draw[|-|](-2,0)--(-2,2.8)node[midway,left]{$H$};
\draw[guide](O)--(2,3.2)node[midway,right]{$R$};
\end{tikzpicture}
\caption{引导段只用于建立入环速度。高度 H 与圆环半径 R 分别标出。}\label{sup2:fig:d2-h4-q}
\end{figure}
\par\medskip

\begin{solution}
由能量守恒，进入 B 时 $u_B^2=2gH$。进入内环后，用从最低半径起算的角度 $\theta$，
\[
v^2=2g[H-R(1-\cos\theta)],\qquad
N=\frac{mv^2}{R}+mg\cos\theta
=mg\left(\frac{2H}{R}-2+3\cos\theta\right).
\]
内环支持力朝向圆心，其允许条件与绳摆的张力条件相同。因此可直接比较 $v^2=0$、$N=0$ 的先后，但须保留本题的入环高度含义。

\begin{figure}[H]\centering
\begin{tikzpicture}[scale=.95]
\draw[axis](-.15,0)--(7,0)node[right]{$H/R$};\draw[thick](0,.35)--(6.5,.35);
\draw(1.7,.2)--(1.7,.5);\draw(4.25,.2)--(4.25,.5);\node[below]at(1.7,.2){$1$};\node[below]at(4.25,.2){$5/2$};
\node at(.8,1){转回};\node at(2.95,1){中途分离};\node at(5.5,1){完整周转};
\node at(3.2,-.8){边界 $H/R=1$：速度与支持力同时为零};
\end{tikzpicture}
\caption{入环高度决定的三种运动。完成周转的临界高度包含等号。}\label{sup2:fig:d2-h4-s}
\end{figure}

当 $0<H<R$，在 $\cos\theta_t=1-H/R>0$ 处先转向，且 $N_t=mg\cos\theta_t>0$。H=R 时在水平半径处同时有 $v=N=0$；随后沿环返回，支持力重新为正。水平点是这一运动的转向点。

当 $R<H<5R/2$，支持力先在
\[
\cos\theta_s=\frac{2-2H/R}{3}
\]
处为零，且 $v_s^2=-gR\cos\theta_s>0$。继续上行要求轨道向外拉质点，故在该处进入圆环内部的自由抛体运动。这里 $\cos\theta_s<0$，分离点在上半圆。

当 $H\ge5R/2$，环顶速度满足 $v_{\rm top}^2=2g(H-2R)\ge gR$，最小支持力 $N_{\rm top}=mg(2H/R-5)\ge0$，可以完成一周。等号 H=5R/2 时顶点速率非零、N 只瞬时触零，仍能通过。图\ref{sup2:fig:d2-h4-s}将三个高度区间及边界分开。H=0 的质点则停在最低点。

对图\ref{sup2:fig:d2-h4-s}中的两个触零边界，可考察解除约束后的自由运动。在 $H=R$ 的水平转向点，若质点自由下落，经过短时间 $\tau$ 后，到圆心的距离满足
\[
r^2=R^2+\frac14g^2\tau^4>R^2.
\]
在 $H=5R/2$ 的环顶，初始切向速率为 $\sqrt{gR}$；自由抛体同样给出
\[
r^2=gR\tau^2+(R-g\tau^2/2)^2
=R^2+\frac14g^2\tau^4.
\]
内环允许质点处于圆内，不能穿到圆外。因此这两个边界都需轨道继续约束，支持力在触零后重新变正；不能仅由瞬间 $N=0$ 判为分离。

\end{solution}


\begin{exercise}[不等端点质量的离墙角]\label{sup2:problem-exercise-5}\label{sup2:exercise:d2-h5}

沿用例\ref{ex:hw2-1}的装置，但上端 A 的质量为 m，下端 B 的质量为 5m。仍从竖直附近静止倒下，见图\ref{sup2:fig:d2-h5-q}。求保持双接触时的角速度和墙面反力，确定首次离墙角，并求该时刻的地面反力。

\end{exercise}
\begin{figure}[H]\centering
\begin{tikzpicture}[scale=.95]
\coordinate(O)at(0,0);\coordinate(A)at(0,2.7);\coordinate(B)at(2.1,0);
\draw[thick](0,-.3)--(0,3.4);\draw[thick](-.3,0)--(3.5,0);\draw[curve](A)--(B);
\fill(A)circle(2pt)node[left]{$A,m$};\fill(B)circle(2pt)node[below]{$B$};\node[below left]at(0,0){$O$};
\draw(0,2.05)arc(-90:-52.125:.65);\node at(.28,1.8){$\theta$};
\node at(1.85,2.05){$AB=L$};
\node[right]at(2.25,.2){$5m$};
\end{tikzpicture}
\caption{下端质量增大后，动能中的角度系数随构形变化。}\label{sup2:fig:d2-h5-q}
\end{figure}
\par\medskip

\begin{solution}
取 $c=\cos\theta$、$s=\sin\theta$。两端速度关联不变，但质量不同，故
\[
K=\frac12mL^2(s^2+5c^2)\dot\theta^2
=\frac12mL^2(1+4c^2)\dot\theta^2.
\]
势能仍只有上球贡献，记 $A(c)=1+4c^2$，由能量守恒及求导得
\[
\dot\theta^2=\frac{2g}{L}\frac{1-c}{A},\qquad
\ddot\theta=\frac gL\frac{s(1+8c-4c^2)}{A^2}.
\]
求角加速度时，因 $\dot c=-s\dot\theta$，有
\[
\ddot\theta=-\frac s2\frac{\dd(\dot\theta^2)}{\dd c},\qquad
\frac{\dd}{\dd c}\frac{1-c}{A(c)}
=\frac{-A(c)-8c(1-c)}{A(c)^2}
=-\frac{1+8c-4c^2}{A(c)^2}.
\]
分母 $A(c)$ 的导数贡献了上式中的 $8c(1-c)$ 项。

\begin{figure}[H]\centering
\begin{tikzpicture}[scale=.95]
\coordinate(O)at(0,0);\coordinate(A)at(0,2.7);\coordinate(B)at(2.1,0);
\draw[thick](0,-.3)--(0,3.4);\draw[thick](-.3,0)--(3.5,0);\draw[curve](A)--(B);
\fill(A)circle(2pt)node[left]{$A$};\fill(B)circle(2pt)node[below]{$B$};\node[below left]at(0,0){$O$};
\draw(0,2.05)arc(-90:-52.125:.65);\node at(.28,1.8){$\theta$};
\node at(1.85,2.05){$AB=L$};

\draw[axis](A)--(1.05,2.7)node[above]{$N_w$};\draw[axis](B)--(2.1,1.15)node[right]{$N_f$};
\draw[axis](A)--++(0,-.8)node[left]{$mg$};\draw[axis](B)--++(0,-.8)node[below]{$5mg$};
\end{tikzpicture}
\caption{两种反力均由整体动量方程求出，墙面分离时地面反力仍为正。}\label{sup2:fig:d2-h5-s}
\end{figure}

图\ref{sup2:fig:d2-h5-s}中总水平动量为 $5mL c\dot\theta$，因此
\[
N_w=5mL(c\ddot\theta-s\dot\theta^2)
=\frac{5mg\,s}{A^2}(4c^3+3c-2).
\]
把角速度、角加速度代入水平动量导数时，分子的整理过程是
\[
c(1+8c-4c^2)-2(1-c)(1+4c^2)=4c^3+3c-2.
\]
图\ref{sup2:fig:d2-h5-s}中反力表达式的其余因子满足 $s>0$、$A^2>0$。因此，首次离墙的位置由这一多项式首次变号确定，角加速度变号不构成离墙条件。

括号在 $[0,1]$ 严格递增，且在 c=1/2 为零。真实运动中 c 从 1 向下减小，所以墙面力第一次从正到负发生在
\[
\cos\theta_s=\frac12,\quad \theta_s=60^\circ,
\quad \dot\theta_s^2=\frac g{2L},\quad
\ddot\theta_s=\frac{\sqrt3g}{2L}.
\]
此时上端竖直加速度为
\[
a_{Ay}=-L(s\ddot\theta+c\dot\theta^2)=-g.
\]
由整体竖直方程 $N_f-6mg=m a_{Ay}$，得到 $N_f=5mg>0$。

还须确认地面不会在此前先分离。对本题可将地面反力化简为
\[
N_f=5mg\frac{(1-c)^2+10c^2+20c^4}{(1+4c^2)^2}>0
\qquad(0\le c\le1).
\]
因此离墙确实是首个接触改变事件。相同的几何关联下，质量分布改变了动能系数和离墙角，不能直接沿用等质量时的 $\cos\theta_s=2/3$。

\end{solution}


\begin{exercise}[滑动与倾覆的先后]\label{sup2:problem-exercise-8}\label{sup2:exercise:d2-h8}

均匀长方体随水平车厢一起向右加速，保持无相对滑动和无转动。长方体在运动平面内的底面宽为 $2b$、高为 $2h$，静摩擦因数为 $\mu$，见图\ref{sup2:fig:d2-h8-q}。求允许的加速度范围，并比较滑动、倾覆的先后。取 $b=h/2,\mu=3/4$ 时，哪个边界先到达？

\end{exercise}
\begin{figure}[H]\centering
\begin{tikzpicture}[scale=.95]
\draw[thick](-1,0)--(4.4,0);\draw[curve](0,0)rectangle(2,3);
\coordinate(C)at(1,1.5);\fill(C)circle(1.5pt)node[right]{$C$};\draw[guide](C)--(1,0);
\draw[|-|](0,-.45)--(2,-.45)node[midway,below]{$2b$};\draw[|-|](-.5,0)--(-.5,3)node[midway,left]{$2h$};
\draw[axis](2.6,.4)--(4,.4)node[midway,above]{$a$};\node at(3.4,2){车厢地板};
\end{tikzpicture}
\caption{长方体质心位于底面上方 h。加速度沿水平向右。}\label{sup2:fig:d2-h8-q}
\end{figure}
\par\medskip

\begin{solution}
质心平动方程要求 $f=ma$、$N=mg$。取底面中心为水平原点，支持力作用点横坐标为 $x_N$，向右为正。图\ref{sup2:fig:d2-h8-s}中摩擦作用在质心下方，对质心产生转动力矩；支持力必须移向后方才能抵消。
\[
Nx_N+hf=0,\qquad x_N=-\frac{ha}{g}.
\]
接触压力不能为负，因此合力作用点必须位于底面内，$-b\le x_N\le b$。这给出 $a\le gb/h$，而静摩擦给出 $a\le\mu g$。

\begin{figure}[H]\centering
\begin{tikzpicture}[scale=.95]
\draw[curve](0,0)rectangle(2,3);\coordinate(C)at(1,1.5);\fill(C)circle(1.4pt)node[right]{$C$};
\draw[axis](C)--(1,.5)node[right]{$mg$};\draw[axis](.5,0)--(.5,1.2)node[left]{$N$};
\draw[axis](.5,0)--(1.55,0)node[below]{$f$};\draw[guide](1,1.5)--(1,0);
\draw[|-|](.5,-.5)--(1,-.5)node[midway,below]{$ha/g$};
\node[anchor=west]at(2.8,2.6){$N_L+N_R=mg$};\node[anchor=west,align=left]at(2.8,1.35){$N_R\to0$ 时\\绕左边缘倾覆};
\end{tikzpicture}
\caption{支持力向左移位以平衡摩擦力矩。接触合力到达左边缘时开始失去整面接触。}\label{sup2:fig:d2-h8-s}
\end{figure}

也可把支持合力等效分到左、右边缘，写为
\[
N_L+N_R=mg,\qquad b(N_R-N_L)+hma=0.
\]
因此
\[
N_L=\frac{mg}{2}\left(1+\frac{ha}{bg}\right),\qquad
N_R=\frac{mg}{2}\left(1-\frac{ha}{bg}\right).
\]
加速度增大时右边缘载荷先减到零，物体开始绕左侧后缘倾覆，而总支持力仍为 mg。这是接触区域缩到边缘的临界，不是总反力 N=0。

综合两项限制，保持原运动的范围为
\[
0\le a\le g\min\left(\mu,\frac bh\right).
\]
若 $\mu<b/h$，先滑动；若 $\mu>b/h$，先倾覆；两者相等时两个候选边界同时出现，后续运动须联合分析。对给定特例，$b/h=1/2<\mu=3/4$，在 $a=g/2$ 先达到倾覆边界。此时所需摩擦为 mg/2，仍小于 $3mg/4$，证实尚未达到滑动极限。

\end{solution}


\begin{exercise}[固定初速的斜面最大射程]\label{sup2:problem-exercise-9}\label{sup2:exercise:d2-h9}

从倾角 $\alpha\in(0,\pi/2)$ 的斜面上 O 点以固定速率 $v_0$ 向右上方抛出质点，忽略空气阻力。发射角 $\theta$ 从水平线起算，限定 $\alpha<\theta<\pi/2$，见图\ref{sup2:fig:d2-h9-q}。用速度分量的条件极值求沿斜面的最大射程及对应发射角，并检查允许区间的端点。

\end{exercise}
\begin{figure}[H]\centering
\begin{tikzpicture}[scale=.95]
\draw[axis](-.2,0)--(4.1,0)node[right]{$x$};\draw[axis](0,-.2)--(0,2.4)node[above]{$y$};
\pgfmathsetmacro{\flightend}{4*cos(20)/(1+sin(20))}
\draw[thick](0,0)--(3.6,{3.6*tan(20)});\draw[curve,domain=0:\flightend,samples=80]plot(\x,{tan(55)*\x-(\x)^2/(8*cos(55)*cos(55))});
\draw[axis](0,0)--({1.1*cos(55)},{1.1*sin(55)})node[above]{$v_0$};
\draw[domain=0:20]plot({.6*cos(\x)},{.6*sin(\x)});\node at(.8,.13){$\alpha$};
\draw[domain=0:55]plot({1*cos(\x)},{1*sin(\x)});\node at(.97,.53){$\theta$};
\fill(\flightend,{\flightend*tan(20)})circle(1.2pt)node[right]{$B$};\node[below left]at(0,0){$O$};
\end{tikzpicture}
\caption{斜面上的抛射。射程沿 OB 度量，两个角均以水平线为基准。}\label{sup2:fig:d2-h9-q}
\end{figure}
\par\medskip

\begin{solution}
取初速度分量 $v_x>0$、$v_y>v_x\tan\alpha$。按习题\ref{sup2:exercise:d2-h1}中已积分得到的自由轨迹，令起点坐标为零，即有 $x=v_xt,y=v_yt-gt^2/2$。着落斜面要求 $y=x\tan\alpha$，除去 $t=0$，飞行时间为
\[
t_B=\frac{2(v_y-v_x\tan\alpha)}g.
\]
沿斜面射程是 $R=x_B/\cos\alpha$，故
\[
R=\frac{2}{g\cos\alpha}
(v_xv_y-v_x^2\tan\alpha),\qquad v_x^2+v_y^2=v_0^2.
\]
图\ref{sup2:fig:d2-h9-s}中的圆弧就是速度可行集的一部分。只需极大化 $F=v_xv_y-v_x^2\tan\alpha$。

\begin{figure}[H]\centering
\begin{tikzpicture}[scale=.95]
\draw[axis](-.15,0)--(3.5,0)node[right]{$v_x$};\draw[axis](0,-.15)--(0,3)node[above]{$v_y$};
\draw[curve](2,0)arc(0:90:2);\coordinate(P)at({2*cos(55)},{2*sin(55)});\draw[guide](0,0)--(P);
\fill(P)circle(1.5pt)node[above left]{$P_*$};
\draw[axis](P)--++({.8*cos(55)},{.8*sin(55)})node[above]{$\nabla F$};
\draw[guide]($(P)+({-.7*sin(55)},{.7*cos(55)})$)--($(P)+({.7*sin(55)},{-.7*cos(55)})$);
\node at(3,1){$v_x^2+v_y^2=v_0^2$};
\end{tikzpicture}
\caption{速度圆上的驻值条件。极值点处目标函数梯度与圆的法线平行，沿可行切线的一阶变化为零。}\label{sup2:fig:d2-h9-s}
\end{figure}

令 $\nabla F=\lambda\nabla(v_x^2+v_y^2)$，得到
\[
v_y-2v_x\tan\alpha=2\lambda v_x,
\qquad v_x=2\lambda v_y.
\]
由于 $v_x,v_y>0$，可消去 $\lambda$，得到
\[
v_y^2-2v_xv_y\tan\alpha-v_x^2=0.
\]
取正根，
\[
\tan\theta_* =\frac{v_y}{v_x}
=\tan\alpha+\sec\alpha
=\tan\left(\frac\pi4+\frac\alpha2\right).
\]
因此 $\theta_* =\pi/4+\alpha/2$，确在允许区间内。将速度分量写回 $v_x=v_0\cos\theta$、$v_y=v_0\sin\theta$，
\[
R(\theta)=\frac{2v_0^2}{g\cos^2\alpha}
\cos\theta\sin(\theta-\alpha).
\]
其中 $2\cos\theta\sin(\theta-\alpha)=\sin(2\theta-\alpha)-\sin\alpha$。对射程求导，
\[
R'(\theta)=\frac{2v_0^2}{g\cos^2\alpha}\cos(2\theta-\alpha),
\qquad
R''(\theta)=-\frac{4v_0^2}{g\cos^2\alpha}\sin(2\theta-\alpha).
\]
允许区间内只有 $2\theta-\alpha=\pi/2$ 这一个驻点，且二阶导数为负。在最优角处，$\sin(2\theta_* -\alpha)=1$，故
\[
R_{\max}=\frac{v_0^2(1-\sin\alpha)}{g\cos^2\alpha}.
\]
利用 $\cos^2\alpha=(1-\sin\alpha)(1+\sin\alpha)$，可约去分子中的因子。题设 $\alpha<\pi/2$ 保证此因子非零。

因此
\[
\boxed{R_{\max}=\frac{v_0^2}{g(1+\sin\alpha)}.}
\]
当 $\theta$ 趋于 $\alpha$ 时飞行时间趋零，当 $\theta$ 趋于 $\pi/2$ 时水平射程趋零，两个端点都不优于内部解。若 $\alpha\to0$，结果恢复水平地面的 $45^\circ$ 发射和 $v_0^2/g$ 射程。

\end{solution}


\begin{exercise}[松绳后的下落与再次绷紧]\label{deep2:problem-exercise-2}\label{deep2:exercise:d2-h2}

绳摆在右上方 $\theta_0=120^\circ$ 处由静止释放，角度从向下半径起算，绳长 R，见图\ref{deep2:fig:d2-h2-q}。求再次绷紧的位置与时间。规定绳绷紧瞬间消去沿绳向外的速度，切向速度连续；求约束冲量、损失动能和绷紧后的张力。

\end{exercise}
\begin{figure}[H]\centering
\begin{tikzpicture}[scale=.95]
\coordinate(O)at(0,1.8);\draw[guide](O)circle(1.5);
\coordinate(A)at({1.5*sqrt(3)/2},2.55);\coordinate(B)at({1.5*sqrt(3)/2},1.05);
\draw[thick](O)--(A);\draw[guide](A)--(B);\fill(O)circle(1.4pt)node[left]{$O$};
\fill(A)circle(1.5pt)node[above right]{$A$};\fill(B)circle(1.3pt)node[right]{$B$};
\node at(2.8,2.7){$\theta_0=120^\circ$};\node at(2.7,1.9){$OA=R$};
\end{tikzpicture}
\caption{初始位于上半圆的静止质点。虚线 AB 是松绳后保持水平位置不变的下落路径。}\label{deep2:fig:d2-h2-q}
\end{figure}
\par\medskip

\begin{solution}
若初始仍沿圆轨道运动，所需张力为 $T=mg\cos120^\circ=-mg/2$，不允许。因此绳立即松弛。以 O 为原点、向上为 $y$ 正向，释放点为 $(\sqrt3R/2,R/2)$。在习题\ref{sup2:exercise:d2-h1}已积分的自由轨迹中取初速度为零，松绳后
\[
x=\frac{\sqrt3}2R,\qquad y=\frac R2-\frac12gt^2.
\]
再次达到绳长要求 $x^2+y^2=R^2$。除去 $t=0$，另一根为
\[
t_B=\sqrt{\frac{2R}g},\qquad y_B=-\frac R2,
\qquad\theta_B=60^\circ.
\]
在 $0<t<t_B$ 内 $x^2+y^2<R^2$，确实是松绳阶段，没有遗漏更早的绷紧。

\begin{figure}[H]\centering
\begin{tikzpicture}[scale=.95]
\coordinate(O)at(0,1.8);\coordinate(B)at({1.5*sqrt(3)/2},1.05);
\draw[guide](O)circle(1.5);\draw[thick](O)--(B);\fill(O)circle(1.3pt)node[left]{$O$};\fill(B)circle(1.5pt)node[right]{$B$};
\draw[axis](B)--++(0,-1.25)node[below]{$\vct v^-$};
\draw[axis](B)--++(-.65,.375)node[midway,above right]{$\vct J$};
\draw[axis,thick](B)--++(-.54,-.935)node[left]{$\vct v^+$};
\node at(3,2.4){$\vct v^+=v_t^-\vct e_t$};
\end{tikzpicture}
\caption{绷紧瞬间的速度投影。冲量沿 BO，消去向外径向速度，保留切向分量。}\label{deep2:fig:d2-h2-s}
\end{figure}

图\ref{deep2:fig:d2-h2-s}中的绷紧前速度为 $\vct v^-=(0,-\sqrt{2gR})$。在 B 点取
\[
\vct e_r=(\sqrt3/2,-1/2),\qquad
\vct e_t=(1/2,\sqrt3/2),
\]
其中径向向外，切向沿角度增大。投影得
\[
v_r^- =\sqrt{\frac{gR}{2}}>0,\qquad
v_t^-=-\sqrt{\frac{3gR}{2}}.
\]
根据题设的绷紧模型，$v_r^+=0$、$v_t^+=v_t^-$，于是
\[
\vct J=-m\sqrt{\frac{gR}{2}}\,\vct e_r,\qquad
\Delta K_{\rm loss}=\frac12m(v_r^-)^2=\frac14mgR.
\]
把投影结果换回图\ref{deep2:fig:d2-h2-s}的地面坐标，绷紧后有
\[
\vct v^+=v_t^-\vct e_t
=\left(-\sqrt{\frac{3gR}{8}},-3\sqrt{\frac{gR}{8}}\right).
\]
沿绳向内的冲量使水平速度从零变为负值，同时减小竖直速度的大小。碰前动能为 $mgR$，碰后为 $3mgR/4$，损失为 $mgR/4$。冲量消去了径向速度，切向速度则保持不变。

机械能在自由下落阶段守恒，在绳绷紧的冲量阶段减少。

绷紧后的速率平方为 $3gR/2$。沿绳向内列式，
\[
T^+-mg\cos60^\circ=\frac{m(v_t^+)^2}{R},
\qquad T^+=2mg>0.
\]
所以绷紧后的局部圆弧运动满足非负张力条件。冲量给出的速度跳变和之后有限张力维持的圆弧运动，是两个不同阶段。

\end{solution}


\begin{exercise}[端部加重后的危险截面（选做）]\label{deep2:problem-exercise-7}\label{deep2:exercise:d2-h7}

例\ref{ex:hw2-4}的均匀杆顶端另固定一个质量 m 的质点，杆本身质量仍为 m、长度 L，见图\ref{deep2:fig:d2-h7-q}。从竖直附近静止倒下。求截面剪力、弯矩及全过程最大弯矩的位置，比较它与无端部质量时的位置。

\end{exercise}
\begin{figure}[H]\centering
\begin{tikzpicture}[scale=.95]
\coordinate(O)at(0,0);\coordinate(A)at({3.4*sin(32)},{3.4*cos(32)});
\draw[thick](-.5,0)--(3.4,0);\draw[fill=white](-.14,0)--(.14,0)--(0,.22)--cycle;
\draw[line width=2pt](O)--(A);\draw[guide](0,0)--(0,3.2);
\fill(O)circle(1.3pt)node[below]{$O$};\node[right]at(A){$A$};
\coordinate(S)at({1.2*sin(32)},{1.2*cos(32)});\draw[thick]($(S)+(-.22,.14)$)--($(S)+(.22,-.14)$);\node[left,xshift=-4pt]at(S){截面};\node[above right,xshift=3pt]at(S){$S$};
\node at(3.3,3.15){$OA=L$};\node at(3.3,2.2){$OS=\xi L$};\draw(0,.7)arc(90:58:.7);\node at(.18,.53){$\theta$};
\fill(A)circle(3pt)node[above right]{$m$};\node at(3,1.1){杆质量 $m$};
\end{tikzpicture}
\caption{顶端质点与杆刚性连接，改变角加速度及各截面所需传递的惯性载荷。}\label{deep2:fig:d2-h7-q}
\end{figure}
\par\medskip

\begin{solution}
总转动惯量为 $I_O=mL^2/3+mL^2=4mL^2/3$。重力矩为 $(mgL/2+mgL)\sin\theta$，所以
\[
\ddot\theta=\frac{9g}{8L}\sin\theta.
\]
取切口到固定端的距离为$s=\xi L$，其中$0\le\xi\le1$。截取切口以上的杆段并包含端部质点。其总切向惯性力按各部分的质心距离相加，故切面力在 $\vct e_t$ 方向的分量为
\begin{align*}
F_t={}&\ddot\theta\left[m(1-\xi)\frac{L(1+\xi)}2+mL\right]
-mg(2-\xi)\sin\theta\\
={}&\frac{mg\sin\theta}{16}(-5+16\xi-9\xi^2).
\end{align*}
对上段杆的切向惯性项，可从质量分布直接积分：
\[
\int_{\xi L}^{L}s\ddot\theta\,\frac mL\dd s
=\frac{m\ddot\theta}L\left[\frac{s^2}2\right]_{\xi L}^{L}
=\frac{mL}{2}(1-\xi^2)\ddot\theta.
\]
加上端部小球的 $mL\ddot\theta$，就是方括号乘角加速度。向心加速度沿杆轴线，不进入此切向方程；这也是剪力式中没有单独出现 $\dot\theta^2$ 的原因。

\begin{figure}[H]\centering
\begin{tikzpicture}[scale=.95]
\draw[axis](-.1,0)--(4,0)node[right]{$\xi$};\draw[axis](0,-.15)--(0,1.7)node[above]{$M/(mgL\sin\theta)$};
\draw[curve,domain=0:1,samples=100]plot({3.4*\x},{10*\x*(1-\x)*(5-3*\x)/16});
\draw[guide]({3.4*(8-sqrt(19))/9},0)--({3.4*(8-sqrt(19))/9},{10*((8-sqrt(19))/9)*(1-(8-sqrt(19))/9)*(5-3*(8-sqrt(19))/9)/16});
\node[below]at({3.4*(8-sqrt(19))/9},0){$\xi_*$};\node[below]at(3.4,0){$1$};
\draw(-.07,.5)--(.07,.5)node[left,xshift=-3pt]{$1/20$};
\node at(2.4,1.45){$\xi_*=(8-\sqrt{19})/9$};
\end{tikzpicture}
\caption{弯矩分布的形状及内部最大值位置。两端弯矩均为零。}\label{deep2:fig:d2-h7-s}
\end{figure}

按例\ref{ex:hw2-4}相同的符号，以沿杆的长度坐标 $s=\xi L$ 求导，截面平衡给出 $\dd M/\dd s=-F_t$。端点的集中质点可传递力但不传递额外力偶，故 $M(\xi=1)=0$。积分得
\[
M(\xi)=L\int_\xi^1F_t(u)\dd u
=\frac{mgL\sin\theta}{16}\xi(1-\xi)(5-3\xi).
\]
积分的多项式可写成
\[
\int_\xi^1(-5+16u-9u^2)\dd u
=[-5u+8u^2-3u^3]_\xi^1
=\xi(1-\xi)(5-3\xi).
\]
它在两端为零、在内部为正，和图中弯矩曲线的端点及符号相符。

图\ref{deep2:fig:d2-h7-s}给出 $M/(mgL\sin\theta)$ 随截面位置的变化。

M 在杆两端为零，内部极值由剪力为零确定，
\[
9\xi^2-16\xi+5=0.
\]
两根中只有
\[
\xi_* =\frac{8-\sqrt{19}}9\approx0.405
\]
位于 $[0,1]$，另一根大于 1。该点的导数由正变负，故为最大弯矩截面。时间上的最大值在杆水平时取得，数值为
\[
M_{\max}=\frac{mgL}{16}\xi_*(1-\xi_*)(5-3\xi_*)
=\frac{28+19\sqrt{19}}{1944}mgL.
\]
相比无端部质量的 $L/3$，危险截面向顶端移动。端部载荷同时改变重力矩和转动惯量，因此不能只在原弯矩公式上叠加一个静态端部重力矩。

\end{solution}


\begin{exercise}[悬点可移动后的两个自由度（选做）]\label{deep2:problem-exercise-12}\label{deep2:exercise:d2-h12}

将例\ref{ex:hw2-8}的 A、B 改为套在水平杆上的光滑轻环，两环之间连接劲度系数 $k=mg/L$、原长 $D_{\rm rest}=(1+3\sqrt3/2)L$ 的轻弹簧。其余三球质量均为 m，杆长均为 L，只讨论左右对称、$0<\alpha,\beta<\pi/2$ 的构形，见图\ref{deep2:fig:d2-h12-q}。求平衡角、弹簧状态及对称构形内的稳定性，说明为何此时两个角可独立求偏导。

\end{exercise}
\begin{figure}[H]\centering
\begin{tikzpicture}[scale=.95]
\coordinate(A)at(-2.732,0);\coordinate(B)at(2.732,0);\coordinate(C)at(-1.732,-1.732);\coordinate(D)at(1.732,-1.732);\coordinate(E)at(0,-2.732);
\draw[thick](-3.1,0)--(3.1,0);\draw[curve](A)--(C)--(E)--(D)--(B);
\foreach \p in {A,B,C,D,E}{\fill(\p)circle(1.5pt);}
\node[above]at(A){$A$};\node[above]at(B){$B$};\node[left]at(C){$C,m$};\node[right]at(D){$D,m$};\node[below]at(E){$E,m$};
\draw[guide](C)--++(1,0);\draw[domain=0:30]plot({-1.732+.65*cos(\x)},{-1.732-.65*sin(\x)});\node at(-.85,-1.93){$\beta$};
\draw[domain=0:60]plot({-2.732+.7*cos(\x)},{-.7*sin(\x)});\node at(-1.9,-.5){$\alpha$};
\draw[|-|](-2.732,.55)--(2.732,.55)node[midway,above]{$D=2L(\cos\alpha+\cos\beta)$};

\draw[fill=white](A)circle(.11);\draw[fill=white](B)circle(.11);
\draw(-2.6,.18)--(-2.3,.18)--(-2.1,.34)--(-1.7,.02)--(-1.3,.34)--(-.9,.02)--(-.5,.34)--(-.1,.02)--(.3,.34)--(.7,.02)--(1.1,.34)--(1.5,.02)--(1.9,.34)--(2.3,.18)--(2.6,.18);
\end{tikzpicture}
\caption{两悬点可以沿水平杆移动，间距 D 随两角改变，弹簧提供水平约束力。}\label{deep2:fig:d2-h12-q}
\end{figure}
\par\medskip

\begin{solution}
悬点间距不再固定，$D=2L(\cos\alpha+\cos\beta)$ 只是由两角计算弹簧长度的几何关系。整体水平平移不改变势能；固定整体位置后，$\alpha,\beta$ 是两个独立的形状坐标。总势能为
\[
U=-mgL(3\sin\alpha+\sin\beta)
+\frac12k(D-D_{\rm rest})^2.
\]
记压缩弹簧向外推两环的力为 $H=k(D_{\rm rest}-D)$。分别对两个独立角度求偏导，
\[
U_\alpha=-3mgL\cos\alpha+2HL\sin\alpha=0,
\]
\[
U_\beta=-mgL\cos\beta+2HL\sin\beta=0.
\]
由于角度在第一象限，两式要求 H>0，所以平衡时弹簧必为压缩状态。消去 H 后仍有 $\tan\alpha=3\tan\beta$，但另一条方程现在来自弹簧力，不再是固定间距条件。

\begin{figure}[H]\centering
\begin{tikzpicture}[scale=.95]
\coordinate(C)at(0,0);\coordinate(E)at(4.3,-.4);
\fill(C)circle(1.7pt)node[left]{$C$};\draw[axis](C)--(-.65,1.126)node[above]{$T_u$};\draw[axis](C)--(1.126,-.65)node[right]{$T_l$};\draw[axis](C)--(0,-1.15)node[below]{$mg$};
\fill(E)circle(1.7pt)node[below]{$E$};\draw[axis](E)--++(-1.126,.65)node[left]{$T_l$};\draw[axis](E)--++(1.126,.65)node[right]{$T_l$};\draw[axis](E)--++(0,-1.05)node[below]{$mg$};
\node at(2.1,1.6){$T_u\cos\alpha=T_l\cos\beta=H$};
\node at(2,-2.75){$H=k(D_{\rm rest}-D)>0$};
\end{tikzpicture}
\caption{杆力的水平分量由压缩弹簧平衡。中央球与侧球的竖直平衡仍分别成立。}\label{deep2:fig:d2-h12-s}
\end{figure}

图\ref{deep2:fig:d2-h12-s}中各杆传递相同的水平分力 H。代入
\[
\alpha=60^\circ,\qquad\beta=30^\circ
\]
可得 $D=(1+\sqrt3)L$，因此弹簧压缩量为 $\sqrt3L/2$，
\[
H=\frac{mg}{L}\frac{\sqrt3L}{2}=\frac{\sqrt3}{2}mg,
\]
满足两个平衡方程。相应杆力仍为 $T_l=mg,T_u=\sqrt3mg$。

唯一性可由 $\beta$ 检查：由 $\alpha=\arctan(3\tan\beta)$，$\beta$ 增大时 D 严格减小，弹簧压缩力 H 严格增大；平衡所需的 $mg\cot\beta/2$ 严格减小。两者仅有一个交点。

由 $D_\alpha=-2L\sin\alpha$、$D_\beta=-2L\sin\beta$，有 $H_\alpha=2kL\sin\alpha$、$H_\beta=2kL\sin\beta$。继续对势能的一阶偏导求导，
\[
\begin{aligned}
U_{\alpha\alpha}&=3mgL\sin\alpha+2HL\cos\alpha+4kL^2\sin^2\alpha,\\
U_{\beta\beta}&=mgL\sin\beta+2HL\cos\beta+4kL^2\sin^2\beta,\\
U_{\alpha\beta}&=4kL^2\sin\alpha\sin\beta.
\end{aligned}
\]
将三项组成二次型并配方，在平衡点处，总势能的二阶变化为
\begin{align*}
\delta^2U={}&(3mgL\sin\alpha+2HL\cos\alpha)(\delta\alpha)^2\\
&+(mgL\sin\beta+2HL\cos\beta)(\delta\beta)^2\\
&+4kL^2(\sin\alpha\,\delta\alpha+\sin\beta\,\delta\beta)^2.
\end{align*}
前两项的系数都为正，最后一项非负，所以任意非零形状扰动都有 $\delta^2U>0$。平衡在题设的对称构形内稳定。整体水平平移只是中性自由度，不影响这一形状稳定性的判断。

固定悬点时，两角的虚位移受间距条件约束。悬点可移动后，弹簧形变计入势能，两角仍可独立变化，所以应分别令总势能对两角的偏导数为零。

\end{solution}


\chapter{质心运动问题}


合外力决定质心的运动，约束关系和能量方程则用来求系统内部的相对运动。总动能包含质心平动动能与相对动能；系统不断拾取或抛出物质时，还须计入进出物质携带的动量。下面从固定物质系统出发，逐步讨论这些情形。

除相对论部分外，均采用牛顿力学，以地面为惯性参考系。绳链弯折、物质黏附等过程所用的理想模型会在题目中说明；若模型允许动能耗散，就不再额外要求机械能守恒。图中力、速度与动量箭头主要表示方向，长度不统一表示大小。



\section{固定物质系统的质心与整体反冲}

内部相互作用可以改变相对运动，而封闭系统的质心运动由合外力决定。系统包含哪些物质须预先固定，质心平动与相对运动的能量再分别计算。



\subsection{固定物质系统的质心与能量分解}
\subsubsection{先确定系统由哪些物质组成}
设系统含有质量不变的 $N$ 个质点，总质量为 $M$。质心位置、速度与总动量为
\begin{equation}
\vct R=\frac1M\sum_i m_i\vct r_i,\qquad
\vct V=\dot{\vct R},\qquad \vct P=M\vct V.
\label{sup3:eq:d3-cm}
\end{equation}
对连续体，将求和改为积分：$M=\int\dd m$，$\vct R=M^{-1}\int\vct r\,\dd m$。积分必须覆盖同一批物质。静止的绳段、已抛出的燃气若属于所选系统，也必须计入。

在牛顿第三定律适用的经典模型中，内力成对抵消，故
\begin{equation}
M\ddot{\vct R}=\vct F_{\rm ext}.
\label{sup3:eq:d3-cm-eom}
\end{equation}
各质点的速度可以不同。合外力为零时，质心保持匀速直线运动，系统内部仍可振动、转动或发生碰撞。若合外力只在某一方向上为零，动量也只在该方向上守恒。

\begin{figure}[H]\centering
\begin{tikzpicture}[scale=.92]
\coordinate(O)at(0,0);\coordinate(C)at(3,1.3);\coordinate(P)at(5,2.5);
\draw[axis](O)--(6,0)node[right]{$x$};\draw[axis](O)--(0,3.2)node[above]{$y$};
\draw[axis](O)--(C)node[midway,below]{$\vct R$};
\draw[axis](C)--(P)node[midway,above left]{$\vct\rho_i$};
\draw[guide](O)--(P);\node[left]at(O){$O$};
\fill(C)circle(2pt)node[below right]{$C$};\fill(P)circle(2pt)node[right]{$m_i$};
\foreach \p in {(1.8,1.9),(2.2,.7),(4.1,.5),(3.5,2.6)}\fill\p circle(1.2pt);
\draw[axis](3,1.3)--(4.1,1.3)node[below]{$\vct V$};
\end{tikzpicture}
\caption{将质点位置分为质心位置与相对位置。相对位矢的质量加权和为零。}\label{sup3:fig:d3-cm}
\end{figure}

\subsubsection{总动能中还有相对运动}
按图\ref{sup3:fig:d3-cm}写 $\vct r_i=\vct R+\vct\rho_i$，则
$\sum m_i\vct\rho_i=0$，求导后有 $\sum m_i\dot{\vct\rho}_i=0$。展开总动能，交叉项因而消失：
\begin{align}
K&=\frac12\sum_i m_i(\vct V+\dot{\vct\rho}_i)^2\notag\\
 &=\underbrace{\frac12MV^2}_{K_C}
 +\underbrace{\frac12\sum_i m_i\dot{\vct\rho}_i^{\,2}}_{K_{\rm rel}}.
\label{sup3:eq:d3-konig}
\end{align}
对平面运动的刚体，$K_{\rm rel}=I_C\omega^2/2$；对一般质点组，它还可包含振动等内部运动。若只保留位于 C 的等效质点，其动能中便不再包含 $K_{\rm rel}$。

两质点问题尤其适合使用相对坐标。令 $\vct r=\vct r_2-\vct r_1$，约化质量 $\mu=m_1m_2/(m_1+m_2)$，则
\begin{equation}
K=\frac12MV^2+\frac12\mu\dot{\vct r}^{\,2}.
\label{sup3:eq:d3-reduced}
\end{equation}
证明只需将 $\dot{\vct\rho}_1=-m_2\dot{\vct r}/M$、
$\dot{\vct\rho}_2=m_1\dot{\vct r}/M$ 代入式\eqref{sup3:eq:d3-konig}。

\subsubsection{质心动能定理中的位移}
将式\eqref{sup3:eq:d3-cm-eom}与 $\dd\vct R=\vct V\dd t$ 点乘，左侧为 $M\dot{\vct V}\cdot\vct V\dd t=\dd(MV^2/2)$。沿实际质心轨迹积分，得到
\begin{equation}
\Delta K_C=\int\vct F_{\rm ext}\cdot\dd\vct R.
\label{sup3:eq:d3-pseudowork}
\end{equation}
积分使用的是\emph{质心位移}。计算实际外力做功时，应使用各力作用点的位移：
\begin{equation}
W_{\rm ext}=\sum_i\int\vct F_i^{\rm ext}\cdot\dd\vct r_i.
\end{equation}
两式一般不相等。合外力为恒力时，$\vct F_{\rm ext}\cdot\Delta\vct R$ 仍只等于 $\Delta K_C$；系统是否耗散，还要根据实际功与总能量判断。

将各作用点速度分解后，实际外力功率为
\[
\mathcal P_{\rm ext}=\vct F_{\rm ext}\cdot\vct V
 +\sum_i\vct F_i^{\rm ext}\cdot\dot{\vct\rho}_i.
\]
后一项可改变相对运动。内力的功率也未必相消：一对大小相等、方向相反的内力作用于速度不同的两点，其总功率一般不为零。

\begin{figure}[H]\centering
\begin{tikzpicture}[scale=.95]
\draw[curve](0,0)rectangle(4.4,.55);\fill(2.2,.28)circle(1.5pt)node[above=5pt]{$C$};
\fill(4.4,.28)circle(1.5pt)node[below=3pt]{$A$};
\draw[axis](4.4,.28)--(5.7,.28)node[right]{$F$};
\draw[axis](1.3,-.55)--(2.6,-.55)node[midway,below]{$\dd X_C$};
\draw[axis](3.6,-1.25)--(5.4,-1.25)node[midway,below]{$\dd x_A$};
\node[align=left]at(2.5,1.6){质心平动：$F\,\dd X_C=\dd K_C$\\实际输入：$F\,\dd x_A$};
\end{tikzpicture}
\caption{可形变物体的恒力作用点 A 与质心 C 可以有不同位移。实际功与质心平动动能使用不同位移计算。}\label{sup3:fig:d3-work}
\end{figure}

\subsubsection{质心系中的碰撞能量}
在一维两体碰撞中，质心速度由外冲量决定。若外冲量可忽略，$V$ 不变，碰撞损失只来自相对运动。将碰前速度记为$u_1,u_2$，碰后速度记为$v_1,v_2$。恢复系数$e$使$(v_1-v_2)^2=e^2(u_1-u_2)^2$。在式\eqref{sup3:eq:d3-reduced}中相减，质心动能项抵消，留下
\[
K_{\rm before}-K_{\rm after}
=\frac\mu2\big[(u_1-u_2)^2-(v_1-v_2)^2\big].
\]
代入恢复系数关系，得到
\begin{equation}
K_{\rm before}-K_{\rm after}
=\frac12\mu(1-e^2)(u_1-u_2)^2.
\label{sup3:eq:d3-collision-loss}
\end{equation}
完全非弹性碰撞使两物体的相对速度降为零，故在总动量一定时留下的平动动能最小，为 $P^2/(2M)$。若物体碰后还有转动，转动能仍须另计。


\subsection{例题与模型变式}

\begin{example}[光滑墙角中的正方形]\label{ex:hw3-4}

边长$l$、质量$m$的均匀薄正方形，一对相邻顶点分别接触光滑竖直墙和光滑水平地面，正方形在墙角外侧，构形如图\ref{fig:hw3-4-q}。令连接两接触顶点的边与水平夹角为$\phi$，以$\theta=\phi-\pi/4$描述相对对称构形的倾斜。求瞬时角速度为零时的角加速度。

\end{example}


\begin{figure}[H]\centering
\begin{tikzpicture}[>=Stealth,line cap=round,line join=round]
\draw[thick](0,3)--(0,0)--(3.5,0);\coordinate(A)at(0,1.4);\coordinate(B)at(1.9,0);\coordinate(D)at(3.3,1.9);\coordinate(E)at(1.4,3.3);\draw[curve](A)--(B)--(D)--(E)--cycle;\fill(1.65,1.65)circle(2pt)node[right]{$C$};\node[left]at(A){$A$};\node[below]at(B){$B$};\draw[guide](0,0)--(3,3);\draw(1.35,0)arc(180:143.61565:.55)node[midway,above left]{$\phi$};
\end{tikzpicture}
\caption{相邻顶点接触墙与地面的正方形}\label{fig:hw3-4-q}
\end{figure}

\begin{solution}
先按两点都保持接触计算，再判断反力是否允许这一运动。接触点为$A=(0,l\sin\phi)$、$B=(l\cos\phi,0)$，由正方形几何可得
\[
 x_C=y_C=\frac l2(\sin\phi+\cos\phi)=\frac l{\sqrt2}\cos\theta.
\]

\begin{figure}[H]\centering
\begin{tikzpicture}[>=Stealth,line cap=round,line join=round]
\draw[thick](0,3)--(0,0)--(3.5,0);\coordinate(A)at(0,1.4);\coordinate(B)at(1.9,0);\coordinate(D)at(3.3,1.9);\coordinate(E)at(1.4,3.3);\draw[curve](A)--(B)--(D)--(E)--cycle;\fill(1.65,1.65)circle(2pt)node[right]{$C$};\node[left]at(A){$A$};\node[below]at(B){$B$};\draw[axis](0,1.4)--(.7,1.4)node[above]{$a_{Ax}$};\node[left]at(0,2){$N_A=0$};\draw[axis](1.9,0)--(1.9,.8)node[right,yshift=-3pt]{$N_B$};\draw[axis](1.65,1.65)--(1.65,.8)node[left]{$mg$};\draw(1.35,0)arc(180:143.61565:.55)node[midway,above left]{$\phi$};\draw[guide](1.65,0)--(1.65,-.6);\draw[<->](1.65,-.45)--(1.9,-.45)node[midway,below]{$|d|$};
\end{tikzpicture}
\caption{离墙瞬间只保留地面支持力，墙上接触点向右加速}\label{fig:hw3-4-s}
\end{figure}

双接触约束下，$I_C=ml^2/6$，动能与势能为
\[
 T=\frac12ml^2\left(\frac16+\sin^2\theta\right)\dot\theta^2,
 \qquad V=\frac{mgl}{\sqrt2}\cos\theta.
\]
令 $A(\theta)=ml^2(1/6+\sin^2\theta)$，则
\[
\frac{\dd}{\dd t}\frac{\partial T}{\partial\dot\theta}
=A\ddot\theta+2ml^2\sin\theta\cos\theta\,\dot\theta^2,
\quad
\frac{\partial T}{\partial\theta}=ml^2\sin\theta\cos\theta\,\dot\theta^2,
\]
且 $V'=-mgl\sin\theta/\sqrt2$。代入拉格朗日方程，在 $\dot\theta=0$ 时得到形式解
\[
 \beta_{\text{双}}:=\ddot\theta=
 \frac{g\sin\theta}{\sqrt2l(1/6+\sin^2\theta)}.
\]
墙面对薄片的反力应向右；上述约束解要求
\[
 N_A=m\ddot x_C
 =-\frac{ml}{\sqrt2}\sin\theta\,\beta_{\text{双}}
 =-\frac{3mg\sin^2\theta}{1+6\sin^2\theta}<0
 \qquad(\theta\ne0).
\]
非对称构形从静止释放时，墙面接触立即解除。因此，双接触形式解不适用于这一瞬间的实际运动。

离墙后，薄片只受重力和地面支持力，见图\ref{fig:hw3-4-s}。令$d=x_C-x_B=(l/\sqrt2)\sin\theta$为质心相对地面接触点的水平距离，取$\beta=\ddot\theta$；刚体的逆时针角加速度为$-\beta$。由于地面光滑，$\ddot x_C=0$。在当前零角速度瞬间，保持$B$接地给出$\ddot y_C=-d\beta$。竖直平动与绕质心转动方程为
\[
 N_B-mg=-md\beta,\qquad I_C\beta=dN_B.
\]
联立两式，得到实际角加速度
\[
 \boxed{\beta=\frac{mgd}{I_C+md^2}
 =\frac{g\sin\theta}{\sqrt2l(1/6+\tfrac12\sin^2\theta)}}.
\]
此时$N_B=mg/(1+3\sin^2\theta)>0$，地面接触仍然保持。刚体运动关系给出$a_{Ax}=d\beta\ge0$，与$A$向右离墙的运动相符。对称位置$\theta=0$有$\beta=0$，正方形可以静止；它对倾斜扰动不稳定。


\end{solution}



\begin{example}[恒力拉动弹簧两体系统]\label{sup3:problem-example-1}\label{sup3:ex:d3-e1}
光滑水平面上，质量分别为 $m_1,m_2$ 的两小块由劲度系数为 $k$ 的轻弹簧连接，初始静止且弹簧处于原长。自 $t=0$ 起，对右侧小块施加向右恒力 $F$。以初始质心为原点，求两小块的运动，并比较恒力实际做功与质心平动动能的增量。
\end{example}
\begin{figure}[H]\centering
\begin{tikzpicture}
\draw[guide](-.5,0)--(8,0);\draw[curve](0,0)rectangle(1.3,.7);\node at(.65,.35){$m_1$};
\draw[spring](1.3,.35)--(4.8,.35);\node[above]at(3, .6){$k,\ \ell_0$};
\draw[curve](4.8,0)rectangle(6.1,.7);\node at(5.45,.35){$m_2$};
\draw[axis](6.1,.35)--(7.5,.35)node[above]{$F$};
\draw[axis](0,-.7)--(7.5,-.7)node[right]{$x$};\node[below]at(3,-.7){$C_0$};
\end{tikzpicture}
\caption{恒力只作用于右块。弹簧伸长与系统整体平移是两个独立的运动量。}\label{sup3:fig:d3-e1-q}
\end{figure}
\begin{solution}
令 $M=m_1+m_2$，$\mu=m_1m_2/M$。用 $X$ 表示质心坐标，$q=x_2-x_1-\ell_0$ 表示弹簧伸长。这里 $x_1,x_2$ 是两块的质点坐标。初态坐标须满足
$m_1x_{10}+m_2x_{20}=0$ 与 $x_{20}-x_{10}=\ell_0$。故
\[
x_1=X-\frac{m_2}{M}(\ell_0+q),\qquad
x_2=X+\frac{m_1}{M}(\ell_0+q).
\]
图\ref{sup3:fig:d3-e1-s}给出两块相对 C 的速度；对应的相对动量大小相等、方向相反。
\begin{figure}[H]\centering
\begin{tikzpicture}
\draw[axis](-.4,0)--(7.8,0)node[right]{$x$};
\fill(1,0)circle(2pt)node[below]{$m_1$};\fill(5.8,0)circle(2pt)node[below]{$m_2$};
\fill(3,0)circle(1.6pt)node[below]{$C$};
\draw[axis](3,.4)--(4.2,.4)node[right]{$\dot X$};
\draw[axis](1,.7)--(-.1,.7)node[above right]{$-\dfrac{m_2}{M}\dot q$};
\draw[axis](5.8,.7)--(7,.7)node[above left]{$\dfrac{m_1}{M}\dot q$};
\draw[<->](1,-.65)--(5.8,-.65)node[midway,below]{$\ell_0+q$};
\end{tikzpicture}
\caption{质心平移与相对运动的分解。图中相对速度箭头对应 $\dot q>0$ 的时刻。}\label{sup3:fig:d3-e1-s}
\end{figure}
系统合外力为 $F$，所以 $M\ddot X=F$。由零初始条件先积分速度，再积分位置：
\[
\dot X=\frac FM\int_0^t\dd u=\frac{Ft}M,
\qquad X=\frac FM\int_0^tu\dd u=\frac FM[u^2/2]_0^t.
\]
于是
\[
X=\frac{Ft^2}{2M},\qquad \dot X=\frac{Ft}{M}.
\]
两块各自的方程为 $m_1\ddot x_1=kq$、$m_2\ddot x_2=F-kq$。后一式除以 $m_2$，减去前一式除以 $m_1$，得到
\[
\ddot q+\omega^2q=\frac F{m_2},\qquad
\omega^2=k\left(\frac1{m_1}+\frac1{m_2}\right)=\frac k\mu.
\]
图\ref{sup3:fig:d3-e1-q}中，弹簧对两块的力等大反向，但两块速度不同，两力的总功率为 $kq\dot x_1-kq\dot x_2=-kq\dot q$，恰为弹性势能变化率的负值。采用质心和相对坐标后，可分别求解整体平移与弹簧伸缩。

定义 $q_*=F/(m_2\omega^2)=m_1F/(Mk)$，令 $Q=q-q_*$，则 $\ddot Q+\omega^2Q=0$。试取 $Q=\ee^{rt}$，得特征方程 $r^2+\omega^2=0$，根为 $r=\pm\mathrm i\omega$，实通解为 $Q=A\cos\omega t+B\sin\omega t$。初始条件给出 $A=-q_*$、$\omega B=0$，故
\[
q=q_*(1-\cos\omega t),\qquad \dot q=\omega q_*\sin\omega t.
\]
两块的速度也可明确写成
\[
\dot x_1=\frac{Ft}{M}-\frac{m_2}{M}\omega q_*\sin\omega t,
\quad
\dot x_2=\frac{Ft}{M}+\frac{m_1}{M}\omega q_*\sin\omega t.
\]
图\ref{sup3:fig:d3-e1-s}中的相对速度方向每半个周期改变一次；质心速度始终向右增大。

将它与 $X(t)$ 代回坐标式，即得两块的完整运动。

恒力作用点是右块，故实际功为
\[
W=F(x_2-x_{20})=FX+\frac{m_1F}{M}q.
\]
质心动能却只有
\[
K_C=\frac12M\dot X^2=\frac{F^2t^2}{2M}=FX.
\]
余下的功进入相对动能与弹簧势能。将相对运动方程乘以 $\mu\dot q$，
\[
\frac{\dd}{\dd t}\left(\frac12\mu\dot q^2+\frac12kq^2\right)
=\mu\dot q\ddot q+kq\dot q=\frac{m_1F}{M}\dot q.
\]
从初态积分，右侧为 $\frac{m_1F}{M}\int_0^{q(t)}\dd\xi=\frac{m_1F}{M}[\xi]_0^{q(t)}$，有
\[
\frac12\mu\dot q^2+\frac12kq^2=\frac{m_1F}{M}q.
\]
也可代入 $q(t)$ 直接验证：左边为 $kq_*^2(1-\cos\omega t)=kq_*q$。实际功中，$FX$ 增加质心平动动能，其余部分增加相对动能与弹性势能；本题没有耗散。
\end{solution}



\begin{example}[可动斜面的反冲]\label{deep3:problem-example-5}\label{deep3:ex:d3-e5}
质量为 $M$、倾角为 $\alpha\in(0,\pi/2)$ 的斜面放在光滑水平面上。质量为 $m$ 的小块从斜面上静止释放，接触面光滑。斜面向右下方倾斜，初始整个系统静止。设小块相对斜面沿坡向下走过 $s$，下降高度 $h=s\sin\alpha$。求斜面的位移和速度、小块的相对速度与加速度，并求支持力。
\end{example}
\begin{figure}[H]\centering
\begin{tikzpicture}
\draw[curve](-.6,0)--(7.5,0);\draw[curve](0,.05)--(0,3.2)--(5.4,.05)--cycle;
\node at(2,1){$M$};\begin{scope}[shift={(1.7,2.2)},rotate=-30.26]\draw[curve,fill=white](-.35,0)rectangle(.35,.5);\node at(0,.25){$m$};\end{scope}
\draw[guide](1.7,2.2)--(4.1,2.2);\draw[<->](4.2,2.2)--(4.2,.8)node[midway,right]{$h$};
\draw(4.6,.05)arc(180:149.74:.8);\node at(4.2,.42){$\alpha$};
\draw[axis](5.8,-.6)--(7,-.6)node[right]{$x$};
\end{tikzpicture}
\caption{小块向右下方滑动，斜面向左反冲。高度差由相对位移 $s$ 决定。}\label{deep3:fig:d3-e5-q}
\end{figure}
\begin{solution}
取斜面位移 $X$ 向右为正，相对速度 $u=\dot s$ 向坡下为正。小块地面速度的分量为
\[
v_x=\dot X+u\cos\alpha,\qquad v_y=-u\sin\alpha.
\]
图\ref{deep3:fig:d3-e5-s}将水平反冲与相对滑动分别画出。系统在水平方向不受外力，所以
\[
(M+m)\dot X+mu\cos\alpha=0,
\qquad \dot X=-\frac{m\cos\alpha}{M+m}u.
\]
由 $X(0)=s(0)=0$，同式积分给出
\[
X=-\frac{m\cos\alpha}{M+m}\int_0^{s(t)}\dd\xi
=-\frac{m\cos\alpha}{M+m}[\xi]_0^{s(t)}
=-\frac{ms\cos\alpha}{M+m}.
\]
它同时满足质心水平坐标不变的要求。
\begin{figure}[H]\centering
\begin{tikzpicture}
\coordinate(P)at(1.4,2);\draw[curve](-.3,3)--(4.3,.3);
\fill(P)circle(2pt)node[above left]{$P$};
\draw[axis](P)--(2.7,1.24)node[right]{$u$};
\draw[axis](1.4,2.6)--(.1,2.6)node[midway,above]{$\dot X<0$};
\draw[axis](P)--(2.1,3.2)node[right]{$N$};
\draw[axis](P)--(1.4,.65)node[left]{$mg$};
\begin{scope}[shift={(5.1,0)}]
\draw[axis](0,.5)--(3.2,.5);\fill(1.3,.5)circle(1.6pt)node[above]{$C$};
\draw[guide](1.3,0)--(1.3,2.8);\node[align=center]at(1.4,2){$\Delta X_C=0$\\$M\dot X+m v_x=0$};
\end{scope}
\end{tikzpicture}
\caption{斜面反冲使小块的地面水平速度小于相对速度的水平分量；小块受力只有重力与法向支持力。}\label{deep3:fig:d3-e5-s}
\end{figure}
固定地面不做功，两接触面的内力总功率为零，系统机械能守恒。代入速度分量，有
\begin{align*}
K&=\frac12(M+m)\dot X^2+mu\dot X\cos\alpha+\frac12mu^2\\
 &=\frac12m\frac{M+m\sin^2\alpha}{M+m}u^2=mgh.
\end{align*}
令 $D=M+m\sin^2\alpha$，得到
\[
u^2=\frac{2gh(M+m)}D,\qquad
\dot X=-\frac{m\cos\alpha}{M+m}\sqrt{\frac{2gh(M+m)}D}.
\]
因为 $h=s\sin\alpha$，对 $u^2(s)$ 求导可得恒定相对加速度
\[
\begin{aligned}
a_s&=u\frac{\dd u}{\dd s}=\frac12\frac{\dd u^2}{\dd s}
=\frac12\frac{\dd}{\dd s}\frac{2gs\sin\alpha(M+m)}D\\
&=\frac{g(M+m)\sin\alpha}D,\qquad
\ddot X=-\frac{mg\sin\alpha\cos\alpha}D.
\end{aligned}
\]
小块的竖直加速度为 $-a_s\sin\alpha$。竖直受力方程
$N\cos\alpha-mg=-ma_s\sin\alpha$ 给出
\[
N=\frac{mgM\cos\alpha}D>0.
\]
也可由斜面的水平受力 $M\ddot X=-N\sin\alpha$ 验证符号。$M\to\infty$ 时恢复固定斜面的 $u^2=2gh$、$N=mg\cos\alpha$；$M$ 有限时，重力释放的能量同时进入斜面的反冲运动。
相对加速度也可由力方程直接求出，从而避免在释放瞬间约去零速度。图\ref{deep3:fig:d3-e5-s}中沿坡向下投影，小块满足
\[
m(\ddot X\cos\alpha+\dot u)=mg\sin\alpha.
\]
把水平动量关系求导所得 $\ddot X=-m\dot u\cos\alpha/(M+m)$ 代入，得
\[
m\frac{M+m\sin^2\alpha}{M+m}\dot u=mg\sin\alpha.
\]
这与能量法一致，并在 $u=0$ 的释放时刻仍成立。支持力始终为正，因此只要小块尚未滑出斜面，所假定的接触运动就能维持。

\end{solution}



\section{质量输运、动量与能量收支}

提绳、拾取颗粒和火箭喷气都涉及物质穿过所选边界。进出速度必须注明参考系；动量输运与能量收支分别列写，黏附损失不能由动量守恒代替。



\subsection{拾取物质时怎样选择系统}
\subsubsection{固定物质系统与变质量物体}
若研究对象是“当前正在运动的部分”，其成员随时间改变，不能直接把固定系统公式写成 $F=M\dot v$。设该物体质量为 $M$、速度为 $v$，吸收物质的地面速度为 $w$，吸收速率 $\dot M>0$。在 $\dd t$ 内，将物体与将被吸收的微元 $\dd M$ 共同作为固定物质系统：
\[
Mv+w\dd M+F\dd t=(M+\dd M)(v+\dd v).
\]
略去二阶小量后得到
\begin{equation}
M\dot v=F+(w-v)\dot M,
\qquad \frac{\dd(Mv)}{\dd t}=F+w\dot M.
\label{sup3:eq:d3-accretion}
\end{equation}
所有速度都相对同一惯性系。若流入物质原来静止，$w=0$，就有 $\dd(Mv)/\dd t=F$；但加速度仍不等于 $F/M$。

\begin{figure}[H]\centering
\begin{tikzpicture}[scale=.95]
\draw[curve](0,0)rectangle(2,1);\node at(1,.5){$M,\ v$};
\draw[curve](3.2,.2)rectangle(3.7,.8);\node[above]at(3.45,.8){$\dd M$};
\draw[axis](3.7,.5)--(4.6,.5)node[above]{$w$};
\draw[guide](-.2,-.2)rectangle(4.8,1.4);
\node at(2.3,-.65){吸收前：同一批物质};
\begin{scope}[shift={(6,0)}]
\draw[curve](0,0)rectangle(3,1);\node at(1.5,.5){$M+\dd M$};
\draw[axis](3,.5)--(4,.5)node[right]{$v+\dd v$};
\node at(1.5,-.65){吸收后};
\end{scope}
\end{tikzpicture}
\caption{用同一批物质比较吸收前后的动量，流入动量为 $w\dd M$。}\label{sup3:fig:d3-flux}
\end{figure}

\subsubsection{黏附时损失多少机械能}
微元黏附前后若没有可恢复的形变能，损失的机械能为
\begin{equation}
\dd E_{\rm diss}=\frac12(v-w)^2\dd M.
\label{sup3:eq:d3-diss}
\end{equation}
可在物体的瞬时共动惯性系中看出：微元原有相对动能为 $(v-w)^2\dd M/2$，黏附后这一相对运动消失到一阶精度。该损失由相对速度决定，伽利略变换不改变它。

如果物质经过弯折区后改变方向，进入移动部分的速率可能不等于牵引端速度。绳长约束先给出弯折点速度，据此才能确定质量流率。下面的折绳和堆绳虽然都在启动静止绳元，流率却不同。


\subsection{开放系统的能量收支}
\label{sec:complete3}
取质量为 $m$、速度为 $\vct v$ 的本体，外界物质以质量速率 $q$ 被完全黏入，流入前的地面速度为 $\vct u$。短时间内选取本体和即将流入的同一批物质，由动量平衡得到
\[
 m\dot{\vct v}=\vct F+q(\vct u-\vct v),\qquad \dot m=q.
\]
与 $\vct v$ 点乘，并计入 $m$ 的变化：
\[
\begin{aligned}
\frac\dd{\dd t}\left(\frac12mv^2\right)
&=m\dot{\vct v}\cdot\vct v+\frac12\dot m\,v^2\\
&=\vct F\cdot\vct v+q\vct u\cdot\vct v-\frac12qv^2.
\end{aligned}
\]
用 $|\vct u-\vct v|^2=u^2+v^2-2\vct u\cdot\vct v$ 配方，有
\[
 \frac\dd{\dd t}\left(\frac12mv^2\right)
 =\vct F\cdot\vct v+\frac12qu^2
   -\frac12q|\vct u-\vct v|^2.
\]
右边三项分别计入外力对本体做功、流入物质携带的动能，以及黏入过程的非弹性损失。静止物质被拾取时，最后一项为$qv^2/2$。动量与能量的通量须分别列出，二者由不同的守恒关系确定。例\ref{ex:hw3-1}与例\ref{ex:hw3-2}中的弯折区或拾取区，正是这种局部损失发生的位置。

\subsubsection{运动点取矩的两个定义}
对总质量恒定的物质系统，取运动点$Q$，定义采用绝对速度的角动量
\[
 \vct H_Q=\sum_i(\vct r_i-\vct r_Q)\times m_i\vct v_i.
\]
求导时，位置矢量的导数为 $\vct v_i-\vct v_Q$，所以
\[
\begin{aligned}
\dot{\vct H}_Q
&=\sum_i(\vct v_i-\vct v_Q)\times m_i\vct v_i
+\sum_i(\vct r_i-\vct r_Q)\times m_i\vct a_i\\
&=-\vct v_Q\times\vct P+\vct M_Q,
\end{aligned}
\]
其中 $\vct v_i\times\vct v_i=0$，内力矩在本题采用的经典刚体模型中成对抵消。因而
\[
 \dot{\vct H}_Q=\vct M_Q-\vct v_Q\times\vct P.
\]
若另定义 $\widetilde{\vct H}_Q=\sum_i(\vct r_i-\vct r_Q)\times m_i(\vct v_i-\vct v_Q)$，求导中第一组叉乘为零，剩下
\[
\sum_i(\vct r_i-\vct r_Q)\times m_i(\vct a_i-\vct a_Q)
=\vct M_Q-\left[\sum_i m_i(\vct r_i-\vct r_Q)\right]\times\vct a_Q.
\]
方括号等于 $M(\vct R_C-\vct r_Q)$，则
\[
 \dot{\widetilde{\vct H}}_Q
 =\vct M_Q-M(\vct R_C-\vct r_Q)\times\vct a_Q.
\]
质心与固定点可分别使相应附加项消失。材料点某一瞬间静止，只说明$\vct v_Q=0$，不能推出$\vct a_Q=0$。例\ref{ex:hw3-4}应先判定哪些接触能够保持，再选取角动量的参考点。


\subsection{例题与模型变式}

\begin{example}[水平折绳的匀速与恒力牵引]\label{ex:hw3-1}

一根柔软、不可伸长的均匀绳，长$l$、质量$m$，初始拉直放在光滑水平面上。将左端向右牵引，绳在运动端与未动绳段之间折回。不计两段非弯折处的相互作用。(1)绳端以$v_0$匀速运动，求牵引力与端点位移$x$的关系；(2)牵引力为恒力$F_0$，求绳端运动。只讨论拾取完成前。

\end{example}


\begin{figure}[H]\centering
\begin{tikzpicture}[>=Stealth,line cap=round,line join=round]
\draw[guide](0,0)--(5,0);\draw[curve](2,.24)--(1,.24)arc(90:270:.12)--(5,0);\draw[axis](2,.24)--(3,.24)node[above]{$F,v$};\draw[<->](0,.85)--(2,.85)node[midway,above]{$x$};\node[below]at(0,0){初始端点};\node at(3.5,-.7){$\lambda=m/l$};
\end{tikzpicture}
\caption{水平柔绳的运动端与弯折区}\label{fig:hw3-1-q}
\end{figure}

\begin{solution}
绳端前进$x$时，弯折处前进$x/2$，被拉动的绳长为$x/2$，整段速度为$v=\dot x$。图\ref{fig:hw3-1-s}中其质量为$\lambda x/2$，而不是$\lambda x$。

\begin{figure}[H]\centering
\begin{tikzpicture}[>=Stealth,line cap=round,line join=round]
\draw[guide](0,0)--(5,0);\draw[curve](2,.24)--(1,.24)arc(90:270:.12)--(5,0);\draw[axis](2,.24)--(3,.24)node[above]{$F,v$};\draw[<->](0,.85)--(2,.85)node[midway,above]{$x$};\node[below]at(0,0){初始端点};\node at(3.5,-.7){$\lambda=m/l$};\draw[<->](1,1.5)--(2,1.5)node[midway,above]{$x/2$};\draw[guide](1,.24)--(1,1.6);\draw[guide](2,.24)--(2,1.6);\node at(3.8,-1.1){下段速度为零};
\end{tikzpicture}
\caption{绳端前进$x$，弯折区只前进$x/2$}\label{fig:hw3-1-s}
\end{figure}

绳的总动量为$P=\lambda xv/2$，所以
\[
 F=\dot P=\frac\lambda2(v^2+x\dot v).
\]
匀速时$\boxed{F(x)=\lambda v_0^2/2=mv_0^2/(2l)}$，与$x$无关。牵引功率为$\lambda v_0^3/2$，动能增加率为$\lambda v_0^3/4$；另一半在绳元突然改变速度的区域耗散。

恒力时，由 $P(0)=0$，$P(t)=\int_0^t F_0\dd\tau=F_0t$，即 $\lambda xv/2=F_0t$。利用 $v=\dot x$ 分离变量，
\[
\frac\lambda2\int_0^{x(t)}\xi\dd\xi=F_0\int_0^t\tau\dd\tau,
\qquad \frac\lambda2[\xi^2/2]_0^{x(t)}=F_0[\tau^2/2]_0^t.
\]
整理得
\[
 x^2=\frac{2F_0}{\lambda}t^2,\qquad
 \boxed{x(t)=\sqrt{\frac{2F_0l}{m}}t,\quad
 v(t)=\sqrt{\frac{2F_0l}{m}}\quad(t>0)}.
\]
理想模型的初始运动质量为零，启动解因而具有奇性：有限恒力对应非零的右极限速度，与额外指定的$v(0^+)=0$不相容。若端点带质量$\varepsilon>0$，方程改为$(\varepsilon+\lambda x/2)v=F_0t$，启动便连续；令$\varepsilon\to0$才得到上式。拾取完成对应$x/2=l$，即端点已前进$2l$。匀速牵引的适用时间为$0<t<2l/v_0$；恒力解中则为$0<t<2l\sqrt{\lambda/(2F_0)}$。

绳全部运动后应改用整绳的定质量方程。


弯折处的质量流率可以直接给出耗散功率。运动绳段质量$M_v=\lambda x/2$，故$\dot M_v=\lambda v/2$。其动能为$T=\lambda xv^2/4$，于是
\[
 \dot T=\frac\lambda4v^3+\frac\lambda2xv\dot v,
 \qquad Fv=\frac\lambda2v^3+\frac\lambda2xv\dot v,
 \qquad \dot E_{\text{耗}}=\frac\lambda4v^3.
\]
耗散功率等于$\tfrac12\dot M_vv^2$，来自静止绳元并入运动绳段。对匀速或恒力分支，这一结果都与图\ref{fig:hw3-1-s}所示的拾取过程相容；不能对整个过程直接使用机械能守恒。

\end{solution}



\begin{example}[竖直提绳及带球提绳]\label{ex:hw3-2}

均匀柔绳长$l$、质量$m$，最初堆在地面同一点，绳的空间体积忽略。(1)将一端竖直提起：分别求匀速$v_0$所需的力，以及恒力$F_0$时端点的运动规律。(2)若绳端固定一个质量也为$m$的小球，初始给予小球向上速度$v_0$，之后不加牵引力，求小球的最大高度，并求整绳能离地的条件。

\end{example}


\begin{figure}[H]\centering
\begin{tikzpicture}[>=Stealth,line cap=round,line join=round]
\draw[thick](-1,0)--(1.4,0);\draw[curve](0,0)--(0,2.4);
\draw[decorate,decoration={snake,amplitude=2pt,segment length=5pt}](-.6,.05)--(.6,.05);
\draw[<->](.65,0)--(.65,2.4)node[midway,right]{$x$};\draw[axis](0,2.4)--(0,3.3)node[above]{$F$};\node at(0,-.55){无球提绳};\begin{scope}[shift={(5.5,0)}]\draw[thick](-1,0)--(1.4,0);\draw[curve](0,0)--(0,2.4);
\draw[decorate,decoration={snake,amplitude=2pt,segment length=5pt}](-.6,.05)--(.6,.05);
\draw[<->](.65,0)--(.65,2.4)node[midway,right]{$x$};\draw[fill=white](0,2.4)circle(.16);\node[left]at(-.2,2.4){$m$};
\draw[axis](.6,2.6)--(.6,3.3)node[right]{$v$};\node at(0,-.55){带球自由上升};\node at(0,-1.1){$v(0)=v_0$};\end{scope}
\end{tikzpicture}
\caption{无球提绳与带球自由上升采用不同的运动质量和外力}\label{fig:hw3-2-q}
\end{figure}

\begin{solution}
取绳的线密度$\lambda=m/l$，竖直向上为正。图\ref{fig:hw3-2-s}中，运动绳长为$x$，速度为$v=\dot x$；新绳元从地面静止部分进入运动段，质量进入率为$\lambda v$，进入前的动量为零。因此无球时
\[
 F-\lambda xg=\frac{\dd}{\dd t}(\lambda xv)
 =\lambda x\dot v+\lambda v^2.
\]
\begin{figure}[H]\centering
\begin{tikzpicture}[>=Stealth,line cap=round,line join=round]
\draw[thick](-1,0)--(1.4,0);\draw[curve](0,0)--(0,2.4);
\draw[decorate,decoration={snake,amplitude=2pt,segment length=5pt}](-.6,.05)--(.6,.05);
\draw[<->](.65,0)--(.65,2.4)node[midway,right]{$x$};\draw[axis](0,2.4)--(0,3.3)node[above]{$F$};\node at(0,-.55){无球提绳};\node[right]at(0,1.4){$\lambda x$};\draw[axis](-.5,1.9)--(-.5,.9)node[left]{$\lambda xg$};\draw[axis](0,.15)--(0,.8)node[right]{$v$};\node at(.1,-1.1){$\dot m=\lambda v$};\begin{scope}[shift={(5.5,0)}]\draw[thick](-1,0)--(1.4,0);\draw[curve](0,0)--(0,2.4);
\draw[decorate,decoration={snake,amplitude=2pt,segment length=5pt}](-.6,.05)--(.6,.05);
\draw[<->](.65,0)--(.65,2.4)node[midway,right]{$x$};\draw[fill=white](0,2.4)circle(.16);\node[left]at(-.2,2.4){$m$};
\draw[axis](.6,2.6)--(.6,3.3)node[right]{$v$};\node at(0,-.55){带球自由上升};\node[right]at(.9,1.8){$m+\lambda x$};\draw[axis](-.5,1.9)--(-.5,.9)node[left]{$(m+\lambda x)g$};\draw[axis](0,.15)--(0,.8)node[right]{$v$};\node at(.1,-1.1){$\dot m=\lambda v$};\end{scope}
\end{tikzpicture}
\caption{无球提绳与带球自由上升采用不同的运动质量和外力}\label{fig:hw3-2-s}
\end{figure}
匀速提绳时$\dot v=0$，所需拉力为
$\boxed{F=\lambda gx+\lambda v_0^2}$。
前一项克服运动段重力，后一项使新加入的绳元获得动量。

恒力$F=F_0$时，令$w(x)=v^2$，用$\dot v=v\,\dd v/\dd x=w'/2$化为
\[
 w'+\frac2xw=\frac{2F_0}{\lambda x}-2g.
\]
在$x>0$处乘积分因子$x^2$并积分：
\[
 (x^2w)'=\frac{2F_0}{\lambda}x-2gx^2,\qquad
 x^2w=\int\left(\frac{2F_0}{\lambda}x-2gx^2\right)\dd x
 =\frac{F_0}{\lambda}x^2-\frac{2g}{3}x^3+C.
\]
要求$x\to0^+$时速度有限，须取$C=0$。记$A=F_0/\lambda$、$b=2g/3$，在上升支$v=\sqrt{A-bx}$上，由$x(0)=0$，
\[
 t=\int_0^x\frac{\dd\xi}{\sqrt{A-b\xi}}
 =\left[-\frac2b\sqrt{A-b\xi}\right]_0^x
 =\frac2b\left(\sqrt A-\sqrt{A-bx}\right).
\]
解出根式后，
\[
 \boxed{v=\sqrt{F_0/\lambda}-\frac{gt}{3},\qquad
 x=\sqrt{F_0/\lambda}\,t-\frac{gt^2}{6}}.
\]
该式用于$x\le l$且$v\ge0$的拾绳阶段。其初速极限为$\sqrt{F_0/\lambda}$，不能同时任意规定为零；这是$x=0$处运动质量为零的启动奇性。若从有限绳长启动，应使用相应初始条件确定$C$。

带球且无外加拉力时，图中运动质量改为$M(x)=m+\lambda x$。动量方程为
\[
 M\dot v=-Mg-\lambda v^2,\qquad
 w'+\frac{2\lambda}{m+\lambda x}w=-2g.
\]
积分因子由 $z=m+\lambda x$、$\dd z=\lambda\dd x$ 求得：
\[
\int\frac{2\lambda\dd x}{m+\lambda x}=2\int\frac{\dd z}z
=2\ln z+C,
\qquad \exp(2\ln z)=z^2=(m+\lambda x)^2.
\]
从$x=0,w=v_0^2$积分：
\[
 \begin{aligned}
 (m+\lambda x)^2w-m^2v_0^2
 &=-2g\int_0^x(m+\lambda\xi)^2\dd\xi\\
 &=-\frac{2g}{3\lambda}[(m+\lambda\xi)^3]_0^x
 =-\frac{2g}{3\lambda}\big[(m+\lambda x)^3-m^3\big].
 \end{aligned}
\]
所以
\[
 v^2(x)=
 \frac{m^2v_0^2-\frac{2g}{3\lambda}[(m+\lambda x)^3-m^3]}
 {(m+\lambda x)^2}.
\]
若在绳未全部拾起时速度降至零，分子为零。代入$\lambda=m/l$，
\[
 \left(1+\frac hl\right)^3=1+\frac{3v_0^2}{2gl},\qquad
 \boxed{h=l\left[\left(1+\frac{3v_0^2}{2gl}\right)^{1/3}-1\right]}.
\]
这一分支要求$h\le l$，等价于$v_0^2\le14gl/3$。
若$v_0^2>14gl/3$，先在$x=l$处求交接速度：
\[
 v_l^2=\frac{m^2v_0^2-(2gl/3m)(8m^3-m^3)}{4m^2}
 =\frac14\left(v_0^2-\frac{14}{3}gl\right)>0.
\]
此后整绳与球质量恒定，只受重力。设继续上升高度为 $\Delta h$，由 $v\,\dd v=-g\,\dd h$ 积分，$[v^2/2]_{v_l}^{0}=-g[h]_0^{\Delta h}$，得 $\Delta h=v_l^2/(2g)$，故
$\boxed{h=l+v_l^2/(2g)}$。
等号$v_0^2=14gl/3$只表示在最高点恰好拾起最后绳元；要使整绳以非零速度离地，须满足严格不等式。
\end{solution}



\begin{example}[恒力推动拾尘板]\label{ex:hw3-3}

初始质量$M_0$、迎风面积$S$的板，在恒力$F>0$作用下从静止开始向右运动。右侧静止尘埃均匀分布，质量密度为$\rho$，被板完全黏附。忽略其他阻力，求$v(t),x(t)$。

\end{example}


\begin{figure}[H]\centering
\begin{tikzpicture}[>=Stealth,line cap=round,line join=round]
\fill[gray!10](0,0)rectangle(2.2,1.7);\fill[pattern=dots](2.25,0)rectangle(4.8,1.7);
\draw[guide](0,0)--(0,1.7);\draw[fill=white](2,0)rectangle(2.2,1.7);
\draw[axis](-.3,-.15)--(5.1,-.15)node[right]{$x$};\node[below left]at(0,-.15){$0$};
\node at(.9,.8){已扫过};\node at(3.7,2.1){$\rho,S$};\node[above]at(0,1.9){初态$M_0$};
\draw[axis](2.2,1.05)--(3.1,1.05)node[above]{$F$};\draw[<->](0,-.6)--(2.2,-.6)node[midway,below]{$x$};
\node at(2.2,-1.3){$M(x)=M_0+\rho Sx$};
\end{tikzpicture}
\caption{板从尘埃区的边界出发，位移决定已黏附质量}\label{fig:hw3-3-q}
\end{figure}

\begin{solution}
图\ref{fig:hw3-3-s}中板扫过$x$后质量为$M=M_0+qx$，其中$q=\rho S$。每秒黏附质量为$qv$，并非一个预先给定的常量。

\begin{figure}[H]\centering
\begin{tikzpicture}[>=Stealth,line cap=round,line join=round]
\fill[gray!10](0,0)rectangle(2.2,1.7);\fill[pattern=dots](2.25,0)rectangle(4.8,1.7);
\draw[guide](0,0)--(0,1.7);\draw[fill=white](2,0)rectangle(2.2,1.7);
\draw[axis](-.3,-.15)--(5.1,-.15)node[right]{$x$};\node[below left]at(0,-.15){$0$};
\node at(.9,.8){已扫过};\node at(3.7,2.1){$\rho,S$};\node[above]at(0,1.9){初态$M_0$};
\draw[axis](2.2,1.05)--(3.1,1.05)node[above]{$F$};\draw[<->](0,-.6)--(2.2,-.6)node[midway,below]{$x$};
\node at(2.2,-1.3){$M(x)=M_0+\rho Sx$};\draw[guide](2.6,-.05)--(2.6,1.85);\draw[<->](2.2,2.55)--(2.6,2.55)node[midway,above]{$\dd x$};
\node at(4.4,2.8){$\dd m=\rho S\dd x$};
\end{tikzpicture}
\caption{板前方薄层提供下一时刻黏入的尘埃}\label{fig:hw3-3-s}
\end{figure}

进入尘埃的水平动量为零，故 $\dd(Mv)/\dd t=F$。从初态积分，$Mv=\int_0^tF\dd u=Ft$，即 $(M_0+qx)\dot x=Ft$。分离变量并再次积分：
\[
\int_0^{x(t)}(M_0+q\xi)\dd\xi=F\int_0^tu\dd u,
\qquad [M_0\xi+q\xi^2/2]_0^{x(t)}=F[u^2/2]_0^t.
\]
代入上下限，得
\[
 M_0x+\frac q2x^2=\frac F2t^2.
\]
取满足$x(0)=0$的根，得到
\[
 \boxed{x(t)=\frac{\sqrt{M_0^2+qFt^2}-M_0}{q},\qquad
 v(t)=\frac{Ft}{\sqrt{M_0^2+qFt^2}}}.
\]
对速度求导，使用乘积与链式法则，
\[
\dot v=\frac F{\sqrt{M_0^2+qFt^2}}
-\frac{Ft}{2}(M_0^2+qFt^2)^{-3/2}(2qFt)
=\frac{FM_0^2}{(M_0^2+qFt^2)^{3/2}}>0.
\]
速度趋于 $\sqrt{F/q}$。这一极限来自不断拾取静止尘埃所需的动量通量。

展开动量导数也可写成$M\dot v=F-qv^2$。式中的$qv^2$是使新黏附尘埃获得板速所需的动量通量。由$x(t)$消去时间，累计耗散为
\[
 D=Fx-\frac12Mv^2=\frac{qFx^2}{2(M_0+qx)}\ge0.
\]
图\ref{fig:hw3-3-s}中从初始位置到板前缘的扫过长度为$x$，其右侧薄层$\dd x$则提供下一时刻黏入的质量$q\dd x$。

实际外力功为$Fx$，动能为$F^2t^2/(2M)$。瞬时耗散率为$qv^3/2$，可由单个尘埃微元的完全非弹性黏附求得；因此不能直接令$Fx=Mv^2/2$。


\end{solution}



\begin{example}[带水平速度的颗粒流进入小车]\label{sup3:problem-example-8}\label{sup3:ex:d3-e8}
质量 $M_0$ 的小车在光滑水平面上以 $v_0$ 运动。自 $t=0$ 起，颗粒以恒定质量速率 $\sigma$ 落入并留在车中，其入车前的水平地面速度为 $w$。颗粒的竖直运动由竖直冲量处理，与本题水平运动无关。求小车的水平速度、位移，以及因水平相对运动消失而产生的耗散。
\end{example}
\begin{figure}[H]\centering
\begin{tikzpicture}
\draw[curve](.5,.4)--(.5,1.5);\draw[curve](.5,.4)--(4,.4)--(4,1.5);\draw(.9,.15)circle(.2);\draw(3.6,.15)circle(.2);
\foreach \xx in {1.2,2,2.8,3.6}\fill(\xx,2.8)circle(1.5pt);
\draw[axis](2,2.6)--(2.7,1.7);\node[right]at(3.8,2.8){$\sigma,\ w$};
\draw[axis](4.1,.9)--(5.5,.9)node[right]{$v(t)$};\node at(2.2,.9){$M_0+\sigma t$};
\draw[curve](-.3,-.1)--(6,-.1);
\end{tikzpicture}
\caption{颗粒进入前的水平动量通量为 $\sigma w$，并非总是零。}\label{sup3:fig:d3-e8-q}
\end{figure}
\begin{solution}
水平方向无外力，当前车内质量为 $M=M_0+\sigma t$。按图\ref{sup3:fig:d3-e8-s}比较接收前后的同一批物质，有
\[
\frac{\dd(Mv)}{\dd t}=\sigma w.
\]
结合 $Mv|_{t=0}=M_0v_0$，积分得 $Mv-M_0v_0=\int_0^t\sigma w\dd u=\sigma wt$，所以
\[
v(t)=\frac{M_0v_0+\sigma wt}{M_0+\sigma t}
=w+\frac{M_0(v_0-w)}{M_0+\sigma t}.
\]
因此 $v$ 总趋向入流速度 $w$，不一定趋向零。求导有
$\dot v=-\sigma M_0(v_0-w)/(M_0+\sigma t)^2$，也符合
$M\dot v=\sigma(w-v)$。
\begin{figure}[H]\centering
\begin{tikzpicture}
\draw[axis](0,0)--(4.2,0)node[right]{$t$};\draw[axis](0,-.2)--(0,2.8)node[above]{$v$};
\draw[guide](0,1)--(4,1)node[right]{$w$};
\draw[curve,domain=0:4,samples=60]plot(\x,{1+1.4/(1+\x)});\node[left]at(0,2.4){$v_0$};
\begin{scope}[shift={(6,0)}]
\draw[curve](0,.2)rectangle(2.5,1.2);\node at(1.25,.7){小车静止};
\fill(3.5,.7)circle(2pt)node[below]{$\dd M$};\draw[axis](3.5,.7)--(2.65,.7)node[midway,above]{$w-v$};
\node[align=center]at(1.5,-.6){水平相对动能\\$\tfrac12\dd M\,(v-w)^2$};
\end{scope}
\end{tikzpicture}
\caption{左图取 $v_0>w$；右图在小车的瞬时共动惯性系中表示入流的水平相对运动。}\label{sup3:fig:d3-e8-s}
\end{figure}
取 $x(0)=0$，将速度积分。使用 $z=M_0+\sigma u$、$\dd z=\sigma\dd u$，
\[
\begin{aligned}
x(t)&=\int_0^t\left[w+\frac{M_0(v_0-w)}{M_0+\sigma u}\right]\dd u\\
&=wt+\frac{M_0(v_0-w)}\sigma[\ln z]_{M_0}^{M_0+\sigma t}.
\end{aligned}
\]
代入上下限，得到
\[
x(t)=wt+\frac{M_0(v_0-w)}\sigma
\ln\left(1+\frac{\sigma t}{M_0}\right).
\]
耗散可从一次微小黏附求出。图\ref{sup3:fig:d3-e8-s}右图中，质量为 $M$ 的小车与质量 $\dd M$ 的颗粒黏合，相对动能减少
\[
\dd E_{\rm diss,x}=\frac12\frac{M\dd M}{M+\dd M}(v-w)^2
=\frac12\dd M(v-w)^2+O((\dd M)^2).
\]
用 $\dd M=\sigma\dd t$ 除以时间，就得到以下功率；它不依赖选取哪一个匀速运动的惯性系。

水平相对运动造成的耗散功率为
\[
\dot E_{\rm diss,x}=\frac12\sigma(v-w)^2
=\frac{\sigma M_0^2(v_0-w)^2}{2(M_0+\sigma t)^2}.
\]
从零到 $t$ 积分，同样令 $z=M_0+\sigma u$，
\[
\begin{aligned}
E_{\rm diss,x}
&=\frac{\sigma M_0^2(v_0-w)^2}{2}\int_0^t\frac{\dd u}{(M_0+\sigma u)^2}\\
&=\frac{M_0^2(v_0-w)^2}{2}[-z^{-1}]_{M_0}^{M_0+\sigma t}.
\end{aligned}
\]
代入上下限，可得
\[
E_{\rm diss,x}=\frac12M_0(v_0-w)^2
\left(1-\frac{M_0}{M_0+\sigma t}\right).
\]
在地面系中，将初始小车动能与入流颗粒动能相加，再减去末态总动能，同样得到
\[
E_{\rm diss,x}=\frac12M_0v_0^2+\frac12\sigma t\,w^2
-\frac{(M_0v_0+\sigma wt)^2}{2(M_0+\sigma t)}
=\frac{M_0\sigma t}{2(M_0+\sigma t)}(v_0-w)^2.
\]
其中第二项是入流颗粒带来的水平动能。若只取小车前后的动能差，就会漏去这部分能量输入。颗粒竖直落下的能量变化不包含在此式中；$v_0=w$时水平耗散为零，竖直方向仍可能有非弹性损失。

\end{solution}



\begin{example}[恒定喷气率的火箭]\label{ex:hw3-6}

火箭空体质量$M_s$、初始燃料质量$M_f$，每秒喷出燃料质量$\alpha$，喷气相对火箭速度大小恒为$u$，方向向下。初态$v=y=0$。(1)忽略重力，求$v(t)$；(2)在向下均匀重力场$g$中求$v(t)$；(3)求最大上升高度。忽略空气阻力，燃料耗尽后停止喷气。

\end{example}


\begin{figure}[H]\centering
\begin{tikzpicture}[>=Stealth,line cap=round,line join=round]
\draw[curve](0,0)--(.7,0)--(.7,2)--(.35,2.6)--(0,2)--cycle;\node at(.35,1){$M(t)$};\draw[axis](.35,2.6)--(.35,3.3)node[right]{$v$};\draw[axis](.35,0)--(.35,-.8)node[right]{$-u\ \text{（相对火箭）}$};
\end{tikzpicture}
\caption{火箭与向后喷出的燃料}\label{fig:hw3-6-q}
\end{figure}

\begin{solution}
置$M_0=M_s+M_f$、$M(t)=M_0-\alpha t$，燃烧时间$t_b=M_f/\alpha$。图\ref{fig:hw3-6-s}中向下的箭头表示相对火箭的喷气速度$-u$，喷出微元的地面速度为$v-u$；它的符号会随$v$改变。取喷出质量为正值$\alpha\dd t$，故
\[
 M\dot v=\alpha u-Mg.
\]

\begin{figure}[H]\centering
\begin{tikzpicture}[>=Stealth,line cap=round,line join=round]
\draw[curve](0,0)--(.7,0)--(.7,2)--(.35,2.6)--(0,2)--cycle;\node at(.35,1){$M(t)$};\draw[axis](.35,2.6)--(.35,3.3)node[right]{$v$};\draw[axis](.35,0)--(.35,-.8)node[right]{$-u\ \text{（相对火箭）}$};\draw[axis](.8,1.6)--(.8,2.5)node[right]{$\alpha u$};\draw[axis](1.5,1.6)--(1.5,.5)node[right]{$Mg$};\node at(3.7,1){$M=M_0-\alpha t$};
\end{tikzpicture}
\caption{推力、重力与随时间减小的质量}\label{fig:hw3-6-s}
\end{figure}

将质量表达式代入，并从初态积分：
\[
v(t)-v(0)=\alpha u\int_0^t\frac{\dd s}{M_0-\alpha s}-gt
=-u\bigl[\ln(M_0-\alpha s)\bigr]_0^t-gt.
\]
使用$v(0)=0$，忽略重力时$\boxed{v=u\ln[M_0/M(t)]}$；计入重力则
\[
 \boxed{v(t)=u\ln\frac{M_0}{M_0-\alpha t}-gt},\qquad 0\le t\le t_b.
\]
为求高度，将速度再积分一次。令积分变量中的质量为$\widetilde M=M_0-\alpha s$，则
\[
\int_0^t\ln\frac{M_0}{M_0-\alpha s}\,\dd s
=\frac1\alpha\int_{M(t)}^{M_0}\ln\frac{M_0}{\widetilde M}\,\dd\widetilde M
=t-\frac{M(t)}\alpha\ln\frac{M_0}{M(t)}.
\]
这里用了原函数$\widetilde M[\ln(M_0/\widetilde M)+1]$。因此
\[
y(t)=u\left[t-\frac{M(t)}\alpha\ln\frac{M_0}{M(t)}\right]-\frac12gt^2.
\]
若从地面发射且$\alpha u\ge M_0g$，加速度随燃烧增加，燃烧段不会先向下。燃尽时
\begin{align*}
 v_b&=u\ln\frac{M_0}{M_s}-gt_b,\\
 y_b&=u\left[t_b-\frac{M_s}{\alpha}\ln\frac{M_0}{M_s}\right]-\frac12gt_b^2.
\end{align*}
燃烧结束后质量恒为$M_s$，火箭继续作竖直上抛，故
\[
 \boxed{h=y_b+\frac{v_b^2}{2g}}.
\]
若火箭初始放在地面支架上，且允许燃烧时保持原位，则起飞问题还有两种情形：$\alpha u\le M_sg$时燃料耗尽前不能起飞；$M_sg<\alpha u<M_0g$时，起飞发生在
\[
 t_{\rm l}=\frac{M_0-\alpha u/g}{\alpha},\qquad
 M_{\rm l}=\frac{\alpha u}{g}.
\]
把上式中的初始质量改成$M_{\rm l}$，燃烧时间改成$(M_{\rm l}-M_s)/\alpha$，并从起飞时重新计时、计高，即可使用同一组$v_b,y_b,h$公式。起飞前的燃料已被消耗，却没有贡献上升高度。

若火箭从初始时刻起自由运动，推力不足时可以先向下运动，上述$v(t),y(t)$仍成立。这时最高点也能明确判定。燃烧段的加速度为
\[
\dot v=\frac{\alpha u}{M_0-\alpha t}-g,\qquad
\ddot v=\frac{\alpha^2u}{(M_0-\alpha t)^2}>0.
\]
加速度严格增大，故速度至多先减后增。除初态外，燃烧段内若出现$v=0$，必是由下落转为上升的最低点，不是最高点。若$v_b\le0$，整个燃烧段及燃尽后都没有高于初态的位置，故$h=0$。若$v_b>0$，燃尽后上升到$y_b+v_b^2/(2g)$，还须与初态高度零比较。因此不受地面阻挡的自由运动有
\[
\boxed{h=\begin{cases}
0,&v_b\le0,\\
\max\{0,\ y_b+v_b^2/(2g)\},&v_b>0.
\end{cases}}
\]
式中$v_b,y_b$仍取从初始质量$M_0$计算的燃尽值。若存在地面支架，应使用前述起飞条件，不能将先向下的自由运动解用于地面以下。


\end{solution}



\section{固有加速度与相对论推进}

高速推进需要区分地面时间和固有时间，固有加速度也不同于地面系加速度。燃料消耗及航时条件由四动量和快度处理，不能继续使用牛顿速度叠加。



\subsection{固有加速度与相对论推进}
\label{sec:complete3-relativity}
在一维运动中令快度为 $\eta$，定义 $v=c\tanh\eta$，则 $\gamma=\cosh\eta$。地面时间与固有时间满足 $\dd t=\cosh\eta\,\dd\tau$，所以
\[
\frac{\dd v}{\dd t}
=\frac{c\operatorname{sech}^2\eta\,\dd\eta/\dd\tau}{\cosh\eta}
=\frac c{\gamma^3}\frac{\dd\eta}{\dd\tau}.
\]
固有加速度定义为瞬时共动系中测得的加速度，满足 $a_0=\gamma^3\dd v/\dd t$，故
\[
 a_0=c\frac{\dd\eta}{\dd\tau},\qquad
 \frac{\dd v}{\dd t}=\frac{a_0}{\gamma^3},\qquad
 \frac{\dd}{\dd t}(m_0\gamma v)=m_0a_0.
\]
恒定固有加速度对应恒定纵向三维力；加速度变换与动量求导在例\ref{ex:hw3-7}中展开。地面坐标加速度则随速率增大而减小。从初态静止 $\eta(0)=0$ 积分，$\eta=\int_0^\tau(a_0/c)\dd s=a_0\tau/c$。再用 $\dd\tau=(c/a_0)\dd\eta$，
\[
t=\frac c{a_0}\int_0^\eta\cosh u\dd u
=\frac c{a_0}[\sinh u]_0^\eta,
\qquad
x=\frac{c^2}{a_0}\int_0^\eta\sinh u\dd u
=\frac{c^2}{a_0}[\cosh u]_0^\eta.
\]
这里位移积分使用 $\dd x/\dd\tau=v\gamma=c\sinh\eta$，并取 $t(0)=x(0)=0$，得
\[
 \eta=\frac{a_0\tau}{c},\qquad
 t=\frac c{a_0}\sinh\eta,\qquad
 x=\frac{c^2}{a_0}(\cosh\eta-1).
\]
其中$\tau$是船上钟的固有时间，$t$是地面惯性系的时间。

对理想光子火箭，在瞬时共动系中喷出向后的光能$\dd E_\gamma$。保留一阶小量的能量、动量关系为
\[
 -c^2\dd m=\dd E_\gamma,\qquad
 mc\,\dd\eta=\frac{\dd E_\gamma}{c},
 \qquad\frac{\dd m}{m}=-\dd\eta.
\]
对单向加速段分离变量并积分，
\[
\int_{m_i}^{m_f}\frac{\dd m}m=-\int_{\eta_i}^{\eta_f}\dd\eta,
\qquad [\ln m]_{m_i}^{m_f}=-[\eta]_{\eta_i}^{\eta_f}.
\]
向前加速时 $\eta_f>\eta_i$，故 $m_i/m_f=\exp(\eta_f-\eta_i)$；反向发光制动时快度减小，质量仍减少，统一写为 $m_i/m_f=\exp|\Delta\eta|$。加速后还要减速到静止时，质量消耗由两段快度变化的绝对值相加。例\ref{ex:hw3-8}将巡航快度与固有时间限制结合；把出发、制动阶段理想化为瞬时过程所得的是燃料比例下确界，有限推力航程须另外计入两段加减速期间的距离与船上时间。



\LectureNeedspace{4\baselineskip}


\subsection{例题与模型变式}

\begin{example}[恒定固有加速度]\label{ex:hw3-7}

静质量$m_0$保持不变的飞船，从静止沿直线出发，在瞬时共动惯性系中测得的加速度恒为$a_0$。求地面惯性系中的作用力$F(t)$与速度$v(t)$。

\end{example}


\begin{figure}[H]\centering
\begin{tikzpicture}[>=Stealth,line cap=round,line join=round]
\draw[axis](0,0)--(5,0)node[right]{$x$};\draw[curve](2,.3)--(3.5,.3)--(4,.8)--(3.5,1.3)--(2,1.3)--cycle;\node at(2.8,.8){$m_0$};\node at(1,2){地面时$t$};\node at(4,2){船钟$\tau$};
\end{tikzpicture}
\caption{直线运动飞船及两种计时方式}\label{fig:hw3-7-q}
\end{figure}

\begin{solution}
图\ref{fig:hw3-7-s}的横轴采用地面时间$t$的无量纲组合，右侧给出它与飞船固有时$\tau$的关系。先固定当前速率$v$，取此刻的共动惯性系；飞船在该系中获得微小速度增量$\delta v'$后，速度合成式为
\[
v_{\rm new}=\frac{v+\delta v'}{1+v\delta v'/c^2}.
\]
对$\delta v'$求导并取$\delta v'=0$，
\[
\left.\frac{\partial v_{\rm new}}{\partial(\delta v')}\right|_0
=\left.\frac{1+v\delta v'/c^2-(v+\delta v')v/c^2}
{(1+v\delta v'/c^2)^2}\right|_0
=1-\frac{v^2}{c^2}=\frac1{\gamma^2}.
\]
因此一阶增量满足$\dd v=\dd v'/\gamma^2$。在同一瞬间，$\dd v'=a_0\dd\tau$、$\dd t=\gamma\dd\tau$，于是
\[
\dot v=\frac{\dd v'/\gamma^2}{\gamma\dd\tau}
=\frac{a_0}{\gamma^3}.
\]

\begin{figure}[H]\centering
\begin{tikzpicture}[>=Stealth,line cap=round,line join=round]
\draw[axis](0,0)--(4.5,0)node[right]{$a_0t/c$};\draw[axis](0,0)--(0,2.5)node[above]{$v/c$};
\draw[guide](0,2)--(4.2,2);\node[left]at(0,2){$1$};
\draw[curve,domain=0:4.2,samples=90]plot(\x,{2*\x/sqrt(1+\x*\x)});
\node[align=left]at(6.7,1.5){$t=\dfrac c{a_0}\sinh\dfrac{a_0\tau}{c}$\\[8pt]$\dd t=\gamma\dd\tau$};
\node at(2.1,-.65){$F=m_0a_0$};
\end{tikzpicture}
\caption{地面系速度趋于光速，船钟时间由双曲函数参数化}\label{fig:hw3-7-s}
\end{figure}

相对论动量为$p=\gamma m_0v$。先对速度求导：
\[
\frac{\dd\gamma}{\dd v}
=\frac{v}{c^2}(1-v^2/c^2)^{-3/2}
=\frac{\gamma^3v}{c^2},\qquad
\frac{\dd(\gamma v)}{\dd v}
=\gamma+\frac{\gamma^3v^2}{c^2}
=\gamma^3.
\]
最后一步用了$\gamma^{-2}=1-v^2/c^2$。再按链式法则对时间求导，得到
\[
 \boxed{F=\frac{\dd p}{\dd t}=m_0\gamma^3\dot v=m_0a_0}.
\]
从 $p(0)=0$ 积分，$p(t)=m_0a_0\int_0^t\dd s=m_0a_0t$，即 $\gamma v=a_0t$。平方并解出 $v^2$，$v^2=(a_0t)^2[1-v^2/c^2]$；选取向前运动的正根，进而
\[
 \boxed{v(t)=\frac{a_0t}{\sqrt{1+(a_0t/c)^2}}}.
\]
将第\ref{sec:complete3-relativity}节已积分求得的快度关系改写为固有时，得
\[
 t=\frac c{a_0}\sinh\frac{a_0\tau}{c},\qquad
 v=c\tanh\frac{a_0\tau}{c}.
\]
前一式将船钟时间换算为地面时间，后一式给出相应速度。

由 $\sinh\eta=a_0t/c$，得 $\cosh\eta=\sqrt{1+(a_0t/c)^2}$。代回已求出的位移式，
\[
x=\frac{c^2}{a_0}\left[\sqrt{1+(a_0t/c)^2}-1\right],
\qquad \tau=\frac c{a_0}\operatorname{arsinh}\frac{a_0t}c.
\]


地面系中的力保持恒定，共动系中测得的加速度也始终不变。当 $a_0t\ll c$ 时，$v\simeq a_0t$、$x\simeq a_0t^2/2$；当 $a_0t\gg c$ 时，$v$ 趋于 $c$，而地面加速度 $\dot v=a_0[1+(a_0t/c)^2]^{-3/2}$ 趋于零。

\end{solution}



\begin{example}[光子火箭的最低燃料比例]\label{ex:hw3-8}

地球与目标星在同一惯性系中相距$D=4.2$光年。飞船用正反物质完全湮灭产生的定向光辐射推进、减速，出发与到达时均相对两星静止，飞船自身经历的总时间为$\tau_f=42$年。求初始燃料质量占初始总质量的理论下限。假设全部燃料能量都能形成所需方向的光束。

\end{example}


\begin{figure}[H]\centering
\begin{tikzpicture}[>=Stealth,line cap=round,line join=round]
\draw[axis](0,0)--(6,0);\fill(0,0)circle(2pt)node[below]{地球};\fill(5.5,0)circle(2pt)node[below]{目标星};\draw[<->](0,.7)--(5.5,.7)node[midway,above]{$D=4.2\,\mathrm{ly}$};\node at(2.8,-.8){船钟总时$\tau_f=42\,\mathrm{yr}$};
\end{tikzpicture}
\caption{两星间的距离与飞船固有时间}\label{fig:hw3-8-q}
\end{figure}

\begin{solution}
以快度$\eta=\operatorname{artanh}(v/c)$描述运动。在飞船瞬时共动系中，向后发出能量$\dd E$的光，使飞船静质量减少$-\dd M$、快度增加$\dd\eta$。一阶能量与动量守恒给出
\[
 \dd E=-c^2\dd M=Mc^2\dd\eta,\quad
 \frac{\dd M}{M}=-\dd\eta.
\]

\begin{figure}[H]\centering
\begin{tikzpicture}[>=Stealth,line cap=round,line join=round]
\draw[axis](0,0)--(6,0);\fill(0,0)circle(2pt)node[below]{地球};\fill(5.5,0)circle(2pt)node[below]{目标星};\draw[<->](0,.7)--(5.5,.7)node[midway,above]{$D=4.2\,\mathrm{ly}$};\node at(2.8,-.8){船钟总时$\tau_f=42\,\mathrm{yr}$};\draw[axis](0,1.6)--(1.1,1.6)node[above]{加速};\draw[axis](2,1.6)--(3.5,1.6)node[above]{巡航};\draw[axis](5.5,1.6)--(4.4,1.6)node[above]{减速};
\end{tikzpicture}
\caption{理想瞬时加减速之间的巡航段}\label{fig:hw3-8-s}
\end{figure}

一般推进历程中，发射光束使质量变化满足$-\dd\ln M=|\dd\eta|$。从静止加速到$\eta_{\max}$再回到静止，快度的总变差至少为$2\eta_{\max}$，因此
\[
 \ln\frac{M_0}{M_f}=\int|\dd\eta|\ge2\eta_{\max}.
\]
若途中反复加减速，不等式通常为严格不等式；只有沿同一方向单调加速、随后单调减速，才取等号。
总质量比于是满足$M_0/M_f\ge\exp(2\eta_{\max})$。图\ref{fig:hw3-8-s}的巡航段在理想瞬时加、减速极限中占据全部固有时间，且
\[
 D=\int_0^{\tau_f}c\sinh\eta\,\dd\tau
 \le c\tau_f\sinh\eta_{\max}.
\]
故令$q=D/(c\tau_f)=0.1$，有$\eta_{\max}\ge\operatorname{arsinh}q$，从而
\[
 \boxed{\frac{M_{\text{燃料}}}{M_0}\ge
 1-\exp[-2\operatorname{arsinh}(0.1)]
 =1-(\sqrt{1.01}-0.1)^2\approx0.1810}.
\]
相应巡航速度为$v/c=q/\sqrt{1+q^2}=1/\sqrt{101}$。42年按船钟计时。直接取$v=4.2/42=0.1c$等于把它当成地面时间，忽略了两种时间的差别。有限功率的加减速过程要占用固有时间，燃料比例将高于这里的下限。


位移不等式$D\le c\tau_f\sinh\eta_{\max}$要求几乎全部固有时间都保持最大快度。两式同时取等，要求在起点和终点瞬时完成速度改变。所求数值是这一极限的下确界。

\end{solution}



\begin{example}[限定地面航时的光子火箭]\label{deep3:problem-example-10}\label{deep3:ex:d3-e10}
相对地面静止的目的地距地球为 $D$，要求飞船在地面系时间 $T$ 内抵达，抵达时相对目的地静止。飞船干质量为 $M_d$，以完全转化为定向光子的理想燃料推进，不计其他损失。先采用出发瞬间加速、随后匀速、抵达瞬间制动的脉冲模型，求燃料质量与干质量之比。当 $D=4.2$ 光年、$T=42$ 年时求数值。再说明有限推力时该结果的意义。
\end{example}
\begin{figure}[H]\centering
\begin{tikzpicture}
\draw[axis](0,0)--(8,0)node[right]{$x$};\fill(0,0)circle(2pt)node[below]{地球};\fill(7,0)circle(2pt)node[below]{目的地};
\draw[<->](0,-.8)--(7,-.8)node[midway,below]{$D$};
\draw[axis](.3,.8)--(1.7,.8)node[midway,above]{加速};\draw[axis](6.8,.8)--(5.4,.8)node[midway,above]{制动};
\node at(3.5,1.6){匀速 $\beta c$，地面航时 $T$};
\node[left]at(.3,.8){$0$};\node[right]at(6.8,.8){$0$};
\end{tikzpicture}
\caption{两次脉冲改变快度，航程中的匀速段承担全部地面位移。制动时光子发射方向朝前。}\label{deep3:fig:d3-e10-q}
\end{figure}
\begin{solution}
脉冲的时间与位移均取零，因此匀速速度由航程直接确定：
\[
\beta=\frac D{cT},\qquad 0<\beta<1.
\]
可行条件为 $T>D/c$。相应快度为
\[
\eta=\operatorname{artanh}\beta
=\frac12\ln\frac{1+\beta}{1-\beta}.
\]
图\ref{deep3:fig:d3-e10-s}中的两次脉冲分别使快度由 $0$ 变为 $\eta$、再由 $\eta$ 变为 $0$，各需相同质量比 $\exp\eta$。
\begin{figure}[H]\centering
\begin{tikzpicture}
\draw[axis](0,0)--(7.5,0)node[right]{$t$};\draw[axis](0,0)--(0,2.5)node[above]{$\beta$};
\draw[curve](0,0)--(0,1.8)--(6.5,1.8)--(6.5,0);\node[left]at(0,1.8){$D/(cT)$};\node[below]at(6.5,0){$T$};
\node at(3.3,.8){面积 $=D/c$};
\node[above]at(1,2.2){$M_0\to M_1$};\node[above]at(5.9,2.2){$M_1\to M_d$};
\end{tikzpicture}
\caption{脉冲模型的速度历程。航程等于速度曲线与时间轴之间的面积乘以 $c$。}\label{deep3:fig:d3-e10-s}
\end{figure}
若中间质量为 $M_1$，两次守恒关系分别给出
\[
M_0=M_1e^\eta,\qquad M_1=M_de^\eta.
\]
加速时还须携带制动燃料。两段分别从 $M_0$、$M_1$ 开始消耗质量，因此总质量比等于两段质量比的乘积。

消去中间质量 $M_1$，得到
\[
\frac{M_0}{M_d}=\exp(2\eta)=\frac{1+\beta}{1-\beta}.
\]
燃料质量 $M_f=M_0-M_d$，故
\[
\frac{M_f}{M_d}=\frac{2\beta}{1-\beta},\qquad
\frac{M_f}{M_0}=\frac{2\beta}{1+\beta}.
\]
前一比值以干质量为分母，后一比值以出发总质量为分母。代入 $\beta=4.2/42=0.1$，得到
\[
\frac{M_0}{M_d}=\frac{11}9,\qquad
\frac{M_f}{M_d}=\frac29,\qquad
\frac{M_f}{M_0}=\frac2{11}.
\]

有限推力使加速和制动占据时间。由图\ref{deep3:fig:d3-e10-s}的面积关系，令平均速度参数为$\bar\beta=D/(cT)$、峰值速度为$\beta_{\max}c$，任意允许历程都满足
\[
\bar\beta=\frac1T\int_0^T\beta(t)\dd t\le\beta_{\max}.
\]
快度从零到峰值再回到零，总变化量至少为两倍峰值快度；其余往返变化只会继续消耗质量。于是
\[
\frac{M_0}{M_d}\ge
\exp\!\left(2\operatorname{artanh}\beta_{\max}\right)
\ge\frac{1+\bar\beta}{1-\bar\beta}.
\]
要取等号，速度须在除脉冲以外的全部时间内等于平均速度。有限时长的加、减速阶段不能满足这个等号条件。因此$2\bar\beta/(1-\bar\beta)$是燃料与干质量之比的下确界；要确定有限推力下的实际燃料量，还须给出推力或固有加速度的限制。

\end{solution}



\section{流束偏转中的动量通量与功率}
\subsection{流入率由相对速度决定}
质量输运还包括物质流入以后再流出的情形。设一股截面积$S$、密度$\rho$的水平水流以实验室速度$u$射向小车，小车速度为$v<u$。随车控制面在单位时间内截获的体积为$S(u-v)$，故
\[
 \dot m=\rho S(u-v).
\]
喷嘴每秒排出$\rho Su$的水，小车则以较小的流率截获。两者之差积存在喷嘴与小车之间不断增长的流束中，因此车的受力和功率均应按随车控制面的流率计算。

假定车上的光滑导流器使水相对于车保持速率$w=u-v$，出流方向与来流方向夹角为$\theta$。忽略导流器内存水质量，并采用水流穿过导流器所需时间远小于小车速度变化时间的近似，则
\[
 u_{{\rm out},x}=v+w\cos\theta,\qquad
 F_x=\dot m(u-u_{{\rm out},x})
 =\rho S(u-v)^2(1-\cos\theta).
\]
这是水对小车的推力。入、出口压强均取环境压强时，没有额外的压强推力。对完整返向$\theta=\pi$，推力为$2\rho S(u-v)^2$，比同流率下黏附捕获的瞬时推力大一倍。
\ExpansionDiagram{fig:dep-d3-flux}{随车控制面中的流率和出流速度}{
\draw[dashed](-.3,-1.1)rectangle(6.2,1.3);
\draw[->,very thick](-1,.7)--(2.2,.7)node[midway,above]{$u$};
\draw[thick](2.2,.7)..controls(4.4,.7)and(4.4,-.6)..(2.2,-.6);
\draw[->,very thick](2.3,-.6)--(.6,-.6)node[midway,below]{$u_{\rm out}=2v-u$};
\draw[thick](2,-1.3)--(5.2,-1.3);\draw(2.5,-1.5)circle(.2);\draw(4.7,-1.5)circle(.2);
\draw[->](4.5,1)--(5.8,1)node[midway,above]{$v$};
\node[right]at(4.2,-.2){光滑返流器};}
图\ref{fig:dep-d3-flux}中的出流箭头画的是$v<u/2$时的实验室方向；$v>u/2$时水在实验室中仍向右，只是向右速度低于小车，因而相对小车向后流出。

\subsection{水流能量的减少怎样转为车的功}
完整返向时，流入和流出水的实验室速度分别为$u$和$2v-u$。忽略高度差，水流减少的动能率为
\[
 \frac{\dot m}{2}\bigl[u^2-(2v-u)^2\bigr]
 =2\dot m\,v(u-v)=F_xv.
\]
理想导流器中没有黏附损失，这部分能量全部成为小车的机械功。若水与车黏附，车载物质质量增长，除平动动能变化外还有$\dot m(u-v)^2/2$转为内能。两种模型的能量方程应分别建立。

\begin{example}[水流推动返流小车]
质量$M$的小车在无摩擦水平轨道上从静止出发。恒定水流参数为$\rho,S,u$，车上装有图\ref{fig:dep-d3-flux}所示的理想返流器。水不积存在车内。求$v(t)$、位移$x(t)$，以及机械功率相对于喷嘴出流动能功率的最大比例。取$x(0)=0$。
\ExpansionDiagram{fig:dep-d3-cart-q}{小车的质量保持不变，推力随追上水流而减小}{
\draw[thick](0,0)--(6,0);\draw[thick](2,.3)rectangle(4,1.2);
\draw(2.4,.15)circle(.15);\draw(3.6,.15)circle(.15);\node at(3,.75){$M$};
\draw[->,thick](.1,.8)--(1.8,.8)node[midway,above]{$u$};
\draw[->,thick](4.1,.8)--(5.8,.8)node[midway,above]{$F_x$};\node[above]at(3,1.4){$v(0)=0$};}
\end{example}
\begin{solution}
图\ref{fig:dep-d3-cart-q}中车的质量没有增长，所以用$M\dot v=F_x$，不再加$\dot Mv$。代入截获流率，
\[
 M\frac{\dd v}{\dd t}=2\rho S(u-v)^2.
\]
令 $w=u-v$，则 $\dot w=-2\rho Sw^2/M$。分离变量，使用 $w(0)=u$，
\[
\int_u^{w(t)}\frac{\dd\zeta}{\zeta^2}
=-\frac{2\rho S}{M}\int_0^t\dd\tau,
\qquad [-\zeta^{-1}]_u^{w(t)}=-\frac{2\rho S}{M}t.
\]
因此
\[
 \frac1w-\frac1u=\frac{2\rho S}{M}t.
\]
定义 $t_c=M/(2\rho Su)$，则 $w=u/(1+t/t_c)$。由 $v=u-w$ 恢复速度，位移则由
\[
x=u\int_0^t\left(1-\frac{t_c}{\tau+t_c}\right)\dd\tau
=u[\tau-t_c\ln(\tau+t_c)]_0^t
\]
求得。代入上下限，
\[
 \boxed{v(t)=u\frac{t}{t+t_c}},\qquad
 \boxed{x(t)=u\left[t-t_c\ln(1+t/t_c)\right]}.
\]
初始加速度为$u/t_c$；随着$t$增大，$u-v$下降，小车在有限时刻始终达不到$u$。

\ExpansionDiagram{fig:dep-d3-cart-a}{返流器的机械功率在$v/u=1/3$处最大}{
\draw[->](0,0)--(5.8,0)node[right]{$z=v/u$};\draw[->](0,0)--(0,2.4)node[above]{$P/P_0$};
\draw[thick,domain=0:1,samples=70]plot({5*\x},{3.5*4*\x*(1-\x)^2});
\draw[dashed](1.6667,0)--(1.6667,2.0741)--(0,2.0741);
\node[below]at(1.6667,0){$1/3$};\node[below]at(5,0){$1$};\node[left]at(0,2.0741){$16/27$};}
喷嘴的恒定动能输出功率为$P_0=\rho Su^3/2$。令$z=v/u$，其中$0\le z<1$。车获得的功率由图\ref{fig:dep-d3-flux}中出、入流能量差给出：
\[
 P=F_xv=2\rho S(u-v)^2v,\qquad
 \eta(z)=\frac P{P_0}=4z(1-z)^2.
\]
求导得$\eta'(z)=4(1-z)(1-3z)$。在$0\leq z<1$内，唯一的内部极大值位于$z=1/3$，故
\[
 \boxed{\eta_{\max}=\frac{16}{27}},\qquad t_{\max}=\frac{t_c}{2}.
\]
图\ref{fig:dep-d3-cart-a}画的是相对于喷嘴功率的比例。若只以车正在截获的水流入射动能率$\dot m u^2/2$作分母，则比例为$4z(1-z)$，其极大值为1，位于$z=1/2$。此时出流水在实验室中瞬时静止。这两个分母对应不同的控制面，不能混用。
\end{solution}

\section{习题与解答}

\begin{exercise}[非均匀杆的质心与偏心受力]\label{sup3:problem-exercise-1}\label{sup3:ex:d3-h1}
长 $L$ 的细杆沿 $x$ 轴放置，线密度为 $\lambda(x)=\lambda_0(1+x/L)$，$0\le x\le L$。求总质量 $M$、质心坐标及绕质心垂直杆的转动惯量。杆在光滑水平面内自由运动，初始静止；若在右端沿平面内垂直于杆的方向施加恒力 $F$，求初始质心加速度与角加速度。力改作用于何处可使杆保持纯平移？
\end{exercise}
\begin{figure}[H]\centering
\begin{tikzpicture}
\draw[axis](-.5,0)--(7.6,0)node[right]{$x$};\draw[curve,line width=3pt](0,.35)--(6,.35);\node[below]at(0,0){$0$};\node[below]at(6,0){$L$};
\draw[axis](6,.35)--(6,1.1)node[above]{$F$};\node at(2.4,1.3){$\lambda(x)=\lambda_0(1+x/L)$};
\draw[axis](-.2,2.35)--(7,2.35)node[right]{$x$};
\draw[axis](0,2.35)--(0,4.15)node[above]{$\lambda$};
\draw[guide](0,3)--(6,3);\draw[guide](0,3.65)--(6,3.65)--(6,2.35);
\draw[curve](0,3)--(6,3.65);\fill(0,3)circle(1pt);\fill(6,3.65)circle(1pt);
\node[left]at(0,3){$\lambda_0$};\node[left]at(0,3.65){$2\lambda_0$};
\node[below left]at(0,2.35){$0$};\node[below]at(6,2.35){$L$};
\end{tikzpicture}
\caption{上图为线密度分布，$\lambda(L)=2\lambda(0)$；下图标出细杆及右端受力。}\label{sup3:fig:d3-h1-q}
\end{figure}
\begin{solution}
图\ref{sup3:fig:d3-h1-q}中质量元为 $\dd m=\lambda_0(1+x/L)\dd x$，所以
\[
M=\lambda_0\int_0^L\left(1+\frac xL\right)\dd x
=\lambda_0[x+x^2/(2L)]_0^L=\frac32\lambda_0L.
\]
一阶质量矩为
\[
\int_0^Lx\dd m=\lambda_0\int_0^L(x+x^2/L)\dd x
=\lambda_0[x^2/2+x^3/(3L)]_0^L=\frac56\lambda_0L^2,
\]
故
\[
x_C=\frac1M\int_0^Lx\dd m=\frac59L.
\]
所得 $x_C>L/2$，与右端较密的分布相符。绕左端的惯量
\[
I_O=\lambda_0\int_0^Lx^2\left(1+\frac xL\right)\dd x
=\lambda_0[x^3/3+x^4/(4L)]_0^L=\frac7{12}\lambda_0L^3.
\]
用平行轴定理减去 $Mx_C^2$，得
\[
I_C=\left(\frac7{12}-\frac{25}{54}\right)\lambda_0L^3
=\frac{13}{108}\lambda_0L^3=\frac{13}{162}ML^2.
\]
\begin{figure}[H]\centering
\begin{tikzpicture}
\draw[curve,line width=2pt](0,0)--(6,0);\fill(3.333,0)circle(2.4pt)node[below]{$C$};
\draw[axis](6,0)--(6,1.5)node[above]{$F$};\draw[axis](3.333,.3)--(3.333,1.25)node[left]{$a_C$};
\draw[<->](0,-.65)--(3.333,-.65)node[midway,below]{$5L/9$};\draw[<->](3.333,-.65)--(6,-.65)node[midway,below]{$4L/9$};
\draw[axis](4.2,1.5)arc(0:110:.55);\node at(3.8,2.4){$\alpha$};
\end{tikzpicture}
\caption{受力点相对质心的力臂为 $4L/9$，不是 $L/2$。}\label{sup3:fig:d3-h1-s}
\end{figure}
按图\ref{sup3:fig:d3-h1-s}，质心沿 $F$ 方向的加速度为 $a_C=F/M$，力对 C 的逆时针力矩为 $F(4L/9)$，因此
\[
\alpha=\frac{4FL}{9I_C}=\frac{72F}{13ML}.
\]
若力始终作用在杆上 $x=5L/9$ 的质心处，力矩为零，初始角速度也为零，杆就保持纯平移。质心加速度由合力确定，是否转动则由关于质心的力矩确定。
图\ref{sup3:fig:d3-h1-s}中的力向上，对 C 的力矩沿逆时针方向，初始角加速度也沿逆时针方向。力移到几何中点仍有非零力臂，只有移到质量中心才能从静止保持纯平移。

\end{solution}


\begin{exercise}[抛体在最高点分裂]\label{sup3:problem-exercise-2}\label{sup3:ex:d3-h2}
总质量 $3m$ 的小球在抛体最高点、距地面高度 $H$ 处以水平速度 $U$ 向右运动，瞬间分裂为质量 $m$、$2m$ 的两块。轻块分裂后相对地面瞬时静止，两块随后仅受重力。以分裂点正下方为地面坐标原点，求两块落点、分裂增加的动能，并验证落地前系统质心的轨迹。
\end{exercise}
\begin{figure}[H]\centering
\begin{tikzpicture}
\draw[axis](-.4,0)--(7.7,0)node[right]{$x$};\draw[axis](0,-.2)--(0,4)node[above]{$y$};
\fill(0,3)circle(2pt)node[above left]{$3m$};\draw[axis](.2,3)--(1.7,3)node[above]{$U$};
\draw[guide](0,3)--(0,0);\draw[<->](-.9,0)--(-.9,3)node[midway,left]{$H$};
\draw[curve,domain=0:4,samples=60]plot(\x,{3-.1875*(\x)^2});\node at(3,3.5){分裂时 $t=0$};
\end{tikzpicture}
\caption{最高点分裂。分裂前竖直动量为零，轻块分裂后地面速度为零。}\label{sup3:fig:d3-h2-q}
\end{figure}
\begin{solution}
分裂冲量持续时间忽略，重力冲量可忽略，所以瞬时总动量守恒：
$3mU=2m v_{2x}$，$0=2m v_{2y}$。故重块初速度为 $(3U/2,0)$，轻块初速度为零。两块具有相同初始高度与竖直速度，因而同时经过
\[
t_f=\sqrt{\frac{2H}g}
\]
落地。图\ref{sup3:fig:d3-h2-s}中的落点 A、B 分别为
\[
x_A=0,\qquad x_B=\frac32U\sqrt{\frac{2H}g}.
\]
\begin{figure}[H]\centering
\begin{tikzpicture}
\draw[curve](-.8,0)--(7,0);\fill(0,3)circle(1.8pt)node[left]{分裂点};
\draw[curve,->](0,3)--(0,0);\draw[curve,domain=0:6,samples=70]plot(\x,{3-(\x)^2/12});
\draw[guide,domain=0:4,samples=60]plot(\x,{3-.1875*(\x)^2});
\fill(0,0)circle(1.6pt)node[below]{$A:m$};\fill(6,0)circle(1.6pt)node[below]{$B:2m$};\fill(4,0)circle(1.6pt)node[below]{$C_f$};
\node at(4.9,2.2){重块};\node[align=left]at(2,1.4){虚线：质心\\$x_C=Ut$};
\end{tikzpicture}
\caption{轻块竖直下落，重块射程增大；质心仍沿原抛体轨迹运动，$AC_f=2AB/3$。}\label{sup3:fig:d3-h2-s}
\end{figure}
分裂前动能为 $3mU^2/2$，分裂后为
$\tfrac12(2m)(3U/2)^2=9mU^2/4$，所以增加
\[
\Delta K=\frac34mU^2.
\]
增加的动能由分裂时释放的内能提供，与瞬时总动量守恒相容。

落地前两块的竖直初速度均为零，故
\[
v_y=-\int_0^t g\dd s=-gt,
\qquad y-H=\int_0^t(-gs)\dd s=[-gs^2/2]_0^t=-\frac12gt^2.
\]
水平方向速度不变，位移由 $\int_0^t v_x\dd s=v_xt$ 求得。因此两块坐标为 $(0,H-gt^2/2)$ 与 $(3Ut/2,H-gt^2/2)$，质量加权得到
\[
x_C=Ut,\qquad y_C=H-\frac12gt^2.
\]
这正是未分裂时原抛体的轨迹。第一次接触地面后会有地面冲量，届时不能再将相同抛体方程继续用于整个系统。
从分裂后到首次落地前，图\ref{sup3:fig:d3-h2-s}中质心到两碎块的水平距离之比为 $2:1$。以未分裂物体的质心轨迹为参照，两块的相对速度分别为 $-U$、$U/2$，相对动能为
\[
\frac12mU^2+\frac12(2m)(U/2)^2=\frac34mU^2.
\]
分裂前相对动能为零，所以新增动能全部进入相对运动，质心平动部分没有改变。两块竖直相对速度均为零，因而落地前质心的竖直运动也不受分裂影响。

\end{solution}


\begin{exercise}[人在船上行走后跳离]\label{sup3:problem-exercise-3}\label{sup3:ex:d3-h3}
质量 $M$ 的船在静水中，质量 $m$ 的人起初与船共同静止。不计水的水平阻力。人相对船向右走 $L$ 后停住，再从高于水面 $H$ 的船边水平向右跳出，跳出瞬间相对船的水平速度为 $u$，竖直速度为零。求行走阶段船的位移、起跳后人和船的水平速度，以及人落水前一瞬船的总位移。跳跃与飞行中均不计空气阻力，起跳时的高度已包含人体质心高度。
\end{exercise}
\begin{figure}[H]\centering
\begin{tikzpicture}
\draw[curve](0,.7)--(.5,0)--(6,0)--(6.5,.7)--cycle;\node at(3,.35){$M$};
\draw[guide](.2,-.25)--(7.8,-.25);\node[right]at(7.8,-.25){水面};
\fill(1.2,1.25)circle(2pt);\fill(5.4,1.25)circle(2pt);\draw[axis](1.2,1.65)--(5.4,1.65)node[midway,above]{相对船走 $L$};
\draw[axis](5.4,1.25)--(7,1.25)node[above]{起跳 $u$};\draw[<->](7.4,-.25)--(7.4,1.25)node[midway,right]{$H$};
\end{tikzpicture}
\caption{行走位移和起跳速度均相对船给出，必须转换为地面量。}\label{sup3:fig:d3-h3-q}
\end{figure}
\begin{solution}
取向右为正，船位移记为 $X$。人相对船走 $L$，相对地面位移为 $X+L$。水平方向无合外力，初始质心静止，因此
$MX+m(X+L)=0$，解得行走阶段
\[
X_1=-\frac m{M+m}L.
\]
人停住后与船速度相同，而总水平动量仍为零，所以二者在起跳前都静止。

起跳冲量是人船系统的内冲量。图\ref{sup3:fig:d3-h3-s}中人的地面水平速度为 $v$，船速为 $V$，动量与相对速度条件分别是
\[
MV+mv=0,\qquad v-V=u.
\]
由第二式 $v=V+u$ 代入第一式，得到
\[
V=-\frac{m}{M+m}u,\qquad v=\frac{M}{M+m}u.
\]
\begin{figure}[H]\centering
\begin{tikzpicture}
\draw[curve](0,.5)--(.4,0)--(4,0)--(4.5,.5);\draw[axis](2,.9)--(.7,.9)node[midway,above]{$|V|$};
\fill(4.5,1.7)circle(2pt)node[above]{$m$};\draw[axis](4.5,1.7)--(6,1.7)node[above]{$v$};
\draw[curve,domain=0:2,samples=50]plot({4.5+\x},{1.7-.4625*(\x)^2});
\fill(6.5,-.15)circle(1pt)node[below right]{落水};
\draw[guide](-.4,-.15)--(7,-.15);\node at(3.2,-.9){$v-V=u$，相对水平间隔以速率 $u$ 增大};
\end{tikzpicture}
\caption{起跳后船匀速向左，人做抛体运动。水平动量关系同时决定两个地面速度。}\label{sup3:fig:d3-h3-s}
\end{figure}
沿用习题\ref{sup3:ex:d3-h2}已积分得到的零竖直初速下落式，$H-g\tau^2/2=0$，得落水时间 $\tau=\sqrt{2H/g}$。人在空中时船不受水平力，继续以 $V$ 匀速运动，位移为 $\int_0^\tau V\dd s=V\tau$，因此
\[
X_{\rm total}=X_1+V\tau
=-\frac m{M+m}\left(L+u\sqrt{\frac{2H}g}\right).
\]
上式只到人首次落水前为止。落水后人与水交换水平动量，若仍把“人和船”作为系统，水的水平作用就不能忽略。行走和起跳会把人体内能转化为机械能。本题据水平质心与动量关系求解，机械能不守恒。
还可把图\ref{sup3:fig:d3-h3-s}中飞行阶段的两条水平位移代回质心关系。自初态至落水前，人的总水平位移为 $X_1+L+v\tau$，船为 $X_1+V\tau$，于是
\[
M(X_1+V\tau)+m(X_1+L+v\tau)
=(M+m)X_1+mL+(MV+mv)\tau=0.
\]
行走阶段给出 $(M+m)X_1+mL=0$，起跳后则有 $MV+mv=0$。两式分别消去上式中的常数项和含 $\tau$ 的项，说明从行走到飞行，总质心的水平位置始终不变。

\end{solution}


\begin{exercise}[有限尘区与离开后的运动]\label{sup3:problem-exercise-8}\label{sup3:ex:d3-h8}
沿用例\ref{ex:hw3-3}的拾尘板，令 $\kappa=\rho S$，但尘埃只占据 $0\le x\le L$。板从 $x=0$ 静止起受恒力 $F$，离开尘区后恒力仍继续作用。求出区时间、出区速度、出区后的位移及整个拾尘阶段的机械能损失。
\end{exercise}
\begin{figure}[H]\centering
\begin{tikzpicture}
\draw[axis](-.3,0)--(8.6,0)node[right]{$x$};\draw[guide](.6,.2)rectangle(5.7,2.1);
\foreach \xx in {1,1.8,2.6,3.4,4.2,5}\foreach \yy in {.5,1.1,1.7}\fill(\xx,\yy)circle(.7pt);
\draw[curve,line width=2pt](.6,.1)--(.6,2.3);\node[below]at(.6,0){$0$};\node[below]at(5.7,0){$L$};\node at(7,1.5){无尘区};
\draw[axis](.7,2.6)--(2,2.6)node[above]{$F$};
\end{tikzpicture}
\caption{尘区边界 $x=L$ 是质量增长规律改变的位置，恒力本身不发生变化。}\label{sup3:fig:d3-h8-q}
\end{figure}
\begin{solution}
尘区内有 $M=M_0+\kappa x$，由例\ref{ex:hw3-3}的积分式
$M_0x+\kappa x^2/2=Ft^2/2$。令 $x=L$，出区时间为
\[
t_L=\sqrt{\frac{2M_0L+\kappa L^2}F}.
\]
最终质量为 $M_L=M_0+\kappa L$，由 $M_Lv_L=Ft_L$，得
\[
v_L=\frac{\sqrt{F(2M_0L+\kappa L^2)}}{M_L}.
\]
出区以后没有流入项，板以固定质量 $M_L$ 运动，满足 $\dot v=F/M_L$。以出区时刻的状态为初值，
\[
v-v_L=\frac F{M_L}\int_{t_L}^t\dd s=\frac F{M_L}(t-t_L),
\]
\[
x-L=\int_{t_L}^t\left[v_L+\frac F{M_L}(s-t_L)\right]\dd s
=\left[v_L(s-t_L)+\frac F{2M_L}(s-t_L)^2\right]_{t_L}^t.
\]
故对 $t\ge t_L$ 有
\[
v=v_L+\frac F{M_L}(t-t_L),\qquad
x=L+v_L(t-t_L)+\frac F{2M_L}(t-t_L)^2.
\]
图\ref{sup3:fig:d3-h8-s}中的速度和位移在边界连续；加速度则从
$[F-\kappa v_L^2]/M_L$ 跳到 $F/M_L$。尘埃密度在边界突然降为零，使加速度发生有限跳变，速度仍保持连续。
\begin{figure}[H]\centering
\begin{tikzpicture}
\draw[axis](0,0)--(7,0)node[right]{$t$};\draw[axis](0,0)--(0,3.4)node[above]{$v$};
\draw[curve,domain=0:3,samples=60]plot(\x,{1.6*\x/sqrt(1+(\x)^2)});
\draw[curve](3,1.517893)--(6,3.035786);\draw[guide](3,0)--(3,3);\node[below]at(3,0){$t_L$};\node[right]at(3,1.3){$v_L$};
\node at(1.3,2.6){质量增长};\node at(5.3,.65){固定质量 $M_L$};
\end{tikzpicture}
\caption{速度在出区时连续，切线斜率增大。左段的水平渐近趋势在出区后不再适用。}\label{sup3:fig:d3-h8-s}
\end{figure}
拾尘阶段输入功为 $FL$，末动能为
\[
K_L=\frac12M_Lv_L^2
=\frac{F(2M_0L+\kappa L^2)}{2M_L}.
\]
两者之差为
\[
E_{\rm diss}=FL-K_L=\frac{F\kappa L^2}{2M_L}.
\]
它只在尘区内增加，出区后不再有黏附耗散。如果把尘区内的渐近速度 $\sqrt{F/\kappa}$ 当成全程极限速度，就会漏掉质量增长停止后的加速阶段。
图\ref{sup3:fig:d3-h8-s}的边界加速度可进一步化简。由出区速度可得
\[
F-\kappa v_L^2
=F\frac{M_L^2-\kappa(2M_0L+\kappa L^2)}{M_L^2}
=\frac{FM_0^2}{M_L^2}.
\]
因此边界前后分别为 $a_-=FM_0^2/M_L^3$、$a_+=F/M_L$，且 $a_+>a_-$。边界没有额外的集中冲量，故只改变加速度，不改变当时的速度。

\end{solution}


\begin{exercise}[漏砂小车有无喷气推力]\label{sup3:problem-exercise-9}\label{sup3:ex:d3-h9}
总初质量 $M_0$ 的小车在光滑水平面上从静止起受向右恒力 $F$。车内砂从底孔以恒定质量速率 $\alpha$ 漏出，离车瞬间砂的水平速度与车相同，漏出后也不受水平力。令当前车与车内砂的质量为 $M=M_0-\alpha t>0$。求车的水平速度和位移，并通过“车与已经漏出的全部砂”的总水平动量解释变质量项。
\end{exercise}
\begin{figure}[H]\centering
\begin{tikzpicture}
\draw[curve](0,1)rectangle(4,2.2);\node at(2,1.7){$M_0-\alpha t$};\draw(.6,.7)circle(.25);\draw(3.4,.7)circle(.25);
\draw[axis](4,1.6)--(5.5,1.6)node[above]{$F$};
\foreach \yy in {.7,.35,0}\fill(2,\yy)circle(1.2pt);\draw[axis](2.2,-.05)--(3.4,-.05)node[right]{$w_x=v$};
\draw[guide](-.4,-.55)--(5.8,-.55);\node at(.1,.15){$\alpha$};
\end{tikzpicture}
\caption{砂没有相对车的水平喷出速度，竖直漏出不自动产生水平推力。}\label{sup3:fig:d3-h9-q}
\end{figure}
\begin{solution}
图\ref{sup3:fig:d3-h9-q}中流出砂的地面水平速度为 $w_x=v$，所以流出砂相对车的水平速度为零，水平运动方程是
\[
(M_0-\alpha t)\dot v=F.
\]
把即将漏出的砂也计入同一批物质，在 $\dd t$ 内写出
\[
Mv+F\dd t=(M-\alpha\dd t)(v+\dd v)+\alpha\dd t\,v.
\]
展开后，流出质量乘原速度的两项抵消，只剩 $M\dd v=F\dd t$。这与火箭中流出物具有非零水平相对速度的情形不同。

砂带走动量，但不会额外提供水平喷气推力。由 $v(0)=0$ 积分，
\[
v(t)=F\int_0^t\frac{\dd s}{M_0-\alpha s}
=-\frac F\alpha[\ln(M_0-\alpha s)]_0^t,
\]
即
\[
v(t)=\frac F\alpha\ln\frac{M_0}{M_0-\alpha t}.
\]
令 $M=M_0-\alpha t$，再积分一次并取 $x(0)=0$。所需的对数积分与例\ref{ex:hw3-6}求高度时相同，沿用其中的质量代换和原函数，
\[
x=\frac F{\alpha^2}
\left[\widetilde M\left(\ln\frac{M_0}{\widetilde M}+1\right)\right]_{M}^{M_0},
\]
得到
\[
x(t)=\frac F\alpha\left[t-\frac M\alpha\ln\frac{M_0}M\right].
\]
初期对数展开为 $v\simeq Ft/M_0$；随着质量减小，加速度 $F/M$ 增大。
\begin{figure}[H]\centering
\begin{tikzpicture}
\draw[guide](-.3,-.2)rectangle(8.5,2.4);\draw[curve](5.6,.5)rectangle(7.8,1.5);\node at(6.7,1){当前车 $M,v$};
\foreach \xx/\vv in {.5/.4,1.8/.6,3.3/.9}{\fill(\xx,.7)circle(2pt);\draw[axis](\xx,.7)--({\xx+\vv},.7);}
\node at(2.6,1.6){已漏砂各保持离车时的水平速度};
\node at(4,-.8){$P_{\rm total}=Mv+\displaystyle\int_0^t\alpha v(s)\,\dd s=Ft$};
\end{tikzpicture}
\caption{扩大成固定物质系统后，漏出的砂仍须计入动量。虚线只表示系统成员，不是容器壁。}\label{sup3:fig:d3-h9-s}
\end{figure}
按图\ref{sup3:fig:d3-h9-s}，若已漏砂后续不受水平力，时刻 $s$ 漏出的砂在时刻 $t$ 仍有水平速度 $v(s)$。其总动量为
\[
P_{\rm sand}=\int_0^t\alpha v(s)\,\dd s=\alpha x(t).
\]
因此固定系统总动量应满足 $Mv+\alpha x=Ft$。代入上面两个显式解，含对数的两项恰好抵消，确实得到 $Ft$。

只对当前小车求导，则
\[
\frac{\dd(Mv)}{\dd t}=M\dot v-\alpha v=F-\alpha v.
\]
式中的 $\alpha v$ 是砂带走的水平动量流率，故当前小车的动量变化率为 $F-\alpha v$。若 $F=0$ 而初始速度为某个 $v_0$，车速恒为 $v_0$，仍会持续失去当前物体的动量；其动量由离车砂携走。
\end{solution}


\begin{exercise}[两级火箭的分级收益]\label{sup3:problem-exercise-10}\label{sup3:ex:d3-h10}
在无重力、无阻力的一维空间中，火箭初始静止。有效载荷质量为 $m$；外级空壳和内级空壳质量各为 $m$，每级燃料各为 $2m$，所有燃料相对当前火箭向后喷出的速率均为 $u$。先烧尽外级燃料，抛掉外级空壳，再烧尽内级燃料，最后抛掉内级空壳。抛壳瞬间无相对分离速度。求载荷最终速度，并与两级空壳均保留到燃料全部耗尽的方案比较。
\end{exercise}
\begin{figure}[H]\centering
\begin{tikzpicture}
\draw[curve](0,0)rectangle(9,1.65);\draw[curve](3,0)--(3,1.65);\draw[curve](6,0)--(6,1.65);
\node[align=center]at(1.5,.825){外级\\空壳 $m$\\燃料 $2m$};\node[align=center]at(4.5,.825){内级\\空壳 $m$\\燃料 $2m$};\node[align=center]at(7.5,.825){载荷\\$m$};
\draw[axis](9,.825)--(10.5,.825)node[above]{前进};\node at(4.5,2.25){初始总质量 $7m$};
\end{tikzpicture}
\caption{分级收益来自尽早抛掉后续不再需要加速的空壳。}\label{sup3:fig:d3-h10-q}
\end{figure}
\begin{solution}
无外力时，喷气方程 $M\dd v=-u\dd M$ 的积分已在例\ref{ex:hw3-6}求出。每个燃烧段按各自的初末质量取界，$\Delta v=-u[\ln M]_{M_i}^{M_f}=u\ln(M_i/M_f)$。图\ref{sup3:fig:d3-h10-s}中的质量阶梯必须区分“燃烧”与“抛壳”：燃烧有相对喷气速度，抛壳没有相对分离速度，后者不使剩余火箭速度瞬间跳变。
\begin{figure}[H]\centering
\begin{tikzpicture}
\draw[axis](0,0)--(8,0)node[right]{过程};\draw[axis](0,0)--(0,4.2)node[above]{$M/m$};
\draw[curve](.3,3.5)--(2.2,2.5)--(2.2,2)--(4.8,1)--(4.8,.5)--(7,.5);
\foreach \yy/\la in {3.5/7,2.5/5,2/4,1/2,.5/1}{\draw[guide](0,\yy)--(.25,\yy);\node[left]at(0,\yy){$\la$};}
\node at(1.1,3.4){燃烧};\node[right]at(2.2,2.35){抛外壳};\node at(3.7,1.9){燃烧};\node[right]at(4.8,.75){抛内壳};
\end{tikzpicture}
\caption{质量在燃烧中连续下降，抛壳时离散下降。只有燃烧段提供本题的速度增量。}\label{sup3:fig:d3-h10-s}
\end{figure}
第一段质量由 $7m$ 降至 $5m$，所以 $\Delta v_1=u\ln(7/5)$。抛外壳后剩 $4m$，速度不变。第二段由 $4m$ 降至 $2m$，故 $\Delta v_2=u\ln2$。最后抛内壳也不改变速度，最终
\[
v_f=u\ln\frac{14}5.
\]
若不提前抛外壳，燃料全部消耗使质量由 $7m$ 降至 $3m$，随后才抛空壳，载荷速度为
\[
v_{\rm keep}=u\ln\frac73.
\]
分级增加的速度为
\[
v_f-v_{\rm keep}=u\ln\frac65>0.
\]
各燃烧段的速度增量相加，对应的质量比相乘。抛壳过程虽然使当前火箭质量减少，却因相对分离速度为零而没有推进作用；它的收益体现在后续燃烧需要加速的质量更小。
图\ref{sup3:fig:d3-h10-s}中各阶段的质量为：出发时，载荷与两个空壳共 $3m$，两份燃料共 $4m$；第一段燃烧后剩 $5m$，抛掉外壳后剩 $4m$；第二段燃烧后，内壳与载荷共 $2m$。抛壳的相对速度为零，设抛壳前速度为 $V$，瞬时动量关系为
\[
M_iV=M_fV+(M_i-M_f)V,
\]
所以质量虽跳变，速度不跳变。两级收益来自第二段的起始、终止质量比由保留外壳时的 $5/3$ 改成 $4/2$。

\end{solution}


\begin{exercise}[有有限启动长度的水平折绳]\label{deep3:problem-exercise-6}\label{deep3:ex:d3-h6}
沿用例\ref{ex:hw3-1}的水平折绳模型，本题将绳的全长记为$L$、线密度记为$\lambda$。从 $x=x_0>0$、端点速度 $v_0\ge0$ 的状态开始计时，以恒力 $F_0>0$ 向右牵引。设 $x_0<2L$。求 $x(t)$、$v(t)$、整绳启动所需时间，说明 $x_0\to0^+$ 与直接指定零端点初速度之间的关系。
\end{exercise}
\begin{figure}[H]\centering
\begin{tikzpicture}
\draw[curve](7,0)--(1.75,0)arc(-90:-270:.15)--(3.2,.3);
\fill(3.2,.3)circle(1.5pt)node[above]{$A_0:x_0$};\fill(1.6,.15)circle(1pt)node[left]{$B_0$};
\draw[axis](3.2,.3)--(4.6,.3)node[above]{$F_0,\ v_0$};\node[below]at(7,0){$L$};
\draw[guide](1.6,.15)--(1.6,-.65);
\draw[<->](0,-.65)--(1.6,-.65)node[midway,below]{$x_0/2$};\node[above left]at(0,-.65){$0$};
\end{tikzpicture}
\caption{初始已有长 $x_0/2$ 的绳在运动，避免了零运动质量的奇异初态。}\label{deep3:fig:d3-h6-q}
\end{figure}
\begin{solution}
总动量仍为 $P=\lambda xv/2$，故
\[
\frac\lambda2\frac{\dd(xv)}{\dd t}=F_0.
\]
定义 $A=2F_0/\lambda$，由初始动量积分，$xv-x_0v_0=\int_0^t A\dd s=At$。再用 $v=\dot x$ 分离变量，
\[
\int_{x_0}^{x(t)}\xi\dd\xi=\int_0^t(x_0v_0+As)\dd s,
\qquad [\xi^2/2]_{x_0}^{x(t)}=[x_0v_0s+As^2/2]_0^t.
\]
整理得到
\[
x^2=x_0^2+2x_0v_0t+At^2.
\]
因为 $x>0$，取正根：
\[
x(t)=\sqrt{x_0^2+2x_0v_0t+At^2},\qquad
v(t)=\frac{x_0v_0+At}{\sqrt{x_0^2+2x_0v_0t+At^2}}.
\]
在 $x_0>0$ 时，该解确实满足 $v(0)=v_0$。相应的速度位置关系为
\[
v^2=A+\frac{x_0^2(v_0^2-A)}{x^2}.
\]
图\ref{deep3:fig:d3-h6-s}显示：若初速小于 $\sqrt A$，端点加速；若大于 $\sqrt A$，端点减速，但因外力为正，全绳总动量仍增加。
\begin{figure}[H]\centering
\begin{tikzpicture}
\draw[axis](0,0)--(4.2,0)node[right]{$x$};\draw[axis](0,0)--(0,3)node[above]{$v^2$};
\draw[guide](.5,1.5)--(4,1.5)node[right]{$A$};\draw[guide](1,0)--(1,2.5);\node[below]at(1,0){$x_0$};
\draw[curve,domain=1:4,samples=60]plot(\x,{1.5-1/(\x)^2});
\draw[curve,domain=1:4,samples=60]plot(\x,{1.5+.8/(\x)^2});
\node at(2.7,.45){$v_0^2<A$};\node at(2.7,2.5){$v_0^2>A$};
\begin{scope}[shift={(6.4,.3)}]
\draw[curve](0,0)--(2.6,0);\fill(0,0)circle(1.3pt)node[below]{$B:L$};\fill(2.6,0)circle(1.3pt)node[below]{$A:2L$};\node[above]at(1.3,1){完成时 $x=2L$};\draw[axis](1.1,.4)--(2,.4)node[above]{$v$};\node at(1.3,-1){静止段长度 $L-x/2=0$};
\end{scope}
\end{tikzpicture}
\caption{速度平方随牵引位移趋向 $A$；达到 $x=2L$ 后，模型转为固定质量整绳平移。}\label{deep3:fig:d3-h6-s}
\end{figure}
令 $x=2L$，由二次方程得到正的完成时间
\[
t_f=\frac{-x_0v_0+\sqrt{x_0^2v_0^2+A(4L^2-x_0^2)}}A.
\]
完成以前的耗散功率仍是 $\lambda v^3/4$，与是否匀速无关。完成后若继续施加 $F_0$，整绳质量为 $\lambda L$，加速度改为 $F_0/(\lambda L)$，不再有新启动绳元。

固定 $t>0$，令 $x_0\to0^+$，上面的解趋于 $x=\sqrt A\,t$、$v=\sqrt A$，与例\ref{ex:hw3-1}一致。但在 $t=0$，有限 $x_0$ 的解始终有 $v(0)=v_0$。两个极限一般不能交换；尤其当 $v_0=0$ 时，极限过程含有越来越短的启动阶段，不能将极限解解释为有限质量端点在零时刻无冲量地获得非零速度。
初始加速度由已经积分的动量式 $xv=x_0v_0+At$ 再求导：$v^2+x\dot v=A$。解出加速度，并取初值，
\[
\dot v=\frac{A-v^2}{x},\qquad
\dot v(0)=\frac{A-v_0^2}{x_0}.
\]
当 $v_0=0$ 而 $x_0\to0^+$ 时，初始加速度发散。图\ref{deep3:fig:d3-h6-s}的不同初速分支由此在一个很短的时间内趋近同一速度。典型启动时间量级为 $x_0/\sqrt A$；将这个时间段压缩为零，才得到零启动长度的奇异解。

\end{solution}


\begin{exercise}[偏心冲量与作用点的有效质量]\label{deep3:problem-exercise-11}\label{deep3:ex:d3-h11}
均匀细杆长 $\ell$、质量 $m$，在光滑水平面内静止。于距质心 C 的有向距离 $a$ 处的点 P 施加垂直于杆的短时冲量 $J>0$，其中 $|a|\le\ell/2$。取杆向右为有向距离正向，冲量向图中上方。求冲量后质心速度、角速度、点 P 的速度与总动能，并定义满足 $v_P=J/m_{\rm eff}$ 的有效质量。
\end{exercise}
\begin{figure}[H]\centering
\begin{tikzpicture}
\draw[curve,line width=2pt](0,0)--(7,0);\fill(3.5,0)circle(2pt)node[below]{$C$};\fill(5.5,0)circle(2pt)node[below]{$P$};
\draw[axis](5.5,0)--(5.5,1.8)node[above]{$J$};\draw[<->](3.5,-.7)--(5.5,-.7)node[midway,below]{$a$};\draw[<->](0,-1.5)--(7,-1.5)node[midway,below]{$\ell$};
\end{tikzpicture}
\caption{图示 $a>0$；一般公式同时适用于质心左侧的作用点。}\label{deep3:fig:d3-h11-q}
\end{figure}
\begin{solution}
外冲量决定质心动量增量，因此 $V_C=J/m$，方向与冲量相同。关于质心的角冲量为 $aJ$，而 $I_C=m\ell^2/12$，故取逆时针为正有
\[
\omega=\frac{aJ}{I_C}=\frac{12aJ}{m\ell^2}.
\]
图\ref{deep3:fig:d3-h11-s}中点 P 的速度由质心平移与绕质心转动叠加，初始杆水平，所以两个速度分量均沿冲量轴：
\[
v_P=V_C+\omega a
=J\left(\frac1m+\frac{a^2}{I_C}\right).
\]
即使 $a<0$，转动方向也同时改变，使 $\omega a\ge0$。
\begin{figure}[H]\centering
\begin{tikzpicture}
\draw[curve,line width=1.5pt](0,0)--(7,0);\fill(3.5,0)circle(2pt)node[below]{$C$};\fill(5.5,0)circle(2pt)node[below]{$P$};
\draw[axis](3.5,.1)--(3.5,1.5)node[left]{$J/m$};\draw[axis](5.5,.1)--(5.5,1.5)node[left]{$V_C$};\draw[axis](5.5,1.5)--(5.5,2.5)node[left]{$\omega a$};
\draw[axis](6.6,.1)--(6.6,2.5)node[right]{$v_P$};\draw[axis](2.2,1.6)arc(0:110:.55);\node at(1.4,2.6){$\omega$};
\end{tikzpicture}
\caption{点 P 的速度分解。作用点远离质心时，转动使该点速度增量更大。}\label{deep3:fig:d3-h11-s}
\end{figure}
按 $v_P=J/m_{\rm eff}$ 定义，有
\[
m_{\rm eff}=\left(\frac1m+\frac{a^2}{I_C}\right)^{-1}
=\frac{m}{1+12a^2/\ell^2}.
\]
总动能应同时包含平动与转动：
\[
K=\frac{J^2}{2m}+\frac{a^2J^2}{2I_C}
=\frac{J^2}{2m_{\rm eff}}=\frac12Jv_P.
\]
当 $a=0$ 时，$m_{\rm eff}=m$，只有平动；当 $|a|=\ell/2$ 时，$m_{\rm eff}=m/4$，总动能为 $2J^2/m$，其中平动动能仍仅为 $J^2/(2m)$。本题从静止开始，速度随累积冲量线性增加，因而得到 $K=Jv_P/2$；有初速度或响应非线性时须重新计算。
式 $K=Jv_P/2$ 也有直接的功解释。令 $j$ 表示从零累积到 $J$ 的冲量，图\ref{deep3:fig:d3-h11-s}中作用点沿冲量方向的速度为 $v_P(j)=j/m_{\rm eff}$。由于 $\dd j=F\dd t$，
\[
W=\int Fv_P\dd t=\int_0^J\frac{j}{m_{\rm eff}}\dd j
=\frac1{m_{\rm eff}}[j^2/2]_0^J=\frac{J^2}{2m_{\rm eff}}.
\]
作用点的速度从零随累积冲量线性增长，积分后的功为最终速度与总冲量乘积的一半。

\end{solution}


\begin{exercise}[有限固有加速度的对称航行，选做]\label{deep3:problem-exercise-12}\label{deep3:ex:d3-h12}
静质量保持不变的飞船从地面系静止出发，前半程以恒定固有加速度 $a_0$ 向前加速，后半程以相同大小的固有加速度制动，恰在距离 $D$ 的静止目的地停下，无匀速段。求地面航时 $T$、船内经历的总固有时 $\tau$ 和中点最大速度，验证低速极限。
\end{exercise}
\begin{figure}[H]\centering
\begin{tikzpicture}
\draw[axis](0,0)--(8,0)node[right]{$x$};\fill(0,0)circle(1.8pt)node[below]{$0$};\fill(3.5,0)circle(1.8pt)node[below]{$D/2$};\fill(7,0)circle(1.8pt)node[below]{$D$};
\draw[axis](.5,1)--(2,1)node[above]{$+a_0$};\draw[axis](6.5,1)--(5,1)node[above]{$-a_0$};
\node at(3.5,2){中点换向推力，速度连续};
\end{tikzpicture}
\caption{两段固有加速度大小相同，出发和到达都静止。中点切换的是加速度方向。}\label{deep3:fig:d3-h12-q}
\end{figure}
\begin{solution}
将例\ref{ex:hw3-7}用于前半程，令地面时间为 $t_h$，末态洛伦兹因子为 $\gamma_h$。位移条件为
\[
\frac D2=\frac{c^2}{a_0}(\gamma_h-1),\qquad
\gamma_h=1+\frac{a_0D}{2c^2}.
\]
由前半程动量关系 $\gamma_hv_h=a_0t_h$，或由 $\gamma_h^2=1+(a_0t_h/c)^2$，可得
\[
t_h=\sqrt{\frac D{a_0}+\frac{D^2}{4c^2}}.
\]
后半程是前半程在地面系中的时间反演平移，所用时间与位移相同，故
\[
T=2t_h=2\sqrt{\frac D{a_0}+\frac{D^2}{4c^2}}.
\]
\begin{figure}[H]\centering
\begin{tikzpicture}
\draw[axis](0,0)--(7.6,0)node[right]{$t$};\draw[axis](0,0)--(0,3)node[above]{$v$};
\draw[curve,domain=0:3.2,samples=60]plot(\x,{2.4*\x/sqrt(1+(\x)^2)});
\draw[curve,domain=3.2:6.4,samples=60]plot(\x,{2.4*(6.4-\x)/sqrt(1+(6.4-\x)^2)});
\draw[guide](3.2,0)--(3.2,2.6);\node[below]at(3.2,0){$T/2$};\node[below]at(6.4,0){$T$};
\node[above]at(3.2,2.6){$v_h=c\sqrt{1-\gamma_h^{-2}}$};\node at(3.2,.8){曲线下面积为 $D$};
\end{tikzpicture}
\caption{有限加速度下的对称速度历程。中点速度最大，出发与到达速度均为零。}\label{deep3:fig:d3-h12-s}
\end{figure}
图\ref{deep3:fig:d3-h12-s}中峰值速度为
\[
v_h=c\sqrt{1-\left(1+\frac{a_0D}{2c^2}\right)^{-2}}<c.
\]
船内前半程固有时由 $\gamma_h=\cosh\eta_h$ 与
$\eta_h=a_0\tau_h/c$ 给出。所以
\[
\tau=2\tau_h=\frac{2c}{a_0}\operatorname{arcosh}\left(1+\frac{a_0D}{2c^2}\right)
=\frac{4c}{a_0}\operatorname{arsinh}\sqrt{\frac{a_0D}{4c^2}}.
\]
后一个等式用到了 $\cosh(2z)=1+2\sinh^2z$。

当 $a_0D/c^2\ll1$ 时，$T\simeq2\sqrt{D/a_0}$、$v_h\simeq\sqrt{a_0D}$，并由 $\operatorname{arsinh}z\simeq z$ 得 $\tau\simeq2\sqrt{D/a_0}$，正好恢复先匀加速、后匀减速的牛顿结果。任意有限 $a_0$ 下都有 $T>D/c$，与光速限制一致。
对后半程，仍记 $\ell=c^2/a_0$，以距抵达时刻的剩余地面时间 $s=T-t$ 表示后半程。其速度和位置为
\[
v(t)=\frac{a_0s}{\sqrt{1+(a_0s/c)^2}},\qquad
x(t)=D-\ell\left[\sqrt{1+(a_0s/c)^2}-1\right].
\]
在中点 $s=t_h$，两段的位置都为 $D/2$、速度都为 $v_h$；在终点 $s=0$，有 $x=D,v=0$。图\ref{deep3:fig:d3-h12-s}中两段世界线因此在中点有共同切线，推进方向在此改变。

\end{solution}


\chapter{简正模问题}


平衡附近的运动常可用线性方程近似。在单一振动模态中，各坐标以同一频率运动，振幅之比保持不变，这组比例称为振型。稳定平衡附近的任意线性小振动由这些模态叠加而成，初始条件决定各模态的份额。



\section{平衡构形、二次能量与固有频率}

小振动围绕满足约束的平衡构形展开。独立位移确定后，动能与势能的二次项在同一组坐标中给出质量矩阵和刚度矩阵，由此求固有频率。



\subsection{从平衡位置到二次能量}
设独立坐标为 $q_i$，平衡位置为 $q_{i*}$，小偏移为 $\eta_i=q_i-q_{i*}$。以下假定约束不显含时间、真实力可由势能表示，且平衡态静止。若平衡态本身持续转动，还须保留转动参考系中的陀螺项，方程形式会有所不同。

在平衡点附近，将能量保留到小量的二阶：
\begin{equation}
T_2=\frac12\dot{\vct\eta}^{\mathsf T}M\dot{\vct\eta},\qquad
U_2=U_*+\frac12\vct\eta^{\mathsf T}K\vct\eta.
\label{sup4:eq:d7-energy}
\end{equation}
$M$ 是平衡位置处的质量矩阵，$K_{ij}=\partial_i\partial_jU|_*$ 是刚度矩阵，二者均对称。使用独立且非奇异的坐标时，$M$ 正定；$K$ 则由平衡性质决定。

对单坐标系统，若 $T_2=A_*\dot\eta^2/2$、$U_2=U_*+B\eta^2/2$，有
\[
A_*\ddot\eta+B\eta=0,\qquad \omega^2=\frac B{A_*}.
\]
$A_*$ 由所选坐标决定：直线位移对应等效质量，角度对应等效转动惯量。由角度改用弧长时，动能与势能的系数须同时变换，所得物理频率保持不变。
\begin{figure}[H]\centering
\begin{tikzpicture}
\draw[axis](-2.5,0)--(2.7,0)node[right]{$\eta$};\draw[axis](0,0)--(0,2.7)node[above]{$U-U_*$};
\draw[curve,domain=-1.8:1.8,samples=90]plot(\x,{.5*(\x)^2+.08*(\x)^4});\draw[guide,domain=-1.8:1.8,samples=80]plot(\x,{.5*(\x)^2});\node[right]at(1.6,.7){二次近似};
\end{tikzpicture}
\caption{势阱在平衡附近的二次近似。振幅增大后，高阶项可改变周期。}\label{sup4:fig:d7-potential}
\end{figure}
图\ref{sup4:fig:d7-potential}中，二次系数为正时存在非零的小振动频率。若二次项消失，不能把零频率当作稳定性结论，须保留高阶势能；若存在整体平移自由度，则零频率也可能对应匀速漂移。


\subsection{质量矩阵与固有振型}
将二次能量写为分量和式，利用 $M,K$ 的对称性，
\[
\frac{\partial T_2}{\partial\dot\eta_i}
=\frac12\sum_j(M_{ij}+M_{ji})\dot\eta_j=\sum_jM_{ij}\dot\eta_j,
\qquad
\frac{\partial U_2}{\partial\eta_i}=\sum_jK_{ij}\eta_j.
\]
质量矩阵在本阶取平衡点处常量，将动量对时间求导，拉格朗日方程给出 $\sum_jM_{ij}\ddot\eta_j+\sum_jK_{ij}\eta_j=0$，写成矩阵即
\begin{equation}
M\ddot{\vct\eta}+K\vct\eta=0.
\label{sup4:eq:d7-matrix}
\end{equation}
例如 $T_2$ 中若有 $c\dot\eta_1\dot\eta_2$，则 $M_{12}=M_{21}=c$，不是 $2c$；若势能含 $\kappa(\eta_2-\eta_1)^2/2$，展开后 $K_{12}=K_{21}=-\kappa$。

取所有坐标具有同一时间因子的解 $\vct\eta=\vct u\cos(\omega t+\phi)$，得到广义特征方程
\begin{equation}
(K-\omega^2M)\vct u=0,\qquad
\det(K-\omega^2M)=0.
\label{sup4:eq:d7-eigen}
\end{equation}
非零矢量 $\vct u$ 给出振型：各分量之比决定同一模态中各物体的振幅比例和相位关系。整体倍乘可并入模态振幅，故振型本身没有固定长度。
\begin{figure}[H]\centering
\begin{tikzpicture}
\draw[curve](0,0)--(7,0);\draw[curve](0,0)--(0,1.1);\draw[curve](7,0)--(7,1.1);\draw[spring](0,.5)--(2,.5);\draw[curve](2,.2)rectangle(2.7,.8);\draw[spring](2.7,.5)--(4.3,.5);\draw[curve](4.3,.2)rectangle(5,.8);\draw[spring](5,.5)--(7,.5);
\node at(2.35,.5){$m$};\node at(4.65,.5){$m$};\node[above]at(1,.65){$k$};\node[above]at(3.5,.65){$\kappa$};\node[above]at(6,.65){$k$};\draw[axis](2.1,-.5)--(3,-.5)node[right]{$x_1$};\draw[axis](4.4,-.5)--(5.3,-.5)node[right]{$x_2$};
\end{tikzpicture}
\caption{双振子以各自平衡位置为位移原点，两个正方向均向右。}\label{sup4:fig:d7-pair}
\end{figure}
图\ref{sup4:fig:d7-pair}的方程为
\[
m\ddot x_1+(k+\kappa)x_1-\kappa x_2=0,\qquad
m\ddot x_2+(k+\kappa)x_2-\kappa x_1=0.
\]
令 $x_+=x_1+x_2$、$x_-=x_1-x_2$。相加与相减分别给出
\[
m\ddot x_++kx_+=0,\qquad
m\ddot x_-+(k+2\kappa)x_-=0.
\]
所以同相振型 $(1,1)$ 和反相振型 $(1,-1)$ 的频率分别为 $\sqrt{k/m}$、$\sqrt{(k+2\kappa)/m}$。同相运动不拉伸中间弹簧，因此该模态的频率与 $\kappa$ 无关。


\subsection{例题与模型变式}

\begin{example}[圆环上的四个等质量球]\label{ex:hw4-1}

四个质量为$m$的小球可在固定光滑圆环上滑动，相邻两球用沿圆环布置的相同轻弹簧连接，劲度系数均为$k$。平衡时四球等距。求沿环方向小振动的全部简正模式。

\end{example}


\begin{figure}[H]\centering
\begin{tikzpicture}[>=Stealth,line cap=round,line join=round]
\draw(0,0)circle(1.5);\foreach \a/\i in {90/1,0/2,270/3,180/4}{\fill(\a:1.5)circle(3pt)node[shift={(\a:.35)}]{$\i$};}\foreach \a in {0,90,180,270}{\draw[spring,domain=\a+10:\a+80,samples=25]plot({1.5*cos(\x)},{1.5*sin(\x)});}\node at(0,0){$m,k$};
\end{tikzpicture}
\caption{由四个相同弹簧连接的环上质点}\label{fig:hw4-1-q}
\end{figure}

\begin{solution}
取图\ref{fig:hw4-1-s}中各球沿同一绕环正向的弧长位移$q_i$，下标按模4循环。弹簧伸长量为$q_{i+1}-q_i$，因此
\[
 m\ddot q_i=k(q_{i+1}+q_{i-1}-2q_i).
\]

\begin{figure}[H]\centering
\begin{tikzpicture}[>=Stealth,line cap=round,line join=round]
\node at(1.7,1.5){各$q_i$以顺时针为正，箭头表示模态位移};\begin{scope}[shift={(0,0)}]\draw(0,0)circle(.8);\fill(0.000000,0.800000)circle(2pt)node[shift={(90:.22)}]{$ 1 $};\draw[axis](0.000000,0.800000)--(0.500000,0.800000);\fill(0.800000,0.000000)circle(2pt)node[shift={(0:.22)}]{$ 2 $};\draw[axis](0.800000,0.000000)--(0.800000,-0.500000);\fill(-0.000000,-0.800000)circle(2pt)node[shift={(270:.22)}]{$ 3 $};\draw[axis](-0.000000,-0.800000)--(-0.500000,-0.800000);\fill(-0.800000,0.000000)circle(2pt)node[shift={(180:.22)}]{$ 4 $};\draw[axis](-0.800000,0.000000)--(-0.800000,0.500000);\node at(0,-1.55){$\omega^2=0$};\end{scope}\begin{scope}[shift={(3.5,0)}]\draw(0,0)circle(.8);\fill(0.000000,0.800000)circle(2pt)node[shift={(90:.22)}]{$ 1 $};\draw[axis](0.000000,0.800000)--(0.500000,0.800000);\fill(0.800000,0.000000)circle(2pt)node[shift={(0:.22)}]{$ 2 $};\fill(-0.000000,-0.800000)circle(2pt)node[shift={(270:.22)}]{$ 3 $};\draw[axis](-0.000000,-0.800000)--(0.500000,-0.800000);\fill(-0.800000,0.000000)circle(2pt)node[shift={(180:.22)}]{$ 4 $};\node at(0,-1.55){$\omega^2=2k/m$};\end{scope}\begin{scope}[shift={(0,-3)}]\draw(0,0)circle(.8);\fill(0.000000,0.800000)circle(2pt)node[shift={(90:.22)}]{$ 1 $};\fill(0.800000,0.000000)circle(2pt)node[shift={(0:.22)}]{$ 2 $};\draw[axis](0.800000,0.000000)--(0.800000,-0.500000);\fill(-0.000000,-0.800000)circle(2pt)node[shift={(270:.22)}]{$ 3 $};\fill(-0.800000,0.000000)circle(2pt)node[shift={(180:.22)}]{$ 4 $};\draw[axis](-0.800000,0.000000)--(-0.800000,-0.500000);\node at(0,-1.55){$\omega^2=2k/m$};\end{scope}\begin{scope}[shift={(3.5,-3)}]\draw(0,0)circle(.8);\fill(0.000000,0.800000)circle(2pt)node[shift={(90:.22)}]{$ 1 $};\draw[axis](0.000000,0.800000)--(0.500000,0.800000);\fill(0.800000,0.000000)circle(2pt)node[shift={(0:.22)}]{$ 2 $};\draw[axis](0.800000,0.000000)--(0.800000,0.500000);\fill(-0.000000,-0.800000)circle(2pt)node[shift={(270:.22)}]{$ 3 $};\draw[axis](-0.000000,-0.800000)--(-0.500000,-0.800000);\fill(-0.800000,0.000000)circle(2pt)node[shift={(180:.22)}]{$ 4 $};\draw[axis](-0.800000,0.000000)--(-0.800000,-0.500000);\node at(0,-1.55){$\omega^2=4k/m$};\end{scope}
\end{tikzpicture}
\caption{四个独立简正模及各球的相对位移方向}\label{fig:hw4-1-s}
\end{figure}

代入四个独立向量，可得：
\[
 \begin{array}{c|c}
 \vct a&\omega^2\\\hline
 (1,1,1,1)&0\\
 (1,0,-1,0)&2k/m\\
 (0,1,0,-1)&2k/m\\
 (1,-1,1,-1)&4k/m
 \end{array}
\]
第一种使各弹簧伸长均不变，是整体绕环的零频运动；其时间系数可为$A+Bt$。中间两个模式退化，其任意独立线性组合仍属于同一频率。最后一个模式中相邻球反向，回复作用最大。

四个向量构成完整基底，给定初始条件后可逐项投影。两个非零频率中有一个二重简并，因此仍须保留四个独立模式。


四个位移组成列向量后，质量矩阵为$m\vct 1$，刚度矩阵为
\[
 \vct K=k\begin{pmatrix}
 2&-1&0&-1\\-1&2&-1&0\\0&-1&2&-1\\-1&0&-1&2
 \end{pmatrix}.
\]
令$a_j=\exp[i\kappa(j-1)]$，周期条件要求$\kappa=s\pi/2$，$s=0,1,2,3$。代入方程即得$m\omega^2=2k(1-\cos\kappa)$，从而得到表中的零频、二重简并频率和最高频率。复指数的实部、虚部分别对应两个实的简并模式。

若初位移、初速度分别为$\vct q_0,\vct w_0$，表中任一非零频模式$\vct a_s$的系数为
\[
 Q_s(t)=\frac{\vct a_s\cdot\vct q_0}{\vct a_s\cdot\vct a_s}\cos\omega_st
 +\frac{\vct a_s\cdot\vct w_0}{\omega_s\vct a_s\cdot\vct a_s}\sin\omega_st.
\]
零频模式则有$Q_0(t)=\tfrac14\sum q_{i0}+\tfrac14\sum w_{i0}\,t$。把$Q_s(t)\vct a_s$相加，便得到满足任意小初态的完整运动。图\ref{fig:hw4-1-s}中的箭头均按绕环的同一正向判断，不能直接比较箭头在纸面上的左右。

\end{solution}



\begin{example}[两定滑轮与三等质量物块]\label{ex:hw4-2}

图\ref{fig:hw4-2-q}中两个光滑定滑轮等高，三物块质量均为$m$，中间物块由两条不可伸长轻绳悬挂，两端各挂一个物块。平衡时两段斜绳长均为$l$，与竖直夹角$60^\circ$。求平面内的两个小振动模式。

\end{example}


\begin{figure}[H]\centering
\begin{tikzpicture}[>=Stealth,line cap=round,line join=round]
\coordinate(A)at(-2,0);\coordinate(B)at(2,0);\coordinate(C)at(0,-1.15);\draw(A)circle(.15);\draw(B)circle(.15);\draw[curve](-2.15,-1.8)--(-2.15,0)--(A)--(C)--(B)--(2.15,0)--(2.15,-1.8);\draw[fill=white](-2.4,-2.2)rectangle(-1.9,-1.8);\node at(-2.15,-2){$m$};\draw[fill=white](1.9,-2.2)rectangle(2.4,-1.8);\node at(2.15,-2){$m$};\draw[fill=white](-0.25,-1.3499999999999999)rectangle(0.25,-0.95);\node at(0,-1.15){$m$};\node[above]at(A){$A$};\node[above]at(B){$B$};\node at(-.75,-.3){$l$};\node at(.75,-.3){$l$};
\end{tikzpicture}
\caption{两定滑轮悬挂三个等质量物块}\label{fig:hw4-2-q}
\end{figure}

\begin{solution}
取中间物块相对平衡位置的位移为$(x,y)$，$y$向下。两斜段长度为
\[
 r_{1,2}=\sqrt{(\pm a+x)^2+(h+y)^2},\quad
 a=\sqrt3l/2,\ h=l/2.
\]
图\ref{fig:hw4-2-s}两端物块向下位移$z_i=-(r_i-l)$。一阶关系给出
\[
 z_1=-\frac{\sqrt3}{2}x-\frac12y,\quad
 z_2=\frac{\sqrt3}{2}x-\frac12y.
\]

\begin{figure}[H]\centering
\begin{tikzpicture}[>=Stealth,line cap=round,line join=round]
\coordinate(A)at(-2,0);\coordinate(B)at(2,0);\coordinate(C)at(0,-1.15);\draw(A)circle(.15);\draw(B)circle(.15);\draw[curve](-2.15,-1.8)--(-2.15,0)--(A)--(C)--(B)--(2.15,0)--(2.15,-1.8);\draw[fill=white](-2.4,-2.2)rectangle(-1.9,-1.8);\node at(-2.15,-2){$m$};\draw[fill=white](1.9,-2.2)rectangle(2.4,-1.8);\node at(2.15,-2){$m$};\draw[fill=white](-0.25,-1.3499999999999999)rectangle(0.25,-0.95);\node[left]at(-.4,-1.45){$m$};\node[above]at(A){$A$};\node[above]at(B){$B$};\node at(-.75,-.3){$l$};\node at(.75,-.3){$l$};\draw[axis](0,-1.15)--(1,-1.15)node[above]{$x$};\draw[axis](0,-1.15)--(0,-2.2)node[left]{$y$};\node at(0,.65){$a=\sqrt3l/2,\quad h=l/2$};
\end{tikzpicture}
\caption{中央物块的两个独立小位移}\label{fig:hw4-2-s}
\end{figure}

把根号写为$r_{1,2}=l\sqrt{1+\varepsilon_{1,2}}$，其中
\[
\varepsilon_{1,2}=\frac{\pm\sqrt3x+y}{l}+\frac{x^2+y^2}{l^2}.
\]
用$\sqrt{1+\varepsilon}=1+\varepsilon/2-\varepsilon^2/8+O(\varepsilon^3)$展开；二次项中的$\varepsilon^2$只需保留$(\pm\sqrt3x+y)^2/l^2$。合并同次项后，绳长为
\[
 r_{1,2}=l\pm\frac{\sqrt3}{2}x+\frac12y
 +\frac{x^2+3y^2\mp2\sqrt3xy}{8l}+O(3).
\]
动能只需一阶位移关系；势能的一次项互相抵消后，二阶绳长才产生回复作用。

由此$T=\frac m2(\frac52\dot x^2+\frac32\dot y^2)$。势能须把绳长展开到二阶，得到
\[
 U=-mg(y+z_1+z_2)
 =\frac{mg}{4l}x^2+\frac{3mg}{4l}y^2+\text{常量}.
\]
分别对二次动能、势能求导，有
\[
\frac{\dd}{\dd t}\frac{\partial T}{\partial\dot x}
=\frac{5m}{2}\ddot x,\quad
\frac{\partial U}{\partial x}=\frac{mg}{2l}x;
\qquad
\frac{\dd}{\dd t}\frac{\partial T}{\partial\dot y}
=\frac{3m}{2}\ddot y,\quad
\frac{\partial U}{\partial y}=\frac{3mg}{2l}y.
\]
因此两方程为$\ddot x+gx/(5l)=0$与$\ddot y+gy/l=0$。各自的特征根为$\pm\mathrm i\sqrt{g/(5l)}$和$\pm\mathrm i\sqrt{g/l}$，故
\[
 \boxed{\omega_x=\sqrt{g/(5l)},\qquad\omega_y=\sqrt{g/l}}.
\]
水平模式的两端物块反向升降；竖直模式的两端物块同向，且中间物块位移与它们的位移方向相反、大小为两倍。若将绳长只展开到一阶，势能中这些回复项就会遗漏。


对水平纯模式取$y=0$，两端速度为$\dot z_{1,2}=\mp\sqrt3\dot x/2$；对竖直纯模式取$x=0$，两端速度均为$-\dot y/2$。这两组位移比例给出了图\ref{fig:hw4-2-s}中三个物块在每个模式下的实际运动。

\end{solution}



\begin{example}[半密度矩形长木条的浮动]\label{ex:hw4-4}

均匀矩形长木条密度为水的一半，横截面水平宽度$a$、竖直高度$b$，长度$L$。木条浮在水面。(1)求水平平衡稳定的条件；(2)求上下平动与绕长轴摆动的周期。忽略流体附加质量及波浪辐射，按静水压力计算回复作用。

\end{example}


\begin{figure}[H]\centering
\begin{tikzpicture}[>=Stealth,line cap=round,line join=round]
\draw[thick](-2,0)--(2.5,0);\fill[pattern=north east lines](-2,-1.6)rectangle(2.5,-1.5);\draw[curve](-1,-.7)rectangle(1,.7);\draw[guide](-1,0)--(1,0);\fill(0,0)circle(1pt)node[above left]{$C$};\node at(2,0.35){水面};\draw[<->](-1,-1.05)--(1,-1.05)node[midway,below]{$a$};\draw[<->](1.4,-.7)--(1.4,.7)node[midway,right]{$b$};
\end{tikzpicture}
\caption{半密度木条的平衡吃水}\label{fig:hw4-4-q}
\end{figure}

\begin{solution}
平衡时吃水$b/2$，质量$m=\rho_wabL/2$。图\ref{fig:hw4-4-s}中下沉位移$z$使排水体积增加$aLz$，故
\[
 m\ddot z=-\rho_wgaLz,\qquad
 \boxed{T_z=2\pi\sqrt{b/(2g)}}.
\]

\begin{figure}[H]\centering
\begin{tikzpicture}[>=Stealth,line cap=round,line join=round]
\draw[thick](-2.1,0)--(2.1,0);\node[right]at(2.1,0){水面};
\begin{scope}\clip(-2,-1.6)rectangle(2,0);\begin{scope}[rotate=15]\fill[gray!18](-1.2,-.7)rectangle(1.2,.7);\end{scope}\end{scope}
\begin{scope}[rotate=15]\draw[curve](-1.2,-.7)rectangle(1.2,.7);\draw[guide](0,-.85)--(0,1.7);\end{scope}
\fill(0,0)circle(1.5pt)node[above right,fill=white,inner sep=1pt]{$C$};
\fill(-.09,-.35)circle(1.2pt)node[left,fill=white,inner sep=1pt]{$B$};\fill(-.09,.336)circle(1.2pt)node[left,fill=white,inner sep=1pt]{$M$};
\draw[axis](-.09,-.35)--(-.09,1.25)node[right]{$F_b$};\draw[axis](0,0)--(0,-1.15)node[right]{$mg$};
\draw[guide](0,0)--(0,1.65);\draw(0,1.25)arc(90:105:1.25);\node at(-.5,1.45){$\theta$};
\node at(0,-1.85){$x_B\simeq-CM\theta$};
\begin{scope}[shift={(5.1,0)}]\draw[thick](-1.2,0)--(1.2,0);\draw[curve](-.8,-1.2)rectangle(.8,.2);
\draw[guide](-.8,-.7)--(.8,-.7);\draw[<->](1.2,-.7)--(1.2,-1.2)node[midway,right]{$z$};
\node at(0,.9){下沉分量};\node at(0,-1.6){$\Delta V=aLz$};\end{scope}
\end{tikzpicture}
\caption{倾斜时浮心偏移产生回复力矩，下沉时排水体积增加}\label{fig:hw4-4-s}
\end{figure}

小倾角 $\theta$ 下，浮心初始在质心下方 $b/4$。水线面的面积二阶矩为 $L\int_{-a/2}^{a/2}x^2\dd x=L[x^3/3]_{-a/2}^{a/2}=La^3/12$，除以排水体积 $abL/2$，得到稳心半径 $a^2/(6b)$。因此稳心高为
\[
 CM=\frac{a^2}{6b}-\frac b4.
\]
浮心的横向偏移也可由排水区域积分求出。若截面逆时针倾斜$\theta>0$，在随木条转动的坐标中，水面为$y=-\theta x+O(\theta^2)$，浮心的横坐标为
\[
 x_B^{\rm body}=\frac{-\theta\int_{-a/2}^{a/2}x^2\dd x}{ab/2}
 =-\frac{a^2}{6b}\theta.
\]
转换回地面坐标时还要计入浮心原在质心下方$b/4$，故$x_B=-CM\theta$。图\ref{fig:hw4-4-s}中的浮力作用线通过偏移后的$B$，其关于$C$的力矩为$-mgCM\theta$。$I_C=m(a^2+b^2)/12$，故
\[
 \boxed{\frac ba<\sqrt{\frac23},\qquad
 T_\theta=2\pi\sqrt{\frac{b(a^2+b^2)}{g(2a^2-3b^2)}}}.
\]
这里的开区间保证非零的线性回复频率。等号处可继续保留高阶项：只要水线仍切过两侧边，水面在随体坐标中为 $y=-x\tan\theta$，排水截面面积仍为 $ab/2$。横向一阶矩只需将前式的 $\theta$ 换为 $\tan\theta$；竖向一阶矩为
\[
\int_{-a/2}^{a/2}\int_{-b/2}^{-x\tan\theta}y\dd y\dd x
=\frac12\int_{-a/2}^{a/2}\left(x^2\tan^2\theta-\frac{b^2}4\right)\dd x
=\frac{a^3}{24}\tan^2\theta-\frac{ab^2}8.
\]
这里的内层原函数为 $y^2/2$，外层使用上文已算出的 $x^2$ 积分。除以排水截面面积，浮心在随体坐标中的位置为
\[
 x_B^{\rm body}=-\frac{a^2}{6b}\tan\theta,\qquad
 y_B^{\rm body}=-\frac b4+\frac{a^2}{12b}\tan^2\theta.
\]
将其转回地面系，浮力矩为
\[
 \tau_C=-mg\sin\theta\left[CM+\frac{a^2}{12b}\tan^2\theta\right].
\]
因此 $CM=0$ 时，$\tau_C=-mga^2\theta^3/(12b)+O(\theta^5)$。由 $\dd U/\dd\theta=-\tau_C$，积分得
\[
U(\theta)-U(0)=\frac{mga^2}{12b}\int_0^\theta\xi^3\dd\xi+O(\theta^6)
=\frac{mga^2}{12b}[\xi^4/4]_0^\theta+O(\theta^6).
\]
势能的首个转动项为 $mga^2\theta^4/(48b)>0$。在题设静水模型中，完整的局部稳定条件为 $b/a\le\sqrt{2/3}$；等号处没有有限的线性周期，不能套用上面的 $T_\theta$。超过该比值时水平构形不稳定。

对称性使上下与转动在一阶近似中解耦。实际流体的附加惯性会改变周期，上述周期取决于题设采用的静水模型。


\end{solution}



\begin{example}[两均匀杆组成的复摆]\label{ex:hw4-6}

两根均匀细杆长度均为$l$、质量均为$m$，第一杆上端铰接于固定点$O$，下端与第二杆上端在$A$处铰接。两杆在同一竖直面内绕向下平衡位置作小振动，求简正模式。

\end{example}


\begin{figure}[H]\centering
\begin{tikzpicture}[>=Stealth,line cap=round,line join=round]
\draw[thick](-.5,0)--(.7,0);\coordinate(O)at(0,0);\coordinate(A)at(.4,-1.7);\coordinate(B)at(1.1,-3.3);\draw[curve](O)--(A)--(B);\fill(A)circle(2pt)node[right]{$A$};\node[above]at(O){$O$};\node at(-.3,-.8){$\theta_1$};\node at(.2,-2.5){$\theta_2$};\draw[guide](O)--(0,-1.7);\draw[guide](A)--(.4,-3.3);
\end{tikzpicture}
\caption{两均匀杆铰接成的复摆}\label{fig:hw4-6-q}
\end{figure}

\begin{solution}
图\ref{fig:hw4-6-s}中$\theta_1,\theta_2$均相对竖直向下量取，为两杆的绝对倾角。两杆质心位置分别为$(l\sin\theta_1/2,-l\cos\theta_1/2)$与$(l\sin\theta_1+l\sin\theta_2/2,-l\cos\theta_1-l\cos\theta_2/2)$。

\begin{figure}[H]\centering
\begin{tikzpicture}[>=Stealth,line cap=round,line join=round]
\draw[thick](-.5,0)--(.7,0);\coordinate(O)at(0,0);\coordinate(A)at(.4,-1.7);\coordinate(B)at(1.1,-3.3);\draw[curve](O)--(A)--(B);\fill(A)circle(2pt)node[right]{$A$};\node[above]at(O){$O$};\node at(-.3,-.8){$\theta_1$};\node at(.2,-2.5){$\theta_2$};\draw[guide](O)--(0,-1.7);\draw[guide](A)--(.4,-3.3);\draw[axis](.2,-.85)--(1,-.85)node[above]{$v_{C1,x}^{(1)}$};\draw[axis](.75,-2.5)--(1.65,-2.5)node[above]{$v_{C2,x}^{(1)}$};\node at(3,-1.5){$I_C=ml^2/12$};
\end{tikzpicture}
\caption{质心速度取到位移一阶，水平分量进入二阶动能}\label{fig:hw4-6-s}
\end{figure}

第一杆关于固定上端的动能为$ml^2\dot\theta_1^2/6$。第二杆质心的水平速度在一阶近似中为$l\dot\theta_1+l\dot\theta_2/2$，再加上自身转动动能，得
\[
 T_2=\frac{ml^2}{2}\dot\theta_1^2+
 \frac{ml^2}{2}\dot\theta_1\dot\theta_2+
 \frac{ml^2}{6}\dot\theta_2^2.
\]
交叉项来自第二杆质心同时受到上杆运动和自身转动的影响。

将质心平动及各杆绕质心的转动相加，保留二次项：
\[
 T=\frac{ml^2}{2}\left(\frac43\dot\theta_1^2+\dot\theta_1\dot\theta_2+\frac13\dot\theta_2^2\right),
 \quad U=\frac{mgl}{2}\left(\frac32\theta_1^2+\frac12\theta_2^2\right).
\]
由于二次动能不显含角坐标，拉格朗日方程中的惯性项为
\[
\frac{\dd}{\dd t}\frac{\partial T}{\partial\dot\theta_1}
=ml^2\left(\frac43\ddot\theta_1+\frac12\ddot\theta_2\right),\qquad
\frac{\dd}{\dd t}\frac{\partial T}{\partial\dot\theta_2}
=ml^2\left(\frac12\ddot\theta_1+\frac13\ddot\theta_2\right).
\]
势能导数分别为$3mgl\theta_1/2$、$mgl\theta_2/2$。令$\theta_j=a_j\cos\omega t$，将$\ddot\theta_j=-\omega^2\theta_j$代入两个方程。再置$\lambda=\omega^2l/g$，非零振幅要求
\[
 \det\begin{pmatrix}3/2-4\lambda/3&-\lambda/2\\-\lambda/2&1/2-\lambda/3\end{pmatrix}=0
 \quad\Longrightarrow\quad7\lambda^2-42\lambda+27=0.
\]
故
\[
 \boxed{\omega_\pm^2=\left(3\pm\frac6{\sqrt7}\right)\frac gl,
 \qquad\frac{a_2}{a_1}=\frac{3\lambda_\pm}{3-2\lambda_\pm}}.
\]
低频根给出同向摆动，高频根给出反向摆动。若把$\theta_2$误当成两杆夹角，质心速度和动能交叉项都会改变；方程与图中的角度定义须始终一致。


展开行列式，得
\[
 \frac34-\frac76\lambda+\left(\frac49-\frac14\right)\lambda^2=0.
\]
乘以36即为$7\lambda^2-42\lambda+27=0$。求振幅比时采用第二行$-(\lambda/2)a_1+(1/2-\lambda/3)a_2=0$，便得$3\lambda/(3-2\lambda)$。低频根$\lambda_-\simeq0.7322$的振幅比约为$1.4305$，高频根$\lambda_+\simeq5.2678$的振幅比约为$-2.0972$。图\ref{fig:hw4-6-s}中的两杆倾角都以固定竖直线为零点，故正比值表示同向，负比值表示反向。

若给定初态，也应使用质量矩阵$ml^2\left(\begin{smallmatrix}4/3&1/2\\1/2&1/3\end{smallmatrix}\right)$投影。第二杆随第一杆运动，使质量矩阵具有非零的非对角元；仅按两杆质量相等作普通投影并不成立。

\end{solution}



\begin{example}[摆杆、弹簧与悬重的耦合]\label{ex:hw4-7}

长$l$、质量$3m$的均匀杆上端固定铰接，平衡时竖直。杆下端左侧通过劲度系数$k$的水平弹簧连接墙，右侧通过同样弹簧及绕定滑轮的轻绳连接质量$m$的悬重。弹簧预伸长使题设构形平衡，求小振动简正模式。

\end{example}


\begin{figure}[H]\centering
\begin{tikzpicture}[>=Stealth,line cap=round,line join=round]
\draw[thick](-.5,0)--(.5,0);\draw[curve](0,0)--(0,-2);\draw[spring](-2,-2)--(0,-2);\draw[spring](0,-2)--(1.75,-2);\draw(2,-2.15)circle(.15);\draw(1.75,-2)--(2,-2)arc(90:0:.15)--(2.15,-3.3);\draw[thick](-2,-2.5)--(-2,-1.5);\draw[fill=white](1.9,-3.7)rectangle(2.4,-3.3);\node at(2.15,-3.5){$m$};\node at(.6,-1){$3m,l$};\node at(-1,-1.5){$k$};\node at(1,-1.5){$k$};
\end{tikzpicture}
\caption{弹簧连接的均匀杆与悬挂物块}\label{fig:hw4-7-q}
\end{figure}

\begin{solution}
令杆向右小转角为$\theta$、悬重相对平衡位置向下位移为$x=l\varphi$。图\ref{fig:hw4-7-s}右弹簧的伸长增量为$x-l\theta$，左弹簧为$l\theta$。

\begin{figure}[H]\centering
\begin{tikzpicture}[>=Stealth,line cap=round,line join=round]
\draw[thick](-.5,0)--(.5,0);\draw[curve](0,0)--(0,-2);\draw[spring](-2,-2)--(0,-2);\draw[spring](0,-2)--(1.75,-2);\draw(2,-2.15)circle(.15);\draw(1.75,-2)--(2,-2)arc(90:0:.15)--(2.15,-3.3);\draw[thick](-2,-2.5)--(-2,-1.5);\draw[fill=white](1.9,-3.7)rectangle(2.4,-3.3);\node at(2.15,-3.5){$m$};\node at(.6,-1){$3m,l$};\node at(-1,-1.5){$k$};\node at(1,-1.5){$k$};\draw[axis](0,-2.35)--(.8,-2.35)node[below]{$l\theta$};\draw[axis](2.8,-3)--(2.8,-3.8)node[right]{$x$};
\end{tikzpicture}
\caption{杆端水平位移与物块竖直位移}\label{fig:hw4-7-s}
\end{figure}

杆的惯量为$(3m)l^2/3=ml^2$。以平衡位置为展开点，一次项由预拉力和重力抵消，二次式为
\[
 T=\frac12ml^2(\dot\theta^2+\dot\varphi^2),\quad
 U=\frac{3mgl}{4}\theta^2+\frac{kl^2}{2}[\theta^2+(\varphi-\theta)^2].
\]
对势能求导，
\[
\frac{\partial U}{\partial\theta}
=\frac{3mgl}{2}\theta+kl^2(2\theta-\varphi),\qquad
\frac{\partial U}{\partial\varphi}=kl^2(\varphi-\theta).
\]
动能分别给出$ml^2\ddot\theta$、$ml^2\ddot\varphi$，故设$a=3g/(2l)$、$b=k/m$后，两方程为
\[
\ddot\theta+(a+2b)\theta-b\varphi=0,\qquad
\ddot\varphi+b\varphi-b\theta=0.
\]
取$\theta=A\cos\omega t$、$\varphi=B\cos\omega t$，代入得
\[
 \begin{pmatrix}a+2b-\omega^2&-b\\-b&b-\omega^2\end{pmatrix}
 \binom{A}{B}=0.
\]
令$\lambda=\omega^2$，行列式逐项展开为
\[
(a+2b-\lambda)(b-\lambda)-b^2
=\lambda^2-(a+3b)\lambda+ab+b^2=0.
\]
二次方程的判别式为$(a+3b)^2-4(ab+b^2)=(a+b)^2+4b^2$。用求根公式，并从矩阵第二行$-bA+(b-\lambda)B=0$求振幅比，得到
\[
 \boxed{\omega_\pm^2=\frac{a+3b\pm\sqrt{(a+b)^2+4b^2}}2,
 \qquad B/A=\frac{b}{b-\omega_\pm^2}}.
\]
两个频率平方之和为$a+3b>0$，乘积为$ab+b^2>0$，且判别式为$(a+b)^2+4b^2$，故$k>0$时两根均为正。低频模式中$B/A>0$，杆端向右时悬重向下；高频模式中两者反向。

这里$B/A$是$\varphi/\theta$的振幅比，实际悬重位移与角振幅之比为$lB/A$。只有再给定$k/m=g/l$之类的附加关系，才能进一步简化频率；这里保留两个独立参数。


由特征方程$(a+2b-\omega^2)(b-\omega^2)-b^2=0$，可知低频根小于$b$而高频根大于$b$，因而上述振幅比的符号与图示一致。求弱弹簧极限时，令$\varepsilon=b/a\ll1$。频率公式的根式为
\[
\sqrt{(a+b)^2+4b^2}=a\sqrt{1+2\varepsilon+5\varepsilon^2}.
\]
在$\sqrt{1+z}=1+z/2-z^2/8+O(z^3)$中代入$z=2\varepsilon+5\varepsilon^2$，二次系数为$5/2-4/8=2$，故根式等于$a+b+2b^2/a+O(b^3/a^2)$。分别代入两根，得到
\[
 \omega_-^2=b-\frac{b^2}{a}+O(b^3/a^2),\qquad
 \omega_+^2=a+2b+\frac{b^2}{a}+O(b^3/a^2).
\]
前者以悬重运动为主，杆的振幅较小；后者趋近单独均匀杆的物理摆频率$\sqrt{3g/(2l)}$。这里极限$b\to0^+$只描述二次振动方程的参数变化；若实际装置不能维持题设预拉力，则须另行确定新的平衡构形。

\end{solution}



\begin{example}[正方形顶点电荷的平面稳定性]\label{ex:hw4-8}

边长$L$的正方形两条对角线端点上，分别放置两对固定正电荷$Q$与$\alpha Q$，$\alpha>0$。质量$m$、电荷$q>0$的粒子静止在中心，运动限制在正方形平面内。求稳定平衡的$\alpha$范围及两个方向的小振动周期。

\end{example}


\begin{figure}[H]\centering
\begin{tikzpicture}[>=Stealth,line cap=round,line join=round]
\draw[curve](-1.6,0)--(0,1.6)--(1.6,0)--(0,-1.6)--cycle;\foreach \x/\y/\lab in {-1.6/0/Q,1.6/0/Q,0/1.6/\alpha Q,0/-1.6/\alpha Q}{\fill(\x,\y)circle(2pt)node[shift={({\x*.18},{\y*.18})}]{$\lab$};}\fill(0,0)circle(2pt)node[below right]{$q,m$};
\end{tikzpicture}
\caption{正方形四顶点的交替电荷}\label{fig:hw4-8-q}
\end{figure}

\begin{solution}
图\ref{fig:hw4-8-s}沿两条对角线取$x,y$轴，两对电荷坐标分别为$(\pm a,0)$和$(0,\pm a)$，其中$a=L/\sqrt2$。为展开每一对电荷的势，令$\vct e$沿该对电荷的轴线、$r^2=x^2+y^2$。对每一项使用$(1+z)^{-1/2}=1-z/2+3z^2/8+O(z^3)$，有
\[
\frac1{|a\vct e\mp\vct r|}
=\frac1a\left[1\pm\frac{\vct e\cdot\vct r}{a}
+\frac{3(\vct e\cdot\vct r)^2-r^2}{2a^2}\right]+O(r^3).
\]
正负两项相加消去奇次项，再分别取$\vct e=\vct e_x,\vct e_y$并乘相应电荷量，得到
\[
 U=\frac{2k_eqQ(1+\alpha)}a
 +\frac{k_eqQ}{a^3}[(2-\alpha)x^2+(2\alpha-1)y^2]+O(|\vct r|^4).
\]

\begin{figure}[H]\centering
\begin{tikzpicture}[>=Stealth,line cap=round,line join=round]
\draw[curve](-1.6,0)--(0,1.6)--(1.6,0)--(0,-1.6)--cycle;\foreach \x/\y in {-1.6/0,1.6/0,0/1.6,0/-1.6}{\fill(\x,\y)circle(2pt);}\node at(-1.95,.2){$Q$};\node at(1.95,.2){$Q$};\node[right]at(.16,1.95){$\alpha Q$};\node at(0,-2){$\alpha Q$};\fill(0,0)circle(2pt)node[below right]{$q,m$};\draw[axis](0,0)--(2.5,0)node[above]{$x$};\draw[axis](0,0)--(0,2.4)node[right]{$y$};\node at(3,-1){$a=L/\sqrt2$};
\end{tikzpicture}
\caption{平面内两条对角方向取为坐标轴}\label{fig:hw4-8-s}
\end{figure}

一次项因成对对称抵消，$xy$项也为零。两个线性回复系数都为正的条件是$1/2<\alpha<2$。这个开区间给出两个非零的小振动频率；判断稳定平衡的完整范围还须检查端点。
对二次势能分别求导，线性运动方程为
\[
m\ddot x=-\frac{\partial U_2}{\partial x}
=-\frac{2k_eqQ(2-\alpha)}{a^3}x,\qquad
m\ddot y=-\frac{\partial U_2}{\partial y}
=-\frac{2k_eqQ(2\alpha-1)}{a^3}y.
\]
各自除以$m$，括号内位移的系数就是$\omega_x^2,\omega_y^2$。在此开区间内，由$T=2\pi/\omega$得到
\[
 \boxed{T_x=2\pi\sqrt{\frac{ma^3}{2k_eqQ(2-\alpha)}},\quad
 T_y=2\pi\sqrt{\frac{ma^3}{2k_eqQ(2\alpha-1)}}}.
\]
沿$x$轴，四电荷的精确势能可先合并为
\[
U(x,0)=k_eqQ\left[\frac{2a}{a^2-x^2}
+\frac{2\alpha}{\sqrt{a^2+x^2}}\right].
\]
第一项用$(1-z)^{-1}=1+z+z^2+O(z^3)$，第二项用上述二项式，均取$z=x^2/a^2$。减去中心值后，
\[
U(x,0)-U(0,0)=\frac{k_eqQ}{a}
\left[(2-\alpha)\frac{x^2}{a^2}
+\left(2+\frac{3\alpha}{4}\right)\frac{x^4}{a^4}\right]+O(x^6).
\]
在$\alpha=2$时，二次项消失，留下
\[
 U(x,0)-U(0,0)=\frac{7k_eqQ}{2a^5}x^4+O(x^6)>0.
\]
沿$y$轴则有$U(0,y)=k_eqQ[2\alpha a/(a^2-y^2)+2/\sqrt{a^2+y^2}]$；按相同的两种二项式展开，二次项系数为$2\alpha-1$、四次项系数为$2\alpha+3/4$。代入$\alpha=1/2$，得到
\[
 U(0,y)-U(0,0)=\frac{7k_eqQ}{4a^5}y^4+O(y^6)>0.
\]
还须证明任意平面扰动都会增加势能，而不只检查两条坐标轴。以$\alpha=2$为例，记$A_y=3k_eqQ/a^3>0$、$B_x=7k_eqQ/(2a^5)>0$。势能关于$x,y$分别为偶函数，故在原点邻域可写成
\[
U-U(0,0)=A_y y^2+B_x x^4+R(x,y),
\]
其中未写出的四次项只有$x^2y^2,y^4$，更高项从六次开始。因此存在有限常数$C_1,C_2,C_3$，使
\[
|R|\le C_1x^2y^2+C_2y^4+C_3(x^2+y^2)^3.
\]
在$|x|,|y|\le\delta$内，分别有$x^2y^2\le\delta^2y^2$、$y^4\le\delta^2y^2$，并且
\[
(x^2+y^2)^3=x^6+3x^4y^2+3x^2y^4+y^6
\le4\delta^2x^4+4\delta^4y^2.
\]
取$\delta$足够小，可使$|R|\le B_xx^4/2+A_yy^2/2$，于是
\[
U-U(0,0)\ge\frac{A_y}{2}y^2+\frac{B_x}{2}x^4>0
\quad\text{（非零扰动）}.
\]
这证明$\alpha=2$仍为严格局部极小。$\alpha=1/2$时交换$x,y$，并将势能乘一个正常数，即回到同一论证。因此完整的平面稳定范围为
\[
 \boxed{\tfrac12\le\alpha\le2}.
\]
端点的退化方向由三次回复力控制，振动周期随振幅趋于零而发散；上述有限周期公式只用于开区间。本题明确限制在平面内；离开平面时四个斥力势都随$|z|$增加而降低，中心并非三维稳定平衡。


\end{solution}



\begin{example}[圆弧槽内的纯滚动振动]\label{deep4:problem-example-1}\label{deep4:ex:d7-e1}
半径为 $R$ 的固定圆弧槽开口向上，半径为 $a<R$、质量为 $m$ 的圆形刚体在槽内无滑动滚动，绕自身轴的转动惯量为 $I=\beta ma^2$。以最低点为 $\theta=0$，刚体中心与槽圆心的连线相对竖直向下方向的转角为 $\theta$。求小振动频率。若从 $\theta=\Theta$ 静止释放，$0<\Theta<\pi/2$，求维持全过程纯滚动所需的最小静摩擦系数。
\end{example}
\begin{figure}[H]\centering
\begin{tikzpicture}
\draw[curve](-2.954,-.521)arc(-170:-10:3);\fill(0,0)circle(1.5pt)node[above]{$O$};\draw[guide](0,0)--(0,-1.25);\coordinate(C)at({2.5*sin(30)},{-2.5*cos(30)});\coordinate(S)at({3*sin(30)},{-3*cos(30)});\draw[curve](C)circle(.5);\fill(C)circle(1.7pt);\draw[guide](C)--(.1,-1.9)node[left]{$C$};\fill(S)circle(1.5pt)node[below right]{$S$};\draw[guide](0,0)--(S);\node[left]at(.65,-1.15){$\ell$};\draw[<->](C)--++(.5,0)node[midway,above]{$a$};\draw(0,-.65)arc(-90:-60:.65);\node at(.5,-.7){$\theta$};\node at(-1.9,-.6){固定圆弧槽};
\end{tikzpicture}
\caption{中心C沿半径 $\ell=R-a$ 的圆运动，接触点S位于OC的延长线上。}\label{deep4:fig:d7-e1-q}
\end{figure}
\begin{solution}
图\ref{deep4:fig:d7-e1-q}中，中心速率为 $\ell|\dot\theta|$，其中 $\ell=R-a$。接触点相对固定槽瞬时静止，因此刚体自转角速度大小为 $\ell|\dot\theta|/a$。总动能与势能为
\[
T=\frac12m\ell^2\dot\theta^2+\frac12I\frac{\ell^2\dot\theta^2}{a^2}
=\frac12m(1+\beta)\ell^2\dot\theta^2,
\qquad U=mg\ell(1-\cos\theta).
\]
为求转折点处也适用的运动方程，按图\ref{deep4:fig:d7-e1-s}取增大 $\theta$ 的切向为正，用滚动约束消去自转角加速度，平动与转动方程成为
\[
m\ell\ddot\theta=-mg\sin\theta+f,\qquad
\beta m\ell\ddot\theta=-f.
\]
静摩擦同时改变平动和自转，虽不耗散能量，仍会影响振动频率。
两式相加得到
\[
(1+\beta)\ell\ddot\theta+g\sin\theta=0.
\]
在线性近似中 $\sin\theta\simeq\theta$，故
\[
\boxed{\omega=\sqrt{\frac{g}{(1+\beta)(R-a)}}.}
\]
等效惯量包含平动与自转两部分，所得频率是同轨道上光滑滑动质点频率的 $1/\sqrt{1+\beta}$。
\begin{figure}[H]\centering
\begin{tikzpicture}
\fill(0,0)circle(2pt)node[below left]{$C$};\draw[guide](-1,1.732)--(1,-1.732);\draw[axis](0,0)--(-.9,1.559)node[above]{$N$};\draw[axis](0,0)--(0,-1.8)node[below]{$mg$};\draw[axis](0,0)--(1.2,.693)node[right]{$f$};\node[align=left]at(4.1,.7){$\theta>0$ 时\\摩擦沿槽向上\\正切向取增大 $\theta$ 的方向};
\end{tikzpicture}
\caption{法向力指向槽圆心，摩擦提供滚动所需的自转角加速度。}\label{deep4:fig:d7-e1-s}
\end{figure}
按图\ref{deep4:fig:d7-e1-s}取沿槽向上为切向正方向。平动方程为 $m\ell\ddot\theta=-mg\sin\theta+f$，代入已求出的 $\ddot\theta$，得
\[
f=\frac{\beta}{1+\beta}mg\sin\theta.
\]
中心的向心加速度给出 $N-mg\cos\theta=m\ell\dot\theta^2$。从转折点释放后的能量关系为
\[
\ell\dot\theta^2=\frac{2g}{1+\beta}(\cos\theta-\cos\Theta),
\quad
N=\frac{mg}{1+\beta}[(\beta+3)\cos\theta-2\cos\Theta].
\]
由于 $|\theta|\le\Theta<\pi/2$，$N>0$，接触始终维持。摩擦比值为
\[
\frac{|f|}{N}=\frac{\beta|\sin\theta|}{(\beta+3)\cos\theta-2\cos\Theta}.
\]
该比值与习题\ref{deep1:exercise:d1-p5}的摩擦比是同一函数，只需取 $\kappa=\beta$、$\theta_0=\Theta$。使用那里已完整展开的商导数，在 $0\le\theta\le\Theta$ 上，
\[
\frac{\dd}{\dd\theta}\frac{|f|}{N}
=\frac{\beta[(\beta+3)-2\cos\Theta\cos\theta]}
{[(\beta+3)\cos\theta-2\cos\Theta]^2}>0.
\]
其余半程关于最低点对称，故最大值在两个转折点取得。因此
\[
\boxed{\mu_{s,\min}=\frac{\beta}{1+\beta}\tan\Theta.}
\]
摩擦条件在给定的 $0<\Theta<\pi/2$ 范围内未使用小角近似，频率公式则要求 $\Theta\ll1$。对实心球，$\beta=2/5$，小振动频率为 $\sqrt{5g/[7(R-a)]}$。
\end{solution}



\section{振型分解、简正坐标与初始条件}

频率不能单独确定一次实际运动，初位移和初速度还要分别分解到振型中。不同质量的系统采用质量加权的正交关系，再由模态运动恢复各物体的位移。



\subsection{振型正交与初始条件分解}
若 $K\vct u_r=\omega_r^2M\vct u_r$，由矩阵对称性可得
\[
\vct u_s^{\mathsf T}K\vct u_r
=\omega_r^2\vct u_s^{\mathsf T}M\vct u_r
=\omega_s^2\vct u_s^{\mathsf T}M\vct u_r.
\]
当两频率不同，
\begin{equation}
\vct u_s^{\mathsf T}M\vct u_r=0.
\label{sup4:eq:d7-orthogonal}
\end{equation}
正交关系的权重是质量矩阵。质量不相等或动能含交叉项时，投影必须使用这个权重。选取归一化振型 $\vct e_r$，使 $\vct e_s^{\mathsf T}M\vct e_r=\delta_{sr}$；对简并频率，也可在相应子空间内按此内积选取正交基。

写成 $\vct\eta=\sum_r\vct e_r Q_r$，则
\[
\ddot Q_r+\omega_r^2Q_r=0,\qquad
Q_r(0)=\vct e_r^{\mathsf T}M\vct\eta(0),\quad
\dot Q_r(0)=\vct e_r^{\mathsf T}M\dot{\vct\eta}(0).
\]
对 $\omega_r>0$，特征方程 $\lambda^2+\omega_r^2=0$ 有根 $\pm\mathrm i\omega_r$。实通解为 $Q_r=A_r\cos\omega_rt+B_r\sin\omega_rt$，由初态得 $A_r=Q_r(0)$、$B_r=\dot Q_r(0)/\omega_r$，即
\[
Q_r(t)=Q_r(0)\cos\omega_rt+\frac{\dot Q_r(0)}{\omega_r}\sin\omega_rt.
\]
若未归一化，投影式须再除以 $\vct u_r^{\mathsf T}M\vct u_r$。每个模态能量独立守恒，因为 $\dd E_r/\dd t=\dot Q_r(\ddot Q_r+\omega_r^2Q_r)=0$：
\[
E_2-U_* =\frac12\sum_r(\dot Q_r^2+\omega_r^2Q_r^2).
\]
\begin{figure}[H]\centering
\begin{tikzpicture}
\begin{scope}
\draw[axis](-1.7,0)--(1.9,0)node[right]{$x_1$};\draw[axis](0,-1.4)--(0,1.6)node[above]{$x_2$};\draw[axis](0,0)--(1.2,1.2)node[above right]{$\vct u_+$};\draw[axis](0,0)--(1.2,-1.2)node[below right]{$\vct u_-$};\node at(0,-2){等质量的正交振型};
\end{scope}
\begin{scope}[xshift=5.5cm]
\draw[axis](-1.8,0)--(2,0)node[right]{$Q_+$};\draw[axis](0,-1.5)--(0,1.7)node[above]{$Q_-$};\draw[curve](0,0)ellipse(1.3 and .85);\node at(0,-2){模态坐标中的能量椭圆};
\end{scope}
\end{tikzpicture}
\caption{左图以等质量双振子示意振型方向；右图为给定能量、速度为零截面上的椭圆。}\label{sup4:fig:d7-modes}
\end{figure}
图\ref{sup4:fig:d7-modes}的模态分解也可用于冲量后的运动。广义冲量为 $\vct J_q$ 时，$M\Delta\dot{\vct\eta}=\vct J_q$。$\vct J_q$ 须由作用点的虚位移求出；对角坐标，它还包含相应的力臂。


\subsection{简正坐标的构造与退化情形}
\label{sec:complete4}
质量矩阵正定时，可写成$M=S^{\mathsf T}S$，令$\vct z=S\vct\eta$，则
\[
 T_2=\frac12\dot{\vct z}^{\mathsf T}\dot{\vct z},\qquad
 U_2-U_* =\frac12\vct z^{\mathsf T}
 (S^{-\mathsf T}KS^{-1})\vct z.
\]
括号内的实对称矩阵可经正交变换对角化，因而$\omega_j^2$为实数，振型可选成完备的质量正交基。坐标由长度改成角度时，质量矩阵和刚度矩阵须同时变换；单独比较某个矩阵的对角元没有坐标不变的意义。

\subsubsection{由运动和初态求模态能量}
若采用未归一化振型$\vct u_j$，记$M_j=\vct u_j^{\mathsf T}M\vct u_j$，并写$\vct\eta=\sum_j\vct u_jQ_j$。利用正交关系，得到
\[
 E_j=\frac12M_j(\dot Q_j^2+\omega_j^2Q_j^2),\quad
 Q_j(0)=\frac{\vct u_j^{\mathsf T}M\vct\eta(0)}{M_j}.
\]
初速度按相同权重投影。若系统原来静止，受到广义冲量$\vct J_q$，则$\dot Q_j(0^+)=\vct u_j^{\mathsf T}\vct J_q/M_j$。若受冲量的位置是某一振型的节点，该模态的初速度投影为零。

\subsection{例题与模型变式}

\begin{example}[弹簧耦合的双单摆]\label{ex:hw4-3}

两个质量$m$的小球分别由长$l$的轻杆悬挂在等高、相距$l$的两铰点下，球之间用原长$l$、劲度系数$k$的轻弹簧相连。初态为静止平衡，瞬间给$A$球水平速度$v_0$。求小振动范围内两球的运动。

\end{example}


\begin{figure}[H]\centering
\begin{tikzpicture}[>=Stealth,line cap=round,line join=round]
\draw[thick](-.6,0)--(2.6,0);\draw[curve](0,0)--(0,-2);\draw[curve](2,0)--(2,-2);\fill(0,-2)circle(2pt)node[below]{$A,m$};\fill(2,-2)circle(2pt)node[below]{$B,m$};\draw[spring](0,-2)--(2,-2);\node at(1,-1.6){$k,l$};\node at(-.3,-1){$l$};\draw[<->](0,.45)--(2,.45)node[midway,above]{$l$};
\end{tikzpicture}
\caption{等长双摆在平衡位置由弹簧连接}\label{fig:hw4-3-q}
\end{figure}

\begin{solution}
用水平小位移$x_1,x_2$，图\ref{fig:hw4-3-s}中杆角约为$x_i/l$，弹簧伸长为$x_2-x_1$。运动方程为
\[
 m\ddot x_1=-\frac{mg}{l}x_1-k(x_1-x_2),\quad
 m\ddot x_2=-\frac{mg}{l}x_2-k(x_2-x_1).
\]

\begin{figure}[H]\centering
\begin{tikzpicture}[>=Stealth,line cap=round,line join=round]
\draw[thick](-.6,0)--(2.6,0);\draw[curve](0,0)--(0,-2);\draw[curve](2,0)--(2,-2);\fill(0,-2)circle(2pt)node[below]{$A,m$};\fill(2,-2)circle(2pt)node[below]{$B,m$};\draw[spring](0,-2)--(2,-2);\node at(1,-1.6){$k,l$};\node at(-.3,-1){$l$};\draw[<->](0,.45)--(2,.45)node[midway,above]{$l$};\draw[axis](0,-2.4)--(.7,-2.4)node[below]{$x_1$};\draw[axis](2,-2.4)--(2.7,-2.4)node[below]{$x_2$};
\end{tikzpicture}
\caption{两个水平位移确定弹簧伸长量}\label{fig:hw4-3-s}
\end{figure}

相加、相减得$q_s=x_1+x_2$、$q_a=x_1-x_2$，满足
\[
 \ddot q_s+\omega_s^2q_s=0,\quad\omega_s^2=g/l,
 \qquad\ddot q_a+\omega_a^2q_a=0,\quad\omega_a^2=g/l+2k/m.
\]
两方程的特征根分别为$\pm\mathrm i\omega_s$、$\pm\mathrm i\omega_a$，通解写成
\[
q_s=A_s\cos\omega_st+B_s\sin\omega_st,\qquad
q_a=A_a\cos\omega_at+B_a\sin\omega_at.
\]
初态$q_s(0)=q_a(0)=0$给出$A_s=A_a=0$；对通解求导并代入$\dot q_s(0)=\dot q_a(0)=v_0$，得$\omega_sB_s=\omega_aB_a=v_0$。再由$x_1=(q_s+q_a)/2$、$x_2=(q_s-q_a)/2$还原，得到
\[
 \boxed{x_1=\frac{v_0}{2}\left(\frac{\sin\omega_st}{\omega_s}
 +\frac{\sin\omega_at}{\omega_a}\right),\quad
 x_2=\frac{v_0}{2}\left(\frac{\sin\omega_st}{\omega_s}
 -\frac{\sin\omega_at}{\omega_a}\right)}.
\]
由上式可直接验证$x_1(0)=x_2(0)=0$、$\dot x_1(0)=v_0$、$\dot x_2(0)=0$。一个足够保证小振幅的条件是
\[
 \frac{v_0}{2l}\left(\frac1{\omega_s}+\frac1{\omega_a}\right)\ll1.
\]
若不满足这一条件，球的水平位移不再等于$l\theta_i$，弹簧长度也不能只保留$x_2-x_1$。

同向模式不改变弹簧长度，反向模式则同时受到重力和弹簧回复。两频接近时可见缓慢的振幅交换，但精确解仍是上式的两个正弦叠加。


\end{solution}



\begin{example}[三弹簧与不等质量双振子]\label{ex:hw4-5}

两小球在光滑水平面上由三个相同弹簧依次与两墙连接，弹簧劲度系数为$k$。分别求$m_1=m_2=m$及$m_1=m,m_2=2m$时的简正频率与振幅比。

\end{example}


\begin{figure}[H]\centering
\begin{tikzpicture}[>=Stealth,line cap=round,line join=round]
\draw[fill=white](0.75,-0.2)rectangle(1.25,0.2);\node at(1,0){$m_1$};\draw[fill=white](2.75,-0.2)rectangle(3.25,0.2);\node at(3,0){$m_2$};\draw[spring](-.5,0)--(.75,0);\draw[spring](1.25,0)--(2.75,0);\draw[spring](3.25,0)--(4.5,0);\draw[thick](-.5,-.5)--(-.5,.5);\draw[thick](4.5,-.5)--(4.5,.5);\node at(2,.7){三个弹簧均为$k$};
\end{tikzpicture}
\caption{三弹簧连接的两个质量块}\label{fig:hw4-5-q}
\end{figure}

\begin{solution}
图\ref{fig:hw4-5-s}的$x_1,x_2$从各自平衡位置向右量起，回复力为$-2kx_1+kx_2$及$kx_1-2kx_2$。因此
\[
 \det\begin{pmatrix}2k-m_1\omega^2&-k\\-k&2k-m_2\omega^2\end{pmatrix}=0.
\]

\begin{figure}[H]\centering
\begin{tikzpicture}[>=Stealth,line cap=round,line join=round]
\draw[fill=white](0.75,-0.2)rectangle(1.25,0.2);\node at(1,0){$m_1$};\draw[fill=white](2.75,-0.2)rectangle(3.25,0.2);\node at(3,0){$m_2$};\draw[spring](-.5,0)--(.75,0);\draw[spring](1.25,0)--(2.75,0);\draw[spring](3.25,0)--(4.5,0);\draw[thick](-.5,-.5)--(-.5,.5);\draw[thick](4.5,-.5)--(4.5,.5);\node at(2,.7){三个弹簧均为$k$};\draw[axis](1,-.5)--(1.8,-.5)node[below]{$x_1$};\draw[axis](3,-.5)--(3.8,-.5)node[below]{$x_2$};
\begin{scope}[shift={(1,-1.8)}]
\fill(0,0)circle(2pt)node[below]{$m$};\fill(2,0)circle(2pt)node[below]{$2m$};
\draw[axis](0,0)--(.55,0);\draw[axis](2,0)--(2.7513,0);\node[left]at(-.7,0){低频};
\node at(1,-.6){$a_2/a_1=(1+\sqrt3)/2$};\end{scope}
\begin{scope}[shift={(1,-3.2)}]
\fill(0,0)circle(2pt)node[below]{$m$};\fill(2,0)circle(2pt)node[below]{$2m$};
\draw[axis](0,0)--(.55,0);\draw[axis](2,0)--(1.7987,0);\node[left]at(-.7,0){高频};
\node at(1,-.6){$a_2/a_1=(1-\sqrt3)/2$};\end{scope}
\end{tikzpicture}
\caption{不等质量系统中，两模式的位移方向和相对振幅}\label{fig:hw4-5-s}
\end{figure}

等质量时，两模式为$(1,1)$、$(1,-1)$，分别有$\boxed{\omega_s^2=k/m,\ \omega_a^2=3k/m}$。

不等质量时置$\lambda=m\omega^2/k$，行列式化为$2\lambda^2-6\lambda+3=0$，故
\[
 \boxed{\omega_\pm^2=\frac{3\pm\sqrt3}{2}\frac km,\qquad
 \frac{a_2}{a_1}=2-\lambda_\pm=\frac{1\mp\sqrt3}{2}}.
\]
取$\vct a_\pm=(1,(1\mp\sqrt3)/2)$。它们的普通内积不为零，但
\[
 \vct a_-^{\mathsf T}\begin{pmatrix}m&0\\0&2m\end{pmatrix}\vct a_+
 =m\left[1+2\frac{1+\sqrt3}{2}\frac{1-\sqrt3}{2}\right]=0.
\]
因此展开初态时应采用质量加权投影，才能得到独立模式的振幅。

低频模式中两球同向，高频模式中反向。此时简单的$x_1\pm x_2$不再自动解耦；可令$\vct x=c_-\vct a_-\cos(\omega_-t+\delta_-)+c_+\vct a_+\cos(\omega_+t+\delta_+)$，再用初态求系数。


设$\vct M=\operatorname{diag}(m,2m)$，给定初态为$\vct x_0,\vct v_0$。用质量内积将两式分别投影，得到
\[
 q_s(0)=\frac{\vct a_s^{\mathsf T}\vct M\vct x_0}{\vct a_s^{\mathsf T}\vct M\vct a_s},\qquad
 \dot q_s(0)=\frac{\vct a_s^{\mathsf T}\vct M\vct v_0}{\vct a_s^{\mathsf T}\vct M\vct a_s},
 \quad s=\pm.
\]
于是$\vct x(t)=\sum_s\vct a_s[q_s(0)\cos\omega_st+\dot q_s(0)\sin\omega_st/\omega_s]$。例如只把第一球拉开$d$再静止释放，$q_s(0)=d/[1+2r_s^2]$，其中$r_s=(1\mp\sqrt3)/2$。两个系数之和为$d$，而$r_-q_-(0)+r_+q_+(0)=0$，恰好还原两球的初位移。

\end{solution}



\begin{example}[弹簧小车上的单摆]\label{sup4:problem-example-2}\label{sup4:ex:d7-e2}
质量为 $M$ 的小车在光滑水平面上运动，通过劲度系数为 $k$ 的弹簧连接固定墙。车上悬挂长为 $\ell$、摆球质量为 $m$ 的单摆。车的平衡位移记为 $x$，摆角 $\theta$ 从竖直向下方向向右量起。求两个小振动频率和振型，讨论 $k\to\infty$ 的极限。
\end{example}
\begin{figure}[H]\centering
\begin{tikzpicture}
\draw[curve](0,0)--(6,0);\draw[curve](0,0)--(0,1.5);\draw[spring](0,.6)--(2,.6)node[midway,above]{$k$};\draw[curve](2,.2)rectangle(3.8,1);\node at(2.9,.55){$M$};\draw[curve](2.9,1)--(2.9,2.8);\fill(2.9,2.8)circle(1.5pt)node[above]{$P$};\draw[curve](2.9,2.8)--(3.7,1.3)node[midway,right]{$\ell$};\fill(3.7,1.3)circle(2.7pt)node[right]{$m$};\draw[guide](2.9,2.8)--(2.9,1.3);\draw(2.9,2.15)arc(-90:-62:.65);\node at(3.15,1.9){$\theta$};\draw[axis](2.3,-.45)--(3.4,-.45)node[right]{$x$};
\end{tikzpicture}
\caption{P随小车平移。摆球的地面水平位移一阶为 $x+\ell\theta$。}\label{sup4:fig:d7-e2-q}
\end{figure}
\begin{solution}
图\ref{sup4:fig:d7-e2-q}中，摆球地面坐标为 $(x+\ell\sin\theta,-\ell\cos\theta)$。因此二次动能和势能为
\[
T_2=\frac12(M+m)\dot x^2+m\ell\dot x\dot\theta
+\frac12m\ell^2\dot\theta^2,
\qquad U_2=\frac12kx^2+\frac12mg\ell\theta^2.
\]
摆球的精确速率平方为 $\dot x^2+2\ell\cos\theta\,\dot x\dot\theta+\ell^2\dot\theta^2$。其中交叉项说明：小车速度与摆动速度同向时，摆球动能增加，反向时则减小。保留二次项后，由动能、势能逐项列出
\[
\begin{aligned}
(M+m)\ddot x+m\ell\ddot\theta+kx&=0,\\
m\ell\ddot x+m\ell^2\ddot\theta+mg\ell\theta&=0.
\end{aligned}
\]
第一式是整个车摆系统的水平动量方程，第二式描述加速悬点下的摆动。对一支模态取 $x=X\cos\omega t$、$\theta=\Theta\cos\omega t$，将二阶导数换成 $-\omega^2$ 倍的位移，即得下面的矩阵方程。
对应矩阵为
\[
M_q=\begin{pmatrix}M+m&m\ell\\m\ell&m\ell^2\end{pmatrix},
\qquad K_q=\begin{pmatrix}k&0\\0&mg\ell\end{pmatrix}.
\]
记 $\lambda=\omega^2$，特征行列式除以 $m\ell^2$ 后得到
\[
M\lambda^2-\left[k+\frac{(M+m)g}{\ell}\right]\lambda
+\frac{kg}{\ell}=0.
\]
行列式展开的具体形式为
\[
[k-(M+m)\lambda](g-\ell\lambda)-m\ell\lambda^2=0.
\]
其判别式可写成 $[k-(M+m)g/\ell]^2+4mkg/\ell>0$，两根之和、之积又分别为正，故确有两个不同的正频率平方。
故两频率平方为
\[
\boxed{\omega_{\pm}^2=
\frac{k+(M+m)g/\ell\ \pm\sqrt{[k+(M+m)g/\ell]^2-4Mkg/\ell}}{2M}.}
\]
将 $\lambda=g/\ell$ 代入上面的二次多项式，值为 $-m(g/\ell)^2<0$；而两个根均为正，所以 $\omega_-^2<g/\ell<\omega_+^2$。第二行模态方程给出
\[
\frac{\Theta}{X}=\frac{\omega^2}{\ell(g/\ell-\omega^2)}.
\]
\begin{figure}[H]\centering
\begin{tikzpicture}
\begin{scope}
\draw[curve](0,0)rectangle(1.7,.65);\draw[curve](.15,.65)--(.15,2.85)--(.85,2.85);\fill(.85,2.85)circle(1.5pt)node[above right]{$P$};
\draw[guide](.85,2.85)--(.85,1.45);\draw[curve](.85,2.85)--(1.3,1.6);\fill(1.3,1.6)circle(2pt);
\draw[axis](.2,3.45)--(1.3,3.45)node[right]{$X$};\draw[axis](.85,1.3)--(1.3,1.3)node[midway,below]{$\ell\Theta$};\node at(.9,-.5){低频：$X\Theta>0$};
\end{scope}
\begin{scope}[xshift=4.7cm]
\draw[curve](0,0)rectangle(1.7,.65);\draw[curve](.15,.65)--(.15,2.85)--(.85,2.85);\fill(.85,2.85)circle(1.5pt)node[above right]{$P$};
\draw[guide](.85,2.85)--(.85,1.45);\draw[curve](.85,2.85)--(.4,1.6);\fill(.4,1.6)circle(2pt);
\draw[axis](.2,3.45)--(1.3,3.45)node[right]{$X$};\draw[axis](.85,1.3)--(.4,1.3)node[midway,below]{$\ell\Theta$};\node at(.9,-.5){高频：$X\Theta<0$};
\end{scope}
\end{tikzpicture}
\caption{两种模态的位移方向。$X$ 为小车位移，$\ell\Theta$ 为摆球相对悬点P的水平位移；箭头长度仅作示意。}\label{sup4:fig:d7-e2-s}
\end{figure}
图\ref{sup4:fig:d7-e2-s}中，低频模态的小车位移与摆角同号，高频模态中两者反号。摆球在地面系中的水平位移为 $x+\ell\theta$，其中 $\ell\theta$ 只表示相对悬点的位移。

求$k\to\infty$的极限时，不直接计算低频根中两个大数的差。将高频根的根式除以$k$，可得
\[
\frac{\omega_+^2}{k}
=\frac{1+(M+m)g/(k\ell)+
\sqrt{[1+(M+m)g/(k\ell)]^2-4Mg/(k\ell)}}{2M}
\longrightarrow\frac1M.
\]
再由两根之积$\omega_-^2\omega_+^2=kg/(M\ell)$，得到
\[
\omega_-^2=\frac{kg}{M\ell\omega_+^2}\longrightarrow\frac g\ell.
\]
对低频模态保持有限的摆角振幅，振幅比公式给出$X\to0$，恢复固定支点单摆。高频模态则有$\ell\Theta/X=\omega_+^2/(g/\ell-\omega_+^2)\to-1$，摆球的地面水平位移近乎抵消。弹簧在这一快速模态中主要驱动小车自身，所以$\omega_+^2\sim k/M$。
\end{solution}



\begin{example}[等长等质量双摆]\label{sup4:problem-example-6}\label{sup4:ex:d7-e6}
双摆两杆均轻且长为 $\ell$，两质点质量均为 $m$，在同一竖直平面运动。两摆角 $\theta_1,\theta_2$ 均从竖直向下方向量起。求小振动频率及振型，并求初始 $\theta_1=\Theta,\theta_2=0$、两角速度均为零时的运动。
\end{example}
\begin{figure}[H]\centering
\begin{tikzpicture}
\fill(0,0)circle(1.6pt)node[above]{$O$};\coordinate(A)at(.8,-1.8);\coordinate(B)at(.3,-3.7);\draw[curve](0,0)--(A)--(B);\fill(A)circle(2.5pt)node[right]{$m$};\fill(B)circle(2.5pt)node[left]{$m$};\draw[guide](0,0)--(0,-1.5);\draw[guide](A)--(.8,-3.3);\draw(0,-.7)arc(-90:-66:.7);\node at(.22,-1.25){$\theta_1$};\draw(.8,-2.5)arc(-90:-105:.7);\node at(.66,-3.1){$\theta_2$};\node[right]at(.5,-1){$\ell$};\node[left]at(.4,-2.7){$\ell$};\node[align=left]at(4,-1.8){两角均为绝对摆角\\第二角不从第一杆量起};
\end{tikzpicture}
\caption{第二杆图示向左偏，故 $\theta_2<0$。两个正角均取向右偏转。}\label{sup4:fig:d7-e6-q}
\end{figure}
\begin{solution}
按图\ref{sup4:fig:d7-e6-q}，两质点水平位移的一阶式分别为 $\ell\theta_1$、$\ell(\theta_1+\theta_2)$。竖直速度只从二阶开始，因此
\[
T_2=\frac12m\ell^2(2\dot\theta_1^2+2\dot\theta_1\dot\theta_2+\dot\theta_2^2),
\qquad U_2=\frac12mg\ell(2\theta_1^2+\theta_2^2).
\]
竖直位移对势能仍须保留二次项。上质点上升 $\ell\theta_1^2/2$，下质点上升 $\ell(\theta_1^2+\theta_2^2)/2$，相加便得到 $U_2$。由二次能量列出的两个角坐标方程为
\[
2\ddot\theta_1+\ddot\theta_2+\frac{2g}{\ell}\theta_1=0,\qquad
\ddot\theta_1+\ddot\theta_2+\frac g\ell\theta_2=0.
\]
图\ref{sup4:fig:d7-e6-q}中的第二角是第二杆相对竖直线的绝对角；若改成两杆夹角，坐标关系和矩阵都会改变。
两个矩阵为
\[
M=m\ell^2\begin{pmatrix}2&1\\1&1\end{pmatrix},\qquad
K=mg\ell\begin{pmatrix}2&0\\0&1\end{pmatrix}.
\]
令 $\eta=\omega^2\ell/g$，特征方程为
\[
\det\begin{pmatrix}2-2\eta&-\eta\\-\eta&1-\eta\end{pmatrix}
=\eta^2-4\eta+2=0.
\]
故
\[
\boxed{\omega_-^2=(2-\sqrt2)g/\ell,\qquad
\omega_+^2=(2+\sqrt2)g/\ell.}
\]
由第二行 $-\eta\Theta_1+(1-\eta)\Theta_2=0$，得到低频振型 $(1,\sqrt2)$、高频振型 $(1,-\sqrt2)$。
\begin{figure}[H]\centering
\begin{tikzpicture}
\begin{scope}
\fill(0,0)circle(1.5pt);\draw[curve](0,0)--(.45,-1.5)--(1.07,-2.95);\fill(.45,-1.5)circle(2pt);\fill(1.07,-2.95)circle(2pt);\node at(.5,.55){低频};\node at(.5,-3.5){$\Theta_2=\sqrt2\Theta_1$};
\end{scope}
\begin{scope}[xshift=4.8cm]
\fill(0,0)circle(1.5pt);\draw[curve](0,0)--(.45,-1.5)--(-.17,-2.95);\fill(.45,-1.5)circle(2pt);\fill(-.17,-2.95)circle(2pt);\node at(.3,.55){高频};\node at(.3,-3.5){$\Theta_2=-\sqrt2\Theta_1$};
\end{scope}
\end{tikzpicture}
\caption{两模态中，两杆绝对摆角同号或反号。构形只示小角比例。}\label{sup4:fig:d7-e6-s}
\end{figure}
图\ref{sup4:fig:d7-e6-s}的比例针对两杆摆角。第二质点的水平位移为 $\ell(\theta_1+\theta_2)$，故两质点带符号的水平振幅比为 $(\Theta_1+\Theta_2)/\Theta_1$。

若两个余弦模态的系数为 $c_-,c_+$，初态分别要求 $c_-+c_+=\Theta$、$\sqrt2c_--\sqrt2c_+=0$，所以两者均为 $\Theta/2$。两初速度为零，又使两个正弦系数同时为零。
初始构形分解为 $(\Theta,0)=\frac\Theta2(1,\sqrt2)+\frac\Theta2(1,-\sqrt2)$，因此
\[
\boxed{\theta_1=\frac\Theta2(\cos\omega_-t+\cos\omega_+t),\qquad
\theta_2=\frac\Theta{\sqrt2}(\cos\omega_-t-\cos\omega_+t).}
\]
本题势能矩阵已对角，耦合来自动能交叉项。第一杆偏转还同时抬高两个质点，这使 $\theta_1^2$ 的势能系数为单个摆球的两倍。
\end{solution}



\begin{example}[三振子链的节点与冲量选择]\label{deep4:problem-example-9}\label{deep4:ex:d7-e9}
三个相同质量 $m$ 的滑块依次由四根相同弹簧 $k$ 与两端固定墙连接，在光滑直线上运动。求全部小振动模态。系统初始位于平衡，对中间块施加向右冲量 $J$，求此后的运动及各模态能量。
\end{example}
\begin{figure}[H]\centering
\begin{tikzpicture}
\draw[curve](0,0)--(9,0);\draw[curve](0,0)--(0,1.2);\draw[curve](9,0)--(9,1.2);\foreach\x in{1.6,4.1,6.6}{\draw[curve](\x,.2)rectangle++(.8,.7);\node at(\x+.4,.55){$m$};}\draw[spring](0,.55)--(1.6,.55);\draw[spring](2.4,.55)--(4.1,.55);\draw[spring](4.9,.55)--(6.6,.55);\draw[spring](7.4,.55)--(9,.55);\foreach\x in{.8,3.25,5.75,8.2}{\node[above]at(\x,.7){$k$};}\draw[axis](4.5,1.2)--(5.8,1.2)node[right]{$J$};\foreach\x/\i in{2/1,4.5/2,7/3}{\node[below]at(\x,0){$x_\i$};}
\end{tikzpicture}
\caption{中间块受冲量；左右两块在初态关于中点对称。}\label{deep4:fig:d7-e9-q}
\end{figure}
\begin{solution}
图\ref{deep4:fig:d7-e9-q}的刚度矩阵为
\[
K=k\begin{pmatrix}2&-1&0\\-1&2&-1\\0&-1&2\end{pmatrix},\qquad M=mI_3.
\]
矩阵中的对角元2来自每个物块连接的两根弹簧，邻接元 $-1$ 来自中间弹簧形变量之差。若直接求特征行列式，令 $\eta=m\omega^2/k$，有
\[
\det(K/k-\eta I_3)=(2-\eta)[(2-\eta)^2-2].
\]
利用图\ref{deep4:fig:d7-e9-s}的镜面对称性，可以先分出反对称模态，再在对称子空间求另外两根，避免展开三次多项式。
反对称振型可先由镜面对称性找出：$(1,0,-1)$，中间块不动，频率平方为 $2k/m$。其余模态左右同向，置振型为 $(1,b,1)$，外块和中间块方程给出
\[
2-b=\eta,\qquad 2b-2=\eta b,\qquad \eta=m\omega^2/k.
\]
消元得 $b^2=2$。按频率升序，
\[
\begin{array}{c|c|c}
&\omega_j^2&\vct v_j\ \text{（普通单位长度）}\\\hline
1&(2-\sqrt2)k/m&(1,\sqrt2,1)/2\\
2&2k/m&(1,0,-1)/\sqrt2\\
3&(2+\sqrt2)k/m&(1,-\sqrt2,1)/2
\end{array}
\]
\begin{figure}[H]\centering
\begin{tikzpicture}
\foreach\y/\j in{2/1,0/2,-2/3}{\node[left]at(0,\y){模态\j};\foreach\x in{1,3.5,6}{\fill(\x,\y)circle(2pt);}}
\draw[axis](1,2)--(1.8,2);\draw[axis](3.5,2)--(4.631,2);\draw[axis](6,2)--(6.8,2);
\draw[axis](1,0)--(1.8,0);\node[above]at(3.5,0){节点};\draw[axis](6,0)--(5.2,0);
\draw[axis](1,-2)--(1.8,-2);\draw[axis](3.5,-2)--(2.369,-2);\draw[axis](6,-2)--(6.8,-2);
\end{tikzpicture}
\caption{中间块在模态2中为节点，打在此处的冲量不激起该模态。}\label{deep4:fig:d7-e9-s}
\end{figure}
冲量后的速度为 $(0,J/m,0)$。令 $\vct x=\sum\vct v_j a_j(t)$，因振型已按普通长度归一化且三质量相同，模态等效质量均为m。投影得到
\[
\dot a_1(0)=J/(\sqrt2m),\quad \dot a_2(0)=0,\quad
\dot a_3(0)=-J/(\sqrt2m),\qquad a_j(0)=0.
\]
冲量先确定模态初速度。对普通单位振型 $\vct v_j$，投影给出 $m\dot a_j(0)=\vct v_j^{\mathsf T}(0,J,0)^{\mathsf T}$；再解 $\ddot a_j+\omega_j^2a_j=0$，得到 $a_j(t)=\dot a_j(0)\sin\omega_jt/\omega_j$。位移振幅中的频率分母由这一步产生。
所以
\[
\boxed{\begin{aligned}
x_1=x_3&=\frac{J}{2\sqrt2m}\left(\frac{\sin\omega_1t}{\omega_1}-\frac{\sin\omega_3t}{\omega_3}\right),\\
x_2&=\frac{J}{2m}\left(\frac{\sin\omega_1t}{\omega_1}+\frac{\sin\omega_3t}{\omega_3}\right).
\end{aligned}}
\]
图\ref{deep4:fig:d7-e9-s}的模态2具有中间节点，初始位移、模态速度均为零，此后始终不出现。各模态能量为 $E_1=E_3=J^2/(4m)$、$E_2=0$，总和为中间块初始动能 $J^2/(2m)$。
\end{solution}



\section{拍振、失谐与外力响应}

近邻频率的叠加产生慢包络，周期外力则按其施加方式激发不同振型。自由拍振和强迫响应分别讨论，以区分能量交换、频率避交及反共振。



\subsection{拍振、周期外力与反共振}
两个相近频率叠加时，利用
\[
\cos\omega_1t+\cos\omega_2t
=2\cos\frac{\omega_2-\omega_1}{2}t\,
\cos\frac{\omega_2+\omega_1}{2}t,
\]
可把运动分为快速振动和缓慢包络。只有 $|\omega_2-\omega_1|\ll(\omega_1+\omega_2)/2$ 时，才有清楚分离的快、慢时间尺度。包络过零的时刻，位移为零，速度未必也为零；中间耦合弹簧还可能储存能量。

若外力为 $\vct f\cos\nu t$，周期特解 $\vct X\cos\nu t$ 满足
\begin{equation}
(K-\nu^2M)\vct X=\vct f.
\label{sup4:eq:d7-forced}
\end{equation}
当 $\nu$ 接近固有频率且外力对该模态的投影非零，会出现共振。若某个坐标的响应分子为零，而矩阵仍可逆，该坐标的周期响应可为零，称为反共振。无阻尼系统从任意初态出发还包含不衰减的自由振动；上述周期特解须与初始条件一起使用。
\begin{figure}[H]\centering
\begin{tikzpicture}
\draw[axis](0,0)--(6.1,0)node[right]{$\nu^2$};\draw[axis](0,-1.6)--(0,2)node[above]{有符号响应};
\draw[guide](1.5,-1.5)--(1.5,1.8);\draw[guide](4.5,-1.5)--(4.5,1.8);\node[below]at(1.5,-1.5){$\omega_1^2$};\node[below]at(4.5,-1.5){$\omega_2^2$};
\draw[curve](2,-1.3)..controls(2.6,-.5)and(3.3,.5)..(4,1.3);\fill(3,0)circle(1.8pt)node[below right]{反共振};\node at(3,2){两共振频率之间的一种情形};
\end{tikzpicture}
\caption{反共振是某个观测坐标的周期响应过零，其他坐标仍可运动。曲线仅示意零点与极点的次序。}\label{sup4:fig:d7-antiresonance}
\end{figure}
图\ref{sup4:fig:d7-antiresonance}中的极点还须结合外力和观测位置判断。若外力恰与某模态正交，就不能激起该模态；若测点位于该振型的节点，该频率也可能不出现在测点响应中。


\subsection{例题与模型变式}

\begin{example}[对称双振子与拍振包络]\label{sup4:problem-example-3}\label{sup4:ex:d7-e3}
两个质量均为 $m$ 的滑块在光滑直线上运动，左右外侧分别由弹簧 $k$ 连到固定墙，中间由弹簧 $\kappa$ 相连，三弹簧在平衡态均无形变。初始 $x_1=A,x_2=0$，速度均为零。求运动与模态能量，讨论弱耦合时的拍振，并判断包络过零是否代表第一块已经静止。
\end{example}
\begin{figure}[H]\centering
\begin{tikzpicture}
\draw[curve](0,0)--(7.5,0);\draw[curve](0,0)--(0,1.2);\draw[curve](7.5,0)--(7.5,1.2);\draw[spring](0,.55)--(2.3,.55);\draw[curve](2.3,.2)rectangle(3,.9);\draw[spring](3,.55)--(4.5,.55);\draw[curve](4.5,.2)rectangle(5.2,.9);\draw[spring](5.2,.55)--(7.5,.55);\node at(2.65,.55){$m$};\node at(4.85,.55){$m$};\node[above]at(1.1,.65){$k$};\node[above]at(3.75,.65){$\kappa$};\node[above]at(6.3,.65){$k$};\node at(2.65,-.55){$x_1(0)=A$};\node at(5.2,-.55){$x_2(0)=0$};
\end{tikzpicture}
\caption{两个位移均向右为正，图示初始位移条件。}\label{sup4:fig:d7-e3-q}
\end{figure}
\begin{solution}
图\ref{sup4:fig:d7-e3-q}中间弹簧的形变量为 $x_2-x_1$。它对第一块的力为 $\kappa(x_2-x_1)$，对第二块为相反数；外侧弹簧各提供 $-kx_i$。所以
\[
\begin{aligned}
m\ddot x_1&=-kx_1+\kappa(x_2-x_1),\\
m\ddot x_2&=-kx_2-\kappa(x_2-x_1).
\end{aligned}
\]
两式相加时耦合力抵消，相减时耦合力加倍，这就是选取同相、反相坐标的依据。
图\ref{sup4:fig:d7-e3-q}的同相与反相坐标取为 $s=(x_1+x_2)/2$、$d=(x_1-x_2)/2$。由两物块方程相加、相减，
\[
\ddot s+\omega_s^2s=0,\quad \omega_s^2=k/m;
\qquad \ddot d+\omega_a^2d=0,\quad \omega_a^2=(k+2\kappa)/m.
\]
两个方程的通解为$s=C_s\cos\omega_st+D_s\sin\omega_st$、$d=C_a\cos\omega_at+D_a\sin\omega_at$。初始$s(0)=d(0)=A/2$，给出$C_s=C_a=A/2$；求导后代入零初速度，得到$\omega_sD_s=\omega_aD_a=0$。因此$D_s=D_a=0$，还原$x_1=s+d$、$x_2=s-d$，得到
\[
\boxed{x_1=\frac A2(\cos\omega_st+\cos\omega_at),\quad
x_2=\frac A2(\cos\omega_st-\cos\omega_at).}
\]
用 $x_1=s+d$、$x_2=s-d$ 代回能量，动能中的交叉项相消，三根弹簧的势能化为
\[
T=m(\dot s^2+\dot d^2),\qquad U=ks^2+(k+2\kappa)d^2.
\]
因此s、d各自对应的等效质量为 $2m$。图\ref{sup4:fig:d7-e3-s}中的两条包络来自这两个独立简谐运动的叠加。
总能量写成 $E=m(\dot s^2+\omega_s^2s^2)+m(\dot d^2+\omega_a^2d^2)$，故
\[
E_s=\frac14kA^2,\qquad E_a=\frac14(k+2\kappa)A^2,
\quad E=\frac12(k+\kappa)A^2.
\]
两模态的初始位移系数均为 $A/2$，能量仍因频率不同而有差别。
\begin{figure}[H]\centering
\begin{tikzpicture}
\draw[axis](0,0)--(6.6,0)node[right]{$\tau=\Delta\omega t/2$};\draw[axis](0,-1.6)--(0,1.8)node[above]{包络系数};\draw[curve,domain=0:180,samples=90]plot({\x/30},{1.3*cos(\x)});\draw[guide,domain=0:180,samples=90]plot({\x/30},{1.3*sin(\x)});\node[below]at(3,0){$\pi/2$};\node[below]at(6,0){$\pi$};\node[left]at(0,1.3){$1$};\node[left]at(0,-1.3){$-1$};\node at(5.2,1.4){$x_2$：虚线};\node at(5,-1.8){$x_1$：实线};
\end{tikzpicture}
\caption{有符号慢包络；各自的振幅包络还须取绝对值。}\label{sup4:fig:d7-e3-s}
\end{figure}
令 $\bar\omega=(\omega_a+\omega_s)/2$、$\Delta\omega=\omega_a-\omega_s$，则
\[
x_1=A\cos\bar\omega t\cos\frac{\Delta\omega t}{2},\qquad
x_2=A\sin\bar\omega t\sin\frac{\Delta\omega t}{2}.
\]
图\ref{sup4:fig:d7-e3-s}给出两个慢因子。若$\kappa\ll k$，将$\omega_a=\omega_s\sqrt{1+2\kappa/k}$展开到一阶，得$\omega_a\simeq\omega_s(1+\kappa/k)$，所以$\Delta\omega\simeq\kappa/(m\omega_s)$。首次包络交接的时间为$t_b=\pi/\Delta\omega$。对$x_1$按乘积法则求导，
\[
\dot x_1=-A\bar\omega\sin\bar\omega t\cos\frac{\Delta\omega t}{2}
-\frac{A\Delta\omega}{2}\cos\bar\omega t\sin\frac{\Delta\omega t}{2}.
\]
代入$\Delta\omega t_b/2=\pi/2$，得到
\[
x_1(t_b)=0,\qquad
\dot x_1(t_b)=-\frac{A\Delta\omega}{2}\cos(\bar\omega t_b),
\]
第一块一般仍有非零速度。即使某次恰有 $x_1=\dot x_1=0$，中间弹簧也可能储能，见习题\ref{sup4:ex:d7-h4}。弱耦合时可近似说能量主要转移到另一物块，精确分配则须将三根弹簧全部计入。
\end{solution}



\begin{example}[失谐双摆与频率避交]\label{sup4:problem-example-5}\label{sup4:ex:d7-e5}
两单摆摆球质量均为 $m$，摆长为 $\ell_1,\ell_2$，支点高度使两球在静止时等高。两球由水平轻弹簧 $\kappa$ 连接，平衡时无形变。令球的水平小位移为 $x_1,x_2$。求小振动频率和振型；保持 $\omega_c^2=(g/\ell_1+g/\ell_2)/2$ 不变，改变失谐量 $d=(g/\ell_1-g/\ell_2)/2$，讨论两频率是否相交。
\end{example}
\begin{figure}[H]\centering
\begin{tikzpicture}
\fill(0,3)circle(1.6pt);\fill(4.5,2.3)circle(1.6pt);\draw[curve](0,3)--(0,0)node[midway,left]{$\ell_1$};\draw[curve](4.5,2.3)--(4.5,0)node[midway,right]{$\ell_2$};\fill(0,0)circle(3pt)node[below]{$m$};\fill(4.5,0)circle(3pt)node[below]{$m$};\draw[spring](.15,0)--(4.35,0)node[midway,above]{$\kappa$};\draw[axis](-.3,-.7)--(.8,-.7)node[right]{$x_1$};\draw[axis](4.2,-.7)--(5.3,-.7)node[right]{$x_2$};
\end{tikzpicture}
\caption{平衡时两摆球等高，支点高度不同。弹簧的一阶伸长为 $x_2-x_1$。}\label{sup4:fig:d7-e5-q}
\end{figure}
\begin{solution}
图\ref{sup4:fig:d7-e5-q}中，角度近似为 $x_i/\ell_i$。二次动能为 $m(\dot x_1^2+\dot x_2^2)/2$，势能为
\[
U_2=\frac{mg}{2\ell_1}x_1^2+\frac{mg}{2\ell_2}x_2^2
+\frac\kappa2(x_2-x_1)^2.
\]
势能的两个偏导分别为$\partial U_2/\partial x_1=mgx_1/\ell_1+\kappa(x_1-x_2)$、$\partial U_2/\partial x_2=mgx_2/\ell_2+\kappa(x_2-x_1)$。各令其与$m\ddot x_i$之和为零，记$c=\kappa/m>0$，得
\[
\ddot x_1+(\omega_c^2+d+c)x_1-cx_2=0,\qquad
\ddot x_2+(\omega_c^2-d+c)x_2-cx_1=0.
\]
代入$x_i=X_i\cos\omega t$，频率平方是矩阵
$\left(\begin{smallmatrix}\omega_c^2+d+c&-c\\-c&\omega_c^2-d+c\end{smallmatrix}\right)$ 的特征值，故
\[
\boxed{\omega_\pm^2=\omega_c^2+c\pm\sqrt{d^2+c^2}.}
\]
令 $\lambda=\omega^2$，特征行列式可先写成
\[
[(\omega_c^2+c)-\lambda]^2-d^2-c^2=0.
\]
这显示频率分支关于平均值 $\omega_c^2+c$ 对称。其两根之积为 $(g/\ell_1+c)(g/\ell_2+c)-c^2>0$，两根之和也为正，所以允许参数下两根均为正。
图\ref{sup4:fig:d7-e5-s}的纵轴是频率平方；实际频率差为 $2\sqrt{d^2+c^2}/(\omega_++\omega_-)$，不能把图上的间隔 $2c$ 当成频率差。
两摆长均为正还限制了 $|d|<\omega_c^2$。在允许范围内，两频率平方之差为 $2\sqrt{d^2+c^2}\ge2c$，非零耦合使两个分支保持分离。
\begin{figure}[H]\centering
\begin{tikzpicture}
\draw[axis](-3,0)--(3.2,0)node[right]{$d/c$};\draw[axis](0,0)--(0,4.7)node[above]{$\omega^2/c$};\draw[curve,domain=-2:2,samples=100]plot({1.3*\x},{.6*(5+sqrt((\x)^2+1))});\draw[curve,domain=-2:2,samples=100]plot({1.3*\x},{.6*(5-sqrt((\x)^2+1))});\draw[<->](.3,2.4)--(.3,3.6)node[midway,right]{$2$};\node[left]at(0,2.4){$4$};\node[left]at(0,3.6){$6$};\node at(-2,4.8){$\omega_c^2=4c$};
\end{tikzpicture}
\caption{有限耦合使频率平方分支避交；纵轴标出的间隔对应实际频率平方差 $2c$。}\label{sup4:fig:d7-e5-s}
\end{figure}
对低频模态，第一行方程给出
\[
\frac{X_2}{X_1}=\frac{d+\sqrt{d^2+c^2}}{c}>0,
\]
高频则为 $(d-\sqrt{d^2+c^2})/c<0$。在 $d=0$ 处恢复同相、反相振型。若 $d\gg c$ 且仍满足摆长允许范围，低频主要由第二摆运动；若 $d\ll-c$，低频主要由第一摆运动。图\ref{sup4:fig:d7-e5-s}的分支连续，跨过失谐零点后，低频模态中运动较强的摆球由第一摆换成第二摆。
\end{solution}



\begin{example}[双振子的反共振]\label{sup4:problem-example-8}\label{sup4:ex:d7-e8}
沿用例\ref{sup4:ex:d7-e3}的双振子，对第一块施加周期力 $F_0\cos\nu t$，第二块无外加驱动力。求无阻尼系统的周期特解，确定第一块发生反共振的频率及此时第二块的运动，讨论反共振与共振的差别。
\end{example}
\begin{figure}[H]\centering
\begin{tikzpicture}
\draw[curve](0,0)--(7.5,0);\draw[curve](0,0)--(0,1.2);\draw[curve](7.5,0)--(7.5,1.2);\draw[spring](0,.55)--(2,.55);\draw[curve](2,.2)rectangle(2.8,.9);\draw[spring](2.8,.55)--(4.6,.55);\draw[curve](4.6,.2)rectangle(5.4,.9);\draw[spring](5.4,.55)--(7.5,.55);\node at(2.4,.55){$m$};\node at(5,.55){$m$};\node[above]at(1,.7){$k$};\node[above]at(3.7,.7){$\kappa$};\node[above]at(6.4,.7){$k$};\draw[axis](2.4,1.3)--(4,1.3)node[right]{$F_0\cos\nu t$};
\end{tikzpicture}
\caption{驱动力只作用在第一块，两个位移正向均向右。}\label{sup4:fig:d7-e8-q}
\end{figure}
\begin{solution}
先在地面系列出受迫方程：
\[
\begin{aligned}
m\ddot x_1+(k+\kappa)x_1-\kappa x_2&=F_0\cos\nu t,\\
m\ddot x_2+(k+\kappa)x_2-\kappa x_1&=0.
\end{aligned}
\]
图\ref{sup4:fig:d7-e8-q}中外力只作用于第一块，却能通过中间弹簧驱动第二块。求同频周期特解时，两位移可取与外力同一余弦因子，但振幅系数允许为负，以表示反相。
图\ref{sup4:fig:d7-e8-q}中取 $x_i=X_i\cos\nu t$。令 $b=k+\kappa-m\nu^2$，幅值方程为
\[
\begin{pmatrix}b&-\kappa\\-\kappa&b\end{pmatrix}
\begin{pmatrix}X_1\\X_2\end{pmatrix}=
\begin{pmatrix}F_0\\0\end{pmatrix}.
\]
消元时不除以$b$，以免漏掉反共振点。将第一行乘$b$、第二行乘$\kappa$再相加，得到$(b^2-\kappa^2)X_1=bF_0$；将第一行乘$\kappa$、第二行乘$b$再相加，得到$(b^2-\kappa^2)X_2=\kappa F_0$。非共振时行列式$b^2-\kappa^2\ne0$，所以
\[
\boxed{X_1=\frac{F_0b}{b^2-\kappa^2},\qquad
X_2=\frac{F_0\kappa}{b^2-\kappa^2}.}
\]
分母在$\nu^2=k/m$、$(k+2\kappa)/m$为零，正是两固有频率。为求此时的响应，取$s=(x_1+x_2)/2$、$d=(x_1-x_2)/2$，相加、相减得到
\[
\ddot s+\omega_s^2s=\frac{F_0}{2m}\cos\nu t,\qquad
\ddot d+\omega_a^2d=\frac{F_0}{2m}\cos\nu t.
\]
若$\nu$等于其中一支固有频率，设该支的特解为$q_p=Ct\sin\nu t$。求导有
$\dot q_p=C\sin\nu t+C\nu t\cos\nu t$、
$\ddot q_p=2C\nu\cos\nu t-C\nu^2t\sin\nu t$，故
\[
\ddot q_p+\nu^2q_p=2C\nu\cos\nu t
=\frac{F_0}{2m}\cos\nu t,
\qquad q_p=\frac{F_0}{4m\nu}t\sin\nu t.
\]
它的振幅随时间增长，不能用有界周期特解代替；另一支模态仍按非共振方程求解。
\begin{figure}[H]\centering
\begin{tikzpicture}
\draw[axis](0,0)--(6.2,0)node[right]{$q=m\nu^2/k$};\draw[axis](0,-2.5)--(0,2.6)node[above]{$kX_1/F_0$};\draw[guide](1.4,-2.2)--(1.4,2.2);\draw[guide](4.2,-2.2)--(4.2,2.2);
\foreach\lo/\hi in{0/.85,1.15/2.85,3.15/4}{\draw[curve,domain=\lo:\hi,samples=100]plot({1.4*\x},{.55*(2-\x)/((1-\x)*(3-\x))});}
\fill(2.8,0)circle(1.7pt)node[below right,yshift=-5pt]{反共振};\node[below]at(1.4,0){$1$};\node[above left]at(2.8,0){$2$};\node[below]at(4.2,0){$3$};\node at(5.1,2.3){$\kappa=k$};\node[left]at(0,.55){$1$};
\end{tikzpicture}
\caption{第一块有符号周期响应。两个极点之间，响应在 $q=2$ 处过零。}\label{sup4:fig:d7-e8-s}
\end{figure}
图\ref{sup4:fig:d7-e8-s}的零点对应 $b=0$，此时
\[
\boxed{\nu_{\rm ar}^2=(k+\kappa)/m,\qquad X_1=0,\quad X_2=-F_0/\kappa.}
\]
此处 $b=0$ 而行列式为 $-\kappa^2\ne0$，因此反共振是响应分子的零点，不是固有频率处的奇点。由第二行方程 $-\kappa X_1=0$ 得 $X_1=0$，再由第一行 $-\kappa X_2=F_0$ 求出第二块振幅，既不需要也不允许令两块同时静止。
第一块保持在平衡位置，第二块的弹簧作用力 $\kappa x_2=-F_0\cos\nu t$ 恰好抵消驱动力。第二块则以第一块固定时的固有频率振动，故能维持这种周期运动。

这组周期解对应的初态为 $x_1=0,x_2=-F_0/\kappa$、两速度为零。从两块均静止于平衡位置开始驱动，还须叠加自由振动，详见习题\ref{sup4:ex:d7-h9}。无阻尼、非共振周期解的平均输入功率为零；加入阻尼后，响应幅值与耗散功率随之改变。
\end{solution}



\section{零频模态与线性近似的边界}

二次势能为零的方向，可能对应整体平移，也可能仍有高阶回复力。零频率的来源决定线性化是否足够，以及是否需要保留非线性项。



\subsection{零频率与小振动近似的范围}
$K$ 正定时，所有 $\omega_r^2>0$，线性化后的小振动有界。若 $\omega_r^2=-\sigma^2<0$，特征方程为 $\lambda^2-\sigma^2=0$，通解为 $Q_r=A\ee^{\sigma t}+B\ee^{-\sigma t}$；存在增长分支，原平衡不稳定。零特征值须直接解 $\ddot Q_0=0$：先积分得 $\dot Q_0=\dot Q_0(0)$，再积分得 $Q_0=Q_0(0)+\dot Q_0(0)t$。非零初速度产生漂移，不能将零频率直接代入含 $1/\omega$ 的简谐运动公式。

展开动能和势能时，须采用同一小量阶数。例如 $\rho\theta^2$、$\rho\dot\theta^2$ 通常为三阶能量项，在线性方程中可舍去，在更高阶运动中却会产生耦合。振动解的适用范围还受接触、纯滚动及弹簧线性等条件限制，各题中另行讨论。


\subsection{零频模态与四次回复势}
当势能沿某方向保持不变，零频模态给出匀速漂移。若二次项恰好消失，而最低非零项为$U-U_* =\beta q^4/4$，$\beta>0$，则运动方程变为
\[
 M_*\ddot q+\beta q^3=0.
\]
从振幅 $A>0$ 处静止释放，能量为 $\beta A^4/4$，因此
\[
\frac12M_*\dot q^2=\frac\beta4(A^4-q^4),\qquad
|\dot q|=\sqrt{\frac\beta{2M_*}}\sqrt{A^4-q^4}.
\]
利用四分之一周期的对称性，并令 $q=As$、$\dd q=A\dd s$，上下限 $q=0,A$ 对应 $s=0,1$，得到
\[
\frac T4=\sqrt{\frac{2M_*}\beta}\int_0^A\frac{\dd q}{\sqrt{A^4-q^4}}
=\frac1A\sqrt{\frac{2M_*}\beta}\int_0^1\frac{\dd s}{\sqrt{1-s^4}}.
\]
最后的积分是与振幅无关的有限常数，不能用初等函数给出原函数。于是
\[
 T(A)=\frac4A\sqrt{\frac{2M_*}{\beta}}
 \int_0^1\frac{\dd s}{\sqrt{1-s^4}}.
\]
回复运动的周期随振幅而变；$A\to0$时周期趋于无穷，故不能赋予它非零的线性固有频率。例\ref{ex:hw4-4}的临界木条和例\ref{ex:hw4-8}的两个端点电荷比，都须保留四次势能才能判断。若多坐标系统还有非退化方向，应同时考虑耦合高阶项，不能只沿一条坐标轴画势能曲线就结束稳定性分析。



\LectureNeedspace{4\baselineskip}


\subsection{例题与模型变式}

\begin{example}[自由双质点的零频模态]\label{sup4:problem-example-10}\label{sup4:ex:d7-e10}
光滑直线上，质量为 $m_1,m_2$ 的两个物体由劲度系数为 $\kappa$、原长为 $\ell_0$ 的弹簧连接，两端均不固定。初始弹簧无形变、两物体静止，对第一物体施加向右冲量 $J$。求两个物体相对初始位置的位移及能量分配，说明零频模态的含义。
\end{example}
\begin{figure}[H]\centering
\begin{tikzpicture}
\draw[curve](-.5,0)--(7.5,0);\draw[curve](0,.2)rectangle(1,.9);\draw[curve](5.5,.2)rectangle(6.9,.9);\draw[spring](1,.55)--(5.5,.55)node[midway,above]{$\kappa,\ \ell_0$};\node at(.5,.55){$m_1$};\node at(6.2,.55){$m_2$};\draw[axis](.5,1.3)--(2,1.3)node[right]{$J$};\draw[axis](2,-.6)--(3.3,-.6)node[right]{位移正向};
\end{tikzpicture}
\caption{两端没有固定墙，弹簧只改变相对运动，不限制整体平移。}\label{sup4:fig:d7-e10-q}
\end{figure}
\begin{solution}
令总质量 $M=m_1+m_2$、约化质量 $\mu=m_1m_2/M$。在图\ref{sup4:fig:d7-e10-q}中，以质心位置 $X$ 和弹簧伸长 $r=x_2-x_1-\ell_0$ 为坐标，有
\[
T=\frac12M\dot X^2+\frac12\mu\dot r^2,\qquad U=\frac12\kappa r^2.
\]
$x_1,x_2$ 表示地面位置，所求位移为各自相对初始位置的 $u_1,u_2$。反解坐标变换，
\[
x_1=X-\frac{m_2}{M}(\ell_0+r),\qquad
x_2=X+\frac{m_1}{M}(\ell_0+r).
\]
对时间求导再代入两质点动能，交叉项系数为 $-m_1m_2/M+m_2m_1/M=0$，径向项的系数为约化质量 $\mu$。这解释了质心平移与弹簧伸缩为何能够分离。
方程完全分离为 $\ddot X=0$、$\ddot r+\omega^2r=0$，其中
\[
\omega^2=\frac\kappa\mu=\kappa\left(\frac1{m_1}+\frac1{m_2}\right).
\]
质心方程积分两次，通解为$X=C_0+C_1t$；相对方程的特征根为$\pm\mathrm i\omega$，通解为$r=A\cos\omega t+B\sin\omega t$。冲量后$X(0)=X_0$、$\dot X(0)=J/M$，所以$C_0=X_0,C_1=J/M$；又$r(0)=0$给出$A=0$，$\dot r(0)=\omega B=-J/m_1$给出$B=-J/(m_1\omega)$。因此
\[
X-X_0=\frac JM t,\qquad r=-\frac{J}{m_1\omega}\sin\omega t.
\]
\begin{figure}[H]\centering
\begin{tikzpicture}
\begin{scope}
\draw[axis](0,0)--(3.8,0)node[right]{$t$};\draw[axis](0,0)--(0,2.2)node[above]{$X-X_0$};\draw[curve](0,0)--(3.2,1.7)node[right]{$Jt/M$};\node at(1.5,-.6){整体平移};
\end{scope}
\begin{scope}[xshift=5.4cm]
\draw[axis](0,0)--(3.8,0)node[right]{$\omega t$};\draw[axis](0,-1.4)--(0,1.4)node[above]{$r$};\draw[curve,domain=0:360,samples=90]plot({\x/110},{-sin(\x)});\node at(1.6,-1.9){内部振动};
\end{scope}
\end{tikzpicture}
\caption{质心匀速漂移，弹簧伸长保持有界。右图纵向以 $J/(m_1\omega)$ 为单位。}\label{sup4:fig:d7-e10-s}
\end{figure}
由质心分配关系，两物体相对初始位置的位移为
\[
\boxed{u_1=\frac JM t+\frac{m_2J}{Mm_1\omega}\sin\omega t,
\qquad u_2=\frac JM t-\frac{J}{M\omega}\sin\omega t.}
\]
图\ref{sup4:fig:d7-e10-s}中的零频分支对应同向平移振型 $(1,1)$，其时间因子为常量加线性项。扣除这一漂移后，内部振动的两位移满足质量加权和为零。

能量分为
\[
E_{\rm cm}=\frac{J^2}{2M},\qquad
E_{\rm int}=\frac12\mu\frac{J^2}{m_1^2}
=\frac{m_2J^2}{2m_1M},
\]
两部分相加为 $J^2/(2m_1)$。内部运动保持有界，两物体的地面位置则随质心漂移；讨论内部平衡的稳定性时，应先扣除整体平移。
\end{solution}



\begin{example}[弹簧摆的线性分离与内共振]\label{deep4:problem-example-7}\label{deep4:ex:d7-e7}
一轻弹簧劲度系数为 $k$、原长为 $\ell_0>0$，上端固定，下端连接质量为 $m$ 的小球，小球可在竖直平面运动。取弹簧长度 $r$、偏离竖直向下的角度 $\theta$ 为坐标。求平衡位置附近的径向、摆动频率，并求两者满足 $\omega_r=2\omega_\theta$ 时的参数条件。再令首阶摆动为 $\theta=\Theta\cos\omega_\theta t$，求它引起的最低阶径向驱动项。
\end{example}
\begin{figure}[H]\centering
\begin{tikzpicture}
\fill(0,0)circle(1.6pt)node[above]{$O$};\draw[guide](0,0)--(0,-3.6);\draw[spring](0,0)--(1,-3);\fill(1,-3)circle(3pt)node[right]{$m$};\node[right]at(.7,-1.6){$r$};\draw(0,-.85)arc(-90:-72:.85);\node at(.18,-1.6){$\theta$};\node[align=left]at(4,-1.4){$k,\ \ell_0$\\平衡长度为 $r_0$};
\end{tikzpicture}
\caption{径向伸缩与角向摆动使用不同坐标；平衡态为竖直悬垂。}\label{deep4:fig:d7-e7-q}
\end{figure}
\begin{solution}
图\ref{deep4:fig:d7-e7-q}中的精确能量为
\[
T=\frac12m(\dot r^2+r^2\dot\theta^2),\qquad
U=\frac12k(r-\ell_0)^2-mgr\cos\theta.
\]
先列出能量中的各项导数：
\[
\frac{\partial U}{\partial r}=k(r-\ell_0)-mg\cos\theta,
\quad\frac{\partial U}{\partial\theta}=mgr\sin\theta,
\quad\frac{\partial T}{\partial r}=mr\dot\theta^2,
\]
\[
\frac{\dd}{\dd t}\frac{\partial T}{\partial\dot r}=m\ddot r,
\qquad
\frac{\dd}{\dd t}\frac{\partial T}{\partial\dot\theta}
=\frac{\dd}{\dd t}(mr^2\dot\theta)
=mr^2\ddot\theta+2mr\dot r\dot\theta.
\]
平衡时令两个势能偏导为零，取向下悬垂分支，得$r_0=\ell_0+mg/k$、$\theta=0$。置 $\rho=r-r_0$，按 $\rho/r_0$ 与 $\theta$ 同为小量展开，二次项为
\[
T_2=\frac12m\dot\rho^2+\frac12mr_0^2\dot\theta^2,
\quad U_2=U_*+\frac12k\rho^2+\frac12mgr_0\theta^2.
\]
例如动能中的 $mr_0\rho\dot\theta^2$、势能中的 $mg\rho\theta^2/2$ 都是三阶量，在线性方程所需的二次能量中应舍去。径向一次势能项则由 $k(r_0-\ell_0)=mg$ 抵消。两个方向在线性阶分离，并不表示精确运动没有耦合。
对二次能量求导，径向方程为$m\ddot\rho+k\rho=0$，角向方程为$mr_0^2\ddot\theta+mgr_0\theta=0$。分别除以各自的惯性系数，两个线性方程独立，其频率为
\[
\boxed{\omega_r=\sqrt{k/m},\qquad \omega_\theta=\sqrt{g/r_0}.}
\]
条件 $\omega_r^2=4\omega_\theta^2$ 等价于 $kr_0=4mg$，利用平衡长度，得到
\[
\boxed{k\ell_0=3mg.}
\]
\begin{figure}[H]\centering
\begin{tikzpicture}
\begin{scope}
\draw[axis](-.5,0)--(4,0)node[right]{驱动频率};\draw[axis](-.5,-1.5)--(-.5,1.5)node[above]{二阶径向驱动};\draw[curve](0,0)--(0,.4);\node[below]at(0,0){$0$};\draw[curve](2.5,0)--(2.5,-1.2);\node[above]at(2.5,0){$2\omega_\theta$};\node[right]at(0,.45){$g\Theta^2/4$};\node[right]at(2.5,-1.2){$-3g\Theta^2/4$};
\end{scope}
\end{tikzpicture}
\caption{角向运动产生恒定项和倍频项；竖线表示傅里叶系数及符号。}\label{deep4:fig:d7-e7-s}
\end{figure}
精确径向方程为 $\ddot r-r\dot\theta^2=g\cos\theta-(k/m)(r-\ell_0)$。保留 $\Theta^2$ 阶，得
\[
\ddot\rho+\omega_r^2\rho
=r_0\dot\theta^2-\frac g2\theta^2
=\frac{g\Theta^2}{4}-\frac{3g\Theta^2}{4}\cos2\omega_\theta t.
\]
代入时，$r_0\omega_\theta^2=g$，而 $\sin^2\omega_\theta t=(1-\cos2\omega_\theta t)/2$、$\cos^2\omega_\theta t=(1+\cos2\omega_\theta t)/2$。离心项与重力二次项共同给出常量 $g\Theta^2/4$ 和倍频项 $-3g\Theta^2\cos2\omega_\theta t/4$，两项都应保留。常量只平移径向振动中心，倍频项才可能发生共振。
图\ref{deep4:fig:d7-e7-s}中，第二根谱线在$2\omega_\theta=\omega_r$时与径向固有频率相合。常量驱动的特解为$g\Theta^2/(4\omega_r^2)$。对余下的倍频驱动，设$\rho_{\rm sec}=Ct\sin\omega_rt$，逐次求导：
\[
\dot\rho_{\rm sec}=C\sin\omega_rt+C\omega_rt\cos\omega_rt,
\qquad
\ddot\rho_{\rm sec}=2C\omega_r\cos\omega_rt-C\omega_r^2t\sin\omega_rt.
\]
代入左端后含$t\sin\omega_rt$的两项抵消，故$2C\omega_r=-3g\Theta^2/4$，得到
\[
\rho_{\rm sec}=-\frac{3g\Theta^2}{8\omega_r}t\sin\omega_rt.
\]
在这一阶近似内，径向通解还包括$A\cos\omega_rt+B\sin\omega_rt$及上述常量特解，$A,B$由径向初态确定。用前述角坐标的导数代入拉格朗日方程，并除以$mr$，精确角向方程为
\[
r\ddot\theta+2\dot r\dot\theta+g\sin\theta=0.
\]
径向变化会通过这两个含r的项反过来影响摆角；图\ref{deep4:fig:d7-e7-s}中的“驱动”是同一系统内部的能量交换。
随时间增长的项说明首阶角向解只适用于早期运动。实际系统总能量守恒，径向能量增加会反过来改变摆动；当 $|\rho|/r_0$ 或角向振幅的变化已不可忽略时，须将两个非线性方程联立求解。

还可以由能量估计早期近似的时间范围。在精确内共振$\omega_r=2\omega_\theta$处，利用$g=r_0\omega_\theta^2$，久期项的振幅为
\[
 \frac{A_\rho(t)}{r_0}=\frac{3\Theta^2\omega_\theta t}{16}.
\]
取多个振动周期以后的缓变包络，径向振动能约为$m\omega_r^2A_\rho^2/2$，角向初始能量约为$mr_0^2\omega_\theta^2\Theta^2/2$，两者之比为
\[
 \frac{E_r}{E_{\theta,0}}\simeq
 \frac{4A_\rho^2}{r_0^2\Theta^2}
 =\frac9{64}\Theta^2(\omega_\theta t)^2.
\]
忽略角向振幅的变化至少要求$\Theta\omega_\theta t\ll1$。当$t$达到$1/(\Theta\omega_\theta)$量级，能量交换已不能视为微小修正，即使径向位移相对于$r_0$仍小，固定角向振幅的驱动近似也会失效。此估计给出近似失效的尺度，并不预测能量随后继续无界增长。
\end{solution}



\section{习题与解答}

\begin{exercise}[双摆下端受冲量]\label{sup4:problem-exercise-7}\label{sup4:ex:d7-h7}
例\ref{sup4:ex:d7-e6}的等长等质量双摆初始竖直静止，对下端质点施加向右水平冲量 $J$。求冲量后的两个角速度、此后的小振动及各模态能量。设冲量足够小，使随后摆角均可线性化。
\end{exercise}
\begin{figure}[H]\centering
\begin{tikzpicture}
\fill(0,0)circle(1.6pt)node[above]{$O$};\draw[curve](0,0)--(0,-3.6);\fill(0,-1.8)circle(2.5pt)node[left]{$m$};\fill(0,-3.6)circle(2.5pt)node[left]{$m$};\node[right]at(0,-.9){$\ell$};\node[right]at(0,-2.7){$\ell$};\draw[axis](0,-3.6)--(1.8,-3.6)node[right]{$J$};\node[align=left]at(3,-1.4){两角从竖直向下量起\\正向均向右};
\end{tikzpicture}
\caption{冲量作用在下端质点，两个角坐标的广义冲量都须由作用点位移求出。}\label{sup4:fig:d7-h7-q}
\end{figure}
\begin{solution}
图\ref{sup4:fig:d7-h7-q}中，下端水平位置为 $x_B=\ell(\sin\theta_1+\sin\theta_2)$。在竖直构形，虚位移为 $\delta x_B=\ell(\delta\theta_1+\delta\theta_2)$，故广义冲量为 $(J\ell,J\ell)$。利用双摆的质量矩阵，
\[
m\ell^2\begin{pmatrix}2&1\\1&1\end{pmatrix}
\begin{pmatrix}\dot\theta_1(0^+)\\\dot\theta_2(0^+)\end{pmatrix}
=J\ell\begin{pmatrix}1\\1\end{pmatrix},
\]
竖直构形处，质量矩阵的无量纲部分逆矩阵为
\[
\begin{pmatrix}2&1\\1&1\end{pmatrix}^{-1}
=\begin{pmatrix}1&-1\\-1&2\end{pmatrix}.
\]
乘上广义冲量后，第一行两个分量抵消，第二行剩一个单位分量。这一消元说明实际施加在下端的冲量，不能只记在第二个角坐标上。
解得 $\dot\theta_1(0^+)=0$、$\dot\theta_2(0^+)=J/(m\ell)$。上端质点初始水平速度仍为零，下端速度为 $J/m$，与初始两杆均竖直、约束冲量不产生水平分量相符。
\begin{figure}[H]\centering
\begin{tikzpicture}
\node at(1.2,2.3){冲量后瞬间};\fill(0,1)circle(2pt)node[left]{上端};\fill(0,-.4)circle(2pt)node[left]{下端};\node[right]at(0,1){$v_A=0$};\draw[axis](0,-.4)--(1.6,-.4)node[right]{$v_B=J/m$};
\draw[axis](4,0)--(7,0);\draw[axis](4,0)--(4,2.4)node[above]{能量比例};\draw[curve,fill=gray!15](4.4,0)rectangle(5.3,1.707);\draw[curve,fill=gray!15](6,0)rectangle(6.9,.293);\node[above]at(4.85,1.707){$(2+\sqrt2)/4$};\node[above]at(6.45,.35){$(2-\sqrt2)/4$};\node[below]at(4.85,0){低频};\node[below]at(6.45,0){高频};
\end{tikzpicture}
\caption{初始水平速度及两模态的能量份额，低频模态获得较多能量。}\label{sup4:fig:d7-h7-s}
\end{figure}
两振型仍为 $(1,\sqrt2)$、$(1,-\sqrt2)$。由初始位移为零，两模态时间因子均为正弦。先求模态初速度 $b_-,b_+$：由图\ref{sup4:fig:d7-h7-s}的实际初速度，
\[
b_-+b_+=0,\qquad
\sqrt2(b_--b_+)=\frac{J}{m\ell}.
\]
第一式给 $b_+=-b_-$，代入第二式得 $2\sqrt2b_-=J/(m\ell)$，所以
\[
b_-=\frac{J}{2\sqrt2m\ell},\qquad b_+=-\frac{J}{2\sqrt2m\ell}.
\]
模态坐标分别为 $q_\pm=(b_\pm/\omega_\pm)\sin\omega_\pm t$；再由 $\theta_1=q_-+q_+$、$\theta_2=\sqrt2(q_--q_+)$，得到
\[
\boxed{\begin{aligned}
\theta_1&=\frac{J}{2\sqrt2m\ell}\left(\frac{\sin\omega_-t}{\omega_-}-\frac{\sin\omega_+t}{\omega_+}\right),\\
\theta_2&=\frac{J}{2m\ell}\left(\frac{\sin\omega_-t}{\omega_-}+\frac{\sin\omega_+t}{\omega_+}\right).
\end{aligned}}
\]
未归一化振型的等效惯量为 $m\ell^2(4+2\sqrt2)$、$m\ell^2(4-2\sqrt2)$，两模态初始速度系数的大小均为 $J/(2\sqrt2m\ell)$。因此
\[
E_- =\frac{J^2}{8m}(2+\sqrt2),\qquad
E_+ =\frac{J^2}{8m}(2-\sqrt2).
\]
相加为 $J^2/(2m)$，等于图\ref{sup4:fig:d7-h7-s}中下端质点的初始动能。若只给第二角坐标一个广义冲量，会错误地算出上端质点也获得水平初速度。
\end{solution}


\begin{exercise}[弹簧能否稳定两个倒立摆（选做）]\label{sup4:problem-exercise-12}\label{sup4:ex:d7-h12}
两个相同倒立摆各由轻杆和端点质量 $m$ 组成，杆长为 $\ell$，下端支点固定且等高，间距为 $D$。两球在竖直向上的平衡态由水平弹簧 $\kappa$ 连接，弹簧原长为 $D$，可伸可压。两杆可独立在同一竖直平面转动，摆角均从竖直向上向右量起。求线性模态，判断增大 $\kappa$ 能否使整个平衡稳定，再讨论严格反对称运动在临界刚度处的稳定性。
\end{exercise}
\begin{figure}[H]\centering
\begin{tikzpicture}
\fill(0,0)circle(1.7pt)node[below]{$P_1$};\fill(4.5,0)circle(1.7pt)node[below]{$P_2$};\draw[curve](0,0)--(0,2.5);\draw[curve](4.5,0)--(4.5,2.5);\fill(0,2.5)circle(2.7pt)node[above]{$m$};\fill(4.5,2.5)circle(2.7pt)node[above]{$m$};\draw[spring](.15,2.5)--(4.35,2.5)node[midway,above]{$\kappa$};\node[left]at(0,1.2){$\ell$};\node[right]at(4.5,1.2){$\ell$};\draw[<->](0,-.65)--(4.5,-.65)node[midway,below]{$D$};
\end{tikzpicture}
\caption{竖直向上的构形是相对位移零点，两杆均允许独立向左右偏转。}\label{sup4:fig:d7-h12-q}
\end{figure}
\begin{solution}
图\ref{sup4:fig:d7-h12-q}中，球的水平位移一阶为 $\ell\theta_i$，弹簧伸长为 $\ell(\theta_2-\theta_1)$。因此
\[
T_2=\frac12m\ell^2(\dot\theta_1^2+\dot\theta_2^2),\quad
U_2=U_*-\frac12mg\ell(\theta_1^2+\theta_2^2)
+\frac12\kappa\ell^2(\theta_2-\theta_1)^2.
\]
取同相、反相坐标 $u=(\theta_1+\theta_2)/2$、$v=(\theta_1-\theta_2)/2$，得到
\[
\ddot u-\frac g\ell u=0,\qquad
\ddot v+\left(\frac{2\kappa}{m}-\frac g\ell\right)v=0.
\]
\begin{figure}[H]\centering
\begin{tikzpicture}
\begin{scope}
\foreach\x in{0,2}{\fill(\x,0)circle(1.5pt);\draw[curve](\x,0)--++(.5,1.8);\fill(\x+.5,1.8)circle(2pt);}\draw[spring](.5,1.8)--(2.5,1.8);\draw[axis](.5,2.2)--(1.3,2.2);\draw[axis](2.5,2.2)--(3.3,2.2);\node at(1.3,-.65){同相：弹簧不变长};
\end{scope}
\begin{scope}[xshift=5cm]
\fill(0,0)circle(1.5pt);\fill(2.5,0)circle(1.5pt);\draw[curve](0,0)--(.45,1.8);\draw[curve](2.5,0)--(2.05,1.8);\fill(.45,1.8)circle(2pt);\fill(2.05,1.8)circle(2pt);\draw[spring](.45,1.8)--(2.05,1.8);\node at(1.25,-.65){反相：弹簧提供回复项};
\end{scope}
\end{tikzpicture}
\caption{同相倾斜保留原来的球间距离，耦合弹簧无法抑制该不稳定方向。}\label{sup4:fig:d7-h12-s}
\end{figure}
图\ref{sup4:fig:d7-h12-s}中的同相模态始终有$\omega_u^2=-g/\ell<0$。令$\gamma=\sqrt{g/\ell}$，特征方程$s^2-\gamma^2=0$有两个实根$\pm\gamma$，所以$u=C_+\ee^{\gamma t}+C_-\ee^{-\gamma t}$。例如初始$u(0)=\varepsilon,\dot u(0)=0$，则$C_++C_-=\varepsilon$、$\gamma(C_+-C_-)=0$，解得$C_+=C_-=\varepsilon/2$，即$u=\varepsilon\cosh(\gamma t)$。任意小偏差均会增长，所以任何有限$\kappa$都不能使整个平衡稳定。

同相运动时，$\theta_1=\theta_2$，两球具有相同的水平、竖直位移，相对位置始终为 $(D,0)$，弹簧长度保持D。图\ref{sup4:fig:d7-h12-s}左图的共同偏转因此始终不受弹簧的回复作用，这一结论在有限小偏角下也成立。
严格限定 $u=0$ 的反对称运动则在 $\kappa>mg/(2\ell)$ 时具有正线性频率平方。在临界 $\kappa=mg/(2\ell)$，线性频率为零，须继续展开势能。令 $\theta_1=v,\theta_2=-v$，小角范围内两球等高，弹簧伸长为 $-2\ell\sin v$，故
\[
U(v)=2mg\ell\cos v+2\kappa\ell^2\sin^2v.
\]
在临界刚度处，第二项系数为$mg\ell$。将$\cos v=1-v^2/2+v^4/24+O(v^6)$、$\sin^2v=v^2-v^4/3+O(v^6)$分别代入，得
\[
\begin{aligned}
U(v)&=2mg\ell-mg\ell v^2+\frac{mg\ell}{12}v^4
+mg\ell v^2-\frac{mg\ell}{3}v^4+O(v^6)\\
&=U_*-\frac{mg\ell}{4}v^4+O(v^6).
\end{aligned}
\]
四次项为负，反对称子空间内的平衡也不稳定。即使刚度超过此临界值、反相模态转为稳定，同相模态的负频率平方仍使整个平衡不稳定。
在临界刚度、严格反对称的子空间内，动能为 $m\ell^2\dot v^2$，等效惯量为 $2m\ell^2$。由四次势能项求导，
\[
2m\ell^2\ddot v=-U'(v)=mg\ell v^3+O(v^5),
\qquad \ddot v=\frac g{2\ell}v^3+O(v^5).
\]
小的正偏角获得正加速度，小的负偏角获得负加速度，运动都远离平衡。它与第六讲临界圆环的正四次势阱符号相反，零线性频率本身不足以给出稳定性结论。
\end{solution}


\begin{exercise}[包络交接时能量在哪里]\label{sup4:problem-exercise-4}\label{sup4:ex:d7-h4}
例\ref{sup4:ex:d7-e3}中取 $\kappa=3k/2$，仍从 $x_1=A,x_2=0$、速度均为零释放。求 $t_* =\pi/\omega_s$ 时两块的位置、速度及三根弹簧各自储能。是否可以说此时全部能量都在第二块及其外侧弹簧中？
\end{exercise}
\begin{figure}[H]\centering
\begin{tikzpicture}
\draw[curve](0,0)--(7,0);\foreach\x in{0,7}{\draw[curve](\x,0)--(\x,1.1);}\draw[spring](0,.5)--(2,.5);\draw[curve](2,.2)rectangle(2.7,.8);\draw[spring](2.7,.5)--(4.2,.5);\draw[curve](4.2,.2)rectangle(4.9,.8);\draw[spring](4.9,.5)--(7,.5);\node[above]at(1,.6){$k$};\node[above]at(3.45,.6){$3k/2$};\node[above]at(5.9,.6){$k$};\node[below]at(2.35,0){初始 $A$};\node[below]at(4.55,0){初始 $0$};
\end{tikzpicture}
\caption{耦合不弱，中间弹簧的储能不能忽略。}\label{sup4:fig:d7-h4-q}
\end{figure}
\begin{solution}
图\ref{sup4:fig:d7-h4-q}的反相频率为 $\omega_a=\sqrt{(k+3k)/m}=2\omega_s$。在 $t_* =\pi/\omega_s$，有 $\cos\omega_st_*=-1$、$\cos\omega_at_*=1$，两正弦均为零，所以
\[
x_1(t_*)=0,\qquad x_2(t_*)=-A,\qquad
\dot x_1(t_*)=\dot x_2(t_*)=0.
\]
此时左弹簧无形变，右弹簧形变量大小为 $A$，中间弹簧形变量为 $x_2-x_1=-A$。
\begin{figure}[H]\centering
\begin{tikzpicture}
\draw[axis](0,0)--(6.5,0)node[right]{储能部分};\draw[axis](0,0)--(0,2.7)node[above]{$E_i/E$};\draw[curve](1,0)rectangle(2,0);\draw[curve,fill=gray!15](3,0)rectangle(4,2.1);\draw[curve,fill=gray!15](5,0)rectangle(6,1.4);\node[below]at(1.5,0){左弹簧};\node[below]at(3.5,0){中间弹簧};\node[below]at(5.5,0){右弹簧};\node[above]at(1.5,0){$0$};\node[above]at(3.5,2.1){$3/5$};\node[above]at(5.5,1.4){$2/5$};
\end{tikzpicture}
\caption{指定时刻的能量分布。两物块动能均为零，中间弹簧仍储存五分之三的总能量。}\label{sup4:fig:d7-h4-s}
\end{figure}
由图\ref{sup4:fig:d7-h4-s}，
\[
U_L=0,\qquad U_C=\frac12\kappa A^2=\frac34kA^2,\qquad U_R=\frac12kA^2.
\]
总能量为 $5kA^2/4$，等于初态的 $\frac12kA^2+\frac12\kappa A^2$。第一块虽瞬时静止于平衡位置，右外侧弹簧中仍只有总能量的五分之二，其余五分之三储存在耦合弹簧中，不能全归于某一块独立振子的能量。
瞬时静止也不表示此后停在原处。由第一块受力方程，在 $t_*$ 有
\[
\ddot x_1(t_*)=\frac{\kappa x_2(t_*)}{m}=-\frac{\kappa A}{m}\ne0.
\]
耦合弹簧的力随后使两块重新运动。图\ref{sup4:fig:d7-h4-s}记录的是一个瞬时能量分布；两速度及弹簧能量变化率在此刻为零，受力却未平衡。
\end{solution}


\begin{exercise}[反共振频率下从静止启动]\label{sup4:problem-exercise-9}\label{sup4:ex:d7-h9}
例\ref{sup4:ex:d7-e8}中取 $\kappa=k$，从两块均静止于平衡位置开始，对第一块施加 $F_0\cos\nu t$，其中 $\nu=\sqrt{2k/m}$。求完整运动，并解释为什么第一块仍会运动。
\end{exercise}
\begin{figure}[H]\centering
\begin{tikzpicture}
\draw[axis](0,0)--(6.2,0)node[right]{频率};\draw[axis](0,0)--(0,1.9);\draw[curve](1.3,0)--(1.3,.8);\draw[curve](3,0)--(3,1.4);\draw[curve](4.8,0)--(4.8,.8);\node[below]at(1.3,0){$\omega_s$};\node[below]at(3,0){$\nu$};\node[below]at(4.8,0){$\omega_a$};\node[above]at(3,1.4){驱动频率};\node at(3,-1){初态：$x_1=x_2=\dot x_1=\dot x_2=0$};
\end{tikzpicture}
\caption{驱动频率位于两固有频率之间，频率位置仅按次序示意。}\label{sup4:fig:d7-h9-q}
\end{figure}
\begin{solution}
图\ref{sup4:fig:d7-h9-q}中 $\omega_s=\sqrt{k/m}$、$\omega_a=\sqrt{3k/m}$。反共振周期特解为 $x_{1p}=0$、$x_{2p}=-(F_0/k)\cos\nu t$。它在 $t=0$ 的第二块位移不为零，因此尚未满足题设初态。

令$s=(x_1+x_2)/2$、$d=(x_1-x_2)/2$，两个模态方程均具有形式
\[
\ddot q+\omega^2q=\frac{F_0}{2m}\cos\nu t,
\qquad q(0)=\dot q(0)=0,
\]
其中$\omega$分别取$\omega_s,\omega_a$。由于$\nu\ne\omega$，设特解$q_p=C\cos\nu t$，代入得$C(\omega^2-\nu^2)=F_0/(2m)$。通解是
\[
q=A\cos\omega t+B\sin\omega t+C\cos\nu t.
\]
在$t=0$代入位移条件，得$A+C=0$；求导并代入速度条件，得$B\omega=0$。所以$A=-C,B=0$，两模态的完整解为
\[
\begin{aligned}
s&=\frac{F_0}{2m(\omega_s^2-\nu^2)}(\cos\nu t-\cos\omega_st),\\
d&=\frac{F_0}{2m(\omega_a^2-\nu^2)}(\cos\nu t-\cos\omega_at).
\end{aligned}
\]
本题两个分母分别为 $-2k$、$2k$，相加求 $x_1=s+d$ 时，驱动频率的余弦项恰好相消；固有频率项则继续保留。这与图\ref{sup4:fig:d7-h9-s}的三条频率标记对应。
自由解在初始时须补入位移 $(0,F_0/k)$，初速度仍为零。将这一位移分解成同相与反相部分，得到
\[
\boxed{\begin{aligned}
x_1&=\frac{F_0}{2k}(\cos\omega_st-\cos\omega_at),\\
x_2&=-\frac{F_0}{k}\cos\nu t+\frac{F_0}{2k}(\cos\omega_st+\cos\omega_at).
\end{aligned}}
\]
\begin{figure}[H]\centering
\begin{tikzpicture}
\draw[axis](0,0)--(6.2,0)node[right]{频率};\draw[axis](0,-1.5)--(0,1.7)node[above]{第一块余弦系数};\draw[curve](1.3,0)--(1.3,1);\fill(3,0)circle(1.7pt);\draw[curve](4.8,0)--(4.8,-1);\node[above]at(1.3,1){$F_0/(2k)$};\node[below]at(4.8,-1){$-F_0/(2k)$};\node[below]at(1.3,0){$\omega_s$};\node[above]at(3,0){$0$};\node[below]at(3,-.35){$\nu$};\node[above]at(4.8,0){$\omega_a$};
\end{tikzpicture}
\caption{第一块在驱动频率处的系数为零，但两个自由频率分量均不为零。}\label{sup4:fig:d7-h9-s}
\end{figure}
图\ref{sup4:fig:d7-h9-s}中，第一块的 $\nu$ 分量为零，两个自由频率分量则因初始条件而保留。无阻尼时它们不衰减，这就是启动后第一块继续运动的原因。代入 $t=0$，得 $x_1=0$、$\ddot x_1=F_0/m$，与初始外力相符。
\end{solution}


\begin{exercise}[由频率求弹簧与模态能量]\label{deep4:problem-exercise-3}\label{deep4:ex:d7-h3}
例\ref{sup4:ex:d7-e3}的对称双振子已测得同相、反相频率分别为 $\omega_s,\omega_a$，且 $\omega_a>\omega_s$，两质量 $m$ 已知。求 $k,\kappa$。若初始 $x_1=A>0,x_2=B$，$0<B<A$，两速度为零，求使两个模态能量相等的 $B$。
\end{exercise}
\begin{figure}[H]\centering
\begin{tikzpicture}
\foreach\y/\txt in{1/同相,-1/反相}{\node[left]at(0,\y){\txt};\fill(1,\y)circle(2.3pt);\fill(4,\y)circle(2.3pt);}\draw[axis](1,1)--(2,1);\draw[axis](4,1)--(5,1);\draw[axis](1,-1)--(2,-1);\draw[axis](4,-1)--(3,-1);\node[above]at(2.5,1.5){频率 $\omega_s$};\node[below]at(2.5,-1.5){频率 $\omega_a$};
\end{tikzpicture}
\caption{测得的两频率分别对应中间弹簧不伸长与有伸长的模态。}\label{deep4:fig:d7-h3-q}
\end{figure}
\begin{solution}
图\ref{deep4:fig:d7-h3-q}的两模态满足 $m\omega_s^2=k$、$m\omega_a^2=k+2\kappa$，所以
\[
\boxed{k=m\omega_s^2,\qquad \kappa=\frac m2(\omega_a^2-\omega_s^2)>0.}
\]
振型 $(1,1)$ 与 $(1,-1)$ 都未按单位长度归一化。代入 $x_1=s+d$、$x_2=s-d$ 后，
\[
T=m(\dot s^2+\dot d^2)
=\frac12(2m)\dot s^2+\frac12(2m)\dot d^2.
\]
两个模态的惯性系数均为 $2m$。若误取单块质量m，每个模态的能量与总能量都会少算一半，而能量比仍保持不变。
初始同相、反相坐标为 $s_0=(A+B)/2$、$d_0=(A-B)/2$。相应等效质量均为 $2m$，故
\[
E_s=\frac m4\omega_s^2(A+B)^2,\qquad
E_a=\frac m4\omega_a^2(A-B)^2.
\]
\begin{figure}[H]\centering
\begin{tikzpicture}
\draw[axis](0,0)--(5.2,0)node[right]{$B/A$};\draw[axis](0,0)--(0,3.3)node[above]{模态能量};\draw[curve,domain=0:1,samples=70]plot({4.3*\x},{.65*(1+\x)^2});\draw[guide,domain=0:1,samples=70]plot({4.3*\x},{2.6*(1-\x)^2});\fill(1.433,1.156)circle(1.7pt);\draw[guide](1.433,1.156)--(2,2.2);\node[above]at(2.2,2.4){能量相等};\node[below]at(4.3,0){$1$};\node at(3.6,3){$\omega_a=2\omega_s$};
\end{tikzpicture}
\caption{同相能量随B增大而增加，反相能量减小；图示一组频率比。}\label{deep4:fig:d7-h3-s}
\end{figure}
由于 $0<B<A$，两式开正平方根即可得 $\omega_s(A+B)=\omega_a(A-B)$，因而
\[
\boxed{B=A\frac{\omega_a-\omega_s}{\omega_a+\omega_s}.}
\]
图\ref{deep4:fig:d7-h3-s}的交点唯一位于 $0<B/A<1$。模态能量相等并不要求模态位移系数相等，较高频模态在相同振幅下储存的能量较多。
\end{solution}


\begin{exercise}[失谐对拍振包络的抑制（选做）]\label{deep4:problem-exercise-6}\label{deep4:ex:d7-h6}
沿用例\ref{sup4:ex:d7-e5}的失谐双摆，取初始 $x_1=A,x_2=0$、两速度为零。记 $c=\kappa/m>0$，两频率为 $\omega_-<\omega_+$。求第二摆的位移及慢包络上界，并求包络上界降至 $A/2$ 时的失谐量。说明包络上界与实际位移最大值的区别。
\end{exercise}
\begin{figure}[H]\centering
\begin{tikzpicture}
\draw[axis](-1.8,0)--(5,0)node[right]{$x_1$};\draw[axis](0,-.5)--(0,3.5)node[above]{$x_2$};\draw[axis](0,0)--(2.8,1.5)node[right]{$\vct u_-$};\draw[axis](0,0)--(-1.5,2.8)node[above]{$\vct u_+$};\draw[curve](.8,0)arc(0:28:.8);\node at(1.1,.3){$\psi$};\fill(3.5,0)circle(2pt)node[below]{$(A,0)$};
\end{tikzpicture}
\caption{等质量时，可用相差直角的两个单位振型分解初始位移。}\label{deep4:fig:d7-h6-q}
\end{figure}
\begin{solution}
在图\ref{deep4:fig:d7-h6-q}中取 $\vct u_-=(\cos\psi,\sin\psi)$、$\vct u_+=(-\sin\psi,\cos\psi)$，其中 $0<\psi<\pi/2$。低频振型关系给出混合角的正弦和余弦：
\[
\cos2\psi=-\frac d{\sqrt{d^2+c^2}},\qquad
\sin2\psi=\frac c{\sqrt{d^2+c^2}}.
\]
它们等价于低频振型的 $\tan\psi=(d+\sqrt{d^2+c^2})/c$，使图\ref{deep4:fig:d7-h6-q}中的低频向量始终落在第一象限。仅用正弦关系会留下把两个摆球混合程度互换的歧义。
初态在两模态上的投影为 $A\cos\psi$、$-A\sin\psi$，所以
\[
x_2=A\cos\psi\sin\psi(\cos\omega_-t-\cos\omega_+t).
\]
令 $\bar\omega=(\omega_++\omega_-)/2$、$\Delta\omega=\omega_+-\omega_-$，得
\[
\boxed{x_2=\frac{Ac}{\sqrt{d^2+c^2}}\sin\bar\omega t\,
\sin\frac{\Delta\omega t}{2}.}
\]
\begin{figure}[H]\centering
\begin{tikzpicture}
\draw[axis](-3.4,0)--(3.6,0)node[right]{$d/c$};\draw[axis](0,0)--(0,2.7)node[above]{$A_{2,\rm env}^{\max}/A$};\draw[curve,domain=-3:3,samples=120]plot(\x,{2/sqrt(1+(\x)^2)});\draw[guide](-3,1)--(3,1);\draw[guide](1.732,0)--(1.732,1);\draw[guide](-1.732,0)--(-1.732,1);\node[below]at(1.732,0){$\sqrt3$};\node[below]at(-1.732,0){$-\sqrt3$};\node[left]at(0,2){$1$};\node[left]at(0,1){$1/2$};
\end{tikzpicture}
\caption{失谐使混合减弱，两侧对称地降低第二摆的包络上界。}\label{deep4:fig:d7-h6-s}
\end{figure}
由图\ref{deep4:fig:d7-h6-s}，包络上界为 $A_{2,\rm env}^{\max}=Ac/\sqrt{d^2+c^2}$。令它为 $A/2$，得到
\[
\boxed{|d|=\sqrt3c,}
\]
该失谐量还须满足 $|d|<\omega_c^2$。当 $|d|\gg c$，第二摆的响应明显减弱。两频率足够接近时，第二个正弦因子才描述慢包络；实际位移还乘有快速因子 $\sin\bar\omega t$。包络到达峰值时，快速因子未必恰为 $\pm1$，因此一拍内实际达到的最大位移可以小于这个上界。
\end{solution}


\begin{exercise}[中间质量改变后的三振子链]\label{deep4:problem-exercise-10}\label{deep4:ex:d7-h10}
例\ref{deep4:ex:d7-e9}的三振子链中，两外块质量仍为 $m$，中间块改为 $2m$，四根弹簧仍均为 $k$。求全部频率和振型，确定只对中间块施加冲量时不被激起的模态。
\end{exercise}
\begin{figure}[H]\centering
\begin{tikzpicture}
\foreach\x/\mass in{0/m,3/2m,6/m}{\draw[curve](\x,0)rectangle++(1,.7);\node at(\x+.5,.35){$\mass$};}\draw[spring](1,.35)--(3,.35);\draw[spring](4,.35)--(6,.35);\draw[spring](-1.5,.35)--(0,.35);\draw[spring](7,.35)--(8.5,.35);\draw[curve](-1.5,-.3)--(-1.5,1);\draw[curve](8.5,-.3)--(8.5,1);\draw[axis](3.5,1.2)--(5,1.2)node[right]{$J$};\node at(3.5,-.7){四弹簧均为 $k$};
\end{tikzpicture}
\caption{质量分布仍左右对称，中间质量改变不破坏镜面对称性。}\label{deep4:fig:d7-h10-q}
\end{figure}
\begin{solution}
图\ref{deep4:fig:d7-h10-q}中，反对称振型仍为 $(1,0,-1)$。中间块固定不动，两外块各受两根弹簧的回复力，因此 $\omega_a^2=2k/m$，与中间质量无关。

对称模态令 $x_1=x_3=x$、$x_2=y$，方程为
\[
m\ddot x+2kx-ky=0,\qquad
2m\ddot y+2ky-2kx=0.
\]
置 $\eta=m\omega^2/k$，得 $(2-\eta)x-y=0$、$-x+(1-\eta)y=0$。故
\[
\eta^2-3\eta+1=0,
\quad \boxed{\omega_\pm^2=\frac{k}{m}\frac{3\pm\sqrt5}{2}.}
\]
取外块振幅为1，对应振型为
\[
\vct u_-=(1,(1+\sqrt5)/2,1),\qquad
\vct u_+=(1,(1-\sqrt5)/2,1).
\]
\begin{figure}[H]\centering
\begin{tikzpicture}
\foreach\y/\txt in{2/低频,0/反对称,-2/高频}{\node[left]at(0,\y){\txt};\foreach\x in{1,3.5,6}{\fill(\x,\y)circle(2pt);}}
\draw[axis](1,2)--(1.8,2);\draw[axis](3.5,2)--(4.794,2);\draw[axis](6,2)--(6.8,2);
\draw[axis](1,0)--(1.8,0);\node[above]at(3.5,0){$u_2=0$};\draw[axis](6,0)--(5.2,0);
\draw[axis](1,-2)--(1.8,-2);\draw[axis](3.5,-2)--(3.006,-2);\draw[axis](6,-2)--(6.8,-2);
\end{tikzpicture}
\caption{改变中间质量后，两个对称模态的比例改变，反对称节点保留。}\label{deep4:fig:d7-h10-s}
\end{figure}
图\ref{deep4:fig:d7-h10-s}按频率从低到高排列：$(3-\sqrt5)/2<2<(3+\sqrt5)/2$。两个对称振型满足质量正交，因为令 $b_-=(1+\sqrt5)/2$、$b_+=(1-\sqrt5)/2$ 后，$b_-b_+=-1$，故
\[
\vct u_-^{\mathsf T}\operatorname{diag}(m,2m,m)\vct u_+
=2m+2mb_-b_+=0.
\]
中间冲量在反对称振型上的投影为零，该模态仍不出现。左右对称性保留了中间节点，中间质量的改变不影响这一选择规则。
把中间冲量记作 $\vct J=(0,J,0)^{\mathsf T}$，质量矩阵为 $M=\operatorname{diag}(m,2m,m)$。任一未归一化模态的初速度系数应由
\[
\dot q_j(0^+)=\frac{\vct u_j^{\mathsf T}\vct J}{\vct u_j^{\mathsf T}M\vct u_j}
\]
求出。图\ref{deep4:fig:d7-h10-s}中反对称振型的中间分量为零，分子随之为零；两个对称振型的中间分量非零，均会被激起。该判断同时用到了题设的零初位移，若初态已有反对称形变，该模态不会被一次中间冲量消除。
\end{solution}


\chapter{刚体碰撞进阶}



\section{碰撞守恒量、冲量与有效质量}

碰撞中可用的守恒量取决于短时间内的外冲量。恢复系数约束接触点的法向相对速度；刚体发生转动时，该点的有效质量通常不同于物体质量。



\subsection{碰撞过程中的守恒量}
碰撞持续时间很短，接触力却可能很大。求碰后的运动时，通常只需计算接触力的冲量，无须求出它随形变变化的全过程。以下采用经典力学中的刚体近似，碰撞前后的量分别加上标 $-$、$+$。

\subsubsection{短时近似与外冲量}
质点在接触时间 $\tau$ 内满足
\begin{equation}
 m(\vct v^+-\vct v^-)=\vct J_{\rm c}+\int_{t_-}^{t_+}\vct F_{\rm ext}\dd t.
\end{equation}
其中 $\vct J_{\rm c}$ 是接触冲量。若重力与接触冲量相比很小，即 $mg\tau\ll |\vct J_{\rm c}|$，可忽略重力冲量。这个条件只适用于碰撞阶段，随后的飞行或摆动仍须计入重力。

支座反力、绳的张力和轨道约束力可能随接触力一起增大。它们的冲量并不因 $\tau$ 小而自动消失。图\ref{sup5:fig:d4-t1}中的两球若都自由运动，总动量守恒；若其中一球被固定，固定装置会向两球系统传递冲量。

\begin{figure}[H]\centering
\begin{tikzpicture}[x=1cm,y=1cm]
\draw[curve] (0,0) circle(.48) (3,0) circle(.64);
\node at(0,0){$m_1$};\node at(3,0){$m_2$};
\draw[axis](-1.5,.9)--(-.3,.9) node[midway,above]{$\vct v_1^-$};
\draw[axis](.6,0)--(1.8,0) node[midway,below]{$\vct n$};
\draw[axis](-.15,-.95)--(-1.25,-.95) node[midway,below]{$-J\vct n$};
\draw[axis](3.15,-.95)--(4.25,-.95) node[midway,below]{$J\vct n$};
\draw[guide](-1.8,-1.5) rectangle(4.8,1.4);
\node at(1.5,1.8){系统边界包含两个物体};
\end{tikzpicture}
\caption{内冲量相互抵消。法向单位矢量 $\vct n$ 从物体1指向物体2。}\label{sup5:fig:d4-t1}
\end{figure}

对于由多个物体组成的系统，先把各物体的动量方程相加，再判断哪一个分量没有外冲量。光滑水平面上的可动斜面与小球碰撞时，系统水平动量守恒，竖直动量通常不守恒。固定铰链处的冲量对铰点力矩为零，因此关于铰点的角动量可守恒，系统的线动量却可改变。

刚体的冲量方程应写成
\begin{equation}\label{sup5:eq:d4-rigid}
 M\Delta\vct V_C=\sum\vct J_i,\qquad
 I_C\Delta\omega=\sum(\vct r_i\times\vct J_i)_z.
\end{equation}
这里 $\vct r_i$ 从质心指向受冲量点，$I_C$ 是关于质心且垂直于运动平面的转动惯量。碰撞期间构形近似不变，速度和角速度可以有有限跃变。

\subsubsection{恢复系数约束的是接触点的相对速度}
取 $\vct n$ 从物体1指向物体2，令两接触点的速度为 $\vct v_{P1}$、$\vct v_{P2}$，定义
\begin{equation}
 g_n=(\vct v_{P2}-\vct v_{P1})\cdot\vct n.
\end{equation}
接近时 $g_n^-<0$。牛顿恢复系数定义为
\begin{equation}\label{sup5:eq:d4-e}
 g_n^+=-e g_n^-,\qquad 0\le e\le1.
\end{equation}
本讲把 $e$ 视为给定的碰撞参数。$e=1$ 表示法向相对速度反号；$e=0$ 表示碰撞结束时法向相对速度为零，并不保证两个物体永久黏合。黏合还要求接触处的切向相对运动消失，并在后续运动中保持连接。

若物体发生转动，\eqref{sup5:eq:d4-e}必须使用接触点速度，不能直接用质心速度代替。光滑接触没有切向冲量；只有对不受其他切向冲量的质点，才可直接断言其切向速度不变。


\subsection{冲量、约化质量与能量损失}
\subsubsection{质点正碰}
两个自由质点沿同一直线接触，设碰前速度为 $u_1,u_2$，正方向与 $\vct n$ 相同，且 $u_1>u_2$。物体1受到冲量 $-J$，物体2受到 $J$，于是
\begin{align}
 v_1&=u_1-J/m_1,&v_2&=u_2+J/m_2,\\
 v_2-v_1&=e(u_1-u_2).
\end{align}
消去速度得
\begin{equation}\label{sup5:eq:d4-Jpoint}
 J=(1+e)\mu(u_1-u_2),\qquad
 \mu=\frac{m_1m_2}{m_1+m_2}.
\end{equation}
代回即有
\begin{align}\label{sup5:eq:d4-vpoint}
 v_1&=\frac{(m_1-em_2)u_1+(1+e)m_2u_2}{m_1+m_2},\\
 v_2&=\frac{(1+e)m_1u_1+(m_2-em_1)u_2}{m_1+m_2}.
\end{align}
质心速度 $V=(m_1u_1+m_2u_2)/(m_1+m_2)$ 不变。把总动能分解为质心平动和相对运动两部分，得到
\begin{equation}
 T=\frac12(m_1+m_2)V^2+\frac12\mu(u_1-u_2)^2.
\end{equation}
恢复系数只把相对速度的大小变为原来的 $e$ 倍。因此能量损失为
\begin{equation}\label{sup5:eq:d4-loss}
 Q=T^--T^+=\frac12\mu(1-e^2)(u_1-u_2)^2.
\end{equation}

\begin{figure}[H]\centering
\begin{tikzpicture}
\draw[axis](-.3,0)--(6.6,0) node[right]{$e$};
\draw[axis](0,-.2)--(0,2.8) node[above]{$Q/T_{\rm rel}^-$};
\draw[curve,domain=0:1,samples=60] plot({6*\x},{2.2*(1-\x*\x)});
\draw[guide](0,2.2)--(6,2.2);\node[left]at(0,2.2){$1$};
\node[below]at(0,0){$0$};\node[below]at(6,0){$1$};
\node[above right]at(2.6,1.79){$1-e^2$};
\end{tikzpicture}
\caption{自由质点正碰的相对动能损失，$T_{\rm rel}^-=\mu(u_1-u_2)^2/2$。}\label{sup5:fig:d4-t2}
\end{figure}

图\ref{sup5:fig:d4-t2}纵轴以碰前相对动能为单位。$e=0$ 时相对动能全部损失，质心平动动能仍保留；只有在质心原来静止时，总动能才会全部消失。

\subsubsection{含转动时的有效质量}
考虑两个在平面内自由运动的刚体，接触冲量沿固定法线。接触点速度为
\begin{equation}
 \vct v_{Pi}=\vct V_i+\omega_i\vct k\times\vct r_i.
\end{equation}
由式\eqref{sup5:eq:d4-rigid}，物体2有 $\Delta\vct V_2=J\vct n/m_2$、$\Delta\omega_2=J(\vct r_2\times\vct n)_z/I_2$。再用 $\vct n\cdot(\vct k\times\vct r_2)=(\vct r_2\times\vct n)_z$，可得接触点沿法向的速度增量
\begin{equation}
 J\left[\frac1{m_2}+\frac{(\vct r_2\times\vct n)_z^2}{I_2}\right].
\end{equation}
物体1受相反冲量，其接触点法向速度减小。两者相减可写为
\begin{equation}\label{sup5:eq:d4-kappa}
 g_n^+=g_n^-+\kappa J,\quad
 \kappa=\frac1{m_1}+\frac1{m_2}
 +\frac{(\vct r_1\times\vct n)_z^2}{I_1}
 +\frac{(\vct r_2\times\vct n)_z^2}{I_2}.
\end{equation}
因而
\begin{equation}\label{sup5:eq:d4-general}
 J=-\frac{(1+e)g_n^-}{\kappa},\qquad
 Q=\frac{(1-e^2)(g_n^-)^2}{2\kappa}.
\end{equation}
量 $m_{\rm eff}=1/\kappa$ 是此次接触的有效质量。它取决于冲量方向和作用点，不能只凭刚体总质量确定。

这一能量关系也可沿累积冲量直接推导。记碰撞中已经传递的法向冲量为 $j$，当前相对速度为 $g_n(j)=g_n^-+\kappa j$。作用与反作用两力的总功率为 $Fg_n$，且 $\dd j=F\dd t$，故
\[
T^+-T^-=\int_0^J(g_n^-+\kappa j)\dd j
=[g_n^-j+\kappa j^2/2]_0^J.
\]
于是
\begin{equation}
 T^+-T^-=Jg_n^-+\tfrac12\kappa J^2.
\end{equation}
代入恢复关系即可得到式\eqref{sup5:eq:d4-general}。有固定铰链或可动支座时，应先用约束消去不允许的速度，再求 $\kappa$；不可照搬两个自由刚体的表达式。支座运动且对系统做功时，还须另算支座的能量交换。


\subsection{例题与模型变式}

\begin{example}[刚体弹性碰撞的法向相对速度]\label{ex:hw5-1}

两个平面刚体在$P$处发生无摩擦、完全弹性碰撞。证明：碰后接触点法向相对速度与碰前大小相等、方向相反。说明证明中为何不能直接使用质心相对速度。

\end{example}


\begin{figure}[H]\centering
\begin{tikzpicture}[>=Stealth,line cap=round,line join=round]
\draw[curve](0,0)..controls(0,.9)and(-1.5,1.1)..(-2,.3)..controls(-2.5,-1)and(0,-1.1)..(0,0);
\draw[curve](0,0)..controls(0,-1)and(1.8,-1.1)..(2.2,-.2)..controls(2.5,1.1)and(0,1.2)..(0,0);
\coordinate(P)at(0,0);\coordinate(Ca)at(-1.1,.3);\coordinate(Cb)at(1.2,-.25);
\fill(P)circle(2pt);\draw[thin](P)--(0,-1.2)node[below]{$P$};\fill(Ca)circle(1.3pt)node[below left]{$C_1$};\fill(Cb)circle(1.3pt)node[above right]{$C_2$};
\draw[guide](Ca)--(P)node[midway,above]{$\vct r_1$};\draw[guide](Cb)--(P)node[midway,below]{$\vct r_2$};\draw[axis](0,0)--(-.65,0)node[below]{$\vct n$};
\end{tikzpicture}
\caption{接触法线不通过质心时，碰撞同时改变平动与转动}\label{fig:hw5-1-q}
\end{figure}

\begin{solution}
取图\ref{fig:hw5-1-s}中由刚体2指向刚体1的单位法线$\vct n$，冲量$J\vct n$作用于刚体1，反冲量作用于刚体2。定义$g^-=(\vct v_{P1}^--\vct v_{P2}^-)\cdot\vct n<0$，有
\[
 g^+=g^-+WJ,\quad W=\frac1{m_1}+\frac1{m_2}
 +\frac{(\vct r_1\times\vct n)_z^2}{I_1}
 +\frac{(\vct r_2\times\vct n)_z^2}{I_2}>0.
\]

\begin{figure}[H]\centering
\begin{tikzpicture}[>=Stealth,line cap=round,line join=round]
\draw[curve](0,0)..controls(0,.9)and(-1.5,1.1)..(-2,.3)..controls(-2.5,-1)and(0,-1.1)..(0,0);
\draw[curve](0,0)..controls(0,-1)and(1.8,-1.1)..(2.2,-.2)..controls(2.5,1.1)and(0,1.2)..(0,0);
\coordinate(P)at(0,0);\coordinate(Ca)at(-1.1,.3);\coordinate(Cb)at(1.2,-.25);
\fill(P)circle(2pt);\draw[thin](P)--(0,-1.2)node[below]{$P$};\fill(Ca)circle(1.3pt)node[below left]{$C_1$};\fill(Cb)circle(1.3pt)node[above right]{$C_2$};
\draw[guide](Ca)--(P)node[midway,above]{$\vct r_1$};\draw[guide](Cb)--(P)node[midway,below]{$\vct r_2$};\draw[axis](0,0)--(-.9,0)node[below]{$J\vct n$};
\draw[axis](0,0)--(.9,0)node[above]{$-J\vct n$};
\end{tikzpicture}
\caption{同一对接触冲量关于两质心分别具有非零力臂}\label{fig:hw5-1-s}
\end{figure}

记$\vct r_i=\overrightarrow{C_iP}$。刚体1的变化为$\Delta\vct V_1=J\vct n/m_1$、$\Delta\omega_1=J(\vct r_1\times\vct n)_z/I_1$，刚体2则取相反冲量。代入
\[
 \Delta\vct v_{Pi}=\Delta\vct V_i+
 \Delta\omega_i\vct e_z\times\vct r_i
\]
并沿$\vct n$投影，每个转动项都产生一个非负的平方项，于是得到上式的$W$。图中的两个质心偏离接触法线，因此两个转动力臂均不为零。

令$a_i=(\vct r_i\times\vct n)_z$。分别用$\tfrac12m[(\vct V+\Delta\vct V)^2-\vct V^2]=m\vct V\cdot\Delta\vct V+\tfrac12m|\Delta\vct V|^2$展开平动动能，转动动能同样展开，得到
\[
\begin{aligned}
\Delta T_1&=J(\vct V_1^-\cdot\vct n+\omega_1^-a_1)
+\frac{J^2}{2}\left(\frac1{m_1}+\frac{a_1^2}{I_1}\right),\\
\Delta T_2&=-J(\vct V_2^-\cdot\vct n+\omega_2^-a_2)
+\frac{J^2}{2}\left(\frac1{m_2}+\frac{a_2^2}{I_2}\right).
\end{aligned}
\]
因$(\vct e_z\times\vct r_i)\cdot\vct n=a_i$，两式的线性项之和正是$Jg^-$，二次项之和是$WJ^2/2$。完全弹性要求总动能不变，故
\[
 0=\Delta T=Jg^-+\tfrac12WJ^2.
\]
排除未发生碰撞的$J=0$根，得$J=-2g^-/W$，代回即
\[
 \boxed{g^+=-g^-}.
\]
接触点偏离质心时，冲量还改变角速度，$W$中的转动项不可删去。本结论只针对法向分量；光滑碰撞没有切向冲量，但偏心冲量引起的转动仍可能改变接触点的切向速度。


\end{solution}



\begin{example}[正碰中的冲量与能量损失]\label{sup5:problem-example-1}
光滑水平面上，质量为 $m$ 的小球A以速度 $u>0$ 追上静止的质量为 $2m$ 的小球B，二者沿球心连线碰撞，恢复系数为 $e$。求两球末速度、相互作用的冲量与能量损失，并讨论A何时反弹。
\end{example}
\begin{figure}[H]\centering
\begin{tikzpicture}
\draw[guide](-1.5,-.5)--(5,-.5);
\draw[curve](0,0)circle(.5)(3,0)circle(.5);
\node at(0,0){A};\node at(3,0){B};
\node[above]at(0,.6){$m$};\node[above]at(3,.6){$2m$};
\draw[axis](.7,0)--(1.9,0)node[midway,above]{$u$};
\node[right]at(3.7,0){静止};\draw[axis](4,-1)--(5,-1)node[right]{$x$};
\end{tikzpicture}
\caption{两球沿同一水平直线正碰。}\label{sup5:fig:d4-e1-q}
\end{figure}
\begin{solution}
沿图\ref{sup5:fig:d4-e1-q}的 $x$ 方向取正。水平面支持力没有水平冲量，故两球系统满足
\begin{equation}
 mu=mv_A+2mv_B.
\end{equation}
碰前A相对B的接近速度为 $u$，碰后B相对A的分离速度为 $v_B-v_A$，所以
\begin{equation}
 v_B-v_A=eu.
\end{equation}
由第二式 $v_A=v_B-eu$，代入第一式得 $3v_B=(1+e)u$，于是
\begin{equation}
 v_A=\frac{1-2e}{3}u,\qquad v_B=\frac{1+e}{3}u.
\end{equation}
图\ref{sup5:fig:d4-e1-s}画出了冲量方向，并以 $e=1/2$ 标出碰后状态。B的冲量为
\begin{equation}
 J=2mv_B=\frac{2(1+e)}3mu,
\end{equation}
A所受冲量是 $-J$。即使 $v_A>0$，A的速度也比原来小，所受冲量仍向左。
\begin{figure}[H]\centering
\begin{tikzpicture}
\draw[curve](0,0)circle(.42)(3,0)circle(.42);\node at(0,0){A};\node at(3,0){B};
\draw[axis](-.55,0)--(-1.65,0)node[midway,above]{$-J$};
\draw[axis](3.55,0)--(4.65,0)node[midway,above]{$J$};
\node[below]at(0,-.6){$e=\tfrac12:\ v_A=0$};
\draw[axis](2.8,-1.45)--(4,-1.45)node[midway,below]{$v_B=u/2$};
\end{tikzpicture}
\caption{冲量总是使两球的相对接近速度减小；$e=1/2$ 时A恰好停止。}\label{sup5:fig:d4-e1-s}
\end{figure}
本题质心速度为 $u/3$，碰撞不能改变其平动动能。把总动能逐项分解为
\[
T^- =\frac16mu^2+\frac13mu^2,
\qquad
T^+ =\frac16mu^2+\frac13me^2u^2.
\]
每式第一项是质心平动能，第二项是相对运动能。图\ref{sup5:fig:d4-e1-s}中A停止时，B仍在运动，并携带系统的全部动量。

约化质量为 $\mu=2m/3$，由相对动能的变化得
\begin{equation}
 Q=\frac m3(1-e^2)u^2.
\end{equation}
当 $e=1/2$ 时，$Q=mu^2/4$；初动能为 $mu^2/2$，末动能为 $\tfrac12(2m)(u/2)^2=mu^2/4$，两种算法相符。A在 $e>1/2$ 时向左反弹，$e<1/2$ 时继续向右，$e=1/2$ 时停止。对所有 $e\ge0$，都有 $v_B\ge v_A$，符合碰后不再相互穿入的条件。
\end{solution}



\begin{example}[自由杆端部受撞]\label{sup5:problem-example-5}
长 $L$、质量 $M$ 的均匀细杆静止在光滑水平面上。质量 $m$ 的小球以速度 $u$ 垂直于杆，撞击杆的右端P。小球与杆光滑接触，恢复系数为 $e$。求小球、杆质心和杆的碰后运动；再取 $m=M$、$e=1$ 讨论。
\end{example}
\begin{figure}[H]\centering
\begin{tikzpicture}
\draw[line width=3pt](-2,0)--(2,0);\fill(0,0)circle(2pt);\node[above]at(0,.1){C};\node[above]at(2,.1){P};
\draw[<->](-2,.7)--(2,.7)node[midway,above]{$L$};
\draw[curve](2,-1.65)circle(.25);\node[right]at(2.3,-1.65){$m$};
\draw[axis](2,-1.3)--(2,-.35)node[midway,right]{$u$};
\node[left]at(-2.3,0){$M$};
\end{tikzpicture}
\caption{小球从下方撞击杆端P，碰撞法线竖直向上。}\label{sup5:fig:d4-e5-q}
\end{figure}
\begin{solution}
图\ref{sup5:fig:d4-e5-q}中取向上为速度正方向，逆时针为角速度正方向。设杆受到向上的冲量 $J$，杆质心速度为 $V_C$，角速度为 $\omega$，球的末速度为 $v$。分别写出冲量方程：
\begin{equation}
 v=u-J/m,\qquad MV_C=J,\qquad
 \frac{ML^2}{12}\omega=\frac L2J.
\end{equation}
于是 $\omega=6J/(ML)$。杆端P的向上速度由平动和转动叠加，图\ref{sup5:fig:d4-e5-s}给出
\begin{equation}
 v_P=V_C+\omega L/2=\frac{4J}{M}.
\end{equation}
\begin{figure}[H]\centering
\begin{tikzpicture}
\draw[line width=2.5pt](-2,0)--(2,0);\node[below]at(0,-.1){C};\node[below]at(2,-.1){P};
\draw[axis](0,0)--(0,.65)node[above left]{$V_C=J/M$};
\draw[axis](2,0)--(2,2.6)node[above]{$v_P=4J/M$};
\draw[axis](-.7,-.7)arc(225:315:.98995);\node[below]at(0,-1.1){$\omega>0$};
\draw[curve](3.8,-.8)circle(.2);\draw[axis](3.8,-.5)--(3.8,.475)node[above]{$v$};
\end{tikzpicture}
\caption{杆端速度包含质心平动与绕质心转动两部分。小球箭头按 $m=M,e=1$ 的特例绘制。}\label{sup5:fig:d4-e5-s}
\end{figure}
恢复关系为 $v_P-v=eu$，因而
\begin{equation}
 J\left(\frac4M+\frac1m\right)=(1+e)u,
 \qquad J=\frac{(1+e)mM}{M+4m}u.
\end{equation}
所以
\begin{equation}
 v=\frac{4m-eM}{M+4m}u,\quad
 V_C=\frac{(1+e)m}{M+4m}u,\quad
 \omega=\frac{6(1+e)m}{M+4m}\frac uL.
\end{equation}
杆端的冲量响应系数为 $4/M$，相应的有效质量为 $M/4$。冲量同时引起杆的平动和转动，接触点的速度变化由这两部分共同决定。

取 $m=M,e=1$，可得 $J=2Mu/5$，$v=3u/5$，$V_C=2u/5$，$\omega=12u/(5L)$。此时P点速度为 $8u/5$，大于入射球速率，但 $v_P-v=u$，恢复关系仍成立。杆的动能为
\begin{equation}
 T_{\rm rod}=\frac{J^2}{2M}+\frac{(JL/2)^2}{2(ML^2/12)}
 =\frac{2J^2}{M}.
\end{equation}
上述特例中，球和杆的末动能分别为 $9Mu^2/50$、$8Mu^2/25$，总和仍为 $Mu^2/2$。一般情况下，损失为 $(1-e^2)u^2/[2(1/m+4/M)]$。
小球是否反弹，应从已求得的带符号速度判断：$4m<eM$ 时才有 $v<0$。即使碰撞完全弹性，等质量的小球和杆也不会交换质心速度，因为接触点还包含转动速度。

图\ref{sup5:fig:d4-e5-s}中杆所获动能占入射动能的比例为
\[
\frac{T_{\rm rod}}{mu^2/2}
=\frac{4(1+e)^2mM}{(M+4m)^2}.
\]
在 $e=1$ 时，只有 $M=4m$ 才使该比值达到1，此时球恰好停止。对应的接触点有效质量为 $M/4=m$，与两相等有效质量的弹性碰撞相符。

\end{solution}



\begin{example}[铰接杆与打击中心]\label{sup5:problem-example-6}
均匀细杆质量为 $M$、长为 $L$，上端O用固定光滑铰链悬挂，初始竖直静止。小球质量为 $m$，以水平速度 $u$ 撞击距O为 $a$ 的点P，$0<a\le L$。接触光滑，恢复系数为 $e$。求小球末速度、杆角速度和铰链冲量，并求使水平铰链冲量为零的打击位置。
\end{example}
\begin{figure}[H]\centering
\begin{tikzpicture}
\draw[line width=3pt](1,2.5)--(1,-1);\draw[curve](1,2.5)circle(.12);\node[right]at(1.2,2.5){O};
\fill(1,.4)circle(2pt);\node[right]at(1.15,.4){P};\node[right]at(1.15,-.95){$L$};
\draw[<->](2.1,2.5)--(2.1,.4)node[midway,right]{$a$};
\draw[curve](-1,.4)circle(.22);\draw[axis](-.7,.4)--(.6,.4)node[midway,above]{$u$};
\node[below]at(-1,.05){$m$};\node[left]at(.8,1.5){$M$};
\end{tikzpicture}
\caption{杆只绕固定点O转动，冲量作用点为P。}\label{sup5:fig:d4-e6-q}
\end{figure}
\begin{solution}
在图\ref{sup5:fig:d4-e6-q}中，向右为速度和冲量的正方向，杆向右摆起的角速度记为正。杆关于O的转动惯量为 $I_O=ML^2/3$。铰链冲量对O无力矩，所以
\begin{equation}
 I_O\omega=Ja,\qquad v_P=a\omega=\frac{Ja^2}{I_O}.
\end{equation}
小球有 $v=u-J/m$。将接触点P的速度代入恢复关系 $v_P-v=eu$，得
\begin{equation}
 J=\frac{(1+e)u}{1/m+a^2/I_O},\quad
 \omega=\frac{(1+e)mau}{I_O+ma^2},\quad
 v=\frac{ma^2-eI_O}{I_O+ma^2}u.
\end{equation}
杆质心C距O为 $L/2$，故 $V_C=\omega L/2$。图\ref{sup5:fig:d4-e6-s}中杆的线动量方程给出水平铰链冲量
\begin{equation}
 H_x+J=MV_C,\qquad
 H_x=J\left(\frac{3a}{2L}-1\right).
\end{equation}
碰撞前后C的瞬时速度均无竖直分量，忽略重力冲量后有 $H_y=0$。这不表示碰撞以后的竖直铰链力为零。
\begin{figure}[H]\centering
\begin{tikzpicture}
\draw[line width=2.5pt](0,2.5)--(0,-2.5);\node[right]at(.15,2.65){O};
\fill(0,0)circle(2pt);\node[left]at(-.15,0){C};
\draw[axis](0,2.5)--(-.3,2.5);\node[above left]at(-.35,2.5){$H_x$};
\fill(0,-.5)circle(1.5pt);\node[left]at(-.15,-.5){P};
\draw[axis](0,-.5)--(3,-.5);\node[above]at(2.6,-.5){$J$};
\draw[axis](0,0)--(1,0)node[right]{$V_C$};
\draw[guide](0,-.833333)--(3.8,-.833333);
\draw[<->](3.8,2.5)--(3.8,-.833333)node[midway,right]{$a_*=2L/3$};
\end{tikzpicture}
\caption{铰接杆的冲量与质心速度。图示 $L/2<a<2L/3$，P位于C与打击中心之间，铰链冲量向左。}\label{sup5:fig:d4-e6-s}
\end{figure}
令 $H_x=0$，得到 $a_*=2L/3$。在这一位置打击，P处冲量恰好等于杆的线动量增量，铰链无需再提供水平冲量。$a<a_*$ 时铰链冲量向左，$a>a_*$ 时向右。

损失为
\begin{equation}
 Q=\frac{(1-e^2)u^2}{2(1/m+a^2/I_O)}.
\end{equation}
因铰点固定，铰链冲量不做功，$e=1$ 时球与杆的总动能守恒。小球反弹的条件是 $ma^2<eI_O$，由此也能看出，靠近铰点打击时杆的接触点难以获得较大速度。

在打击中心，杆沿P点法线的有效质量为
\[
\frac{I_O}{a_*^2}
=\frac{ML^2/3}{(2L/3)^2}=\frac34M.
\]
代回冲量式时，应使用这一数值。铰链水平冲量为零，不表示杆没有转动，也不表示小球可以按碰上总质量为 $M$ 的自由质点处理。

\end{solution}



\section{斜碰、速度空间与允许末态}

守恒量限定末速度的几何轨迹，接近与分离条件决定物理分支。接触法线和速度方向须在同一图中标明，再将相应分量代入冲量方程。



\subsection{斜碰中的几何关系}
两光滑球碰撞时，接触法线沿两球球心连线。将相对速度分成法向和切向分量，切向分量保持不变，法向分量按式\eqref{sup5:eq:d4-e}改变。若球在接触前无自转，沿球心连线的冲量也不会使其自转。

设质量 $m$ 的入射球以速度 $u\vct e_x$ 撞击静止的质量 $M$ 的球，半径和为 $d$，瞄准距离为 $b$。图\ref{sup5:fig:d4-t3}中 $b>0$，故
\begin{equation}
 \beta=\frac b d,\quad c=\sqrt{1-\beta^2},\quad
 \vct n=c\vct e_x-\beta\vct e_y,\quad 0\le\beta<1.
\end{equation}
入射球位于靶球左上方，$g_n^-=-uc$，冲量和速度为
\begin{align}\label{sup5:eq:d4-oblique}
 J&=(1+e)\frac{mM}{m+M}uc,\\
 \vct v_1&=u\vct e_x-\frac Jm\vct n,
 &\vct v_2&=\frac JM\vct n.
\end{align}

\begin{figure}[H]\centering
\begin{tikzpicture}
\coordinate(A)at(0,1);\coordinate(B)at(1.732,0);
\draw[curve](A)circle(1)(B)circle(1);
\fill(A)circle(1.3pt)(B)circle(1.3pt);
\node[above left]at(A){$m$};\node[below right]at(B){$M$};
\draw[guide](-2.4,0)--(3.4,0);\draw[guide](-2.4,1)--(-1,1);
\draw[<->](-2.1,0)--(-2.1,1)node[midway,left]{$b$};
\draw[axis](-2.2,2.4)--(-.6,2.4)node[midway,above]{$u$};
\draw[guide](A)--(B);
\draw[guide](A)--(-.7,-.212)(B)--(1.032,-1.212);
\draw[<->](-.7,-.212)--(1.032,-1.212)node[midway,below left]{$d$};
\draw[axis](1.6,1.8)--(2.466,1.3)node[midway,above right]{$\vct n$};
\draw[axis](3.4,0)--(4.4,0)node[right]{$x$};
\end{tikzpicture}
\caption{瞄准距离和接触法线，图示 $b=d/2$。接触时两球球心距离为 $d$。}\label{sup5:fig:d4-t3}
\end{figure}

对于弹性碰撞，在质心系中两球速度的大小分别不变。回到实验室系时，必须给速度矢量加上同一个质心速度。入射球的末速度端点在速度平面内满足
\begin{equation}\label{sup5:eq:d4-circle}
 |\vct v_1-\vct V|=\frac{M}{m+M}u,\qquad
 \vct V=\frac{m}{m+M}u\vct e_x.
\end{equation}
若 $m>M$，原点在此圆外，最大偏转角由原点向圆的切线确定；若 $m<M$，圆包围原点，入射球可向后反弹。等质量弹性斜碰且两球末速度均不为零时，由动量和能量关系可得 $\vct v_1\cdot\vct v_2=0$。


\subsection{例题与模型变式}

\begin{example}[两个自由球的速度空间]\label{ex:hw5-6}

球$A$质量$3m$、初速度$(v_0,2v_0,0)$，球$B$质量$m$、初速度$(v_0,0,0)$，两球完全弹性碰撞，碰撞方向可变。(1)求$A$碰后可能速率范围；(2)求其速度方向可能分布的立体角。

\end{example}


\begin{figure}[H]\centering
\begin{tikzpicture}[>=Stealth,line cap=round,line join=round]
\draw[axis](0,0)--(4,0)node[right]{$v_x$};\draw[axis](0,0)--(0,3)node[above]{$v_y$};\fill(1,1.5)circle(2pt)node[right]{$\vct V$};\draw[curve](1,1.5)circle(.5);\draw[guide](0,0)--(1,1.5);
\end{tikzpicture}
\caption{两自由球碰撞后的速度球截面}\label{fig:hw5-6-q}
\end{figure}

\begin{solution}
质心速度为$\vct V=(v_0,3v_0/2,0)$。在质心系中，$A$初速度为$(0,v_0/2,0)$，弹性碰撞后大小仍为$v_0/2$。图\ref{fig:hw5-6-s}是速度空间球面的一个截面：球心为$\vct V$，半径为$v_0/2$。

\begin{figure}[H]\centering
\begin{tikzpicture}[>=Stealth,line cap=round,line join=round]
\draw[axis](0,0)--(4,0)node[right]{$v_x$};\draw[axis](0,0)--(0,3)node[above]{$v_y$};\fill(1,1.5)circle(2pt)node[right]{$\vct V$};\draw[curve](1,1.5)circle(.5);\draw[guide](0,0)--(1,1.5);\draw[guide](0,0)--(.523372891,1.651084740);\draw[guide](0,0)--(1.322780956,1.118146030);\node at(2.8,2.7){$R=v_0/2$};\node at(2.8,2.2){$d=\sqrt{13}v_0/2$};
\draw[guide](1,1.5)--(.523372891,1.651084740);
\draw(.61017,.915255)arc(56.30993:72.4105:1.1);\node at(.49,1.01){$\Theta$};\node[below left]at(0,0){$O$};
\end{tikzpicture}
\caption{原点向速度球的两条切线限定方向范围}\label{fig:hw5-6-s}
\end{figure}

球心距原点$d=\sqrt{13}v_0/2$，因此
\[
 \boxed{\frac{\sqrt{13}-1}{2}v_0\le|\vct v_A'|
 \le\frac{\sqrt{13}+1}{2}v_0}.
\]
由原点向速度球面作切线，半锥角满足$\sin\Theta=R/d=1/\sqrt{13}$。以$O\vct V$为极轴，图\ref{fig:review5-cap}中允许的方向为$0\le\vartheta\le\Theta$、$0\le\varphi<2\pi$。因此立体角为
\[
\begin{aligned}
\Omega&=\int_0^{2\pi}\dd\varphi\int_0^\Theta\sin\vartheta\,\dd\vartheta
=\int_0^{2\pi}\bigl[-\cos\vartheta\bigr]_0^\Theta\,\dd\varphi\\
&=(1-\cos\Theta)\bigl[\varphi\bigr]_0^{2\pi}
=2\pi(1-\cos\Theta).
\end{aligned}
\]
将$\cos\Theta=\sqrt{1-1/13}$代入，得到
\[
 \boxed{\Omega=2\pi(1-\cos\Theta)
 =2\pi\left(1-\sqrt{\frac{12}{13}}\right)}.
\]
此处不限定在某一平面内碰撞，才使用立体角。若运动限制在平面内，可能方向改为一个平面角区间$2\Theta$。


下面用碰撞运动学确定速度球面的可达性。两球的初始相对速度为$\vct g=2v_0\vct e_y$。对光滑弹性球，选择接触法线$\vct n$后有$\vct g'=\vct g-2(\vct g\cdot\vct n)\vct n$，从而$|\vct g'|=2v_0$。让碰撞参数和接触法线变化，可使相对速度方向遍及球面；无偏转方向作为擦边极限包含在内。又
\[
 \vct v_A'=\vct V+\frac{m}{3m+m}\vct g',
\]
因此速度端点遍历半径$v_0/2$的速度球面，而不是球内体积。

沿轴方向的两交点分别给出最小、最大速率，切锥边界则给出允许方向的边界；速率极值点与允许方向的切点是速度球上的不同位置。
\begin{figure}[H]\centering
\begin{tikzpicture}[>=Stealth,line cap=round,line join=round]
\coordinate(O)at(0,0);\coordinate(V)at(3,0);\draw[curve](V)circle(.9);\draw[axis](O)--(4.4,0)node[right]{$O\vct V$};\fill(O)circle(1.5pt)node[below left]{$O$};\fill(V)circle(1.5pt)node[below]{$\vct V$};
\draw(O)--(2.73,.8585);\draw(O)--(2.73,-.8585);\draw[guide](V)--(2.73,.8585)node[midway,right]{$R$};\draw(.85,0)arc(0:17.4576:.85);\node at(1.12,.16){$\Theta$};
\draw[guide](2.73,-.8585)arc(-90:90:.25 and .8585);\node[align=left]at(2,-1.65){绕极轴转过全部方位角\\每条允许射线穿过速度球};
\end{tikzpicture}
\caption{速度球的切锥；空间方向由锥内射线给出}\label{fig:review5-cap}
\end{figure}


\end{solution}



\begin{example}[三个等质量球的守恒约束]\label{ex:hw5-7}

三个质量$m$的球在光滑水平面上碰撞，只有$A$初速度为$(v_0,0)$，其余静止，碰撞完全弹性，不限定碰撞次序及具体几何。(1)求守恒关系允许的$A$末速度区域；(2)求$A$反向最大速率；(3)求$A$偏转$90^\circ$时的最大速率；(4)已知$\vct v_A'=(2v_0/3,v_0/3)$，求$B$相对原方向的最大偏转角。

\end{example}


\begin{figure}[H]\centering
\begin{tikzpicture}[>=Stealth,line cap=round,line join=round]
\begin{scope}[scale=2]\draw[axis](0,0)--(1.5,0)node[right]{$v_x/v_0$};\draw[axis](0,0)--(0,1)node[above]{$v_y/v_0$};\draw[curve](.333333,0)circle(.666667);\fill(.333333,0)circle(1pt)node[below]{$1/3$};\node at(.5,1.5){$\vct v_A$的允许区域};\end{scope}
\end{tikzpicture}
\caption{守恒关系允许的首球末速度圆盘}\label{fig:hw5-7-q}
\end{figure}

\begin{solution}
令$\vct s=(v_0,0)-\vct v_A'$，则$\vct v_B'+\vct v_C'=\vct s$。对给定$\vct s$，后两球动能最小时两速度相同，所以守恒的充要代数条件为
\[
 |\vct v_A'|^2+\tfrac12|\vct s|^2\le v_0^2.
\]

\begin{figure}[H]\centering
\begin{tikzpicture}[>=Stealth,line cap=round,line join=round]
\begin{scope}[scale=2.25]
\fill[gray!10](.333333,0)circle(.666667);\draw[curve](.333333,0)circle(.666667);
\draw[axis](-.5,0)--(1.25,0)node[right]{$v_{Ax}'/v_0$};\draw[axis](0,-.8)--(0,.95)node[above]{$v_{Ay}'/v_0$};
\fill(-.333333,0)circle(.7pt)node[below]{$-1/3$};\fill(0,.57735)circle(.7pt)node[right]{$1/\sqrt3$};
\fill(.333333,0)circle(.6pt)node[below]{$1/3$};\end{scope}
\begin{scope}[shift={(6.3,0)},scale=2.5]
\draw[curve](.166667,-.166667)circle(.40824829);
\draw[axis](-.45,0)--(.9,0)node[right]{$v_{Bx}'/v_0$};\draw[axis](0,-.7)--(0,.7)node[above]{$v_{By}'/v_0$};
\fill(.166667,-.166667)circle(.6pt)node[right]{$\vct V_{BC}$};\fill(-.206011,0)circle(.8pt)node[below left]{$B_-$};
\node[above right]at(0,0){$O$};\end{scope}
\end{tikzpicture}
\caption{左图为守恒允许圆盘，右图速度圆上的负轴点实现完全反向}\label{fig:hw5-7-s}
\end{figure}

充分性也可构造说明：令$\vct v_{B,C}'=\vct s/2\pm\vct w$，则
\[
 |\vct w|^2=\frac12\left(v_0^2-|\vct v_A'|^2-\frac12|\vct s|^2\right).
\]
右端非负时，在平面内任选具有该大小的$\vct w$，就能同时满足动量和能量。这一构造满足守恒方程；能否由指定的碰撞路径产生，还要检验接触条件。

令$\vct v_A'=(x,y)$，把$\vct s=(v_0-x,-y)$代入约束并乘以2，有
\[
2(x^2+y^2)+(v_0-x)^2+y^2\le2v_0^2,
\]
即$3x^2+3y^2-2v_0x\le v_0^2$。除以3并对$x$配方，得到图\ref{fig:hw5-7-s}左侧的圆盘：
\[
 \boxed{(v_{Ax}'-v_0/3)^2+v_{Ay}'^2\le4v_0^2/9}.
\]
反向取$x$轴负半轴，最大速率为$\boxed{v_0/3}$；偏转$90^\circ$令$v_{Ax}'=0$，最大速率为$\boxed{v_0/\sqrt3}$。

给定$A$末速度后，$|\vct v_A'|^2=5v_0^2/9$、$\vct s=(v_0/3,-v_0/3)$，故$BC$质心速度为$\vct V_{BC}=\vct s/2=(v_0/6,-v_0/6)$。写$\vct v_{B,C}'=\vct V_{BC}\pm\vct w$，代入前述能量关系，
\[
|\vct w|^2=\frac12\left(v_0^2-\frac59v_0^2-\frac19v_0^2\right)
=\frac{v_0^2}{6}.
\]
圆心到原点的距离为$v_0/\sqrt{18}$，小于半径$v_0/\sqrt6$，所以速度圆包围原点，$B$可沿任意方向。要使$B$完全反向，取$v_{By}'=0$，即$w_y=v_0/6$；再取负的水平分量
$w_x=-\sqrt{v_0^2/6-v_0^2/36}=-\sqrt5v_0/6$，于是
\[
 v_{Bx}'=\frac{1-\sqrt5}{6}v_0<0,
 \qquad\boxed{\text{最大偏转角}=180^\circ}.
\]
这些是未限定碰撞构形时的守恒允许范围。若另给固定的入射位置、碰撞法线或次序，还须将相应接触条件与此区域取交集。


\end{solution}



\begin{example}[等质量光滑球的斜碰]\label{sup5:problem-example-2}
两个相同光滑球的质量均为 $m$，半径均为 $R$。A沿水平直线以速度 $u$ 入射，B原来静止。瞄准距离 $b=R$，碰撞完全弹性。求两球末速度的分量、大小和夹角。
\end{example}
\begin{figure}[H]\centering
\begin{tikzpicture}[scale=.9]
\draw[curve](0,1)circle(1)(1.732,0)circle(1);
\fill(0,1)circle(1.3pt)(1.732,0)circle(1.3pt);
\node[above left]at(0,1){A};\node[below right]at(1.732,0){B};
\draw[guide](0,1)--(1.732,0);
\draw[guide](0,1)--(-.7,-.212)(1.732,0)--(1.032,-1.212);
\draw[<->](-.7,-.212)--(1.032,-1.212)node[midway,below left]{$2R$};
\draw[guide](-2.4,0)--(3,0);\draw[guide](-2.4,1)--(-1,1);
\draw[<->](-2.1,0)--(-2.1,1)node[midway,left]{$R$};
\draw[axis](-2.3,2.25)--(-.7,2.25)node[midway,above]{$u$};
\end{tikzpicture}
\caption{接触时两球球心距离为 $2R$，球心连线与水平方向夹角为 $30^\circ$。}\label{sup5:fig:d4-e2-q}
\end{figure}
\begin{solution}
图\ref{sup5:fig:d4-e2-q}中，A在左上方，B在右下方。法线从A指向B，因此
\begin{equation}
 \vct n=\frac{\sqrt3}{2}\vct e_x-\frac12\vct e_y.
\end{equation}
碰前A沿法线的速度分量为 $u\sqrt3/2$，B为零。等质量弹性正碰会交换这两个法向分量；切向分量没有改变。也可直接由恢复关系求得
\begin{equation}
 J=2\frac m2\frac{\sqrt3}2u=\frac{\sqrt3}2mu.
\end{equation}
由图\ref{sup5:fig:d4-e2-s}的矢量分解，
\begin{align}
 \vct v_A&=u\vct e_x-\frac Jm\vct n
 =\frac u4\vct e_x+\frac{\sqrt3u}4\vct e_y,\\
 \vct v_B&=\frac Jm\vct n
 =\frac{3u}4\vct e_x-\frac{\sqrt3u}4\vct e_y.
\end{align}
\begin{figure}[H]\centering
\begin{tikzpicture}
\coordinate(O)at(0,0);\coordinate(A)at(1,1.732);\coordinate(U)at(4,0);
\draw[axis](O)--(A)node[midway,left]{$\vct v_A$};
\draw[axis](A)--(U)node[midway,above right]{$\vct v_B$};
\draw[axis](O)--(U)node[midway,below]{$u\vct e_x$};
\draw[guide](A)--(1,0);\node[below]at(1,0){$u/4$};
\draw (A)++(-.15,-.26)--++(.26,-.15)--++(.15,.26);
\node[right]at(4.2,1.2){$|\vct v_A|=u/2$};
\node[right]at(4.2,.5){$|\vct v_B|=\sqrt3u/2$};
\end{tikzpicture}
\caption{末速度首尾相接等于初速度；直角位于两条末速度的方向之间。}\label{sup5:fig:d4-e2-s}
\end{figure}
沿接触法线垂直方向取单位矢量 $\vct t=(1/2,\sqrt3/2)$。图\ref{sup5:fig:d4-e2-q}中，入射速度可分成
\[
u\vct e_x=\frac{\sqrt3u}{2}\vct n+\frac u2\vct t.
\]
法向分量在等质量弹性碰撞中传给B，切向分量仍留在A。因此 $\vct v_A=(u/2)\vct t$、$\vct v_B=(\sqrt3u/2)\vct n$，也就得到图\ref{sup5:fig:d4-e2-s}中的两个末速度。这种分解解释了两方向为何垂直。

由分量平方和可得 $|\vct v_A|=u/2$、$|\vct v_B|=\sqrt3u/2$。A向水平线上方偏转 $60^\circ$，B向下偏转 $30^\circ$，两末速度垂直。直接计算内积也有
\begin{equation}
 \vct v_A\cdot\vct v_B
 =\frac{3u^2}{16}-\frac{3u^2}{16}=0
\end{equation}
末动能为 $\tfrac12m(u^2/4+3u^2/4)=mu^2/2$。若 $b=0$，A停止而B沿原方向运动，零速度没有方向，不能再给两个末速度指定夹角。
\end{solution}



\begin{example}[重球的最大偏转角]\label{deep5:problem-example-3}
质量为 $m$ 的光滑球以速度 $u$ 撞击静止的质量为 $M$ 的光滑球，$m>M$，碰撞完全弹性。两球半径和为 $d$。求入射球的最大偏转角及达到该角度时的瞄准距离；给出 $m=2M$ 的结果。
\end{example}
\begin{figure}[H]\centering
\begin{tikzpicture}
\draw[curve](0,.9)circle(.7)(2,0)circle(.5);
\draw[guide](-1.6,.9)--(0,.9);\draw[guide](-1.6,0)--(3,0);
\draw[axis](-1.6,1.9)--(-.2,1.9)node[midway,above]{$u$};
\draw[<->](-1.3,0)--(-1.3,.9)node[midway,left]{$b$};
\node at(0,.9){$m$};\node at(2,0){$M$};
\node[right]at(3.1,.8){$m>M$};
\end{tikzpicture}
\caption{改变入射直线的偏置 $b$，比较入射球的碰后方向。图示碰前尚未接触。}\label{deep5:fig:d4-e3-q}
\end{figure}
\begin{solution}
图\ref{deep5:fig:d4-e3-q}中靶球静止，系统质心速度为 $V=mu/(m+M)$。在质心系中，入射球碰前速率为 $u-V=Mu/(m+M)$，弹性碰撞保持这一速率。因此实验室系中的末速度端点落在以 $C=(V,0)$ 为圆心、以 $r=Mu/(m+M)$ 为半径的圆上。

图\ref{deep5:fig:d4-e3-s}中，末速度矢量从速度原点O出发。因 $V>r$，偏转角的最大值由切点T确定，且 $CT\perp OT$。直角三角形OCT给出
\begin{equation}
 \sin\theta_{\max}=\frac rV=\frac Mm,\qquad
 v_{1,\rm tan}=\sqrt{V^2-r^2}
 =u\sqrt{\frac{m-M}{m+M}}.
\end{equation}
\begin{figure}[H]\centering
\begin{tikzpicture}[scale=.8]
\draw[axis](-.3,0)--(7,0)node[right]{$v_x$};
\draw[axis](0,-2.5)--(0,2.8)node[above]{$v_y$};
\draw[curve](4,0)circle(2);\fill(4,0)circle(1.5pt)(3,1.732)circle(1.5pt);
\draw[axis](0,0)--(3,1.732)node[midway,above left]{$\vct v_1$};
\draw[guide](4,0)--(3,1.732)node[midway,right]{$r$};
\node[below left]at(0,0){O};\node[below]at(4,0){C};\node[above]at(3,1.732){T};
\draw (1,0)arc(0:30:1);\node at(1.35,.35){$\theta_{\max}$};
\draw(3,1.732)++(-.26,-.15)--++(.15,-.26)--++(.26,.15);
\node[below]at(2,-.65){$OC=V$};
\end{tikzpicture}
\caption{速度圆与切线；图中按 $m=2M$ 绘制，圆心到原点的距离为半径的两倍。}\label{deep5:fig:d4-e3-s}
\end{figure}
切点对应的瞄准距离可由接触几何求出。令 $c=\sqrt{1-b^2/d^2}$，式\eqref{sup5:eq:d4-oblique}在 $e=1$ 时给出
\begin{equation}
 \frac{v_{1x}}u=1-\frac{2M}{m+M}c^2,\qquad
 \frac{v_{1y}}u=\frac{2M}{m+M}c\sqrt{1-c^2}.
\end{equation}
由速度分量先算出
\[
|\vct v_1|^2=u^2\left[1-\frac{4mM}{(m+M)^2}c^2\right],
\qquad
Vv_{1x}=\frac{mu^2}{m+M}\left[1-\frac{2M}{m+M}c^2\right].
\]
两式相减为 $Mu^2[1-2mc^2/(m+M)]/(m+M)$。图\ref{deep5:fig:d4-e3-s}中的切线条件是 $CT\perp OT$，其标量形式恰为 $|\vct v_1|^2-Vv_{1x}=0$。

切点满足 $\vct v_1\cdot(\vct v_1-V\vct e_x)=0$。将上述分量代入，化简得
\begin{equation}
 c^2=\frac{m+M}{2m},\qquad
 \frac b d=\sqrt{\frac{m-M}{2m}}.
\end{equation}
当 $m=2M$ 时，$\theta_{\max}=30^\circ$，$b=d/2$，末速率为 $u/\sqrt3$。极值在斜碰中取得，对应非零的瞄准距离和接触冲量。$m\to M^+$ 时切点趋近速度原点，极限偏角趋于 $90^\circ$，但原点本身没有速度方向。
\end{solution}



\section{多约束、摩擦与连续碰撞}

支座、绳和粗糙接触可以同时传递冲量，相继碰撞还会改变约束状态。各阶段分别联立冲量关系与有效的接触条件，碰撞结束后的运动另用动力学方程求解。



\subsection{约束、摩擦与连续碰撞}
碰撞后的速度必须同时满足恢复关系和运动约束。铰接杆上的接触点只能沿圆周切线运动，轨道上的小珠只能沿轨道切线运动。约束冲量先由这些条件求出，再判断相应方向能否使用动量守恒。

粗糙接触还可能产生切向冲量。只给法向恢复系数不足以确定转动和切向速度；还需要摩擦模型。第\ref{sup5:sec:d4-rough}节采用刚性接触下的库仑冲量模型：接触点持续滑动时，切向冲量的增量与法向冲量增量成 $\mu$ 倍；一旦切向相对速度降至零，若维持零滑动无需额外切向冲量，则不再增加切向冲量。这个模型只描述短时速度跃变，不把碰后瞬时的零切向速度解释为持续滚动。

\begin{figure}[H]\centering
\begin{tikzpicture}
\draw[axis](0,0)--(6.4,0)node[right]{$t$};
\foreach \x/\txt in {1/{碰撞1},3.2/{碰撞2},5.4/{碰撞3}}{
\draw[curve](\x,-.12)--(\x,.12);\node[above]at(\x,.25){\txt};}
\node[below,text width=2cm,align=center]at(2.1,-.2){自由运动\\更新位置};
\node[below,text width=2cm,align=center]at(4.3,-.2){自由运动\\再求接触时刻};
\end{tikzpicture}
\caption{相互分离的多次碰撞逐次处理，每次都以此前的末状态为初状态。}\label{sup5:fig:d4-t4}
\end{figure}

图\ref{sup5:fig:d4-t4}的每次接触之间均有自由运动阶段。求出某次末速度后，应重新计算下一次接触的位置和时刻；若接触点已在分离，就不能继续套用碰撞公式。对于同时存在多个接触点的系统，须给出各接触的冲量关系和相容条件。若有限形变过程影响接触先后，还需要具体的接触模型，不能任意指定二体碰撞的计算顺序。

以下例题先求碰撞造成的速度变化，再处理随后的运动。每题从冲量方程和接触点速度列式，所得能量损失可再用式\eqref{sup5:eq:d4-general}核算。


\subsection{多约束碰撞的联立方法}
\label{sec:complete5}
设碰前构形已知，用独立广义速度$\dot{\vct q}$描述系统。碰撞时间内构形不变，将动能写成$T=\dot{\vct q}^{\mathsf T}M\dot{\vct q}/2$。第$a$个接触的法向相对速度可写成$g_a=\vct c_a^{\mathsf T}\dot{\vct q}$，对应压缩冲量为$J_a\ge0$，于是
\[
 M(\dot{\vct q}^+-\dot{\vct q}^-)=\sum_a\vct c_aJ_a,
 \qquad
 g_a^+=g_a^-+\sum_b W_{ab}J_b,
\]
\[
 W_{ab}=\vct c_a^{\mathsf T}M^{-1}\vct c_b.
\]
向量$\vct c_a$是接触间隙对坐标的梯度，其正向与法向相对速度的定义一致。$W$对称，其非对角元表示某处冲量对另一处接触速度的影响。单接触时，它退化为前文有效质量的倒数。

杆的连接方式决定允许的独立速度，碰撞矩阵也随之改变。例\ref{ex:hw5-3}中把铰接模型当成整根刚体，会额外强加不应存在的角速度约束。例\ref{ex:hw5-4}则须在杆长约束之外保留地面接触条件，并在求解后检查地面冲量能否由向上的支持作用提供。

若所有列入的接触均按给定恢复关系闭合，可联立$g_a^+=-e_ag_a^-$求冲量；求得负$J_a$表明该组接触假设不成立。将这一接触撤去后，应核对其碰后间隙速度是否非负。多个同时接触的恢复关系还要与能量及具体接触模型相容，逐对套用二体弹性碰撞公式并任意规定顺序，通常不能确定同一次多点碰撞。

\subsubsection{能量约束与可达运动}
以各接触向量$\vct c_a$为列组成矩阵$C$，则$\vct g^-=C^{\mathsf T}\dot{\vct q}^-$、$W=C^{\mathsf T}M^{-1}C$。碰撞期间构形不变，质量矩阵$M$也不变。给定冲量向量$\vct J$后，广义速度增量为$\Delta\dot{\vct q}=M^{-1}C\vct J$。将$\dot{\vct q}^+=\dot{\vct q}^-+\Delta\dot{\vct q}$代入动能，展开二次型：
\begin{align*}
 T^+-T^-&=(\dot{\vct q}^-)^{\mathsf T}M\Delta\dot{\vct q}
 +\frac12\Delta\dot{\vct q}^{\mathsf T}M\Delta\dot{\vct q}\\
 &=\vct J^{\mathsf T}C^{\mathsf T}\dot{\vct q}^-
 +\frac12\vct J^{\mathsf T}C^{\mathsf T}M^{-1}C\vct J\\
 &=\vct J^{\mathsf T}\vct g^-+\frac12\vct J^{\mathsf T}W\vct J.
\end{align*}
将联立求出的冲量代入此式，可直接计算动能变化。它还说明，$e=0$只使法向分离速度为零，永久黏合需要额外的连接条件；库仑摩擦则需要独立的切向冲量模型。

对于例\ref{ex:hw5-6}和例\ref{ex:hw5-7}的速度空间问题，动量、能量守恒给出可能末态的代数边界。具体球面接触还要求碰撞法线、相对接近速度以及碰撞次序能够实现这些末态。因此，守恒关系允许的区域与指定几何下的实际可达集合应分别表述。自由碰撞中“相对速率不变”的结论适用于质心系；回到实验室系必须把同一个质心速度加回每个物体。

\subsubsection{粗糙碰撞的切向冲量}\label{sup5:sec:d4-rough}
对撞击固定水平面的轴对称圆盘，取切向接触速度为$g_t$。法向冲量由恢复系数给出；切向冲量则须同时考虑转动。单位切向冲量改变$g_t$的系数为$W_t=1/m+R^2/I_C$，故消除初始滑动所需冲量为$J_t^*=-g_t^-/W_t$。在法、切向运动解耦、滑动过程中摩擦方向不反转的库仑冲量模型下，若$|J_t^*|\le\mu J_n$，切向冲量取$J_t^*$；否则取$J_t=-\mu J_n\,\operatorname{sgn}g_t^-$。

上述规则给出碰后瞬间的速度。若随后继续接触地面，还须重新求支持力和静摩擦，才能判断是否持续滚动。一般非圆刚体的法、切向速度变化可能耦合，不能直接套用圆盘的两条独立冲量公式。



\LectureNeedspace{4\baselineskip}


\subsection{例题与模型变式}

\begin{example}[三球刚杆撞墙]\label{ex:hw5-2}

三个质量均为$m$的小球$A,B,C$等距固定在长$2l$的轻杆上，$B$为中点。系统无初始转动，以速度$v_0$水平撞向光滑竖直墙，$C$首先碰墙，杆与墙法线夹角$45^\circ$。碰撞完全弹性，求碰后运动。

\end{example}


\begin{figure}[H]\centering
\begin{tikzpicture}[>=Stealth,line cap=round,line join=round]
\coordinate(A)at(0,0);\coordinate(B)at(1.2,1.2);\coordinate(C)at(2.4,2.4);\draw[curve](A)--(C);\foreach \p in {A,B,C}{\fill(\p)circle(2pt)node[left]{$\p$};}\draw[thick](2.4,-0.5)--(2.4,3);\fill[pattern=north east lines](2.4,-0.5)rectangle(2.52,3);\draw[axis](1.2,1.2)--(2.1,1.2)node[above]{$v_0$};
\end{tikzpicture}
\caption{三个等质量质点固定在轻杆上撞墙}\label{fig:hw5-2-q}
\end{figure}

\begin{solution}
质心在$B$，$I_B=2ml^2$。图\ref{fig:hw5-2-s}中墙对$C$的冲量向左，记大小为$J$，逆时针为正，则
\[
 V_x'=v_0-\frac{J}{3m},\quad V_y'=0,\quad
 \omega'=\frac{J\sin45^\circ}{2ml}.
\]

\begin{figure}[H]\centering
\begin{tikzpicture}[>=Stealth,line cap=round,line join=round]
\coordinate(A)at(0,0);\coordinate(B)at(1.2,1.2);\coordinate(C)at(2.4,2.4);\draw[curve](A)--(C);\foreach \p in {A,B,C}{\fill(\p)circle(2pt);}\node[left]at(A){$A$};\node[left]at(B){$B$};\node[right,xshift=4pt]at(C){$C$};\draw[thick](2.4,-0.5)--(2.4,3);\fill[pattern=north east lines](2.4,-0.5)rectangle(2.52,3);\draw[axis](2.4,2.4)--(1.4,2.4)node[above]{$J$};\node at(4,1){$I_B=2ml^2$};
\draw[axis](0,0)--(.64,-.77)node[right]{$\vct v_A'$};
\draw[axis](2.4,2.4)--(1.5,3.17)node[above left]{$\vct v_C'$};
\end{tikzpicture}
\caption{墙冲量改变角速度，两端球速度均须合成平动与转动}\label{fig:hw5-2-s}
\end{figure}

$C$点法向速度为$V_x'-\omega'l\sin45^\circ$，弹性条件要求它等于$-v_0$。所以$J[1/(3m)+1/(4m)]=2v_0$，得到
\[
 \boxed{J=\frac{24}{7}mv_0,\quad
 \vct V_B'=-\frac17v_0\vct e_x,\quad
 \omega'=\frac{6\sqrt2}{7}\frac{v_0}{l}}.
\]
按$\vct v_i'=\vct V_B'+\omega'\vct e_z\times\vct r_{Bi}$合成三个球的速度，取向右、向上为坐标正向，结果如下，与图\ref{fig:hw5-2-s}的箭头方向一致：
\[
 \begin{array}{c|cc}
 &v_x'/v_0&v_y'/v_0\\\hline
 A&5/7&-6/7\\ B&-1/7&0\\ C&-1&6/7
 \end{array}
\]
质心$B$以$v_0/7$向左反弹，接触球$C$的水平分量反向，竖直分量则由转动产生，其速率为$\sqrt{85}v_0/7$。完全弹性条件只要求法向分量反向，切向分量须由刚体运动求出。

三球的水平总动量为$-3mv_0/7$，比碰前$3mv_0$减少$24mv_0/7$，等于向左的墙冲量；竖直总动量仍为零。关于碰撞瞬间的$B$点，墙冲量的力臂为$l/\sqrt2$，所以角冲量为$Jl/\sqrt2=I_B\omega'$。碰后平动动能为$3mv_0^2/98$，转动动能为$72mv_0^2/49$，二者之和为$3mv_0^2/2$，与碰前相同。系统的大部分动能进入了绕质心的转动。

\end{solution}



\begin{example}[同一构形的刚接与铰接碰撞]\label{ex:hw5-3}

三个质量$m$的小球由两根长$l$的轻杆$AB,BC$连接，$BC$水平，$BA$相对向左水平线向下倾斜$45^\circ$。系统以$v_0$向右平移，$C$与光滑竖直墙完全弹性碰撞。(1)两杆在$B$刚接；(2)两杆在$B$铰接。分别求碰后$C$的速度。

\end{example}


\begin{figure}[H]\centering
\begin{tikzpicture}[>=Stealth,line cap=round,line join=round]
\coordinate(A)at(-1.2,-1.2);\coordinate(B)at(0,0);\coordinate(C)at(1.7,0);\draw[curve](A)--(B)--(C);\foreach \p in {A,B,C}{\fill(\p)circle(2pt)node[above left]{$\p$};}\draw[thick](1.7,-1.8)--(1.7,1);\fill[pattern=north east lines](1.7,-1.8)rectangle(1.8199999999999998,1);\draw[axis](-1.3,.7)--(-.3,.7)node[above]{$v_0$};\node at(-.6,-.2){$45^\circ$};
\end{tikzpicture}
\caption{相同几何构形可采用刚接或铰接}\label{fig:hw5-3-q}
\end{figure}

\begin{solution}
取$x$向右、$y$向上，图\ref{fig:hw5-3-s}中的墙冲量为$-J\vct e_x$，其中$J>0$。墙光滑且恢复系数为$1$，两种连接方式都满足$v_{Cx}'=-v_0$；竖直速度须由各自内部约束确定。
\begin{figure}[H]\centering
\begin{tikzpicture}[>=Stealth,line cap=round,line join=round]
\coordinate(A)at(-1.2,-1.2);\coordinate(B)at(0,0);\coordinate(C)at(1.7,0);\draw[curve](A)--(B)--(C);\foreach \p in {A,B,C}{\fill(\p)circle(2pt)node[above left]{$\p$};}\draw[thick](1.7,-1.8)--(1.7,1);\fill[pattern=north east lines](1.7,-1.8)rectangle(1.8199999999999998,1);\draw[axis](-1.3,.7)--(-.3,.7)node[above]{$v_0$};\node at(-.6,-.2){$45^\circ$};\draw[axis](1.7,0)--(.8,0)node[above]{$J$};\draw[axis](3.8,.4)--(3.1,.85);\node[right]at(3.8,.4){刚接：$\vct v_C'$};
\draw[axis](3.8,-.8)--(3.1,-.8);\node[right]at(3.8,-.8){铰接：$\vct v_C'$};
\end{tikzpicture}
\caption{两种连接方式对碰后速度的约束不同}\label{fig:hw5-3-s}
\end{figure}
刚接时，取$B=(0,0)$、$C=(l,0)$、$A=(-l/\sqrt2,-l/\sqrt2)$。质心坐标及转动惯量为
\[
 G=\left(\frac{l-l/\sqrt2}{3},-\frac{l}{3\sqrt2}\right),\qquad
 I_G=\frac m3(AB^2+BC^2+CA^2)
 =\frac{ml^2}{3}(4+\sqrt2).
\]
其中$CA^2=(2+\sqrt2)l^2$。设$C$相对$G$的位置为$(d_x,d_y)$，则
$d_x=l(2+1/\sqrt2)/3$、$d_y=l/(3\sqrt2)$。
质心动量和绕质心的角动量方程分别给出
\[
 \vct v_G'=\left(v_0-\frac{J}{3m},0\right),\qquad
 I_G\omega'=Jd_y.
\]
将$\vct v_C'=\vct v_G'+\omega'\vct k\times(d_x,d_y)$分解，
\[
 v_{Cx}'=v_0-J\left(\frac1{3m}+\frac{d_y^2}{I_G}\right)
 =v_0-JW,\qquad
 v_{Cy}'=\frac{Jd_xd_y}{I_G},
\]
\[
 W=\frac1{3m}+\frac1{6m(4+\sqrt2)}
 =\frac{9+2\sqrt2}{6m(4+\sqrt2)}.
\]
由$v_{Cx}'=-v_0$得$J=2v_0/W$，代入竖直分量，得到
\[
 \boxed{\vct v_C'=
 \left(-v_0,\frac{2+4\sqrt2}{9+2\sqrt2}v_0\right)}
 \qquad\text{（刚接）}.
\]
竖直分量来自整体转动，并非墙提供了竖直冲量。

铰接时，轻杆只能沿杆方向传递冲量。令$AB$对$A$的冲量为
$p(1,1)/\sqrt2$，$BC$对$B$的冲量为$q\vct e_x$；
$p,q$允许为负。三个球的分量动量方程为
\[
 \begin{aligned}
 m(v_{Ax}'-v_0)&=p/\sqrt2,&mv_{Ay}'&=p/\sqrt2,\\
 m(v_{Bx}'-v_0)&=-p/\sqrt2+q,&mv_{By}'&=-p/\sqrt2,\\
 m(v_{Cx}'-v_0)&=-q-J,&mv_{Cy}'&=0.
 \end{aligned}
\]
图中的$BC$水平，所以其长度约束给$v_{Bx}'=v_{Cx}'=-v_0$。
杆$AB$的长度约束为
$(v_{Bx}'-v_{Ax}')+(v_{By}'-v_{Ay}')=0$。
将上列速度分量代入，后者给出$q=2\sqrt2p$，而$B$的水平动量方程化为
\[
 -2mv_0=-p/\sqrt2+q=3p/\sqrt2.
\]
因此$p=-2\sqrt2mv_0/3$、$q=-8mv_0/3$。再代入$C$的水平动量方程求$J$，得
\[
 \begin{gathered}
 \vct v_A'=v_0(1/3,-2/3),\qquad
 \vct v_B'=v_0(-1,2/3),\\
 J=\frac{14}{3}mv_0,\qquad
 \boxed{\vct v_C'=(-v_0,0)}\quad\text{（铰接）}.
 \end{gathered}
\]
这些速度给出
$T'=\tfrac12mv_0^2(5/9+13/9+1)=3mv_0^2/2=T$，
与完全弹性碰撞一致。$B$处刚接能够传递冲量矩，铰接只保持两杆长度，因而两种模型的接触点切向速度不同。
\end{solution}



\begin{example}[两杆下落时端球与墙碰撞]\label{ex:hw5-4}

三球$A,B,C$分别为$2m,m,m$，两根长$l$的轻杆铰接成$A$、$C$接触光滑地面、$B$在上方的构形。初始两杆近乎竖直，因微小扰动分别向两侧展开。右墙距初始底端位置$a=5\sqrt2l/8$，$C$与墙完全弹性碰撞。证明$C$碰撞前后的速率相等。

\end{example}


\begin{figure}[H]\centering
\begin{tikzpicture}[>=Stealth,line cap=round,line join=round]
\draw[thick](-1,0)--(4.3,0);\coordinate(A)at(0,0);\coordinate(B)at(1.3,1.3);\coordinate(C)at(2.6,0);\draw[curve](A)--(B)--(C);\fill(A)circle(2pt)node[below]{$A,2m$};\fill(B)circle(2pt)node[above]{$B,m$};\fill(C)circle(2pt)node[below]{$C,m$};\draw[thick](2.6,0)--(2.6,2.2);\fill[pattern=north east lines](2.6,0)rectangle(2.72,2.2);
\end{tikzpicture}
\caption{两杆展开到正交位置时端球碰墙}\label{fig:hw5-4-q}
\end{figure}

\begin{solution}
碰墙前系统水平质心不动。两杆等长，$B$的水平坐标是$A,C$坐标的平均值，故$5x_A+3x_C=0$。当$x_C=a$时，$AC=8a/5=\sqrt2l$，两杆在$B$处夹角为$90^\circ$，如图\ref{fig:hw5-4-s}。

\begin{figure}[H]\centering
\begin{tikzpicture}[>=Stealth,line cap=round,line join=round]
\begin{scope}
\draw[thick](-.4,0)--(3,0);\draw[thick](2.6,0)--(2.6,2.1);
\draw[curve](0,0)--(1.3,1.3)--(2.6,0);
\fill(0,0)circle(2pt)node[below]{$A$};\fill(1.3,1.3)circle(2pt)node[above right]{$B$};\fill(2.6,0)circle(2pt)node[below]{$C$};
\draw(1.15,1.15)--(1.3,1)--(1.45,1.15);
\draw[axis](1.3,1.3)--(.75,1.85)node[left]{$Q$};\node at(.4,-.65){$P=0$};
\end{scope}
\begin{scope}[shift={(5.8,.7)}]
\fill(0,0)circle(2pt)node[above right]{$C$};
\draw[axis](0,0)--(-1.2,0)node[left]{$J_w=3mu$};
\draw[axis](0,0)--(0,1.15)node[above]{$J_g=mu$};
\draw[axis](0,0)--(.8,-.8)node[right]{$\vct J_{BC\to C}$};
\node at(.1,-1.3){$\vct J_{BC\to C}=(mu,-mu)$};
\end{scope}
\end{tikzpicture}
\caption{杆的正交关系给出零冲量，墙与地面的冲量均指向允许方向}\label{fig:hw5-4-s}
\end{figure}

设碰前$v_C=u>0$。由上述水平关系与杆长约束，
\[
 v_A=-3u/5,\qquad (v_{Bx},v_{By})=(u/5,-4u/5).
\]
先假设两个底端碰后仍接地，并在求解后检查地面冲量。令$\vct e_{AB}=(1,1)/\sqrt2$，$AB$对$A$的内冲量为$P\vct e_{AB}$；令$BC$对$B$的内冲量为$Q(-1,1)/\sqrt2$。由于两杆正交，$BC$的冲量在$AB$方向投影为零。$A$只能水平运动，因此
\[
 \Delta\vct v_A\cdot\vct e_{AB}=\frac P{4m},\qquad
 \Delta\vct v_B\cdot\vct e_{AB}=-\frac Pm.
\]
两端相对速度沿杆的分量必须相等，故$P=0$，$v_A$不变。若$v_C'=w$，则
\[
 v_{Bx}'=\tfrac12(w-3u/5),\quad
 v_{By}'=-\tfrac12(w+3u/5).
\]
代入完全弹性条件，消去$A$不变的动能：
\[
 u^2+\frac{17}{25}u^2
 =w^2+\frac12w^2+\frac{9}{50}u^2.
\]
由此$w^2=u^2$。碰后$C$离开右墙，故$\boxed{w=-u}$，速率不变。碰后$\vct v_B'=u(-4/5,1/5)$，故$Q=\sqrt2mu$。$BC$对$C$的冲量为$(mu,-mu)$；地面给$C$的竖直冲量为$mu$，墙给$C$的向左冲量为$3mu$，均满足单侧接触方向。因此碰后$C$保持接地的假设成立，所求速率确为$u$。

给定墙距使$AB$的冲量恰好为零。若改变墙距，须重新求两杆内冲量并判断$C$是否离地；只要碰后仍保持地面接触，墙的弹性条件仍使其水平速度反向。可能改变速率的是脱离地面后出现的竖直分量，不能把特殊墙距解释为法向弹性关系成立的条件。


\end{solution}



\begin{example}[直角轻绳系统中的黏合]\label{ex:hw5-5}

三个球$A,B,C$用绷直、不可伸长轻绳$AB,BC$连接，$AB\perp BC$，系统在光滑水平面上静止。第四球$D$沿该直角的角平分线以$v_0$撞向$B$并完全黏合，之后两绳仍绷紧。质量比$m_A:m_B:m_C:m_D=1:1:2:1$。求黏合体速度及损失动能占$D$初动能的比例。

\end{example}


\begin{figure}[H]\centering
\begin{tikzpicture}[>=Stealth,line cap=round,line join=round]
\coordinate(B)at(0,0);\coordinate(A)at(-2,0);\coordinate(C)at(0,-2);\draw[curve](A)--(B)--(C);\fill(A)circle(2pt)node[above left]{$A,m$};\fill(B)circle(2pt)node[above right]{$B,m$};\fill(C)circle(2pt)node[below]{$C,2m$};\fill(-1.5,-1.5)circle(2pt)node[left]{$D,m$};\draw[axis](-1.5,-1.5)--(-.45,-.45)node[midway,above left]{$v_0$};
\end{tikzpicture}
\caption{第四球沿角平分线撞向直角绳结}\label{fig:hw5-5-q}
\end{figure}

\begin{solution}
取图\ref{fig:hw5-5-s}中$x$沿$AB$向右、$y$向上，$D$初速度为$(v_0/\sqrt2,v_0/\sqrt2)$。绳约束要求$v_{Ax}=V_x$、$v_{Cy}=V_y$；$A$没有竖直冲量，$C$没有水平冲量，所以$v_{Ay}=v_{Cx}=0$。

\begin{figure}[H]\centering
\begin{tikzpicture}[>=Stealth,line cap=round,line join=round]
\draw[curve](-2,0)--(0,0)--(0,-2);\fill(-2,0)circle(2pt)node[above]{$A,m$};\fill(0,0)circle(2pt)node[above left]{$BD,2m$};\fill(0,-2)circle(2pt)node[below]{$C,2m$};\draw[axis](0,0)--(1.1,0)node[below]{$V_x$};\draw[axis](0,0)--(0,1)node[right]{$V_y$};\node at(2,-.8){$m_A+m_{BD}=3m$};\node at(2,-1.5){$m_C+m_{BD}=4m$};
\end{tikzpicture}
\caption{黏合体与另两球共同分配两个方向的动量}\label{fig:hw5-5-s}
\end{figure}

黏合后的$BD$质量为$2m$。分别写全系统的水平与竖直动量：
\[
 3mV_x=\frac{mv_0}{\sqrt2},\quad
 4mV_y=\frac{mv_0}{\sqrt2}.
\]
因此
\[
 \boxed{\vct V_{BD}=v_0(\sqrt2/6,\sqrt2/8)}.
\]
碰后总动能为$T'=\tfrac12[3mV_x^2+4mV_y^2]=7mv_0^2/48$，而$D$初动能为$T_0=mv_0^2/2$，故
\[
 \boxed{\frac{T_0-T'}{T_0}=\frac{17}{24}}.
\]
绳$AB$传递的拉冲量为$mV_x=mv_0/(3\sqrt2)>0$，绳$BC$传递的拉冲量为$2mV_y=mv_0/(2\sqrt2)>0$。两根绳都只受拉，与题设碰后仍绷紧的条件相容。

绳传递的冲量使$A,C$也获得动能。计算损失时须计入这两个球的末态动能，单独对$B,D$套用两球黏合公式会漏掉这部分能量。


设两绳拉冲量分别为$J_1,J_2$。对图\ref{fig:hw5-5-s}中的$A,C$分别列冲量方程：$J_1=mV_x$、$J_2=2mV_y$；对黏合体列式则为
\[
 2mV_x=\frac{mv_0}{\sqrt2}-J_1,\qquad
 2mV_y=\frac{mv_0}{\sqrt2}-J_2.
\]
消去$J_1,J_2$便分别出现惯性系数$3m$和$4m$。碰后各部分动能为
\[
 T_A=\frac{mv_0^2}{36},\qquad
 T_C=\frac{mv_0^2}{32},\qquad
 T_{BD}=\frac{25mv_0^2}{288},
\]
相加仍得$7mv_0^2/48$。两个方向的等效质量不同，原因是水平冲量要带动$A$，竖直冲量则要带动质量为$2m$的$C$；它们都不是整个系统的共同平动。

\end{solution}



\begin{example}[小球与可动斜面的碰撞]\label{sup5:problem-example-4}
质量为 $M$ 的斜面放在光滑水平地面上，斜面角为 $\theta$，向右升高。质量为 $m$ 的小球以竖直向下的速度 $u$ 撞击光滑斜面，斜面原来静止，法向恢复系数为 $e$。求碰后小球速度与斜面速度。碰撞期间斜面只做水平平动。
\end{example}
\begin{figure}[H]\centering
\begin{tikzpicture}
\draw[curve](0,0)--(4,0)--(4,2.5)--cycle;
\draw[guide](-.5,-.1)--(5,-.1);\node at(3,.6){$M$};
\draw[curve](1.8834,1.4366)circle(.22);\node[left]at(1.55,1.4366){$m$};
\draw[axis](1.8834,3.1)--(1.8834,1.85)node[midway,right]{$u$};
\draw(1,0)arc(0:32:1);\node at(1.25,.32){$\theta$};
\draw[axis](3.4,2.125)--(2.87,2.973)node[above left]{$\vct n$};
\end{tikzpicture}
\caption{斜面向右升高，指向小球的法线为 $(-\sin\theta,\cos\theta)$。}\label{sup5:fig:d4-e4-q}
\end{figure}
\begin{solution}
取图\ref{sup5:fig:d4-e4-q}中的法线 $\vct n=(-\sin\theta,\cos\theta)$，它从斜面指向小球。设小球受冲量 $J\vct n$，斜面碰后水平速度为 $V$。小球的两分量满足
\begin{equation}
 v_x=-\frac{J\sin\theta}{m},\qquad
 v_y=-u+\frac{J\cos\theta}{m}.
\end{equation}
斜面受到相反冲量，水平动量方程为
\begin{equation}
 MV=J\sin\theta.
\end{equation}
图\ref{sup5:fig:d4-e4-s}中，地面竖直支持冲量为 $J\cos\theta$，抵消小球向斜面传递的向下冲量。斜面的竖直速度因此不变，两物体系统的竖直动量则随外冲量改变。

\begin{figure}[H]\centering
\begin{tikzpicture}
\fill(0,1.2)circle(2pt);\node[left]at(-.15,1.2){小球};
\draw[axis](0,1.2)--(-.8,2.5)node[above]{$J\vct n$};
\draw[curve](2.7,0)rectangle(4.6,.9);\node at(3.65,.45){斜面};
\draw[axis](3.65,.9)--(4.45,-.4)node[right]{$-J\vct n$};
\draw[axis](3.3,-1)--(3.3,.05)node[midway,left]{$J\cos\theta$};
\draw[axis](4.9,.45)--(6.1,.45)node[midway,above]{$V$};
\end{tikzpicture}
\caption{分别画出小球与斜面的冲量；地面支持冲量是外冲量。}\label{sup5:fig:d4-e4-s}
\end{figure}
恢复关系使用小球相对于斜面的法向速度。碰前该速度为 $-u\cos\theta$，碰后为
\begin{align}
 (\vct v-V\vct e_x)\cdot\vct n
 &=-u\cos\theta+J\left(\frac1m+\frac{\sin^2\theta}{M}\right)\\
 &=eu\cos\theta.
\end{align}
故
\begin{equation}
 J=\frac{(1+e)mM u\cos\theta}{M+m\sin^2\theta},\qquad
 V=\frac{(1+e)m u\sin\theta\cos\theta}{M+m\sin^2\theta}.
\end{equation}
把$J$分别代回小球的两个分量式，得到
\[
v_x=-\frac{(1+e)M u\sin\theta\cos\theta}{M+m\sin^2\theta},
\]
\[
\begin{aligned}
v_y&=-u+\frac{(1+e)M u\cos^2\theta}{M+m\sin^2\theta}\\
&=u\,\frac{eM\cos^2\theta-(M+m)\sin^2\theta}{M+m\sin^2\theta}.
\end{aligned}
\]
在$M=m$、$\theta=45^\circ$、$e=1$时，
\begin{equation}
 v_x=-\frac23u,\quad v_y=-\frac13u,\quad V=\frac23u.
\end{equation}
小球在实验室中仍向下运动，但相对斜面的法向速度已经向外。不能以 $v_y<0$ 判定它仍在压入斜面。

上述特例中，小球相对于斜面的速度为 $(-4u/3,-u/3)$。沿外法线投影，分离速度为
\[
\left(-\frac43u,-\frac13u\right)\cdot
\left(-\frac1{\sqrt2},\frac1{\sqrt2}\right)=\frac u{\sqrt2}>0.
\]
图\ref{sup5:fig:d4-e4-q}中斜面同时向右移动，这个相对速度才决定小球是否离开斜面。竖直动量不能用于整体守恒，但水平动量 $mv_x+MV$ 仍为零。

本题地面约束点竖直位移为零，支持冲量不做功；由有效质量可得
\begin{equation}
 Q=\frac{(1-e^2)u^2\cos^2\theta}{2(1/m+\sin^2\theta/M)}.
\end{equation}
在上述弹性特例中，末动能为 $\tfrac12m(5u^2/9)+\tfrac12m(4u^2/9)=mu^2/2$。$M\to\infty$ 时斜面趋于固定，冲量趋于 $(1+e)mu\cos\theta$。
\end{solution}



\begin{example}[滑动终止的冲量条件]\label{sup5:problem-example-9}
均匀圆盘质量为 $m$、半径为 $R$，在竖直平面内运动。碰前质心速度为 $(u,-w)$，$u,w>0$，角速度为零。圆盘撞击固定水平粗糙平面，法向恢复系数为 $e$，摩擦因数为 $\mu$。采用本讲的库仑冲量模型，求碰后质心速度、角速度及能量损失。
\end{example}
\begin{figure}[H]\centering
\begin{tikzpicture}
\draw[curve](0,1)circle(1);\fill(0,1)circle(2pt);\node[left]at(-.15,1){C};
\draw[curve](-2,0)--(4,0);\node[below]at(0,-.1){P};
\draw[axis](0,1)--(1.4,1)node[right]{$u$};
\draw[axis](0,1)--(0,.2)node[midway,right]{$w$};
\draw[guide](0,1)--(-1,1)node[midway,above]{$R$};
\node[right]at(2.1,2){$\omega^-=0$};
\end{tikzpicture}
\caption{碰前P点向右滑动，摩擦冲量向左。}\label{sup5:fig:d4-e9-q}
\end{figure}
\begin{solution}
图\ref{sup5:fig:d4-e9-q}中取向右、向上和逆时针为正。法向冲量 $J_n$ 通过圆心，没有力矩，因此
\begin{equation}
 J_n=m(1+e)w,\qquad v_y=ew.
\end{equation}
设向左的摩擦冲量大小为 $P\ge0$。它同时减小质心水平速度并使圆盘顺时针转动：
\begin{equation}
 v_x=u-\frac Pm,\qquad
 I\omega=-RP,\quad I=\frac12mR^2,\quad
 \omega=-\frac{2P}{mR}.
\end{equation}
接触点P相对平面的切向速度为 $g_t=v_x+R\omega$，所以
\begin{equation}
 g_t=u-\frac{3P}{m}.
\end{equation}
令 $g_t=0$ 得停止滑动所需的冲量 $P_*=mu/3$。库仑摩擦能提供的冲量上限是 $\mu J_n$。按图\ref{sup5:fig:d4-e9-s}，两种分支可合写为
\begin{equation}\label{sup5:eq:d4-friction}
 P=\min\left\{\frac{mu}{3},\ \mu m(1+e)w\right\}.
\end{equation}
\begin{figure}[H]\centering
\begin{tikzpicture}
\draw[axis](0,0)--(5.4,0)node[right]{$P$};
\draw[axis](0,-.2)--(0,2.7)node[above]{$g_t$};
\draw[curve](0,2.2)--(3.5,0);\node[left]at(0,2.2){$u$};
\draw[guide](2,0)--(2,.943);\node[below]at(2,-.1){$\mu J_n<P_*$};
\fill(3.5,0)circle(2pt);\node[below]at(4.15,-.1){$P_*=mu/3$};
\node[right]at(2.15,1.4){持续滑动};
\end{tikzpicture}
\caption{摩擦冲量随接触过程累积；到达 $P_*$ 后切向速度为零，不沿直线继续延长到负值。}\label{sup5:fig:d4-e9-s}
\end{figure}
若 $u>3\mu(1+e)w$，接触结束时仍向右滑动，应取 $P=\mu m(1+e)w$。若 $u\le3\mu(1+e)w$，滑动可在接触期间终止，应取 $P=mu/3$，得到 $v_x=2u/3$、$\omega=-2u/(3R)$。后一个分支仍可能有 $v_y>0$，圆盘随后会离开平面，并非沿平面持续滚动。

切向损失还可沿累积摩擦冲量积分。令 $p$ 从零增加到最终的 $P$，图\ref{sup5:fig:d4-e9-s}中的切向滑动速度为 $u-3p/m$，因此
\[
Q_t=\int_0^P\left(u-\frac{3p}{m}\right)\dd p
=\left[up-\frac{3p^2}{2m}\right]_0^P
=uP-\frac{3P^2}{2m}.
\]
在滑动终止以前，被积函数非负；一旦到达 $P_*$ 就应停止增加摩擦冲量，而不让它沿图中直线继续越过零点。

法向损失为 $m(1-e^2)w^2/2$。水平方向的平动与转动动能合计为
\begin{equation}
 \frac12m\left(u-\frac Pm\right)^2
 +\frac12\frac{mR^2}{2}\left(\frac{2P}{mR}\right)^2
 =\frac12mu^2-uP+\frac{3P^2}{2m}.
\end{equation}
因此总损失为
\begin{equation}
 Q=\frac12m(1-e^2)w^2+uP-\frac{3P^2}{2m}.
\end{equation}
在实际取值 $0\le P\le mu/3$ 内，切向损失非负。若不判断滑动是否已停止，直接令 $P=\mu J_n$，就可能算出不符合所用摩擦模型的反向滑动。
\end{solution}



\begin{example}[圆轨道小珠受到斜向碰撞]\label{deep5:problem-example-10}
固定光滑圆轨道位于水平面内，质量为 $M$ 的小珠原来静止于轨道最右端P。质量为 $m$ 的光滑小球以速度 $u\vct n$ 撞击小珠，接触法线为
\begin{equation*}
 \vct n=\cos\beta\,\vct t+\sin\beta\,\vct e_r,
 \qquad 0\le\beta\le\pi/2,
\end{equation*}
其中 $\vct t$ 在P点向上，$\vct e_r$ 沿径向向右。恢复系数为 $e$，忽略两物体的尺寸。求小球与小珠末速度及轨道冲量。
\end{example}
\begin{figure}[H]\centering
\begin{tikzpicture}
\draw[curve](0,0)circle(1.5);\fill(0,0)circle(1.3pt)(1.5,0)circle(2pt);
\node[left]at(0,0){O};\node[below right]at(1.5,0){P};
\draw[axis](1.5,0)--(1.5,1.4)node[above]{$\vct t$};
\draw[axis](1.5,0)--(3,0)node[right]{$\vct e_r$};
\draw[axis](.4,-1.1)--(1.3,-.2)node[midway,below left]{$u\vct n$};
\draw[axis](1.5,0)--(2.5,1)node[right]{$\vct n$};
\draw(1.5,.65)arc(90:45:.65);\node at(1.9,.82){$\beta$};
\end{tikzpicture}
\caption{小珠只能沿P点切线运动；接触法线与切线成 $\beta$ 角。}\label{deep5:fig:d4-e10-q}
\end{figure}
\begin{solution}
小珠受球的冲量 $J\vct n$。图\ref{deep5:fig:d4-e10-q}中径向运动被轨道限制，切向分量则改变小珠速度。令小珠末速度为 $V\vct t$，轨道给小珠的冲量为 $K\vct e_r$，则
\begin{equation}
 MV=J\cos\beta,\qquad 0=J\sin\beta+K.
\end{equation}
小球所受冲量为 $-J\vct n$，其末速度仍沿 $\vct n$ 的直线：
\begin{equation}
 \vct v=\left(u-\frac Jm\right)\vct n.
\end{equation}
图\ref{deep5:fig:d4-e10-s}的径向冲量互相抵消，切向速度用于恢复关系，
\begin{equation}
 (V\vct t-\vct v)\cdot\vct n
 =-u+J\left(\frac1m+\frac{\cos^2\beta}{M}\right)=eu.
\end{equation}
\begin{figure}[H]\centering
\begin{tikzpicture}
\fill(0,0)circle(2pt);\node[below]at(0,-.15){P};
\draw[axis](0,0)--(1.6,1.6)node[above right]{$J\vct n$};
\draw[guide](1.6,1.6)--(0,1.6);\draw[guide](1.6,1.6)--(1.6,0);
\draw[axis](0,0)--(-1.6,0)node[above]{$K\vct e_r$};
\draw[axis](0,0)--(0,1.6)node[above left]{$MV\vct t$};
\node[right]at(2.1,.4){$K=-J\sin\beta$};
\end{tikzpicture}
\caption{小珠的冲量分解：径向分量由轨道抵消，合冲量沿切线。}\label{deep5:fig:d4-e10-s}
\end{figure}
由此得到
\begin{align}
 J&=\frac{(1+e)mM}{M+m\cos^2\beta}u,\\
 V&=\frac{(1+e)m\cos\beta}{M+m\cos^2\beta}u,\\
 \vct v&=\frac{m\cos^2\beta-eM}{M+m\cos^2\beta}u\vct n,\\
 K&=-\frac{(1+e)mM\sin\beta}{M+m\cos^2\beta}u.
\end{align}
有效质量中的 $\cos^2\beta$ 来自两次投影：冲量 $J$ 的切向分量为 $J\cos\beta$，使小珠获得速度 $V=J\cos\beta/M$；恢复关系中再取该速度的法向分量 $V\cos\beta$。图\ref{deep5:fig:d4-e10-s}中的小珠受圆周约束，碰后速度只能沿切线。

直接展开碰后动能，亦有
\[
T^+-T^-=-uJ+\frac{J^2}{2}
\left(\frac1m+\frac{\cos^2\beta}{M}\right).
\]
记$W=1/m+\cos^2\beta/M$，代入$J=(1+e)u/W$，得
\[
T^+-T^-=-\frac{(1+e)u^2}{W}+\frac{(1+e)^2u^2}{2W}
=-\frac{(1-e^2)u^2}{2W}.
\]

轨道冲量没有切向分量，所以小球与小珠的总切向动量守恒；总径向动量却因 $K$ 改变。固定光滑轨道不做功，损失为
\begin{equation}
 Q=\frac{(1-e^2)u^2}{2(1/m+\cos^2\beta/M)}.
\end{equation}
当 $\beta=0$ 时退化为两自由质点正碰；当 $\beta=\pi/2$ 时小珠不动，冲量完全传给轨道，小球末速度为 $-eu\vct n$。轨道对径向运动的限制已包含在有效质量中，碰撞结果随法线方向改变。
\end{solution}



\begin{example}[三个悬殊质量球的连续碰撞]\label{ex:hw5-8}

三球质量满足$m_1\gg m_2\gg m_3$，从高度$h$一起由静止落下，半径与彼此间距相对$h$可忽略。碰撞按大球与地面、球1与球2、球2与球3的顺序完全弹性发生，之后不再发生球间碰撞。球在水平方向可有微小错位。求最后碰撞后瞬间球3速率的可能范围，采用质量比趋于无穷的极限。

\end{example}


\begin{figure}[H]\centering
\begin{tikzpicture}[>=Stealth,line cap=round,line join=round]
\draw[thick](-1,0)--(2,0);\draw(0,1.1)circle(.55);\draw(.08,1.95)circle(.3);\draw(.15,2.43)circle(.18);\node[right]at(.6,1.1){$m_1$};\node[right]at(.5,1.95){$m_2$};\node[right]at(.5,2.43){$m_3$};\draw[<->](1.5,0)--(1.5,.55)node[midway,right]{$h$};
\end{tikzpicture}
\caption{允许微小水平错位的三个悬殊质量球}\label{fig:hw5-8-q}
\end{figure}

\begin{solution}
落地前共同速率为$u=\sqrt{2gh}$。图\ref{fig:hw5-8-s}用向上单位矢量$\vct e$表示方向。球1反弹后速度为$u\vct e$，球2原速度为$-u\vct e$。在球1的共动系中，球2碰后速度大小为$2u$，所以
\[
 \vct v_2=u\vct e+2u\vct a,\qquad |\vct a|=1.
\]

\begin{figure}[H]\centering
\begin{tikzpicture}[>=Stealth,line cap=round,line join=round]
\begin{scope}[scale=.8]
\draw[axis](-2.5,0)--(3.2,0)node[right]{$v_{2x}/u$};\draw[axis](0,-1.4)--(0,3.8)node[above]{$v_{2z}/u$};
\draw[curve](0,1)circle(2);\fill(0,1)circle(1pt)node[left]{$1$};
\fill(0,3)circle(1.5pt)node[right]{$3$};\fill(1.32287566,-.5)circle(1.5pt)node[right]{$s=-1/2$};
\draw[axis](0,0)--(1.32287566,-.5);\node[below left]at(0,-1){$-1$};
\end{scope}
\begin{scope}[shift={(5.2,.1)}]
\draw[axis](0,-.3)--(0,2.5)node[above]{$z$};\fill(0,0)circle(3pt)node[left]{$m_2$};\fill(0,1.3)circle(2pt)node[left]{$m_3$};
\draw[axis](.4,0)--(1.3,0)node[right]{$v_{2x}$};\draw[axis](.4,1.3)--(.4,.5)node[right]{$-u$};
\node[align=left]at(1.1,-1){第二次法线竖直\\$v_{3z}=(2s+1)u$};\end{scope}
\end{tikzpicture}
\caption{选择中球速度圆上的不同点，第二次碰撞连续给出零至七倍落速}\label{fig:hw5-8-s}
\end{figure}

第二次球间碰撞可把球2视为速度不变的重体。球3原速度仍为$-u\vct e$，在球2系中碰后相对速率为$|\vct v_2+u\vct e|$。因此
\[
 |\vct v_3|\le|\vct v_2|+|\vct v_2+u\vct e|\le3u+4u=7u.
\]
三球共线、碰撞法线竖直时，$v_2=3u$、$v_3=7u$，上界可取到。

下界零也能达到：选择第一次碰后$v_{2z}=-u/2$、水平分量大小$\sqrt7u/2$，它满足上述速度圆；再使第二次碰撞法线竖直。球3的竖直速度从$-u$变为$2(-u/2)-(-u)=0$，水平速度仍为零。更明确地，选球2碰后竖直分量$su$，其中$-1/2\le s\le3$，并令水平分量为$u\sqrt{4-(s-1)^2}$。它始终位于图中的速度圆上。第二次碰撞取竖直法线，则球3碰后速度为$(2s+1)u\vct e$，恰好连续覆盖$0$到$7u$。在题设规定的碰撞次序及无后续球间碰撞的构形中，可由水平错位选择相应法线，故
\[
 \boxed{0\le|\vct v_3|\le7\sqrt{2gh}}.
\]
$7u$是共线碰撞的上界。水平错位使碰撞法线改变，末速率因而可以取整个区间中的值。有限质量比会修正该区间；这里已经采用质量逐级悬殊的极限。


\end{solution}



\begin{example}[弹道摆与整周运动]\label{deep5:problem-example-7}
质量为 $M$ 的小摆球由长 $l$ 的轻绳悬挂，初始位于最低点。质量为 $m$ 的子弹以水平速度 $u$ 射入并留在摆球中，摆球及子弹的尺寸均忽略。求碰撞后共同速度、碰撞损失，以及轻绳始终绷紧完成整周运动所需的最小入射速度。若碰后速度平方小于 $2gl$，求最大摆角。
\end{example}
\begin{figure}[H]\centering
\begin{tikzpicture}
\fill(1,2.5)circle(2pt);\node[above]at(1,2.5){O};
\draw[curve](1,2.5)--(1,-.5);\draw[curve](1,-.5)circle(.27);\node[right]at(1.4,-.5){$M$};
\node[right]at(1.1,1){$l$};
\draw[curve](-1.7,-.5)circle(.11);\node[below]at(-1.7,-.7){$m$};
\draw[axis](-1.4,-.5)--(.5,-.5)node[midway,above]{$u$};
\end{tikzpicture}
\caption{子弹在摆球最低点水平入射，轻绳沿竖直方向。}\label{deep5:fig:d4-e7-q}
\end{figure}
\begin{solution}
碰撞期间轻绳可能提供竖直冲量，但没有水平冲量。图\ref{deep5:fig:d4-e7-q}的水平方向因此满足
\begin{equation}
 mu=(M+m)v_0,\qquad v_0=\frac m{M+m}u.
\end{equation}
碰撞阶段子弹嵌入摆球，不能用机械能守恒求 $v_0$。损失为
\begin{equation}
 Q=\frac12mu^2-\frac12(M+m)v_0^2
 =\frac{mM}{2(M+m)}u^2.
\end{equation}
碰撞结束后，系统质量固定为 $M+m$，轻绳张力不做功。摆到与竖直向下方向成 $\theta$ 角时，图\ref{deep5:fig:d4-e7-s}给出升高量 $l(1-\cos\theta)$，故
\begin{equation}\label{deep5:eq:d4-pendspeed}
 v^2=v_0^2-2gl(1-\cos\theta).
\end{equation}
\begin{figure}[H]\centering
\begin{tikzpicture}[scale=.83]
\coordinate(O)at(0,0);\coordinate(B)at(0,-2);\coordinate(P)at(1.732,-1);\coordinate(T)at(0,2);
\draw[guide](O)circle(2);\fill(O)circle(2pt)(B)circle(2pt)(P)circle(2pt)(T)circle(2pt);
\draw[curve](B)--(O)--(P);\node[left]at(O){O};\node[below left]at(B){最低点};\node[right]at(P){P};\node[above]at(T){最高点};
\draw(0,-.6)arc(-90:-30:.6);\node at(.4,-.7){$\theta$};
\draw[axis](B)--(1.23,-2)node[midway,below]{$v_0$};
\draw[axis](T)--(-.55,2)node[left]{$v_{\rm top}$};
\draw[axis](P)--(2.282,-.047)node[right]{$v$};
\draw[guide](P)--(3,-1);\draw[guide](B)--(3,-2);
\draw[<->](2.8,-2)--(2.8,-1)node[midway,right]{$l(1-\cos\theta)$};
\end{tikzpicture}
\caption{摆动阶段的高度与速度，速度箭头按整周临界状态绘制。右侧切向速度向上，最高点速度向左。}\label{deep5:fig:d4-e7-s}
\end{figure}
若 $v_0^2<2gl$，摆球在下半圆内速度降为零，此时张力仍非负。由式\eqref{deep5:eq:d4-pendspeed}得到
\begin{equation}
 \theta_{\max}=\arccos\left(1-\frac{v_0^2}{2gl}\right).
\end{equation}
在这个转向点，速度为零，而图\ref{deep5:fig:d4-e7-s}的径向方程给出
\[
T_t=(M+m)g\cos\theta_{\max}
=(M+m)g\left(1-\frac{v_0^2}{2gl}\right)>0.
\]
上升段中，$v^2$ 和 $\cos\theta$ 都递减。由 $T=(M+m)(v^2/l+g\cos\theta)$ 可知，张力在转向点最小；该处仍有 $T_t>0$，所以在到达转向点之前绳始终绷紧。条件 $v_0^2<2gl$ 限定了这一运动分支。

整周运动还要求轻绳张力在全程非负。沿指向O的方向列式，
\begin{equation}
 T-(M+m)g\cos\theta=(M+m)\frac{v^2}{l}.
\end{equation}
在最高点 $\theta=\pi$，有 $T_{\rm top}+(M+m)g=(M+m)v_{\rm top}^2/l$，所以 $v_{\rm top}^2\ge gl$。最低点到最高点的能量关系为 $v_0^2=v_{\rm top}^2+4gl$，从而
\begin{equation}
 v_{0,\min}=\sqrt{5gl},\qquad
 u_{\min}=\frac{M+m}{m}\sqrt{5gl}.
\end{equation}
把式\eqref{deep5:eq:d4-pendspeed}代回张力公式，$T=(M+m)(v_0^2/l-2g+3g\cos\theta)$，可知最高点确实是整周运动中张力最小的位置。若仅要求能量足够达到高度 $2l$，会漏掉最高点的绷紧条件。
\end{solution}



\begin{example}[上下两球落地后的速度放大]\label{deep5:problem-example-8}
质量为 $M$ 的大球与质量为 $m$ 的小球上下排列，$M>3m$，小球在上。两球从静止同时下落，落地前速率均近似为 $u=\sqrt{2gh}$。大球先与地面弹性碰撞，随后立即与仍在下降的小球发生弹性正碰。求第二次碰撞后两球速度和小球上升高度。

两球接触事件相互分离，但时间间隔远小于 $u/g$；球的半径及两事件间的高度变化相对 $h$ 均忽略。
\end{example}
\begin{figure}[H]\centering
\begin{tikzpicture}
\draw[curve](0,2.5)circle(.3)(0,1)circle(.55);\node[right]at(.65,1){$M$};\node[right]at(.4,2.5){$m$};
\draw[axis](-1,2.8)--(-1,1.9)node[midway,left]{$-u$};
\draw[axis](1.6,1.4)--(1.6,.5)node[midway,right]{$-u$};
\draw[curve](-1.6,0)--(3,0);\fill[pattern=north east lines](-1.6,-.15)rectangle(3,0);
\node[right]at(2,2.6){向上为正};
\end{tikzpicture}
\caption{两球先共同下落，大球先接触地面。}\label{deep5:fig:d4-e8-q}
\end{figure}
\begin{solution}
图\ref{deep5:fig:d4-e8-q}中地面碰撞结束后，大球速度已变为 $+u$，小球尚未碰撞，速度仍为 $-u$。两球第二次碰撞的接近速度是 $2u$。图\ref{deep5:fig:d4-e8-s}将这两个阶段分别画出。
\begin{figure}[H]\centering
\begin{tikzpicture}
\foreach \x in {0,4}{\draw[curve](\x,2)circle(.25)(\x,.7)circle(.45);}
\draw[axis](1,.5)--(1,1.1)node[right]{$+u$};
\draw[axis](1,2.6)--(1,2)node[right]{$-u$};
\draw[axis](5,2)--(5,3.5)node[right]{$v$};
\draw[axis](5,.5)--(5,.8)node[right]{$V$};
\node[below]at(0,0){地面碰撞后};\node[below]at(4,0){球间碰撞后};
\draw[axis](1.7,1.6)--(2.7,1.6);
\end{tikzpicture}
\caption{球间碰撞前后的速度。末态按 $M=7m$ 示意，四个速度箭头采用同一比例。}\label{deep5:fig:d4-e8-s}
\end{figure}
第二次碰撞期间地面不再接触大球，两球竖直动量守恒：
\begin{equation}
 Mu-mu=MV+mv.
\end{equation}
恢复关系给 $v-V=2u$。代入动量方程得
\begin{equation}
 V=\frac{M-3m}{M+m}u,\qquad
 v=\frac{3M-m}{M+m}u.
\end{equation}
在题设 $M>3m$ 下，两球碰后都向上运动，且小球更快。小球从第二次碰撞位置向上升高
\begin{equation}
 H=\frac{v^2}{2g}
 =\left(\frac{3M-m}{M+m}\right)^2h.
\end{equation}
当 $M/m\to\infty$，有 $v\to3u$、$H\to9h$。有限质量比时高度小于 $9h$。能量并没有增加：大球的碰后动能相应减小，且
\begin{equation}
 MV^2+mv^2=(M+m)u^2.
\end{equation}
上述结果要求两次接触在时间上分开。若球间接触与地面接触持续重叠，应同时处理两处接触，不能再按两个独立的二体碰撞计算。
小球上升高度增加所需的能量来自大球。第二次碰撞中，大球减少的动能与小球增加的动能分别为
\[
\frac12M(u^2-V^2)
=\frac{4Mm(M-m)}{(M+m)^2}u^2,
\]
\[
\frac12m(v^2-u^2)
=\frac{4Mm(M-m)}{(M+m)^2}u^2.
\]
两式相等，说明第二次弹性碰撞中小球增加的动能全部来自大球。将图\ref{deep5:fig:d4-e8-s}中的 $M=7m$ 及相应速度代入，也可逐项得到这一结果。$M/m$ 很大时，大球的速度变化很小，但质量与速度变化量的乘积仍有限，因而传递的能量不趋于零。

\end{solution}



\section{习题与解答}

\begin{exercise}[圆盘边缘的偏心黏合]\label{sup5:problem-exercise-7}
质量为 $M$、半径为 $R$ 的均匀圆盘静止在光滑水平气垫台上。质量为 $m$ 的小颗粒以速度 $u$ 沿圆盘边缘切线运动，在P点与圆盘黏合。取圆盘原中心C为原点，P为 $(R,0)$，颗粒初速度沿 $+y$ 方向。求组合体的质心速度、角速度和碰撞损失。
\end{exercise}
\begin{figure}[H]\centering
\begin{tikzpicture}
\draw[curve](0,0)circle(1.5);\fill(0,0)circle(2pt)(1.5,0)circle(2pt);
\draw[guide](0,0)--(1.5,0)node[midway,above]{$R$};\node[left]at(0,0){C};\node[right]at(1.6,0){P};
\draw[curve](1.5,-2.3)circle(.1);\node[right]at(1.7,-2.3){$m$};
\draw[axis](1.5,-2)--(1.5,-.4)node[midway,right]{$u$};
\node[left]at(-1.6,.8){$M$};
\end{tikzpicture}
\caption{颗粒沿边缘切线入射，圆盘不受固定转轴约束。}\label{sup5:fig:d4-h7-q}
\end{figure}
\begin{solution}
图\ref{sup5:fig:d4-h7-q}中的圆盘中心C并不是组合体的质心。设组合体质心为G，$CG=c$，则
\begin{equation}
 c=\frac{mR}{M+m},\qquad
 \vct V_G=\frac{mu}{M+m}\vct e_y.
\end{equation}
先取关于G的角动量。碰前圆盘不动，颗粒对G的力臂为 $R-c$，故
\begin{equation}
 L_G^-=mu(R-c)=\frac{mM}{M+m}Ru.
\end{equation}
组合体关于G的转动惯量应同时计入圆盘中心偏离G的平行轴项和颗粒项：
\begin{align}
 I_G&=\frac12MR^2+Mc^2+m(R-c)^2\\
 &=\frac12MR^2+\frac{mM}{M+m}R^2
 =\frac{MR^2(M+3m)}{2(M+m)}.
\end{align}
\begin{figure}[H]\centering
\begin{tikzpicture}
\draw[curve](0,0)circle(1.5);\fill(0,0)circle(1.5pt)(.6,0)circle(2pt)(1.5,0)circle(2pt);
\node[below]at(0,-.1){C};\node[below right]at(.7,-.1){G};\node[right]at(1.6,0){P};
\draw[<->](0,-.8)--(.6,-.8)node[midway,below]{$c$};
\draw[guide](.6,0)--(.6,-1.8)(1.5,0)--(1.5,-1.8);
\draw[<->](.6,-1.8)--(1.5,-1.8)node[midway,below]{$R-c$};
\draw[axis](.6,0)--(.6,1.4)node[above]{$V_G$};
\draw[axis](-.1524,2.0673)arc(110:220:2.2)node[below left]{$\omega$};
\end{tikzpicture}
\caption{关于组合体质心G分解平动和转动；图示 $m=2M/3$ 的位置比例。}\label{sup5:fig:d4-h7-s}
\end{figure}
由图\ref{sup5:fig:d4-h7-s}所示的质心平动与绕G转动分解，角动量守恒给
\begin{equation}
 \omega=\frac{L_G^-}{I_G}=\frac{2mu}{R(M+3m)},
\end{equation}
方向为逆时针。末动能为
\begin{align}
 T^+&=\frac12(M+m)V_G^2+\frac12I_G\omega^2\\
 &=\frac{m^2u^2}{2(M+m)}
 +\frac{Mm^2u^2}{(M+m)(M+3m)}
 =\frac{3m^2u^2}{2(M+3m)}.
\end{align}
因此
\begin{equation}
 Q=\frac12mu^2-T^+=\frac{mM}{2(M+3m)}u^2.
\end{equation}
圆盘没有固定轴，黏合后的运动包含组合体平动和绕G转动。若只用圆盘关于C的转动惯量，会同时漏掉质心偏移和组合体平动。图中的G随组合体运动；本题外力矩为零，关于G的角动量仍守恒。
图\ref{sup5:fig:d4-h7-s}中，圆盘原中心C也在运动。由刚体点速度关系，可求得碰后瞬间
\[
v_C=V_G-\omega c=\frac{mu}{M+3m},\qquad
v_P=V_G+\omega(R-c)=\frac{3mu}{M+3m}.
\]
两者均沿原颗粒的入射方向。黏合颗粒与圆盘P点具有同一末速度，但这个速度既不等于 $V_G$，也不等于圆盘中心速度；因此用固定轴转动来代替组合体运动会改变动能与角动量的分配。

\end{solution}


\begin{exercise}[两球与固定墙的三次碰撞]\label{sup5:problem-exercise-10}
一条光滑直线上有两个相同小球A、B，质量均为 $m$，右端是固定弹性墙。A、B的恢复系数为 $e$，$0<e\le1$。A以速度 $u$ 向右碰上原来静止的B，第一次碰撞位置取为 $x=0$，墙在 $x=\ell$。忽略小球尺寸，且各次接触相互分离。求此后各次碰撞的速度、第二次球间碰撞的时刻与位置，以及总损失。
\end{exercise}
\begin{figure}[H]\centering
\begin{tikzpicture}
\draw[curve](0,0)circle(.3)(1.6,0)circle(.3);\node at(0,0){A};\node at(1.6,0){B};
\draw[axis](.45,0)--(1.05,0)node[midway,above]{$u$};
\draw[line width=2pt](4,-.8)--(4,1.2);\node[above]at(4,1.2){墙};
\draw[axis](-.6,-1.25)--(4.7,-1.25)node[right]{$x$};
\draw[guide](1.6,-.35)--(1.6,-1.25);
\node[below]at(4,-1.3){$\ell$};\node[below]at(1.6,-1.3){$0$};
\end{tikzpicture}
\caption{碰前的位置关系。静止球B的初始位置取为 $x=0$，第一次接触结束时取 $t=0$。}\label{sup5:fig:d4-h10-q}
\end{figure}
\begin{solution}
第一次球间碰撞的动量方程与恢复关系分别为 $u=a+b$、$b-a=eu$，故
\begin{equation}
 a=\frac{1-e}{2}u,\qquad b=\frac{1+e}{2}u.
\end{equation}
图\ref{sup5:fig:d4-h10-q}中B先到墙，时刻为 $t_w=\ell/b$。墙反射后B速度为 $-b$，A仍以速度 $a$ 向右运动。图\ref{sup5:fig:d4-h10-s}中的两条位置曲线在B返回段相交，满足
\begin{equation}
 x_A=at,\quad x_B=2\ell-bt,\quad
 t_2=\frac{2\ell}{a+b}=\frac{2\ell}{u},\quad
 x_2=(1-e)\ell.
\end{equation}
因 $t_2/t_w=1+e>1$，第二次球间碰撞发生在墙反射之后，位置不超过墙的位置。
\begin{figure}[H]\centering
\begin{tikzpicture}[x=1cm,y=1cm]
\draw[axis](0,0)--(5.4,0)node[right]{$ut/\ell$};\draw[axis](0,0)--(0,3)node[above]{$x/\ell$};
\draw[guide](0,2.4)--(5,2.4)node[right]{$1$};
\draw[curve](0,0)--(2.667,2.4)--(4,1.2)--(5,1.2);
\draw[curve](0,0)--(4,1.2)--(5,.6);
\node[below]at(1.9,.55){A};\node[above left]at(1.35,1.25){B};
\fill(2.667,2.4)circle(2pt)(4,1.2)circle(2pt);
\node[above left]at(2.667,2.4){撞墙};\node[above right]at(4,1.2){再次接触};
\node[below]at(4,0){$2$};\node[left]at(0,1.2){$1-e$};
\node[below]at(2.5,-.55){图按 $e=1/2$ 绘制};
\end{tikzpicture}
\caption{球A与球B的位置随时间变化，墙反射与再次球碰的时刻均由轨迹交点确定。}\label{sup5:fig:d4-h10-s}
\end{figure}
B碰墙时，A的位置为
\[
x_A(t_w)=\frac ab\ell=\frac{1-e}{1+e}\ell<\ell.
\]
所以图\ref{sup5:fig:d4-h10-s}中的A不会先赶到墙边。B反弹到再次碰球的间隔为
\[
t_2-t_w=\frac{2e\ell}{(1+e)u}>0.
\]
当 $e\to0^+$ 时，这一间隔趋于零。在 $e=0$ 处，球间接触与碰墙事件不能按题设的顺序分别处理，因此该解限于 $e>0$。

第二次球间碰撞的初速度为 $a,-b$，两球接近速率仍为 $a+b=u$。应用等质量公式，
\begin{align}
 v_A&=\frac{1-e}{2}a-\frac{1+e}{2}b=-eu,\\
 v_B&=\frac{1+e}{2}a-\frac{1-e}{2}b=0.
\end{align}
A向左离去，B静止，此后没有新的接触。接触总数为三次：球间碰撞、B撞墙、球间再碰。两次球间碰撞各损失 $m(1-e^2)u^2/4$，墙反射不损失能量，因此
\begin{equation}
 Q_{\rm total}=\frac12m(1-e^2)u^2.
\end{equation}
这与初动能减去末动能 $me^2u^2/2$ 一致。$e=0$ 时首次碰撞已无法保证两个独立接触事件，故题目排除了这个边界。
\end{solution}


\begin{exercise}[碰撞后的弹簧运动与再次接触]\label{deep5:problem-exercise-4}
光滑水平面上，质量为 $3m$ 的滑块连接劲度系数 $k$ 的水平轻弹簧，弹簧另一端固定在右侧墙上。弹簧初始无形变，滑块静止。质量为 $m$ 的小球以速度 $u$ 从左撞上滑块，碰撞完全弹性，接触时间远小于弹簧振动周期。求最大压缩量，并判断小球是否会再次碰到滑块。忽略物体尺寸。
\end{exercise}
\begin{figure}[H]\centering
\begin{tikzpicture}
\draw[curve](1.6,0)rectangle(2.6,1);\node at(2.1,.5){$3m$};
\draw[spring](2.6,.5)--(5,.5);\draw[line width=2pt](5,0)--(5,1.5);\node[above]at(3.7,.7){$k$};
\draw[curve](-1,.5)circle(.5);\node[left]at(-1.6,.5){$m$};
\draw[axis](-.35,.5)--(.8,.5)node[midway,above]{$u$};
\draw[guide](-1.8,0)--(5,0);
\end{tikzpicture}
\caption{碰撞发生时弹簧无形变，短时弹簧冲量忽略。}\label{deep5:fig:d4-h4-q}
\end{figure}
\begin{solution}
先处理图\ref{deep5:fig:d4-h4-q}中的碰撞。动量守恒与$e=1$分别给出$u=v+3V$、$V-v=u$。第二式给$v=V-u$，代入第一式得$4V=2u$，所以小球速度$v=-u/2$、滑块速度$V=u/2$。从碰撞结束计时，以原接触位置为$x=0$。小球不受水平力，其位置由恒速积分为$x_{\rm ball}=vt$。

滑块满足$3m\ddot x+kx=0$。令$\omega=\sqrt{k/(3m)}$，特征根为$\pm\mathrm i\omega$，通解为$x=C\cos\omega t+D\sin\omega t$。初始$x(0)=0$给出$C=0$；求导后代入$\dot x(0)=u/2$，得$\omega D=u/2$。因此图\ref{deep5:fig:d4-h4-s}所示两位置函数为
\begin{equation}
 x_{\rm ball}=-\frac u2t,\qquad
 x_{\rm block}=\frac{u}{2\omega}\sin\omega t,
 \quad\omega=\sqrt{\frac{k}{3m}}.
\end{equation}
滑块在 $t=\pi/(2\omega)$ 首次达到最右点，弹簧最大压缩量为
\begin{equation}
 x_{\max}=\frac{u}{2\omega}=\frac u2\sqrt{\frac{3m}{k}}.
\end{equation}
\begin{figure}[H]\centering
\begin{tikzpicture}[x=.85cm,y=.8cm]
\draw[axis](0,0)--(7,0)node[right]{$\omega t$};\draw[axis](0,-3.4)--(0,1.7)node[above]{$2\omega x/u$};
\draw[curve,domain=0:6.283,samples=90]plot(\x,{sin(\x r)});
\draw[guide](0,0)--(3.1,-3.1)node[below right]{$-\omega t$};
\node[above]at(2.4,.9){$\sin\omega t$};
\node[below]at(3.142,0){$\pi$};\node[below]at(6.283,0){$2\pi$};
\node[left]at(0,1){$1$};\node[left]at(0,-1){$-1$};
\end{tikzpicture}
\caption{同一位置尺度下，滑块的位置有界，小球持续向左远离。}\label{deep5:fig:d4-h4-s}
\end{figure}
滑块稍后会向左返回，仍须判断它能否追上小球。令两位置相等，置 $z=\omega t>0$，须有
\begin{equation}
 \sin z=-z.
\end{equation}
在 $0<z<\pi$ 内左边为正，右边为负；在 $z\ge\pi$ 内，$-z<-1\le\sin z$。两区间均无解，因此没有第二次碰撞。能量验算中应只把滑块碰后的动能用于弹簧压缩：$3mu^2/8=kx_{\max}^2/2$，小球携带的动能不参与后续振动。
忽略碰撞期间弹簧冲量的条件也可量化。设接触时间为 $\delta t$，其间滑块速度量级为 $V$。由于碰前弹簧无形变，碰撞中弹簧力从零增至约 $kV\delta t$，故其冲量量级为
\[
J_k\sim\frac12kV(\delta t)^2,
\qquad \frac{J_k}{3mV}\sim\frac12(\omega\delta t)^2\ll1.
\]
因此图\ref{deep5:fig:d4-h4-q}中的弹簧在短时碰撞方程中可略去，在图\ref{deep5:fig:d4-h4-s}的后续振动中却必须保留。

\end{solution}


\begin{exercise}[粗糙碰撞的滑动分支]\label{deep5:problem-exercise-9}
均匀圆盘的参数与例\ref{sup5:problem-example-9}相同，碰前水平速度为 $u=3w$，竖直速度为 $-w$，角速度为零。取 $e=1/2$、$\mu=2/5$。求碰后速度、角速度、接触点切向速度及损失，并求本题恰好终止滑动所需的最小摩擦因数。
\end{exercise}
\begin{figure}[H]\centering
\begin{tikzpicture}
\draw[curve](0,1)circle(1);\fill(0,1)circle(2pt);\draw[curve](-1.8,0)--(4.5,0);
\draw[axis](0,1)--(2.4,1)node[right]{$3w$};
\draw[axis](0,1)--(0,.2)node[midway,right]{$w$};
\node[right]at(2.6,2){$e=1/2,\ \mu=2/5$};\node[below]at(0,-.1){P};
\end{tikzpicture}
\caption{圆盘先向右滑动，摩擦冲量不足时，接触结束仍会保留右向滑动。}\label{deep5:fig:d4-h9-q}
\end{figure}
\begin{solution}
由图\ref{deep5:fig:d4-h9-q}的法向速度，$J_n=m(1+e)w=3mw/2$。向左摩擦冲量的上限为
\begin{equation}
 \mu J_n=\frac35mw.
\end{equation}
终止滑动需要 $P_*=mu/3=mw$。摩擦冲量上限小于这个值，接触结束时仍在滑动，故取 $P=3mw/5$。
\begin{equation}
 v_x=3w-\frac35w=\frac{12}{5}w,\qquad
 v_y=\frac12w,\qquad
 \omega=-\frac{6w}{5R}.
\end{equation}
\begin{figure}[H]\centering
\begin{tikzpicture}
\draw[curve](0,1)circle(1);\draw[curve](-1.8,0)--(4.6,0);\fill(0,0)circle(2pt);
\draw[axis](0,1)--(2.4,1)node[right]{$v_x=12w/5$};
\draw[axis](0,1)--(0,1.5);
\draw[guide](-1.55,1.5)--(0,1.5);\node[left]at(-1.55,1.5){$v_y=w/2$};
\draw[axis](0,0)--(1.2,0)node[below right]{$g_t=6w/5$};
\draw[axis](-1.2235,1.5705)arc(155:25:1.35);\node at(1.7,2.2){$\omega<0$};
\end{tikzpicture}
\caption{碰后质心向右上方运动，圆盘顺时针转动，接触点的切向相对速度仍向右。}\label{deep5:fig:d4-h9-s}
\end{figure}
在图\ref{deep5:fig:d4-h9-s}的P点，转动速度向左，故
\begin{equation}
 g_t=v_x+R\omega=\frac{12}{5}w-\frac65w=\frac65w>0.
\end{equation}
法向损失为 $3mw^2/8$，切向损失为
\begin{equation}
 uP-\frac{3P^2}{2m}
 =\left(\frac95-\frac{27}{50}\right)mw^2
 =\frac{63}{50}mw^2.
\end{equation}
总损失为 $Q=327mw^2/200$。恰好终止滑动要求 $\mu_*J_n=P_*$，所以
\begin{equation}
 \mu_* =\frac{u}{3(1+e)w}=\frac23.
\end{equation}
这一门槛来自接触点速度，并非来自质心水平速度降至零。即使达到门槛，质心仍有 $v_x=2w$。
分别计算碰后的水平平动能、竖直平动能和转动能。由图\ref{deep5:fig:d4-h9-s}中的速度，得到
\[
K_x=\frac{72}{25}mw^2,\quad K_y=\frac18mw^2,
\quad K_{\rm rot}=\frac9{25}mw^2.
\]
三项合为 $673mw^2/200$，初动能为 $5mw^2$，差值恰为 $327mw^2/200$。摩擦减小了水平平动能，同时增加了转动能，只有两者之和的减少才是切向耗散。

\end{solution}


\begin{exercise}[选做：轻绳松弛后的最高点]\label{deep5:problem-exercise-12}
弹道摆的结构与例\ref{deep5:problem-example-7}相同。子弹嵌入后，组合体在最低点的速度满足 $v_0^2=3gl$。求轻绳首次松弛时的摆角、速率，以及松弛后的首次抛体飞行中最高点相对最低点的高度。松弛后忽略空气阻力。
\end{exercise}
\begin{figure}[H]\centering
\begin{tikzpicture}
\draw[guide](0,0)circle(1.8);\fill(0,0)circle(2pt)(0,-1.8)circle(2pt);\node[left]at(0,0){O};\node[below]at(0,-1.8){B};
\draw[curve](0,0)--(0,-1.8);\draw[axis](0,-1.8)--(1.3,-1.8)node[right]{$\sqrt{3gl}$};
\draw[axis](1.8,0)--(1.8,1)node[above]{$\vct e_y$};
\node[right]at(2,0){最低点为高度零点};
\end{tikzpicture}
\caption{初速不足以保持整周绷紧，须先求张力降为零的位置。}\label{deep5:fig:d4-h12-q}
\end{figure}
\begin{solution}
在图\ref{deep5:fig:d4-h12-q}中，令 $\theta$ 从竖直向下方向量起。只要绳仍绷紧，能量关系给
\begin{equation}
 v^2=3gl-2gl(1-\cos\theta)=gl(1+2\cos\theta).
\end{equation}
设组合体质量为 $\mathcal M=M+m$，向心方程给
\begin{equation}
 T=\mathcal M\left(\frac{v^2}{l}+g\cos\theta\right)
 =\mathcal Mg(1+3\cos\theta).
\end{equation}
轻绳不能提供负张力，首次松弛时
\begin{equation}
 \cos\theta_s=-\frac13,\qquad
 \theta_s=\arccos(-1/3),\qquad v_s=\sqrt{\frac{gl}{3}}.
\end{equation}
此点S已位于O的上方。图\ref{deep5:fig:d4-h12-s}把绷紧圆弧与松弛后的抛物线分开画出。
\begin{figure}[H]\centering
\begin{tikzpicture}
\begin{scope}[scale=1.4]
\draw[guide](0,0)circle(1);\fill(0,0)circle(1.3pt);\node[left]at(0,0){O};
\draw[curve](0,-1)arc(-90:19.471:1);
\draw[curve,domain=0:.544331,samples=50]plot({.942809-.192450*\x},{.333333+.544331*\x-.5*\x*\x});
\fill(.942809,.333333)circle(1.5pt)(.838052,.481481)circle(1.5pt);
\node[right]at(.99,.30){S};\node[above left]at(.838052,.53){A};
\draw[guide](0,0)--(.942809,.333333);\node[below]at(0,-1){B};
\draw[guide](0,-1)--(-1.5,-1);\draw[guide](.838052,.481481)--(-1.5,.481481);
\draw[<->](-1.3,-1)--(-1.3,.481481)node[midway,left]{$40l/27$};
\end{scope}
\begin{scope}[shift={(6.5,-1.2)}]
\draw[axis](0,0)--(-2.7,0)node[left]{$-\Delta x$};
\draw[axis](0,0)--(0,3.1)node[above]{$\Delta y$};
\draw[curve,domain=0:.544331,samples=70]plot({-18*.192450*\x},{18*(.544331*\x-.5*\x*\x)});
\fill(0,0)circle(1.8pt)(-1.8856,2.6667)circle(1.8pt);
\node[below right]at(0,0){S};\node[above left]at(-1.8856,2.6667){A};
\draw[guide](-1.8856,2.6667)--(0,2.6667)node[right]{$4l/27$};
\draw[axis](0,0)--(-.8,2.263);\node[right]at(-.8,2.263){$\vct v_s$};
\node at(-1.3,4.1){S附近的轨迹放大};
\end{scope}
\end{tikzpicture}
\caption{S处松绳后进入圆内的抛物线，A为首次抛体飞行的最高点；右图以S为局部坐标原点。}\label{deep5:fig:d4-h12-s}
\end{figure}
取最低点B为高度零点，松弛高度为 $y_s=l(1-\cos\theta_s)=4l/3$。在S点，沿运动方向的单位切向矢量为 $(\cos\theta_s,\sin\theta_s)$，故
\begin{equation}
 v_{sx}=-\frac13\sqrt{\frac{gl}{3}},\qquad
 v_{sy}=\frac{2\sqrt2}{3}\sqrt{\frac{gl}{3}}.
\end{equation}
进入抛体运动后，额外上升量仅由竖直速度决定：
\begin{equation}
 \Delta y=\frac{v_{sy}^2}{2g}=\frac{4l}{27},\qquad
 y_{\max}=\frac{4l}{3}+\frac{4l}{27}=\frac{40l}{27}.
\end{equation}
还须确认到达A前绳不会重新绷紧。以O为原点、$l$ 为长度单位、$\sqrt{l/g}$ 为时间单位，令 $s=t\sqrt{g/l}$，松弛后的轨迹为
\begin{equation}
 x=\frac{2\sqrt2}{3}-\frac{s}{3\sqrt3},\qquad
 y=\frac13+\frac{2\sqrt2}{3\sqrt3}s-\frac{s^2}{2}.
\end{equation}
记初始无量纲位置为 $(x_0,y_0)$、速度为 $(p,q)$：
\[
(x_0,y_0)=(2\sqrt2/3,1/3),\qquad
(p,q)=(-1/(3\sqrt3),2\sqrt2/(3\sqrt3)).
\]
将 $x=x_0+ps$、$y=y_0+qs-s^2/2$ 平方，先按 $s$ 的次数排列：
\[
x^2+y^2-1=(x_0^2+y_0^2-1)+2(x_0p+y_0q)s
+(p^2+q^2-y_0)s^2-qs^3+\frac{s^4}{4}.
\]
其中$x_0^2+y_0^2=1$、$x_0p+y_0q=0$、$p^2+q^2=y_0=1/3$，前三项均抵消，所以
\begin{equation}
 x^2+y^2-1=s^3\left(\frac s4-\frac{2\sqrt2}{3\sqrt3}\right).
\end{equation}
A点对应 $s_A=2\sqrt2/(3\sqrt3)$，在 $0<s\le s_A$ 内右端为负，质点到O的距离始终小于绳长。因此图中的抛物线到A为止满足松绳条件。
图\ref{deep5:fig:d4-h12-s}中的最高点A还可求出完整坐标。松绳后达到A的时间为
\[
\Delta t_A=\frac{v_{sy}}g
=\frac{2\sqrt2}{3\sqrt3}\sqrt{\frac lg}.
\]
取O为原点，代入水平匀速与竖直抛体方程得
\[
x_A=\frac{16\sqrt2}{27}l,\qquad y_A=\frac{13}{27}l.
\]
最低点B在 $y=-l$，所以A相对B高 $40l/27$。A处竖直速度为零，水平速度仍等于 $v_{sx}$；保留的水平动能相当于高度 $v_{sx}^2/(2g)=l/54$，故不能把初动能全部换成重力势能来求最高点。

\end{solution}


\chapter{刚体运动进阶}



\section{惯量张量、主轴与角动量}

空间转动中，角动量一般不沿角速度方向。惯量张量由质量分布及取轴位置决定；只有满足相应主轴条件，才能使用沿轴的标量关系。



\subsection{惯量张量与角动量}
取$O$为刚体上的固定点，或取$O$为质心并使用相对质心的速度。以该点为原点定义
\[
 \mathsf I=\int[(\vct r\cdot\vct r)\mathsf1-\vct r\vct r^{\mathsf T}]\dd m,
 \quad\vct L_O=\mathsf I\vct\omega,\quad T_{\text{转}}=\tfrac12\vct\omega\cdot\mathsf I\vct\omega.
\]
惯量张量的非对角元为$-\int xy\dd m$等。沿单位轴$\vct n$的转动惯量是$\vct n^{\mathsf T}\mathsf I\vct n$；角动量通常不与角速度平行。在主轴坐标系中，惯量张量为对角矩阵；平行轴定理的张量形式为$\mathsf I_O=\mathsf I_C+m(d^2\mathsf1-\vct d\vct d^{\mathsf T})$。

随刚体转动的坐标系中，矢量导数要加运输项：$(\dd\vct L/\dd t)_{\text{地}}=(\dd\vct L/\dd t)_{\text{体}}+\vct\omega\times\vct L$。主轴下因此得到欧拉方程$I_1\dot\omega_1+(I_3-I_2)\omega_2\omega_3=M_1$及循环式。
\subsubsection{圆锥的主惯量与进动取轴}
\label{sec:dynamics-cone-inertia}
均匀圆锥的密度为 $\rho=3m/(\pi r^2h)$。从尖点沿轴量 $z$，薄片半径为 $rz/h$，故质量元为 $\dd m=\rho\pi(rz/h)^2\dd z=3mz^2\dd z/h^3$。薄片绕轴的惯量是其质量乘半径平方的一半，绕直径的惯量则为四分之一；移到过锥尖的横轴还须加 $z^2\dd m$。因此
\[
 \ell=\frac1m\int_0^hz\dd m
 =\frac3{h^3}[z^4/4]_0^h=\frac{3h}{4},
\]
\[
 I_3=\int_0^h\frac12\left(\frac{rz}{h}\right)^2\dd m
 =\frac{3mr^2}{2h^5}[z^5/5]_0^h=\frac{3mr^2}{10},
\]
\[
 \begin{aligned}
 I_{1,O}&=\int_0^h\left[z^2+\frac14\left(\frac{rz}{h}\right)^2\right]\dd m\\
 &=\frac{3m}{h^3}\left(1+\frac{r^2}{4h^2}\right)[z^5/5]_0^h
 =\frac{3mh^2}{5}+\frac{3mr^2}{20}.
 \end{aligned}
\]
关于质心的横向主惯量还应减去$m\ell^2$。固定尖点的题目可绕尖点列角动量方程；悬挂尖点随进动移动的题目，改绕质心计算便可避开运动取矩点的附加项。

\LectureNeedspace{8\baselineskip}


\subsection{惯量张量的分量与主轴}
\label{sec:complete6}
以固定点或质心为原点，刚体质元的相对速度为$\vct\omega\times\vct r$。利用矢量三重积，角动量为
\[
 \vct L=\int\vct r\times(\vct\omega\times\vct r)\dd m
 =\int[r^2\vct\omega-(\vct r\cdot\vct\omega)\vct r]\dd m.
\]
在直角坐标中，这一线性关系写成
\[
 \begin{pmatrix}L_x\\L_y\\L_z\end{pmatrix}
 =\begin{pmatrix}
 \int(y^2+z^2)\dd m&-\int xy\dd m&-\int xz\dd m\\
 -\int xy\dd m&\int(z^2+x^2)\dd m&-\int yz\dd m\\
 -\int xz\dd m&-\int yz\dd m&\int(x^2+y^2)\dd m
 \end{pmatrix}
 \begin{pmatrix}\omega_x\\\omega_y\\\omega_z\end{pmatrix}.
\]
有些记号把$\int xy\dd m$称为惯量积，矩阵的对应元还带负号；另一些记号将负号计入惯量积的定义。计算时应先声明采用哪一种。本讲以矩阵元本身为准。

质量分布关于$xz$平面对称时，$y$为奇函数的积分为零，故$\int xy\dd m=\int yz\dd m=0$。对称面足够时，主轴方向可由对称性确定；其余情形须求惯量矩阵的本征矢量。主轴上的本征值是主惯量，而非对角元为零只是所选主轴坐标的性质。

绕单位轴$\vct n=(n_1,n_2,n_3)$转动时，$\vct\omega=\omega\vct n$，因此
\[
 I_{\vct n}=\vct n^{\mathsf T}\mathsf I\vct n,
 \quad T_{\rm rot}=\frac12I_{\vct n}\omega^2,
 \quad \vct L\cdot\vct n=I_{\vct n}\omega.
\]
最后一式只给轴向角动量。若$\vct n$不是主轴，$\vct L$还含横向分量，不能直接写成$\vct L=I_{\vct n}\vct\omega$。例\ref{ex:hw6-1}至例\ref{ex:hw6-4}分别用对称性、方向余弦、平行轴移动和体积分展示这几种取轴方法。

\subsubsection{常用惯量定理的范围}
平行轴公式的中心轴必须通过质心，两轴须平行。薄片的垂直轴定理$I_z=I_x+I_y$则要求三轴交于同一点，且$x,y$在薄片平面内；对三维实体还包含$z^2$项，一般没有这一等式。均匀细杆关于中心横轴的惯量为$ml^2/12$，关于端点横轴为$ml^2/3$；均匀实心球关于直径为$2mR^2/5$。换用其他质量分布时，应重新作惯量积分。


\subsection{例题与模型变式}

\begin{example}[矩形薄片的各轴惯量]\label{ex:hw6-1}

均匀矩形薄片质量$m$，两边长$a,b$，中心为$O$，$x,y$轴分别平行于$a,b$边，$z$轴垂直薄片。求三条中心主轴及两条中心对角线的转动惯量。

\end{example}


\begin{figure}[H]\centering
\begin{tikzpicture}[>=Stealth,line cap=round,line join=round]
\draw[curve](-2,-1)rectangle(2,1);\draw[axis](-2.6,0)--(2.7,0)node[right]{$x$};\draw[axis](0,-1.5)--(0,1.7)node[above]{$y$};\draw[guide](-2,-1)--(2,1);\node at(1,1.4){$a$};\node at(2.4,.6){$b$};
\end{tikzpicture}
\caption{矩形薄片的中心坐标轴与对角线}\label{fig:hw6-1-q}
\end{figure}

\begin{solution}
由对称性三个惯量积均为零。图\ref{fig:hw6-1-s}中对$x$轴，质元距轴为$|y|$，故
\[
 I_x=\frac m{ab}\int_{-a/2}^{a/2}\dd x\int_{-b/2}^{b/2}y^2\dd y
 =\frac m{ab}\bigl[x\bigr]_{-a/2}^{a/2}
 \left[\frac{y^3}{3}\right]_{-b/2}^{b/2}
 =\frac m{ab}\,a\,\frac{b^3}{12}=\frac{mb^2}{12}.
\]

\begin{figure}[H]\centering
\begin{tikzpicture}[>=Stealth,line cap=round,line join=round]
\draw[curve](-2,-1)rectangle(2,1);\draw[axis](-2.6,0)--(2.7,0)node[right]{$x$};\draw[axis](0,-1.5)--(0,1.7)node[above]{$y$};\draw[guide](-2,-1)--(2,1);\node at(1,1.4){$a$};\node at(2.4,.6){$b$};\draw[axis](0,0)--(2,1)node[above right]{$\vct n$};\node at(0,-2){$I_d=\vct n^{\mathsf T}\mathsf I\vct n$};
\end{tikzpicture}
\caption{对角线方向的惯量由主惯量投影求得}\label{fig:hw6-1-s}
\end{figure}

对$y$轴，把到轴距离改为$|x|$，积分为
\[
I_y=\frac m{ab}\left[\frac{x^3}{3}\right]_{-a/2}^{a/2}
\bigl[y\bigr]_{-b/2}^{b/2}
=\frac m{ab}\,\frac{a^3}{12}\,b=\frac{ma^2}{12}.
\]
对垂直薄片的$z$轴，距离平方为$x^2+y^2$，所以$I_z=I_x+I_y$。三条主轴的结果为
\[
 I_y=ma^2/12,\qquad I_z=m(a^2+b^2)/12.
\]
对角线单位向量为$\vct n=(a,\pm b,0)/\sqrt{a^2+b^2}$，所以
\[
 \boxed{I_d=\frac{a^2}{a^2+b^2}I_x+
 \frac{b^2}{a^2+b^2}I_y=\frac{ma^2b^2}{6(a^2+b^2)}}.
\]
例如绕方向$\vct n=(a,b,0)/\sqrt{a^2+b^2}$以角速度大小$\omega$转动，有
\[
 \vct L=\frac{m\omega}{12\sqrt{a^2+b^2}}(ab^2,ba^2,0).
\]
除$a=b$外，这个矢量与$(a,b,0)$不平行。其沿轴投影仍等于$I_d\omega$，垂直分量则决定维持定轴转动所需的约束力矩。

两对角线的惯量相同，但$a\ne b$时它们不是主轴。绕对角线转动时，$\vct L$可能有垂直于转轴的分量；$I_d\omega$只给出沿轴投影。


也可直接按到对角线的距离积分。以正斜率对角线为例，面内的垂线单位向量为$(-b,a,0)/\sqrt{a^2+b^2}$，故质元到轴的有向距离为$(-bx+ay)/\sqrt{a^2+b^2}$，于是
\[
 I_d=\frac1{a^2+b^2}\int(b^2x^2-2abxy+a^2y^2)\dd m.
\]
关于$x$轴或$y$轴的镜面对称使$\int xy\dd m=0$；具体地，对固定$x$，$\int_{-b/2}^{b/2}y\dd y=[y^2/2]_{-b/2}^{b/2}=0$。再代入已求出的$\int x^2\dd m=ma^2/12$、$\int y^2\dd m=mb^2/12$，就还原上述结果。积分中的距离为面内到斜轴的垂距，与图\ref{fig:hw6-1-s}的轴线一致。

\end{solution}



\begin{example}[正方体的体对角线与面对角线]\label{ex:hw6-2}

均匀正方体质量$m$、边长$l$。求绕一条体对角线和一条面对角线的转动惯量。

\end{example}


\begin{figure}[H]\centering
\begin{tikzpicture}[>=Stealth,line cap=round,line join=round]
\coordinate(A)at(0,0);\coordinate(B)at(2,0);\coordinate(C)at(2,2);\coordinate(D)at(0,2);\coordinate(E)at(.8,.65);\coordinate(F)at(2.8,.65);\coordinate(G)at(2.8,2.65);\coordinate(H)at(.8,2.65);\draw[curve](A)--(B)--(C)--(D)--cycle;\draw[curve](B)--(F)--(G)--(C);\draw[curve](D)--(H)--(G);\draw[guide](A)--(E)--(F);\draw[guide](E)--(H);
\fill(1.4,1.325)circle(1.2pt)node[right]{$C$};
\end{tikzpicture}
\caption{边长相同的均匀正方体}\label{fig:hw6-2-q}
\end{figure}

\begin{solution}
以正方体中心$C$为原点，三条棱平行的轴由镜面对称性构成主轴，三个主惯量因边长相等而相等。令均匀密度为$\rho=m/l^3$，先计算$x$轴惯量。内层积分为
\[
\int_{-l/2}^{l/2}(y^2+z^2)\dd z
=\left[y^2z+\frac{z^3}{3}\right]_{-l/2}^{l/2}
=ly^2+\frac{l^3}{12}.
\]
再对$y$积分，
\[
\int_{-l/2}^{l/2}\left(ly^2+\frac{l^3}{12}\right)\dd y
=\left[\frac{ly^3}{3}+\frac{l^3y}{12}\right]_{-l/2}^{l/2}
=\frac{l^4}{6}.
\]
最后对$x$积分，得到
\[
I_x=\rho\int_{-l/2}^{l/2}\frac{l^4}{6}\dd x
=\rho\frac{l^4}{6}\bigl[x\bigr]_{-l/2}^{l/2}
=\frac{ml^2}{6}=I_y=I_z.
\]
三个惯量积因被积函数对某一坐标为奇函数而为零，故惯量张量是$(ml^2/6)\vct 1$，过中心任意方向都有相同惯量。

\begin{figure}[H]\centering
\begin{tikzpicture}[>=Stealth,line cap=round,line join=round]
\begin{scope}[shift={(0,0)}]\coordinate(A)at(0,0);\coordinate(B)at(2,0);\coordinate(C)at(2,2);\coordinate(D)at(0,2);\coordinate(E)at(.8,.65);\coordinate(F)at(2.8,.65);\coordinate(G)at(2.8,2.65);\coordinate(H)at(.8,2.65);\draw[curve](A)--(B)--(C)--(D)--cycle;\draw[curve](B)--(F)--(G)--(C);\draw[curve](D)--(H)--(G);\draw[guide](A)--(E)--(F);\draw[guide](E)--(H);\draw[very thick](A)--(G);\node at(1.4,-.5){体对角线};\end{scope}\begin{scope}[shift={(4.2,0)}]\coordinate(A)at(0,0);\coordinate(B)at(2,0);\coordinate(C)at(2,2);\coordinate(D)at(0,2);\coordinate(E)at(.8,.65);\coordinate(F)at(2.8,.65);\coordinate(G)at(2.8,2.65);\coordinate(H)at(.8,2.65);\draw[curve](A)--(B)--(C)--(D)--cycle;\draw[curve](B)--(F)--(G)--(C);\draw[curve](D)--(H)--(G);\draw[guide](A)--(E)--(F);\draw[guide](E)--(H);\draw[very thick](B)--(D);\draw[guide](2.4,.325)--(.4,2.325);\draw[<->](1,1)--(1.4,1.325)node[midway,below right]{$l/2$};\node at(1.4,-.5){面对角线与中心平行轴};\end{scope}
\fill(1.4,1.325)circle(1.2pt)node[right]{$C$};\fill(5.6,1.325)circle(1.2pt)node[right]{$C$};\fill(5.2,1)circle(1.2pt)node[below left,fill=white,inner sep=1pt]{$F$};
\end{tikzpicture}
\caption{体对角线通过中心，面对角线须使用平行轴定理}\label{fig:hw6-2-s}
\end{figure}

体对角线通过中心，故$\boxed{I_{\text{体}}=ml^2/6}$。面对角线与一条过中心平行线相距$l/2$，如图\ref{fig:hw6-2-s}，平行轴定理给出
\[
 \boxed{I_{\text{面}}=\frac{ml^2}{6}+m\left(\frac l2\right)^2
 =\frac5{12}ml^2}.
\]
方向余弦给出过中心的平行轴惯量。面对角线不通过中心，计算其惯量时须加上平移项。


任意单位方向$\vct n$都有$\vct n^{\mathsf T}\vct I_C\vct n=(ml^2/6)\vct n^{\mathsf T}\vct n=ml^2/6$；体对角线取$\vct n=(1,1,1)/\sqrt3$只是其中一个方向。

对面对角线，设它所在面为$z=l/2$，轴为$y=x$。质心到该轴的垂足是这个面的中心$F=(0,0,l/2)$，因而$CF=l/2$。图\ref{fig:hw6-2-s}右图中的短距离线就是$CF$，不是正方形中心到某条对角线的面内距离。也可直接积分$(x-y)^2/2+(z-l/2)^2$，其中新增常量项积分为$ml^2/4$，与平行轴定理一致。

\end{solution}



\begin{example}[圆柱轴截面的一条对角线]\label{ex:hw6-3}

均匀实心圆柱质量$m$、半径$R$、高度$H$。$AB$是包含圆柱轴线的矩形轴截面的一条对角线，求绕$AB$的转动惯量。

\end{example}


\begin{figure}[H]\centering
\begin{tikzpicture}[>=Stealth,line cap=round,line join=round]
\draw(0,0)ellipse(.4 and 1);\draw(3,0)ellipse(.4 and 1);\draw[curve](0,1)--(3,1);\draw[curve](0,-1)--(3,-1);\draw[guide](-.5,0)--(3.5,0);\draw[very thick](0,1)--(3,-1);\node[above]at(0,1){$A$};\node[below]at(3,-1){$B$};\node at(1.5,1.4){$H$};
\end{tikzpicture}
\caption{圆柱轴截面内通过质心的对角线}\label{fig:hw6-3-q}
\end{figure}

\begin{solution}
图\ref{fig:hw6-3-s}中$AB$通过圆柱质心，与圆柱轴方向的余弦为$H/\sqrt{H^2+4R^2}$。先计算圆柱的两个中心主惯量。取轴向为$z$，横截面内用极坐标$s,\varphi$，则$\dd m=\rho s\dd s\dd\varphi\dd z$，其中$\rho=m/(\pi R^2H)$。轴向惯量为
\[
\begin{aligned}
I_\parallel&=\rho\int_{-H/2}^{H/2}\dd z\int_0^{2\pi}\dd\varphi\int_0^Rs^3\dd s\\
&=\rho\bigl[z\bigr]_{-H/2}^{H/2}\bigl[\varphi\bigr]_0^{2\pi}
\left[\frac{s^4}{4}\right]_0^R
=\rho H\cdot2\pi\cdot\frac{R^4}{4}=\frac12mR^2.
\end{aligned}
\]
取横轴为$x$轴，到该轴的距离平方为$s^2\sin^2\varphi+z^2$。角向和径向积分中，
\[
\int_0^{2\pi}\sin^2\varphi\,\dd\varphi
=\left[\frac\varphi2-\frac{\sin2\varphi}{4}\right]_0^{2\pi}=\pi,
\qquad \int_0^Rs\dd s=\left[\frac{s^2}{2}\right]_0^R=\frac{R^2}{2}.
\]
结合已求出的$\int_0^Rs^3\dd s=R^4/4$，得到
\[
\begin{aligned}
I_\perp&=\rho\int_{-H/2}^{H/2}\left(\frac{\pi R^4}{4}+\pi R^2z^2\right)\dd z\\
&=\rho\pi\left[\frac{R^4z}{4}+\frac{R^2z^3}{3}\right]_{-H/2}^{H/2}
=m\left(\frac{R^2}{4}+\frac{H^2}{12}\right).
\end{aligned}
\]
因此两个主惯量为
\[
 I_\parallel=\frac12mR^2,\qquad
 I_\perp=\frac{m}{12}(3R^2+H^2).
\]

\begin{figure}[H]\centering
\begin{tikzpicture}[>=Stealth,line cap=round,line join=round]
\draw(0,0)ellipse(.4 and 1);\draw(3,0)ellipse(.4 and 1);\draw[curve](0,1)--(3,1);\draw[curve](0,-1)--(3,-1);\draw[guide](-.5,0)--(3.5,0);\draw[very thick](0,1)--(3,-1);\node[above]at(0,1){$A$};\node[below]at(3,-1){$B$};\node at(1.5,1.4){$H$};\draw[<->](3.7,-1)--(3.7,1)node[midway,right]{$2R$};\node at(1.5,-1.6){$\cos\gamma=H/\sqrt{H^2+4R^2}$};
\end{tikzpicture}
\caption{对角线方向由圆柱高度与直径共同确定}\label{fig:hw6-3-s}
\end{figure}

轴向与横向都是主方向，故没有交叉项：
\begin{align*}
 I_{AB}&=I_\parallel\frac{H^2}{H^2+4R^2}
 +I_\perp\frac{4R^2}{H^2+4R^2}\\
 &=\boxed{\frac{mR^2(R^2+5H^2/6)}{H^2+4R^2}}.
\end{align*}

当$H\to0$时，该轴趋于圆盘直径，结果为$mR^2/4$；在$R/H\to0$的细长极限中，结果趋于$5mR^2/6$。此时尽管轴与长轴夹角很小，横向惯量乘以小角平方的贡献仍与$R^2$同阶，不能提前删去而只取$mR^2/2$。


图\ref{fig:hw6-3-s}中的对角线两端轴向间距为$H$、径向间距为$2R$，所以$\sin^2\gamma=4R^2/(H^2+4R^2)$。在主轴惯量的加权式中，$I_\parallel$应乘$\cos^2\gamma$，$I_\perp$应乘$\sin^2\gamma$。

\end{solution}



\begin{example}[均匀椭球的主惯量]\label{ex:hw6-4}

均匀实心椭球质量$m$，区域为$x^2/a^2+y^2/b^2+z^2/c^2\le1$。求三个主惯量。

\end{example}


\begin{figure}[H]\centering
\begin{tikzpicture}[>=Stealth,line cap=round,line join=round]
\draw[curve](0,0)ellipse(2 and 1.2);\draw[guide](0,0)ellipse(2 and .4);\draw[axis](0,0)--(2.6,0)node[right]{$a,x$};\draw[axis](0,0)--(0,1.8)node[above]{$c,z$};\draw[axis](0,0)--(-1.2,-1)node[left]{$b,y$};
\end{tikzpicture}
\caption{均匀椭球的三个对称方向}\label{fig:hw6-4-q}
\end{figure}

\begin{solution}
三个坐标面都是对称面，故$\int xy\dd m=\int yz\dd m=\int zx\dd m=0$，坐标轴为主轴。图\ref{fig:hw6-4-s}的变量替换$x=a\xi,y=b\eta,z=c\zeta$把椭球映为单位球，体积元乘以$abc$。

\begin{figure}[H]\centering
\begin{tikzpicture}[>=Stealth,line cap=round,line join=round]
\begin{scope}[shift={(0,0)}]\draw[curve](0,0)ellipse(2 and 1.2);\draw[guide](0,0)ellipse(2 and .4);\draw[axis](0,0)--(2.6,0)node[right]{$a,x$};\draw[axis](0,0)--(0,1.8)node[above]{$c,z$};\draw[axis](0,0)--(-1.2,-1)node[left]{$b,y$};\end{scope}\begin{scope}[shift={(6,0)}]\draw(0,0)circle(1.2);\draw[guide](0,0)ellipse(1.2 and .35);\node at(0,-1.7){$\xi^2+\eta^2+\zeta^2\le1$};\draw[axis](-2.5,.5)--(-1.5,.5);\node at(-2,1.1){$x=a\xi$};\end{scope}
\end{tikzpicture}
\caption{椭球通过坐标缩放映为单位球}\label{fig:hw6-4-s}
\end{figure}

单位球的球对称性给出三个平方坐标积分相等，计算其中一个即可。取$\eta=r\sin\vartheta\sin\varphi$，体积元为$r^2\sin\vartheta\dd r\dd\vartheta\dd\varphi$，于是
\[
\int\eta^2\dd V
=\int_0^1r^4\dd r\int_0^\pi\sin^3\vartheta\,\dd\vartheta
\int_0^{2\pi}\sin^2\varphi\,\dd\varphi.
\]
径向积分为$[r^5/5]_0^1=1/5$。第二个积分令$t=\cos\vartheta$，则$\dd t=-\sin\vartheta\dd\vartheta$，上下限由$0,\pi$变为$1,-1$，故
\[
\int_0^\pi\sin^3\vartheta\,\dd\vartheta
=-\int_1^{-1}(1-t^2)\dd t
=\left[t-\frac{t^3}{3}\right]_{-1}^{1}=\frac43.
\]
第三个积分使用$\sin^2\varphi=(1-\cos2\varphi)/2$，
\[
\int_0^{2\pi}\sin^2\varphi\,\dd\varphi
=\left[\frac\varphi2-\frac{\sin2\varphi}{4}\right]_0^{2\pi}=\pi.
\]
三者相乘，得$\int\eta^2\dd V=(1/5)(4/3)\pi=4\pi/15$，$\xi^2,\zeta^2$的积分相同。

单位球中$\int\eta^2\dd V=\int\zeta^2\dd V=4\pi/15$，又$m=\rho abc\,4\pi/3$，所以
\[
 I_x=\rho abc\int(b^2\eta^2+c^2\zeta^2)\dd V
 =\frac m5(b^2+c^2).
\]
循环得到
\[
 \boxed{I_x=\frac m5(b^2+c^2),\quad I_y=\frac m5(c^2+a^2),
 \quad I_z=\frac m5(a^2+b^2)}.
\]
坐标缩放同时改变距离和体积元。积分中保留密度与体积因子的乘积，才能使变换前后的质量相同。


当$a=b=c=R$时，三个结果均为$2mR^2/5$，还原均匀实心球。若仅令$c\to0$并保持质量不变，得到的极薄椭球厚度仍随面内位置变化，面密度并不均匀；因此$I_z\to m(a^2+b^2)/5$，不能改用均匀薄片的惯量。图\ref{fig:hw6-4-s}的缩放保留的是三维均匀椭球的质量分布。

\end{solution}



\begin{example}[转台上的径向收缩]\label{sup6:problem-example-3}\label{sup6:ex:d5-e3}
水平转台绕固定竖直轴O无摩擦转动，本体惯量为 $I_0$。质量m的小滑块由径向导轨与转台保持相同方位角，并由安装在转台上的电机从半径a移到b，$0<b<a$。初、末径向速度均为零，初始角速度为 $\omega_0$。忽略电机及传动件的质量变化，求末角速度和电机对机械系统所做的功，并由径向力积分验证结果。
\end{example}
\begin{figure}[H]\centering
\begin{tikzpicture}
\draw[curve](0,0)circle(2);\draw[guide](0,0)circle(1);\draw[curve](0,0)--(2,0);\fill(0,0)circle(1.5pt)node[below left]{$O$};
\fill(1,0)circle(2pt)node[below]{$b$};\draw[fill=white](1.65,-.16)rectangle(1.97,.16);\node[right]at(2.2,.2){$m,\ r=a$};
\draw[axis](1.6,.55)--(.9,.55)node[midway,above]{收缩};\draw[axis](-1.2,1.2)arc(135:210:1.7);\node[left]at(-1.8,.5){$\omega_0$};
\node[right]at(2.3,-1.3){本体惯量 $I_0$};
\end{tikzpicture}
\caption{导轨给滑块切向约束力，电机沿径向驱动。轴承不提供轴向力矩。}\label{sup6:fig:d5-e3-q}
\end{figure}
\begin{solution}
以转台和滑块为机械系统。图\ref{sup6:fig:d5-e3-q}中的径向驱动力对O的力矩为零，导轨切向力是系统内力，故
\[
H=(I_0+ma^2)\omega_0=(I_0+mr^2)\omega(r).
\]
令 $I_a=I_0+ma^2$、$I_b=I_0+mb^2$，则 $\omega_b=I_a\omega_0/I_b$。收缩使惯量减小，角速度增大。
径向速度为零的两个端态，其动能分别为 $H^2/(2I_a)$、$H^2/(2I_b)$。没有摩擦耗散，电机向机械运动输入的功为
\[
W=\frac{H^2}{2}\left(\frac1{I_b}-\frac1{I_a}\right)\ge0.
\]
初始有转动时$H\ne0$，上式严格为正，收缩所增加的转动能由电机供给；轴向角动量在此过程中仍保持不变。若$\omega_0=0$，整个过程没有转动，两端径向速度又为零，净机械功为零。
\begin{figure}[H]\centering
\begin{tikzpicture}
\draw[curve](0,0)--(3.5,0);\fill(0,0)circle(1.6pt)node[below]{$O$};\fill(2.4,0)circle(2pt)node[above]{$m$};
\draw[axis](2.4,0)--(1.3,0)node[midway,below]{$F_r<0$};\draw[axis](2.4,0)--(2.4,-1)node[right]{$F_\varphi<0$};\node at(1.1,.8){$r$};
\begin{scope}[shift={(5.3,-.5)}]
\draw[axis](0,0)--(4,0)node[right]{$r$};\draw[axis](0,0)--(0,2.4)node[above]{$K_{\rm rot}$};
\draw[curve,domain=.4:3.5,samples=60]plot(\x,{2.2/(1+.6*\x*\x)});
\node at(2,2.2){$H^2/[2(I_0+mr^2)]$};
\end{scope}
\end{tikzpicture}
\caption{滑块收缩时的受力与转动能。左图取 $\dot r<0$、$\omega>0$，径向力示意为向内；右图保持系统角动量H不变。}\label{sup6:fig:d5-e3-s}
\end{figure}
为验证功，取图\ref{sup6:fig:d5-e3-s}中向外的径向为正，电机力的带符号分量为 $F_r$。滑块径向方程为
\[
F_r=m(\ddot r-r\omega^2).
\]
这里的功是 $\int F_r\dot r\,\dd t$。其中加速径向运动的部分满足
\[
\int m\ddot r\dot r\,\dd t
=\left[\frac12m\dot r^2\right]_{\text{初}}^{\text{末}}.
\]
这一项只由两个端态的径向速度决定。题设初、末径向速度均为零，保证它对整个过程的积分为零，收缩途中则仍可有径向动能。
余下的转动贡献用角动量守恒式消去$\omega$，再令$z=I_0+mr^2$，则$\dd z=2mr\dd r$，上下限分别变为$I_a,I_b$。于是
\[
\begin{aligned}
W&=-mH^2\int_a^b\frac{r\dd r}{(I_0+mr^2)^2}\\
&=-\frac{H^2}{2}\int_{I_a}^{I_b}z^{-2}\dd z
=\frac{H^2}{2}\left[\frac1z\right]_{I_a}^{I_b}
=\frac{H^2}{2}\left(\frac1{I_b}-\frac1{I_a}\right).
\end{aligned}
\]
积分中用到的是初、末径向速度为零，因而该结果也适用于非准静态收缩。运动途中，系统动能还包括 $m\dot r^2/2$。
以逆时针为正，求图\ref{sup6:fig:d5-e3-s}中的切向力。对角动量守恒式$H=(I_0+mr^2)\omega$按乘积法则求导，
\[
0=(I_0+mr^2)\dot\omega+2mr\dot r\omega.
\]
移项求出$\dot\omega$，再代入极坐标切向加速度$r\dot\omega+2\dot r\omega$，得
\[
\dot\omega=-\frac{2mr\dot r\omega}{I_0+mr^2},\qquad
F_\varphi=m(r\dot\omega+2\dot r\omega)
=\frac{2mI_0\dot r\omega}{I_0+mr^2}.
\]
收缩时 $\dot r<0$、$\omega>0$，故 $F_\varphi<0$。滑块位于O右侧时，导轨给它向下的切向力。该力使滑块角动量减小，反作用力使转台本体的角动量增加，系统总角动量不变。
\end{solution}



\section{欧拉角、进动与刚体方程}

自转、进动和章动共同描述陀螺姿态。角速度在同一组基矢下合成后，才可按选定取矩点建立刚体方程，稳态进动只是其中一类运动。



\subsection{欧拉角与刚体方程}
\label{sec:complete6-euler}
设对称轴$\vct e_3$与空间竖直轴夹角为$\vartheta$，其方位角为$\varphi$，绕自身轴的相对转角为$\psi$。取$\vct e_1$在竖直轴与对称轴构成的平面内，使$\vct e_z=\cos\vartheta\vct e_3-\sin\vartheta\vct e_1$，再令$\vct e_2=\vct e_3\times\vct e_1$。角速度可分解为
\[
 \vct\omega=-\dot\varphi\sin\vartheta\vct e_1
 +\dot\vartheta\vct e_2
 +(\dot\psi+\dot\varphi\cos\vartheta)\vct e_3.
\]
这组三轴随对称轴的方向变化，却不随$\psi$自转。在完整的随体主轴上，还须将前两个分量绕$\vct e_3$旋转$\psi$。对轴对称刚体，由$I_1=I_2$可直接写
\[
 T_{\rm rot}=\frac12I_1(\dot\vartheta^2+\dot\varphi^2\sin^2\vartheta)
 +\frac12I_3(\dot\psi+\dot\varphi\cos\vartheta)^2.
\]
例\ref{ex:hw6-8}中的角度按该题图中约定；比较不同文献或坐标架时，应先对应倾角、进动角与自转角，再比较分量的正负。

随体主轴中，惯量矩阵为常量。写 $\vct L=\sum_i I_i\omega_i\vct e_i$，基矢满足 $\dot{\vct e}_i=\vct\omega\times\vct e_i$，逐项求导，
\[
\frac{\dd\vct L}{\dd t}
=\sum_i I_i\dot\omega_i\vct e_i
+\sum_i I_i\omega_i(\vct\omega\times\vct e_i).
\]
因此空间中的导数仍要计入基矢转动：
\[
 \vct M=\left(\frac{\dd\vct L}{\dd t}\right)_{\rm body}
 +\vct\omega\times\vct L.
\]
于是
\[
 \begin{aligned}
 I_1\dot\omega_1+(I_3-I_2)\omega_2\omega_3&=M_1,\\
 I_2\dot\omega_2+(I_1-I_3)\omega_3\omega_1&=M_2,\\
 I_3\dot\omega_3+(I_2-I_1)\omega_1\omega_2&=M_3.
 \end{aligned}
\]
匀速进动可使转动坐标架中的角动量分量恒定；基矢仍在转动，空间导数$\dot{\vct L}$通常不为零。例\ref{ex:hw6-5}至例\ref{ex:hw6-7}均由进动造成的角动量方向变化平衡重力矩。悬挂尖端移动时，绕质心取矩可避免把移动尖端误当固定点。


\subsection{例题与模型变式}

\begin{example}[水平圆锥陀螺的进动]\label{ex:hw6-5}

均匀实心圆锥质量$m$、高$h$、底半径$r$，锥尖$O$固定，圆锥绕对称轴以角速度$\omega$自转。存在对称轴保持水平、绕竖直轴匀速进动的状态，求进动角速度$\Omega$。

\end{example}


\begin{figure}[H]\centering
\begin{tikzpicture}[>=Stealth,line cap=round,line join=round]
\coordinate(O)at(0,0);\draw[curve](O)--(3,1)--(3,-1)--cycle;\draw(3,0)ellipse(.35 and 1);\draw[axis](0,0)--(4,0)node[right]{$\vct e_3,\omega$};\draw[axis](0,-.7)--(0,2)node[above]{$\vct e_z,\Omega$};\fill(2.25,0)circle(2pt)node[above]{$C$};\draw[axis](2.25,0)--(2.25,-1.2)node[right]{$mg$};\node[below right]at(O){$O$};
\end{tikzpicture}
\caption{尖端固定且对称轴水平的圆锥}\label{fig:hw6-5-q}
\end{figure}

\begin{solution}
图\ref{fig:hw6-5-s}中的锥尖到质心距离及尖点主惯量，已在前文“\nameref{sec:dynamics-cone-inertia}”中由薄圆盘积分求出：
\[
\ell=\frac{3h}4,\qquad I_3=\frac{3mr^2}{10},\qquad
I_1=\frac{3mh^2}5+\frac{3mr^2}{20}.
\]
这里的横向惯量包含将各薄片移到锥尖横轴的贡献，不能用质心惯量替代。
总角速度同时含自转与进动，角动量关于固定尖点可写成
\[
 \vct L_O=I_3\omega\vct e_3+I_1\Omega\vct e_z.
\]

\begin{figure}[H]\centering
\begin{tikzpicture}[>=Stealth,line cap=round,line join=round]
\coordinate(O)at(0,0);\draw[curve](O)--(3,1)--(3,-1)--cycle;\draw(3,0)ellipse(.35 and 1);\draw[axis](0,0)--(4,0)node[right]{$\vct e_3,\omega$};\draw[axis](0,-.7)--(0,2)node[above]{$\vct e_z,\Omega$};\fill(2.25,0)circle(2pt)node[above]{$C$};\draw[axis](2.25,0)--(2.25,-1.2)node[right]{$mg$};\node[below right]at(O){$O$};\draw[axis](1,1.8)--(2.4,1.8)node[right]{$L_3=I_3\omega$};\node at(3,2.5){$\dot{\vct e}_3=\Omega\vct e_z\times\vct e_3$};
\end{tikzpicture}
\caption{轴向角动量随对称轴进动}\label{fig:hw6-5-s}
\end{figure}

对称轴水平且匀速进动时，第二项在地面系中不变，第一项随轴以$\Omega$转动。因此
\[
 \dot{\vct L}_O=I_3\omega\Omega\,\vct e_z\times\vct e_3,
 \quad\vct M_O=mg\ell\,\vct e_z\times\vct e_3.
\]
相等即得
\[
 \boxed{\Omega=\frac{mg\ell}{I_3\omega}=\frac{5gh}{2r^2\omega}}.
\]
进动方向随自转方向反转。进动角动量的竖直分量在地面系中保持不变，其导数为零。因此，在精确水平的稳态下，上述关系也适用于未作快陀螺近似的情况。


这里$I_1$是过锥尖且垂直于对称轴的惯量。图\ref{fig:hw6-5-s}中$\vct e_3$水平，故$\vct\Omega\perp\vct e_3$；进动不增加自转轴上的角速度分量，$\omega$就是总角速度在$\vct e_3$上的分量。重力矩与$\dot{\vct L}_O$均沿$\vct e_z\times\vct e_3$，比较同一方向的系数即可求出$\Omega$。

\end{solution}



\begin{example}[细杆悬挂圆锥的偏斜]\label{ex:hw6-6}

上一题的圆锥尖端通过长$h/2$的无质量细杆铰接于固定点$O_0$。圆锥仍有水平自转轴，存在绕竖直方向的匀速进动状态。求$\Omega$及细杆与竖直方向夹角$\theta$；并求$h\sim r$、$\omega^2r\gg g$时的近似。

\end{example}


\begin{figure}[H]\centering
\begin{tikzpicture}[>=Stealth,line cap=round,line join=round]
\coordinate(O)at(0,0);\draw[curve](O)--(3,1)--(3,-1)--cycle;\draw(3,0)ellipse(.35 and 1);\draw[axis](0,0)--(4,0)node[right]{$\vct e_3,\omega$};\draw[axis](-.4,-.7)--(-.4,2.15)node[above]{$\vct e_z,\Omega$};\fill(2.25,0)circle(2pt)node[above]{$C$};\draw[axis](2.25,0)--(2.25,-1.2)node[right]{$mg$};\node[below right]at(O){$O$};\draw[curve](-.4,1.44568323)--(0,0);\fill(-.4,1.44568323)circle(2pt)node[left]{$O_0$};\node at(-1,.7){$h/2$};
\end{tikzpicture}
\caption{圆锥尖端由无质量细杆悬挂}\label{fig:hw6-6-q}
\end{figure}

\begin{solution}
取图\ref{fig:hw6-6-s}中圆锥轴向外为径向，细杆向内上方拉尖端，拉力为$T$。质心绕竖直轴的半径为$R_C=3h/4+(h/2)\sin\theta$。竖直与径向运动方程为
\[
 T\cos\theta=mg,\qquad T\sin\theta=m\Omega^2R_C.
\]

\begin{figure}[H]\centering
\begin{tikzpicture}[>=Stealth,line cap=round,line join=round]
\coordinate(O)at(0,0);\draw[curve](O)--(3,1)--(3,-1)--cycle;\draw(3,0)ellipse(.35 and 1);\draw[axis](0,0)--(4,0)node[right]{$\vct e_3,\omega$};\draw[axis](-.4,-.7)--(-.4,2.15)node[above]{$\vct e_z,\Omega$};\fill(2.25,0)circle(2pt)node[above]{$C$};\draw[axis](2.25,0)--(2.25,-1.2)node[right]{$mg$};\node[below right]at(O){$O$};\draw[curve](-.4,1.44568323)--(0,0);\fill(-.4,1.44568323)circle(2pt)node[left]{$O_0$};\node at(-1,.7){$h/2$};\draw[axis](0,0)--(-.4,1.44568323)node[midway,right,xshift=8pt,yshift=-5pt]{$T$};\draw(-.4,.89568323)arc(-90:-74.534:.55)node[midway,right]{$\theta$};\draw[<->](-.4,-1.8)--(2.25,-1.8)node[midway,below]{$R_C$};
\end{tikzpicture}
\caption{竖直进动轴通过悬点，拉力提供向心力}\label{fig:hw6-6-s}
\end{figure}

绕圆锥质心取矩，细杆拉力的竖直分量产生力矩$mg(3h/4)$，与上一题相同。因此
\[
 \boxed{\Omega=\frac{5gh}{2r^2\omega},\qquad
 \tan\theta=\frac{\Omega^2}{g}
 \left(\frac{3h}{4}+\frac h2\sin\theta\right)}.
\]
绕质心的角动量包含$I_{1,C}\Omega\vct e_z$，该项在地面系中的导数为零；产生水平转矩响应的仍是$I_3\omega\vct e_3$。尖端处拉力的水平分量与$\overrightarrow{CO}$共线，不产生质心力矩，只有其竖直分量$mg$参加进动平衡。因此，本题的稳态进动仍满足上一题的$\Omega$关系。

这是偏角的隐式方程。快自转时$\theta\ll1$，且$\Omega^2h/g\ll1$，保留最低阶得
\[
 \boxed{\theta\simeq\frac{3h\Omega^2}{4g}
 =\frac{75gh^3}{16r^4\omega^2}}.
\]
细杆的微小偏斜使拉力具有水平分量，由它提供质心作圆周运动所需的向心力。


令$\epsilon=\Omega^2h/g$，偏角方程写成$\tan\theta=\epsilon(3/4+\sin\theta/2)$。这个方程在$0<\theta<\pi/2$中恰有一解：函数
\[
f(\theta)=\frac{\tan\theta}{3/4+\sin\theta/2}
\]
从零连续增至无穷。按商法则求导，
\[
f'(\theta)=\frac{\sec^2\theta(3/4+\sin\theta/2)-\tan\theta\cos\theta/2}
{(3/4+\sin\theta/2)^2}
=\frac{\tfrac34\sec^2\theta+\tfrac12\sin\theta\tan^2\theta}
{(3/4+\sin\theta/2)^2}>0.
\]

对于小偏角支，设$\theta=A\epsilon+B\epsilon^2+O(\epsilon^3)$。用$\tan\theta=\theta+O(\theta^3)$、$\sin\theta=\theta+O(\theta^3)$代入，比较到$\epsilon^2$阶：
\[
A\epsilon+B\epsilon^2
=\frac34\epsilon+\frac A2\epsilon^2+O(\epsilon^3).
\]
同次项系数相等，得$A=3/4$、$B=A/2=3/8$，即
\[
 \theta=\frac34\epsilon+\frac38\epsilon^2+O(\epsilon^3).
\]
最低阶即为上式；第二项来自细杆偏斜后质心回转半径的增加。图\ref{fig:hw6-6-s}的$R_C$从竖直进动轴量到质心；若只取尖端到质心的$3h/4$，便会遗漏这一反馈。又$T=mg/\cos\theta$，所以拉力比$mg$大的修正从$\theta^2$阶开始，不影响最低阶偏角。

\end{solution}



\begin{example}[倾斜圆锥的稳态进动双支]\label{ex:hw6-7}

均匀实心圆锥的尖端固定，对称轴与水平夹角为$\theta$，相对随轴进动坐标架的自转角速度为$\omega$。求保持倾角不变的匀速进动角速度$\Omega$，说明两支解的条件。

\end{example}


\begin{figure}[H]\centering
\begin{tikzpicture}[>=Stealth,line cap=round,line join=round]
\draw[axis](0,-.4)--(0,2.5)node[above]{$z,\Omega$};\draw[axis](0,0)--(3.5,0)node[right]{水平};\draw[curve](0,0)--(2.83,.2)--(2.17,1.8)--cycle;\draw[axis](0,0)--(3,1.2)node[right]{$\vct e_3,\omega$};\draw(.8,0)arc(0:21.801409:.8);\node at(1.2,.25){$\theta$};\fill(1.875,.75)circle(2pt)node[above]{$C$};
\end{tikzpicture}
\caption{对称轴与水平成固定倾角的圆锥}\label{fig:hw6-7-q}
\end{figure}

\begin{solution}
关于尖点的主惯量为$I_3=3mr^2/10$、$I_1=3mr^2/20+3mh^2/5$，质心距尖点$\ell=3h/4$。图\ref{fig:hw6-7-s}中总角速度为$\vct\omega_{\text{总}}=\omega\vct e_3+\Omega\vct e_z$，沿对称轴的总分量是$\omega+\Omega\sin\theta$。

\begin{figure}[H]\centering
\begin{tikzpicture}[>=Stealth,line cap=round,line join=round]
\draw[axis](0,-.4)--(0,2.5)node[above]{$z,\Omega$};\draw[axis](0,0)--(3.5,0)node[right]{水平};\draw[curve](0,0)--(2.83,.2)--(2.17,1.8)--cycle;\draw[axis](0,0)--(3,1.2)node[right]{$\vct e_3,\omega$};\draw(.8,0)arc(0:21.801409:.8);\node at(1.2,.25){$\theta$};\fill(1.875,.75)circle(2pt)node[above]{$C$};\draw[axis](1.875,.75)--(1.875,-.5)node[right]{$mg$};\node at(3,2.4){$\omega_3=\omega+\Omega\sin\theta$};
\end{tikzpicture}
\caption{总轴向角速度包含进动的投影}\label{fig:hw6-7-s}
\end{figure}

由惯量张量
\[
 \vct L=I_1\Omega\vct e_z+
 [I_3\omega+(I_3-I_1)\Omega\sin\theta]\vct e_3.
\]
稳态中方括号内的系数恒定，$\vct e_z$也固定，只有$\vct e_3$随进动改变。由$\dot{\vct e}_3=\Omega\vct e_z\times\vct e_3$，
\[
\dot{\vct L}
=\Omega[I_3\omega+(I_3-I_1)\Omega\sin\theta]
\vct e_z\times\vct e_3.
\]
重力矩为$\vct M=\ell\vct e_3\times(-mg\vct e_z)=mg\ell\vct e_z\times\vct e_3$。在$\cos\theta\ne0$时这个共同方向矢量非零，令两系数相等，得到
\[
 (I_3-I_1)\sin\theta\,\Omega^2+I_3\omega\Omega-mg\ell=0.
\]
故在$A=(I_3-I_1)\sin\theta\ne0$时
\[
 \boxed{\Omega_\pm=\frac{-I_3\omega\pm
 \sqrt{I_3^2\omega^2+4Amg\ell}}{2A}}.
\]
须满足判别式非负。$A=0$且$\omega\ne0$时，方程退化为线性方程，解为$\Omega=mg\ell/(I_3\omega)$；若$A=0$且$\omega=0$，则不能平衡非零重力矩，不存在这一非竖直稳态。

当$A\ne0$固定，且$|4Amg\ell/(I_3^2\omega^2)|\ll1$时，可分别求两支的快自转近似。慢支取$\Omega_s=O(\omega^{-1})$，则$A\Omega_s^2$相对于常数项是高阶小量，方程最低阶给出
\[
 I_3\omega\Omega_s\simeq mg\ell,
 \qquad \Omega_s\simeq\frac{mg\ell}{I_3\omega}.
\]
两根之和为$\Omega_s+\Omega_f=-I_3\omega/A$，所以快支满足
\[
 \Omega_f=-\frac{I_3\omega}{A}-\Omega_s
 \simeq-\frac{I_3\omega}{A}.
\]

例如细长圆锥常有$I_1>I_3$，且轴在水平面上方时$\sin\theta>0$，于是$A<0$。存在两实支须有
\[
 I_3^2\omega^2\ge4(I_1-I_3)mg\ell\sin\theta.
\]
等号时两支合并；低于该自转阈值时，不存在这一固定倾角的稳态。这些根给出固定倾角稳态。要判断稳态对章动扰动是否稳定，还须列出倾角的扰动方程。

如果改把$\omega$定义为总角速度沿对称轴的投影，二次项就应相应改写。两种自转角速度的定义不能混用；本题在题设中已固定为相对进动架的自转。
上述推导适用于$\cos\theta\ne0$。轴竖直时，进动角不再是独立坐标，应直接用欧拉方程描述。

\end{solution}



\begin{example}[翻身陀螺的欧拉角与接触动力学]\label{ex:hw6-8}

把翻身陀螺理想化为轴对称刚体，外形接触面为半径$R$的球面，质量$m$。球心$O$相对质心$C$沿对称轴$\vct e_3$偏移$\alpha R$，主惯量为$I_1=I_2,I_3$。接触点$A$在水平地面上滑动，动摩擦为$-\mu N\vct v_A/|\vct v_A|$。用欧拉角描述运动，求外力、质心力矩、接触约束、角速度、能量及竖直角动量的变化率，并证明$\vct L\times\vct e_3=I_1(\vct\omega\times\vct e_3)$。

\end{example}


\begin{figure}[H]\centering
\begin{tikzpicture}[>=Stealth,line cap=round,line join=round]
\draw(0,0)circle(1.3);\draw[thick](-2,-1.3)--(2,-1.3);\coordinate(C)at(-.2,-.25);\coordinate(A)at(0,-1.3);\fill(C)circle(2pt)node[left]{$C$};\fill(0,0)circle(1pt)node[below right]{$O$};\fill(A)circle(2pt)node[below]{$A$};\draw[axis](C)--(.8,1)node[above]{$\vct e_3$};\draw[axis](C)--(-.2,1.6)node[left]{$\vct e_z$};\draw[guide](C)--(A);
\draw(-.2,.4)arc(90:51.34019:.65)node[midway,above]{$\theta$};
\end{tikzpicture}
\caption{球面接触的翻身陀螺及偏置质心}\label{fig:hw6-8-q}
\end{figure}

\begin{solution}
图\ref{fig:hw6-8-s}固定竖直单位矢量$\vct e_z$，以$\theta$为$\vct e_3$对竖直的夹角、$\phi$为绕竖直的方位角、$\psi$为绕对称轴的自转角。接触几何为
\[
 \vct a:=\overrightarrow{CA}=R(\alpha\vct e_3-\vct e_z),\quad
 z_C=R(1-\alpha\cos\theta),\quad \vct v_A=\vct v_C+\vct\omega\times\vct a.
\]

\begin{figure}[H]\centering
\begin{tikzpicture}[>=Stealth,line cap=round,line join=round]
\draw(0,0)circle(1.3);\draw[thick](-2,-1.3)--(2,-1.3);\coordinate(C)at(-.2,-.25);\coordinate(A)at(0,-1.3);\fill(C)circle(2pt)node[left]{$C$};\fill(0,0)circle(1pt)node[below right]{$O$};\fill(A)circle(2pt)node[below]{$A$};\draw[axis](C)--(.8,1)node[above]{$\vct e_3$};\draw[axis](C)--(-.2,1.6)node[left]{$\vct e_z$};\draw[guide](C)--(A);\draw[axis](0,-1.3)--(0,-.4)node[below right]{$N$};\draw[axis](0,-1.3)--(1,-1.3)node[below]{$\vct f$};\draw[axis](-.2,-.25)--(-.2,-1.1)node[midway,left]{$mg$};
\draw(-.2,.4)arc(90:51.34019:.65)node[midway,above]{$\theta$};\draw[axis](0,-1.95)--(-.9,-1.95)node[below]{$\vct v_A$};
\end{tikzpicture}
\caption{接触点的法向力、摩擦与质心重力}\label{fig:hw6-8-s}
\end{figure}

\textbf{外力与约束。}令$\vct f=-\mu N\vct v_A/|\vct v_A|$，则
\[
 m\dot{\vct v}_C=(N-mg)\vct e_z+\vct f,\quad
 \dot{\vct L}_C=\vct a\times(N\vct e_z+\vct f).
\]
接触要求$\vct v_A\cdot\vct e_z=0$，故$\vct v_A=v_x\vct e_x+v_y\vct e_y$。对$z_C=R(1-\alpha\cos\theta)$逐次求导，
\[
\dot z_C=R\alpha\sin\theta\dot\theta,\qquad
\ddot z_C=R\alpha(\cos\theta\dot\theta^2+\sin\theta\ddot\theta).
\]
代入竖直平动方程$N-mg=m\ddot z_C$，得
\[
 N=mg+mR\alpha(\cos\theta\dot\theta^2+\sin\theta\ddot\theta).
\]
这些式子只适用于$N\ge0$的接触阶段。$\vct v_A=0$时动摩擦方向不再有定义，须改用静摩擦条件，不能继续除以零。

\textbf{欧拉角与角动量。}在由竖直轴、节点线和对称轴依次构造的坐标架中，角速度矢量为
\[
 \vct\omega=\dot\phi\vct e_z+\dot\theta\vct e_\phi+\dot\psi\vct e_3,
 \quad\omega_3=\dot\psi+\dot\phi\cos\theta.
\]
这里$\vct e_\phi=(\vct e_z\times\vct e_3)/\sin\theta$沿方位增加方向，$\dot{\vct e}_3=\vct\omega\times\vct e_3$。欧拉角在$\sin\theta=0$处坐标退化，但矢量形式的接触几何和动力学仍可继续使用。

取随体主轴$\vct e_1,\vct e_2,\vct e_3$，并约定$\psi=0$时$\vct e_1=\vct e_\phi$、$\vct e_2=\vct e_3\times\vct e_\phi$；随后两轴绕$\vct e_3$正向转过$\psi$。按这一零位约定，角速度的随体分量为
\[
 \omega_1=\dot\theta\cos\psi+\dot\phi\sin\theta\sin\psi,\quad
 \omega_2=-\dot\theta\sin\psi+\dot\phi\sin\theta\cos\psi.
\]
轴对称性给出$\vct L=I_1\vct\omega+(I_3-I_1)\omega_3\vct e_3$，两边与$\vct e_3$叉乘，立即得到所要求的正比式。它表明垂直于对称轴的角动量与角速度具有相同方向。

\textbf{能量与耗散。}保留质心平动，总机械能为
\[
 E=\frac m2|\vct v_A-\vct\omega\times\vct a|^2
 +\frac{I_1}{2}(\dot\theta^2+\dot\phi^2\sin^2\theta)
 +\frac{I_3}{2}(\dot\psi+\dot\phi\cos\theta)^2
 +mgR(1-\alpha\cos\theta).
\]
重力已计入势能，平动与转动动能的变化率相加为
\[
\begin{aligned}
\dot E&=(N\vct e_z+\vct f)\cdot\vct v_C
+\vct\omega\cdot[\vct a\times(N\vct e_z+\vct f)]\\
&=(N\vct e_z+\vct f)\cdot(\vct v_C+\vct\omega\times\vct a)
=(N\vct e_z+\vct f)\cdot\vct v_A.
\end{aligned}
\]
第二步用了混合积的循环关系。接触约束给出$\vct e_z\cdot\vct v_A=0$，再代入摩擦定律，因此
\[
 \boxed{\dot E=\vct f\cdot\vct v_A=-\mu N|\vct v_A|\le0}.
\]
竖直角动量为$L_z=I_1\dot\phi\sin^2\theta+I_3\omega_3\cos\theta$，其变化率为
\[
 \boxed{\dot L_z=R\alpha(\vct e_3\times\vct f)\cdot\vct e_z}.
\]
还可直接分解竖直角动量的变化：法向力矩的竖直分量为零，摩擦力矩中的$-R\vct e_z\times\vct f$也没有竖直分量，因此只剩$R\alpha\vct e_3\times\vct f$。机械能减少与$L_z$变化来自同一个接触摩擦，但二者没有固定的比例关系。

最后考虑$\lambda_J=-\vct L\cdot\vct a=R(L_z-\alpha I_3\omega_3)$。按乘积法则求导，
\[
\dot\lambda_J=-\dot{\vct L}\cdot\vct a-\vct L\cdot\dot{\vct a}.
\]
第一项中$\dot{\vct L}=\vct a\times(N\vct e_z+\vct f)$与$\vct a$垂直，故为零。第二项用$\dot{\vct a}=R\alpha\vct\omega\times\vct e_3$和轴对称角动量表达式，
\[
\vct L\cdot\dot{\vct a}
=R\alpha[I_1\vct\omega+(I_3-I_1)\omega_3\vct e_3]
\cdot(\vct\omega\times\vct e_3)=0.
\]
因两个基矢$\vct\omega,\vct e_3$都与其叉积垂直，得到$\dot\lambda_J=0$。摩擦耗能时，这一积分仍保持常量，并与能量变化共同限制运动。是否发生翻转还取决于初态及其后的动力学演化。


\end{solution}



\begin{example}[翻转势垒与临界时间]\label{deep6:problem-example-10}\label{deep6:ex:d5-e10}
质量M、宽 $2b$、高 $2h$ 的匀质矩形薄板在竖直平面内绕右下角O的固定光滑铰轴转动。初始板竖直，给定顺时针角速度 $\omega_0>0$，此后只受重力与铰链力。求质心越过O正上方的角速度阈值，并判断在阈值等号处能否有限时间越过。仅讨论越过最高质心位置之前的运动。
\end{example}
\begin{figure}[H]\centering
\begin{tikzpicture}
\draw[guide](-2.2,0)rectangle(0,3);\draw[curve,rotate=-20](-2.2,0)rectangle(0,3);
\fill(0,0)circle(2pt)node[below right]{$O$};\fill(-.521,1.786)circle(1.7pt)node[left]{$C$};
\draw[guide](0,0)--(-.521,1.786);\draw[axis](-.521,1.786)--(-.521,.8)node[left]{$Mg$};
\draw[guide](0,0)--(0,3.3);\draw(0,.65)arc(90:70:.65);\node at(.25,.9){$\varphi$};
\node[left]at(-2.2,1.3){$2h$};\node[below]at(-1.1,0){$2b$};\node[right]at(1.2,2.3){虚线为初始构形};
\end{tikzpicture}
\caption{O为固定铰轴，不把它等同于仅靠地面摩擦维持的自由接触角点。}\label{deep6:fig:d5-e10-q}
\end{figure}
\begin{solution}
按图\ref{deep6:fig:d5-e10-q}，顺时针转角为 $\varphi$，质心坐标为
\[
x_C=-b\cos\varphi+h\sin\varphi,\qquad
y_C=b\sin\varphi+h\cos\varphi.
\]
令 $d=\sqrt{b^2+h^2}$，最高点满足 $x_C=0$，故 $\varphi_c=\arctan(b/h)$、$y_{C,\max}=d$。用平行轴定理，
\[
I_O=\frac M{12}\big[(2b)^2+(2h)^2\big]+Md^2
=\frac43Md^2.
\]
取质心初始高度h为势能零点，能量方程为
\[
\frac12I_O\dot\varphi^2+Mg(y_C-h)=\frac12I_O\omega_0^2.
\]
势垒高度为 $Mg(d-h)$，因此阈值为
\[
\omega_c=\sqrt{\frac{2Mg(d-h)}{I_O}}
=\sqrt{\frac{3g(d-h)}{2d^2}}.
\]
若 $\omega_0<\omega_c$，板在最高点之前转向；若 $\omega_0>\omega_c$，则以非零角速度越过。
\begin{figure}[H]\centering
\begin{tikzpicture}
\draw[axis](0,0)--(6,0)node[right]{$\varphi$};\draw[axis](0,0)--(0,3.1)node[above]{$U$};
\draw[curve,domain=0:1.2,samples=80]plot({4*\x},{2*(cos((\x-.6)*180/pi)-cos(.6*180/pi))/(1-cos(.6*180/pi))});
\draw[guide](0,2)--(5.5,2)node[right]{$E_c$};\draw[guide](2.4,0)--(2.4,2);\node[below]at(2.4,0){$\varphi_c$};
\node at(4.1,2.9){临界能量};
\end{tikzpicture}
\caption{临界能量在势垒顶与势能曲线相切，速度随距顶点的角距离线性趋零。}\label{deep6:fig:d5-e10-s}
\end{figure}
在等号处，图\ref{deep6:fig:d5-e10-s}中 $E=E_c$。置 $\delta=\varphi_c-\varphi>0$，利用 $y_C=d\cos\delta$，
\[
\frac12I_O\dot\varphi^2=Mg d(1-\cos\delta)
=2Mgd\sin^2\frac\delta2.
\]
记$\lambda=\sqrt{Mgd/I_O}$，上升阶段$\dot\varphi>0$，而$\dot\delta=-\dot\varphi$，故$\dot\delta=-2\lambda\sin(\delta/2)$。令$q=\tan(\delta/4)$，用
$\dot\delta=4\dot q/(1+q^2)$、$\sin(\delta/2)=2q/(1+q^2)$代入，方程化为$\dot q=-\lambda q$。初始$\delta_0=\varphi_c$，所以$q_0=\tan(\varphi_c/4)$。分离变量积分，
\[
\int_{q_0}^{q(t)}\frac{\dd q}{q}=-\lambda\int_0^t\dd s,
\qquad \ln\frac{q(t)}{q_0}=-\lambda t.
\]
于是临界解及到达任一剩余角距离$\delta>0$的时间为
\[
\delta(t)=4\arctan\!\left[\tan\frac{\varphi_c}{4}\,\ee^{-\lambda t}\right],
\qquad
t=\frac1\lambda\ln\frac{\tan(\varphi_c/4)}{\tan(\delta/4)}.
\]
当$\delta\to0^+$时，$\tan(\delta/4)\sim\delta/4$，时间按对数趋于无穷。这也与图\ref{deep6:fig:d5-e10-s}中势垒顶附近剩余动能正比于$\delta^2$一致。因此$\omega_c$是有限时间翻转所需初始角速度的下确界；取等号时板渐近接近最高点，不能在有限时间越过。
\end{solution}



\section{平面刚体、运动点取矩与滚动}

平面刚体须同时满足质心平动和绕质心转动方程。运动点取矩需要相应修正，滚动约束则直接联系两侧接触材料点的速度。



\subsection{平面运动中的运动点取矩}
若$A$是刚体材料点，记$\vct r_{CA}=\vct r_C-\vct r_A$，方向由$A$指向质心$C$。采用相对$A$速度定义的角动量为$I_A\omega$，且$I_A$不随时间改变。将第三讲已推导的运动点相对角动量定理用于平面刚体，得到
\[
 \label{sup6:eq:d5-material-origin} I_A\dot\omega=M_A-[\vct r_{CA}\times m\vct a_A]_z.
\]
这给出本讲拓展例题中移动杆端的求解方法。瞬时转动中心可用于求速度和动能；要直接绕该点列角动量定理，还须满足相应取矩点的运动条件。



\LectureNeedspace{4\baselineskip}


\subsection{例题与模型变式}

\begin{example}[运动杆端取矩]\label{sup6:problem-example-1}\label{sup6:ex:d5-e1}
质量M、长L的匀质细杆AB静置在光滑水平桌面上。现对B施加桌面内垂直于杆、大小为F的力。求开始作用瞬间的质心加速度、角加速度及两端加速度，并分别绕C、A列转动方程。忽略空气阻力。
\end{example}
\begin{figure}[H]\centering
\begin{tikzpicture}
\draw[curve](0,0)--(6,0);\foreach\x/\s in{0/A,3/C,6/B}{\fill(\x,0)circle(1.6pt)node[below]{$\s$};}
\draw[axis](6,0)--(6,1.6)node[right]{$F$};\draw[<->](0,-.8)--(6,-.8)node[midway,below]{$L$};
\draw[axis](0,1.8)--(1.2,1.8)node[right]{$x$};\draw[axis](0,1.8)--(0,2.8)node[above]{$y$};
\end{tikzpicture}
\caption{俯视桌面。力沿正y方向，重力与桌面支持力不影响面内运动。}\label{sup6:fig:d5-e1-q}
\end{figure}
\begin{solution}
以图\ref{sup6:fig:d5-e1-q}中的正y方向为加速度正向，转动取逆时针为正。杆受合力F，故 $a_{Cy}=F/M$。绕质心的力矩为 $FL/2$，而 $I_C=ML^2/12$，所以
\[
\alpha=\frac{FL/2}{ML^2/12}=\frac{6F}{ML}.
\]
开始作用瞬间 $\omega=0$，刚体点加速度关系中没有向心项。由 $\vct a_P=\vct a_C+\vct\alpha\times\vct r_{PC}$，
\[
a_{Ay}=\frac FM-\frac L2\alpha=-\frac{2F}M,
\qquad a_{By}=\frac FM+\frac L2\alpha=\frac{4F}M.
\]
两端的x分量均为零。A向下、B向上，这一加速度分布由质心平移与绕质心转动叠加而成。
\begin{figure}[H]\centering
\begin{tikzpicture}
\draw[curve](0,0)--(6,0);\node[above]at(0,0){$A$};\node[below]at(3,0){$C$};\node[below]at(6,0){$B$};
\draw[axis](0,0)--(0,-1)node[left]{$-2F/M$};\draw[axis](3,0)--(3,.5)node[right]{$F/M$};\draw[axis](6,0)--(6,2)node[right]{$4F/M$};
\draw[axis](3.983,.688)arc(35:145:1.2);\node at(3,1.6){$\alpha>0$};
\end{tikzpicture}
\caption{释放瞬间的三点加速度。材料点A虽瞬时静止，加速度不为零。}\label{sup6:fig:d5-e1-s}
\end{figure}
图\ref{sup6:fig:d5-e1-s}中的A随杆运动，关于A的力矩要保留原点加速度修正。释放瞬间 $\omega=0$，将质心方程移到A，可写成
\[
\begin{aligned}
\tau_A&=I_C\alpha+\big(\vct r_{CA}\times M\vct a_C\big)_z,\\
\vct a_C&=\vct a_A+\vct\alpha\times\vct r_{CA}.
\end{aligned}
\]
代入后，与 $\alpha$ 成正比的两项合为 $(I_C+ML^2/4)\alpha=I_A\alpha$；含 $\vct a_A$ 的一项保留下来，反映A本身的加速运动。
绕图\ref{sup6:fig:d5-e1-s}的A取矩时，$I_A=ML^2/3$，须用式\eqref{sup6:eq:d5-material-origin}：
\[
I_A\alpha=FL-\left(\frac L2\vct e_x\times M\vct a_A\right)_z
=FL+FL=2FL.
\]
由此仍得 $\alpha=6F/(ML)$。若误把A当作固定铰点，写成 $I_A\alpha=FL$，角加速度就会少一半。上述结果只针对开始作用的瞬间；随后杆转动，必须保留速度平方项并说明F的方向是否随杆改变。
\end{solution}



\begin{example}[铰接杆的速度与支座力]\label{sup6:problem-example-2}\label{sup6:ex:d5-e2}
质量M、长L的匀质杆一端O固定铰接，在竖直平面内从水平位置静止释放。以杆从初始位置向下转过的角为 $\theta$，求 $0\le\theta\le\pi/2$ 内的角速度、角加速度和铰链力，并求铰链力大小的最大值。
\end{example}
\begin{figure}[H]\centering
\begin{tikzpicture}
\coordinate(O)at(0,0);\coordinate(B)at(4,-2.8);\draw[guide](0,0)--(4.7,0);\draw[curve](O)--(B);
\fill(O)circle(2pt)node[above left]{$O$};\fill(2,-1.4)circle(1.6pt)node[above]{$C$};\node[right]at(B){$B$};
\draw(1,0)arc(0:-35:1);\node at(1.3,-.35){$\theta$};\draw[axis](2,-1.4)--(2,-2.6)node[left]{$Mg$};
\draw[axis](-.2,-.2)--(-.2,-1.3)node[left]{$y$};\draw[axis](.2,.3)--(1.3,.3)node[right]{$x$};
\end{tikzpicture}
\caption{取x向右、y向下；$\theta$ 的正向为顺时针。}\label{sup6:fig:d5-e2-q}
\end{figure}
\begin{solution}
图\ref{sup6:fig:d5-e2-q}中，质心下降 $L\sin\theta/2$。固定铰链不做功，能量守恒给出
\[
\frac12\frac{ML^2}{3}\omega^2=Mg\frac L2\sin\theta,
\qquad \omega^2=\frac{3g}L\sin\theta.
\]
重力对O的顺时针力矩为 $MgL\cos\theta/2$，因此
\[
\alpha=\frac{3g}{2L}\cos\theta.
\]
此式由力矩方程得到，在刚释放、$\omega=0$ 时也成立。
图\ref{sup6:fig:d5-e2-s}中，切向加速度垂直OC，法向加速度沿CO。两者都要投影到图\ref{sup6:fig:d5-e2-q}的固定坐标轴；随杆转动的方向本身也随时间改变。
质心坐标为 $x_C=(L/2)\cos\theta$、$y_C=(L/2)\sin\theta$。一次求导给出 $\dot x_C=-(L/2)\omega\sin\theta$、$\dot y_C=(L/2)\omega\cos\theta$。再用 $\dot\omega=\alpha$ 对乘积求导，得到含 $\alpha$ 的切向项与含 $\omega^2$ 的法向项：
\[
\begin{aligned}
a_{Cx}&=-\frac L2(\alpha\sin\theta+\omega^2\cos\theta)
=-\frac94g\sin\theta\cos\theta,\\
a_{Cy}&=\frac L2(\alpha\cos\theta-\omega^2\sin\theta)
=\frac34g(1-3\sin^2\theta).
\end{aligned}
\]
\begin{figure}[H]\centering
\begin{tikzpicture}
\draw[curve](0,0)--(4,-2.8);\fill(0,0)circle(2pt)node[above right]{$O$};\fill(2,-1.4)circle(1.6pt)node[right]{$C$};
\draw[axis](0,0)--(-1.2,0)node[below]{$H_x<0$};\draw[axis](0,0)--(0,1.1)node[left]{$H_y<0$};
\draw[axis](2,-1.4)--(.8,-.56);\node at(1.65,-.32){$\omega^2L/2$};
\draw[axis](2,-1.4)--(1.35,-2.328571)node[below left]{$\alpha L/2$};
\node[align=left]at(6,-1){法向加速度指向O\\切向加速度沿下降方向};
\end{tikzpicture}
\caption{铰链力及质心的法向、切向加速度。图中为中间构形。}\label{sup6:fig:d5-e2-s}
\end{figure}
设铰链对杆的力分量为 $H_x,H_y$。按图\ref{sup6:fig:d5-e2-s}的坐标正向，
\[
H_x=Ma_{Cx}=-\frac94Mg\sin\theta\cos\theta,
\qquad H_y+Mg=Ma_{Cy}.
\]
所以 $H_y=-Mg(1+9\sin^2\theta)/4$，负号表示向上。力的大小为
\[
H^2=\frac{M^2g^2}{16}
\left[81\sin^2\theta\cos^2\theta+(1+9\sin^2\theta)^2\right]
=\frac{M^2g^2}{16}(1+99\sin^2\theta).
\]
它在此区间单调增大，最大值在杆竖直向下时达到，为 $H_{\max}=5Mg/2$；刚释放时为 $Mg/4$。最低点虽有 $\alpha=0$，向心加速度仍为 $3g/2$，因此支座力并不等于Mg。
\end{solution}



\begin{example}[绕线方向与线轴的运动]\label{sup6:problem-example-5}\label{sup6:ex:d5-e5}
线轴质量M，外半径R，绕线轴半径r，$0<r<R$，转动惯量 $I_C=\beta MR^2$。线从内轴下侧切向引出，以大小F、与水平向右方向成 $\theta$ 角的力拉动，$0\le\theta\le\pi/2$。线轴在水平地面上纯滚动。求质心加速度、静摩擦力及纯滚动条件，并说明何时线轴向左开始运动。
\end{example}
\begin{figure}[H]\centering
\begin{tikzpicture}
\draw[curve](-2,-1.3)--(3,-1.3);\draw[curve](0,0)circle(1.3);\draw[curve](0,0)circle(.65);\fill(0,0)circle(1.5pt)node[above right]{$C$};
\coordinate(Q)at({.65*sin(35)},{-.65*cos(35)});\fill(Q)circle(1.5pt)node[below left]{$Q$};\draw[guide](0,0)--(Q)node[midway,left]{$r$};
\draw[axis](Q)--++(35:2.6)node[right]{$F$};\draw[guide](Q)--++(1.1,0);\draw($(Q)+(0:.6)$)arc(0:35:.6);\node at($(Q)+(17.5:.75)$){$\theta$};
\draw[guide](0,0)--(-.919,.919);\node at(-.45,.85){$R$};\node[below]at(0,-1.3){$P$};
\end{tikzpicture}
\caption{力沿内轴下侧的切线，因而对C的力矩为逆时针 $Fr$。}\label{sup6:fig:d5-e5-q}
\end{figure}
\begin{solution}
按图\ref{sup6:fig:d5-e5-q}，取质心加速度a、地面摩擦分量f向右为正，角加速度取顺时针为正。拉力对C的力矩与顺时针正向相反，故
\[
Ma=F\cos\theta+f,\qquad
\beta MR^2\alpha=-Fr-fR,\qquad a=R\alpha.
\]
图\ref{sup6:fig:d5-e5-q}中，半径CQ与拉线垂直，力臂就是内轴半径r，拉力矩大小为 $Fr$。拉线与水平线的夹角 $\theta$ 用于分解拉力的水平、竖直分量。
将第二式除以R，并与第一式相加，消去f：
\[
(1+\beta)Ma=F(\cos\theta-r/R).
\]
于是
\[
a=\frac{F(\cos\theta-r/R)}{(1+\beta)M},\qquad
f=-\frac{F(\beta\cos\theta+r/R)}{1+\beta}.
\]
本题范围内f总为负，即摩擦向左；质心加速度却可以向右或向左。临界角为 $\theta_*=\arccos(r/R)$。从静止开始时，$\theta<\theta_*$ 向右运动，$\theta>\theta_*$ 向左运动。
\begin{figure}[H]\centering
\begin{tikzpicture}
\draw[axis](0,0)--(4.6,0)node[right]{$\theta$};\draw[axis](0,-1.3)--(0,1.8)node[above]{$(1+\beta)Ma/F$};
\draw[curve,domain=0:90,samples=70]plot({.045*\x},{2*(cos(\x)-.5)});
\node[below]at(2.7,0){$\theta_*$};\node[below]at(4.05,0){$\pi/2$};\node[left]at(0,1){$1-r/R$};
\begin{scope}[shift={(7.4,.2)}]
\draw[curve](0,0)circle(1);\draw[curve](-1.5,-1)--(1.5,-1);
\draw[axis](0,-1)--(-1,-1)node[below]{$f<0$};\draw[axis](0,-1)--(0,-.2)node[left]{$N$};\draw[axis](0,0)--(0,-.7)node[right]{$Mg$};
\end{scope}
\end{tikzpicture}
\caption{左图示意 $r/R=1/2$ 时的加速度变号；右图的地面摩擦始终向左。}\label{sup6:fig:d5-e5-s}
\end{figure}
竖直方向没有加速度，$N=Mg-F\sin\theta$。因此图\ref{sup6:fig:d5-e5-s}所用模型还要求
\[
N>0,\qquad
\frac{F(\beta\cos\theta+r/R)}{1+\beta}
\le\mu_s(Mg-F\sin\theta).
\]
接触一旦失效，或摩擦超过极限，就须按新的运动状态重列方程。
还可核对拉力功率。由图\ref{sup6:fig:d5-e5-q}中的切点位置
$\vct r_{QC}=(r\sin\theta,-r\cos\theta)$，纯滚动时
\[
v_{Qx}=v_C(1-r\cos\theta/R),\qquad
v_{Qy}=-v_Cr\sin\theta/R.
\]
故拉力功率为 $Fv_C(\cos\theta-r/R)$，与
$\dd K/\dd t=(1+\beta)Mv_Ca$ 相等。从静止受力后向左运动时，$v_C$ 和括号均为负，输入功率仍为正。
\end{solution}



\begin{example}[圆柱在被拉动的板上滚动]\label{sup6:problem-example-7}\label{sup6:ex:d5-e7}
质量m的长板可在光滑水平面上平移，板上圆柱质量M、半径R、惯量 $I_C=\beta MR^2$。系统初始静止，对板施加向右恒力F。圆柱与板无滑动，接触静摩擦系数为 $\mu_s$。求板和圆柱的加速度、角加速度、摩擦力及无滑动条件，并讨论摩擦所做的功。板长足够，运动中圆柱未到板端。
\end{example}
\begin{figure}[H]\centering
\begin{tikzpicture}
\draw[curve](-2,-1.3)rectangle(4,-1);\draw[curve](-2,-1.55)--(4,-1.55);\draw[curve](0,0)circle(1);
\fill(0,0)circle(1.5pt)node[above]{$C,M$};\draw[axis](4,-1.15)--(5.3,-1.15)node[right]{$F$};\node[below]at(2,-1.55){光滑地面};\node at(2.8,-.7){板m};
\draw[axis](.2,1.3)--(1.3,1.3)node[right]{$x$};
\end{tikzpicture}
\caption{板与地面无摩擦；圆柱与板的接触材料点具有相同水平速度。}\label{sup6:fig:d5-e7-q}
\end{figure}
\begin{solution}
图\ref{sup6:fig:d5-e7-q}中，板向右加速时，摩擦使圆柱向右平移并逆时针转动。本题取逆时针角加速度为正。设板加速度A、圆柱加速度a、板对圆柱的摩擦f均以向右为正，则
\[
Ma=f,\qquad \beta MR^2\alpha=fR,\qquad
a+R\alpha=A.
\]
最后一式由接触点速度 $v_C+R\omega=V$ 求导得到，加号与本题逆时针正向一致。前两式代入约束，得到
\[
a=\frac{\beta}{1+\beta}A,\quad
R\alpha=\frac{A}{1+\beta},\quad
f=\frac{\beta M}{1+\beta}A.
\]
板受到F与反向摩擦，故 $mA=F-f$。令 $D=(1+\beta)m+\beta M$，消元后
\[
\boxed{A=\frac{(1+\beta)F}D,\quad
a=\frac{\beta F}D,\quad
\alpha=\frac F{RD},\quad f=\frac{\beta MF}D.}
\]
圆柱没有竖直加速度，$N=Mg$，因此须有 $F\le\mu_sgD/\beta$。
\begin{figure}[H]\centering
\begin{tikzpicture}
\draw[curve](-1.6,-1.25)rectangle(2.5,-1);\draw[curve](0,0)circle(1);\fill(0,-1)circle(1.5pt)node[below]{$P$};
\draw[axis](0,-1)--(1.15,-1)node[right]{$f$};\draw[axis](0,-1)--(-1.15,-1)node[midway,below]{$-f$};
\draw[axis](0,.15)--(1.1,.15)node[right]{$v_C$};\draw[axis](-.8,.8)arc(135:210:1.1);\node[left]at(-1.2,.3){$\omega$};
\node[align=left]at(5.7,.15){$v_P=v_C+R\omega=V$\\圆柱：$\mathcal P_f=fV$\\板：$\mathcal P_f=-fV$};
\end{tikzpicture}
\caption{相同的接触点速度使两个摩擦功率相消；圆柱单独得到正的摩擦功。}\label{sup6:fig:d5-e7-s}
\end{figure}
由图\ref{sup6:fig:d5-e7-s}，圆柱摩擦功率也可拆成 $fv_C+(fR)\omega=fV$。它同时增加平动能和转动能；板的摩擦功率为 $-fV$，整个系统的外力功率只剩FV。
由速度约束，$v_C=\beta V/(1+\beta)$、$R\omega=V/(1+\beta)$。系统总动能为
\[
K=\frac12\left(m+\frac{\beta M}{1+\beta}\right)V^2.
\]
对时间求导，$FV=\dot K$ 给出相同的板加速度。圆柱的接触点随板运动，计算圆柱自身的能量变化时，必须计入所获摩擦功。
\end{solution}



\begin{example}[圆柱同时与上下两面无滑动接触]\label{sup6:problem-example-8}\label{sup6:ex:d5-e8}
质量M、半径R、惯量 $\beta MR^2$ 的圆柱置于固定水平地面与质量m的水平长板之间，$0<\beta\le1$。上板由理想导向机构保持水平平移，导向机构没有水平力、不做功；已知板对圆柱的法向压力恒为N。对板施加向右力F。两处均无滑动，求上板和圆柱的加速度及两处摩擦，写出摩擦系数分别为 $\mu_1,\mu_2$ 时的成立条件。
\end{example}
\begin{figure}[H]\centering
\begin{tikzpicture}
\draw[curve](-2,-1)--(3,-1);\draw[curve](0,0)circle(1);\draw[curve](-2,1)rectangle(3,1.25);
\draw[guide](-2,1.6)--(3,1.6)node[right]{水平导向};\draw[axis](3,1.12)--(4.3,1.12)node[right]{$F$};
\fill(0,1)circle(1.5pt)node[above left]{$Q$};\fill(0,-1)circle(1.5pt)node[below]{$P$};\fill(0,0)circle(1.5pt)node[left]{$C$};
\node[right]at(1.3,.5){$N$已给定};\node[above]at(-1.4,1.25){板m};
\end{tikzpicture}
\caption{上板的竖直载荷由导向与接触共同承担，不能自行令给定压力N等于mg。}\label{sup6:fig:d5-e8-q}
\end{figure}
\begin{solution}
图\ref{sup6:fig:d5-e8-q}中，下接触点P的速度为零。取圆柱向右平动、顺时针转动为正，$a=R\alpha$。上接触点Q速度为 $v_Q=v_C+R\omega=2v_C$，故上板加速度 $A=2a$。
这里对两个接触位置的水平速度求导，得到 $A=2a$。接触材料点还可有竖直的向心加速度，故不能把“接触点瞬时速度为零”解释为该材料点的整个加速度为零。图\ref{sup6:fig:d5-e8-q}中两个约束使用的都是水平方向分量。
设上、下两处作用于圆柱的摩擦分别为 $f_1,f_2$，均先取向右为正。平动和转动方程为
\[
Ma=f_1+f_2,\qquad \beta Ma=f_1-f_2.
\]
相加、相减，
\[
f_1=\frac{1+\beta}{2}Ma,\qquad
f_2=\frac{1-\beta}{2}Ma.
\]
板的水平运动满足 $mA=F-f_1$，结合 $A=2a$，得到
\[
D=4m+(1+\beta)M,\qquad a=\frac{2F}D,\quad A=\frac{4F}D,
\]
\[
f_1=\frac{(1+\beta)MF}D,\qquad
f_2=\frac{(1-\beta)MF}D.
\]
当 $\beta<1$ 时，地面对圆柱的静摩擦也向右。上接触的驱动力矩使轮缘具有向左滑动的趋势，下接触摩擦由此指向右方。
\begin{figure}[H]\centering
\begin{tikzpicture}
\draw[curve](0,0)circle(1);\draw[curve](-1.5,-1)--(1.8,-1);\draw[curve](-1.5,1)--(1.8,1);
\draw[axis](0,1)--(1.3,1)node[above]{$f_1$};\draw[axis](0,-1)--(.8,-1)node[below]{$f_2$};
\draw[axis](0,1)--(0,.25)node[left]{$N$};\draw[axis](0,-1)--(0,-.25)node[left]{$Mg+N$};
\begin{scope}[shift={(5,0)}]
\draw[axis](0,-1)--(0,1.5)node[above]{高度};\draw[axis](0,-1)--(3,-1)node[right]{水平速度};
\draw[curve](0,-1)--(1.2,0)--(2.4,1);\fill(0,-1)circle(1.5pt)node[left]{$P:0$};\fill(1.2,0)circle(1.5pt)node[right]{$C:v_C$};\fill(2.4,1)circle(1.5pt)node[right]{$Q:2v_C$};
\end{scope}
\end{tikzpicture}
\caption{两处摩擦的力矩方向相反；右图的速度分布给出上板速率为圆柱质心速率的两倍。}\label{sup6:fig:d5-e8-s}
\end{figure}
图\ref{sup6:fig:d5-e8-s}中的上接触压力为N，下接触压力由圆柱竖直平衡得 $Mg+N$。因此两个独立条件为
\[
\frac{(1+\beta)MF}D\le\mu_1N,\qquad
\frac{(1-\beta)MF}D\le\mu_2(Mg+N).
\]
任一处先达到极限，都须更换约束。对于薄圆环 $\beta=1$，下接触摩擦恰为零；它仍满足无滑动运动，并不需要假设两处摩擦都达到极限。
用上板速度V表示总动能，还可得
\[
K=\frac12mV^2+\frac12(1+\beta)M(V/2)^2=\frac D8V^2.
\]
整个系统中，上接触的摩擦功率相消，下接触点静止，导向机构不做功，故 $FV=\dot K=(D/4)V\dot V$，与 $A=4F/D$ 相符。
\end{solution}



\begin{example}[球沿凸圆面滚下，何时先打滑]\label{deep6:problem-example-9}\label{deep6:ex:d5-e9}
匀质实心球半径r、质量m，在半径R的固定凸圆弧外侧运动，球心轨迹半径为 $a=R+r$。球从靠近顶点的角度 $\theta_0>0$ 处静止释放，取 $\theta_0\to0^+$ 的极限；$\theta$ 从向上的竖直半径量起。静摩擦系数 $\mu_s=1/5$。求纯滚动阶段的速度、反力和摩擦，并求首次打滑的位置。判断能否直接用纯滚动方程的 $N=0$ 求实际离面角。
\end{example}
\begin{figure}[H]\centering
\begin{tikzpicture}
\draw[curve](-2.5,0)arc(180:0:2.5);\fill(0,0)circle(1.5pt)node[below]{$O$};\draw[guide](0,0)--(0,3.7);
\coordinate(C)at(1.87,2.67);\draw[curve](C)circle(.76);\fill(C)circle(1.4pt)node[above]{$C$};\draw[guide](0,0)--(C)node[midway,right]{$a$};
\draw(0,.9)arc(90:55:.9);\node at(.45,1.1){$\theta$};\node[left]at(-2,1.2){固定圆弧R};
\end{tikzpicture}
\caption{球心沿半径 $a=R+r$ 的圆运动。角度从顶点向右侧增大。}\label{deep6:fig:d5-e9-q}
\end{figure}
\begin{solution}
先保留释放角 $\theta_0$。纯滚动的能量关系为
\[
mg a(\cos\theta_0-\cos\theta)=\frac12(1+\beta)mv^2,
\qquad \beta=\frac25.
\]
对给定位置再取 $\theta_0\to0^+$，得到
\[
v^2=\frac{2ga}{1+\beta}(1-\cos\theta).
\]
严格位于顶点且完全静止的球会保持不动。这里先在 $\theta_0>0$ 处释放，再令 $\theta_0\to0^+$，以保留离开不稳定平衡所需的初始偏离。
沿图\ref{deep6:fig:d5-e9-q}的切向向下列式，$m\dot v=mg\sin\theta-f$。以球滚下时的自转方向为正，转动方程为$\beta mr^2\dot\omega=fr$，滚动约束给出$\dot v=r\dot\omega$，所以$f=\beta m\dot v$。代回平动方程并合并，
\[
(1+\beta)m\dot v=mg\sin\theta,\qquad
\dot v=\frac{g\sin\theta}{1+\beta},\qquad
f=\frac{\beta}{1+\beta}mg\sin\theta.
\]
这一步没有约去速度，在静止释放瞬间同样成立。摩擦方向沿坡向上。再沿指向O的方向列式，
\[
mg\cos\theta-N=\frac{mv^2}a,
\quad N=\frac{mg}{1+\beta}\big[(3+\beta)\cos\theta-2\big].
\]
\begin{figure}[H]\centering
\begin{tikzpicture}
\draw[curve](0,0)circle(.8);\draw[axis](0,0)--(0,-1.3)node[left]{$mg$};\draw[axis](0,0)--(.75,1.07)node[right]{$N$};\draw[axis](0,0)--(-.98,.69)node[left]{$f$};
\begin{scope}[shift={(4,-.3)}]
\draw[axis](0,0)--(5,0)node[right]{$\theta$};\draw[axis](0,0)--(0,2.8)node[above]{力除以mg};
\draw[curve,domain=0:54,samples=70]plot({.075*\x},{2.3*(17*cos(\x)-10)/7});
\draw[guide,domain=0:54,samples=70]plot({.075*\x},{2.3*(2/7)*sin(\x)*5});
\node[right]at(3.7,.45){$N/(mg)$};\node at(2.8,2.5){$f/(\mu_smg)$};
\pgfmathsetmacro{\xs}{.075*acos(340/389)}
\draw[guide](\xs,0)--(\xs,{2.3*270/389});\fill(\xs,{2.3*270/389})circle(1.5pt);\node[below]at(\xs,0){$\theta_s$};
\end{scope}
\end{tikzpicture}
\caption{首次 $f/\mu_s=N$ 即到达滚动极限。相交点右侧的曲线仅是继续强加滚动约束的形式结果。}\label{deep6:fig:d5-e9-s}
\end{figure}
图\ref{deep6:fig:d5-e9-s}中，首次打滑满足
\[
\beta\sin\theta_s=\mu_s\big[(3+\beta)\cos\theta_s-2\big].
\]
代入参数，化为 $\sin\theta_s=(17/10)\cos\theta_s-1$。在 $0<\theta_s<\arccos(10/17)$ 内，右端为正，平方后
\[
1-c^2=(17c/10-1)^2,\qquad
c\left(\frac{389}{100}c-\frac{17}5\right)=0.
\]
符合区间条件的根为
\[
\cos\theta_s=\frac{340}{389},\quad
\sin\theta_s=\frac{189}{389},\quad
v_s^2=\frac{70}{389}ga.
\]
此时 $N_s=270mg/389>0$、$f_s=54mg/389=\mu_sN_s$，球仍与圆面接触。还须确认这是沿运动方向遇到的第一个摩擦边界。令摩擦比值为
\[
q(\theta)=\frac{\beta\sin\theta}{(3+\beta)\cos\theta-2}.
\]
记分母$D=(3+\beta)\cos\theta-2$，则$D'=-(3+\beta)\sin\theta$。按商法则，
\[
\begin{aligned}
q'(\theta)&=\frac{\beta\cos\theta D-\beta\sin\theta D'}{D^2}\\
&=\frac{\beta[(3+\beta)(\cos^2\theta+\sin^2\theta)-2\cos\theta]}{D^2}
=\frac{\beta[(3+\beta)-2\cos\theta]}{D^2}>0.
\end{aligned}
\]
分母为正的接触区内，$q$ 从零增至无穷，因此图\ref{deep6:fig:d5-e9-s}中摩擦边界必先于纯滚动模型的离面边界被碰到。平方求根时排除 $c=0$，也是因为它已超出这个接触区。
若继续强加纯滚动约束，会在 $\cos\theta_d=10/17$ 处得到 $N=0$。但球已在更早的 $\theta_s$ 处打滑，此后的平动与自转不再满足 $v=r\omega$。实际离面位置须由动摩擦规律和两套运动方程重新求出。
\end{solution}



\section{加速与转动参考系中的约束运动}

支承加速或转动时，惯性力由地面系与动系之间的加速度变换得到。平衡、稳定性和摩擦极限都须说明参考系，并限制在实际允许的位移内。



\subsection{加速参考系与旋转约束}
车厢加速或导轨转动时，物体相对它们的加速度与地面加速度不同。将加速度变换代入牛顿方程，便得到相应的惯性力。以下先推导这些变换，再用于加速容器、旋转约束和地球附近的落体运动。速度与加速度的参考系须随符号一并确定。

\subsubsection{平移参考系与有效重力}
设惯性系为 $S$，参考系 $S'$ 的原点位置为 $\vct R_0(t)$，两系坐标轴始终平行。质点位置满足
\[
\vct r=\vct R_0+\vct r',\qquad
\vct a=\vct a_0+\vct a'.
\]
将牛顿方程 $m\vct a=\vct F$ 改写为
\begin{equation}
m\vct a'=\vct F-m\vct a_0.
\label{rot:eq:d6-translation}
\end{equation}
右端的 $-m\vct a_0$ 是平移惯性力。它与参考系加速度反向，不是另一物体施加的作用，也不与某个真实作用力构成作用力、反作用力对。
\begin{figure}[H]\centering
\begin{tikzpicture}
\coordinate(O)at(0,0);\coordinate(A)at(2.2,.7);\coordinate(P)at(5.8,2.5);
\draw[axis](O)--(6.7,0)node[right]{$x$};\draw[axis](O)--(0,3.2)node[above]{$y$};\node[below left]at(O){$O$};
\fill(A)circle(1.6pt)node[below]{$O'$};\fill(P)circle(2pt)node[above]{$m$};
\draw[axis](O)--(A)node[midway,above left]{$\vct R_0$};\draw[axis](A)--(P)node[midway,above]{$\vct r'$};
\draw[guide](O)--(P);\draw[axis](A)--++(1.4,0)node[right]{$\vct a_0$};
\end{tikzpicture}
\caption{平移系原点与质点的相对位置。坐标轴未转动，故没有科里奥利项。}\label{rot:fig:d6-translation}
\end{figure}
图\ref{rot:fig:d6-translation}中若 $\vct a_0$ 恒定，均匀重力可与惯性力合成
\[
\vct g_{\rm eff}=\vct g-\vct a_0.
\]
相对静止的悬绳沿有效重力方向，静止液面与有效重力垂直。处理接触问题时，还须另列 $N\ge0$：接触面只能推物体，所求反力超出这一范围时，原接触状态便无法维持。

若真实保守力在 $S'$ 中具有不显含时间的势能 $U(\vct r')$，且 $\vct a_0$ 恒定，惯性力势能为 $m\vct a_0\cdot\vct r'$，从而
\begin{equation}
E'=\frac12m|\vct v'|^2+U+m\vct a_0\cdot\vct r'
\label{rot:eq:d6-translation-energy}
\end{equation}
在无功理想约束下守恒。对上式求导，保守力为 $-\nabla U$，约束力的相对功率为零，因此
\[
\begin{aligned}
\dot E'&=m\vct v'\cdot\vct a'+\nabla U\cdot\vct v'
+m\vct a_0\cdot\vct v'+m\dot{\vct a}_0\cdot\vct r'\\
&=\vct v'\cdot(-m\vct a_0)+m\vct a_0\cdot\vct v'
+m\dot{\vct a}_0\cdot\vct r'.
\end{aligned}
\]
前两项相消。若 $\vct a_0$ 随时间改变，则
\[
\frac{\dd E'}{\dd t}=m\dot{\vct a}_0\cdot\vct r'
\]
一般不为零。$E'$ 描述相对运动，地面系机械能须用地面速度另行计算。

\subsubsection{转动参考系中的动力学}
\paragraph{加速度变换}
对任意矢量 $\vct B$，它在两系中的时间导数满足
\[
\left(\frac{\dd\vct B}{\dd t}\right)_S
=\left(\frac{\dd\vct B}{\dd t}\right)_{S'}+\vct\Omega\times\vct B.
\]
先对相对位置求导，得 $\vct v=\vct V_0+\vct v'+\vct\Omega\times\vct r'$。第二次求导时分别有
\[
\left(\frac{\dd\vct v'}{\dd t}\right)_S=\vct a'+\vct\Omega\times\vct v',
\]
\[
\left(\frac{\dd}{\dd t}\vct\Omega\times\vct r'\right)_S
=\dot{\vct\Omega}\times\vct r'
+\vct\Omega\times(\vct v'+\vct\Omega\times\vct r').
\]
再加原点加速度，得到
\begin{equation}
\begin{aligned}
\vct a={}&\vct a_0+\vct a'
+2\vct\Omega\times\vct v'
+\dot{\vct\Omega}\times\vct r'\\
&+\vct\Omega\times(\vct\Omega\times\vct r').
\end{aligned}
\label{rot:eq:d6-acceleration}
\end{equation}
两个 $\vct\Omega\times\vct v'$ 分别来自相对速度的惯性系导数和 $\vct\Omega\times\vct r'$ 的求导，相加得到系数2。
因此在 $S'$ 中，
\begin{equation}
m\vct a'=\vct F+\vct F_{\rm tr}+\vct F_{\rm cor}
+\vct F_{\rm E}+\vct F_{\rm cf},
\end{equation}
其中
\[
\begin{aligned}
\vct F_{\rm tr}&=-m\vct a_0,&
\vct F_{\rm cor}&=-2m\vct\Omega\times\vct v',\\
\vct F_{\rm E}&=-m\dot{\vct\Omega}\times\vct r',&
\vct F_{\rm cf}&=-m\vct\Omega\times(\vct\Omega\times\vct r').
\end{aligned}
\]
后三项依次为科里奥利力、欧拉惯性力和离心惯性力。

\paragraph{方向与相对速度}
\begin{figure}[H]\centering
\begin{tikzpicture}
\node at(0,0){$\odot\,\Omega$};\draw[guide](.4,0)--(2,0);\fill(2,0)circle(1.6pt)node[above left]{$P$};
\draw[axis](2,0)--(3.4,0)node[right]{$\vct v'$};\draw[axis](2,0)--(2,-1.4)node[below]{$\vct F_{\rm cor}$};\node at(1,-2.1){沿半径向外运动};
\begin{scope}[shift={(6,0)}]
\node at(0,0){$\odot\,\Omega$};\draw[guide](.4,0)--(2,0);\fill(2,0)circle(1.6pt)node[below]{$P$};
\draw[axis](2,0)--(2,1.4)node[above]{$\vct v'$};\draw[axis](2,0)--(3.4,0)node[right]{$\vct F_{\rm cor}$};\node at(1,-2.1){沿正方位角运动};
\end{scope}
\end{tikzpicture}
\caption{角速度垂直纸面向外。先算 $\vct\Omega\times\vct v'$，再取负号。}\label{rot:fig:d6-coriolis}
\end{figure}
图\ref{rot:fig:d6-coriolis}的 $\vct v'$ 是物体相对转动系的速度。相对转台静止的物体有 $\vct v'=0$，科里奥利力为零；它在地面系中仍可作圆周运动。反过来，地面系静止的物体在转动系中通常有非零相对速度。
对固定转轴，离心力始终背离转轴，大小为 $m\Omega^2\rho$。若转速增大，欧拉惯性力与转动的切向加速度反向。再计入图\ref{rot:fig:d6-coriolis}中的科里奥利力，合惯性力一般不沿径向；各项的方向须分别确定。

\subsubsection{离心势与旋转系能量}
当原点固定、$\vct\Omega$ 恒定时，离心力可写成势能的负梯度，
\[
U_{\rm cf}=-\frac12m|\vct\Omega\times\vct r'|^2
=-\frac12m\Omega^2\rho^2.
\]
离心势随距轴距离增大而降低，其负梯度正好指向外侧。
\begin{figure}[H]\centering
\begin{tikzpicture}
\node at(0,0){$\odot\,\Omega$};\fill(2,0)circle(1.6pt)node[above]{$P$};\draw[guide](.4,0)--(2,0)node[midway,below]{$\rho$};\draw[axis](2,0)--(3.3,0)node[right]{$m\Omega^2\rho$};
\begin{scope}[shift={(6,1)}]
\draw[axis](0,0)--(3.6,0)node[right]{$\rho$};\draw[axis](0,-2.4)--(0,.6)node[above]{$U_{\rm cf}$};
\draw[curve,domain=0:3,samples=60]plot(\x,{-.22*\x*\x});\node at(2,-2.6){$-m\Omega^2\rho^2/2$};
\end{scope}
\end{tikzpicture}
\caption{离心力向外，离心势随半径降低。势能零点取在转轴上。}\label{rot:fig:d6-centrifugal}
\end{figure}
科里奥利力与相对速度垂直，$\vct F_{\rm cor}\cdot\vct v'=0$。若真实势能 $U$ 在旋转系中不显含时间，约束也相对该系静止，且没有摩擦耗散，则
\begin{equation}
E'=\frac12m|\vct v'|^2+U-\frac12m\Omega^2\rho^2
=\text{常量}.
\label{rot:eq:d6-jacobi}
\end{equation}
图\ref{rot:fig:d6-centrifugal}中的负二次项只属于旋转系能量。取转轴为 $z$ 轴，将 $\vct v=\vct v'+\vct\Omega\times\vct r'$ 代入惯性系动能，
\[
E=\frac12m|\vct v'|^2+m\vct v'\cdot(\vct\Omega\times\vct r')
+\frac12m\Omega^2\rho^2+U.
\]
惯性系角动量的轴向分量满足
\[
\Omega L_z=m\vct\Omega\cdot(\vct r'\times\vct v')+m\Omega^2\rho^2
=m\vct v'\cdot(\vct\Omega\times\vct r')+m\Omega^2\rho^2.
\]
两式相减，交叉项消去，有
\begin{equation}
E'=E-\Omega L_z.
\label{rot:eq:d6-energy-relation}
\end{equation}
即使 $E'$ 守恒，维持转速的电机仍可改变地面系机械能 $E$，变化量满足 $\Delta E=\Omega\Delta L_z$。角速度随时间变化时，还须计入欧拉惯性力和离心势的时间变化，式\eqref{rot:eq:d6-jacobi}的守恒条件已不成立。

\subsubsection{相对平衡、稳定性与约束}
约束相对旋转系固定、运动只需一个独立坐标 $q$ 时，能量方程常可写为
\[
\frac12A(q)\dot q^2+U_{\rm eff}(q)=E',\qquad A(q)>0.
\]
相对平衡满足 $U_{\rm eff}'(q_*)=0$。若二阶导数为正，小振动角频率为 $\sqrt{U_{\rm eff}''(q_*)/A(q_*)}$；若二阶导数为零，则须继续展开高阶项。
在两个或更多自由度中，科里奥利项会耦合速度分量。有效势的严格极小值可给出稳定性的充分条件，但有效势出现负曲率并不总意味着运动发散；应保留速度耦合项，求线性方程的特征频率。
由运动方程求出相对平衡后，还须核对约束力：静摩擦不能超过极限，单侧接触的压力和轻绳张力不能为负。

\subsubsection{地球附近的局部坐标}
取 $x$ 向东、$y$ 向北、$z$ 向上，三轴构成右手系。在纬度 $\lambda$ 处，地球角速度的局部分量为
\[
\vct\Omega=(0,\Omega\cos\lambda,\Omega\sin\lambda).
\]
\begin{figure}[H]\centering
\begin{tikzpicture}
\draw[curve](0,0)circle(1.8);\draw[guide](-2,0)--(2,0);\draw[axis](0,-1.9)--(0,2.6)node[above]{$\vct\Omega$};
\coordinate(P)at(1.474,1.032);\fill(P)circle(1.6pt)node[right]{$P$};\draw[guide](0,0)--(P);\draw(.7,0)arc(0:35:.7);\node at(.9,.3){$\lambda$};\node[below left]at(0,0){$O$};
\begin{scope}[shift={(5.3,-.5)}]
\draw[axis](0,0)--(3.2,0)node[right]{北y};\draw[axis](0,0)--(0,2.7)node[above]{上z};\node[below left]at(0,0){$\odot$东x};
\draw[axis](0,0)--(2.25,1.575)node[above right]{$\vct\Omega$};\draw[guide](2.25,0)--(2.25,1.575)--(0,1.575);
\node[below]at(1.2,0){$\Omega\cos\lambda$};\node[left]at(0,1.575){$\Omega\sin\lambda$};\draw(.7,0)arc(0:35:.7);\node at(.95,.3){$\lambda$};
\end{scope}
\end{tikzpicture}
\caption{左图以地心纬度示意局部方向；右图为北、上平面，东向垂直纸面向外。}\label{rot:fig:d6-earth}
\end{figure}
将地球近似为球体，并在图\ref{rot:fig:d6-earth}的局部区域内取均匀重力。$g$ 为当地有效重力加速度，已计入静止物体的离心效应，列式时不再另加这一静态离心项。
短时运动可用零阶速度代入科里奥利力，求一阶自转修正。所需尺度条件为运动距离远小于地球半径、$\Omega t\ll1$；忽略项包括二阶自转效应以及按此模型舍去的重力和局部坐标梯度。

\LectureNeedspace{4\baselineskip}


\subsection{例题与模型变式}

\begin{example}[加速斜面上的相对运动]\label{rot:problem-example-1}\label{rot:ex:d6-e1}
光滑斜面向右上方升高，与水平面成 $\gamma$ 角。斜面保持姿态，水平加速度为 $A$，向右为正。质量为 $m$ 的小块初始相对斜面静止。求接触阶段的相对加速度、支持力和地面加速度，确定相对静止条件及接触能够维持的范围。
\end{example}
\begin{figure}[H]\centering
\begin{tikzpicture}
\draw[curve](0,0)--(5,0)--(5,2.9)--cycle;\draw[fill=white,rotate around={30:(2.7,1.56)}](2.4,1.56)rectangle(3,2.06);\node at(2.9,2.35){$m$};
\draw[axis](1,3)--(3,3)node[right]{$A$};\draw(.9,0)arc(0:30:.9);\node at(1.2,.32){$\gamma$};\draw[axis](3.8,2.4)--(4.9,3.04)node[right]{$s$};
\end{tikzpicture}
\caption{相对坐标 $s$ 沿坡向上为正，水平加速度 $A$ 可取正或负。}\label{rot:fig:d6-e1-q}
\end{figure}
\begin{solution}
在图\ref{rot:fig:d6-e1-q}的随斜面平移系中，惯性力为 $-mA\vct e_x$。沿坡和向外法线的单位矢量分别为
\[
\vct e_s=(\cos\gamma,\sin\gamma),\qquad
\vct e_n=(-\sin\gamma,\cos\gamma).
\]
图\ref{rot:fig:d6-e1-s}中的合有效力为 $(-mA,-mg)$。沿坡取点积，重力和惯性力都可能使物块下滑；向外法线取点积时，惯性力却有向外分量：
\[
m\ddot s=(-mA,-mg)\cdot\vct e_s,\qquad
0=N+(-mA,-mg)\cdot\vct e_n.
\]
法向相对加速度为零来自仍在平面上运动的约束，不能把它直接用于离面以后的阶段。
将重力与惯性力分别投影，有
\[
\ddot s=-g\sin\gamma-A\cos\gamma,\qquad
N=m(g\cos\gamma-A\sin\gamma).
\]
相对静止要求 $\ddot s=0$，即 $A=-g\tan\gamma$：斜面须向左加速。此时 $N=mg/\cos\gamma>0$，因而该相对平衡能满足接触条件。
\begin{figure}[H]\centering
\begin{tikzpicture}
\fill(0,0)circle(1.8pt)node[above right]{$m$};\draw[guide](-1.7,-.981)--(1.7,.981);
\draw[axis](0,0)--(0,-1.6)node[below]{$mg$};\draw[axis](0,0)--(-1.5,0)node[left]{$mA$};\draw[axis](0,0)--(-.7,1.212)node[above]{$N$};
\node[align=left]at(4.3,.1){图示 $A>0$\\惯性力向左\\法向反力随A增大而减小};
\end{tikzpicture}
\caption{平移系受力图。惯性力的箭头只对应图示的正A情形。}\label{rot:fig:d6-e1-s}
\end{figure}
图\ref{rot:fig:d6-e1-s}中，保持接触要求 $A\le g\cot\gamma$。当 $A>g\cot\gamma$，有效重力的向外法向分量为 $A\sin\gamma-g\cos\gamma>0$，自由运动使物块离开斜面，所需的负支持力无法提供。
等号处 $N=0$，有效重力恰与斜面平行。若平面无限延伸且初始法向相对速度为零，物块的自由轨迹可以继续与斜面重合，法向间隙仍为零。因此 $N=0$ 给出接触的临界条件，是否分离还须考察随后的法向相对运动。

保持接触时，小块地面加速度为 $\vct a=A\vct e_x+\ddot s\vct e_s$，逐项整理得
\[
\begin{aligned}
a_x&=A\sin^2\gamma-g\sin\gamma\cos\gamma,\\
a_y&=-g\sin^2\gamma-A\sin\gamma\cos\gamma.
\end{aligned}
\]
这也等于 $-g\vct e_y+(N/m)\vct e_n$，与地面系中只受重力和支持力的方程一致。有限斜面还须在物块到达端点时终止这一阶段。
\end{solution}



\begin{example}[加速水槽的液面与压力]\label{rot:problem-example-2}\label{rot:ex:d6-e2}
开口矩形水槽沿 $x$ 方向长 $2\ell$、高 $H$、垂直纸面宽 $B$，初始平均液深为 $h_0$。水槽以恒加速度 $A$ 向右运动，密度为 $\rho$ 的液体相对槽静止，忽略毛细效应。取槽中心为 $x=0$，底面为 $z=0$。求液面、底部表压及两端壁对液体的水平合力，确定不溢流且槽底全部被液体覆盖的条件。
\end{example}
\begin{figure}[H]\centering
\begin{tikzpicture}
\fill[gray!12](0,0)--(6,0)--(6,1)--(0,2.2)--cycle;\draw[curve](0,2.8)--(0,0)--(6,0)--(6,2.8);\draw[curve](0,2.2)--(6,1);\draw[guide](0,1.6)--(6,1.6)node[right]{$h_0$};
\draw[axis](2,3.2)--(4,3.2)node[right]{$A$};\draw[<->](-.55,0)--(-.55,2.8)node[midway,left]{$H$};
\node[below]at(0,0){$-\ell$};\node[below]at(3,0){$0$};\node[below]at(6,0){$\ell$};\node at(3,.65){$\rho$，宽B};
\end{tikzpicture}
\caption{液面左高右低，虚线为原来的平均液面高度。}\label{rot:fig:d6-e2-q}
\end{figure}
\begin{solution}
在随槽系中，有效重力为 $(-A,-g)$。流体静力方程给出
\[
\frac{\partial p}{\partial x}=-\rho A,\qquad
\frac{\partial p}{\partial z}=-\rho g.
\]
沿图\ref{rot:fig:d6-e2-q}的自由液面，表压恒为零，故 $\dd p=-\rho A\dd x-\rho g\dd z_s=0$，液面斜率为 $\dd z_s/\dd x=-A/g$。加速后两端液深都会改变，积分留下的高度常数须由体积守恒确定。
先对压强方程积分：$\partial p/\partial x=-\rho A$给$p=-\rho Ax+f(z)$，再由$f'(z)=-\rho g$得$p=-\rho Ax-\rho gz+C$。自由液面上$p=p_{\rm atm}$，故液面可写为$z_s(x)=c-Ax/g$。高度常数$c$由体积守恒确定：
\[
B\int_{-\ell}^{\ell}\left(c-\frac Agx\right)\dd x
=B\left[cx-\frac{A}{2g}x^2\right]_{-\ell}^{\ell}
=2B\ell c=2B\ell h_0.
\]
所以$c=h_0$、$C=p_{\rm atm}+\rho gh_0$，得到
\[
z_s(x)=h_0-\frac Agx,\qquad
p(x,z)-p_{\rm atm}=\rho g(h_0-z)-\rho Ax.
\]
图\ref{rot:fig:d6-e2-q}中，两端液深分别为
$h_L=h_0+A\ell/g$、$h_R=h_0-A\ell/g$，底部表压为 $p_b(x)=\rho gh_0-\rho Ax$。
\begin{figure}[H]\centering
\begin{tikzpicture}
\draw[curve](0,2.4)--(0,0)--(4,0)--(4,2.4);\draw[guide](0,2)--(4,.8);
\draw[axis](0,.67)--(1.4,.67)node[above]{$F_L$};\draw[axis](4,.27)--(3.3,.27)node[above]{$F_R$};
\begin{scope}[shift={(6,0)}]
\draw[axis](0,0)--(3.8,0)node[right]{$x$};\draw[axis](0,0)--(0,2.7)node[above]{$p_b$};\draw[curve](.3,2.2)--(3.2,1);
\node[below]at(.3,0){$-\ell$};\node[below]at(3.2,0){$\ell$};\node at(1.7,2.7){线性压强分布};
\end{scope}
\end{tikzpicture}
\caption{端壁作用于液体的力方向相反，左端合力较大。底部表压从左向右线性降低。}\label{rot:fig:d6-e2-s}
\end{figure}
端壁上的表压随深度线性增大。分别积分，图\ref{rot:fig:d6-e2-s}中的两个力为
\[
\begin{aligned}
F_L&=\rho gB\int_0^{h_L}(h_L-z)\dd z
=\rho gB\left[h_Lz-\frac{z^2}{2}\right]_0^{h_L}
=\frac12\rho gBh_L^2,\\
F_R&=\rho gB\int_0^{h_R}(h_R-z)\dd z
=\rho gB\left[h_Rz-\frac{z^2}{2}\right]_0^{h_R}
=\frac12\rho gBh_R^2.
\end{aligned}
\]
于是净水平力为
\[
F_L-F_R=\frac12\rho gB(h_L-h_R)(h_L+h_R)
=2\rho B\ell h_0A=M_{\rm liq}A.
\]
这一合力使全部液体获得水平加速度$A$。底面对液体的竖直合力为
\[
B\int_{-\ell}^{\ell}p_b\dd x
=\rho B\left[gh_0x-\frac{Ax^2}{2}\right]_{-\ell}^{\ell}
=2\rho gB\ell h_0=M_{\rm liq}g,
\]
恰好平衡液体重力。
取 $A\ge0$，槽底不露出要求 $h_R\ge0$，不溢出要求 $h_L\le H$，故
\[
0\le A\le\frac g\ell\min(h_0,H-h_0).
\]
右端液深恰为零或左端液面恰到槽口时取等号。若右端已经露底，体积积分的范围须改为实际湿润区域，见习题\ref{rot:ex:d6-h2}。
\end{solution}



\begin{example}[基座振动引起的弹簧响应]\label{rot:problem-example-3}\label{rot:ex:d6-e3}
水平小车上固定一轻弹簧，劲度系数为 $k$，连接质量为 $m$、可无摩擦滑动的物块。弹簧无形变时，物块中心对应的地面坐标为 $X(t)=X_0\cos\nu t$。令物块相对该位置的位移为 $x$，初始 $x=\dot x=0$。求相对运动及共振响应，分别计算相对运动能量的变化率和支点传入物块、弹簧系统的功率。
\end{example}
\begin{figure}[H]\centering
\begin{tikzpicture}
\draw[curve](0,-.6)rectangle(6,1.2);\draw[curve](.5,-.5)--(.5,1);\draw[spring](.5,.3)--(3.9,.3);\draw[fill=white](3.9,-.1)rectangle(4.7,.7);\node at(4.3,.3){$m$};
\draw[guide](2.8,-.5)--(2.8,1);\node[above]at(2.8,1){$O':X(t)$};\draw[<->](2.8,-.9)--(4.3,-.9)node[midway,below]{$x$};\draw[axis](5,1.7)--(6.4,1.7)node[right]{基座运动};
\end{tikzpicture}
\caption{物块地面坐标为 $q=X+x$，$x$ 同时是弹簧伸长量。}\label{rot:fig:d6-e3-q}
\end{figure}
\begin{solution}
图\ref{rot:fig:d6-e3-q}中，小车系的惯性力为$-m\ddot X$。由基座位置逐次求导，$\dot X=-X_0\nu\sin\nu t$、$\ddot X=-X_0\nu^2\cos\nu t$。置$\omega_0=\sqrt{k/m}$，有
\[
\ddot x+\omega_0^2x=-\ddot X=X_0\nu^2\cos\nu t.
\]
地面坐标与形变量的关系是 $q=X+x$，弹簧力为 $-kx$，不是 $-kq$。从地面系写 $m(\ddot X+\ddot x)=-kx$，也得到同一方程。题设的初始静止指相对基座的 $x(0)=\dot x(0)=0$；在本题基座相位下，$q(0)=X_0$、$\dot q(0)=0$。
当$\nu\ne\omega_0$，先取特解$x_p=C\cos\nu t$，其二阶导数为$-C\nu^2\cos\nu t$。代入方程并比较余弦系数，得$C(\omega_0^2-\nu^2)=X_0\nu^2$。加上齐次通解，
\[
x=C\cos\nu t+A\cos\omega_0t+B\sin\omega_0t,
\qquad C=\frac{X_0\nu^2}{\omega_0^2-\nu^2}.
\]
初始位移给$C+A=0$，初始速度给$B\omega_0=0$，因此$A=-C,B=0$，得到
\[
x(t)=\frac{X_0\nu^2}{\omega_0^2-\nu^2}
\big(\cos\nu t-\cos\omega_0t\big).
\]
地面坐标为
\[
q(t)=\frac{X_0}{\omega_0^2-\nu^2}
\big(\omega_0^2\cos\nu t-\nu^2\cos\omega_0t\big).
\]
接近共振时，两个余弦项的频率相近。第二项由初始条件确定，在无阻尼模型中始终保留，响应须由两项叠加求得。
\begin{figure}[H]\centering
\begin{tikzpicture}
\draw[axis](0,0)--(7,0)node[right]{$\omega_0t$};\draw[axis](0,-2.3)--(0,2.8)node[above]{$x/X_0$};
\draw[curve,domain=0:11.5,samples=180]plot({.55*\x},{.19*\x*sin(\x*180/pi)});
\draw[guide](0,0)--(6.325,2.185);\draw[guide](0,0)--(6.325,-2.185);\node[right]at(4,2.5){包络 $\pm\omega_0t/2$};
\node[left]at(0,1.52){$4$};\node[left]at(0,-1.52){$-4$};
\end{tikzpicture}
\caption{实线为共振位移，虚线为线性增长的正、负包络。}\label{rot:fig:d6-e3-s}
\end{figure}
当 $\nu=\omega_0$，对上式取极限，利用
\[
\cos\nu t-\cos\omega_0t\sim-(\nu-\omega_0)t\sin\omega_0t,
\quad \omega_0^2-\nu^2\sim-2\omega_0(\nu-\omega_0),
\]
得到图\ref{rot:fig:d6-e3-s}所示的
\[
x(t)=\frac12X_0\omega_0t\sin\omega_0t.
\]
这一共振解要求基座持续按给定规律运动。振幅增长到弹簧超出线性范围或物块触及行程端点时，模型便须改变。

相对运动能量$E_r=m\dot x^2/2+kx^2/2$求导后，用相对运动方程消去加速度：
\[
\dot E_r=m\dot x\ddot x+kx\dot x
=\dot x(m\ddot x+kx)=-m\ddot X\dot x.
\]
地面系中，物块与弹簧的机械能为 $E=m(\dot X+\dot x)^2/2+kx^2/2$。由 $m\ddot q=-kx$，逐项求导，
\[
\dot E=(-kx)\dot q+kx(\dot q-\dot X)=-kx\dot X.
\]
后一式用支点的地面速度 $\dot X$ 计算实际传入功率；前一式用相对速度 $\dot x$ 计算惯性力对相对运动的功率。二者对应不同的能量方程。
\end{solution}



\begin{example}[匀速旋转径向杆上的滑珠]\label{rot:problem-example-4}\label{rot:ex:d6-e4}
轻直杆在水平面内绕固定端 $O$ 以恒角速度 $\Omega>0$ 逆时针转动。质量为 $m$ 的滑珠可沿杆无摩擦滑动，初始半径为 $a$、相对速度为零，杆端半径为 $R$，$0<a<R$。求到达杆端前的运动、杆对珠的切向作用力、到端点的时间，以及维持转速的电机对滑珠所做的功。
\end{example}
\begin{figure}[H]\centering
\begin{tikzpicture}
\draw[curve](0,0)--(5.6,0);\fill(0,0)circle(1.7pt)node[below]{$O$};\draw[fill=white](1.8,-.18)rectangle(2.2,.18);\node[above]at(2,.2){$m$};\node[below]at(1,-.1){$a$};\node[below]at(5.6,0){$R$};
\draw[axis](.7,.7)arc(45:135:1);\node at(.3,1.7){$\Omega$};\draw[axis](3,.5)--(4.4,.5)node[right]{$\vct e_r$};
\end{tikzpicture}
\caption{俯视水平转杆；珠的径向相对速度初始为零，方位角随杆转动。}\label{rot:fig:d6-e4-q}
\end{figure}
\begin{solution}
图\ref{rot:fig:d6-e4-q}中，杆不施加沿杆作用力。随杆系的径向方程为
\[
\ddot r=\Omega^2r.
\]
令$r\propto e^{st}$，特征方程为$s^2-\Omega^2=0$，故通解为$r=C_1e^{\Omega t}+C_2e^{-\Omega t}$。其导数为$\dot r=\Omega C_1e^{\Omega t}-\Omega C_2e^{-\Omega t}$。由$r(0)=a,\dot r(0)=0$，得到$C_1+C_2=a$、$\Omega(C_1-C_2)=0$，所以两系数均为$a/2$，即
\[
r=a\cosh\Omega t,\qquad
\dot r=a\Omega\sinh\Omega t.
\]
杆对滑珠的切向力记为 $N_\varphi$，正向取逆时针。由于相对切向加速度为零，它须平衡科里奥利力，
\[
N_\varphi=2m\Omega\dot r.
\]
该力也正好提供地面系极坐标加速度中的 $2\Omega\dot r$ 项。
\begin{figure}[H]\centering
\begin{tikzpicture}
\fill(0,0)circle(1.7pt)node[above right]{$m$};\draw[guide](-1.4,0)--(2.1,0);
\draw[axis](0,0)--(1.7,0)node[right]{$m\Omega^2r$};\draw[axis](0,0)--(0,1.4)node[above]{$N_\varphi$};\draw[axis](0,0)--(0,-1.4)node[below]{$2m\Omega\dot r$};
\node[align=left]at(5,.2){向下：科里奥利力\\向上：杆的真实作用力\\径向：相对加速度不为零};
\end{tikzpicture}
\caption{随杆系的水平力图。竖直重力与支持力已相互平衡。}\label{rot:fig:d6-e4-s}
\end{figure}
到达杆端时，$\cosh\Omega t_R=R/a$，所以
\[
t_R=\frac1\Omega\operatorname{arcosh}\frac Ra,\qquad
v_r=\Omega\sqrt{R^2-a^2},\qquad v_\varphi=\Omega R.
\]
图\ref{rot:fig:d6-e4-s}中的出杆时间还可写为普通对数。由 $\cosh z=(e^z+e^{-z})/2$，解正的z得到
\[
t_R=\frac1\Omega\ln\frac{R+\sqrt{R^2-a^2}}a.
\]
初始滑珠虽相对杆静止，地面切向速度却为 $\Omega a$，所以计算地面动能增量时不能把初动能置零。
由式\eqref{rot:eq:d6-jacobi}，相对运动满足
\[
\frac12m\dot r^2-\frac12m\Omega^2r^2=-\frac12m\Omega^2a^2.
\]
地面系机械能不守恒。图\ref{rot:fig:d6-e4-s}中的切向力对珠做功，功率为
$N_\varphi\Omega r=2m\Omega^2r\dot r$。维持转速的电机提供相等的功率，因此
\[
W=\int_a^R2m\Omega^2r\dd r
=m\Omega^2\bigl[r^2\bigr]_a^R=m\Omega^2(R^2-a^2).
\]
直接计算地面动能变化，初态为 $m\Omega^2a^2/2$，末态为 $m\Omega^2(2R^2-a^2)/2$，差值相同。杆端之后若珠自由离开，其地面初速度须同时保留径向和切向分量。
在图\ref{rot:fig:d6-e4-s}所示的到端点过程，径向动能增加 $m\Omega^2(R^2-a^2)/2$，切向动能也增加同样大小，两者之和正好为电机所做的功。杆对珠的力垂直其相对速度，但不垂直其地面速度，这正是两种能量描述不同的原因。
\end{solution}



\begin{example}[无摩擦转盘上的自由质点]\label{rot:problem-example-5}\label{rot:ex:d6-e5}
半径为 $R$ 的水平圆盘绕中心 $O$ 以恒角速度 $\Omega>0$ 逆时针转动。盘面完全光滑，质量为 $m$ 的质点在 $t=0$ 位于盘系 $x'=a,y'=0$，并相对盘静止，$0<a<R$。此后仅受竖直重力和支持力。求两个参考系中的轨迹、相对速度、到盘边的时间与旋转系能量。
\end{example}
\begin{figure}[H]\centering
\begin{tikzpicture}
\draw[curve](0,0)circle(2.2);\fill(0,0)circle(1.5pt)node[below left]{$O$};\fill(1,0)circle(1.8pt)node[below]{$P_0:a$};\draw[axis](1,0)--(1,1.4)node[right]{$\Omega a$};
\draw[axis](0,0)--(2.7,0)node[right]{$x=x'$};\draw[axis](0,0)--(0,2.7)node[above]{$y=y'$};\node at(-1.6,-1.7){初始时刻};
\end{tikzpicture}
\caption{相对盘静止的质点，在地面系中已有切向初速度 $\Omega a$。}\label{rot:fig:d6-e5-q}
\end{figure}
\begin{solution}
地面系中水平合力为零。由图\ref{rot:fig:d6-e5-q}的初态，
\[
x=a,\qquad y=\Omega at.
\]
相对盘静止只说明初始相对速度为零。由速度变换，地面初速度为 $\vct\Omega\times(a,0,0)=(0,\Omega a,0)$，而不是零。盘面光滑，此后没有水平力把质点拉成圆周运动。
盘系相对于地面逆时针转过 $\tau=\Omega t$，故将地面位置乘逆转矩阵，
\[
\begin{aligned}
x'&=a(\cos\tau+\tau\sin\tau),\\
y'&=a(-\sin\tau+\tau\cos\tau).
\end{aligned}
\]
这里使用的逆转矩阵为
\[
\begin{pmatrix}x'\\y'\end{pmatrix}
=\begin{pmatrix}\cos\tau&\sin\tau\\-\sin\tau&\cos\tau\end{pmatrix}
\begin{pmatrix}a\\a\tau\end{pmatrix}.
\]
这个矩阵改变位置矢量的坐标分量。图\ref{rot:fig:d6-e5-s}中，质点在地面系仍沿直线运动，转盘系中的曲线来自坐标轴随时间转动。
用$\dd/\dd t=\Omega\dd/\dd\tau$，对正弦、余弦以及含$\tau$的乘积分别求导，
\[
\dot x'=a\Omega(-\sin\tau+\sin\tau+\tau\cos\tau),\qquad
\dot y'=a\Omega(-\cos\tau+\cos\tau-\tau\sin\tau).
\]
括号内前两项抵消，得到
\[
\dot x'=a\Omega\tau\cos\tau,\qquad
\dot y'=-a\Omega\tau\sin\tau,
\quad |\vct v'|^2=a^2\Omega^2\tau^2.
\]
\begin{figure}[H]\centering
\begin{tikzpicture}[scale=.67]
\pgfmathsetmacro{\tauR}{sqrt(8)}
\draw[guide](0,0)circle(3);\node at(-1.8,3.35){地面系};
\draw[axis](-3.25,0)--(3.4,0)node[right]{$x/a$};\draw[axis](0,-3.25)--(0,3.4)node[above]{$y/a$};
\draw[curve](1,0)--(1,\tauR);\draw[axis](1,1)--(1,1.8);
\fill(1,0)circle(1.6pt)node[below right]{$P_0$};\fill(1,\tauR)circle(1.8pt)node[above right]{$P_R$};
\node at(-1.25,-1.6){$R/a=3$};
\begin{scope}[xshift=8.7cm]
\draw[guide](0,0)circle(3);\node at(-1.8,3.35){转盘系};
\draw[axis](-3.25,0)--(3.4,0)node[right]{$x'/a$};\draw[axis](0,-3.4)--(0,3.4)node[above]{$y'/a$};
\draw[curve,domain=0:\tauR,samples=150]plot({cos(\x*180/pi)+\x*sin(\x*180/pi)},{-sin(\x*180/pi)+\x*cos(\x*180/pi)});
\draw[axis,domain=1.4:1.7,samples=16]plot({cos(\x*180/pi)+\x*sin(\x*180/pi)},{-sin(\x*180/pi)+\x*cos(\x*180/pi)});
\fill(1,0)circle(1.6pt)node[above right]{$P_0$};
\fill({cos(\tauR*180/pi)+\tauR*sin(\tauR*180/pi)},{-sin(\tauR*180/pi)+\tauR*cos(\tauR*180/pi)})circle(1.8pt)node[below right]{$P_R$};
\end{scope}
\end{tikzpicture}
\caption{图取 $R=3a$，画出质点到达盘边前的轨迹。虚线为盘边；两图中的 $P_R$ 对应同一离盘时刻，$\Omega t_R=\sqrt8$。}\label{rot:fig:d6-e5-s}
\end{figure}
图\ref{rot:fig:d6-e5-s}的两种坐标给出同一到轴距离：
\[
r^2=x^2+y^2=x'^2+y'^2=a^2(1+\Omega^2t^2).
\]
所以首次到达盘边的时刻为
\[
t_R=\frac{\sqrt{R^2-a^2}}{\Omega a}.
\]
地面系速率保持 $\Omega a$，相对盘的速率却随时间增大。旋转系能量为
\[
E'=\frac12ma^2\Omega^2\tau^2
-\frac12m\Omega^2a^2(1+\tau^2)
=-\frac12m\Omega^2a^2.
\]
离心力对相对运动做正功，科里奥利力只改变相对速度方向。地面机械能 $E=m\Omega^2a^2/2$、轴向角动量 $L_z=m\Omega a^2$ 都保持不变，也满足 $E'=E-\Omega L_z$。
\end{solution}



\begin{example}[加速转盘上的静摩擦极限]\label{rot:problem-example-6}\label{rot:ex:d6-e6}
水平转盘从静止开始，以恒定角加速度 $\alpha>0$ 逆时针加速。质量为 $m$ 的小物体最初位于距轴 $r>0$ 处，相对盘静止，静摩擦系数为 $\mu$。求不打滑时的摩擦矢量、首次达到摩擦极限的角速度与时间，确定不打滑阶段存在的条件。再取 $\alpha=3\mu g/(5r)$，求达到极限时两摩擦分量的比例。
\end{example}
\begin{figure}[H]\centering
\begin{tikzpicture}
\draw[curve](0,0)circle(2);\fill(0,0)circle(1.5pt)node[below left]{$O$};\fill(1.3,0)circle(2pt)node[right]{$m$};\draw[guide](0,0)--(1.3,0)node[midway,below]{$r$};
\draw[axis](-1.2,1.2)arc(135:225:1.7);\node[anchor=east]at(-2.2,1.3){$\Omega=\alpha t$};\draw[axis](1.3,.2)--(1.3,1)node[left]{$\vct e_\varphi$};
\end{tikzpicture}
\caption{物体相对盘静止，但地面系加速度同时含向心与切向分量。}\label{rot:fig:d6-e6-q}
\end{figure}
\begin{solution}
图\ref{rot:fig:d6-e6-q}中，物体地面加速度为 $-r\Omega^2\vct e_r+r\alpha\vct e_\varphi$。唯一水平真实力为静摩擦，因此
\[
\vct f=-mr\Omega^2\vct e_r+mr\alpha\vct e_\varphi,
\qquad |\vct f|=mr\sqrt{\Omega^4+\alpha^2}.
\]
在盘系中，静摩擦的径向和切向分量分别抵消离心力和欧拉惯性力。物体相对盘静止，$\vct v'=0$，科里奥利力为零。
无滑动条件为
\[
\Omega^4+\alpha^2\le(\mu g/r)^2.
\]
若 $0<\alpha<\mu g/r$，摩擦随时间严格增大，首次达到极限时
\[
\Omega_s=\big[(\mu g/r)^2-\alpha^2\big]^{1/4},
\qquad t_s=\frac{\Omega_s}{\alpha}.
\]
若 $\alpha>\mu g/r$，初始切向需求已超过最大静摩擦，不能先假设存在不打滑阶段；若 $\alpha=\mu g/r$，只有 $t=0$ 瞬间达到等号，任意正时间均无法维持原约束。
\begin{figure}[H]\centering
\begin{tikzpicture}
\fill(0,0)circle(1.6pt)node[below right]{$m$};\draw[axis](0,0)--(-1.6,1.2)node[above left]{$|\vct f|=\mu mg$};\draw[guide](0,0)--(-1.6,0)--(-1.6,1.2)--(0,1.2)--cycle;
\node[below]at(-.8,0){$4\mu mg/5$};\node[right]at(0,.6){$3\mu mg/5$};\node[align=left]at(4,.7){径向向内\\切向沿加速方向\\两分量按矢量相加};
\end{tikzpicture}
\caption{指定角加速度下的3、4、5直角力三角形。}\label{rot:fig:d6-e6-s}
\end{figure}
对指定参数，令 $K=\mu g/r$，有 $\alpha=3K/5$、$\Omega_s^4=16K^2/25$，因此
\[
\Omega_s^2=\frac45K,\quad
|f_r|=\frac45\mu mg,\quad f_\varphi=\frac35\mu mg,
\quad t_s=\frac{2\sqrt5}{3}\sqrt{\frac r{\mu g}}.
\]
图\ref{rot:fig:d6-e6-s}的摩擦同时含径向、切向分量。达到极限后若继续加速，物体开始相对盘滑动；随后的运动须由动摩擦规律求解，盘的转速 $\Omega=\alpha t$ 已不等于物体自身的方位角速度。
\end{solution}



\begin{example}[绕竖直直径旋转的圆环]\label{rot:problem-example-7}\label{rot:ex:d6-e7}
半径为 $a$ 的光滑圆环以竖直直径为固定轴，由电机维持恒角速度 $\Omega$ 转动。质量为 $m$ 的小珠穿在环上，$\theta$ 从最低点量起，向图中右侧为正。求相对运动方程、全部相对平衡及其稳定性，求稳定平衡附近的小振动频率，并讨论临界转速下的运动。
\end{example}
\begin{figure}[H]\centering
\begin{tikzpicture}
\draw[curve](0,0)circle(2);\draw[axis](0,-2.4)--(0,2.8)node[above]{$z,\ \vct\Omega$};\fill(0,0)circle(1.5pt)node[left]{$O$};\coordinate(P)at(1.414,-1.414);\fill(P)circle(2.2pt)node[right]{$m$};
\draw[guide](0,0)--(P)node[midway,above]{$a$};\draw(0,-.8)arc(-90:-45:.8);\node at(.55,-.95){$\theta$};\draw[guide](P)--(0,-1.414);\draw[guide](1.414,-1.414)--(1.414,-2.55);\draw[<->](0,-2.45)--(1.414,-2.45)node[midway,below]{$a\sin\theta$};
\end{tikzpicture}
\caption{圆环随竖直轴转动，小珠到该轴的距离为 $a\sin\theta$。}\label{rot:fig:d6-e7-q}
\end{figure}
\begin{solution}
在图\ref{rot:fig:d6-e7-q}的随环系中，小珠相对动能为 $ma^2\dot\theta^2/2$，重力势能为 $-mga\cos\theta$。离心势只用到轴距离，故
\[
U_{\rm eff}=-mga\cos\theta-\frac12m\Omega^2a^2\sin^2\theta.
\]
科里奥利力垂直圆环平面，由约束抵消，不进入沿环方程。使用$\dd(\sin^2\theta)/\dd\theta=2\sin\theta\cos\theta$，有效势的一阶导数为
\[
U_{\rm eff}'=mga\sin\theta-m\Omega^2a^2\sin\theta\cos\theta.
\]
代入$ma^2\ddot\theta=-U_{\rm eff}'$，得到
\[
ma^2\ddot\theta=-mga\sin\theta
+m\Omega^2a^2\sin\theta\cos\theta,
\]
即 $\ddot\theta=(\Omega^2\cos\theta-g/a)\sin\theta$。相对平衡角满足
$\sin\theta=0$，或 $\cos\theta=g/(\Omega^2a)$。
\begin{figure}[H]\centering
\begin{tikzpicture}
\draw[axis](-2.2,0)--(2.3,0)node[right]{$\theta$};\draw[axis](0,-.3)--(0,2.3)node[above]{$U_{\rm eff}-U_0$};
\draw[curve,domain=-1.6:1.6,samples=100]plot(\x,{2*(1-cos(\x*180/pi)-.25*sin(\x*180/pi)^2)});\node at(0,-.8){$\Omega^2a/g=1/2$};
\begin{scope}[shift={(6,0)}]
\draw[axis](-2.2,0)--(2.3,0)node[right]{$\theta$};\draw[axis](0,-1.3)--(0,1.3)node[above]{$U_{\rm eff}-U_0$};
\draw[curve,domain=-1.6:1.6,samples=100]plot(\x,{2*(1-cos(\x*180/pi)-sin(\x*180/pi)^2)});
\fill(-1.047,-.5)circle(1.4pt);\fill(1.047,-.5)circle(1.4pt);\node at(0,-1.8){$\Omega^2a/g=2$};
\end{scope}
\end{tikzpicture}
\caption{转速越过临界值后，最低点由势能极小变为极大，两侧出现新的稳定平衡。}\label{rot:fig:d6-e7-s}
\end{figure}
最低点 $\theta=0$ 的线性方程为
$\ddot\eta+(g/a-\Omega^2)\eta=0$。因此当 $\Omega^2<g/a$ 时稳定，小振动角频率为 $\sqrt{g/a-\Omega^2}$；超过该转速则不稳定。环顶 $\theta=\pi$ 的线性系数为 $-(g/a+\Omega^2)$，始终不稳定。
当 $\Omega^2>g/a$，图\ref{rot:fig:d6-e7-s}右图两侧的平衡角为
\[
\theta_* =\pm\arccos\frac g{\Omega^2a}.
\]
在平衡处，运动方程右端对角度的导数为 $-\Omega^2\sin^2\theta_*$，所以两者均稳定，
\[
\omega_*^2=\Omega^2\sin^2\theta_*
=\Omega^2\left[1-\left(\frac g{\Omega^2a}\right)^2\right].
\]
稳定性也可统一由 $U_{\rm eff}''/(ma^2)$ 判断。对有效势求两次导数，
\[
\frac{U_{\rm eff}''}{ma^2}
=\frac ga\cos\theta-\Omega^2(\cos^2\theta-\sin^2\theta).
\]
在非零平衡角处利用 $g/a=\Omega^2\cos\theta_*$，前两项相消，只剩 $\Omega^2\sin^2\theta_*>0$，从而得到两侧平衡的小振动频率。
在临界 $\Omega^2=g/a$ 处，二次势能项为零，但
\[
U_{\rm eff}=U_0+\frac{mga}{8}\theta^4+O(\theta^6).
\]
此时还有精确恒等式
\[
U_{\rm eff}(\theta)-U_{\rm eff}(0)
=\frac{mga}{2}(1-\cos\theta)^2.
\]
最低点附近的小偏离使势能升高，因此它仍稳定。图\ref{rot:fig:d6-e7-s}中临界势阱的井底变平，回复作用从三次项开始，周期随振幅改变，见习题\ref{rot:ex:d6-h7}。

随环系的守恒量是 $E'=ma^2\dot\theta^2/2+U_{\rm eff}$。地面能量满足
$E-E'=\Omega L_z=m\Omega^2a^2\sin^2\theta$，可随珠的位置改变，差额由维持转速的电机提供或回收。
\end{solution}



\begin{example}[高塔落体的东偏]\label{rot:problem-example-8}\label{rot:ex:d6-e8}
在北纬 $0\le\lambda<\pi/2$ 处，质点从高为 $H$ 的塔顶相对地面静止释放。取当地有效重力加速度为 $g$，地球角速度为 $\Omega$，忽略空气阻力。设 $H\ll R_\oplus$、$\Omega\sqrt{H/g}\ll1$。求落地时相对于塔脚的东西偏移，保留到 $\Omega$ 的一阶。
\end{example}
\begin{figure}[H]\centering
\begin{tikzpicture}
\draw[curve](-1,0)--(4,0);\draw[curve](0,0)--(0,3);\fill(0,3)circle(2pt)node[above]{$m$};\node[below left]at(0,0){塔脚};
\draw[axis](0,3)--(0,2)node[right]{$g$};\draw[<->](-.6,0)--(-.6,3)node[midway,left]{$H$};\draw[axis](1,.4)--(3,.4)node[right]{东x};
\node[align=left]at(6.1,2.2){上z：竖直向上\\北y：入纸面\\$\Omega_y=\Omega\cos\lambda$};
\end{tikzpicture}
\caption{局部东、北、上坐标。初始相对地面速度为零。}\label{rot:fig:d6-e8-q}
\end{figure}
\begin{solution}
按图\ref{rot:fig:d6-e8-q}，零阶速度为 $\vct v^{(0)}=(0,0,-gt)$。将它代入科里奥利加速度，
\[
-2\vct\Omega\times\vct v^{(0)}
=(2\Omega gt\cos\lambda,0,0).
\]
因而东西向的一阶方程及初态为
\[
\ddot x=2\Omega gt\cos\lambda,\qquad x(0)=\dot x(0)=0.
\]
用给定的零初速度作第一次定积分，再用零初位移作第二次定积分：
\[
\begin{aligned}
\dot x(t)&=2\Omega g\cos\lambda\int_0^ts\dd s
=2\Omega g\cos\lambda\left[\frac{s^2}{2}\right]_0^t,\\
x(t)&=\Omega g\cos\lambda\int_0^ts^2\dd s
=\Omega g\cos\lambda\left[\frac{s^3}{3}\right]_0^t.
\end{aligned}
\]
因此
\[
\dot x=\Omega gt^2\cos\lambda,\qquad
x=\frac13\Omega gt^3\cos\lambda.
\]
零阶落地时间为 $T=\sqrt{2H/g}$。将一阶东西速度再次代入科里奥利力，得到的竖直和南北修正均为二阶，故求一阶落点偏移时可使用这个 $T$。
\begin{figure}[H]\centering
\begin{tikzpicture}
\draw[axis](0,0)--(4,0)node[right]{$x/\delta$};\draw[axis](0,0)--(0,3.1)node[above]{$z/H$};
\draw[curve,domain=0:1,samples=100]plot({3*\x^3},{2.5*(1-\x^2)});
\fill(0,2.5)circle(1.5pt)node[left]{$1$};\fill(3,0)circle(1.5pt)node[below]{$1$};\node[align=left]at(7,1.8){$x/\delta=(t/T)^3$\\$z/H=1-(t/T)^2$\\东西位移已放大显示};
\end{tikzpicture}
\caption{落体逐渐获得向东速度，落地位置在塔脚东侧。}\label{rot:fig:d6-e8-s}
\end{figure}
图\ref{rot:fig:d6-e8-s}中，落地东偏为
\[
\boxed{\delta=\frac23\Omega H\sqrt{\frac{2H}g}\cos\lambda.}
\]
落地东西速度为 $2\Omega H\cos\lambda$。在给定模型中，赤道附近一阶东偏最大，靠近极点趋零。
也可从绕地轴的角动量理解该方向。记到地轴的距离为 $\rho\simeq(R_\oplus+z)\cos\lambda$，初态为 $\rho_0$。忽略小区域内高阶影响，有
\[
\rho(\Omega\rho+v_x)=\Omega\rho_0^2,
\quad v_x=\Omega\left(\frac{\rho_0^2}\rho-\rho\right)
\simeq2\Omega(H-z)\cos\lambda.
\]
代入 $H-z=gt^2/2$，与前述积分结果相同。下落时到地轴的距离持续减小，质点的地面东向速度相应增加；塔顶与塔脚原有的线速度差只反映了初态的一部分信息。
\end{solution}



\begin{example}[旋转斜管中的不稳定平衡]\label{rot:problem-example-9}\label{rot:ex:d6-e9}
光滑直管与水平面成 $\gamma$ 角，$0<\gamma<\pi/2$，下端 $O$ 位于竖直转轴上。管以恒角速度 $\Omega>0$ 绕该轴转动。质量为 $m$ 的珠沿管坐标为 $s$，$0<s<L$，初始 $s=s_0,\dot s=0$。求相对运动、平衡位置及其稳定性，判断珠先到达哪一端，并求管的两个横向约束力分量。
\end{example}
\begin{figure}[H]\centering
\begin{tikzpicture}
\draw[axis](0,-.3)--(0,4)node[above]{$\vct\Omega,z$};\draw[curve](0,0)--(5,3);\draw[guide](0,0)--(5.4,0);\fill(2.5,1.5)circle(2.2pt)node[above]{$m$};
\draw[guide](2.5,1.5)--(2.5,0);\node[below]at(1.3,0){$s\cos\gamma$};\draw(.9,0)arc(0:31:.9);\node at(1.2,.35){$\gamma$};\node[below left]at(0,0){$O$};\node[right]at(5,3){$L$};
\end{tikzpicture}
\caption{珠到转轴的距离为 $s\cos\gamma$，高度为 $s\sin\gamma$。管提供双侧约束。}\label{rot:fig:d6-e9-q}
\end{figure}
\begin{solution}
图\ref{rot:fig:d6-e9-q}的沿管方向为
$\vct e_s=\cos\gamma\vct e_r+\sin\gamma\vct e_z$。重力沿管分量为 $-mg\sin\gamma$，离心力沿管分量为 $m\Omega^2s\cos^2\gamma$，科里奥利力垂直管所在的竖直平面，因此
\[
\ddot s=\Omega^2s\cos^2\gamma-g\sin\gamma.
\]
令 $\kappa=\Omega\cos\gamma$，平衡位置为
\[
s_* =\frac{g\sin\gamma}{\Omega^2\cos^2\gamma}.
\]
只有$0<s_*<L$时它才位于管内。置$u=s-s_*$，方程化为$\ddot u=\kappa^2u$。特征方程的两根为$\pm\kappa$，故$u=C_1\ee^{\kappa t}+C_2\ee^{-\kappa t}$。代入$u(0)=s_0-s_*$、$\dot u(0)=0$，得到
\[
C_1+C_2=s_0-s_*,\qquad \kappa(C_1-C_2)=0,
\quad C_1=C_2=\frac{s_0-s_*}{2}.
\]
合并指数并对时间求导，得到
\[
s(t)=s_*+(s_0-s_*)\cosh\kappa t,
\qquad \dot s=(s_0-s_*)\kappa\sinh\kappa t.
\]
\begin{figure}[H]\centering
\begin{tikzpicture}
\draw[axis](0,0)--(5.2,0)node[right]{$s$};\draw[axis](0,-1)--(0,2.4)node[above]{$U_{\rm eff}$};
\draw[curve,domain=0:4.8,samples=70]plot(\x,{1.9-.34*(\x-2.4)^2});\draw[guide](2.4,0)--(2.4,1.9);\node[below]at(2.4,0){$s_*$};
\draw[axis](1.6,1.2)--(.8,.6);\draw[axis](3.2,1.2)--(4,.6);\node[align=left]at(7.4,1.1){向内偏离后到O\\向外偏离后到管端\\平衡位于势能极大处};
\end{tikzpicture}
\caption{有效势为开口向下的二次函数，平衡位置不稳定。}\label{rot:fig:d6-e9-s}
\end{figure}
图\ref{rot:fig:d6-e9-s}的有效势为
$U_{\rm eff}=mgs\sin\gamma-m\kappa^2s^2/2$。若 $s_0<s_*$，珠沿管向O运动，到端点时间为
\[
t_O=\frac1\kappa\operatorname{arcosh}\frac{s_*}{s_*-s_0}.
\]
若 $s_0>s_*$，则向外运动并到达L，
\[
t_L=\frac1\kappa\operatorname{arcosh}\frac{L-s_*}{s_0-s_*}.
\]
这两个到端点时间的反双曲余弦参数均大于1。例如 $s_0<s_*$ 时，$s_*/(s_*-s_0)>1$，故存在唯一的正到达时间。由 $\dot s$ 的表达式，其符号在 $t>0$ 始终与 $s_0-s_*$ 相同；图\ref{rot:fig:d6-e9-s}的两支轨道因此不会先转向再到另一端。
只有从 $s_0=s_*$ 相对静止释放时，珠才保持相对平衡，任意小的位置偏差都会增长；若 $s_*\ge L$，管内静止释放的珠均向 $O$ 运动。

求管的约束分量时，取竖直截面内法向
$\vct e_n=-\sin\gamma\vct e_r+\cos\gamma\vct e_z$，方位角正向取 $\vct e_\varphi$。地面加速度为
\[
\vct a=(\ddot s\cos\gamma-\Omega^2s\cos\gamma)\vct e_r
+2\Omega\dot s\cos\gamma\vct e_\varphi
+\ddot s\sin\gamma\vct e_z.
\]
投影到 $\vct e_n$ 时，两个 $\ddot s$ 项抵消，得到
\[
N_n=m(g\cos\gamma+\Omega^2s\sin\gamma\cos\gamma),
\qquad N_\varphi=2m\Omega\dot s\cos\gamma.
\]
管的方位角作用力在地面系做功，电机功率为
$\Omega(s\cos\gamma)N_\varphi=2m\Omega^2s\dot s\cos^2\gamma$。它在珠向外时输入能量、向内时回收能量，与旋转系有效能量守恒并不冲突。
\end{solution}



\begin{example}[科里奥利耦合与表观势能]\label{rot:problem-example-10}\label{rot:ex:d6-e10}
水平面内的质点受各向同性线性回复力 $\vct F=-k\vct r$，固定力心为 $O$，质量为 $m$，记 $\omega_0=\sqrt{k/m}$。采用绕 $O$ 以恒角速度 $\Omega>0$ 逆时针转动的坐标系。初态为 $x'=a,y'=0,\dot x'=\dot y'=0$。求相对运动，判断 $\Omega>\omega_0$ 时轨迹是否发散，并讨论 $\Omega=\omega_0$ 的情形。
\end{example}
\begin{figure}[H]\centering
\begin{tikzpicture}
\draw[axis](-.5,0)--(4,0)node[right]{$x'$};\draw[axis](0,-.7)--(0,2.8)node[above]{$y'$};\fill(0,0)circle(1.5pt)node[below left]{$O$};\fill(2.5,0)circle(2pt)node[above]{$m:a$};
\draw[axis](2.5,0)--(1.1,0)node[midway,below]{$-k\vct r$};\draw[axis](.7,.7)arc(45:135:1);\node at(-1,1.6){$\Omega$};
\end{tikzpicture}
\caption{线性回复力在两个坐标系中都指向O，初始静止条件仅针对转动系。}\label{rot:fig:d6-e10-q}
\end{figure}
\begin{solution}
图\ref{rot:fig:d6-e10-q}中，离心力为 $m\Omega^2(x',y')$，科里奥利力分量为 $(2m\Omega\dot y',-2m\Omega\dot x')$，故
\[
\begin{aligned}
\ddot x'-2\Omega\dot y'+(\omega_0^2-\Omega^2)x'&=0,\\
\ddot y'+2\Omega\dot x'+(\omega_0^2-\Omega^2)y'&=0.
\end{aligned}
\]
当 $\Omega>\omega_0$，有效势 $m(\omega_0^2-\Omega^2)r^2/2$ 的曲率为负。判断运动是否发散，还须解出含科里奥利耦合的两个方程。

置 $\zeta=x'+iy'$，两式合并为
\[
\ddot\zeta+2i\Omega\dot\zeta+(\omega_0^2-\Omega^2)\zeta=0.
\]
再令$\zeta=e^{-i\Omega t}w$，按乘积法则求两次导数：
\[
\dot\zeta=e^{-i\Omega t}(\dot w-i\Omega w),\qquad
\ddot\zeta=e^{-i\Omega t}(\ddot w-2i\Omega\dot w-\Omega^2w).
\]
代入原方程并除以非零的$e^{-i\Omega t}$，
\[
\ddot w-2i\Omega\dot w-\Omega^2w
+2i\Omega(\dot w-i\Omega w)
+(\omega_0^2-\Omega^2)w=0.
\]
两个$\dot w$项抵消，$w$的系数为$-\Omega^2+2\Omega^2+\omega_0^2-\Omega^2=\omega_0^2$，所以$\ddot w+\omega_0^2w=0$。通解为$w=C\cos\omega_0t+D\sin\omega_0t$。由$\zeta(0)=a$得$C=w(0)=a$；由$\dot\zeta(0)=0$得$\dot w(0)=i\Omega a$，即$\omega_0D=i\Omega a$。代回变换，得到
\[
\boxed{\zeta=a e^{-i\Omega t}
\left(\cos\omega_0t+i\frac\Omega{\omega_0}\sin\omega_0t\right).}
\]
\begin{figure}[H]\centering
\begin{tikzpicture}
\draw[axis](0,0)--(5.2,0)node[right]{$\Omega/\omega_0$};\draw[axis](0,0)--(0,3.8)node[above]{$r_{\max}/a$};
\draw[curve](0,1)--(1.5,1)--(4.5,3);\draw[guide](1.5,0)--(1.5,1);\node[below]at(1.5,0){$1$};\node[left]at(0,1){$1$};\node[align=left]at(3.3,.6){负二次有效势\\仍对应有界运动};
\end{tikzpicture}
\caption{角速度可以大于固有频率，轨道半径仍保持有限。图示初始相对速度为零的解。}\label{rot:fig:d6-e10-s}
\end{figure}
图\ref{rot:fig:d6-e10-s}中的半径上界由复坐标的模直接得到：
\[
r^2=|\zeta|^2=a^2\left(\cos^2\omega_0t
+\frac{\Omega^2}{\omega_0^2}\sin^2\omega_0t\right).
\]
故 $r\le a\max(1,\Omega/\omega_0)$，任意有限 $\Omega$ 下都不发散。特征指数 $-i\Omega\pm i\omega_0$ 也均为纯虚数。
当 $\Omega=\omega_0$ 时，括号为 $e^{i\omega_0t}$，因而 $\zeta=a$，质点相对旋转系静止；地面系中它以角速度 $\omega_0$ 作半径为 $a$ 的圆周运动，回复力提供向心力。

本例的地面力场为保守中心力，$E$ 与 $L_z$ 各自守恒。旋转系中的科里奥利项虽不做功，却改变速度方向，使负曲率有效势下仍存在上述有界轨道。
还可用地面守恒量说明半径为何有界。初态给出
\[
L_z=m\Omega a^2\ne0,\qquad E=\frac12m(\omega_0^2+\Omega^2)a^2.
\]
把地面切向动能用 $L_z$ 消去，得
\[
E=\frac12m\dot r^2+\frac{L_z^2}{2mr^2}+\frac12m\omega_0^2r^2.
\]
角动量项在 $r\to0$ 时趋于无穷，回复力势能在 $r\to\infty$ 时趋于无穷，有限能量把半径限制在两者之间。图\ref{rot:fig:d6-e10-s}中的半径上界因此也能由地面守恒量解释，旋转系负二次势能本身未包含这些限制。
\end{solution}



\section{习题与解答}

\begin{exercise}[三角薄板的惯量与支座力]\label{sup6:problem-exercise-1}\label{sup6:ex:d5-h1}
匀质直角等腰三角薄板质量M，顶点为 $O=(0,0)$、$A=(a,0)$、$B=(0,a)$。薄板在光滑水平面内绕O的固定竖直轴转动。在A沿正y方向施力F，初始薄板静止。由面积积分求 $I_O,I_C$，再求初始角加速度和O处的水平支座力。
\end{exercise}
\begin{figure}[H]\centering
\begin{tikzpicture}
\draw[curve](0,0)--(3.5,0)--(0,3.5)--cycle;\node[below left]at(0,0){$O$};\node[below]at(3.5,0){$A$};\node[left]at(0,3.5){$B$};
\draw[axis](3.5,0)--(3.5,1.5)node[right]{$F$};\draw[guide](1.2,0)--(1.2,2.3);\draw[guide](1.4,0)--(1.4,2.1);\node at(1.9,1.8){$\dd x$};
\draw[axis](-.3,-.4)--(4.2,-.4)node[right]{$x$};\draw[axis](-.4,-.3)--(-.4,4.1)node[above]{$y$};
\end{tikzpicture}
\caption{面积积分区域为 $0\le x\le a$、$0\le y\le a-x$。}\label{sup6:fig:d5-h1-q}
\end{figure}
\begin{solution}
图\ref{sup6:fig:d5-h1-q}的面密度为 $\sigma=2M/a^2$。先积分x方向的二次矩，再利用对称性，
\[
\begin{aligned}
I_O&=\sigma\int_0^a\int_0^{a-x}(x^2+y^2)\dd y\dd x\\
&=2\sigma\int_0^a(ax^2-x^3)\dd x
=2\sigma\left[\frac{ax^3}{3}-\frac{x^4}{4}\right]_0^a\\
&=2\sigma\left(\frac{a^4}3-\frac{a^4}4\right)=\frac13Ma^2.
\end{aligned}
\]
积分区域由图\ref{sup6:fig:d5-h1-q}中的直线AB给出：$x+y=a$，因此先固定x时，y的范围为 $0\le y\le a-x$。同一区域内交换x、y不改变积分，才可把两个二次矩合并。质心也可由一次矩求得：
\[
x_C=\frac{\sigma}{M}\int_0^a x(a-x)\,\dd x
=\frac{\sigma}{M}\left[\frac{ax^2}{2}-\frac{x^3}{3}\right]_0^a
=\frac{\sigma a^3}{6M}=\frac a3,
\qquad y_C=x_C.
\]
固定轴O的支持力对O力矩为零，故先由取矩求角加速度，再由质心平动求支座力，可以避免把支座力作为多余未知量引入惯量积分。
质心为 $C=(a/3,a/3)$，故 $OC^2=2a^2/9$，
\[
I_C=I_O-MOC^2=\frac19Ma^2.
\]
力F对O的力矩为aF，逆时针角加速度为 $\alpha=aF/I_O=3F/(Ma)$。初始没有向心加速度，因此
\[
\vct a_C=\vct\alpha\times\overrightarrow{OC}
=\left(-\frac{\alpha a}3,\frac{\alpha a}3\right)
=(-F/M,F/M).
\]
\begin{figure}[H]\centering
\begin{tikzpicture}
\draw[curve](0,0)--(3.3,0)--(0,3.3)--cycle;\fill(1.1,1.1)circle(1.6pt)node[below right]{$C$};\node[below]at(0,0){$O$};
\draw[axis](1.1,1.1)--(.2,2);\node at(.8,1.9){$\vct a_C$};\draw[axis](0,0)--(-1.3,0)node[above]{$H_x=-F$};\draw[axis](3.3,0)--(3.3,1.3)node[right]{$F$};
\draw[axis](4.1,1.6)arc(0:100:.65);\node at(4,2.7){$\alpha>0$};
\end{tikzpicture}
\caption{质心沿垂直OC的方向开始加速，支座只需提供向左的水平分量。}\label{sup6:fig:d5-h1-s}
\end{figure}
图\ref{sup6:fig:d5-h1-s}中，由 $\vct H+(0,F)=M\vct a_C=(-F,F)$，得 $H_x=-F,H_y=0$。本题y轴位于桌面内；垂直桌面的重力另由桌面支持力平衡。
\end{solution}


\begin{exercise}[自由径向滑块与转台，选做]\label{sup6:problem-exercise-3}\label{sup6:ex:d5-h3}
转台本体惯量为 $I_0$，可绕固定竖直轴无摩擦转动。质量m的滑块可沿转台上的光滑径向直导轨自由移动，无电机或径向弹簧。初始半径 $r_0>0$、径向速度零、共同角速度 $\omega_0>0$。求 $\omega(r)$ 和向外运动的径向速度。若 $I_0=mr_0^2$，求滑块首次到达 $2r_0$ 时的两种速度。
\end{exercise}
\begin{figure}[H]\centering
\begin{tikzpicture}
\draw[curve](0,0)circle(2.2);\draw[curve](0,0)--(2.2,0);\fill(0,0)circle(1.5pt)node[below]{$O$};\draw[fill=white](.9,-.15)rectangle(1.2,.15);\node[above]at(1.05,.2){$m$};\node[below]at(.6,-.1){$r_0$};
\draw[axis](-1.4,1.4)arc(135:210:2);\node[left]at(-2,.4){$\omega_0$};\node[right]at(2.5,.7){导轨足够长};
\end{tikzpicture}
\caption{滑块径向无驱动力，导轨仅施加切向约束。}\label{sup6:fig:d5-h3-q}
\end{figure}
\begin{solution}
图\ref{sup6:fig:d5-h3-q}中系统没有轴向外力矩，记 $I_a=I_0+mr_0^2$，则
\[
H=I_a\omega_0,
\qquad \omega(r)=\frac H{I_0+mr^2}.
\]
机械能守恒还需要核对接触力的功率。设导轨对滑块的切向力为 $N_\varphi$。滑块与导轨接触处的切向速度相同，均为 $r\omega$，径向相对速度则与导轨力垂直。于是
\[
\mathcal P_{\text{块}}=N_\varphi r\omega,\qquad
\mathcal P_{\text{台}}=(-rN_\varphi)\omega,
\qquad \mathcal P_{\text{块}}+\mathcal P_{\text{台}}=0.
\]
图\ref{sup6:fig:d5-h3-s}表示机械能从转台传给滑块。两个接触力各自做功，总功率相消依靠的是接触处相同的切向速度和相反的切向作用力。
导轨力对滑块与转台的总功率为零，轴承不做功，所以机械能守恒：
\[
\frac12m\dot r^2+\frac{H^2}{2(I_0+mr^2)}=\frac{H^2}{2I_a}.
\]
向外运动取正根，得到
\[
\dot r=\frac H{\sqrt m}\sqrt{\frac1{I_a}-\frac1{I_0+mr^2}},\qquad r\ge r_0.
\]
\begin{figure}[H]\centering
\begin{tikzpicture}
\draw[axis](0,0)--(5,0)node[right]{$r/r_0$};\draw[axis](0,0)--(0,2.8)node[above]{归一化速度};
\draw[curve,domain=1:3.1,samples=70]plot({1.4*\x},{4/(1+\x*\x)});
\draw[guide,domain=1:3.1,samples=90]plot({1.4*\x},{2*sqrt(2-4/(1+\x*\x))});
\node[below]at(1.4,0){$1$};\node[below]at(2.8,0){$2$};\node[right]at(4.3,.5){$\omega/\omega_0$};\node at(3.4,3.2){$\dot r/(r_0\omega_0)$};
\node[draw,minimum width=3.8cm,minimum height=.65cm]at(8.5,2.2){滑块：$+N_\varphi r\omega$};
\node[draw,minimum width=3.8cm,minimum height=.65cm]at(8.5,.2){转台：$-N_\varphi r\omega$};
\draw[axis](8.5,.6)--(8.5,1.8)node[midway,right]{$N_\varphi r\omega$};
\node at(8.5,-.7){两接触力的总功率为零};
\end{tikzpicture}
\caption{左图取 $I_0=mr_0^2$，转动减慢而径向速率增加；右图表示向外滑动时由转台传给滑块的功率。}\label{sup6:fig:d5-h3-s}
\end{figure}
在图\ref{sup6:fig:d5-h3-s}所用参数下，$I_a=2mr_0^2$、$H=2mr_0^2\omega_0$。到 $r=2r_0$ 时，
\[
\omega_2=\frac25\omega_0,
\qquad \dot r_2^2=4mr_0^4\omega_0^2
\left(\frac1{2mr_0^2}-\frac1{5mr_0^2}\right)
=\frac65r_0^2\omega_0^2.
\]
滑块切向速度为 $r\omega=4r_0\omega_0/5$，径向速度为 $\sqrt{6/5}\,r_0\omega_0$。释放瞬间径向速率虽然为零，径向方程 $\ddot r=r\omega^2$ 仍给出正的初始加速度，故存在向外运动分支。
\end{solution}


\begin{exercise}[水平拉力的作用高度]\label{deep6:problem-exercise-5}\label{deep6:ex:d5-h5}
质量M、半径R、惯量 $\beta MR^2$ 的轮静置在粗糙水平面上。对轮施加向右水平力F，其初始作用线在轮心C上方带符号距离q处，$-R\le q\le R$。求开始作用瞬间纯滚动的加速度和摩擦力，讨论摩擦方向及摩擦为零的作用高度。仅讨论初始瞬间。
\end{exercise}
\begin{figure}[H]\centering
\begin{tikzpicture}
\draw[curve](0,0)circle(1.4);\draw[curve](-2,-1.4)--(3,-1.4);\fill(0,0)circle(1.5pt)node[below right]{$C$};
\draw[axis](0,.8)--(2.8,.8)node[right]{$F$};\draw[<->](-.5,0)--(-.5,.8)node[midway,left]{$q$};\draw[guide](0,0)--(2,0);\node[below]at(0,-1.4){$P$};
\end{tikzpicture}
\caption{q向上为正，图示 $q>0$。正力矩按顺时针取。}\label{deep6:fig:d5-h5-q}
\end{figure}
\begin{solution}
图\ref{deep6:fig:d5-h5-q}中，设摩擦分量f向右为正。水平平动、顺时针转动及无滑动约束分别为
\[
Ma=F+f,\qquad \beta MR^2\alpha=Fq-fR,
\qquad a=R\alpha.
\]
消去角加速度并将前两式相加，
\[
(1+\beta)Ma=F(1+q/R).
\]
所以
\[
a=\frac{F(1+q/R)}{(1+\beta)M},\qquad
f=\frac{F(q/R-\beta)}{1+\beta}.
\]
\begin{figure}[H]\centering
\begin{tikzpicture}
\draw[axis](0,0)--(6,0)node[right]{$q/R$};\draw[axis](1.5,-1.6)--(1.5,1.8)node[above]{$(1+\beta)f/F$};
\draw[curve](0,-1.5)--(3,.5);\node[below]at(0,0){$-1$};\node[below]at(2.25,0){$\beta$};\node[below]at(3,0){$1$};\node at(3.7,1.7){$f>0$};\node at(.5,-1.9){$f<0$};
\end{tikzpicture}
\caption{示意 $\beta=1/2$ 的摩擦变号，横轴零点位于纵轴处。}\label{deep6:fig:d5-h5-s}
\end{figure}
图\ref{deep6:fig:d5-h5-s}中的摩擦变号也可从初始滑动趋势看出。若暂不计摩擦，轮心加速度为 $F/M$，顺时针角加速度为 $Fq/(\beta MR^2)$，下接触点的水平加速度为
\[
a_{P,x}^{(f=0)}=\frac FM\left(1-\frac{q}{\beta R}\right).
\]
它向右时，维持无滑动需要向左的摩擦；它向左时，需要向右的摩擦。$q=\beta R$ 恰使这项为零，与联立方程得到的零摩擦位置相同。
当 $q<\beta R$ 时摩擦向左，$q>\beta R$ 时向右，$q=\beta R$ 时为零。图\ref{deep6:fig:d5-h5-s}的变号对应拉力矩与纯滚动所需角加速度之间的不同关系。
竖直平衡给 $N=Mg$，还须满足
\[
\frac{F|q/R-\beta|}{1+\beta}\le\mu_sMg.
\]
若 $\beta=1/2$，力沿轮心上方 $R/2$ 作用便能在初始瞬间无需摩擦地实现 $a=R\alpha$。这一结果不保证后续仍无需摩擦，因为题设未规定力的作用线如何随轮转动。
\end{solution}


\begin{exercise}[突然获得速度的板]\label{deep6:problem-exercise-7}\label{deep6:ex:d5-h7}
质量M的长板在光滑水平面上，以初速度U向右运动；其上质量同为M、半径R的匀质实心圆柱初始不平移、不转动。板与圆柱间动摩擦系数为 $\mu$，此后无水平外力。求滑动终止时两者的速度、圆柱角速度、经过时间与系统能量损失。板足够长。
\end{exercise}
\begin{figure}[H]\centering
\begin{tikzpicture}
\draw[curve](-2,-1.3)rectangle(4,-1);\draw[curve](0,0)circle(1);\node at(0,1.4){圆柱静止};\draw[axis](2,-.8)--(3.3,-.8)node[right]{$U$};\node[below]at(1,-1.5){板M，地面光滑};\node at(0,0){$M$};
\end{tikzpicture}
\caption{初始板向右滑过圆柱的底部，摩擦作用使圆柱向右并逆时针加速。}\label{deep6:fig:d5-h7-q}
\end{figure}
\begin{solution}
按图\ref{deep6:fig:d5-h7-q}，取圆柱向右平动、逆时针转动为正。设从初态累积的摩擦冲量为J。由于两者质量相同，
\[
v_C=\frac JM,\quad V=U-\frac JM,\quad
\omega=\frac{JR}{MR^2/2}=\frac{2J}{MR}.
\]
图\ref{deep6:fig:d5-h7-q}中，圆柱底部材料点因逆时针自转而向右运动，所以它的地面速度为 $v_C+R\omega$。还要减去板的速度V，才是决定摩擦方向的相对速度。只比较两个质心速度，会漏掉自转对滑动终止条件的贡献。
相对于板的接触滑动速度是 $u_P=v_C+R\omega-V$，所以
\[
u_P=\frac{4J}M-U.
\]
首次无滑动时 $J_r=MU/4$，得到
\[
v_C=\frac U4,\qquad V=\frac{3U}4,
\qquad \omega=\frac U{2R}\quad\text{（逆时针）}.
\]
两者总水平动量为 $M(v_C+V)=MU$，与初态相同。
\begin{figure}[H]\centering
\begin{tikzpicture}
\draw[axis](0,0)--(4.7,0);\node at(3.8,-.8){$J/(MU)$};\draw[axis](0,-2)--(0,1.3)node[above]{$u_P/U$};
\draw[curve](0,-1.5)--(3,0);\node[left]at(0,-1.5){$-1$};\node[below right]at(3,0){$1/4$};
\begin{scope}[shift={(7,0)}]
\draw[curve](0,0)circle(.85);\draw[curve](-1.3,-.85)--(1.5,-.85);\draw[axis](0,-.85)--(1.35,-.85)node[right]{$3U/4$};\draw[axis](0,0)--(.45,0);\node[above]at(.2,.08){$U/4$};\node at(0,1.3){$R\omega=U/2$};
\end{scope}
\end{tikzpicture}
\caption{滑动速度由负值增至零；右图的下接触点速度为 $U/4+U/2=3U/4$。}\label{deep6:fig:d5-h7-s}
\end{figure}
图\ref{deep6:fig:d5-h7-s}所示过程内 $N=Mg$，动摩擦恒为 $\mu Mg$。因此 $J=\mu Mgt$，
\[
t_r=\frac U{4\mu g}.
\]
无滑动后系统无水平外力，保持当前速度只需零摩擦。滑动阶段两接触材料点的相对速率为$U-4\mu gt$，因此耗散功率为$\mu Mg(U-4\mu gt)$。对时间积分，
\[
\begin{aligned}
Q&=\int_0^{t_r}\mu Mg(U-4\mu gt)\dd t
=\mu Mg\bigl[Ut-2\mu gt^2\bigr]_0^{t_r}\\
&=\mu Mg\left(\frac{U^2}{4\mu g}-\frac{U^2}{8\mu g}\right)
=\frac18MU^2.
\end{aligned}
\]
所以末态总动能为$K_r=MU^2/2-Q=3MU^2/8$。热量来自接触点的相对滑动；摩擦对板和圆柱各自的功不能分别当成全部耗散。
\end{solution}


\begin{exercise}[两个固定轴圆盘的摩擦传动]\label{deep6:problem-exercise-8}\label{deep6:ex:d5-h8}
两个相同圆盘半径R、轴向惯量I，轴心 $O_1,O_2$ 固定，两轮缘在P处外切。最初左盘逆时针角速度为 $\Omega_0$，右盘静止；随后接触处发生滑动摩擦，直至两轮缘无滑动。轴承无摩擦。求末角速度、能量损失，并判断两个圆盘的总角动量是否守恒。若滑动摩擦大小恒为f，求所需时间。
\end{exercise}
\begin{figure}[H]\centering
\begin{tikzpicture}
\draw[curve](-1.3,0)circle(1.3);\draw[curve](1.3,0)circle(1.3);\fill(-1.3,0)circle(1.5pt)node[below]{$O_1$};\fill(1.3,0)circle(1.5pt)node[below]{$O_2$};\fill(0,0)circle(1.5pt)node[below]{$P$};
\draw[axis](-2.1,.8)arc(135:220:1.13);\node[left]at(-2.7,.4){$\Omega_0$};\node at(1.3,1.7){初始静止};
\end{tikzpicture}
\caption{正角速度统一取逆时针。接触前左轮缘在P处向上运动。}\label{deep6:fig:d5-h8-q}
\end{figure}
\begin{solution}
按图\ref{deep6:fig:d5-h8-q}的符号，左盘P点速度为 $R\omega_1$ 向上，右盘P点速度为 $-R\omega_2$ 向上。初始相对滑动向上，摩擦对左盘向下、对右盘向上。两力对各自轴心的力矩都为顺时针，故
\[
I\dot\omega_1=-fR,\qquad I\dot\omega_2=-fR.
\]
若把滑动摩擦冲量记作$J_f=\int_0^t f(s)\dd s$，对两个角方程分别积分，
\[
I\int_{\Omega_0}^{\omega_1}\dd\omega
=-R\int_0^tf(s)\dd s=-RJ_f,\qquad
I\int_0^{\omega_2}\dd\omega=-RJ_f.
\]
左端分别为$I(\omega_1-\Omega_0)$和$I\omega_2$，所以两盘的角速度为
\[
\omega_1=\Omega_0-\frac{RJ_f}{I},\qquad
\omega_2=-\frac{RJ_f}{I}.
\]
因此末态角速度不依赖滑动摩擦随时间怎样变化；恒定f只在最后求时间时才需要。图\ref{deep6:fig:d5-h8-s}中的外轴力又说明，角速度之差不变不能被解释为两个圆盘的总角动量守恒。
相减得到 $\omega_1-\omega_2=\Omega_0$。无滑动要求 $R\omega_1=-R\omega_2$，联立这两个线性关系，
\[
\omega_{1f}=\frac{\Omega_0}2,
\qquad \omega_{2f}=-\frac{\Omega_0}2.
\]
两圆盘向相反方向转动，接触材料点的速度却相同。
\begin{figure}[H]\centering
\begin{tikzpicture}
\draw[curve](-1.3,0)circle(1.3);\draw[curve](1.3,0)circle(1.3);
\fill(0,0)circle(1.4pt);\draw[axis](0,0)--(0,-.9);\node at(-.45,-.5){$f$};\draw[axis](0,0)--(0,.9);\node at(.45,.5){$f$};
\draw[axis](-1.3,0)--(-1.3,1);\node at(-1.8,.5){$H_1$};\draw[axis](1.3,0)--(1.3,-1);\node at(1.8,-.5){$H_2$};\node[below]at(0,-1.7){两个轴力构成顺时针外力偶};
\end{tikzpicture}
\caption{图中仅画与滑动摩擦对应的轴力分量，径向压力及其轴力分量未画。}\label{deep6:fig:d5-h8-s}
\end{figure}
初动能为 $I\Omega_0^2/2$，末动能为两项 $I(\Omega_0/2)^2/2$ 之和，故
\[
Q=\frac14I\Omega_0^2.
\]
总自转角动量从 $I\Omega_0$ 减为零。图\ref{deep6:fig:d5-h8-s}中的轴力维持轴心固定，关于P的总外力矩为 $-2fR$，这正是总角动量改变的原因。轮缘摩擦在整个系统中是内力；两条转动方程相减后保持不变的是 $I(\omega_1-\omega_2)$。
当f恒定时，$\omega_1=\Omega_0-fRt/I$、$\omega_2=-fRt/I$，用 $\omega_1+\omega_2=0$ 得
\[
t_r=\frac{I\Omega_0}{2fR}.
\]
此后轴承无摩擦且不需传递负载，维持匀速反向转动所需的接触切向力为零。
\end{solution}


\begin{exercise}[绕绳下降的线轴]\label{deep6:problem-exercise-9}\label{deep6:ex:d5-h9}
线轴质量M，绕中心轴的惯量I，细绳缠绕在半径r的轴上，绳的上端固定。线轴从静止释放，下降时绳保持竖直、无滑动地展开，线轴中心沿竖直线运动。忽略绳质量，求下降加速度、绳张力、下降高度h后的速率和转速，说明张力是否做功。
\end{exercise}
\begin{figure}[H]\centering
\begin{tikzpicture}
\draw[curve](-1,3)--(1.5,3);\draw[curve](.55,3)--(.55,0);\draw[curve](0,0)circle(1.1);\draw[curve](0,0)circle(.55);\fill(0,0)circle(1.5pt)node[below left]{$C$};\draw[guide](0,0)--(-.55,0)node[midway,above]{$r$};
\draw[axis](.55,0)--(.55,1.3)node[right]{$T$};\draw[axis](0,0)--(0,-1.5)node[left]{$Mg$};\node[below right]at(.55,0){$Q$};
\end{tikzpicture}
\caption{绳在内轴右侧切向展开，张力产生逆时针力矩。}\label{deep6:fig:d5-h9-q}
\end{figure}
\begin{solution}
图\ref{deep6:fig:d5-h9-q}中，取向下加速度a和逆时针角加速度 $\alpha$ 为正。绳不滑动，展开长度等于轴转过的弧长，故 $a=r\alpha$。平动与转动方程为
\[
Ma=Mg-T,\qquad I\alpha=Tr.
\]
消去T、$\alpha$，得到
\[
a=\frac{g}{1+I/(Mr^2)},\qquad
T=\frac{MgI}{Mr^2+I}.
\]
两者均为常量，$0<T<Mg$，下降过程加速且绳保持受拉。从静止开始，由 $\dot v=v\dd v/\dd h=a$ 积分，
\[
\int_0^v\xi\dd\xi=\int_0^ha\dd\eta,
\qquad [\xi^2/2]_0^v=a[\eta]_0^h,
\]
即 $v^2=2ah$。再用 $v=r\omega$，
\[
v=\sqrt{\frac{2gh}{1+I/(Mr^2)}},\qquad \omega=\frac vr.
\]
\begin{figure}[H]\centering
\begin{tikzpicture}
\draw[curve](0,0)circle(1.1);\draw[curve](0,0)circle(.55);\fill(.55,0)circle(1.7pt)node[right]{$Q$};
\draw[axis](.55,0)--(.55,-1.2)node[left]{$v_C$};\draw[axis](.55,0)--(.55,1.2)node[right]{$r\omega$};
\node[align=left]at(4.5,0){$v_Q=-v_C+r\omega=0$\\张力作用点瞬时静止\\$\mathcal P_T=0$};
\end{tikzpicture}
\caption{绳接触点的平移速度与转动速度抵消。图中竖直速度以向上为正。}\label{deep6:fig:d5-h9-s}
\end{figure}
图\ref{deep6:fig:d5-h9-s}中张力的作用点Q瞬时静止，张力不做功。质心平动能的增长率包含负项 $-Tv$，自转动能的增长率包含正项 $Tr\omega=Tv$，相加后抵消。因此系统的能量式为
\[
Mgh=\frac12Mv^2+\frac12I\omega^2,
\]
代入 $\omega=v/r$，可得到前述速率。张力的功率由材料点Q的速度确定；质心平动能式中的张力项只计入平动部分。
\end{solution}


\begin{exercise}[有限摩擦下的内环整周滚动，选做]\label{deep6:problem-exercise-10}\label{deep6:ex:d5-h10}
半径r的球在半径R的固定竖直圆形轨道内侧运动，$r<R$，球心轨迹半径 $a=R-r$，球的惯量为 $I_C=\beta mr^2$。底部初始速率为 $v_0$，同时满足纯滚动。静摩擦系数为 $\mu>0$。求保持接触且不打滑地完成整周运动所需的最小 $v_0$，并指出最先达到摩擦极限的位置。
\end{exercise}
\begin{figure}[H]\centering
\begin{tikzpicture}
\draw[curve](0,0)circle(2.5);\draw[guide](0,0)circle(1.9);\fill(0,0)circle(1.5pt)node[left]{$O$};
\draw[curve](1.556,-1.09)circle(.6);\fill(1.556,-1.09)circle(1.5pt)node[above]{$C$};\draw[guide](0,0)--(1.556,-1.09)node[midway,above]{$a$};
\draw[guide](0,0)--(0,-1.9);\draw(0,-.7)arc(-90:-35:.7);\node at(.5,-.8){$\theta$};\draw[axis](0,-1.9)--(1,-1.9)node[below]{$v_0$};
\end{tikzpicture}
\caption{角度从最低点量起，球心轨道半径为 $R-r$。}\label{deep6:fig:d5-h10-q}
\end{figure}
\begin{solution}
图\ref{deep6:fig:d5-h10-q}中，质心上升高度为 $a(1-\cos\theta)$。纯滚动时能量守恒，
\[
v^2=v_0^2-\frac{2ga}{1+\beta}(1-\cos\theta).
\]
沿切向取 $\theta$ 增大的方向为正。平动、转动关系给出
\[
m\dot v=-mg\sin\theta+f,\qquad
\beta m\dot v=-f,
\]
所以 $f=\beta mg\sin\theta/(1+\beta)$。在后半周f会变号，摩擦条件须使用绝对值。向圆心列法向方程，
\[
N-mg\cos\theta=\frac{mv^2}a.
\]
置 $\lambda=v_0^2/(ga)$，则
\[
\frac{N}{mg}=\lambda-\frac2{1+\beta}
+\frac{3+\beta}{1+\beta}\cos\theta.
\]
保持接触且不打滑，需要对所有角度满足 $N\ge |f|/\mu$，即
\[
(1+\beta)\lambda\ge
2-(3+\beta)\cos\theta+\frac\beta\mu|\sin\theta|.
\]
\begin{figure}[H]\centering
\begin{tikzpicture}
\draw[axis](0,0)--(4.5,0)node[right]{$-\cos\theta$};\draw[axis](0,0)--(0,2.6)node[above]{$\sin\theta$};
\draw[guide](0,0)circle(1.6);\draw[axis](0,0)--(3.3,1.2)node[above right]{$(3+\beta,\beta/\mu)$};
\draw[curve](0,0)--(1.504,.547);\fill(1.504,.547)circle(1.8pt)node[above]{$\theta_*$};
\node[align=left]at(6.8,.7){两个向量同向时\\内积取最大值};
\end{tikzpicture}
\caption{在前半周，用单位向量 $(-\cos\theta,\sin\theta)$ 求线性组合的最大值。}\label{deep6:fig:d5-h10-s}
\end{figure}
前半周取 $\vct u=(-\cos\theta,\sin\theta)$，它沿单位圆上半周运动。令 $\vct b=(3+\beta,\beta/\mu)$，则需最大化的角度项为 $\vct b\cdot\vct u$。因为 $\vct b$ 的两个分量都为正，$\vct u=\vct b/|\vct b|$ 对应的角度落在 $\pi/2<\theta<\pi$，确实属于可经过的位置。
图\ref{deep6:fig:d5-h10-s}中，柯西不等式给出右端的最大值为
$2+\sqrt{(3+\beta)^2+(\beta/\mu)^2}$，等号可以在前半周达到。因此
\[
\boxed{v_{0,\min}^2=
\frac{ga}{1+\beta}\left[2+\sqrt{(3+\beta)^2+(\beta/\mu)^2}\right].}
\]
前半周首先达到极限的位置满足
\[
\cos\theta_*=-\frac{3+\beta}{\sqrt{(3+\beta)^2+(\beta/\mu)^2}},\qquad
\sin\theta_*=\frac{\beta/\mu}{\sqrt{(3+\beta)^2+(\beta/\mu)^2}}.
\]
后半周有关于最高点对称的另一个极限位置。这里的最小初速保证最高点 $v^2\ge ga>0$，全程不会先发生速度转向；对有限 $\mu$，最高点反力也严格为正。
取实心球的惯量系数 $\beta=2/5$，结果化为
\[
v_{0,\min}^2=\frac{ga}{7}\left[10+\sqrt{289+4/\mu^2}\right].
\]
当 $\mu\to\infty$ 时趋于 $27ga/7$，即只由最高点接触条件给出的极限。有限摩擦下，最严格的位置在最高点之前，不能只检查环顶反力。
\end{solution}


\begin{exercise}[加速悬点上的复摆]\label{deep6:problem-exercise-11}\label{deep6:ex:d5-h11}
长L、质量M的匀质细杆，一端铰接在作恒定水平加速度A向右运动的车架上。杆在竖直平面内摆动，悬点O始终不发生竖直运动。求相对车架的稳定平衡方向、小振动周期及平衡状态下的铰链力。取杆与向下竖直方向的夹角 $\theta$ 向右为正。
\end{exercise}
\begin{figure}[H]\centering
\begin{tikzpicture}
\draw[curve](-1,1.2)--(2,1.2);\fill(0,1.2)circle(1.7pt)node[above]{$O$};\draw[guide](0,1.2)--(0,-2);
\draw[curve](0,1.2)--(1.6,-1.8);\fill(.8,-.3)circle(1.5pt)node[right]{$C$};\draw[axis](.8,1.8)--(2.3,1.8)node[right]{$A$};
\draw(0,.4)arc(-90:-62:.8);\node at(.35,.15){$\theta$};\node[left]at(-.3,-1.1){竖直线};
\end{tikzpicture}
\caption{图示一般偏角；稳定位置将在竖直线左侧。}\label{deep6:fig:d5-h11-q}
\end{figure}
\begin{solution}
图\ref{deep6:fig:d5-h11-q}中，O为加速材料点，不能直接写成固定轴的重力矩方程。取 $\vct r_{CO}=(L/2)(\sin\theta,-\cos\theta)$，由式\eqref{sup6:eq:d5-material-origin}，
\[
I_O\ddot\theta=-\frac{ML}2g\sin\theta
-\frac{ML}2A\cos\theta,
\qquad I_O=\frac13ML^2.
\]
也可在随车平移的参考系中，在C加向左惯性力MA，得到同一转动方程。令
\[
G=\sqrt{g^2+A^2},\qquad \phi=\arctan(A/g),
\]
则 $g\sin\theta+A\cos\theta=G\sin(\theta+\phi)$。稳定平衡方向为 $\theta_*=-\phi$，杆沿向下且向左的有效重力方向。
\begin{figure}[H]\centering
\begin{tikzpicture}
\fill(0,1.7)circle(1.7pt)node[above]{$O$};\draw[curve](0,1.7)--(-1.5,-1.3);\fill(-.75,.2)circle(1.7pt)node[below right]{$C$};
\draw[axis](-.75,.2)--(-.75,-1.1)node[right]{$Mg$};\draw[axis](-.75,.2)--(-1.4,.2)node[above]{$MA$};\draw[axis,dashed](-.75,.2)--(-1.4,-1.1)node[left]{$MG$};
\node[align=left]at(3.5,.6){随车系：惯性力向左\\地面系：质心加速度向右};
\end{tikzpicture}
\caption{稳定平衡由有效重力确定，惯性力只用于随车系的转动方程。}\label{deep6:fig:d5-h11-s}
\end{figure}
稳定性可由随车系的势能确认。以向上为竖直正向，重力势与平移惯性力势之和为
\[
U_{\rm eff}=-\frac{MLg}{2}\cos\theta+\frac{MLA}{2}\sin\theta
=-\frac{MLG}{2}\cos(\theta+\phi).
\]
对偏角求导时，$\phi$ 是常量，因此
\[
U_{\rm eff}'=\frac{MLG}{2}\sin(\theta+\phi),\qquad
U_{\rm eff}''=\frac{MLG}{2}\cos(\theta+\phi).
\]
在 $\theta_*=-\phi$ 处，$U_{\rm eff}'=0$、$U_{\rm eff}''=MLG/2>0$。图\ref{deep6:fig:d5-h11-s}中杆沿有效重力方向悬垂，正是这一势能极小位置；反向沿有效重力悬立的位置则不稳定。
在图\ref{deep6:fig:d5-h11-s}的稳定方向附近，置 $\eta=\theta+\phi$，小角度方程为
\[
\ddot\eta+\frac{3G}{2L}\eta=0,\qquad
T=2\pi\sqrt{\frac{2L}{3G}}.
\]
求实际铰链力回到地面系。相对车架平衡时杆只平移，质心加速度为 $(A,0)$，所以
\[
H_x=MA,\qquad H_y-Mg=0,
\qquad |\vct H|=M\sqrt{A^2+g^2}.
\]
铰链力向右上方，恰好沿OC的反方向；当 $A\to0$ 时，恢复端点悬杆的小振动周期 $2\pi\sqrt{2L/(3g)}$。
\end{solution}


\begin{exercise}[减速、停止与反向滚动]\label{deep6:problem-exercise-12}\label{deep6:ex:d5-h12}
惯量为 $\beta MR^2$ 的轮在固定斜面上沿坡向下纯滚动，初速度为 $v_0$。从某时刻开始，对轮心施加沿坡向上的恒力F，满足 $F>Mg\sin\gamma$，斜面倾角为 $\gamma$。求纯滚动时的加速度、摩擦方向、停止前位移和时间，并说明停止后摩擦是否立即反向。给出纯滚动所需的摩擦条件。
\end{exercise}
\begin{figure}[H]\centering
\begin{tikzpicture}
\begin{scope}[rotate=-20]
\draw[curve](-3,0)--(4,0);\draw[curve](0,1)circle(1);\fill(0,1)circle(1.5pt)node[above]{$C$};
\draw[axis](0,1)--(-1.6,1)node[left]{$F$};\draw[axis](.3,1.4)--(1.6,1.4)node[right]{$v_0$};\draw[axis](0,0)--(1.1,0)node[below]{$f$};
\end{scope}
\draw[guide](-2.819,1.026)--(-1.6,1.026);\draw(-2.219,1.026)arc(0:-20:.6);\node at(-1.95,.77){$\gamma$};
\end{tikzpicture}
\caption{轮仍沿坡向下运动时，足够大的向上拉力已使摩擦转为沿坡向下。}\label{deep6:fig:d5-h12-q}
\end{figure}
\begin{solution}
以图\ref{deep6:fig:d5-h12-q}中沿坡向下为正，转动正向取向下滚动时的方向。轮心拉力没有绕质心的力矩，因此
\[
Ma=Mg\sin\gamma-F+f,\qquad
\beta MR^2\alpha=-fR,\qquad a=R\alpha.
\]
由后两式 $f=-\beta Ma$，代回平动方程，令 $B=F/M-g\sin\gamma>0$，得到
\[
a=-\frac B{1+\beta},\qquad
f=\frac{\beta MB}{1+\beta}>0.
\]
摩擦沿坡向下，与此时质心速度同向；角加速度为负，因而自转和平动同时减慢。
\begin{figure}[H]\centering
\begin{tikzpicture}
\draw[axis](0,0)--(5.1,0)node[right]{$t$};\draw[axis](0,-1.7)--(0,2.2)node[above]{$v_C$};
\draw[curve](0,1.7)--(4.5,-1.3);\node[left]at(0,1.7){$v_0$};\node[above right]at(2.55,0){$t_s$};
\draw[axis](1,-1.7)--(2,-1.7)node[right]{摩擦沿坡向下};\node at(4.2,1.3){停止后向上滚动};
\end{tikzpicture}
\caption{恒定加速度使速度连续过零，摩擦方向在过零处不改变。}\label{deep6:fig:d5-h12-s}
\end{figure}
从开始施力的时刻计时，积分恒加速度方程，得
\[
v_C(t)-v_0=\int_0^t-\frac{B}{1+\beta}\,\dd\tau
=-\frac{Bt}{1+\beta}.
\]
令 $v_C(t_s)=0$，图\ref{deep6:fig:d5-h12-s}中的停止时刻为 $t_s=(1+\beta)v_0/B$。再对速度积分，停止前位移为
\[
\begin{aligned}
s_s&=\int_0^{t_s}\left(v_0-\frac{B\tau}{1+\beta}\right)\dd\tau
=\left[v_0\tau-\frac{B\tau^2}{2(1+\beta)}\right]_0^{t_s}\\
&=\frac{(1+\beta)v_0^2}{B}-\frac{(1+\beta)v_0^2}{2B}
=\frac{(1+\beta)v_0^2}{2B}.
\end{aligned}
\]
法向力为 $N=Mg\cos\gamma$，保持纯滚动需要
\[
\frac{\beta B}{1+\beta}\le\mu_sg\cos\gamma.
\]
若该条件满足，轮在 $t_s$ 后沿坡向上加速，a、f仍保持原值。纯滚动的接触点相对地面始终静止，质心速度过零并未引起滑动方向改变，因而摩擦不在此时反向。
停止前，静摩擦作用点瞬时静止，重力做功 $Mg\sin\gamma\,s_s$，拉力做功 $-Fs_s$，两者之和为
\[
-(F-Mg\sin\gamma)s_s=-MBs_s
=-\frac12(1+\beta)Mv_0^2,
\]
恰好等于初始平动能和转动能的减少量。
\end{solution}


\begin{exercise}[粗糙斜面的相对静止区间]\label{rot:problem-exercise-1}\label{rot:ex:d6-h1}
斜面向右上方升高，倾角为 $\gamma$，姿态保持不变，水平加速度为 $A$，向右为正。物块与斜面的静摩擦系数为 $\mu$，满足 $\mu\tan\gamma<1$。求物块能相对斜面静止的全部 $A$ 值，并说明摩擦方向随 $A$ 的变化。
\end{exercise}
\begin{figure}[H]\centering
\begin{tikzpicture}
\draw[curve](0,0)--(5,0)--(5,2.887)--cycle;\fill(2.6,1.501)circle(2.4pt)node[above]{$m$};\draw[axis](2.6,1.501)--(3.7,2.136)node[right]{$f>0$};\draw[axis](1,3.2)--(3,3.2)node[right]{$A$};\draw(.8,0)arc(0:30:.8);\node at(1.1,.3){$\gamma$};
\end{tikzpicture}
\caption{摩擦分量 $f$ 沿坡向上为正，负值表示实际摩擦沿坡向下。}\label{rot:fig:d6-h1-q}
\end{figure}
\begin{solution}
在图\ref{rot:fig:d6-h1-q}的平移系中，相对加速度为零。沿坡和法向分别有
\[
f=m(g\sin\gamma+A\cos\gamma),\qquad
N=m(g\cos\gamma-A\sin\gamma).
\]
静摩擦条件为 $|f|\le\mu N$，还须有 $N>0$。将绝对值拆成两个不等式，
\[
\begin{aligned}
A(\cos\gamma+\mu\sin\gamma)&\le g(\mu\cos\gamma-\sin\gamma),\\
A(\cos\gamma-\mu\sin\gamma)&\ge-g(\mu\cos\gamma+\sin\gamma).
\end{aligned}
\]
由 $\mu\tan\gamma<1$，两式中 $A$ 的系数均为正，因此
\[
\boxed{-g\frac{\mu\cos\gamma+\sin\gamma}{\cos\gamma-\mu\sin\gamma}
\le A\le
g\frac{\mu\cos\gamma-\sin\gamma}{\cos\gamma+\mu\sin\gamma}.}
\]
\begin{figure}[H]\centering
\begin{tikzpicture}
\draw[axis](-.5,0)--(7,0)node[right]{$A$};\draw[curve](.5,0)--(5.5,0);\foreach\x/\s in{.5/A_-,3/{-g\tan\gamma},5.5/A_+}{\fill(\x,0)circle(1.6pt)node[below]{$\s$};}
\node at(1.7,.65){$f<0$，沿坡向下};\node at(4.5,.65){$f>0$，沿坡向上};\node at(3,-1){可维持相对静止的区间};
\end{tikzpicture}
\caption{区间图只表示三个加速度的次序，横向未按数值比例绘制。}\label{rot:fig:d6-h1-s}
\end{figure}
记上述上下端点为 $A_-,A_+$。由于
\[
g\cot\gamma-A_+
=\frac{g}{\sin\gamma(\cos\gamma+\mu\sin\gamma)}>0,
\]
整个区间都满足 $N>0$。图\ref{rot:fig:d6-h1-s}中，$A=-g\tan\gamma$ 时摩擦为零；$A$ 比它小时摩擦沿坡向下，比它大时摩擦沿坡向上。两端点分别达到相反方向的静摩擦极限。超出区间便无法维持相对静止；若加速度进一步使接触失效，还须改用离面运动模型。
\end{solution}


\begin{exercise}[水槽右端露底后的液面]\label{rot:problem-exercise-2}\label{rot:ex:d6-h2}
沿用例\ref{rot:ex:d6-e2}的矩形水槽，给定 $h_0<H/2$，向右的加速度 $A$ 已使右端露底，但液体尚未溢出。求槽底湿润长度、左壁液深、自由液面和底部表压，确定这一阶段的 $A$ 范围，并求端壁对液体的水平合力。
\end{exercise}
\begin{figure}[H]\centering
\begin{tikzpicture}
\fill[gray!12](0,0)--(4,0)--(0,2.5)--cycle;\draw[curve](0,3)--(0,0)--(6,0)--(6,3);\draw[curve](0,2.5)--(4,0);\node[below]at(0,0){$-\ell$};\node[below]at(4,0){$x_d$};\node[below]at(6,0){$\ell$};
\draw[<->](0,-.8)--(4,-.8)node[midway,below]{$d$};\draw[<->](-.6,0)--(-.6,2.5)node[midway,left]{$h_L$};\draw[axis](2.5,3.4)--(4.2,3.4)node[right]{$A$};
\end{tikzpicture}
\caption{液体横截面成为三角形，右端干燥部分不计入体积积分。}\label{rot:fig:d6-h2-q}
\end{figure}
\begin{solution}
图\ref{rot:fig:d6-h2-q}中，自由液面斜率仍为 $-A/g$。令湿润长度为 $d$，干湿交点为 $x_d=-\ell+d$，则
\[
h_L=\frac Agd,\qquad z_s(x)=\frac Ag(x_d-x),\quad -\ell\le x\le x_d.
\]
每单位宽度的液体面积由体积守恒保持 $2\ell h_0$。在湿区积分，令 $\xi=x+\ell$，则 $x_d-x=d-\xi$，
\[
\int_{-\ell}^{x_d}z_s(x)\dd x
=\frac Ag\int_0^d(d-\xi)\dd\xi
=\frac Ag[d\xi-\xi^2/2]_0^d=\frac{Ad^2}{2g}.
\]
故
\[
\frac12d h_L=2\ell h_0,
\qquad d=2\sqrt{\frac{g\ell h_0}{A}},\qquad
h_L=2\sqrt{\frac{A\ell h_0}{g}}.
\]
体积只在 $-\ell\le x\le x_d$ 上积分，$x>x_d$ 的干底区域不计入液体截面积。
\begin{figure}[H]\centering
\begin{tikzpicture}
\draw[axis](0,0)--(6.3,0)node[right]{$x$};\draw[axis](0,0)--(0,2.9)node[above]{$p_b$};\draw[curve](.4,2.4)--(4,0)--(5.7,0);
\node[below]at(.4,0){$-\ell$};\node[below]at(4,0){$x_d$};\node at(4.9,.6){干底};\draw[guide](.4,0)--(.4,2.4);\node[right]at(1,2.6){$p_b(-\ell)=\rho gh_L$};
\end{tikzpicture}
\caption{湿底表压线性降至零；干底只承受大气压，表压为零。}\label{rot:fig:d6-h2-s}
\end{figure}
图\ref{rot:fig:d6-h2-s}中的分段压强为
\[
p_b(x)=\begin{cases}\rho A(x_d-x),&-\ell\le x\le x_d,\\0,&x_d<x\le\ell.\end{cases}
\]
露底要求 $d<2\ell$，不溢流要求 $h_L\le H$，即
\[
\frac{gh_0}{\ell}<A\le\frac{gH^2}{4\ell h_0}.
\]
因 $H>2h_0$，这个区间非空。右壁未接触液体，水平压力为零；左壁作用于液体的力为
\[
F_L=\frac12\rho gBh_L^2
=2\rho B\ell h_0A=M_{\rm liq}A.
\]
上式的压力积分为 $F_L=\rho gB\int_0^{h_L}(h_L-z)\dd z=\rho gB[h_Lz-z^2/2]_0^{h_L}$。底部竖直合力使用同一湿区边界，
\[
F_b=\rho AB\int_0^d(d-\xi)\dd\xi
=\rho AB[d\xi-\xi^2/2]_0^d
=\frac{\rho ABd^2}2=2\rho gB\ell h_0=M_{\rm liq}g,
\]
恰好平衡液体重力。
\end{solution}


\begin{exercise}[加速度阶跃下的弹簧压缩]\label{rot:problem-exercise-3}\label{rot:ex:d6-h3}
小车上的水平轻弹簧可受拉也可受压，劲度系数为 $k$，连接质量为 $m$ 的光滑物块。初始小车、物块均静止，弹簧无形变。$t=0$ 起，小车以恒加速度 $A$ 向右运动。取弹簧形变量 $x$ 向右伸长为正，求 $x(t)$、最大压缩量及物块的地面加速度，并写出相对能量守恒式。
\end{exercise}
\begin{figure}[H]\centering
\begin{tikzpicture}
\draw[curve](0,-.5)rectangle(6,1);\draw[curve](.5,-.4)--(.5,.8);\draw[spring](.5,.2)--(3.4,.2);\draw[fill=white](3.4,-.2)rectangle(4.2,.6);\node at(3.8,.2){$m$};\draw[axis](4.8,1.5)--(6.3,1.5)node[right]{$A$};\node at(1.7,1.4){$x(0)=0$};
\end{tikzpicture}
\caption{小车向右加速时，初始惯性力使物块相对车向左压缩弹簧。}\label{rot:fig:d6-h3-q}
\end{figure}
\begin{solution}
图\ref{rot:fig:d6-h3-q}的小车系方程为
\[
\ddot x+\omega_0^2x=-A,\qquad \omega_0=\sqrt{k/m}.
\]
静态平衡形变量为 $x_*=-A/\omega_0^2$。置 $\xi=x-x_*$，得 $\ddot\xi+\omega_0^2\xi=0$。其特征根为 $\pm i\omega_0$，实通解及其导数为
\[
\xi=C\cos\omega_0t+D\sin\omega_0t,\qquad
\dot\xi=-C\omega_0\sin\omega_0t+D\omega_0\cos\omega_0t.
\]
由 $\xi(0)=A/\omega_0^2$ 得 $C=A/\omega_0^2$，由 $\dot\xi(0)=0$ 得 $D=0$。还原 $x=x_*+\xi$，得到
\[
x(t)=-\frac A{\omega_0^2}(1-\cos\omega_0t),\qquad
\dot x=-\frac A{\omega_0}\sin\omega_0t.
\]
恒加速度在 $t=0$ 阶跃，并不等于给物块施加冲量：位置、速度都连续，只有方程中的加速度项改变。新平衡位置在原位置左侧，初态到新平衡位置的距离为 $mA/k$，图\ref{rot:fig:d6-h3-s}中的轨道便以此为振幅绕新平衡振动。
首次最大压缩发生在 $t=\pi/\omega_0$，其大小为
$|x|_{\max}=2A/\omega_0^2=2mA/k$，为静态压缩量的两倍。
\begin{figure}[H]\centering
\begin{tikzpicture}
\draw[axis](0,0)--(7,0)node[right]{$\omega_0t$};\draw[axis](0,-2.7)--(0,.6)node[above]{$kx/(mA)$};
\draw[curve,domain=0:6.283,samples=100]plot(\x,{cos(\x*180/pi)-1});\draw[guide](0,-1)--(6.5,-1);\node[left]at(0,-1){$-1$};\node[left]at(0,-2){$-2$};\node[above]at(3.142,0){$\pi$};\node[above]at(6.283,0){$2\pi$};
\end{tikzpicture}
\caption{弹簧在自然长度与最大压缩状态之间往返，虚线为新的相对平衡位置。}\label{rot:fig:d6-h3-s}
\end{figure}
由 $\dot x$ 再求导，$\ddot x=-(A/\omega_0)\omega_0\cos\omega_0t=-A\cos\omega_0t$，所以地面加速度为 $a_q=A+\ddot x=A(1-\cos\omega_0t)$，范围为 $0\le a_q\le2A$。图\ref{rot:fig:d6-h3-s}的最大压缩处，弹簧对物块向右的力为 $2mA$，与 $a_q=2A$ 相符。
将相对运动方程乘以 $m\dot x$，有
\[
m\dot x\ddot x+kx\dot x+mA\dot x
=\frac{\dd}{\dd t}\left(\frac12m\dot x^2+\frac12kx^2+mAx\right)=0.
\]
各项分别用了 $\dd(\dot x^2)/\dd t=2\dot x\ddot x$、$\dd(x^2)/\dd t=2x\dot x$；$A$ 为常量，故第三项是惯性力势能 $mAx$ 的导数。对时间从 $0$ 积到 $t$，初态 $x=\dot x=0$，得到
\[
\frac12m\dot x^2+\frac12kx^2+mAx=0.
\]
配方后可写为
\[
\frac12m\dot x^2+\frac12k(x+mA/k)^2=\frac{m^2A^2}{2k},
\]
相对能量中须计入惯性力势能$mAx$；相对动能与弹性势能两项之和并不守恒。
\end{solution}


\begin{exercise}[径向滑珠的回转与到轴（选做）]\label{rot:problem-exercise-4}\label{rot:ex:d6-h4}
水平光滑径向杆以恒角速度 $\Omega>0$ 转动，杆从 $O$ 沿正半径延伸。滑珠初始 $r=a>0$，相对杆向内的速度大小为 $u$。求运动，并确定先转向外侧、渐近接近 $O$、在有限时间到达 $O$ 这三种情况各自的条件。杆向外足够长，只讨论到达 $O$ 之前的运动。
\end{exercise}
\begin{figure}[H]\centering
\begin{tikzpicture}
\draw[curve](0,0)--(6,0);\fill(0,0)circle(1.6pt)node[below]{$O$};\fill(2.5,0)circle(2pt)node[above]{$m$};\draw[axis](2.5,0)--(1.3,0)node[midway,below]{$u$};\node[below]at(3.1,0){$a$};\draw[axis](.7,.7)arc(45:135:1);\node at(.1,1.6){$\Omega$};
\end{tikzpicture}
\caption{向内初速度与向外的离心项相竞争，初始径向坐标为a。}\label{rot:fig:d6-h4-q}
\end{figure}
\begin{solution}
沿图\ref{rot:fig:d6-h4-q}的杆方向仍有 $\ddot r=\Omega^2r$。特征方程 $s^2-\Omega^2=0$ 的两个根为 $\pm\Omega$，所以
\[
r=C_1e^{\Omega t}+C_2e^{-\Omega t},\qquad
\dot r=\Omega C_1e^{\Omega t}-\Omega C_2e^{-\Omega t}.
\]
初始条件给出 $C_1+C_2=a$、$\Omega(C_1-C_2)=-u$，相加、相减得
\[
C_1=\frac12\left(a-\frac{u}{\Omega}\right),\qquad
C_2=\frac12\left(a+\frac{u}{\Omega}\right).
\]
用双曲函数合并两个指数项，运动可写为
\[
r=a\cosh\Omega t-\frac u\Omega\sinh\Omega t,
\quad \dot r=a\Omega\sinh\Omega t-u\cosh\Omega t.
\]
指数解中 $C_2$ 始终为正，$C_1$ 在 $u=\Omega a$ 处变号：为正时 $r$ 最终增长，为零时只剩衰减项，为负时 $r$ 会在有限正时间降至零。
相应的旋转系能量为
\[
E'=\frac12m(\dot r^2-\Omega^2r^2)
=\frac12m(u^2-\Omega^2a^2).
\]
若 $0<u<\Omega a$，令 $\dot r=0$ 得
\[
t_* =\frac1\Omega\operatorname{artanh}\frac{u}{\Omega a},\qquad
r_{\min}=\sqrt{a^2-u^2/\Omega^2}>0.
\]
由能量式改写得
\[
\dot r^2=\Omega^2r^2+u^2-\Omega^2a^2.
\]
径向转向必须先满足 $\dot r=0$，因此 $r_{\min}^2=a^2-u^2/\Omega^2$ 只有在 $u<\Omega a$ 时才为正。这同时排除了另外两类轨道中出现正半径转折点的可能。
珠在到轴前径向速度已降至零，随后转而向外加速。
\begin{figure}[H]\centering
\begin{tikzpicture}
\draw[axis](0,0)--(4.6,0)node[right]{$\Omega t$};\draw[axis](0,0)--(0,3.5)node[above]{$r/a$};
\draw[curve,domain=0:2,samples=90]plot({1.8*\x},{1.5*(cosh(\x)-.5*sinh(\x))});
\draw[guide,domain=0:2.2,samples=90]plot({1.8*\x},{1.5*exp(-\x)});
\draw[curve,densely dotted,domain=0:.804719,samples=70]plot({1.8*\x},{1.5*(cosh(\x)-1.5*sinh(\x))});
\node[left]at(0,1.5){$1$};\node[above]at(2.8,3.2){$u=\Omega a/2$};\node at(2.8,.65){$u=\Omega a$};\node at(2,-.65){点线：$u=3\Omega a/2$};
\begin{scope}[shift={(7.4,1.5)},scale=1.1]
\draw[axis](-.15,0)--(2.55,0)node[right]{$r/a$};\draw[axis](0,-1.75)--(0,1.75)node[above]{$\dot r/(\Omega a)$};
\draw[curve,domain=-.549306:1.25,samples=110]plot({sqrt(.75)*cosh(\x)},{sqrt(.75)*sinh(\x)});
\draw[axis,domain=.3:.5,samples=12]plot({sqrt(.75)*cosh(\x)},{sqrt(.75)*sinh(\x)});
\draw[guide](1,-1)--(0,0);\draw[axis,dashed](.65,-.65)--(.45,-.45);
\draw[curve,densely dotted,domain=1:0,samples=80]plot(\x,{-sqrt((\x)^2+1.25)});
\draw[axis,densely dotted](.65,{-sqrt(.65*.65+1.25)})--(.45,{-sqrt(.45*.45+1.25)});
\fill(1,-.5)circle(1.5pt);\fill(1,-1)circle(1.5pt);\fill(1,-1.5)circle(1.5pt);
\fill({sqrt(.75)},0)circle(1.7pt);\draw[guide]({sqrt(.75)},0)--(1.45,.5);\node[right]at(1.45,.5){$r_{\min}/a$};
\node[below left]at(0,0){$O$};\node[below]at(1.3,-1.85){相轨迹，线型同左图};
\end{scope}
\end{tikzpicture}
\caption{左图为三类径向运动，右图为对应相轨迹，箭头沿时间增加方向。临界分支渐近接近原点；实线在正半径处转向，点线以非零速度到轴。}\label{rot:fig:d6-h4-s}
\end{figure}
当 $u=\Omega a$，图\ref{rot:fig:d6-h4-s}的临界解为 $r=ae^{-\Omega t}$，仅在 $t\to\infty$ 时趋零。
若 $u>\Omega a$，解 $r=0$ 得
\[
t_O=\frac1\Omega\operatorname{artanh}\frac{\Omega a}{u},
\qquad \dot r(t_O)=-\sqrt{u^2-\Omega^2a^2}.
\]
到轴时仍有非零向内速度。由于杆只沿正半径延伸，后续运动取决于端部构造，本题的运动阶段在此结束。
若不使用反双曲函数，两种有限时间分别为
\[
\begin{aligned}
t_*&=\frac1{2\Omega}\ln\frac{\Omega a+u}{\Omega a-u},&&0<u<\Omega a,\\
t_O&=\frac1{2\Omega}\ln\frac{u+\Omega a}{u-\Omega a},&&u>\Omega a.
\end{aligned}
\]
两式在各自定义域内分子大于分母，时间均为正。$u$ 从两侧趋于 $\Omega a$ 时，它们都对数发散，与图\ref{rot:fig:d6-h4-s}中临界轨道只在无限时间到轴相符。
\end{solution}


\begin{exercise}[地面静止点的转动系运动]\label{rot:problem-exercise-5}\label{rot:ex:d6-h5}
无摩擦水平面上的一质点在地面系始终静止，与固定点 $O$ 相距 $a$。采用绕 $O$ 以角速度 $\Omega>0$ 逆时针转动的坐标系，初始坐标为 $(a,0)$。求相对轨迹、速度、加速度及各惯性力，说明这些惯性力怎样满足相对运动方程。
\end{exercise}
\begin{figure}[H]\centering
\begin{tikzpicture}
\draw[axis](-.4,0)--(3.8,0)node[right]{$x=x'$};\draw[axis](0,-1.8)--(0,2.4)node[above]{$y=y'$};\fill(2.5,0)circle(2pt)node[above]{$m$};\draw[axis](2.5,0)--(2.5,-1.3)node[right]{$v'_0=\Omega a$};\node at(1.2,1.5){地面速度为零};\node[below left]at(0,0){$O$};
\end{tikzpicture}
\caption{初始两坐标系方向重合，相对转动系的速度沿负 $y$ 方向。}\label{rot:fig:d6-h5-q}
\end{figure}
\begin{solution}
图\ref{rot:fig:d6-h5-q}中的地面位置为 $(a,0)$，乘逆转矩阵即得
\[
x'=a\cos\Omega t,\qquad y'=-a\sin\Omega t.
\]
对这两个坐标分别求一次、两次时间导数，得到
\[
\begin{aligned}
\dot x'&=-a\Omega\sin\Omega t,&\dot y'&=-a\Omega\cos\Omega t,\\
\ddot x'&=-a\Omega^2\cos\Omega t=-\Omega^2x',&
\ddot y'&=a\Omega^2\sin\Omega t=-\Omega^2y'.
\end{aligned}
\]
质点在转动系中沿半径为 $a$ 的圆顺时针运动，以上分量可合为
\[
\vct v'=-\vct\Omega\times\vct r',\qquad
\vct a'=-\Omega^2\vct r'.
\]
在初始位置 $(a,0)$，相对速度为 $(0,-\Omega a)$，图\ref{rot:fig:d6-h5-s}中的科里奥利力因此指向O。一般位置可用三重积恒等式 $\vct\Omega\times(\vct\Omega\times\vct r')=-\Omega^2\vct r'$，其中用到了水平位置与竖直转轴垂直。
角速度恒定、原点固定，平移与欧拉惯性力均为零。离心力为 $m\Omega^2\vct r'$；代入相对速度计算科里奥利力，
\[
\vct F_{\rm cor}=-2m\vct\Omega\times(-\vct\Omega\times\vct r')
=-2m\Omega^2\vct r'.
\]
\begin{figure}[H]\centering
\begin{tikzpicture}
\draw[guide](0,0)circle(1.3);\fill(0,0)circle(1.5pt)node[above left]{$O$};\fill(1.3,0)circle(1.8pt)node[above]{$m$};
\draw[axis](1.3,0)--(2.3,0)node[right]{$m\Omega^2a$};\draw[axis](1.3,0)--(-.7,0);\node at(.2,-.45){$2m\Omega^2a$};\draw[axis](1.3,0)--(1.3,-1)node[right]{$\vct v'$};
\node[align=left]at(5.7,.6){离心力向外\\科里奥利力向内\\合惯性力向内};
\end{tikzpicture}
\caption{初始瞬间的惯性力图，科里奥利力大小是离心力的两倍。}\label{rot:fig:d6-h5-s}
\end{figure}
图\ref{rot:fig:d6-h5-s}中，两个惯性力之和为 $-m\Omega^2\vct r'=m\vct a'$，给出了相对圆周运动所需的向心加速度。地面系中水平真实力与加速度均为零，两种描述通过坐标变换相连。
相对动能为 $m\Omega^2a^2/2$，离心势为其相反数，故 $E'=0$。地面系中 $E=L_z=0$，也符合 $E'=E-\Omega L_z$。
\end{solution}


\begin{exercise}[转盘制动时的摩擦需求]\label{rot:problem-exercise-6}\label{rot:ex:d6-h6}
水平转盘初始以角速度 $\Omega_0>0$ 转动，随后按 $\Omega(t)=\Omega_0-\alpha t$ 匀减速，至 $t_b=\Omega_0/\alpha$ 停止，此后保持静止，$\alpha>0$。距轴 $r$ 处的物体相对盘不动。求全过程所需的最小静摩擦系数，以及制动段和停止后的摩擦。
\end{exercise}
\begin{figure}[H]\centering
\begin{tikzpicture}
\draw[curve](0,0)circle(2);\fill(0,0)circle(1.5pt)node[below left]{$O$};\fill(1.3,0)circle(2pt)node[right]{$m$};\draw[guide](0,0)--(1.3,0)node[midway,below]{$r$};
\draw[axis](-1.2,1.2)arc(135:225:1.7);\node at(-2.5,.7){$\Omega>0$};\node at(3.3,1.5){$\dot\Omega=-\alpha$};
\end{tikzpicture}
\caption{角速度仍为正，角加速度与它反向。}\label{rot:fig:d6-h6-q}
\end{figure}
\begin{solution}
在图\ref{rot:fig:d6-h6-q}的制动段，物体地面加速度同时包括向心分量和反向切向分量，故
\[
\vct f=-mr\Omega^2\vct e_r-mr\alpha\vct e_\varphi,
\qquad |\vct f|=mr\sqrt{\Omega^4+\alpha^2}.
\]
法向力为 $mg$，摩擦系数须满足
\[
\mu\ge\frac rg\sqrt{(\Omega_0-\alpha t)^4+\alpha^2}.
\]
在 $0\le t\le t_b$，$\Omega$ 非负且单调减小，右端最大值在初始时刻。因此
\[
\boxed{\mu_{\min}=\frac rg\sqrt{\Omega_0^4+\alpha^2}.}
\]
\begin{figure}[H]\centering
\begin{tikzpicture}
\draw[axis](0,0)--(5,0)node[right]{$\Omega/\Omega_0$};\draw[axis](0,0)--(0,3)node[above]{$|f|/(mr\Omega_0^2)$};
\draw[curve,domain=0:1,samples=80]plot({4*\x},{2*sqrt(\x^4+.25)});\node[below]at(4,0){$1$};\node[left]at(0,1){$1/2$};\draw[axis](3.4,2.8)--(1.5,2.8)node[midway,above]{时间推进};
\end{tikzpicture}
\caption{示意 $\alpha=\Omega_0^2/2$ 的制动过程，摩擦需求从右向左减小。}\label{rot:fig:d6-h6-s}
\end{figure}
图\ref{rot:fig:d6-h6-s}中，摩擦需求随制动减小，故初始条件满足后，后续制动段也能保持相对静止。

刚停止前，径向分量趋于零，切向摩擦仍为 $-mr\alpha\vct e_\varphi$。$t_b$ 时角加速度改为零，物体已在地面系静止，所需水平摩擦随之变为零。这一有限力的跳变不会产生冲量，物体速度连续。
\end{solution}


\begin{exercise}[临界转速下的非线性振动（选做）]\label{rot:problem-exercise-7}\label{rot:ex:d6-h7}
半径为 $a$ 的光滑圆环绕竖直直径以恒定角速度 $\Omega=\sqrt{g/a}$ 转动。质量为 $m$ 的小珠从 $\theta=\Theta$ 处相对静止释放，以最低点为 $\theta=0$，其中 $0<\Theta\ll1$，角度用弧度表示。保留势能展开中的首个非零项，求周期随振幅的渐近关系，判断最低点的稳定性。
\end{exercise}
\begin{figure}[H]\centering
\begin{tikzpicture}
\draw[curve](0,0)circle(2);\draw[guide](0,-2.05)--(0,2.5)node[above]{转轴};\fill(0,0)circle(1.4pt)node[left]{$O$};
\coordinate(A)at({2*sin(35)},{-2*cos(35)});\coordinate(B)at({-2*sin(35)},{-2*cos(35)});
\draw[guide](B)--(0,0)--(A);\fill(A)circle(2pt)node[below right]{$\theta=\Theta$};\fill(B)circle(1.5pt)node[below left]{$-\Theta$};\draw[axis](0,-.8)arc(-90:-55:.8)node[midway,below right]{$\Theta$};
\node[below]at(0,-2.15){$\theta=0$};\node at(3.5,1){$\Omega^2a=g$};\node at(3.5,.4){$\Theta\ll1$};
\end{tikzpicture}
\caption{两个小振幅转折位置；图中角度为便于辨认而放大。}\label{rot:fig:d6-h7-q}
\end{figure}
\begin{solution}
图\ref{rot:fig:d6-h7-q}中的圆环对小珠只作法向约束，相对运动沿圆周进行。由例\ref{rot:ex:d6-e7}的有效势能，在临界转速下有
\[
U_{\rm eff}=-mga\cos\theta-\frac{mga}{2}\sin^2\theta.
\]
利用 $\cos\theta=1-\theta^2/2+\theta^4/24+O(\theta^6)$ 和 $\sin^2\theta=\theta^2-\theta^4/3+O(\theta^6)$，得
\[
U_{\rm eff}-U_{\rm eff}(0)=\frac{mga}{8}\theta^4+O(\theta^6).
\]
二次项抵消，四次项为正。因此最低点仍是严格局部极小点，只是回复力不再与小位移成正比。振幅为 $\Theta$ 时，相对能量守恒给出
\[
\frac12ma^2\dot\theta^2=\frac{mga}{8}(\Theta^4-\theta^4),
\qquad
|\dot\theta|=\frac12\sqrt{\frac ga}\sqrt{\Theta^4-\theta^4}.
\]
\begin{figure}[H]\centering
\begin{tikzpicture}
\draw[axis](-2.6,0)--(2.7,0)node[right]{$\theta/\Theta$};\draw[axis](0,0)--(0,2.6)node[above]{$(U_{\rm eff}-U_0)/\Delta U$};
\draw[curve,domain=-1.15:1.15,samples=90]plot({2*\x},{1.2*(\x)^4});\draw[guide](-2.3,1.2)--(2.3,1.2);\node[above right]at(2.4,1.1){$\dfrac{E'-U_0}{\Delta U}=1$};\fill(-2,1.2)circle(1.7pt);\fill(2,1.2)circle(1.7pt);\draw[guide](-2,0)--(-2,1.2);\draw[guide](2,0)--(2,1.2);\node[below]at(-2,0){$-1$};\node[below]at(2,0){$1$};\node[left]at(0,1.2){$1$};
\end{tikzpicture}
\caption{四次势阱及两个转折点。$U_0=U_{\rm eff}(0)$，$\Delta U=mga\Theta^4/8$。}\label{rot:fig:d6-h7-s}
\end{figure}
图\ref{rot:fig:d6-h7-s}中，从 $0$ 到 $\Theta$ 的时间为四分之一周期。令 $\theta=\Theta u$，有
\[
T=4\int_0^\Theta\frac{\dd\theta}{|\dot\theta|}
=\frac8\Theta\sqrt{\frac ag}\int_0^1\frac{\dd u}{\sqrt{1-u^4}}.
\]
这个定积分可化成端点没有奇性的形式。取 $u=\sin\phi$，则 $\dd u=\cos\phi\,\dd\phi$，上下限 $u=0,1$ 分别对应 $\phi=0,\pi/2$。在此区间 $\cos\phi\ge0$，且
\[
1-u^4=1-\sin^4\phi=\cos^2\phi(1+\sin^2\phi),
\]
于是分子中的 $\cos\phi$ 与分母中的同一因子约去，得
\[
\int_0^1\frac{\dd u}{\sqrt{1-u^4}}
=\int_0^{\pi/2}\frac{\dd\phi}{\sqrt{1+\sin^2\phi}}.
\]
右侧被积函数在整个区间连续，因而图\ref{rot:fig:d6-h7-s}中的转折点虽速度为零，质点经过转折点附近所需时间仍有限。
该积分没有初等原函数。记 $f(\phi)=(1+\sin^2\phi)^{-1/2}$，将 $[0,\pi/2]$ 等分为偶数 $n$ 段，步长 $h=\pi/(2n)$，复合辛普森公式给出
\[
S_n=\frac h3\left[f(0)+f(\pi/2)
+4\sum_{j=1,3,\ldots,n-1}f(jh)
+2\sum_{j=2,4,\ldots,n-2}f(jh)\right].
\]
取 $n=16,32$，两次均得 $1.311028777$（保留九位小数）。将其乘以周期表达式中的系数 $8$，得
\[
\boxed{T=\frac{10.4882}{\Theta}\sqrt{\frac ag}\,[1+O(\Theta^2)].}
\]
振幅越小，周期越长，$\Theta\to0$ 时振动频率趋于零。线性化只能得到零频率，而四次势阱仍将小能量轨道限制在最低点附近，因此最低点保持稳定。
\end{solution}


\begin{exercise}[竖直上抛后的落点]\label{rot:problem-exercise-8}\label{rot:ex:d6-h8}
在北纬 $\lambda$ 处，以速率 $u$ 相对地面竖直上抛一质点，初始水平速度为零。忽略空气阻力，取当地有效重力加速度为 $g$。仅保留地球角速度 $\Omega$ 的一阶效应，求回到原高度时的东西偏移和东西向速度。设 $0\le\lambda<\pi/2$，飞行高度、水平位移均远小于地球半径。
\end{exercise}
\begin{figure}[H]\centering
\begin{tikzpicture}
\draw[curve](-2,0)--(3,0);\draw[axis](0,0)--(3.3,0)node[right]{$x$（东）};\draw[axis](0,0)--(0,3.3)node[above]{$z$（天顶）};\fill(0,0)circle(2pt)node[below left]{$O$};\draw[axis](.2,.2)--(.2,1.8)node[right]{$u$};\draw[axis](2.1,2.6)--(2.1,1.6)node[right]{$g$};\node[anchor=east]at(-.4,2.5){北向垂直纸面向内};
\end{tikzpicture}
\caption{以发射点为原点；初始速度只有竖直分量。}\label{rot:fig:d6-h8-q}
\end{figure}
\begin{solution}
按图\ref{rot:fig:d6-h8-q}取东为 $x$ 正向、北为 $y$ 正向、天顶为 $z$ 正向。零阶运动为
\[
v_z^{(0)}=u-gt,\qquad z^{(0)}=ut-\frac12gt^2,
\qquad T=\frac{2u}{g}.
\]
由 $\vct\Omega=(0,\Omega\cos\lambda,\Omega\sin\lambda)$，一阶科里奥利加速度的东向分量为
\[
\ddot x=-2\Omega\cos\lambda\,(u-gt).
\]
初始 $x=\dot x=0$。先积加速度，再积速度，两个积分的下限都取发射时刻：
\[
\begin{aligned}
\dot x(t)&=-2\Omega\cos\lambda\int_0^t(u-g\tau)\,\dd\tau
=-2\Omega\cos\lambda\left[u\tau-\frac12g\tau^2\right]_0^t\\
&=-2\Omega\cos\lambda\left(ut-\frac12gt^2\right),\\
x(t)&=-2\Omega\cos\lambda\int_0^t\left(u\tau-\frac12g\tau^2\right)\dd\tau\\
&=-2\Omega\cos\lambda\left[\frac12u\tau^2-\frac16g\tau^3\right]_0^t
=-\Omega ut^2\cos\lambda+\frac13\Omega gt^3\cos\lambda.
\end{aligned}
\]
\begin{figure}[H]\centering
\begin{tikzpicture}
\draw[axis](0,0)--(5,0)node[right]{$\tau=gt/u$};\draw[axis](0,-2.5)--(0,.7)node[above]{$\dot x/(\Omega u^2\cos\lambda/g)$};
\draw[curve,domain=0:2,samples=80]plot({2*\x},{2*(\x^2-2*\x)});\draw[guide](2,0)--(2,-2);\node[above]at(2,0){$1$};\node[above]at(4,0){$2$};\node[left]at(0,-2){$-1$};\node[below]at(2,-2.05){最高点};
\end{tikzpicture}
\caption{全程东向速度不大于零：下降时虽向东加速，质点仍向西漂移。}\label{rot:fig:d6-h8-s}
\end{figure}
积分结果还可整理为 $\dot x=-2\Omega\cos\lambda\,z^{(0)}(t)$。在从抛出到回到原高度的整个开区间，$z^{(0)}>0$，所以东西速度始终向西。上升、下降两段虽然加速度方向相反，向西位移却在两段中持续积累。
图\ref{rot:fig:d6-h8-s}中，最高点 $t=u/g$ 对应东西向加速度为零、向西速度最大，速度为 $-\Omega u^2\cos\lambda/g$。下降时加速度改向东，向西速度逐渐减小，直到落回原高度时
\[
\boxed{x(T)=-\frac{4\Omega u^3\cos\lambda}{3g^2},\qquad \dot x(T)=0.}
\]
落点在发射点西侧。速度曲线全在时间轴下方，曲线与时间轴之间的有向面积为负，给出有限的向西位移。

竖直方向的科里奥利修正含有一阶水平速度，因而属于二阶，求一阶偏移时可直接使用 $T=2u/g$。
\end{solution}


\begin{exercise}[向北平抛时的东西偏移]\label{rot:problem-exercise-9}\label{rot:ex:d6-h9}
在纬度为 $\lambda$ 的地点，从高为 $H$ 处以速率 $u$ 相对地面向正北平抛，初始竖直、东西速度均为零。忽略空气阻力，取当地有效重力加速度为 $g$，只计 $\Omega$ 的一阶效应。求落地时的东西偏移及南半球偏移为零的条件。设 $-\pi/2<\lambda<\pi/2$，$\Omega\sqrt{2H/g}\ll1$，运动范围足够小，可使用固定的当地坐标轴。
\end{exercise}
\begin{figure}[H]\centering
\begin{tikzpicture}
\begin{scope}
\draw[axis](0,0)--(2.7,0)node[right]{东};\draw[axis](0,0)--(0,2.8)node[above]{北};\fill(0,0)circle(2pt)node[below left]{$P$};\draw[axis](0,.2)--(0,1.7)node[right]{$u$};\node at(1.1,-.7){俯视：发射方向};
\end{scope}
\begin{scope}[xshift=5.2cm]
\draw[curve](-.6,0)--(2,0);\draw[curve](0,0)--(0,2);\fill(0,2)circle(2pt)node[above]{$P$};\draw[<->](-.45,0)--(-.45,2)node[midway,left]{$H$};\draw[axis](.2,2)--(1.6,2)node[right]{$u$（北）};\node at(.8,-.7){侧视：初始高度};
\end{scope}
\end{tikzpicture}
\caption{俯视与侧视共用发射点P；东西偏移由北向运动和下落共同产生。}\label{rot:fig:d6-h9-q}
\end{figure}
\begin{solution}
取东、北、天顶为 $x$、$y$、$z$ 正向。图\ref{rot:fig:d6-h9-q}给出的零阶速度为 $(0,u,-gt)$。代入 $-2\vct\Omega\times\vct v^{(0)}$，得
\[
a_x=2\Omega u\sin\lambda+2\Omega gt\cos\lambda,
\qquad a_y=a_z=0\quad\text{（一阶）}.
\]
第一项由向北平抛引起，其符号随半球改变；第二项由下落引起，在本题纬度范围内指东。用 $v_x(0)=x(0)=0$ 依次积分：
\[
\begin{aligned}
v_x(t)&=\int_0^t(2\Omega u\sin\lambda+2\Omega g\tau\cos\lambda)\,\dd\tau\\
&=\left[2\Omega u\tau\sin\lambda+\Omega g\tau^2\cos\lambda\right]_0^t
=2\Omega ut\sin\lambda+\Omega gt^2\cos\lambda,\\
x(t)&=\int_0^t(2\Omega u\tau\sin\lambda+\Omega g\tau^2\cos\lambda)\,\dd\tau\\
&=\left[\Omega u\tau^2\sin\lambda+\frac13\Omega g\tau^3\cos\lambda\right]_0^t
=\Omega ut^2\sin\lambda+\frac13\Omega gt^3\cos\lambda.
\end{aligned}
\]
落地时间 $T=\sqrt{2H/g}$。故所求东向偏移为
\[
\boxed{\delta=\frac{2\Omega uH}{g}\sin\lambda
+\frac{2\Omega H}{3}\sqrt{\frac{2H}{g}}\cos\lambda.}
\]
\begin{figure}[H]\centering
\begin{tikzpicture}
\draw[axis](0,0)--(5.2,0)node[right]{$u/u_*$};\draw[axis](0,-1.7)--(0,2.2)node[above]{$\delta/\delta_0$};\draw[curve](0,1.5)--(4.8,-1.5);\draw[guide](2.4,0)--(2.4,-1.4);\node[below left]at(2.4,0){$1$};\node[left]at(0,1.5){$1$};\node at(1.2,.35){向东};\node at(3.6,-1.3){向西};\node at(3.7,1.5){$\lambda<0$};
\end{tikzpicture}
\caption{南半球的落点偏移。$\delta_0=\Omega gT^3\cos\lambda/3>0$ 是同高度自由落体的东偏。}\label{rot:fig:d6-h9-s}
\end{figure}
在北半球，两项均不小于零；赤道上第一项为零，仍有下落造成的东偏。在南半球 $\sin\lambda<0$，令 $\delta=0$，可得
\[
\boxed{u_*= -\frac{gT}{3}\cot\lambda>0.}
\]
图\ref{rot:fig:d6-h9-s}表示 $\delta=\delta_0(1-u/u_*)$：$u<u_*$ 时落点向东，$u>u_*$ 时向西。$u=u_*$ 只使落地点的东西位移抵消，飞行过程中一般仍有东西运动。这个临界速率还须满足当地小区域近似；若由它算出的射程已不可忽略地球曲率，就须重建模型。
在临界速率 $u_*$ 下，落地的东西速度仍为
\[
v_x(T)=2\Omega u_*T\sin\lambda+\Omega gT^2\cos\lambda
=\frac13\Omega gT^2\cos\lambda>0.
\]
图\ref{rot:fig:d6-h9-s}的零偏移点处，先向西、后向东的位移恰好抵消，落地时仍有向东速度。
\end{solution}


\begin{exercise}[变转速导轨上的匀速外移]\label{rot:problem-exercise-10}\label{rot:ex:d6-h10}
水平径向导轨绕固定竖直轴转动，电机使角速度满足 $\Omega(t)=\Omega_0+\alpha t$，其中 $\Omega_0\ge0$、$\alpha>0$。另一径向执行器控制小珠沿导轨运动，使 $r(t)=a+ut$，$a,u>0$。小珠质量为 $m$，导轨光滑。求径向执行器的力、导轨的方位向力，以及两个机构各自传给小珠的功率。
\end{exercise}
\begin{figure}[H]\centering
\begin{tikzpicture}
\draw[curve](0,-.12)rectangle(4.3,.12);\fill(0,0)circle(1.5pt)node[below left]{$O$};\fill(2.8,0)circle(3pt)node[below]{$m$};\draw[guide](0,-.6)--(2.8,-.6)node[midway,below]{$r=a+ut$};\draw[axis](2.8,.4)--(4.2,.4)node[right]{$u\vct e_r$};\draw[axis](2.8,.1)--(2.8,1.5)node[above]{$\vct e_\varphi$};\draw[axis](-.7,.7)arc(135:225:1);\node at(-1.6,.7){$\Omega,\alpha$};
\end{tikzpicture}
\caption{径向位移和转动均由外部机构规定；竖直支持力与重力抵消。}\label{rot:fig:d6-h10-q}
\end{figure}
\begin{solution}
图\ref{rot:fig:d6-h10-q}中 $\dot r=u$、$\ddot r=0$。地面系极坐标加速度为
\[
a_r=\ddot r-r\Omega^2=-r\Omega^2,
\qquad
a_\varphi=r\dot\Omega+2\dot r\Omega=r\alpha+2u\Omega.
\]
两倍交叉项来自极坐标基矢也在转动。由 $\dot{\vct e}_r=\Omega\vct e_\varphi$、$\dot{\vct e}_\varphi=-\Omega\vct e_r$，对 $\vct v=\dot r\vct e_r+r\Omega\vct e_\varphi$ 求导，$\dot r\vct e_r$ 和 $r\Omega\vct e_\varphi$ 各贡献一个 $\dot r\Omega\vct e_\varphi$，合成 $2u\Omega$。
导轨没有径向摩擦，故径向力全部由执行器提供；方位向力全部由导轨提供。以 $\vct e_r$、$\vct e_\varphi$ 为正向，
\[
\boxed{F_r=-mr\Omega^2,\qquad N_\varphi=m(r\alpha+2u\Omega).}
\]
\begin{figure}[H]\centering
\begin{tikzpicture}
\begin{scope}
\fill(0,0)circle(2pt);\draw[axis](0,0)--(-1.6,0)node[above]{$mr\Omega^2$};\draw[axis](0,0)--(0,1.7)node[above]{$N_\varphi$};\node at(-.1,-.65){真实水平力};
\end{scope}
\begin{scope}[xshift=4.5cm]
\fill(0,0)circle(2pt);\draw[axis](0,0)--(1.6,0)node[above]{$mr\Omega^2$};\draw[axis](0,0)--(0,-1.2)node[right]{$2m\Omega u$};\draw[axis](-.25,0)--(-.25,-1.9)node[left]{$m\alpha r$};\node at(.5,.8){随轨系惯性力};
\end{scope}
\end{tikzpicture}
\caption{向内的执行器力平衡离心力；方位向约束力平衡科里奥利力与欧拉力。右图箭长只示方向。}\label{rot:fig:d6-h10-s}
\end{figure}
随导轨观察，小珠作匀速直线运动，相对加速度为零。图\ref{rot:fig:d6-h10-s}中的真实力与惯性力合力为零，其中 $m\alpha r$ 来自欧拉惯性力，$2m\Omega u$ 来自科里奥利力。

小珠的地面速度为 $u\vct e_r+r\Omega\vct e_\varphi$。径向执行器功率为
\[
\boxed{P_r=F_ru=-mr\Omega^2u.}
\]
执行器施力向内而小珠向外运动，故从小珠回收机械能。导轨对小珠绕 $O$ 的力矩为 $rN_\varphi$，旋转电机经导轨传入的功率为
\[
\boxed{P_{\rm rot}=\Omega rN_\varphi
=m\alpha\Omega r^2+2m\Omega^2ru.}
\]
相加得 $P_r+P_{\rm rot}=m\alpha\Omega r^2+m\Omega^2ru$。另一方面，
\[
K=\frac12m[u^2+(r\Omega)^2],\qquad
\dot K=mr\Omega(u\Omega+r\alpha),
\]
恰与总功率相符。角动量也给出同一方位向力：
\[
\frac{\dd}{\dd t}(mr^2\Omega)=m(2ru\Omega+r^2\alpha)=rN_\varphi.
\]
以上功率仅指传给小珠的部分，不包括导轨、转轴和电机自身转动惯量所需的功率。
\end{solution}


\begin{exercise}[转动径向弹簧的平衡位置]\label{rot:problem-exercise-11}\label{rot:ex:d6-h11}
光滑水平径向导轨以恒定角速度 $\Omega\ge0$ 绕 $O$ 转动。质量为 $m$ 的小珠穿在导轨上，通过劲度系数为 $k$、原长为 $\ell>0$ 的弹簧连接 $O$，弹簧可伸可压。只考虑 $r>0$ 的运动。求相对平衡位置、存在条件及小振动频率，讨论转速达到或超过临界值时的运动。
\end{exercise}
\begin{figure}[H]\centering
\begin{tikzpicture}
\draw[curve](0,-.13)rectangle(4.7,.13);\fill(0,0)circle(1.7pt)node[below left]{$O$};\draw[spring](.1,0)--(3.1,0);\fill(3.2,0)circle(3pt)node[above]{$m$};\draw[<->](0,-.6)--(3.2,-.6)node[midway,below]{$r$};\node at(1.5,.7){$k,\ \ell$};\draw[axis](-.6,.6)arc(135:225:.85);\node at(-1.3,.6){$\Omega$};
\end{tikzpicture}
\caption{弹簧沿导轨，导轨的水平约束力垂直于径向。}\label{rot:fig:d6-h11-q}
\end{figure}
\begin{solution}
记 $\omega_0=\sqrt{k/m}$。图\ref{rot:fig:d6-h11-q}中，弹簧径向力为 $-k(r-\ell)$；离心力为 $m\Omega^2r$。相对运动方程为
\[
m\ddot r=-k(r-\ell)+m\Omega^2r,
\qquad
\ddot r+(\omega_0^2-\Omega^2)r=\omega_0^2\ell.
\]
当 $\Omega<\omega_0$ 时，
\[
\boxed{r_* =\frac{k\ell}{k-m\Omega^2}
=\frac{\omega_0^2\ell}{\omega_0^2-\Omega^2}>0.}
\]
令 $r=r_*+\xi$，消去常数项，得 $\ddot\xi+(\omega_0^2-\Omega^2)\xi=0$。特征根为 $\pm i\sqrt{\omega_0^2-\Omega^2}$，所以
\[
\boxed{\omega_{\rm small}=\sqrt{\omega_0^2-\Omega^2}.}
\]
本题只讨论正半轴，振幅须小于 $r_*$，使整个振动过程保持 $r>0$。
\begin{figure}[H]\centering
\begin{tikzpicture}
\begin{scope}
\draw[axis](0,0)--(3.2,0)node[right]{$r/\ell$};\draw[axis](0,-.8)--(0,2.4)node[above]{$U_{\rm eff}/(k\ell^2)$};\draw[curve,domain=0:2.7,samples=70]plot({\x},{1.4*(.25*\x^2-\x+.5)});\draw[guide](2,0)--(2,-.7);\fill(2,-.7)circle(1.7pt);\node[above]at(2,0){$r_*/\ell$};\node at(1.5,1.8){$\Omega^2/\omega_0^2=1/2$};
\end{scope}
\begin{scope}[xshift=5cm]
\draw[axis](0,0)--(3.2,0)node[right]{$r/\ell$};\draw[axis](0,-2.1)--(0,1.5);\draw[curve](0,.7)--(1.9,-1.96);\node at(1.6,1){$\Omega=\omega_0$};\node at(1.6,-2.5){$U_{\rm eff}$ 单调下降};
\end{scope}
\end{tikzpicture}
\caption{允许区域为 $r\ge0$。左图有势能极小点；右图临界时已无极小点。纵轴使用同一任意正比例缩放。}\label{rot:fig:d6-h11-s}
\end{figure}
有效势能为
\[
U_{\rm eff}=\frac12k(r-\ell)^2-\frac12m\Omega^2r^2.
\]
图\ref{rot:fig:d6-h11-s}左图的二次项系数为正，平衡点是最小值。若 $\Omega=\omega_0$，二次项消失而线性项仍在，运动方程成为
\[
\ddot r=\omega_0^2\ell>0,
\]
在这个临界转速，任取初始半径 $r_i>0$ 和径向初速度 $v_i$，两次积分为
\[
\dot r-v_i=\int_0^t\omega_0^2\ell\,\dd\tau=\omega_0^2\ell t,\qquad
r-r_i=\int_0^t(v_i+\omega_0^2\ell\tau)\,\dd\tau
=\left[v_i\tau+\frac12\omega_0^2\ell\tau^2\right]_0^t.
\]
因此，到碰到 $r=0$ 边界之前的解为
\[
r(t)=r_i+v_it+\frac12\omega_0^2\ell t^2.
\]
若从相对静止释放，半径按时间的平方增大，因而零线性频率在这里不代表中性平衡。
没有相对平衡位置。若 $\Omega>\omega_0$，则对所有 $r>0$，
\[
\ddot r=(\Omega^2-\omega_0^2)r+\omega_0^2\ell>0.
\]
令 $\kappa=\sqrt{\Omega^2-\omega_0^2}$，则方程为 $\ddot r-\kappa^2r=\omega_0^2\ell$。取常数特解 $r_p=-\omega_0^2\ell/\kappa^2$，齐次方程的特征根为 $\pm\kappa$，故通解是
\[
r=-\frac{\omega_0^2\ell}{\kappa^2}+C\cosh\kappa t+D\sinh\kappa t.
\]
例如从 $r_i>0$ 相对静止释放，$r(0)=r_i$ 给出 $C=r_i+\omega_0^2\ell/\kappa^2$，$\dot r(0)=\kappa D=0$ 给出 $D=0$，于是
\[
r(t)=\left(r_i+\frac{\omega_0^2\ell}{\kappa^2}\right)\cosh\kappa t
-\frac{\omega_0^2\ell}{\kappa^2}.
\]
它始终为正并单调向外增长，比临界转速下的平方增长更快。
平衡方程此时只有负的 $r$ 解，落在允许区域之外。

导轨把运动限制为一个自由度，科里奥利力沿方位向，由 $N_\varphi=2m\Omega\dot r$ 平衡，不进入径向方程，所以有效势足以判断本题的稳定性。例\ref{rot:ex:d6-e10}没有这条导轨，两个方向仍由科里奥利项耦合，须联立方程判断。
\end{solution}


\begin{exercise}[旋转坐标中的闭合轨迹（选做）]\label{rot:problem-exercise-12}\label{rot:ex:d6-h12}
水平面内的质点受中心力 $\vct F=-k\vct r$，质量为 $m$，$\omega_0=\sqrt{k/m}$。取绕 $O$ 以 $\Omega=2\omega_0$ 逆时针匀速转动的坐标系。初始时坐标轴与地面轴重合，小珠位于 $(x',y')=(a,0)$，相对速度为零。求相对轨迹的参数方程、最大和最小半径、轨迹周期及相对速度为零的位置，解释有效势为负时轨迹为何仍有界。
\end{exercise}
\begin{figure}[H]\centering
\begin{tikzpicture}
\draw[axis](-1,0)--(4,0)node[right]{$x'$};\draw[axis](0,-1)--(0,2.8)node[above]{$y'$};\fill(0,0)circle(1.5pt)node[below left]{$O$};\fill(2.5,0)circle(2.5pt)node[below]{$(a,0)$};\draw[axis](2.5,0)--(1,0)node[midway,above]{$ka$};\draw[axis](2.5,.15)--(2.5,1.7)node[right]{$\vct v(0)=2\omega_0a\,\vct e_y$};\node at(1.3,-.85){$\vct v'(0)=0$};\node at(1.5,2.7){$\Omega=2\omega_0$};
\end{tikzpicture}
\caption{相对静止的初始条件对应非零地面速度；力始终指向O。}\label{rot:fig:d6-h12-q}
\end{figure}
\begin{solution}
图\ref{rot:fig:d6-h12-q}中，地面初始速度为 $(0,2\omega_0a)$。地面坐标分别满足 $\ddot x+\omega_0^2x=0$、$\ddot y+\omega_0^2y=0$，通解是
\[
x=C_x\cos\omega_0t+D_x\sin\omega_0t,\qquad
y=C_y\cos\omega_0t+D_y\sin\omega_0t.
\]
初位置给出 $C_x=a,C_y=0$；对通解求导后代入初速度，得 $\omega_0D_x=0,\omega_0D_y=2\omega_0a$。因此令 $\tau=\omega_0t$，有
\[
x=a\cos\tau,\qquad y=2a\sin\tau.
\]
将此矢量顺时针转过坐标轴已经转过的角 $2\tau$，可得复坐标
\[
\zeta=x'+iy'=a e^{-2i\tau}(\cos\tau+2i\sin\tau)
=\frac a2(3e^{-i\tau}-e^{-3i\tau}).
\]
故所求参数方程为
\[
\boxed{x'=\frac a2(3\cos\tau-\cos3\tau),\qquad
y'=\frac a2(-3\sin\tau+\sin3\tau).}
\]
\begin{figure}[H]\centering
\begin{tikzpicture}
\draw[axis](-2.1,0)--(2.4,0)node[right]{$x'/a$};\draw[axis](0,-3)--(0,3.1)node[above]{$y'/a$};
\draw[curve,domain=0:360,samples=241]plot({1.25*(1.5*cos(\x)-.5*cos(3*\x))},{1.25*(-1.5*sin(\x)+.5*sin(3*\x))});
\fill(1.25,0)circle(1.7pt);\fill(-1.25,0)circle(1.7pt);\draw[guide](1.25,0)--(2,.7)node[right]{$(1,0)$};\draw[guide](-1.25,0)--(-2,-.7)node[left]{$(-1,0)$};\fill(0,-2.5)circle(1.7pt)node[below right]{$(0,-2)$};\fill(0,2.5)circle(1.7pt)node[above left]{$(0,2)$};\draw[axis,domain=80:100,samples=20]plot({1.25*(1.5*cos(\x)-.5*cos(3*\x))},{1.25*(-1.5*sin(\x)+.5*sin(3*\x))});\draw[axis,domain=260:280,samples=20]plot({1.25*(1.5*cos(\x)-.5*cos(3*\x))},{1.25*(-1.5*sin(\x)+.5*sin(3*\x))});
\end{tikzpicture}
\caption{相对轨迹在 $(\pm a,0)$ 处有尖点，质点在这两点相对瞬时静止。}\label{rot:fig:d6-h12-s}
\end{figure}
坐标转动保持矢量长度，故
\[
r^2=x'^2+y'^2=a^2(\cos^2\tau+4\sin^2\tau)
=a^2(1+3\sin^2\tau).
\]
因此 $r_{\min}=a$、$r_{\max}=2a$；它们在图\ref{rot:fig:d6-h12-s}中分别对应横轴上的两个尖点和纵轴上的两个端点。复坐标按指数函数求导，先写成
\[
\dot\zeta=\frac a2\left[-3i\omega_0e^{-i\tau}+3i\omega_0e^{-3i\tau}\right]
=\frac{3ia\omega_0}{2}e^{-2i\tau}(e^{-i\tau}-e^{i\tau}).
\]
用 $e^{-i\tau}-e^{i\tau}=-2i\sin\tau$，得到
\[
\dot\zeta=3a\omega_0e^{-2i\tau}\sin\tau.
\]
相对速度在 $\tau=n\pi$ 时为零，对应位置 $(a,0)$、$(-a,0)$ 交替出现。坐标中包含频率为 $\omega_0$ 的非零谐波，故完整轨迹的最小周期为
\[
\boxed{T_{\rm orbit}=\frac{2\pi}{\omega_0}.}
\]
从参数式可直接核对最小周期：$\zeta(t+\pi/\omega_0)=-\zeta(t)$，而 $\zeta$ 不为零，所以半个周期尚未回到原位置；再过相同时间，两个非零谐波的相位都恢复。
半径的周期为 $\pi/\omega_0$，经过这段时间位置矢量反向，完整轨迹尚未闭合。

此时 $U_{\rm eff}=\frac12m(\omega_0^2-\Omega^2)r^2=-\frac32m\omega_0^2r^2$。由初始相对速度为零，
\[
E'=-\frac32m\omega_0^2a^2<0.
\]
由已求出的轨迹，质点始终位于 $a\le r\le2a$ 的环带内。有效势向外降低，离心力超过回复力，但科里奥利力持续改变相对速度方向，两个分量耦合后形成图\ref{rot:fig:d6-h12-s}中的闭合轨道。
\end{solution}


\chapter{连续体问题}



\section{连续体的微元、弹性与边界功}

连续体的内部受力和形变随位置变化，仅有整体平衡通常不能确定应力与位移。微元方程结合材料关系求局部解，再由连接条件和端部功确定整体响应。



\subsection{连续体微元与弹性常数}
杨氏模量$Y$联系单向拉伸的应力和应变；剪切模量$G$联系切应力和小剪切角；体积模量$K=-V\dd p/\dd V$联系均匀压强与相对体积变化。三者对应不同形变，不能互换。泊松比描述拉伸时横向与纵向应变的负比值。

对初始位置$x$处的质元，$u(x,t)$表示其位移，纵向应变为$u_x$。截面内力$F(x)$与$F(x+\dd x)$的差是微元合力。用本构关系表示截面内力，再对长度为$\dd x$的质元列牛顿方程，便得到波动方程。以下用未形变时的材料坐标标记质元，当前位置由该坐标加位移得到。
\subsubsection{边界条件、能流与弹性模型}
一维小振动常可写成$w(x)u_{tt}=\partial_x[p(x)u_x]$。$w$表示单位长度惯性，$pu_x$表示与位移共轭的截面力或力矩。其能量满足
\[
 E=\frac12\int_0^L(wu_t^2+pu_x^2)\dd x,\qquad
 \dot E=[pu_xu_t]_0^L.
\]
固定端规定$u=0$，自由端规定$pu_x=0$；无集中质量的连接处要求位移及截面内力连续。自由端通常可以运动，其边界条件来自截面内力为零。分离变量后，空间方程为$-(pX')'=\omega^2wX$，频率由材料、几何和边界共同决定。

对各向同性线弹性材料，泊松比记为$\nu_P$。单向应力产生纵向应变$\sigma/Y$和两个横向应变$-\nu_P\sigma/Y$；将三个正交方向的载荷叠加，可得
\[
 \epsilon_x=\frac{\sigma_x-\nu_P(\sigma_y+\sigma_z)}Y,
\]
另两个方向由下标循环置换得到。均匀压强下$\sigma_x=\sigma_y=\sigma_z=-p$，所以体积应变为
\[
 \epsilon_V=\epsilon_x+\epsilon_y+\epsilon_z
 =-\frac{3(1-2\nu_P)}Yp.
\]
由$K=-p/\epsilon_V$，得$K=Y/[3(1-2\nu_P)]$。再考虑纯剪切应力$\tau$：将坐标轴转过$45^\circ$，两个主应力分别为$\tau,-\tau$，第三个为零。相应主应变为$\epsilon_1=(1+\nu_P)\tau/Y$、$\epsilon_2=-(1+\nu_P)\tau/Y$。小形变时，原坐标的工程剪切应变是$\gamma=\epsilon_1-\epsilon_2=2(1+\nu_P)\tau/Y$，与$\tau=G\gamma$比较，得到
\[
 Y=2G(1+\nu_P).
\]
这些关系适用于同一种各向同性线弹性材料。细杆纵波、体纵波和剪切波的应力约束不同，传播速度仍可以不同。例如体纵波的横向应变为零，令$\epsilon_y=\epsilon_z=0$，由对称性有$\sigma_y=\sigma_z=\nu_P\sigma_x/(1-\nu_P)$。代回$\epsilon_x$，得到
\[
 \sigma_x=\frac{Y(1-\nu_P)}{(1+\nu_P)(1-2\nu_P)}\epsilon_x
 =\left(K+\frac43G\right)\epsilon_x.
\]
因此体纵波的有效纵向模量为$K+4G/3$，而横向可自由收缩的细杆只用$Y$。将相应模量代入微元运动方程，才得到各自的波速。

\LectureNeedspace{8\baselineskip}


\subsection{弹性系统的能量与边界功}
\label{sec:deep7-energy}
细杆的位移是材料点的位移。以未形变杆上的坐标$x$标记材料点，令截面刚度$p(x)=Y(x)S(x)$、线密度$\lambda(x)=\rho(x)S(x)$。在线弹性、小应变近似下，轴力$N=p u_x$以拉伸为正；若单位原长体力为$f$，微段运动方程为
\[
 \lambda u_{tt}=\partial_x(pu_x)+f.
\]
截面有突加集中力时，$N$可以跳变；没有裂隙或脱开的连接处，位移仍须连续。截面力由接头受力决定，位移连续则来自连接处不脱开的条件。

方程乘以 $u_t$，其中 $p,\lambda$ 不随时间变化。分别整理惯性项和内力项，
\[
\lambda u_{tt}u_t=\partial_t(\lambda u_t^2/2),\qquad
\partial_x(pu_x)u_t=\partial_x(pu_xu_t)-pu_xu_{xt},
\]
且 $pu_xu_{xt}=\partial_t(pu_x^2/2)$。于是得到局部能量关系
\[
 \partial_t\left(\frac12\lambda u_t^2+\frac12p u_x^2\right)
 +\partial_x(-Nu_t)=fu_t.
\]
括号内两项分别为单位长度动能和弹性能，沿$x$正向的功率流为$\mathcal P=-Nu_t$。将此式在$[a,b]$上积分，
\[
 \frac{\dd E_{[a,b]}}{\dd t}
 =N(b)u_t(b)-N(a)u_t(a)+\int_a^bfu_t\dd x.
\]
右端的端点项正好是外界对这段杆所做的功率。端点固定时速度为零；自由端轴力为零。这两种边界都没有端点能流，却给出不同的位移反射符号。

对均匀杆，$c=\sqrt{p/\lambda}$，阻抗$Z=\lambda c=p/c$。右行波$u=F(t-x/c)$满足$N=-Zu_t$，左行波$u=G(t+x/c)$满足$N=+Zu_t$。若记两列波的粒子速度为$v_+,v_-$，则
\[
 v=v_++v_-,\qquad N=Z(-v_++v_-),\qquad
 \mathcal E=\lambda(v_+^2+v_-^2).
\]
速度和轴力先叠加，再计算能量。上式中的交叉项在动能与弹性能之间抵消；两列波相遇时，局部粒子速度或应力可以加倍，总能量仍等于两列波各自能量之和。

\subsubsection{由弹性能求位移}
线弹性静力系统的应变能为轴力、弯矩和扭矩各自贡献之和。在细长杆梁的适用范围内，
\[
 U_e=\int\frac{N^2}{2YS}\dd x
 +\int\frac{M_b^2}{2YI}\dd x
 +\int\frac{M_t^2}{2GJ_p}\dd x.
\]
这些积分采用相应的细杆、细梁近似；剪切变形和截面翘曲若不可忽略，须另计其贡献。若外力为缓慢施加的独立广义力$P_i$，其作用点沿相应方向的位移为$q_i$，线性关系写成$q_i=\sum_jC_{ij}P_j$。从无载荷态按同一比例加载，令$P_i(s)=sP_i$、$q_i(s)=sq_i$，其中$0\le s\le1$。没有耗散时，外力功全部储存为应变能，故
\[
 U_e=\int_0^1\sum_i sP_iq_i\,\dd s
 =\left[\frac{s^2}{2}\right]_0^1\sum_iP_iq_i
 =\frac12\sum_{i,j}P_iC_{ij}P_j.
\]
为说明柔度矩阵为何对称，先不预设$C_{ij}=C_{ji}$。将上式对任一载荷$P_k$求偏导，两个因子各贡献一项：
\[
 \frac{\partial U_e}{\partial P_k}
 =\frac12\sum_j(C_{kj}+C_{jk})P_j.
\]
另一方面，对任意缓慢的载荷增量，虚功关系给出
\[
 \dd U_e=\sum_iP_i\dd q_i
 =\sum_k\left(\sum_iP_iC_{ik}\right)\dd P_k.
\]
比较$\dd P_k$的系数，并利用各个$P_j$相互独立，得$C_{kj}=C_{jk}$。再代回能量的偏导，就有
\[
 \frac{\partial U_e}{\partial P_k}=\sum_jC_{kj}P_j=q_k.
\]
求偏导时须保留其他独立载荷。若先把两个力设置成相等再求导，得到的是两处位移之和，未必是某一个作用点的位移。分段杆的两个载荷点可分别说明这两种导数的含义。

\LectureNeedspace{4\baselineskip}


\subsection{例题与模型变式}

\begin{example}[有质量弹簧的自重伸长与纵波]\label{ex:hw7-3}

均匀弹簧原长$l$、质量$m$、整根劲度系数$k$。(1)竖直悬挂，上端固定，下端自由，求静止长度；(2)将其放在光滑水平面上，弹簧足够长，求小纵波传播速度。

\end{example}


\begin{figure}[H]\centering
\begin{tikzpicture}[>=Stealth,line cap=round,line join=round]
\draw[thick](-.5,0)--(.5,0);\draw[spring](0,0)--(0,-3);\draw[axis](.8,0)--(.8,-3.4)node[below]{$x$};\node at(-1,-1.5){$l,m,k$};
\end{tikzpicture}
\caption{自重作用下竖直悬挂的有质量弹簧}\label{fig:hw7-3-q}
\end{figure}

\begin{solution}
取未伸长材料坐标$x$从上端向下。将弹簧等分为 $n$ 段，各段劲度为 $k_n$，串联关系 $1/k=n/k_n$ 给出 $k_n=nk$。一段原长 $\Delta x=l/n$，伸长为 $u_x\Delta x$，故 $T=k_nu_x\Delta x=klu_x$。图\ref{fig:hw7-3-s}中静态截面承担其下方质量，$T(x)=(m/l)(l-x)g$。

\begin{figure}[H]\centering
\begin{tikzpicture}[>=Stealth,line cap=round,line join=round]
\draw[thick](-.5,0)--(.5,0);
\draw[curve,domain=0:1,samples=420]plot({.13*sin(5040*\x)},{-3.5*\x+.5*\x*\x});
\node[left]at(-.3,-.25){$x=0$};\node[left]at(-.3,-3){$x=l$};\node[left]at(-.35,-1.625){截面$x$};
\draw[guide](-.4,-1.625)--(.5,-1.625);\draw[axis](0,-1.625)--(0,-.9)node[right]{$T(x)$};
\node at(0,-3.6){上段伸长率较大};
\begin{scope}[shift={(3.3,-2.8)}]
\draw[axis](0,0)--(2.8,0)node[right]{$x/l$};\draw[axis](0,0)--(0,2.5)node[above]{$T/(mg)$};
\draw[curve](0,2)--(2.2,0);\node[left]at(0,2){$1$};\node[below]at(2.2,0){$1$};\node at(1.4,-.7){按原长标记截面};
\end{scope}
\end{tikzpicture}
\caption{自重使上段伸长率较大，张力沿原长坐标线性减小}\label{fig:hw7-3-s}
\end{figure}

由 $klu_x=T(x)$，先得 $u_x=mg(l-x)/(kl^2)$。用固定端条件 $u(0)=0$，逐段累加伸长：
\[
 u(x)=\int_0^x\frac{mg}{kl^2}(l-s)\,\dd s
 =\frac{mg}{kl^2}\left[ls-\frac{s^2}{2}\right]_0^x
 =\frac{mg}{kl^2}\left(lx-\frac{x^2}{2}\right).
\]
上段承重更多、伸长率更大。图\ref{fig:hw7-3-s}中的截面$x$应按原长坐标识别，不能将变形后到悬点的实际距离直接代入$m(1-x/l)$。

在末端代入 $x=l$，$u(l)=mg(l^2-l^2/2)/(kl^2)=mg/(2k)$，故
\[
 \boxed{l_{\text{静}}=l+mg/(2k)}.
\]
弹簧水平时，对微元列 $kl u_{xx}\dd x=(m/l)\dd x\,u_{tt}$，即 $u_{tt}=(kl^2/m)u_{xx}$。为求传播速度，代入行波 $u=A\cos(qx-\omega t)$。逐次求导得 $u_x=-Aq\sin(qx-\omega t)$、$u_{xx}=-Aq^2\cos(qx-\omega t)$，以及 $u_t=A\omega\sin(qx-\omega t)$、$u_{tt}=-A\omega^2\cos(qx-\omega t)$。比较系数，$\omega^2=(kl^2/m)q^2$，于是
\[
 \boxed{c=l\sqrt{k/m}}.
\]
这一速度按原长坐标定义；若弹簧已经显著预伸长且本构关系改变，应重新线性化。

在竖直静平衡附近研究小振动时，也可令总位移为$u_0(x)+w(x,t)$。静态方程$klu_0''+(m/l)g=0$与动力学方程相减，仍得$w_{tt}=l^2(k/m)w_{xx}$，但顶端、底端的扰动边界条件分别为$w(0,t)=0$、$w_x(l,t)=0$。自重改变平衡形状；在这一本构近似中，它不额外改变材料坐标下的线性波速。

\end{solution}



\begin{example}[圆弧软绳的瞬时冲量分布]\label{ex:hw7-7}

质量$m$、半径$r$、圆心角$\theta_0$的均匀不可伸长柔绳，初始静止置于光滑水平面，围成一段圆弧。于$A$端沿离开绳的切向施加大小$I_0$的瞬时拉冲量，$B$端自由。以$\theta\in[0,\theta_0]$沿绳从$A$向$B$计，求绳各处碰后瞬时速度。

\end{example}


\begin{figure}[H]\centering
\begin{tikzpicture}[>=Stealth,line cap=round,line join=round]
\draw[curve,domain=0:85,samples=70]plot({2.2*cos(\x)},{2.2*sin(\x)});\draw[guide](0,0)--(2.2,0);\draw[guide](0,0)--({2.2*cos(85)},{2.2*sin(85)});\node[below]at(2.2,0){$A,\theta=0$};\node[above]at({2.2*cos(85)},{2.2*sin(85)}){$B,\theta_0$};\draw[axis](2.2,0)--(2.2,-1)node[right]{$I_0$};
\end{tikzpicture}
\caption{圆弧柔绳端点受到向外的切向冲量}\label{fig:hw7-7-q}
\end{figure}

\begin{solution}
设张力对时间的积分为$P(\theta)$，切向单位矢量$\vct t$沿$\theta$增加，外径向单位矢量为$\vct e_r$，有$\partial_\theta\vct t=-\vct e_r$。图\ref{fig:hw7-7-s}对小绳段列冲量方程，线密度$\lambda=m/(r\theta_0)$：
\[
 \lambda r\vct v=P'\vct t-P\vct e_r.
\]

\begin{figure}[H]\centering
\begin{tikzpicture}[>=Stealth,line cap=round,line join=round]
\draw[curve,domain=0:85,samples=70]plot({2.2*cos(\x)},{2.2*sin(\x)});
\draw[guide](0,0)--(2.2,0);\draw[guide](0,0)--({2.2*cos(85)},{2.2*sin(85)});
\node[below left]at(2.2,0){$A$};\node[above]at(.19174,2.19163){$B$};\draw[axis](2.2,0)--(2.2,-.9)node[right]{$I_0$};
\draw[axis](1.555635,1.555635)--(1.06066,2.05061)node[above left]{$\vct t$};
\draw[axis](1.555635,1.555635)--(2.192031,2.192031)node[right]{$\vct e_r$};
\begin{scope}[shift={(5.2,1.5)}]
\fill(0,0)circle(1.5pt);\draw[axis](0,0)--(-.5,-.5)node[left]{$v_r\vct e_r$};
\draw[axis](0,0)--(.75,-.75)node[right]{$v_t\vct t$};
\draw[guide](-.5,-.5)--(.25,-1.25)--(.75,-.75);\draw[axis](0,0)--(.25,-1.25)node[below]{$\vct v$};
\node[align=center]at(.2,.8){$v_t=P'/(\lambda r)<0$\\[4pt]$v_r=-P/(\lambda r)\le0$};\end{scope}
\end{tikzpicture}
\caption{局部基矢与速度分量分开表示，径向分量一般不为零}\label{fig:hw7-7-s}
\end{figure}

将速度分成$v_t=P'/(\lambda r)$与$v_r=-P/(\lambda r)$，并用$\vct e_r'=\vct t$、$\vct t'=-\vct e_r$，得
\[
 \partial_\theta\vct v=\frac1{\lambda r}
 [(P''-P)\vct t-2P'\vct e_r].
\]
绳长的一阶变化率正比于这一导数的切向分量，所以它必须为零。

不可伸长条件要求 $\partial_\theta\vct v\cdot\vct t=0$，因此 $P''-P=0$。特征根为 $\pm1$，可取通解 $P=C\cosh\theta+D\sinh\theta$。边界 $P(0)=I_0$ 给出 $C=I_0$；自由端条件 $P(\theta_0)=0$ 给出 $D=-I_0\coth\theta_0$。用双曲函数差角公式整理，得到
\[
 P(\theta)=I_0\frac{\sinh(\theta_0-\theta)}{\sinh\theta_0}.
\]
再求导，$P'=-I_0\cosh(\theta_0-\theta)/\sinh\theta_0$，其中负号来自内层函数 $\theta_0-\theta$ 的导数。代回速度分量，得到
\[
 \boxed{\vct v(\theta)=-\frac{I_0\theta_0}{m\sinh\theta_0}
 [\cosh(\theta_0-\theta)\vct t+
 \sinh(\theta_0-\theta)\vct e_r]}.
\]
自由端$B$处虽然$P(\theta_0)=0$，其导数却不为零，因此该端仍获得切向速度$-I_0\theta_0\vct t(\theta_0)/(m\sinh\theta_0)$。在受冲击端，径向分量也不为零；端点的切向外冲量经绳内张力传递后，一般会使绳元同时获得切向和径向速度。图\ref{fig:hw7-7-s}把这些分量分别画出。

张力冲量非负，符合柔绳只能拉的条件。全绳动量的积分可利用 $\vct t'=-\vct e_r$ 合并：
\[
\vct p=\int_0^{\theta_0}(P'\vct t-P\vct e_r)\,\dd\theta
=\int_0^{\theta_0}(P\vct t)'\,\dd\theta
=[P\vct t]_0^{\theta_0}=-I_0\vct t(0).
\]
动能用 $\dd m=\lambda r\,\dd\theta$ 计算。由于 $P''=P$，有 $(PP')'=P'^2+P^2$，故
\[
\begin{aligned}
K&=\frac1{2\lambda r}\int_0^{\theta_0}(P'^2+P^2)\,\dd\theta
=\frac1{2\lambda r}[PP']_0^{\theta_0}\\
&=\frac{I_0^2\coth\theta_0}{2\lambda r}
=\frac{I_0^2\theta_0\coth\theta_0}{2m}.
\end{aligned}
\]
当 $\theta_0\to0$，$\theta_0\coth\theta_0\to1$，动能趋于整条直绳取得同一切向速度的结果 $I_0^2/(2m)$。速度的大小和方向均随弧长位置改变，正是柔绳不可伸长约束在弯曲构形中的结果。


\end{solution}



\begin{example}[分段杆的轴力、位移与柔度]\label{deep7:ex-stepped}
两段同材料细杆首尾连接，左段$OA$长$a$、截面积$S_1$，右段$AB$长$b$、截面积$S_2$，杨氏模量均为$Y$。$O$固定，在$A$和$B$处分别缓慢施加向右的恒力$Q$、$P$，$P,Q>0$。忽略自重，求轴力、位移分布和应变能；用应变能的偏导求两载荷点位移，并写出以$u_A,u_B$为坐标的刚度矩阵。
\begin{figure}[H]\centering
\begin{tikzpicture}
\draw[curve](0,-.8)--(0,.8);\draw[curve](0,-.22)rectangle(3,.22);\draw[curve](3,-.36)rectangle(6,-.36+.72);
\fill(3,0)circle(1.6pt);\fill(6,0)circle(1.6pt);
\node[above left]at(0,.8){$O$};\node[below]at(3,-.4){$A$};\node[below]at(6,-.4){$B$};
\draw[axis](3,.65)--(4.3,.65)node[right]{$Q$};\draw[axis](6,0)--(7.3,0)node[right]{$P$};
\draw[<->](0,-.95)--(3,-.95)node[midway,below]{$a,S_1$};\draw[<->](3,-.95)--(6,-.95)node[midway,below]{$b,S_2$};
\draw[axis](.3,1.2)--(1.7,1.2)node[right]{$x$};
\end{tikzpicture}
\caption{载荷$Q$在连接点施加，两段杆承受的轴力不同}\label{deep7:fig-stepped-q}
\end{figure}
\end{example}
\begin{solution}
取图\ref{deep7:fig-stepped-q}中任意截面右侧的部分列平衡方程。$AB$段只传递端力$P$，$OA$段须同时传递$P$和$Q$，因此
\[
 N(x)=\begin{cases}P+Q,&0<x<a,\\P,&a<x<a+b.\end{cases}
\]
在 $A$ 处 $N(a^+)-N(a^-)=-Q$。连接点外力使轴力产生跳变，杆的黏结连接则要求 $u(a^-)=u(a^+)$。分段积分 $u'=N/(YS)$：在左段，$u(x)=\int_0^x(P+Q)/(YS_1)\,\dd s=(P+Q)x/(YS_1)$；在右段，应从连接点接续，$u(x)=u(a)+\int_a^xP/(YS_2)\,\dd s$。因此
\[
 u(x)=\begin{cases}
 (P+Q)x/(YS_1),&0\le x\le a,\\
 (P+Q)a/(YS_1)+P(x-a)/(YS_2),&a\le x\le a+b.
 \end{cases}
\]
\begin{figure}[H]\centering
\begin{tikzpicture}[x=.9cm]
\draw[axis](0,0)--(7,0)node[right]{$x$};\draw[axis](0,0)--(0,2.6)node[above]{$N$};
\draw[curve](0,2.1)--(3,2.1);\draw[curve](3,1.1)--(6,1.1);\draw[guide](3,0)--(3,2.1);\draw[guide](6,0)--(6,1.1);
\node[left]at(0,2.1){$P+Q$};\node[left]at(0,1.1){$P$};\node[below]at(3,0){$a$};\node[below]at(6,0){$a+b$};
\draw[axis](4.1,2.1)--(4.1,1.1)node[midway,right]{$Q$};
\end{tikzpicture}
\caption{轴力在$A$处下降$Q$，两段弹性能须分别积分}\label{deep7:fig-stepped-s}
\end{figure}
按图\ref{deep7:fig-stepped-s}，应变能为 $\int N^2/(2YS)\,\dd x$，两段分别积分：
\[
U_e=\int_0^a\frac{(P+Q)^2}{2YS_1}\,\dd x
+\int_a^{a+b}\frac{P^2}{2YS_2}\,\dd x
=\frac{(P+Q)^2}{2YS_1}[x]_0^a+\frac{P^2}{2YS_2}[x]_a^{a+b},
\]
即
\[
 U_e=\frac{(P+Q)^2a}{2YS_1}+\frac{P^2b}{2YS_2}.
\]
保留 $P,Q$ 为两个独立变量。平方项按链式法则求偏导，例如 $\partial(P+Q)^2/\partial Q=2(P+Q)$、$\partial(P+Q)^2/\partial P=2(P+Q)$，从而
\[
 \frac{\partial U_e}{\partial Q}=\frac{(P+Q)a}{YS_1}=u_A,
 \qquad
 \frac{\partial U_e}{\partial P}=\frac{(P+Q)a}{YS_1}+\frac{Pb}{YS_2}=u_B.
\]
两式均与直接积分一致。若题设数值恰有$P=Q=F$，须在求偏导后再代入；把能量先写成$U_e(F,F)$再求导，链式法则给出$u_A+u_B$。

令$k_1=YS_1/a$、$k_2=YS_2/b$。两段伸长分别是$u_A$与$u_B-u_A$，所以
\[
 U_e=\frac12k_1u_A^2+\frac12k_2(u_B-u_A)^2.
\]
由 $Q=\partial U_e/\partial u_A=k_1u_A-k_2(u_B-u_A)$、$P=\partial U_e/\partial u_B=k_2(u_B-u_A)$，整理为
\[
 \begin{pmatrix}Q\\P\end{pmatrix}
 =\begin{pmatrix}k_1+k_2&-k_2\\-k_2&k_2\end{pmatrix}
 \begin{pmatrix}u_A\\u_B\end{pmatrix}.
\]
矩阵对称，其两个顺序主子式为$k_1+k_2>0$和$k_1k_2>0$，因而固定$O$后不存在零能量的整体平移。这个静态刚度矩阵若与节点质量矩阵相配，便可用于离散弹性系统的小振动。
\end{solution}



\section{行波、反射与动态边界}

波速取决于惯性和弹性，反射取决于端部或界面的条件。入射、反射与透射使用同一种物理量表示，再由位移连续和力平衡确定系数。



\subsection{从微元受力到行波}
\label{sec:complete7}
对截面积$S$、密度$\rho$的均匀细杆，纵向应变为$u_x$，线弹性截面力为$N=YSu_x$。长度$\dd x$微元的两端受力相反，合力为$N(x+\dd x)-N(x)$。因此
\[
 \rho S\dd x\,u_{tt}=YS\,u_{xx}\dd x,
 \qquad u_{tt}=c_L^2u_{xx},\quad c_L=\sqrt{Y/\rho}.
\]
对剪切长条，将$Y$换成$G$；对圆截面杆扭转，将位移换成转角$\theta$，截面力矩为$GJ_p\theta_x$、单位长度转动惯量为$\rho J_p$，得到同样的波速$\sqrt{G/\rho}$。$J_p=\int r_\perp^2\dd S$是截面的极二次矩，不能与整根杆的转动惯量混为一谈。

均匀介质的波动算符可分解为
\[
 (\partial_t-c\partial_x)(\partial_t+c\partial_x)u=0,
 \qquad u=f(x-ct)+g(x+ct).
\]
上述行波通解可在特征坐标 $\xi=x-ct$、$\zeta=x+ct$ 中求得。链式法则给出 $u_{tt}-c^2u_{xx}=-4c^2u_{\xi\zeta}=0$，先对 $\zeta$ 积分，$u_\xi$ 只依赖 $\xi$；再对 $\xi$ 积分，得 $u=f(\xi)+g(\zeta)$。函数 $f$ 沿 $x$ 增大方向传播，$g$ 沿相反方向传播。

无限长介质给定 $u(x,0)=f_0(x)$、$u_t(x,0)=g_0(x)$ 时，初值要求 $f+g=f_0$、$-cf'+cg'=g_0$。将前式求导后与后式联立，
\[
f'=\frac12(f_0'-g_0/c),\qquad
g'=\frac12(f_0'+g_0/c).
\]
分别积分，再将 $\xi=x-ct$、$\zeta=x+ct$ 代回；两原函数的常量之和由 $f+g=f_0$ 确定为零。$g_0$ 的两个积分相减，合成从 $x-ct$ 到 $x+ct$ 的定积分，得到达朗贝尔解
\[
 u(x,t)=\frac{f_0(x-ct)+f_0(x+ct)}2
 +\frac1{2c}\int_{x-ct}^{x+ct}g_0(\xi)\dd\xi.
\]
有界介质还须同时满足端部条件，不能把无限长介质的解直接截断。


\subsection{反射、连接与端部质量}
\label{sec:complete7-boundary}
纵杆或弦的统一截面力写成$pu_x$。固定端由$u=0$要求位移反射反号，自由端则由$u_x=0$要求位移反射同号。应力或截面力的反射符号与位移不完全相同，应从空间导数重新计算。

两个半无限均匀杆在$x=0$连接，取入射、反射、透射位移为
\[
 A_i e^{\mathrm i(\omega t-k_1x)},\quad
 A_r e^{\mathrm i(\omega t+k_1x)},\quad
 A_t e^{\mathrm i(\omega t-k_2x)}.
\]
连续条件为$A_i+A_r=A_t$及$p_1k_1(A_i-A_r)=p_2k_2A_t$。令机械阻抗$Z_j=p_j/c_j$，得到
\[
 \frac{A_r}{A_i}=\frac{Z_1-Z_2}{Z_1+Z_2},\qquad
 \frac{A_t}{A_i}=\frac{2Z_1}{Z_1+Z_2}.
\]
平均能流正比于$Z\omega^2|A|^2$，从而反射、透射能流系数之和为一。连接处无反射的条件是两侧阻抗相等，波速相同还不够。

端部若连接集中质量$M$，该处既非固定端也非无应力自由端。对右端质量列牛顿方程，得到
\[
 -p u_x(L,t)=M u_{tt}(L,t).
\]
力矩问题则将$M$换成端部转动惯量。动态边界条件把连续介质与集中参数系统联在一起，频率方程一般为超越方程。


\subsection{例题与模型变式}

\begin{example}[固定底端长条的剪切驻波]\label{ex:hw7-1}

密度$\rho$、剪切模量$G$、截面积$A$、长度$L$的均匀长条，下端固定，上端无应力。各横截面只沿同一个横向方向平移，截面不旋转，求这种剪切振动的固有周期。

\end{example}


\begin{figure}[H]\centering
\begin{tikzpicture}[>=Stealth,line cap=round,line join=round]
\draw[curve](0,0)rectangle(1.5,3);\draw[thick](-.4,0)--(1.9,0);\draw[guide](0,1.6)--(1.5,1.6);\draw[axis](1.5,1.6)--(2.5,1.6)node[above]{$u(x,t)$};\draw[axis](-.6,0)--(-.6,3.5)node[above]{$x$};\node at(.75,3.4){$A$};
\end{tikzpicture}
\caption{固定底端、自由顶端的剪切长条}\label{fig:hw7-1-q}
\end{figure}

\begin{solution}
图\ref{fig:hw7-1-s}中相邻截面相对位移为$u_x\dd x$，小剪切角为$u_x$，截面切向力为$GAu_x$。对微元列式：
\[
 GA[u_x(x+\dd x)-u_x(x)]=\rho A\dd x\,u_{tt},
 \quad u_{tt}=c_s^2u_{xx},\quad c_s=\sqrt{G/\rho}.
\]

\begin{figure}[H]\centering
\begin{tikzpicture}[>=Stealth,line cap=round,line join=round]
\draw[curve](0,0)rectangle(1.5,3);\draw[thick](-.4,0)--(1.9,0);\draw[guide](0,1.6)--(1.5,1.6);\draw[axis](1.5,1.6)--(2.5,1.6)node[above]{$u(x,t)$};\draw[axis](-.6,0)--(-.6,3.5)node[above]{$x$};\node at(.75,3.4){$A$};\draw[curve,domain=0:3,samples=60]plot({.7+.6*sin(30*\x)},\x);\node at(3.4,2.6){$u_x(L,t)=0$};\node at(3.4,.3){$u(0,t)=0$};
\draw[guide](.7,0)--(.7,3);\draw[<->](.7,2.35)--(1.266,2.35)node[midway,above]{$u$};
\end{tikzpicture}
\caption{剪切位移振型及两端边界条件}\label{fig:hw7-1-s}
\end{figure}

取 $u=X(x)\cos\omega t$，则 $u_{tt}=-\omega^2X\cos\omega t$、$u_{xx}=X''\cos\omega t$。代入波动方程并比较时间因子，得 $X''+k^2X=0$、$k=\omega/c_s$。对 $k>0$，通解及导数为
\[
X=C\sin kx+D\cos kx,\qquad X'=kC\cos kx-kD\sin kx.
\]
固定底端 $X(0)=0$ 给出 $D=0$；自由顶端 $X'(L)=0$ 给出 $kC\cos kL=0$。非零振型须 $C\ne0$，所以 $kL=(2n+1)\pi/2$。于是
\[
 \boxed{\omega_n=\frac{(2n+1)\pi}{2L}\sqrt{G/\rho},\quad
 T_n=\frac{4L}{2n+1}\sqrt{\rho/G},\quad n=0,1,\ldots}.
\]
基频振型为$X_0(x)=C\sin(\pi x/2L)$，顶端是位移腹而非位移节点，底端则相反。其四分之一波长恰等于$L$，所以$T_0=4L/c_s$。图\ref{fig:hw7-1-s}的虚线表示未形变截面的参考位置。

固定端约束位移，自由端约束应力，因而给出不同的边界条件。这里各截面只平移，所求为剪切波；细梁弯曲波还要计入截面转动。


弹性能变分的端点项给出同样的边界条件。由 $U=\tfrac12\int_0^LGAu_x^2\dd x$，先对平方项变分，再分部积分：
\[
\delta U=\int_0^LGAu_x(\delta u)_x\,\dd x
=[GAu_x\delta u]_0^L-\int_0^LGAu_{xx}\delta u\,\dd x.
\]
底端位移固定，$\delta u(0)=0$；顶端可以自由移动，$\delta u(L)$ 任意，故必须有 $GAu_x(L)=0$。这正是图\ref{fig:hw7-1-s}中顶部切线竖直、截面剪切力为零的条件。它不要求 $u(L)=0$。

该模型以截面不转动为题设约束。若将其用于没有这种约束的普通细长悬臂梁，弯曲与截面转动会参与运动，不能把这里的四分之一波长公式直接当作梁的固有周期。

\end{solution}



\begin{example}[自由圆柱的扭转振动]\label{ex:hw7-2}

密度$\rho$、剪切模量$G$、长度$L$、截面积$A$的均匀圆柱，在各截面圆周切线方向受到微扰，端面均无外力矩。求扭转振动频率，说明零频模式。

\end{example}


\begin{figure}[H]\centering
\begin{tikzpicture}[>=Stealth,line cap=round,line join=round]
\draw(0,0)ellipse(1 and .3);\draw(-1,0)--(-1,3);\draw(1,0)--(1,3);\draw(0,3)ellipse(1 and .3);\draw[guide](0,0)--(0,3.6);\draw[axis](.7,3)arc(0:240:.7 and .2)node[left]{$\vartheta$};\node at(2,1.6){$L,\rho,G$};
\end{tikzpicture}
\caption{两端无外力矩的圆柱扭转}\label{fig:hw7-2-q}
\end{figure}

\begin{solution}
令截面转角为$\vartheta(x,t)$，半径为$r$的质元切向位移为$r\vartheta$，剪切应变为$r\vartheta_x$。图\ref{fig:hw7-2-s}中截面力矩为
\[
 M=\int_A Gr\vartheta_x\,r\dd A=GJ_p\vartheta_x,
 \quad J_p=\int_A r^2\dd A.
\]

\begin{figure}[H]\centering
\begin{tikzpicture}[>=Stealth,line cap=round,line join=round]
\draw(0,0)ellipse(1 and .3);\draw(-1,0)--(-1,3);\draw(1,0)--(1,3);\draw(0,3)ellipse(1 and .3);\draw[guide](0,0)--(0,3.6);\draw[axis](.7,3)arc(0:240:.7 and .2)node[left]{$\vartheta$};\node at(2,1.6){$L,\rho,G$};\draw[guide](-1,1.3)arc(180:360:1 and .3);\draw[axis](.7,1.3)arc(0:250:.7 and .2);\node at(2.6,.5){$M=GJ_p\vartheta_x$};\node at(-2,1.6){$\tau=Gr\vartheta_x$};
\node[below]at(0,-.4){$x=0$};\node[above right]at(1,3.3){$x=L$};
\end{tikzpicture}
\caption{截面剪应力积分为扭转力矩}\label{fig:hw7-2-s}
\end{figure}

微段转动惯量为 $\rho J_p\dd x$，两侧内力矩之差为 $M_x\dd x=GJ_p\vartheta_{xx}\dd x$。两者相等，约去 $J_p\dd x$，得 $\vartheta_{tt}=(G/\rho)\vartheta_{xx}$。两端自由给出 $\vartheta_x(0,t)=\vartheta_x(L,t)=0$。
令 $\vartheta=X(x)\cos\omega t$，代入得 $X''+k^2X=0$，其中 $k=\omega\sqrt{\rho/G}$。对 $k>0$，$X=C\cos kx+D\sin kx$，$X'=-kC\sin kx+kD\cos kx$。左端 $X'(0)=0$ 给出 $D=0$；右端 $X'(L)=0$ 给出 $\sin kL=0$，故
\[
 \boxed{\omega_n=\frac{n\pi}{L}\sqrt{G/\rho},\quad n=1,2,\ldots}.
\]
对非零根，振型为$X_n=\cos(n\pi x/L)$，两端均是转角腹；截面力矩正比于$X_n'$，故两端力矩为零。$n=1$时两端转角反向，中间截面的转角为零。

零频情况须另解：空间方程 $X''=0$ 给出 $X=C+Dx$，两端 $X'=0$ 要求 $D=0$；时间方程 $\ddot q=0$ 依次积分为 $\dot q=B$、$q=A+Bt$。因此该模式转角处处相同，代表整体静止转角或匀速转动，不是有限周期振动。圆柱半径在弹性力矩和转动惯量中同时以 $J_p$ 出现，因此波速中消去。


总角动量可以区分整体转动与扭转振动。圆柱关于轴线的角动量为
\[
 L_z=\int_0^L\rho J_p\vartheta_t\dd x.
\]
两端外力矩为零，故 $L_z$ 守恒。对 $\vartheta=A+Bt$，有 $L_z=\rho J_pLB$；而每个 $n\ge1$ 模式的空间积分为
\[
\int_0^L\cos\frac{n\pi x}{L}\,\dd x
=\left[\frac{L}{n\pi}\sin\frac{n\pi x}{L}\right]_0^L=0,
\]
故不携带整体角动量。相应能量为
\[
 E=\frac12\int_0^L[\rho J_p\vartheta_t^2+GJ_p\vartheta_x^2]\dd x.
\]
整体匀速转动只有动能，扭转模态则在两项之间周期交换。图\ref{fig:hw7-2-s}中的$M=GJ_p\vartheta_x$是截面上的内力矩，不能理解为自由端仍受外加力矩。

\end{solution}



\begin{example}[细长弹性杆中的纵向声波]\label{ex:hw7-4}

一根均匀细长杆，密度$\rho$、杨氏模量$Y$、截面积$S$。横向尺寸远小于波长，横截面可自由发生泊松收缩，求纵向波的传播速度。

\end{example}


\begin{figure}[H]\centering
\begin{tikzpicture}[>=Stealth,line cap=round,line join=round]
\draw[curve](0,0)rectangle(4,.7);\draw[guide](1.7,0)--(1.7,.7);\draw[guide](2,0)--(2,.7);\draw[axis](1.7,1)--(2.7,1)node[right]{$u$};\node at(1.85,-.4){$\dd x$};\node at(3,1.6){$Y,\rho,S$};
\end{tikzpicture}
\caption{细长杆内的纵向位移与材料微元}\label{fig:hw7-4-q}
\end{figure}

\begin{solution}
沿图\ref{fig:hw7-4-q}的杆轴取材料微段，先由两端的位移差求纵向应变，再求截面内力之差。

\begin{figure}[H]\centering
\begin{tikzpicture}[>=Stealth,line cap=round,line join=round]
\draw[curve](0,0)rectangle(2,.7);\node at(1,.35){$\rho S\dd x$};
\draw[axis](0,.35)--(-1.2,.35)node[above]{$F(x)$};
\draw[axis](2,.35)--(3.4,.35)node[above]{$F(x+\dd x)$};
\draw[<->](0,-.4)--(2,-.4)node[midway,below]{$\dd x$};\node at(1,1.25){$F=YSu_x$};
\node at(1,-1.2){$F(x+\dd x)-F(x)=\rho S\dd x\,u_{tt}$};
\end{tikzpicture}
\caption{隔离材料微元，截面内力之差驱动纵向加速度}\label{fig:hw7-4-s}
\end{figure}

图\ref{fig:hw7-4-s}中的$x$标记平衡时的材料位置。两端位移差给出微段长度
\[
 \dd x+u(x+\dd x,t)-u(x,t)=(1+u_x)\dd x,
\]
故应变为$u_x$。应力正号表示拉伸；左侧截面对微段的力向左，右侧向右，合力为
\[
 [\sigma(x+\dd x,t)-\sigma(x,t)]S
 =YSu_{xx}\dd x.
\]
令合力等于微段质量$\rho S\dd x$乘以加速度$u_{tt}$，约去$S\dd x$，得$u_{tt}=(Y/\rho)u_{xx}$。代入行波$u=A\cos(kx-\omega t)$，按链式法则有
\[
u_x=-Ak\sin(kx-\omega t),\quad u_{xx}=-Ak^2\cos(kx-\omega t),
\]
\[
u_t=A\omega\sin(kx-\omega t),\quad u_{tt}=-A\omega^2\cos(kx-\omega t).
\]
比较波动方程两侧的系数，得$\omega^2=(Y/\rho)k^2$。取正频率、正波数支，$\omega=\sqrt{Y/\rho}\,k$，所以
\[
 \boxed{c_L=\frac\omega k=\sqrt{Y/\rho}}.
\]
群速度$\dd\omega/\dd k$也等于$c_L$，在这个一维长波近似下没有频散。截面积同时出现在回复力与微段质量中，约去后不影响波速。无限大弹性介质的纵波则用模量$K+4G/3$，其来由已由前文的横向零应变条件推得，速度为$\sqrt{(K+4G/3)/\rho}$；横向可自由收缩的细杆采用本题结果。

例如在自由杆端，轴向力为零，边界是$u_x=0$；若杆端被夹紧，边界则是$u=0$。求有限杆的驻波时，应先按图中的轴力确定端点条件，再用本题得到的波速。

\end{solution}



\begin{example}[空气中的绝热声速]\label{ex:hw7-5}

空气看成摩尔质量$M=29\,\mathrm{g\,mol^{-1}}$的双原子理想气体，$R_g=8.31\,\mathrm{J\,mol^{-1}K^{-1}}$，平衡温度$T_0=300\,\mathrm K$。用微元位移及绝热压缩推导声速并计算数值。

\end{example}


\begin{figure}[H]\centering
\begin{tikzpicture}[>=Stealth,line cap=round,line join=round]
\draw[thick](0,-1.1)--(5,-1.1);\draw[thick](0,1.1)--(5,1.1);\draw[guide](2,-1.1)--(2,1.1);\draw[guide](2.6,-1.1)--(2.6,1.1);\draw[axis](2,1.5)--(3,1.5)node[right]{$u$};\node at(1,-1.45){$p(x,t)$};\node at(4,-1.45){$p(x+\dd x,t)$};
\end{tikzpicture}
\caption{空气层位移与相邻截面的压强}\label{fig:hw7-5-q}
\end{figure}

\begin{solution}
图\ref{fig:hw7-5-s}中气层原体积 $S\dd x$，位移使相对体积改变量为 $u_x$。跟随同一气层写质量守恒：
\[
\rho_0S\dd x=(\rho_0+\delta\rho)S(1+u_x)\dd x.
\]
展开后略去二阶项 $\delta\rho\,u_x$，得 $\delta\rho=-\rho_0u_x$。绝热关系为 $p\rho^{-\gamma}=p_0\rho_0^{-\gamma}$，对乘积求一阶微分，
\[
\rho_0^{-\gamma}\delta p-\gamma p_0\rho_0^{-\gamma-1}\delta\rho=0,
\qquad \frac{\delta p}{p_0}=\gamma\frac{\delta\rho}{\rho_0}.
\]
因而
\[
 \delta p=\gamma p_0\frac{\delta\rho}{\rho_0}=-\gamma p_0u_x.
\]

\begin{figure}[H]\centering
\begin{tikzpicture}[>=Stealth,line cap=round,line join=round]
\fill[gray!10](0,-.5)rectangle(2,.5);\draw[guide](0,-.6)--(0,.6);\draw[guide](2,-.6)--(2,.6);
\draw[axis](0,.1)--(.85,.1);\node[above left]at(0,.3){$p(x,t)S$};
\draw[axis](2,-.1)--(1.15,-.1);\node[below right]at(2,-.3){$p(x+\dd x,t)S$};
\draw[axis](0,.95)--(.6,.95)node[above]{$u(x,t)$};\draw[axis](2,.95)--(2.7,.95)node[above]{$u(x+\dd x,t)$};
\draw[<->](0,-1.05)--(2,-1.05)node[midway,below]{$\dd x$};
\node at(1,-1.9){$\delta\rho=-\rho_0u_x,\quad \delta p=-\gamma p_0u_x$};
\end{tikzpicture}
\caption{同一气层两侧压力相向，位移差决定密度和压强扰动}\label{fig:hw7-5-s}
\end{figure}

气层两侧压强差产生合力$-S\delta p_x\dd x$，故$u_{tt}=(\gamma p_0/\rho_0)u_{xx}$。利用$p_0/\rho_0=R_gT_0/M$及双原子气体$\gamma=7/5$，得到
\[
 \boxed{c=\sqrt{\gamma R_gT_0/M}
 =\sqrt{\frac{(7/5)\times8.31\times300}{0.029}}
 \simeq347\,\mathrm{m\,s^{-1}}}.
\]
代入时，摩尔质量的单位取$\mathrm{kg\,mol^{-1}}$。声振动期间相邻气层间的热交换来不及充分进行，因此采用绝热关系。


图\ref{fig:hw7-5-s}中两侧压力都指向气层内部，所以合力为左侧压力减右侧压力，负号正由此而来。

在题设双原子理想气体近似中，平动和转动共五个有效自由度，摩尔定容热容为$5R_g/2$；由$C_p-C_V=R_g$得$C_p=7R_g/2$，所以$\gamma=C_p/C_V=7/5$。若误用等温关系，便会得到$\sqrt{R_gT_0/M}$，比本题绝热声速小$1/\sqrt\gamma$倍。

\end{solution}



\begin{example}[截面积突变处的纵波反射]\label{ex:sup7-reflection}
两根同材料、同密度的半无限细杆在$x=0$同轴连接，截面积分别为$S$和$4S$，接头无质量、无耗散。一列角频率$\omega$、位移振幅$A$的纵波从细杆射向粗杆。求位移反射系数、透射系数和能流比例，并求透射截面力振幅与入射截面力振幅的比值。
\begin{figure}[H]\centering
\begin{tikzpicture}
\draw[curve](-3,.28)rectangle(0,-.28);\draw[curve](0,.56)rectangle(3,-.56);
\draw[guide](0,-.9)--(0,1.5);\node[below]at(0,-.9){$x=0$};
\node at(-1.5,-.75){$S$};\node at(1.5,-.95){$4S$};
\draw[axis](-2.8,1)--(-.6,1)node[midway,above]{$A_i=A$};
\end{tikzpicture}
\caption{同材料杆在截面积突变处连接，入射方向向右}\label{fig:sup7-reflection-q}
\end{figure}
\end{example}
\begin{solution}
两杆波速相同$c=\sqrt{Y/\rho}$，但图\ref{fig:sup7-reflection-q}中的惯性与截面刚度均相差四倍，机械阻抗满足$Z_2=4Z_1$，其中$Z_1=\rho Sc$。将位移写成同一时间相位的复振幅，连接处的位移连续给出$A+A_r=A_t$。

采用共同时间因子 $e^{i\omega t}$，左、右两侧复位移分别为 $Ae^{-ikx}+A_re^{ikx}$、$A_te^{-ikx}$。空间求导分别得 $-ikAe^{-ikx}+ikA_re^{ikx}$、$-ikA_te^{-ikx}$；在 $x=0$ 令截面力 $YSu_x$ 连续，约去共同因子 $-i$，得到
\[
 YSk(A-A_r)=4YSkA_t.
\]
将 $A_t=A+A_r$ 代入力连续式，$A-A_r=4(A+A_r)$，所以 $-5A_r=3A$，进而
\[
 A_r=-\frac35A,\qquad A_t=\frac25A.
\]
\begin{figure}[H]\centering
\begin{tikzpicture}
\draw[guide](-3,0)--(3,0);\draw[guide](0,-1.3)--(0,1.5);
\draw[curve](-2.8,0)..controls(-2.5,1.2)and(-1.7,1.2)..(-1.4,0);
\draw[curve](-1.3,0)..controls(-1.1,-.72)and(-.5,-.72)..(-.3,0);
\draw[curve](.8,0)..controls(1,.48)and(1.6,.48)..(1.8,0);
\draw[axis](-2.7,1.25)--(-1.5,1.25)node[midway,above]{$A$};
\draw[axis](-.35,-1.05)--(-1.3,-1.05)node[midway,below]{$-3A/5$};
\draw[axis](.8,.85)--(1.9,.85)node[midway,above]{$2A/5$};
\node[below]at(2.5,0){$x$};
\end{tikzpicture}
\caption{位移波包的符号与传播方向示意；三列波分开画出以比较振幅}\label{fig:sup7-reflection-s}
\end{figure}
图\ref{fig:sup7-reflection-s}中反射位移反号。能流须同时计入截面积。以右行实波 $u=A\cos(kx-\omega t)$ 为例，$u_x=-kA\sin(kx-\omega t)$、$u_t=\omega A\sin(kx-\omega t)$，沿正 $x$ 方向的瞬时功率为 $-YSu_xu_t=Z\omega^2A^2\sin^2(kx-\omega t)$。一个周期 $T=2\pi/\omega$ 的平均值中，
\[
\frac1T\int_0^T\sin^2(kx-\omega t)\,\dd t
=\frac1T\left[\frac t2+\frac{\sin(2kx-2\omega t)}{4\omega}\right]_0^T=\frac12.
\]
因此单列波的平均能流大小为 $Z\omega^2|A|^2/2$，反射、透射与入射能流之比为
\[
 R_E=\frac9{25},\qquad
 T_E=\frac{Z_2}{Z_1}\left|\frac{A_t}A\right|^2=\frac{16}{25},
 \qquad R_E+T_E=1.
\]
单列行波的截面力振幅为$YSk|A|=Z\omega|A|$，所以$F_t/F_i=4\times2/5=8/5$。透射侧位移振幅减小，截面力振幅增大；它们共同给出的平均能流仍满足上述守恒关系。
\end{solution}



\begin{example}[端部质量改变弹性杆的频率]\label{ex:sup7-rodmass}
均匀细杆长$L$、截面积$S$、密度$\rho$、杨氏模量$Y$，左端固定，右端连接只能沿杆轴运动的集中质量$M>0$，忽略重力。求纵向固有频率方程，并在$M\gg\rho SL$时求最低频率的一阶杆质量修正。
\begin{figure}[H]\centering
\begin{tikzpicture}
\draw[curve](0,-.7)--(0,.7);\draw[curve](0,-.18)rectangle(4.4,.18);
\draw[curve](4.4,-.5)rectangle(5.2,.5);\node at(4.8,0){$M$};
\draw[<->](0,-.8)--(4.4,-.8)node[midway,below]{$L$};
\draw[axis](4.8,.9)--(6,.9)node[right]{$u(L,t)$};
\node[above]at(2,.25){$Y,S,\rho$};
\end{tikzpicture}
\caption{固定端弹性杆与集中端部质量}\label{fig:sup7-rodmass-q}
\end{figure}
\end{example}
\begin{solution}
杆内仍满足 $u_{tt}=c^2u_{xx}$，$c=\sqrt{Y/\rho}$。令 $u=X(x)\cos\omega t$，代入 $u_{tt}=-\omega^2X\cos\omega t$、$u_{xx}=X''\cos\omega t$，得 $X''+k^2X=0$、$k=\omega/c$。通解 $X=A\sin kx+B\cos kx$ 由左端 $X(0)=0$ 确定 $B=0$，因此 $X'=Ak\cos kx$。图\ref{fig:sup7-rodmass-q}右端质量受杆向左的截面力，故边界方程为
\[
 -YSu_x(L,t)=Mu_{tt}(L,t).
\]
代入模态得$YSk\cos kL=M\omega^2\sin kL$。令$z=kL$、$m_r=\rho SL$，频率方程为
\[
 \boxed{\cot z=\frac{M}{m_r}z},\qquad \omega=\frac cLz.
\]
\begin{figure}[H]\centering
\begin{tikzpicture}
\draw[axis](0,0)--(5,0)node[right]{$x$};\draw[axis](0,0)--(0,1.7)node[above]{$X$};
\draw[curve,domain=0:4.4,samples=80]plot(\x,{1.5*sin(65*\x/4.4)});
\draw[guide](0,0)--(4.4,{1.5*sin(65)});\fill(4.4,{1.5*sin(65)})circle(2pt)node[right]{$X(L)$};
\node[below]at(4.4,0){$L$};\node[above]at(2.5,1.55){$0<k_1L<\pi/2$};
\end{tikzpicture}
\caption{最低纵向模态与端部质量很大时的近似线性位移}\label{fig:sup7-rodmass-s}
\end{figure}
最低根位于 $0<z<\pi/2$。端部质量趋零时，它趋于固定--自由杆的 $\pi/2$；端部质量很大时，图\ref{fig:sup7-rodmass-s}的振型趋近直线。由正弦、余弦展开，
\[
\cot z=\frac{1-z^2/2+O(z^4)}{z(1-z^2/6+O(z^4))}
=\frac1z\bigl(1-z^2/2\bigr)\bigl(1+z^2/6\bigr)+O(z^3)
=\frac1z-\frac z3+O(z^3).
\]
代入频率方程并乘以 $z$，得到
\[
 1=\left(\frac M{m_r}+\frac13\right)z^2+O(z^4),
 \qquad \omega_1^2\simeq\frac{YS/L}{M+m_r/3}.
\]
写成相对杆质量的一阶修正，令 $\varepsilon=m_r/M\ll1$，则
\[
\omega_1^2=\frac{YS}{ML}\left[1-\frac\varepsilon3+O(\varepsilon^2)\right],\qquad
\omega_1=\sqrt{\frac{YS}{ML}}\left[1-\frac\varepsilon6+O(\varepsilon^2)\right].
\]
其中第二式用了 $(1+s)^{1/2}=1+s/2+O(s^2)$。杆身贡献 $m_r/3$ 的等效质量，也可令 $q(t)=u(L,t)$、$u(x,t)\simeq xq(t)/L$，直接计算杆的动能：
\[
K_r\simeq\frac12\int_0^L\rho S\frac{x^2}{L^2}\dot q^2\,\dd x
=\frac{\rho S\dot q^2}{2L^2}\left[\frac{x^3}{3}\right]_0^L
=\frac12\frac{m_r}{3}\dot q^2.
\]
零频率不是固有模态：左端固定后，静态齐次解 $X=Cx$ 在右端力为零时要求 $YSC=0$，只能为零解。
\end{solution}



\begin{example}[压缩脉冲在固定端和自由端的反射]\label{deep7:ex-pulse}
均匀细杆占据$x\le0$，截面积$S$、密度$\rho$、杨氏模量$Y$。一列向右传播的矩形压缩脉冲抵达端点$x=0$，入射粒子速度为$V>0$，持续时间$\tau$。令脉冲前沿到端点的时刻为$t=0$。分别讨论端点固定、端点自由时的速度和轴力反射，并求端部作用在杆上的总冲量、反射后的能量及波重叠区的最大应力或最大粒子速度。
\begin{figure}[H]\centering
\begin{tikzpicture}
\draw[curve](-5,-.18)rectangle(0,.18);\draw[curve](0,-.6)--(0,.6);
\draw[axis](-5,1.1)--(.7,1.1)node[right]{$x$};\node[below right]at(0,-.6){$0$};
\draw[curve](-4,1.1)--(-4,1.8)--(-2,1.8)--(-2,1.1);\node[above]at(-3,1.8){$v_i=V$};
\draw[axis](-3.6,.6)--(-2.3,.6);\draw[<->](-4,-.7)--(-2,-.7)node[midway,below]{$c\tau$};
\node[align=center]at(2,-.05){固定端或\\无载自由端};
\end{tikzpicture}
\caption{右行压缩脉冲，波的传播方向与入射粒子速度方向均向右}\label{deep7:fig-pulse-q}
\end{figure}
\end{example}
\begin{solution}
图\ref{deep7:fig-pulse-q}中的矩形波可写为$f(s)=1$当$0<s<\tau$、其余为零。忽略前沿的理想跃变点，入射波满足
\[
 v_i=Vf(t-x/c),\qquad N_i=-ZVf(t-x/c),
 \qquad Z=\rho Sc.
\]
负轴力对应压缩。反射波沿负$x$传播，故写成$v_r=sVf(t+x/c)$、$N_r=sZVf(t+x/c)$，而不能把速度与轴力的反射系数都取成同一个符号。

固定端满足$v(0,t)=0$，得到$s=-1$；自由端满足$N(0,t)=0$，得到$s=+1$。因此
\[
 \begin{array}{c|cc}
 &v_r/v_i\text{（对应波形）}&N_r/N_i\text{（对应波形）}\\\hline
 \text{固定端}&-1&+1\\
 \text{自由端}&+1&-1
 \end{array}
\]
\begin{figure}[H]\centering
\begin{tikzpicture}
\begin{scope}[shift={(-2.6,0)}]
\draw[axis](-1.8,0)--(.6,0)node[right]{$x$};\draw[curve](0,-.15)--(0,1.3);
\draw[curve](-1.5,.25)rectangle(-.05,.85);\node[above]at(-.9,1){固定端重叠区};
\draw[axis](-.1,-.45)--(-1.3,-.45)node[midway,below]{$F_{\rm end}=-2ZV$};
\node[align=center]at(-.7,-1.5){$v=0$\\$N=-2ZV$};
\end{scope}
\begin{scope}[shift={(2.5,0)}]
\draw[axis](-1.8,0)--(.6,0)node[right]{$x$};\draw[curve](-1.7,.25)--(0,.25);
\draw[axis](-1.4,.7)--(-.1,.7)node[midway,above]{$v=2V$};\node[align=center]at(-.7,-1.5){自由端重叠区\\$N=0$};
\end{scope}
\end{tikzpicture}
\caption{两个等幅脉冲重叠时，固定端附近应力加倍，自由端附近粒子速度加倍}\label{deep7:fig-pulse-s}
\end{figure}
图\ref{deep7:fig-pulse-s}的固定端在 $0<t<\tau$ 对杆施加向左的力 $-2ZV$，故 $J_{\rm fixed}=\int_0^\tau(-2ZV)\,\dd t=[-2ZVt]_0^\tau$；自由端外力恒为零。因此
\[
 J_{\rm fixed}=-2ZV\tau,\qquad J_{\rm free}=0.
\]
入射脉冲所带动量为$\rho Sc\tau V=ZV\tau$。固定端反射后粒子速度反向，动量变为$-ZV\tau$；自由端反射后粒子速度仍向右，动量保持$ZV\tau$。波包向左传播时，粒子速度仍可向右；动量方向由后者决定。

单列行波的单位长度动能为 $\rho SV^2/2$。其单位长度弹性能为 $N_i^2/(2YS)=Z^2V^2/(2YS)=\rho SV^2/2$，最后一步用了 $Z=\rho Sc$、$Y=\rho c^2$。波包长 $c\tau$，以沿波包的长度 $s$ 积分，入射总能量为
\[
 E_i=\int_0^{c\tau}\rho SV^2\,\dd s
 =\rho SV^2[s]_0^{c\tau}=ZV^2\tau.
\]
固定端速度为零，自由端外力为零，两者均不向外界传递边界功，因而$E_r=E_i$。在完全重叠区，固定端速度为零、应力为$-2ZV/S=-2\rho cV$；自由端轴力为零、速度为$2V$。两种情况下局部能量密度均为$2\rho SV^2$，恰为两列行波能量密度相加。应力或速度加倍并不意味着总能量变成四倍。
\end{solution}



\begin{example}[匀速推动有限长弹性杆]\label{deep7:ex-drive}
均匀细杆长$L$、截面积$S$、密度$\rho$、杨氏模量$Y$，初始静止且无应变，右端自由。理想执行器从$t=0$起使左端始终以速度$V>0$向右运动，允许执行器推杆或拉杆。忽略横向运动，$V/c\ll1$。求第一个$4L/c$内的端部驱动力、做功及反射过程中的杆内速度、轴力状态。
\begin{figure}[H]\centering
\begin{tikzpicture}
\draw[curve](0,-.6)--(0,.6);\draw[curve](0,-.2)rectangle(6,.2);\draw[axis](-1.6,0)--(-.15,0)node[midway,above]{$V$};
\draw[<->](0,-.8)--(6,-.8)node[midway,below]{$L$};\node[below left]at(-.2,-.25){执行器};\node[above]at(6,.3){自由端};
\draw[axis](.5,1.1)--(2,1.1)node[right]{$x$};
\end{tikzpicture}
\caption{左端位移规定为$u(0,t)=Vt$，端部力由波的往返反射确定}\label{deep7:fig-drive-q}
\end{figure}
\end{example}
\begin{solution}
令$t_0=L/c$、$Z=\rho Sc$。图\ref{deep7:fig-drive-q}的执行器先发出右行压缩波，波后$v=V$、$N=-ZV$，波前仍为静止未变形状态。对$0<t<t_0$，
\[
 (v,N)=\begin{cases}(V,-ZV),&0<x<ct,\\(0,0),&ct<x<L.\end{cases}
\]
在$t=t_0$，波抵达右端。自由端要求$N=0$，反射波增加$+ZV$的轴力和$+V$的速度。因此$t_0<t<2t_0$时，以反射前沿$x=2L-ct$为界，
\[
 (v,N)=\begin{cases}(V,-ZV),&0<x<2L-ct,\\(2V,0),&2L-ct<x<L.\end{cases}
\]
杆的自由端速度此时是执行器速度的两倍。
\begin{figure}[H]\centering
\begin{tikzpicture}[x=1.3cm,y=.8cm]
\draw[axis](0,0)--(3.9,0)node[right]{$x$};\draw[axis](0,0)--(0,6.5)node[above]{$t$};
\draw[guide](3,0)--(3,6);\draw[curve](0,0)--(3,1.5)--(0,3)--(3,4.5)--(0,6);
\foreach \y/\lab in {1.5/{t_0},3/{2t_0},4.5/{3t_0},6/{4t_0}}{\draw(-.06,\y)--(.06,\y);\node[left]at(-.08,\y){$\lab$};}
\node[below]at(3,0){$L$};
\node[align=center]at(.85,1.35){$V,-ZV$};\node[align=center]at(2.5,.4){$0,0$};\node[align=center]at(2.25,3){$2V,0$};
\node[align=center]at(.85,4.5){$V,+ZV$};\node[align=center]at(2.25,5.85){$0,0$};
\node[right,align=left]at(3.2,2){右端反射：$N=0$};\node[right,align=left]at(3.2,4.8){标注顺序为$(v,N)$};
\end{tikzpicture}
\caption{材料坐标中的波前往返轨迹，各区域标注粒子速度与轴力}\label{deep7:fig-drive-s}
\end{figure}
在$t=2t_0$，反射波抵达左端；若不修正，左端将有$v=2V$。为维持规定的$v=V$，执行器发出一列速度增量$-V$的右行波，其轴力增量为$+ZV$。因此$2t_0<t<3t_0$时，新前沿$x=c(t-2t_0)$左侧为$(V,+ZV)$，右侧为$(2V,0)$。

在$t=3t_0$，这列拉伸波到达自由端。反射后速度再减少$V$，轴力减小$ZV$。于是$3t_0<t<4t_0$时，前沿$x=4L-ct$左侧为$(V,+ZV)$，右侧为$(0,0)$。这四种状态和传播顺序见图\ref{deep7:fig-drive-s}。

左端外力与内轴力相反，$F=-N(0,t)$，故除理想波前到达的瞬间外，
\[
 F(t)=\begin{cases}ZV,&0<t<2t_0,\\-ZV,&2t_0<t<4t_0.\end{cases}
\]
前半周期执行器做正功，后半周期做负功。按力的两段取值积分：
\[
\begin{aligned}
W(2t_0)&=\int_0^{2t_0}ZV^2\,\dd t=ZV^2[t]_0^{2t_0},\\
W(4t_0)&=\int_0^{2t_0}ZV^2\,\dd t+\int_{2t_0}^{4t_0}(-ZV^2)\,\dd t
=2ZV^2t_0-2ZV^2t_0.
\end{aligned}
\]
于是
\[
 W(2t_0)=2ZV^2t_0=2\rho SLV^2,
 \qquad W(4t_0)=0.
\]
在$t=2t_0$反射前的极限状态，整杆速度为$2V$、应变为零，动能为$2\rho SLV^2$，与输入功一致。在$t=4t_0$反射前，杆又全部静止、无应变，但整体已向右平移$4VL/c$。执行器继续运动时，上述波过程重新开始，驱动力以$4L/c$为周期变化。

负驱动力是维持恒速边界所必需的拉力。若装置只能单侧推动，$2L/c$之后将发生脱离，题设边界条件必须改为自由端，后半过程便不再适用。
\end{solution}



\section{驻波、模态正交与初始条件}

有限介质的边界条件选出允许振型，初位移和初速度确定各振型的系数。存在端部集中质量时，该质量也进入模态正交的内积。



\subsection{初始条件与模态系数}
\label{sec:complete7-initial}
两端固定的均匀弦有$X_n=\sin(n\pi x/L)$、$\omega_n=n\pi c/L$。写
\[
 u(x,t)=\sum_{n=1}^\infty
 [A_n\cos\omega_nt+B_n\sin\omega_nt]X_n(x),
\]
利用$\int_0^L X_nX_m\dd x=L\delta_{nm}/2$，得
\[
 A_n=\frac2L\int_0^Lf_0(x)X_n(x)\dd x,\qquad
 B_n=\frac2{L\omega_n}\int_0^Lg_0(x)X_n(x)\dd x.
\]
初位移确定余弦项系数，初速度确定正弦项系数。若存在端部质量，模态正交内积也应包含该集中质量的惯性项；只对杆身积分会漏项。


\subsubsection{加权正交关系与初态投影}
一般的非均匀弦、杆满足 $-(pX_n')'=\omega_n^2wX_n$。将第 $n$ 式乘 $X_m$、第 $m$ 式乘 $X_n$ 后相减，并积分，左侧为
\[
\begin{aligned}
&\int_0^L[-X_m(pX_n')'+X_n(pX_m')']\dd x\\
&=[-X_mpX_n'+X_npX_m']_0^L
+\int_0^L(pX_m'X_n'-pX_n'X_m')\dd x.
\end{aligned}
\]
最后的积分两项相消，得到
\[(\omega_n^2-\omega_m^2)\int_0^LwX_nX_m\dd x
=[p(X_nX_m'-X_mX_n')]_0^L.
\]
对固定端、自由端或相同的齐次线性边界，右端为零。不同频率的模态因而在权$w$下正交，不能在质量密度不均匀时仍使用不带权的普通积分。令$N_n=\int_0^LwX_n^2\dd x$，并设$u(x,0)=f_0(x)$、$u_t(x,0)=g_0(x)$，展开式在$t=0$给
\[f_0=\sum_nA_nX_n,\qquad
g_0=\sum_n\omega_nB_nX_n.
\]
分别乘$wX_m$积分，除所求模态外其余项消去，故
\[A_m=\frac{\int_0^Lwf_0X_m\dd x}{N_m},\qquad
B_m=\frac{\int_0^Lwg_0X_m\dd x}{\omega_mN_m}.
\]
初始位移和初始速度由此分别确定各模态的系数。系统含有哪些振动频率由边值问题决定，各频率以多大振幅参与则取决于初态；基频最低并不保证它占主导。

若$x=L$处有集中质量$M_e$，其边界方程是$pX'(L)=M_e\omega^2X(L)$，边界项不能直接置零。代回前面的正交恒等式，整理为
\[(\omega_n^2-\omega_m^2)
\left[\int_0^LwX_nX_m\dd x+M_eX_n(L)X_m(L)\right]=0.
\]
因此应定义带端部惯性的内积
\[\langle f,g\rangle_M=\int_0^Lwfg\dd x+M_ef(L)g(L),
\]
并将前述投影式中的所有积分统一换成这个内积。分母、初位移和初速度的投影都要同时修改；只改频率方程而不改模态归一化，会算错端部质量参与的振幅。

若系统有零频模态，例如两端自由的杆可整体平移，该模态的时间函数为$A_0+B_0t$。其初速度系数直接投影$g_0$，不除以零频率。有限模态近似还应与初始形状的光滑性相配：初态含尖点时高阶模态衰减较慢，局部曲线的误差可能大于整体能量误差。


\subsection{例题与模型变式}

\begin{example}[两端固定弦的完整运动]\label{ex:hw7-6}

均匀细弦长$l$、质量$m$，两端固定，张力$T$保持不变。忽略阻力，给定小初位移$f(x)$及初速度$g(x)$，求弦的运动规律。

\end{example}


\begin{figure}[H]\centering
\begin{tikzpicture}[>=Stealth,line cap=round,line join=round]
\draw[thick](0,-.4)--(0,.5);\draw[thick](5,-.4)--(5,.5);\draw[curve,domain=0:5,samples=80]plot(\x,{.45*sin(36*\x)});\draw[guide](0,0)--(5,0);\draw[<->](0,-.7)--(5,-.7)node[midway,below]{$l$};
\end{tikzpicture}
\caption{两端固定弦的一种小位移}\label{fig:hw7-6-q}
\end{figure}

\begin{solution}
令线密度$\lambda=m/l$。图\ref{fig:hw7-6-s}中微段两端张力大小相同，竖直合力为$T[u_x(x+\dd x)-u_x(x)]$，所以
\[
 u_{tt}=c^2u_{xx},\quad c=\sqrt{T/\lambda}=\sqrt{Tl/m}.
\]

\begin{figure}[H]\centering
\begin{tikzpicture}[>=Stealth,line cap=round,line join=round]
\draw[curve,domain=0:2,samples=60]plot(\x,{.4-.2*(\x-1)^2});
\draw[axis](.5,.35)--(-.5,.15)node[left]{$-T\vct t_-$};
\draw[axis](1.5,.35)--(2.5,.15)node[right]{$T\vct t_+$};
\draw[guide](.5,-.25)--(.5,.35);\draw[guide](1.5,-.25)--(1.5,.35);
\draw[<->](.5,-.2)--(1.5,-.2)node[midway,below]{$\dd x$};
\node at(1,1){横向合力$=T u_{xx}\dd x$};\node at(1,-1){微元质量$=\lambda\dd x$};
\end{tikzpicture}
\caption{弦微元两端的张力大小相等，切线斜率不同}\label{fig:hw7-6-s}
\end{figure}

两端力在一阶近似中可写为$-T(1,u_x(x))$与$T(1,u_x(x+\dd x))$，水平分量抵消，竖直分量之差才产生横向加速度。使用恒定张力的条件是$|u_x|\ll1$且由弯曲带来的额外伸长可忽略。

对非零模态作分离变量 $u=X(x)q(t)$，代入并除以 $Xq$，得 $\ddot q/q=c^2X''/X=-\omega^2$。空间方程 $X''+k^2X=0$ 的通解是 $C\sin kx+D\cos kx$，$k=\omega/c$。固定端 $X(0)=0$ 给出 $D=0$，$X(l)=0$ 给出 $kl=n\pi$、$n\ge1$。时间方程 $\ddot q+\omega^2q=0$ 的两个特征根为 $\pm i\omega$，通解为 $A\cos\omega t+B\sin\omega t$。叠加全部模态，得到
\[
 \boxed{u(x,t)=\sum_{n=1}^\infty
 [A_n\cos\omega_nt+B_n\sin\omega_nt]\sin\frac{n\pi x}{l},
 \quad\omega_n=\frac{n\pi c}{l}}.
\]
确定初态系数前，先计算空间振型的内积。当 $n\ne j$ 时，用积化和差公式，
\[
\begin{aligned}
\int_0^l\sin\frac{n\pi x}{l}\sin\frac{j\pi x}{l}\,\dd x
&=\frac12\left[\frac{l\sin((n-j)\pi x/l)}{(n-j)\pi}
-\frac{l\sin((n+j)\pi x/l)}{(n+j)\pi}\right]_0^l=0.
\end{aligned}
\]
当 $n=j$ 时，$\sin^2\alpha=(1-\cos2\alpha)/2$，积分为 $[x/2-l\sin(2n\pi x/l)/(4n\pi)]_0^l=l/2$。因此，在 $t=0$ 代入级数及其时间导数，再乘第 $n$ 个振型积分，分别有
\[
\int_0^lf(x)\sin\frac{n\pi x}{l}\,\dd x=\frac l2A_n,\qquad
\int_0^lg(x)\sin\frac{n\pi x}{l}\,\dd x=\frac l2\omega_nB_n.
\]
解出系数，得到
\[
 A_n=\frac2l\int_0^lf(x)\sin\frac{n\pi x}{l}\dd x,\quad
 B_n=\frac2{l\omega_n}\int_0^lg(x)\sin\frac{n\pi x}{l}\dd x.
\]
固有频率$\omega_n$由弦和边界决定，展开系数由初态决定。给定的初位移、初速度须与固定端条件相容，上述级数才描述这一初态随时间的演化。


初速度展开中多出的 $1/\omega_n$，来自时间因子 $\sin\omega_nt$ 求导后产生的 $\omega_n$。

令$q_n=A_n\cos\omega_nt+B_n\sin\omega_nt$，图\ref{fig:hw7-6-s}对应的总能量可逐模态相加：
\[
 E=\frac12\int_0^l(\lambda u_t^2+Tu_x^2)\dd x
 =\frac{\lambda l}{4}\sum_{n=1}^\infty
 (\dot q_n^2+\omega_n^2q_n^2)
 =\frac{\lambda l}{4}\sum_{n=1}^\infty\omega_n^2(A_n^2+B_n^2).
\]
不同模式的交叉项因正交性消去。对有限能量初态，这个恒定值同时说明每个模式的幅度如何进入整根弦的机械能。

\end{solution}



\begin{example}[拨弦初态的谐波成分]\label{ex:sup7-pluck}
均匀弦长$L$、线密度$\lambda$，两端固定，张力$T$。将$x=a$处拨至横向位移$H$后静止释放，$0<a<L$，其余部分由两段直线连接端点，且$H\ll\min(a,L-a)$。求此后的振动；当$a=L/2$时，求基频模态所占初始能量比例。
\begin{figure}[H]\centering
\begin{tikzpicture}
\draw[guide](0,0)--(7,0);\draw[curve](0,0)--(2.5,1.3)--(7,0);
\foreach \x in {0,7}{\fill(\x,0)circle(2pt);\draw[curve](\x,-.3)--(\x,.3);}
\draw[guide](2.5,0)--(2.5,1.3);\draw[<->](2.8,0)--(2.8,1.3)node[midway,right]{$H$};
\node[below]at(0,0){$0$};\node[below]at(2.5,0){$a$};\node[below]at(7,0){$L$};
\draw[axis](.5,-.7)--(2,-.7)node[right]{$x$};
\end{tikzpicture}
\caption{释放前的三角形拨弦位移，两端位置固定}\label{fig:sup7-pluck-q}
\end{figure}
\end{example}
\begin{solution}
图\ref{fig:sup7-pluck-q}的初位移为
\[
 f_0(x)=\begin{cases}Hx/a,&0\le x\le a,\\
 H(L-x)/(L-a),&a\le x\le L,\end{cases}\qquad g_0(x)=0.
\]
波速$c=\sqrt{T/\lambda}$，两端固定给出$\omega_n=n\pi c/L$。因初速度为零，只须展开余弦时间因子。令$k_n=n\pi/L$，对两段分别分部积分：
\[
 \begin{aligned}
 A_n&=\frac{2H}{L}\left[\frac1a\int_0^ax\sin k_nx\dd x
 +\frac1{L-a}\int_a^L(L-x)\sin k_nx\dd x\right].
 \end{aligned}
\]
其中两个分部积分的原函数及端点值分别为
\[
\begin{aligned}
\int_0^ax\sin k_nx\,\dd x
&=\left[-\frac{x\cos k_nx}{k_n}+\frac{\sin k_nx}{k_n^2}\right]_0^a
=-\frac{a\cos k_na}{k_n}+\frac{\sin k_na}{k_n^2},\\
\int_a^L(L-x)\sin k_nx\,\dd x
&=\left[\frac{(x-L)\cos k_nx}{k_n}-\frac{\sin k_nx}{k_n^2}\right]_a^L\\
&=\frac{(L-a)\cos k_na}{k_n}+\frac{\sin k_na}{k_n^2}.
\end{aligned}
\]
各除以 $a$、$L-a$ 后，两个余弦项相消；剩下的正弦项系数为 $(1/a+1/(L-a))/k_n^2=L/[k_n^2a(L-a)]$，所以
\[
A_n=\frac{2H\sin(k_na)}{k_n^2a(L-a)}
=\frac{2HL^2}{n^2\pi^2a(L-a)}\sin\frac{n\pi a}{L}.
\]
完整运动为
\[
 u(x,t)=\sum_{n=1}^{\infty}A_n\sin\frac{n\pi x}L\cos\omega_nt.
\]
\begin{figure}[H]\centering
\begin{tikzpicture}[x=.9cm,y=.95cm]
\draw[axis](0,0)--(7.4,0)node[right]{$x$};
\draw[curve,domain=0:7,samples=90]plot(\x,{.81*sin(180*\x/7)});
\draw[guide,domain=0:7,samples=120]plot(\x,{-.09*sin(540*\x/7)});
\draw[dotted,domain=0:7,samples=150]plot(\x,{.0324*sin(900*\x/7)});
\node[above]at(3.5,.9){$n=1$};\node[below,align=center]at(3.5,-.55){中点拨弦只含奇次模态\\$A_1=8H/\pi^2,\quad A_3=-A_1/9$};
\node[below]at(0,0){$0$};\node[below]at(7,0){$L$};
\end{tikzpicture}
\caption{中点拨弦的前三个非零模态，按其初始带符号振幅绘制}\label{fig:sup7-pluck-s}
\end{figure}
图\ref{fig:sup7-pluck-s}中的偶次模态在中点有节点，其投影为零。初始能量直接由弦的斜率求出：
\[
 E=\frac T2\left[\int_0^a\frac{H^2}{a^2}\,\dd x
 +\int_a^L\frac{H^2}{(L-a)^2}\,\dd x\right]
 =\frac{TH^2}{2}\left(\frac1a+\frac1{L-a}\right)
 =\frac{TH^2L}{2a(L-a)}.
\]
第$n$模态的能量为$E_n=\lambda L\omega_n^2A_n^2/4$。当$a=L/2$时，$E=2TH^2/L$，$A_1=8H/\pi^2$，因此
\[
 \frac{E_1}{E}=\frac8{\pi^2}.
\]
初始折角允许用能量有限的模态级数表示；折角处二阶空间导数不宜按普通光滑函数逐点解释，释放后它对应向两侧传播的斜率变化。
\end{solution}



\section{梁的弯曲、曲率与端部条件}

弯矩引起的曲率沿梁积累为转角和位移。截面内力的正方向统一后，由支承条件确定积分常数；位移条件与受力条件分别列写。



\subsection{弯曲、曲率与边界条件}
\label{sec:complete7-bending}
小弯曲时截面保持平面，距中性轴$z$的纤维纵向应变为$-zy''$。应力与$z$成正比，轴向合力为零，弯矩的大小为
\[
 M_b=YI|y''|,\qquad I=\int z^2\dd S.
\]
符号按弯矩和挠度的约定取定。矩形截面竖直厚度为$a$、水平宽度为$b$时，$I=ba^3/12$。弯曲刚度与厚度的三次方成正比，互换$a,b$对应绕另一截面轴弯曲。

固支端规定$y=y'=0$。自由端若没有集中力偶，弯矩为零；没有横向集中力时，剪力也为零。例\ref{ex:hw7-8}先由截面右侧的端力求弯矩，再积分两次，用固支端条件确定常数。对于分布载荷，应先积分其力矩分布，不能把总载荷集中在自由端。



\LectureNeedspace{4\baselineskip}


\subsection{例题与模型变式}

\begin{example}[悬臂尺的端点位移]\label{ex:hw7-8}

一根均匀尺的一端固定在墙内，露出长度$l$，杨氏模量$Y$，矩形截面宽$b$、竖直厚度$a$。自由端受到竖直力$F$，形变很小。求端点位移。

\end{example}


\begin{figure}[H]\centering
\begin{tikzpicture}[>=Stealth,line cap=round,line join=round]
\draw[fill=gray!15](-.4,-.5)rectangle(0,.5);\draw[curve](0,0)--(4,0);\draw[axis](4,0)--(4,1)node[above]{$F$};\draw[<->](0,-.5)--(4,-.5)node[midway,below]{$l$};\draw[curve](5,-.4)rectangle(6,.4);\node at(5.5,.8){$b$};\node at(6.4,0){$a$};
\end{tikzpicture}
\caption{矩形截面悬臂尺的端点受力}\label{fig:hw7-8-q}
\end{figure}

\begin{solution}
图\ref{fig:hw7-8-s}中中性层既不伸长也不缩短，距中性层$z$的纤维应变为$-z\kappa$，应力为$-Yz\kappa$。对截面应力矩积分，得
\[
 M=YI\kappa,\quad I=b\int_{-a/2}^{a/2}z^2\,\dd z
 =b\left[\frac{z^3}{3}\right]_{-a/2}^{a/2}=\frac{ba^3}{12}.
\]

\begin{figure}[H]\centering
\begin{tikzpicture}[>=Stealth,line cap=round,line join=round]
\draw[fill=gray!15](-.4,-.5)rectangle(0,.5);\draw[guide](0,0)--(4,0);\draw[curve,domain=0:4,samples=70]plot(\x,{.00625*(12*\x*\x-\x*\x*\x)});\draw[axis](4,.8)--(4,1.8)node[above]{$F$};\draw[<->](0,-.5)--(4,-.5)node[midway,below]{$l$};\draw[guide](2,-.3)--(2,.25);\node at(1.7,1.1){$M(x)=F(l-x)$};\draw[curve](6.2,-.4)rectangle(7.2,.4);\node at(6.7,.8){$b$};\node at(7.6,0){$a$};
\draw[guide](4,.8)--(4.7,.8);\draw[guide](4,0)--(4.7,0);\draw[<->](4.45,0)--(4.45,.8)node[midway,right]{$y(l)$};
\end{tikzpicture}
\caption{端力作用于变形后的自由端}\label{fig:hw7-8-s}
\end{figure}

从墙边量 $x$，截面弯矩为 $F(l-x)$。小挠度时 $\kappa\simeq y''$，因此 $YIy''=F(l-x)$。先积分一次，$YIy'=F(lx-x^2/2)+C_1$；用 $y'(0)=0$ 得 $C_1=0$。再积分，$YIy=F(lx^2/2-x^3/6)+C_2$；用 $y(0)=0$ 得 $C_2=0$。整理为
\[
 y(x)=\frac F{6YI}(3lx^2-x^3),\qquad
 \boxed{y(l)=\frac{Fl^3}{3YI}=\frac{4Fl^3}{Yba^3}}.
\]
自由端还满足 $YIy''(l)=0$，与端部没有外加力偶相符，而 $YI\dd^3y/\dd x^3=-F$ 对应恒定剪力。

位移与力同向，端点斜率为$Fl^2/(2YI)$，须远小于1，才与小挠度近似一致。厚度以三次方进入分母，尺翻转截面方向后，位移会有明显变化。


图\ref{fig:review7-stress}进一步画出截面上的应变与应力。中性层上方、下方的纤维分别缩短和伸长，轴向合力$\int_A\sigma\dd A=0$，但力矩不为零。对矩形截面，
\[
 |\sigma(x,z)|=\frac{F(l-x)|z|}{I}.
\]
最大应力位于固支端的最外层纤维，为$|\sigma|_{\max}=6Fl/(ba^2)$。挠度公式除要求端点斜率很小，还要求材料保持线弹性。

\begin{figure}[H]\centering
\begin{tikzpicture}[>=Stealth,line cap=round,line join=round]
\draw[axis](0,-1.4)--(0,1.5)node[above]{$z$};\draw[axis](-1.5,0)--(1.7,0)node[right]{$\sigma$};\draw[curve](-1,1.05)--(1,-1.05);\draw[guide](-1,1.05)--(0,1.05)node[right]{$a/2$};\draw[guide](1,-1.05)--(0,-1.05)node[left]{$-a/2$};\node at(-1.2,1.45){受压};\node at(1.2,-1.45){受拉};
\begin{scope}[shift={(4.7,0)}]\draw[curve](-.65,-1.05)rectangle(.65,1.05);\draw[guide](-.95,0)--(.95,0);\node[right]at(.95,0){中性层};\draw[<->](-.65,-1.5)--(.65,-1.5)node[midway,below]{$b$};\draw[<->](-1,-1.05)--(-1,1.05)node[midway,left]{$a$};\end{scope}
\end{tikzpicture}
\caption{向上弯曲截面的线性应力分布与矩形截面的中性层}\label{fig:review7-stress}
\end{figure}

用弹性储能也可求得端点位移。弯曲能为
\[
 U=\int_0^l\frac{M(x)^2}{2YI}\dd x
 =\frac{F^2}{2YI}\int_0^l(l-x)^2\,\dd x
 =\frac{F^2}{2YI}\left[-\frac{(l-x)^3}{3}\right]_0^l
 =\frac{F^2l^3}{6YI}.
\]
线性加载时外力做功为$Fy(l)/2$，故$y(l)=2U/F=Fl^3/(3YI)$，与逐次积分完全一致。这里的$y$按图\ref{fig:hw7-8-s}取向上为正，正弯矩使$y''>0$；力若反向，位移、曲率与上下层应力符号也同时反向。

\end{solution}



\section{有限夹层中的反射、透射与阻抗匹配}
\subsection{两次反射必须连同相位相加}
有限夹层有两个界面。波在其中往返传播，后界面返回的波与前界面反射波相干叠加，使总反射率随厚度和频率变化。计算时须保留每次传播的相位。

以下考虑细长杆的一维纵波，忽略截面附近的三维修正。材料$j$的线密度为$\mu_j=\rho_jS_j$，波速$c_j=\sqrt{E_j/\rho_j}$，力与速度的特征阻抗为
\[
 Z_j=\rho_jS_jc_j.
\]
取压缩力$P=-E_jS_j\,\partial u/\partial x$为正，质点速度$v=\partial u/\partial t$。向右行波满足$P=Z_jv$，向左行波满足$P=-Z_jv$。连接处无集中质量时，$v$和$P$连续。

采用 $\ee^{-i\omega t}$ 的时间约定，夹层内速度、压缩力可写为
\[
v(x)=A\ee^{ik_2x}+B\ee^{-ik_2x},\qquad
P(x)=Z_2(A\ee^{ik_2x}-B\ee^{-ik_2x}),\qquad k_2=\omega/c_2.
\]
在左端，$A+B=v(0)$、$A-B=P(0)/Z_2$，故 $A=[v(0)+P(0)/Z_2]/2$、$B=[v(0)-P(0)/Z_2]/2$。将这两个系数代入 $x=\ell$ 的表达式，并用 $\ee^{i\delta}+\ee^{-i\delta}=2\cos\delta$、$\ee^{i\delta}-\ee^{-i\delta}=2i\sin\delta$ 合并，得到
\[
 \begin{pmatrix}v(\ell)\\P(\ell)\end{pmatrix}
 =
 \begin{pmatrix}
 \cos\delta&\mathrm i\sin\delta/Z_2\\
 \mathrm i Z_2\sin\delta&\cos\delta
 \end{pmatrix}
 \begin{pmatrix}v(0)\\P(0)\end{pmatrix},
 \qquad\delta=\frac{\omega\ell}{c_2}.
\]
其行列式为 $\cos^2\delta-(i\sin\delta/Z_2)(iZ_2\sin\delta)=1$，可直接求逆；无耗散时相向能流的差保持不变。
\ExpansionDiagram{fig:dep-d7-layer}{夹层中的左右行波以及两端连续条件}{
\fill[gray!15](2,-.5)rectangle(5,1.2);
\draw[thick](0,-.5)--(7,-.5);\draw[thick](0,1.2)--(7,1.2);
\draw[dashed](2,-.8)--(2,1.6);\draw[dashed](5,-.8)--(5,1.6);
\node at(1,1.6){$Z_1$};\node at(3.5,1.6){$Z_2,c_2$};\node at(6,1.6){$Z_3$};
\draw[->](.2,.8)--(1.7,.8)node[midway,above]{$v_{\rm in}$};
\draw[->](1.7,-.05)--(.2,-.05)node[midway,below]{$v_{\rm r}$};
\draw[->](2.4,.7)--(4.5,.7);\draw[->](4.5,.05)--(2.4,.05);
\draw[->](5.3,.5)--(6.8,.5)node[midway,above]{$v_{\rm t}$};
\draw[<->](2,-1)--(5,-1)node[midway,below]{$\ell$};}
图\ref{fig:dep-d7-layer}右侧只有出射波，所以$P(\ell)=Z_3v(\ell)$。由矩阵反解，夹层左端的输入阻抗为
\[
 Z_{\rm in}=
 \frac{Z_3\cos\delta-\mathrm iZ_2\sin\delta}
 {\cos\delta-\mathrm i(Z_3/Z_2)\sin\delta}.
\]
令左侧速度反射振幅比为$r=v_{\rm r}/v_{\rm in}$。左端总速度为$(1+r)v_{\rm in}$、总压缩力为$Z_1(1-r)v_{\rm in}$，因而
\[
 r=\frac{Z_1-Z_{\rm in}}{Z_1+Z_{\rm in}}.
\]
相同材料内速度振幅与位移振幅具有相同反射系数，压缩力的反射系数则相反。

\begin{example}[四分之一波长匹配层]
两根半无限长杆的阻抗满足$Z_3=4Z_1$。在两杆之间接入长度$\ell$、阻抗$Z_2=2Z_1$的无耗散夹层。求能量透射率随频率的变化，并比较$\omega\ell/c_2=\pi/2$及$\pi$时的结果。两个连接面均理想连接。
\ExpansionDiagram{fig:dep-d7-match-q}{以中等阻抗连接两根不同阻抗的杆}{
\draw[thick](0,0)rectangle(2,.8);\draw[thick,fill=gray!15](2,-.2)rectangle(4,1);\draw[thick](4,-.4)rectangle(6,1.2);
\node at(1,.4){$Z_1$};\node at(3,.4){$2Z_1$};\node at(5,.4){$4Z_1$};
\draw[->](-.5,.4)--(-.05,.4);\draw[<->](2,-.8)--(4,-.8)node[midway,below]{$\ell$};}
\end{example}
\begin{solution}
图\ref{fig:dep-d7-match-q}中的粗细只表示阻抗分区，不另假定三个材料的波速相同。设左、右外杆中入射和透射速度振幅为$v_i,v_t$。用传递矩阵的逆矩阵，把右端$(v_t,Z_3v_t)$写回左端，得到
\[
 v(0)=v_t\!\left(\cos\delta-\mathrm i\frac{Z_3}{Z_2}\sin\delta\right),
 \quad
 P(0)=v_t(Z_3\cos\delta-\mathrm iZ_2\sin\delta).
\]
由$v(0)+P(0)/Z_1=2v_i$，有
\[
 \frac{v_t}{v_i}=
 \frac{2Z_1}
 {(Z_1+Z_3)\cos\delta-\mathrm i(Z_2+Z_1Z_3/Z_2)\sin\delta}.
\]
单向行波的平均能流为$Z|v|^2/2$，所以能量透射率还须乘以$Z_3/Z_1$：
\[
 \mathcal T=
 \frac{4Z_1Z_3}
 {(Z_1+Z_3)^2\cos^2\delta+
 (Z_2+Z_1Z_3/Z_2)^2\sin^2\delta}
 =\boxed{\frac{16}{25\cos^2\delta+16\sin^2\delta}}.
\]
\ExpansionDiagram{fig:dep-d7-match-a}{匹配层的透射峰及离开匹配频率后的反射}{
\draw[->](0,0)--(6.8,0)node[right]{$\delta$};\draw[->](0,0)--(0,2.8)node[above]{$\mathcal T$};
\draw[thick,domain=0:6.28318,samples=100]plot(\x,{2.4*16/(25*cos(\x r)^2+16*sin(\x r)^2)});
\draw[dashed](0,2.4)--(6.3,2.4);\node[left]at(0,2.4){$1$};
\node[left]at(0,1.536){$16/25$};\node[below]at(1.5708,0){$\pi/2$};\node[below]at(3.1416,0){$\pi$};\node[below]at(4.7124,0){$3\pi/2$};}
图\ref{fig:dep-d7-match-a}中$\delta=\pi/2$时，$Z_{\rm in}=Z_2^2/Z_3=Z_1$，因此$r=0$、$\mathcal T=1$。虽然每个界面分别都会反射，两界面的返回振幅在左端相消。$\delta=\pi$时，夹层矩阵为负单位阵，输入阻抗又等于$Z_3$，故$\mathcal T=16/25$，与直接连接相同。四分之一波长匹配只在离散频率严格成立。
\end{solution}

\section{弯曲波的色散与能量输运}
\subsection{从静态弯曲方程到动态方程}
设细梁的弯曲刚度为$EI$、线密度为$\mu$。在小挠度、忽略截面转动惯量和剪切变形的Euler--Bernoulli近似中，取截面弯矩$M=EIy_{xx}$，并约定相邻截面剪力之差$Q_x\dd x$为微段的横向合力。微段的力矩平衡给出$Q=-M_x$，平动方程给出$\mu\dd x\,y_{tt}=Q_x\dd x$。代入弯矩与曲率的关系，依次得到
\[
 Q=-EIy_{xxx},\qquad \mu y_{tt}=-EIy_{xxxx},
\]
即
\[
 \mu y_{tt}+EIy_{xxxx}=0.
\]
梁的动态方程含四阶空间导数，有限梁的两端各需两条边界条件。例如固支端要求$y=y_x=0$；自由端要求弯矩和剪力为零，即$y_{xx}=y_{xxx}=0$。

将 $y=a\cos(kx-\omega t)$ 代入方程。空间逐次求导为 $y_x=-ak\sin(kx-\omega t)$、$y_{xx}=-ak^2\cos(kx-\omega t)$、$y_{xxx}=ak^3\sin(kx-\omega t)$、$y_{xxxx}=ak^4\cos(kx-\omega t)$；时间两次求导为 $y_t=a\omega\sin(kx-\omega t)$、$y_{tt}=-a\omega^2\cos(kx-\omega t)$。比较余弦系数，$-\mu\omega^2+EIk^4=0$，所以
\[
 \omega^2=\frac{EI}{\mu}k^4,\qquad
 v_{\rm ph}=\frac{\omega}{k}=\sqrt{\frac{EI}{\mu}}\,k,\qquad
 v_{\rm g}=\frac{\dd\omega}{\dd k}=2v_{\rm ph}.
\]
这与理想弦的$\omega=ck$不同，波包中的不同波数传播速度不同。上述近似限于波长远大于梁的横向尺寸，不能把$k\to\infty$的结果外推为真实材料的无限传播速度。
\ExpansionDiagram{fig:dep-d7-dispersion}{弯曲波的色散曲线；切线斜率为群速度，割线斜率为相速度}{
\draw[->](0,0)--(5.4,0)node[right]{$k$};\draw[->](0,0)--(0,3.5)node[above]{$\omega$};
\draw[thick,domain=0:4.3,samples=60]plot(\x,{.17*\x*\x});
\draw[dashed](0,0)--(4.2,2.142);\draw(1.2,-.306)--(4.1,2.652);
\fill(3,1.53)circle(2pt)node[above left]{$(k_0,\omega_0)$};
\node[right]at(4.05,2.7){切线};\node[right]at(4.2,2.14){割线};}
\subsection{由能量守恒求平均能流}
单位长度的动能与弯曲势能相加，
\[
 e=\frac12\mu y_t^2+\frac12EIy_{xx}^2.
\]
利用运动方程消去$y_{tt}$，再对空间分部整理：
\begin{align*}
 \partial_te
 &=-EIy_ty_{xxxx}+EIy_{xx}y_{xxt}\\
 &=-\partial_x\!\left[EI(y_ty_{xxx}-y_{xt}y_{xx})\right].
\end{align*}
所以向右的瞬时能流为
\[
 J=EI(y_ty_{xxx}-y_{xt}y_{xx}).
\]
其中两项分别来自剪力和弯矩做功。对单色行波，$y_t=a\omega\sin(kx-\omega t)$、$y_{xxx}=ak^3\sin(kx-\omega t)$，而$y_{xx}=-ak^2\cos(kx-\omega t)$、$y_{xt}=a\omega k\cos(kx-\omega t)$，故
\[
 \langle J\rangle=EIa^2\omega k^3,\qquad
 \langle e\rangle=\frac12\mu a^2\omega^2,\qquad
 \frac{\langle J\rangle}{\langle e\rangle}=2\frac\omega k=v_{\rm g}.
\]
能量传播速度与图\ref{fig:dep-d7-dispersion}中的切线斜率一致。

\begin{example}[弯曲波包的传播]
一根足够长的均匀细梁，$EI=16\,\mathrm{N\,m^2}$、$\mu=1\,\mathrm{kg/m}$。中心波数$k_0=2\,\mathrm{m^{-1}}$的窄带小振幅波包向右传播。求载波频率、相速度、群速度，以及载波振幅$a=1\,\mathrm{mm}$时的平均能流。说明波包宽度是否能始终保持不变。
\ExpansionDiagram{fig:dep-d7-packet-q}{载波在缓变包络内振动，波包整体向右传播}{
\draw[->](-.2,0)--(7.2,0)node[right]{$x$};
\draw[thick,domain=0:7,samples=200]plot(\x,{exp(-.4*(\x-3.4)^2)*cos(540*\x)});
\draw[dashed,domain=0:7,samples=60]plot(\x,{exp(-.4*(\x-3.4)^2)});
\draw[dashed,domain=0:7,samples=60]plot(\x,{-exp(-.4*(\x-3.4)^2)});
\draw[->](4.7,1)--(6.3,1)node[midway,above]{$v_{\rm g}$};}
\end{example}
\begin{solution}
图\ref{fig:dep-d7-packet-q}中快速振动由中心波数决定。先算材料组合$\sqrt{EI/\mu}=4\,\mathrm{m^2/s}$，于是
\[
 \omega_0=4k_0^2=16\,\mathrm{s^{-1}},\qquad
 f_0=\frac{\omega_0}{2\pi}=\frac8\pi\,\mathrm{Hz}.
\]
由 $\omega_0/k_0=16/2$，相速度为 $8\,\mathrm{m/s}$。色散关系为 $\omega(k)=4k^2$，求导得 $\dd\omega/\dd k=4\cdot2k=8k$，在 $k_0=2\,\mathrm{m^{-1}}$ 处群速度为 $16\,\mathrm{m/s}$。图中的单个波峰相对包络向后移动，因而不能通过追踪某一个波峰来测量波包能量的传播速度。

把$a=10^{-3}\,\mathrm m$代入能流式：
\[
 \langle J\rangle
 =16(10^{-3})^2\cdot16\cdot2^3
 =2.048\times10^{-3}\,\mathrm W.
\]
另一方面，$\langle e\rangle=\mu a^2\omega_0^2/2
=1.28\times10^{-4}\,\mathrm{J/m}$，乘以$v_{\rm g}=16\,\mathrm{m/s}$也得到同一能流。
\ExpansionDiagram{fig:dep-d7-packet-a}{不同波数的群速度不同，初始窄带包络逐渐展宽}{
\draw[->](0,0)--(7.5,0)node[right]{$x$};
\draw[thick,domain=0:3.5,samples=70]plot(\x,{1.2*exp(-2*(\x-1.6)^2)});
\draw[thick,domain=3.1:7,samples=70]plot(\x,{.7*exp(-.65*(\x-5.2)^2)});
\node[above]at(1.6,1.25){初态};\node[above]at(5.2,.75){较晚时刻};
\draw[->](2.9,1.1)--(4.4,1.1);}
再求一次导数，$\dd^2\omega/\dd k^2=\dd(8k)/\dd k=8\,\mathrm{m^2/s}$。若波数宽度为 $\Delta k$，相应群速度宽度为
\[
 \Delta v_{\rm g}\simeq
 \left|\frac{\dd^2\omega}{\dd k^2}\right|\Delta k
 =8\Delta k.
\]
传播时间$t$后，不同成分的位置差增加约$8\Delta k\,t$；因此图\ref{fig:dep-d7-packet-a}中的展宽是色散引起的，并不表示机械能损失。只有在展宽远小于原包络宽度的时间范围内，才能把包络视为刚性平移。
\end{solution}

\section{分布质量对悬臂梁低频振动的影响}
\subsection{用试探振型建立能量近似}
梁自身的质量不能总压缩为一个位于端点的质点。设长度$L$的均匀悬臂梁端部另接集中质量$M$，梁的线密度为$\mu$。选择满足固支端几何条件、并把端点位移规范为$q(t)$的试探形状
\[
 y(x,t)=q(t)\phi(\xi),\qquad
 \xi=x/L,\qquad
 \phi(\xi)=\frac{\xi^2(3-\xi)}2.
\]
此形状来自端点静力作用的弯曲线，满足$\phi(0)=\phi'(0)=0$、$\phi(1)=1$及端点弯矩为零$\phi''(1)=0$。对带端质量的动态问题，自由端剪力由端质量的惯性力平衡，不能再额外令$\phi'''(1)=0$。

由弯曲势能和动能分别得到
\[
 U=\frac{EIq^2}{2L^3}\int_0^1[\phi''(\xi)]^2\dd\xi
 =\frac12\frac{3EI}{L^3}q^2,
\]
\[
 T=\frac12\left[M+\mu L\int_0^1\phi^2(\xi)\dd\xi\right]\dot q^2.
\]
对试探形状求导，$\phi'=3\xi-3\xi^2/2$、$\phi''=3-3\xi$，故势能所需的积分为
\[
\int_0^1(\phi'')^2\dd\xi=9\int_0^1(1-2\xi+\xi^2)\dd\xi
=9[\xi-\xi^2+\xi^3/3]_0^1=3.
\]
动能中的 $\phi^2=(9\xi^4-6\xi^5+\xi^6)/4$，积分为
\[
 \int_0^1\phi^2\dd\xi
 =\frac14\left[\frac95\xi^5-\xi^6+\frac17\xi^7\right]_0^1
 =\frac14\left(\frac95-1+\frac17\right)=\frac{33}{140}.
\]
单一广义坐标的拉格朗日方程因此为
\[
 \left(M+\frac{33}{140}\mu L\right)\ddot q+
 \frac{3EI}{L^3}q=0.
\]
这里的$33\mu L/140$是相对于端点位移定义的等效质量，来自不同截面的速度平方加权；它不是梁的质心位置，也不是把整根梁替换为端点质点。
\ExpansionDiagram{fig:dep-d7-rayleigh}{悬臂梁的试探振型和沿梁分布的惯性}{
\fill[pattern=north east lines](-.3,-.3)rectangle(0,1.8);\draw[thick](0,0)--(5,0);
\draw[thick,domain=0:5,samples=60]plot(\x,{1.2*(\x/5)^2*(3-\x/5)/2});
\filldraw[fill=white](5,1.2)circle(.17)node[right=4pt]{$M$};
\draw[<->](5.5,0)--(5.5,1.2)node[midway,right]{$q$};
\draw[<->](0,-.5)--(5,-.5)node[midway,below]{$L$};\node[above]at(2.2,.8){$\mu\,\dd x$};}
\subsection{近似频率的适用范围}
上述能量商给出
\[
 \omega_{\rm R}^2=
 \frac{3EI/L^3}{M+(33/140)\mu L}.
\]
线性、无预应力、无耗散的梁具有正定弹性能和质量内积。设真实振型为$X_n$，频率按$\omega_1\le\omega_2\le\cdots$排列。端质量也计入内积，记
\[
 N_n=\int_0^L\mu X_n^2\dd x+M X_n(L)^2>0.
\]
将满足固支端几何条件的试探形状展开为$\phi(x/L)=\sum_n a_nX_n(x)$。模态正交使动能、弯曲能中的交叉项同时消去，能量商因而成为
\[
 \omega_{\rm R}^2
 =\frac{\displaystyle\sum_n a_n^2N_n\omega_n^2}
 {\displaystyle\sum_n a_n^2N_n}.
\]
分母为正，且每项$\omega_n^2\ge\omega_1^2$，所以
\[
 \omega_{\rm R}^2-\omega_1^2
 =\frac{\displaystyle\sum_n a_n^2N_n(\omega_n^2-\omega_1^2)}
 {\displaystyle\sum_n a_n^2N_n}\ge0.
\]
单函数Rayleigh近似由此给出基频的上界；只有试探形状恰为真实第一振型时才取等号。这个结论并不保证高阶模态频率也被准确描述。

\begin{example}[何时可以忽略梁自身质量]
图\ref{fig:dep-d7-rayleigh}中的梁采用上述试探振型。若$\mu L=M$，求计入梁质量后的近似频率与忽略梁质量所得频率的比值。若要求忽略梁质量造成的频率高估不超过$5\%$，端质量至少应为梁质量的多少倍？
\ExpansionDiagram{fig:dep-d7-mass-q}{端点集中质量与梁自身分布质量共同参与振动}{
\fill[pattern=north east lines](-.3,-.3)rectangle(0,1);
\draw[very thick](0,.35)--(5,.35);\filldraw[fill=white](5,.35)circle(.2)node[right=5pt]{$M$};
\draw[<->](0,-.4)--(5,-.4)node[midway,below]{梁质量$\mu L$};}
\end{example}
\begin{solution}
图\ref{fig:dep-d7-mass-q}中端质量的速度为$\dot q$，而梁上截面的速度为$\phi(x/L)\dot q$。因此忽略梁质量时$\omega_0=\sqrt{3EI/(ML^3)}$，计入后
\[
 \frac{\omega_{\rm R}}{\omega_0}
 =\left(1+\frac{33\mu L}{140M}\right)^{-1/2}.
\]
当$\mu L=M$，频率比为$\sqrt{140/173}\simeq0.8996$。忽略梁质量相对于$\omega_{\rm R}$高估约$11.16\%$，这与相对于$\omega_0$计算的$10.04\%$不是同一个百分数。
\ExpansionDiagram{fig:dep-d7-mass-a}{梁的质量贡献由各截面速度平方加权}{
\draw[->](0,0)--(5.8,0)node[right]{$x/L$};\draw[->](0,0)--(0,2.5)node[above]{$\phi^2$};
\fill[gray!15,domain=0:5,samples=70](0,0)--plot(\x,{2*((\x/5)^2*(3-\x/5)/2)^2})--(5,0)--cycle;
\draw[thick,domain=0:5,samples=70]plot(\x,{2*((\x/5)^2*(3-\x/5)/2)^2});
\node[below]at(5,0){$1$};\node[left]at(0,2){$1$};\node at(4.1,.3){$33/140$};}
以计入梁质量的近似频率为基准，要求
\[
 \frac{\omega_0-\omega_{\rm R}}{\omega_{\rm R}}
 =\sqrt{1+\frac{33\mu L}{140M}}-1\leq0.05.
\]
两边先加1再平方，得$33\mu L/(140M)\leq1.05^2-1=0.1025$，故
\[
 \boxed{\frac{M}{\mu L}\geq
 \frac{33}{140(0.1025)}\simeq2.300.}
\]
这个阈值衡量的是在同一试探振型模型中删去梁质量的影响，尚未包含试探振型相对于精确基频的误差。图\ref{fig:dep-d7-mass-a}中的权重解释了为何有效质量小于总质量：靠近固支端的截面速度很小。
\end{solution}

\section{习题与解答}

\begin{exercise}[端部转盘与扭转杆]\label{ex:sup7-torsion}
均匀圆杆长$L$、密度$\rho$、剪切模量$G$，截面极二次矩为$J_p$。左端固定，右端固接绕杆轴转动惯量为$I_d$的薄转盘。求小扭转的频率方程，并求$I_d\gg\rho J_pL$时的最低频率及杆身的等效转动惯量。
\begin{figure}[H]\centering
\begin{tikzpicture}
\draw[curve](0,-.8)--(0,.8);\draw[curve](0,-.15)rectangle(4.1,.15);
\draw[curve](4.15,0)ellipse(.25 and .9);\node[right]at(4.5,.5){$I_d$};
\draw[axis](4.15,0)--(5.6,0)node[right]{$x$};
\draw[<->](0,-1.1)--(4.15,-1.1)node[midway,below]{$L$};
\end{tikzpicture}
\caption{固支扭转杆带动端部转盘}\label{fig:sup7-torsion-q}
\end{figure}
\end{exercise}
\begin{solution}
图\ref{fig:sup7-torsion-q}中，截面力矩为 $GJ_p\theta_x$，单位长度转动惯量为 $\rho J_p$。微段两侧力矩之差为 $GJ_p\theta_{xx}\dd x$，等于 $\rho J_p\dd x\,\theta_{tt}$，得 $\theta_{tt}=(G/\rho)\theta_{xx}$。
取 $\theta=\Theta(x)\cos\omega t$，代入得 $\Theta''+k^2\Theta=0$、$\omega=k\sqrt{G/\rho}$。通解为 $\Theta=A\sin kx+B\cos kx$，固定端 $\Theta(0)=0$ 给出 $B=0$；因此 $\Theta'=Ak\cos kx$。

右端转盘由杆施加的力矩为 $-GJ_p\theta_x(L,t)$，所以 $-GJ_p\theta_x=I_d\theta_{tt}$。代入模态的导数，$GJ_pk\cos kL=I_d\omega^2\sin kL$。约去非零 $k$，再用 $\omega^2=(G/\rho)k^2$，得 $\cot kL=I_dk/(\rho J_p)$。以 $z=kL$ 表示，得到
\[
 \cot z=\frac{I_d}{\rho J_pL}z.
\]
\begin{figure}[H]\centering
\begin{tikzpicture}
\draw[axis](0,0)--(5.2,0)node[right]{$x$};\draw[axis](0,0)--(0,1.6)node[above]{$\Theta$};
\draw[curve](0,0)--(4.4,1.25);\draw[guide](4.4,0)--(4.4,1.25);
\node[right]at(4.4,1.25){$\Theta_L$};\node[below]at(4.4,0){$L$};
\node[above]at(2,1.3){$\Theta(x)\simeq x\Theta_L/L$};
\end{tikzpicture}
\caption{大转盘最低模态中的近似线性扭转角分布}\label{fig:sup7-torsion-s}
\end{figure}
图\ref{fig:sup7-torsion-s}的近似给出杆身动能
\[
 T_{\rm rod}=\frac{\rho J_p\dot\theta_L^2}{2L^2}\int_0^Lx^2\dd x
 =\frac{\rho J_p\dot\theta_L^2}{2L^2}\left[\frac{x^3}{3}\right]_0^L
 =\frac12\frac{\rho J_pL}{3}\dot\theta_L^2,
\]
线性转角分布给出 $\theta_x=\theta_L/L$，故弹性能为 $\tfrac12\int_0^LGJ_p(\theta_L/L)^2\dd x=GJ_p\theta_L^2/(2L)$。将转盘与杆身动能合并，运动方程是 $(I_d+\rho J_pL/3)\ddot\theta_L+(GJ_p/L)\theta_L=0$，因此最低频率为
\[
 \omega_1^2\simeq\frac{GJ_p/L}{I_d+\rho J_pL/3}.
\]
从精确频率方程也可得到这一修正：取 $\cot z=z^{-1}-z/3+O(z^3)$ 并乘以 $z$，得 $1=[I_d/(\rho J_pL)+1/3]z^2+O(z^4)$；再用 $\omega^2=(G/\rho)z^2/L^2$ 即得上式。杆的轴向转动惯量为 $\rho J_pL$。最低模态中，各截面角速度从固支端起逐渐增大，按速度平方积分得到其三分之一的等效量。
\end{solution}


\begin{exercise}[端力与均布载荷共同弯曲悬臂梁]\label{ex:sup7-beam}
均匀悬臂梁长$L$，截面为水平宽$b$、竖直厚$a$的矩形，杨氏模量$Y$。梁承受竖直向下的均布载荷$w$及自由端向下集中力$P$。在小挠度范围内求挠曲线、自由端挠度与转角，并求固支截面最大弯曲正应力。
\begin{figure}[H]\centering
\begin{tikzpicture}
\draw[curve](0,-.6)--(0,.7);\draw[curve](0,0)--(6,0);
\foreach \x in {.6,1.4,2.2,3,3.8,4.6,5.4}{\draw[axis](\x,.8)--(\x,.05);}
\node[above]at(3,.85){$w$};\draw[axis](6,0)--(6,-1.1)node[right]{$P$};
\draw[<->](0,-1.5)--(6,-1.5)node[midway,below]{$L$};
\end{tikzpicture}
\caption{悬臂梁的均布载荷和自由端集中力}\label{fig:sup7-beam-q}
\end{figure}
\end{exercise}
\begin{solution}
取向下挠度为$y>0$，$x$从固支端量起。图\ref{fig:sup7-beam-q}中，截面$x$右侧均布载荷的合力为$w(L-x)$，作用点距该截面为$(L-x)/2$，故
\[
 M_b(x)=P(L-x)+\frac w2(L-x)^2.
\]
在小挠度近似下采用$YIy''=M_b$，矩形截面$I=ba^3/12$。先积分一次：
\[
 YIy'(x)=P\left(Lx-\frac{x^2}{2}\right)
 +\frac w2\left(L^2x-Lx^2+\frac{x^3}{3}\right)+C_1.
\]
固支端的切线方向固定，$y'(0)=0$，所以$C_1=0$。再积分并整理各载荷项：
\[
 YIy(x)=\frac{Px^2(3L-x)}6
 +\frac{wx^2(6L^2-4Lx+x^2)}{24}+C_2.
\]
固支端没有位移，$y(0)=0$，给出$C_2=0$。上式除以$YI$即为全梁挠曲线。
\begin{figure}[H]\centering
\begin{tikzpicture}
\draw[guide](0,0)--(6,0);\draw[curve](0,-.5)--(0,.6);
\draw[curve,domain=0:6,samples=80]plot(\x,{-.9*(\x/6)^2*(3-\x/6)/2-.3*(\x/6)^2*(6-4*\x/6+(\x/6)^2)/3});
\draw[<->](6.4,0)--(6.4,-1.2)node[midway,right]{$y(L)$};
\draw[guide](4.8,-.85)--(6.5,-1.345833);\node[below]at(4.8,-1.6){自由端切线};
\draw[axis](0,-.05)--(0,-1.1)node[left]{$y$};
\end{tikzpicture}
\caption{向下挠曲线及自由端位移，图中变形放大示意}\label{fig:sup7-beam-s}
\end{figure}
在挠度式和斜率式中分别取$x=L$，得到图\ref{fig:sup7-beam-s}中的端点位移和切线小转角：
\[
 \begin{aligned}
 y(L)&=\frac{PL^3(3-1)}{6YI}
 +\frac{wL^4(6-4+1)}{24YI}
 =\boxed{\frac{PL^3}{3YI}+\frac{wL^4}{8YI}},\\
 y'(L)&=\frac{PL^2(1-1/2)}{YI}
 +\frac{wL^3(1-1+1/3)}{2YI}
 =\boxed{\frac{PL^2}{2YI}+\frac{wL^3}{6YI}}.
 \end{aligned}
\]
对于题设向下的$P,w\ge0$，$M_b'(x)=-P-w(L-x)\le0$，最大弯矩位于固支端，为$PL+wL^2/2$。弯曲正应力满足$|\sigma|=M_b|z|/I$，最外层纤维距中性轴$a/2$，因此
\[
 \boxed{|\sigma|_{\max}
 =\frac{(PL+wL^2/2)a}{2I}
 =\frac{6(PL+wL^2/2)}{ba^2}}.
\]
上述线性叠加要求材料处于线弹性范围且全梁$|y'|\ll1$。不能把均布载荷直接移到自由端：这样会改变它相对各截面的力臂，得到不同的弯矩分布和挠度。
\end{solution}


\begin{exercise}[有质量弹簧悬挂端部重物]\label{ex:sup7-heavyspring}
均匀弹簧自然长$L$、质量$m$、整根劲度系数$k$，上端固定、下端悬挂质量$M$的物块。弹簧沿竖直方向的小形变满足线弹性关系，重力加速度为$g$。求静态伸长分布和总伸长；再求平衡附近的纵向频率方程，并求$M\gg m$时的最低频率。
\begin{figure}[H]\centering
\begin{tikzpicture}
\draw[curve](-.8,0)--(.8,0);
\draw[curve,domain=0:1,samples=500]plot({.12*sin(4320*\x)},{-3*(1.4*\x-.2*\x*\x)/1.2});
\draw[curve](-.5,-3)rectangle(.5,-3.7);\node at(0,-3.35){$M$};
\draw[axis](1.4,0)--(1.4,-3.8)node[below]{$x$};
\node[right]at(.4,-1.4){$m,k,L$};
\end{tikzpicture}
\caption{弹簧的材料坐标从固定上端向下量起}\label{fig:sup7-heavyspring-q}
\end{figure}
\end{exercise}
\begin{solution}
图\ref{fig:sup7-heavyspring-q}的$x$为自然态材料坐标。长度$\dd x$弹簧段的劲度为$kL/\dd x$，所以截面张力为$kLu_0'(x)$。静态时该截面悬挂的总质量为$M+m(1-x/L)$，于是
\[
 kLu_0'(x)=[M+m(1-x/L)]g.
\]
由 $u_0(0)=0$，静态位移的积分为
\[
\begin{aligned}
u_0(x)&=\frac g{kL}\int_0^x\left(M+m-\frac{ms}{L}\right)\dd s\\
&=\frac g{kL}\left[(M+m)s-\frac{ms^2}{2L}\right]_0^x.
\end{aligned}
\]
故
\[
 u_0(x)=\frac{(M+m)gx}{kL}-\frac{mgx^2}{2kL^2},\qquad
 u_0(L)=\frac{Mg}{k}+\frac{mg}{2k}.
\]
\begin{figure}[H]\centering
\begin{tikzpicture}
\draw[axis](0,0)--(4.8,0)node[right]{$x$};\draw[axis](0,0)--(0,2.5)node[above]{张力};
\draw[curve](0,2.1)--(4.2,.8);\draw[guide](4.2,0)--(4.2,.8);
\node[left]at(0,2.1){$(M+m)g$};\node[right]at(4.2,.8){$Mg$};\node[below]at(4.2,0){$L$};
\end{tikzpicture}
\caption{静态张力沿弹簧向下线性减小，端部张力为重物重力}\label{fig:sup7-heavyspring-s}
\end{figure}
图\ref{fig:sup7-heavyspring-s}给出的张力对材料长度积分，再除以$kL$，也得到总伸长。对静态解叠加小位移$\eta(x,t)$，重力项与静态张力梯度抵消：
\[
 \frac mL\eta_{tt}=kL\eta_{xx},\qquad
 \eta(0,t)=0,\qquad -kL\eta_x(L,t)=M\eta_{tt}(L,t).
\]
令 $\eta=X(x)\cos\omega t$，代入波动方程得 $X''+q^2X=0$、$q^2=m\omega^2/(kL^2)$。通解 $X=A\sin qx+B\cos qx$ 由上端固定条件确定 $B=0$。下端条件代入 $\eta_x=Aq\cos qx\cos\omega t$、$\eta_{tt}=-A\omega^2\sin qx\cos\omega t$，得
\[
kLq\cos qL=M\omega^2\sin qL.
\]
令 $z=qL$，并用 $\omega^2=kz^2/m$ 约去非零 $z$，得到
\[
 \omega=z\sqrt{k/m},\qquad \cot z=\frac Mm z.
\]
最低模态在 $M\gg m$ 时对应 $z\ll1$。代入 $\cot z=1/z-z/3+O(z^3)$ 并乘以 $z$，得 $1=(M/m+1/3)z^2+O(z^4)$，所以 $\omega_1^2=kz^2/m\simeq k/(M+m/3)$。静态伸长中的弹簧质量以 $m/2$ 出现，最低动态等效质量以 $m/3$ 出现；前者来自静态张力的积分，后者来自振型速度平方的积分。
\end{solution}


\begin{exercise}[均布载荷下的简支梁]\label{deep7:ex-simplebeam}
长$L$的均匀梁两端简支，承受向下的均布载荷$w$。矩形截面水平宽$b$、竖直厚$a$，杨氏模量$Y$。梁的自重也计入$w$。采用小挠度梁模型，求支座反力、剪力、弯矩、挠曲线、最大挠度和最大弯曲正应力，并由弯矩积分求应变能。
\begin{figure}[H]\centering
\begin{tikzpicture}
\draw[curve](0,0)--(6,0);\draw[curve](-.2,-.35)--(0,0)--(.2,-.35)--cycle;\draw[curve](5.8,-.35)--(6,0)--(6.2,-.35)--cycle;\draw(5.84,-.45)circle(.08);\draw(6.16,-.45)circle(.08);
\foreach \x in {.4,1.2,2,2.8,3.6,4.4,5.2,5.8}{\draw[axis](\x,.8)--(\x,.05);}\node[above]at(3,.85){$w$};
\node[left]at(-.25,-.3){$A$};\node[right]at(6.35,-.3){$B$};\draw[<->](0,-1)--(6,-1)node[midway,below]{$L$};
\end{tikzpicture}
\caption{两端简支不限制端部转角，均布载荷遍布全梁}\label{deep7:fig-simplebeam-q}
\end{figure}
\end{exercise}
\begin{solution}
图\ref{deep7:fig-simplebeam-q}中的载荷关于中点对称，总力为$wL$，两支座反力均向上、大小为$wL/2$。从$A$沿梁取$x$，以左段竖直合力为剪力$V_s$、下缘受拉的弯矩为正，得
\[
 V_s(x)=w(L/2-x),\qquad M_b(x)=\frac w2x(L-x),
 \qquad M_b'=V_s.
\]
弯矩在两端为零，中点最大，为$wL^2/8$。

令向下挠度 $y$ 为正，此处正弯矩对应 $y''<0$，所以 $YIy''=-M_b$，其中 $I=ba^3/12$。将 $-M_b=-wLx/2+wx^2/2$ 逐项积分：
\[
YIy'=-\frac{wLx^2}{4}+\frac{wx^3}{6}+C_1,\qquad
YIy=-\frac{wLx^3}{12}+\frac{wx^4}{24}+C_1x+C_2.
\]
由 $y(0)=0$ 得 $C_2=0$；再代 $y(L)=0$，$-wL^4/12+wL^4/24+C_1L=0$，所以 $C_1=wL^3/24$。整理得
\[
 y'(x)=\frac{w}{24YI}(L^3-6Lx^2+4x^3),
 \qquad
 y(x)=\frac{wx}{24YI}(L^3-2Lx^2+x^3).
\]
\begin{figure}[H]\centering
\begin{tikzpicture}
\draw[guide](0,0)--(6,0);\draw[curve,domain=0:6,samples=90]plot(\x,{-3.84*(\x/6)*(1-2*(\x/6)^2+(\x/6)^3)});
\draw[guide](3,0)--(3,-1.2);\draw[<->](3.3,0)--(3.3,-1.2)node[midway,right]{$y_{\max}$};
\draw[guide](-.4,1.1)--(6.4,1.1);\draw[curve,domain=0:6,samples=70]plot(\x,{1.1+1.2*4*(\x/6)*(1-\x/6)});
\node[right]at(6.5,1.6){$M_b$图};\node[left]at(0,1.1){$0$};\node[above]at(3,2.3){$wL^2/8$};\node[below]at(3,-1.45){挠曲线（变形放大）};
\end{tikzpicture}
\caption{弯矩与挠度均在中点达到最大，两端斜率可以不为零}\label{deep7:fig-simplebeam-s}
\end{figure}
对$0<x<L/2$，上式$y'(x)>0$，中点为零；右半段由对称性斜率为负。因此图\ref{deep7:fig-simplebeam-s}的最大向下挠度为
\[
 y_{\max}=y(L/2)=\frac{5wL^4}{384YI}.
\]
最外层纤维距中性轴$a/2$，最大弯曲正应力的绝对值
\[
 |\sigma|_{\max}=\frac{(wL^2/8)(a/2)}I=\frac{3wL^2}{4ba^2}.
\]
按弯矩分布积分，弹性能为
\[
\begin{aligned}
U_e&=\int_0^L\frac{M_b^2}{2YI}\dd x
=\frac{w^2}{8YI}\int_0^L(L^2x^2-2Lx^3+x^4)\dd x\\
&=\frac{w^2}{8YI}\left[\frac{L^2x^3}{3}-\frac{Lx^4}{2}+\frac{x^5}{5}\right]_0^L
=\frac{w^2L^5}{8YI}\left(\frac13-\frac12+\frac15\right)
=\frac{w^2L^5}{240YI}.
\end{aligned}
\]
也可缓慢增加$w$，用$U_e=\tfrac12\int_0^Lwy(x)\dd x$得到相同结果。简支条件是端点挠度与端弯矩为零；若误加$y'(0)=y'(L)=0$，求出的将是另一种支承模型。
\end{solution}


\begin{exercise}[细杆屈曲与初始弯曲的放大（选做）]\label{deep7:ex-buckling}
长$L$的细直杆两端铰支，弯曲刚度为$YI$，沿两端连线施加大小相等的恒定轴向压力$P$。在小斜率模型中，用总势能求直杆失去稳定性的临界压力及各阶屈曲形状。若杆在无载荷时有微小初始弯曲$y_0(x)=a_0\sin(\pi x/L)$，再求低于临界压力时的静态挠度，并说明临界附近线性式的适用限制。
\begin{figure}[H]\centering
\begin{tikzpicture}
\draw[guide](0,0)--(6,0);\draw[curve,domain=0:6,samples=90]plot(\x,{.4*sin(180*\x/6)});
\draw[axis](-1.4,0)--(-.15,0)node[midway,above]{$P$};\draw[axis](7.4,0)--(6.15,0)node[midway,above]{$P$};
\fill(0,0)circle(2pt);\fill(6,0)circle(2pt);\node[below]at(0,0){铰支};\node[below]at(6,0){铰支};
\node[above]at(3,.45){微小初弯$y_0$};\draw[<->](0,-.7)--(6,-.7)node[midway,below]{$L$};
\end{tikzpicture}
\caption{轴向恒力沿端点连线作用，侧向形变会缩短两端间距}\label{deep7:fig-buckling-q}
\end{figure}
\end{exercise}
\begin{solution}
先取理想直杆。小斜率时，中心线弧长与两端距离之差为$\Delta\simeq\tfrac12\int_0^L(y')^2\dd x$。图\ref{deep7:fig-buckling-q}中的恒定压力对这一端缩做正功，故总势能的二次部分为
\[
 \Pi[y]=\frac12YI\int_0^L(y'')^2\dd x
 -\frac12P\int_0^L(y')^2\dd x.
\]
铰支端规定 $y(0)=y(L)=0$，端部不传力偶。对满足端位移条件的 $y$ 变分，两次分部积分为
\[
\begin{aligned}
\delta\Pi&=YI\int_0^Ly''\delta y''\,\dd x-P\int_0^Ly'\delta y'\,\dd x\\
&=[YIy''\delta y'-YIy'''\delta y-Py'\delta y]_0^L
+\int_0^L(YIy''''+Py'')\delta y\,\dd x.
\end{aligned}
\]
两端 $\delta y=0$，但 $\delta y'$ 任意，所以边界项要求 $y''(0)=y''(L)=0$；内部 $\delta y$ 任意，得到
\[
 YI y''''+Py''=0.
\]
当 $P>0$ 时，记 $q=\sqrt{P/(YI)}$。特征方程为 $s^2(s^2+q^2)=0$，空间通解为 $y=C_0+C_1x+C_2\cos qx+C_3\sin qx$。先由 $y''(0)=0$ 得 $C_2=0$，由 $y(0)=0$ 得 $C_0=0$；再由 $y''(L)=0$、$y(L)=0$，依次得 $C_3\sin qL=0$、$C_1L=0$。非零形状须有 $C_3\ne0$，所以 $qL=n\pi$，即 $y_n=A_n\sin(n\pi x/L)$，对应
\[
 P_n=YI\left(\frac{n\pi}L\right)^2,\qquad n=1,2,\ldots.
\]
\begin{figure}[H]\centering
\begin{tikzpicture}[x=.9cm]
\draw[guide](0,0)--(6,0);\draw[curve,domain=0:6,samples=90]plot(\x,{sin(180*\x/6)});\node[right]at(6.2,.5){$n=1$};
\begin{scope}[yshift=-2cm]\draw[guide](0,0)--(6,0);\draw[curve,domain=0:6,samples=90]plot(\x,{.7*sin(360*\x/6)});\node[right]at(6.2,.4){$n=2$};\end{scope}
\node[below]at(3,-3){$P_2=4P_1,\quad P_{\rm cr}=P_1$};
\end{tikzpicture}
\caption{两端铰支的前两阶形状，首先失稳的是无内部节点的第一阶}\label{deep7:fig-buckling-s}
\end{figure}
用正弦级数展开任意允许扰动，各阶正交使二次能量分开：
\[
 \Pi=\frac L4\sum_{n=1}^{\infty}
 \left[YI\left(\frac{n\pi}L\right)^4
 -P\left(\frac{n\pi}L\right)^2\right]A_n^2.
\]
所有系数为正的条件为$P<P_1$，因此图\ref{deep7:fig-buckling-s}的第一阶决定
\[
 \boxed{P_{\rm cr}=\pi^2YI/L^2}.
\]
在临界值处，线性模型中第一阶的二次能量为零；超过此值，该方向的二次能量为负。有限幅度的屈曲后形状和稳定性须保留高阶几何项，不能由二次能量确定。

对初弯杆，弯曲应变由曲率相对自然曲率的变化给出，所以弯曲能改为$\tfrac12YI\int(y''-y_0'')^2\dd x$。令$y=a\sin(\pi x/L)$，忽略与$a$无关的常数，有
\[
 \Pi(a)=\frac{YI\pi^4}{4L^3}(a-a_0)^2
 -\frac{P\pi^2}{4L}a^2.
\]
对 $a$ 求导，两个平方项分别产生 $2(a-a_0)$、$2a$，故
\[
\frac{\dd\Pi}{\dd a}=\frac{YI\pi^4}{2L^3}(a-a_0)-\frac{P\pi^2}{2L}a=0.
\]
除以 $YI\pi^4/(2L^3)$，得 $(1-P/P_{\rm cr})a=a_0$，所以
\[
 \boxed{a=\frac{a_0}{1-P/P_{\rm cr}}},\qquad 0\le P<P_{\rm cr}.
\]
有初弯时，杆在临界压力以下已有侧向挠度，并随压力增大而连续放大。公式所要求的小斜率为$\pi|a|/L\ll1$，即$\pi|a_0|/[L(1-P/P_{\rm cr})]\ll1$；在分母趋零之前，线性近似通常已经失效。
\end{solution}


\begin{exercise}[自重下的等强度悬杆（选做）]\label{deep7:ex-taper}
密度$\rho$、杨氏模量$Y$的竖直细杆自然长$L$，上端固定，下端悬挂质量$M$的重物。允许截面积$S(x)$沿自然态材料坐标$x$变化，$x$从上端向下量起。希望全杆工程拉应力均为给定的$\sigma_*>0$，且$\sigma_*/Y\ll1$。计入杆的自重，求截面积分布、杆质量、静态位移和弹性能；在所有应力不超过$\sigma_*$的同材料悬杆中，说明此设计为何具有最小质量。
\begin{figure}[H]\centering
\begin{tikzpicture}
\draw[curve](-1.2,0)--(1.2,0);
\draw[curve,domain=0:3.2,samples=60]plot({-.7*exp(-.35*\x)},{-\x});\draw[curve,domain=0:3.2,samples=60]plot({.7*exp(-.35*\x)},{-\x});
\draw[curve](-.6,-3.3)rectangle(.6,-4);\draw[curve](-.2284,-3.2)--(-.2284,-3.3);\draw[curve](.2284,-3.2)--(.2284,-3.3);\node at(0,-3.65){$M$};
\draw[axis](1.5,0)--(1.5,-3.4)node[below]{$x$};\draw[<->](-1.2,0)--(-1.2,-3.2)node[midway,left]{$L$};
\draw[guide](-.415,-1.5)--(.415,-1.5);\node[right]at(.6,-1.5){$S(x)$};
\end{tikzpicture}
\caption{上部截面承担更多自重，等强度设计由上向下连续变细；外形仅为示意}\label{deep7:fig-taper-q}
\end{figure}
\end{exercise}
\begin{solution}
图\ref{deep7:fig-taper-q}中的截面张力等于截面以下全部物质的重力，
\[
 N(x)=Mg+\rho g\int_x^LS(\xi)\dd\xi,
 \qquad N'(x)=-\rho gS(x),\quad N(L)=Mg.
\]
等强度条件为 $N=\sigma_*S$，代入微段平衡式得到 $S'=-(\rho g/\sigma_*)S$。记 $\alpha=\rho g/\sigma_*$，底部条件是 $S(L)=Mg/\sigma_*$。分离变量，并从下端积到任意截面：
\[
\int_{S(L)}^{S(x)}\frac{\dd s}{s}=-\alpha\int_L^x\dd\xi,\qquad
\ln\frac{S(x)}{S(L)}=-\alpha(x-L).
\]
两边取指数，得到
\[
 \boxed{S(x)=\frac{Mg}{\sigma_*}
 \exp\left[\frac{\rho g}{\sigma_*}(L-x)\right]}.
\]
截面积指数随$x$下降；若截面保持相似形状，横向线度与$\sqrt{S(x)}$成正比，外形曲线的指数是截面积指数的一半。
\begin{figure}[H]\centering
\begin{tikzpicture}
\draw[axis](0,0)--(5.2,0)node[right]{$x$};\draw[axis](0,0)--(0,3.2)node[above]{$N/(Mg)=S/S(L)$};
\draw[curve,domain=0:4.4,samples=70]plot(\x,{exp(1-.2272727*\x)});
\draw[guide](0,1)--(4.4,1)--(4.4,0);\node[below]at(4.4,0){$L$};\node[left]at(0,1){$1$};\node[left]at(0,2.718){$e^{\rho gL/\sigma_*}$};
\end{tikzpicture}
\caption{张力与截面积具有相同的指数分布，底端张力仍为$Mg$}\label{deep7:fig-taper-s}
\end{figure}
按图\ref{deep7:fig-taper-s}的截面积积分，杆质量为
\[
\begin{aligned}
m_{\rm rod}&=\frac{\rho Mg}{\sigma_*}\int_0^L\ee^{\alpha(L-x)}\dd x
=\frac{\rho Mg}{\sigma_*}\left[-\frac1\alpha\ee^{\alpha(L-x)}\right]_0^L\\
&=\frac{\rho Mg}{\alpha\sigma_*}(\ee^{\alpha L}-1)
=M(\ee^{\alpha L}-1).
\end{aligned}
\]
上端张力 $Mg e^{\alpha L}=(M+m_{\rm rod})g$，满足整体平衡。应变 $u'=N/(YS)=\sigma_*/Y$ 为常量，用 $u(0)=0$ 积分得 $u(x)=\int_0^x(\sigma_*/Y)\dd\xi=\sigma_*x/Y$。能量积分则可利用已经算出的总体积：
\[
U_e=\int_0^L\frac{\sigma_*^2S(x)}{2Y}\,\dd x
=\frac{\sigma_*^2}{2\rho Y}m_{\rm rod}.
\]
故
\[
 u(x)=\frac{\sigma_*}{Y}x,\qquad
 \Delta L=\frac{\sigma_*L}{Y},\qquad
 U_e=\int_0^L\frac{N^2}{2YS}\dd x
 =\frac{M\sigma_*^2}{2\rho Y}(e^{\alpha L}-1).
\]
以上以自然态坐标计算工程应变和单位原长自重，在小应变近似内一致。

对任意满足应力上限的设计，它在每处必须有$S(x)\ge N(x)/\sigma_*$。将平衡方程从下端向上写，以$z=L-x$为自变量，
\[
 \frac{\dd N}{\dd z}=\rho gS\ge\alpha N,
 \qquad N(0)=Mg.
\]
乘以积分因子 $\ee^{-\alpha z}>0$，按乘积求导，
\[
\frac{\dd}{\dd z}(N\ee^{-\alpha z})
=\ee^{-\alpha z}\left(\frac{\dd N}{\dd z}-\alpha N\right)\ge0.
\]
从 $0$ 到 $z$ 积分得 $N(z)\ee^{-\alpha z}-Mg\ge0$，所以 $N(z)\ge Mg\ee^{\alpha z}$，进而 $S(z)\ge Mg\ee^{\alpha z}/\sigma_*$。等强度设计在每一截面都取到这个下界，故总体积和质量也达到最小。该结论只比较纯轴向、小形变且不另受制造尺寸限制的悬杆。
\end{solution}


\chapter{天体运动进阶}



\section{二体运动、有效势与径向分支}

双星运动可以分成质心运动与相对运动，再用角动量把径向方程独立出来。有效势给出可达区域和圆轨道稳定性，零角动量的径向运动另行处理。



\subsection{二体问题与相对运动}
中心力的方向始终沿着质点与力心的连线。若力心固定于惯性系，且力的大小只依赖距离，便可写成
\[
 \vct F=F(r)\vct e_r=-\frac{\dd U}{\dd r}\vct e_r.
\]
约定径向向外为正，吸引力对应 $F(r)<0$。力矩 $\vct r\times\vct F=0$，角动量矢量守恒。角动量不为零时，$\vct r$ 与 $\vct v$ 始终位于垂直于角动量的固定平面内，以下就在这个平面中求运动。零角动量对应直线径向运动，需另列方程；涉及除以角动量的轨道公式不适用于这一退化情形。

\subsubsection{约化质量和引力参数}

两个质量分别为 $M,m$ 的质点只受彼此的引力。令
\[
 \vct r=\vct r_m-\vct r_M,\qquad
 \vct R_C=\frac{M\vct r_M+m\vct r_m}{M+m},\qquad
 \mu_{\rm r}=\frac{Mm}{M+m}.
\]
将两个牛顿方程相减，再按第三讲的质心动能分解式分离整体平动和相对运动，得到
\begin{align}
 \ddot{\vct R}_C&=0, &
 \mu_{\rm r}\ddot{\vct r}&=-\frac{GMm}{r^3}\vct r,\\
 T&=\frac12(M+m)\dot{\vct R}_C^{\,2}+\frac12\mu_{\rm r}\dot{\vct r}^{\,2}.
\end{align}
于是相对坐标满足
\begin{equation}\label{sup8:eq:d8-relative}
 \ddot{\vct r}=-\frac{\mu}{r^3}\vct r,\qquad \mu=G(M+m).
\end{equation}
本讲的记号中，$\mu$ 是引力参数，$\mu_{\rm r}$ 是约化质量。若 $m\ll M$，可忽略大质量天体的运动，取 $\mu\simeq GM$；此时 $\vct r$ 就是小物体相对于固定天体的位置。

\begin{figure}[H]\centering
\begin{tikzpicture}[scale=.95]
 \coordinate (M) at (-2,0);\coordinate (m) at (3,0);\coordinate (C) at (0,0);
 \draw[guide] (-3,0)--(4,0);
 \fill (M) circle (2.5pt) node[above left=3pt] {$M$};
 \fill (m) circle (2pt) node[below=5pt] {$m$};
 \fill (C) circle (1.4pt) node[below=5pt] {$C$};
 \draw[axis] (-1.9,.7)--(2.9,.7) node[midway,above] {$\vct r=\vct r_m-\vct r_M$};
 \draw[<->,>=Stealth] (-2,-1.35)--(0,-1.35) node[midway,below] {$\dfrac{m}{M+m}r$};
 \draw[<->,>=Stealth] (0,-1.35)--(3,-1.35) node[midway,below] {$\dfrac{M}{M+m}r$};
 \draw[axis] (M)--++(.9,0) node[above] {$\vct F_M$};
 \draw[axis] (m)--++(-.9,0) node[above] {$\vct F_m$};
 \draw[axis] (M)--++(0,-.8) node[left] {$\vct v_{M/C}$};
 \draw[axis] (m)--++(0,1.2) node[right] {$\vct v_{m/C}$};
\end{tikzpicture}
\caption{相对坐标与质心坐标。图示 $M>m$，相对运动的正转向为逆时针。}\label{sup8:fig:d8-two-body}
\end{figure}

由图\ref{sup8:fig:d8-two-body} 中质心 $C$ 的位置，两个物体相对于质心的坐标为
\[
 \vct r_{M/C}=-\frac{m}{M+m}\vct r,\qquad
 \vct r_{m/C}=\frac{M}{M+m}\vct r.
\]
它们的轨道是相对轨道按上述比例缩放、再分别取正反方向所得的曲线。两条轨道周期相同，但半长轴一般不同。

将相对运动的能量和角动量除以约化质量，定义
\begin{equation}\label{sup8:eq:d8-invariants}
 \varepsilon=\frac12v^2-\frac{\mu}{r},\qquad
 \vct h=\vct r\times\vct v,\qquad h=r^2\dot\varphi.
\end{equation}
取正转向使 $h>0$。相对运动的实际能量为 $E_{\rm rel}=\mu_{\rm r}\varepsilon$，实际角动量为 $\vct L_{\rm rel}=\mu_{\rm r}\vct h$。二体系统的总能量还包括质心动能。以下凡使用 $\varepsilon,h$，均按单位质量或单位约化质量定义。


\subsection{有效势、转向点与圆轨道}
\subsubsection{把径向运动单独列出}

极坐标下
\[
 v^2=\dot r^2+r^2\dot\varphi^2=\dot r^2+\frac{h^2}{r^2}.
\]
若单位质量势能为 $\Phi(r)$，则
\begin{equation}\label{sup8:eq:d8-effective}
 \frac12\dot r^2+\Phi_{\rm eff}(r)=\varepsilon,\qquad
 \Phi_{\rm eff}(r)=\Phi(r)+\frac{h^2}{2r^2}.
\end{equation}
$h^2/(2r^2)$ 是消去角变量后保留下来的横向动能项。它已包含在有效势中，若再计入完整动能，就会重复计算横向动能。

反平方引力的有效势为
\begin{equation}
 \Phi_{\rm eff}(r)=\frac{h^2}{2r^2}-\frac{\mu}{r}.
\end{equation}
对半径求导，
\[
\Phi_{\rm eff}'=-\frac{h^2}{r^3}+\frac\mu{r^2},\qquad
\Phi_{\rm eff}''=\frac{3h^2}{r^4}-\frac{2\mu}{r^3}.
\]
令一阶导数为零，得 $r_c=h^2/\mu$；在该处二阶导数为 $\mu/r_c^3>0$。对于 $h\ne0$，有效势在 $r\to0$ 时趋于正无穷，在 $r\to\infty$ 时从下方趋近零，并在
\[
 r_c=\frac{h^2}{\mu},\qquad
 \Phi_{\rm eff}(r_c)=-\frac{\mu^2}{2h^2}
\]
处取极小值。图\ref{sup8:fig:d8-effective} 中能量水平线与曲线的交点满足 $\dot r=0$，即径向转向点。

\begin{figure}[H]\centering
\begin{tikzpicture}[x=1.6cm,y=1.8cm]
 \draw[axis] (0,0)--(4.5,0) node[right] {$r/r_c$};
 \draw[axis] (0,-1.2)--(0,1.25) node[above] {$\Phi_{\rm eff}/(\mu/r_c)$};
 \draw[curve,domain=.4:4.2,samples=150] plot (\x,{.5/(\x*\x)-1/\x});
 \draw[guide] (.55,-.32)--(4.2,-.32) node[right] {$\varepsilon<0$};
 \draw[guide] (.4,.35)--(4.2,.35) node[right] {$\varepsilon>0$};
 \draw[guide] (1,0)--(1,-.5);
 \fill (1,-.5) circle (1.1pt) node[below=5pt] {圆轨道};
 \fill (.625,-.32) circle (1.1pt);\fill (2.5,-.32) circle (1.1pt);
 \node[above] at (.625,0) {$r_p$};\node[above] at (2.5,0) {$r_a$};
 \draw[<->,>=Stealth] (.65,-1.02)--(2.48,-1.02) node[midway,below] {束缚运动的径向区间};
\end{tikzpicture}
\caption{固定 $h$ 时的引力有效势。横轴取 $r_c=h^2/\mu$ 为单位，图中的 $r_p,r_a$ 表示相应的无量纲位置。}\label{sup8:fig:d8-effective}
\end{figure}

允许的能量满足 $\varepsilon\ge\min\Phi_{\rm eff}$。在理想点质量引力场中，负能量的非圆轨道有两个转向点，零能量和正能量轨道只有一个有限的最近点。若天体有实体半径 $R$，运动还受到表面限制：实际经过的半径须不小于 $R$，负能量轨道也可能撞击表面。

\subsubsection{圆轨道的稳定性}

一般中心势中的圆轨道由
\[
 \Phi_{\rm eff}'(r_c)=0
 \quad\Longleftrightarrow\quad
 h^2=r_c^3\Phi'(r_c)
\]
确定。一般有效势的二阶导数为 $\Phi_{\rm eff}''=\Phi''+3h^2/r^4$。令 $r=r_c+\xi$，在固定 $h$ 下将径向方程展开为
\[
\ddot\xi=-\Phi_{\rm eff}'(r_c)-\Phi_{\rm eff}''(r_c)\xi+O(\xi^2).
\]
常数项因圆轨道条件而为零，二阶导数中的 $h^2$ 用 $r_c^3\Phi'(r_c)$ 替换，得
\begin{equation}\label{sup8:eq:d8-stability}
 \ddot\xi+\omega_r^2\xi=0,\qquad
 \omega_r^2=\Phi''(r_c)+\frac{3\Phi'(r_c)}{r_c}.
\end{equation}
$\omega_r^2>0$ 时，小径向扰动作振动；$\omega_r^2<0$ 时，扰动指数增长。等于零时，二阶展开不能给出稳定结论，须回到原方程。

引力场中 $\Phi=-\mu/r$，所以
\[
 \omega_r^2=\frac{\mu}{r_c^3},\qquad
 \omega_\varphi=\frac{h}{r_c^2}=\sqrt{\frac{\mu}{r_c^3}}.
\]
径向振动一周与绕力心转动一周相合，这是近圆引力轨道没有近心点进动的结果。对于其他中心势，两个频率通常不同。


\subsection{例题与模型变式}

\begin{example}[质量不可忽略的双星]\label{sup8:problem-example-1}\label{sup8:ex:d8-e1}
两颗可视为质点的恒星质量分别为 $M,m$，系统质心静止。某时刻两星相距 $r_0$，相对速度垂直于连线，大小为
\[
 v_0=\sqrt{\frac{3G(M+m)}{2r_0}}.
\]
求相对轨道的半长轴、偏心率、最近和最远距离，以及两星各自绕质心的轨道尺度和共同周期。
\end{example}
\begin{figure}[H]\centering
\begin{tikzpicture}
 \draw[guide] (-3,0)--(3,0);
 \fill (-1.2,0) circle (2.3pt) node[above left=3pt] {$M$};
 \fill (1.8,0) circle (1.8pt) node[below=4pt] {$m$};
 \fill (0,0) circle (1.1pt) node[below=4pt] {$C$};
 \draw[<->,>=Stealth] (-1.2,.6)--(1.8,.6) node[midway,above] {$r_0$};
 \draw[axis] (-1.2,0)--(-1.2,-1) node[left] {$\vct v_M$};
 \draw[axis] (1.8,0)--(1.8,1.5) node[right] {$\vct v_m$};
 \node[align=left] at (4,0) {$\vct v_0=\vct v_m-\vct v_M$\\$\vct v_0\perp\vct r_0$};
\end{tikzpicture}
\caption{双星的初态。$C$ 静止，图示质量比仅用于示意。}\label{sup8:fig:d8-e1-q}
\end{figure}
\begin{solution}
图\ref{sup8:fig:d8-e1-q} 给出的是相对速度，而不是某一颗恒星的速度。取 $\mu=G(M+m)$，单位约化质量的两个守恒量为
\[
 \varepsilon=\frac{v_0^2}2-\frac{\mu}{r_0}=-\frac{\mu}{4r_0},\qquad
 h=r_0v_0,\quad h^2=\frac32\mu r_0.
\]
负能量说明相对轨道为椭圆。由 $\varepsilon=-\mu/(2a)$ 和式\eqref{sup8:eq:d8-eccentricity}，
\[
 a=2r_0,\qquad e^2=1-\frac34=\frac14,\qquad e=\frac12.
\]
初始径向速度为零。又因 $v_0>\sqrt{\mu/r_0}$，此时的径向加速度
\[
 \ddot r=\frac{v_0^2}{r_0}-\frac{\mu}{r_0^2}=\frac{\mu}{2r_0^2}>0,
\]
故初始点是近心点，不是远心点。因此
\[
 r_p=a(1-e)=r_0,\qquad r_a=a(1+e)=3r_0.
\]
\begin{figure}[H]\centering
\begin{tikzpicture}[scale=.9]
 \begin{scope}[xshift=-2.4cm,scale=.55]
  \draw[curve] (-1,0) ellipse (2 and 1.732);
  \fill (0,0) circle (2pt) node[below] {$O$};
  \fill (1,0) circle (2pt) node[right] {$P$};\fill (-3,0) circle (2pt) node[left] {$A$};
  \draw[axis] (1,0)--(1,1.3) node[right] {$v_0$};
  \node at (-1,-2.3) {相对轨道：$a=2r_0$};
 \end{scope}
 \begin{scope}[xshift=2.8cm,scale=1.15]
  \draw[curve] (-.6,0) ellipse (1.2 and 1.039);
  \draw[curve] (.4,0) ellipse (.8 and .693);
  \fill (0,0) circle (1.5pt) node[above=3pt] {$C$};
  \fill (.6,0) circle (2pt) node[above right] {$m$};
  \fill (-.4,0) circle (2.4pt) node[below left] {$M$};
  \draw[axis] (.6,0)--(.6,.8);\draw[axis] (-.4,0)--(-.4,-.533333);
  \node[align=center] at (0,-2.25) {质心系中的两条轨道\\按质量比分别缩放};
 \end{scope}
\end{tikzpicture}
\caption{左图的 $OP=r_0,OA=3r_0$；右图以同一焦点 $C$ 表示两颗恒星的质心轨道。}\label{sup8:fig:d8-e1-s}
\end{figure}
图\ref{sup8:fig:d8-e1-s} 右图由相对椭圆按质心关系得到，故
\[
 a_M=\frac{m}{M+m}\,2r_0,\qquad
 a_m=\frac{M}{M+m}\,2r_0.
\]
每颗恒星轨道的偏心率都是 $1/2$，近、远质心距离分别为本星半长轴的 $1/2,3/2$。两星始终位于质心两侧，对应的近心点方向相反；它们在同一时刻经过各自的近心点。

共同周期必须由相对半长轴计算：
\[
 T=2\pi\sqrt{\frac{(2r_0)^3}{G(M+m)}}.
\]
质心静止还给出 $\vct v_M=-m\vct v_0/(M+m)$、$\vct v_m=M\vct v_0/(M+m)$。题设的 $v_0$ 因而大于任一恒星的速率；上述能量与周期均使用相对速度计算。
\end{solution}



\begin{example}[逆平方力与逆立方力的轨道]\label{ex:hw8-3}

质量$m$的质点受到固定中心力。(1)$\vct F=-GMm\vct e_r/r^2$时求轨道，并说明何时为椭圆；(2)$\vct F=-\alpha m\vct e_r/r^3$、$\alpha>0$时分类求轨道。取非零比角动量$h$。

\end{example}


\begin{figure}[H]\centering
\begin{tikzpicture}[>=Stealth,line cap=round,line join=round]
\draw[axis](-.2,0)--(4.2,0)node[right]{$r$};\draw[axis](0,-1.05)--(0,1.65)node[above]{$V_{\rm eff}$};
\path[fill=gray!15](.612574,-.3)--plot[domain=.612574:2.720759,samples=80](\x,{.5/(\x*\x)-1/\x})--(2.720759,-.3)--cycle;
\draw[curve,domain=.35:4,samples=120]plot(\x,{.5/(\x*\x)-1/\x});\draw[guide](.35,-.3)--(4,-.3)node[right]{$\varepsilon$};
\draw[guide](.612574,0)--(.612574,-.8);\draw[guide](2.720759,0)--(2.720759,-.8);
\node[below]at(.612574,-.8){$r_p$};\node[below]at(2.720759,-.8){$r_a$};\node at(2.9,1){允许的径向区间};
\end{tikzpicture}
\caption{逆平方引力的比有效势；负能量线两交点给出近、远心距}\label{fig:hw8-3-q}
\end{figure}

\begin{solution}
逆平方力的比奈方程为 $u''+u=\mu/h^2$，这里 $u=1/r$、撇号表示对极角求导。取常数特解 $u_p=\mu/h^2$，齐次方程的特征根为 $\pm i$，故通解为
\[
u=\frac\mu{h^2}+A\cos\theta+B\sin\theta
=\frac\mu{h^2}[1+e\cos(\theta-\theta_p)],
\]
其中 $e=(h^2/\mu)\sqrt{A^2+B^2}\ge0$，$\theta_p$ 由 $A=(\mu e/h^2)\cos\theta_p$、$B=(\mu e/h^2)\sin\theta_p$ 确定。又 $\dot r=-hu'$、$r\dot\theta=hu$，能量为
\[
\begin{aligned}
\varepsilon&=\frac{h^2}{2}(u'^2+u^2)-\mu u\\
&=\frac{\mu^2}{2h^2}\left[e^2\sin^2(\theta-\theta_p)
+(1+e\cos(\theta-\theta_p))^2-2(1+e\cos(\theta-\theta_p))\right]\\
&=\frac{\mu^2}{2h^2}(e^2-1).
\end{aligned}
\]
解出偏心率并取 $r=1/u$，得到
\[
 \boxed{r=\frac{h^2/\mu}{1+e\cos(\theta-\theta_p)},\qquad
 e^2=1+\frac{2\varepsilon h^2}{\mu^2}}.
\]

\begin{figure}[H]\centering
\begin{tikzpicture}[>=Stealth,line cap=round,line join=round]
\begin{scope}
\draw[axis](0,0)--(2.7,0)node[right]{$\theta$};\draw[axis](0,0)--(0,1.9)node[above]{$u$};
\draw[curve,domain=0:2.4,samples=70]plot(\x,{1.5*cos(37.5*\x)});\node at(1.2,-.6){$h^2>\alpha$};\end{scope}
\begin{scope}[shift={(4,0)}]
\draw[axis](0,0)--(2.7,0)node[right]{$\theta$};\draw[axis](0,0)--(0,1.9)node[above]{$u$};
\draw[curve](0,1.5)--(2.4,0);\node at(1.2,-.6){$h^2=\alpha$};\end{scope}
\begin{scope}[shift={(8,0)}]
\draw[axis](0,0)--(2.7,0)node[right]{$\theta$};\draw[axis](0,0)--(0,1.9)node[above]{$u$};
\draw[curve,domain=0:2.5,samples=70]plot(\x,{1.5*exp(-\x)});\node at(1.2,-.6){$h^2<\alpha$};\end{scope}
\end{tikzpicture}
\caption{逆立方力三类轨道各取一支示意，实际运动只取正的倒数半径}\label{fig:hw8-3-s}
\end{figure}

由图\ref{fig:hw8-3-q}的有效势可见，$\varepsilon<0$对应$0\le e<1$，轨道为椭圆；零能量为抛物线，正能量为双曲线。因而椭圆轨道还要求能量为负。

固定$h$时，逆平方有效势的最低值为$-\mu^2/(2h^2)$。实际初态的能量不得更低；等号对应$e=0$的圆轨道，严格介于该值与零之间才是非圆椭圆。图中能量线与有效势的两个交点分别是$r_p,r_a$。

逆立方力则有 $u''+(1-\alpha/h^2)u=0$。特征方程 $s^2+1-\alpha/h^2=0$ 的根随 $h$ 改变：$h^2>\alpha$ 时为 $\pm ik$，其中 $k=\sqrt{1-\alpha/h^2}$；相等时为二重零根，$u''=0$ 依次积分得 $u'=B$、$u=A+B\theta$；$h^2<\alpha$ 时为 $\pm\nu$，其中 $\nu=\sqrt{\alpha/h^2-1}$。三类通解为
\[
 \boxed{\begin{array}{ll}
 h^2>\alpha:&u=A\cos(k\theta)+B\sin(k\theta),\\
 h^2=\alpha:&u=A+B\theta,\\
 h^2<\alpha:&u=Ae^{\nu\theta}+Be^{-\nu\theta},\quad
 \nu=\sqrt{\alpha/h^2-1}.
 \end{array}}
\]
逆立方力的比势能 $\Phi$ 满足 $\Phi'(r)=\alpha/r^3$。取无穷远为零点，积分为 $\Phi(r)=\int_\infty^r\alpha s^{-3}\dd s=[-\alpha/(2s^2)]_\infty^r=-\alpha/(2r^2)$，因而
\[
 V_{\rm eff}=\frac{h^2-\alpha}{2r^2},\qquad
 \varepsilon=\frac{h^2}{2}\left[(u')^2+
 \left(1-\frac\alpha{h^2}\right)u^2\right].
\]
该能量式同时约束各支解的积分常数。$h^2>\alpha$时有效势单调向无穷远降至零，无有限半径的稳定圆轨道；$h^2<\alpha$时向心塌缩可能出现，不能把指数解误认作椭圆。

各式均以$r=1/u>0$的角区间为实际轨道，如图\ref{fig:hw8-3-s}所示。临界情形$B=0,A>0$允许任意半径的圆轨道；它没有径向回复刚度。小于临界的指数解可形成趋向中心的螺旋段。上述解中的常数必须由初始半径、径向速度及角位置确定。


图\ref{fig:review8-potential}将逆立方力与第一问的势阱分开表示。$h^2>\alpha$时$V_{\rm eff}>0$且单调下降，有限运动必须有$\varepsilon>0$，只可能有一个最近距离；$h^2=\alpha$时离心项与势能恰好抵消，径向运动满足$\ddot r=0$；$h^2<\alpha$时有效势向中心降至负无穷，不出现阻止向心运动的势垒。

\begin{figure}[H]\centering
\begin{tikzpicture}[>=Stealth,line cap=round,line join=round]
\begin{scope}\draw[axis](0,0)--(2.6,0)node[right]{$r$};\draw[axis](0,-1.3)--(0,1.65)node[above]{$V_{\rm eff}$};\draw[curve,domain=.55:2.4,samples=60]plot(\x,{.42/(\x*\x)});\node at(1.25,-1.65){$h^2>\alpha$};\end{scope}
\begin{scope}[shift={(3.75,0)}]\draw[axis](0,0)--(2.6,0)node[right]{$r$};\draw[axis](0,-1.3)--(0,1.65);\draw[curve](.1,.025)--(2.4,.025);\node at(1.25,.55){$V_{\rm eff}=0$};\node at(1.25,-1.65){$h^2=\alpha$};\end{scope}
\begin{scope}[shift={(7.5,0)}]\draw[axis](0,0)--(2.6,0)node[right]{$r$};\draw[axis](0,-1.3)--(0,1.65);\draw[curve,domain=.58:2.4,samples=60]plot(\x,{-.42/(\x*\x)});\node at(1.25,-1.65){$h^2<\alpha$};\end{scope}
\end{tikzpicture}
\caption{逆立方力的比有效势为$(h^2-\alpha)/(2r^2)$；三种情形均没有稳定势阱}\label{fig:review8-potential}
\end{figure}

积分常数也可直接从初态求出。令初始角位置为零，$u_0=1/r_0$。由链式法则，$u'=(\dd u/\dd t)/\dot\theta=(-\dot r/r^2)/(h/r^2)=-\dot r/h$，故 $u_0'=-v_{r0}/h$。当 $h^2>\alpha$ 时，$u'=-kA\sin k\theta+kB\cos k\theta$，代入零角得 $A=u_0$、$B=u_0'/k$；临界时直接有 $A=u_0$、$B=u_0'$；当 $h^2<\alpha$ 时，零角条件为 $A+B=u_0$、$\nu(A-B)=u_0'$，相加、相减得 $A=(u_0+u_0'/\nu)/2$、$B=(u_0-u_0'/\nu)/2$。代回相应的解，并只保留 $u>0$ 的区间，即得到由指定初态出发的轨道。

\end{solution}



\begin{example}[从静止开始的径向落下]\label{sup8:problem-example-4}\label{sup8:ex:d8-e4}
半径为 $R$ 的球形天体在外部产生引力参数为 $\mu$ 的场。质点从距球心 $r_0>R$ 处由静止释放。求到达表面的速率和时间，并求理想点质量极限 $R\to0$ 下的落下时间。
\end{example}
\begin{figure}[H]\centering
\begin{tikzpicture}
 \draw[curve,fill=black!5] (0,0) circle (1);
 \fill (0,0) circle (1.5pt) node[below] {$O$};\fill (4,0) circle (1.8pt) node[above] {$P_0$};
 \draw[axis] (3.7,0)--(2.8,0) node[midway,above] {落下};
 \draw[<->,>=Stealth] (0,-1.35)--(4,-1.35) node[midway,below] {$r_0$};
 \draw[curve] (0,0)--(1,0) node[midway,above] {$R$};\node[right] at (4,.7) {$v_0=0$};
\end{tikzpicture}
\caption{零角动量的径向落下，运动在天体表面终止。}\label{sup8:fig:d8-e4-q}
\end{figure}
\begin{solution}
图\ref{sup8:fig:d8-e4-q} 中初速度为零，故 $h=0$。机械能守恒直接给出
\[
 \frac12\dot r^2-\frac\mu r=-\frac\mu{r_0},\qquad
 \dot r=-\sqrt{2\mu\left(\frac1r-\frac1{r_0}\right)}.
\]
负号对应由 $P_0$ 向 $O$ 靠近的分支。表面速率为
\[
 v_R=\sqrt{2\mu\left(\frac1R-\frac1{r_0}\right)}.
\]
飞行时间为
\[
 t_R=\int_R^{r_0}\frac{\dd r}{\sqrt{2\mu(1/r-1/r_0)}}.
\]
令 $r=r_0\cos^2\eta$。从释放点开始，$\eta$ 由 $0$ 增大到 $\eta_R=\arccos\sqrt{R/r_0}$，而
\[
 \dd r=-2r_0\sin\eta\cos\eta\,\dd\eta,\qquad
 \sqrt{2\mu(1/r-1/r_0)}=\sqrt{\frac{2\mu}{r_0}}\tan\eta.
\]
将微分与分母中的 $\tan\eta$ 约分，并反转积分上下限，得到
\[
\begin{aligned}
t_R&=\sqrt{\frac{2r_0^3}{\mu}}\int_0^{\eta_R}\cos^2\eta\,\dd\eta\\
&=\sqrt{\frac{2r_0^3}{\mu}}\left[\frac\eta2+\frac{\sin2\eta}{4}\right]_0^{\eta_R}.
\end{aligned}
\]
用 $\sin2\eta=2\sin\eta\cos\eta$ 整理，得
\begin{equation}
 \boxed{t_R=\sqrt{\frac{r_0^3}{2\mu}}
   \left(\eta_R+\sin\eta_R\cos\eta_R\right)}.
\end{equation}
\begin{figure}[H]\centering
\begin{tikzpicture}[x=4.7cm,y=2.7cm]
 \draw[axis] (0,0)--(1.8,0) node[right] {$\tau=t\sqrt{2\mu/r_0^3}$};
 \draw[axis] (0,0)--(0,1.18) node[above] {$r/r_0$};
 \draw[curve,domain=0:90,samples=120,variable=\s] plot ({\s*pi/180+sin(\s)*cos(\s)},{cos(\s)^2});
 \draw[guide] (0,.5)--(1.2854,.5)--(1.2854,0);
 \fill (1.2854,.5) circle (.9pt) node[above right] {$r=r_0/2$};
 \node[below] at (1.5708,0) {$\pi/2$};\node[left] at (0,1) {$1$};
\end{tikzpicture}
\caption{径向落下的时间曲线。图中 $r=r_0/2$ 对应 $\eta=\pi/4$，不是半个落下时间。}\label{sup8:fig:d8-e4-s}
\end{figure}
图\ref{sup8:fig:d8-e4-s} 中起点切线水平，符合由静止释放；越向内运动，径向速率越大。当 $R\to0$ 时，$\eta_R\to\pi/2$，得到
\[
 t_0=\frac\pi{2\sqrt2}\sqrt{\frac{r_0^3}{\mu}}.
\]
到 $r_0/2$ 的时间为 $\sqrt{r_0^3/(2\mu)}(\pi/4+1/2)$，大于 $t_0/2$。落下距离与时间并不成比例。真实天体表面以内的引力规律一般改变，不能将外部 $-\mu/r$ 势直接延伸到球心；本题的 $t_0$ 只属于点质量模型。
\end{solution}



\section{由初态确定轨道的形状与方向}

位置与速度共同确定轨道能量、轨道平面和近心点方向。径向与切向分量用于辨认圆锥曲线的分支，力心在椭圆中心的情形则作对照。



\subsection{从初态求圆锥曲线}
\subsubsection{轨道方程的推导}

令 $u=1/r$，撇号表示对 $\varphi$ 求导。由 $\dot\varphi=hu^2$，逐次使用链式法则，
\[
\dot r=-u^{-2}u'\dot\varphi=-hu',\qquad
\ddot r=-hu''\dot\varphi=-h^2u^2u''.
\]
同时 $r\dot\varphi^2=u^{-1}(hu^2)^2=h^2u^3$。
将其代入 $\ddot r-r\dot\varphi^2=F(r)/m$，得到 Binet 方程
\begin{equation}\label{sup8:eq:d8-binet}
 u''+u=-\frac{F(1/u)}{mh^2u^2}.
\end{equation}
若用 $f_r=F/m$ 表示单位质量的径向力，右边写作 $-f_r(1/u)/(h^2u^2)$。引力场中右边为常数 $\mu/h^2$。常特解为 $u_p=\mu/h^2$，齐次方程的特征根为 $\pm\mathrm i$，故以初始方位 $\varphi_0$ 为基准的实通解为
\[
u=\frac\mu{h^2}+A\cos(\varphi-\varphi_0)+B\sin(\varphi-\varphi_0).
\]
由 $r(\varphi_0)=r_0$、$\dot r(\varphi_0)=v_{r0}=-hu'(\varphi_0)$，确定 $A=1/r_0-\mu/h^2$、$B=-v_{r0}/h$。将两个三角项合成为 $e\mu\cos(\varphi-\varphi_p)/h^2$，其中 $e\ge0$，再取倒数，故
\begin{equation}\label{sup8:eq:d8-conic}
 r=\frac{p}{1+e\cos f},\qquad
 p=\frac{h^2}{\mu},\qquad f=\varphi-\varphi_p.
\end{equation}
$f$ 称真近点角，从近心点方向沿运动方向量起；$p$ 为半通径，$e\ge0$ 为偏心率。物理轨道只取 $r>0$ 的部分。对双曲线，不能将分母为负的极坐标分支也当作同一运动轨道。

由 $\dot r=-hu'$ 和 $v_\varphi=h/r$，有
\begin{equation}\label{sup8:eq:d8-velocities}
 v_r=\frac{\mu e}{h}\sin f,\qquad
 v_\varphi=\frac{\mu}{h}(1+e\cos f).
\end{equation}
代回能量式，
\[
\begin{aligned}
\varepsilon
&=\frac{\mu^2}{2h^2}\left[e^2\sin^2 f+(1+e\cos f)^2\right]
-\frac{\mu^2}{h^2}(1+e\cos f)\\
&=\frac{\mu^2}{2h^2}(e^2-1).
\end{aligned}
\]
一次余弦项相消，平方项用 $\sin^2 f+\cos^2 f=1$ 合并，得到
\begin{equation}\label{sup8:eq:d8-eccentricity}
 e^2=1+\frac{2\varepsilon h^2}{\mu^2}.
\end{equation}
在 $h\ne0$ 的引力问题中，$0\le e<1$ 对应椭圆，$e=1$ 对应抛物线，$e>1$ 对应双曲线。$e=0$ 是椭圆中的圆。径向直落虽然也使上式形式上得到 $e=1$，却不是通常所说的非退化抛物线。

\subsubsection{椭圆的几何量和速度}

当 $\varepsilon<0$ 时，最近、最远位置分别对应$f=0,\pi$，所以$r_p=p/(1+e)$、$r_a=p/(1-e)$。两者之和为长轴长，因而
\[
 a=\frac{r_p+r_a}{2}
 =\frac p2\left(\frac1{1+e}+\frac1{1-e}\right)
 =\frac p{1-e^2}.
\]
由椭圆几何关系$b_{\rm ell}^2=a^2-(ae)^2$，得到半短轴；再将$p=h^2/\mu=a(1-e^2)$代入式\eqref{sup8:eq:d8-eccentricity}，能量满足
\[
 \varepsilon=\frac{\mu^2(e^2-1)}{2h^2}
 =-\frac{\mu^2(1-e^2)}{2\mu a(1-e^2)}
 =-\frac\mu{2a}.
\]
这些关系汇集为
\begin{align}
 r_p&=a(1-e), &r_a&=a(1+e),\\
 p&=a(1-e^2), &b_{\rm ell}&=a\sqrt{1-e^2},\\
 \varepsilon&=-\frac{\mu}{2a}, &h^2&=\mu a(1-e^2).
\end{align}
这里 $b_{\rm ell}$ 是椭圆半短轴；后文散射问题的瞄准距离用 $b$，不能混淆。由能量立即得到活力公式
\begin{equation}\label{sup8:eq:d8-vis-viva}
 v^2=\mu\left(\frac2r-\frac1a\right).
\end{equation}
在近心点和远心点，速度纯沿切向，故 $r_pv_p=r_av_a=h$。其他位置一般有 $v_r\ne0$，不能把 $rv$ 当作 $h$。

\begin{figure}[H]\centering
\begin{tikzpicture}[scale=.85]
 \draw[curve] (-1.5,0) ellipse (3 and 2.598);
 \draw[guide] (-4.8,0)--(2,0);
 \coordinate (O) at (0,0);\coordinate (C) at (-1.5,0);\coordinate (P) at (-.3,2.381);
 \fill (O) circle (1.8pt) node[below right] {$O$};
 \fill (C) circle (1.2pt) node[below] {$C$};
 \fill (1.5,0) circle (1.5pt) node[right] {$P_p$};
 \fill (-4.5,0) circle (1.5pt) node[left] {$P_a$};
 \fill (P) circle (1.5pt) node[above right] {$P$};
 \draw[axis] (O)--(P) node[midway,right] {$r$};
 \draw (0.7,0) arc[start angle=0,end angle=97.18,radius=.7] node[midway,above right] {$f$};
 \draw[axis] (P)--++(-1.05,.423) node[above left] {$\vct v$};
 \draw[<->,>=Stealth] (-4.5,-.7)--(0,-.7) node[midway,below] {$r_a=a(1+e)$};
 \draw[<->,>=Stealth] (0,-.7)--(1.5,-.7) node[midway,below] {$r_p$};
 \draw[<->,>=Stealth] (-1.5,0)--(-1.5,2.598) node[midway,left] {$b_{\rm ell}$};
\end{tikzpicture}
\caption{力心位于椭圆焦点 $O$，$C$ 为几何中心；真近点角从 $OP_p$ 方向量起。}\label{sup8:fig:d8-ellipse}
\end{figure}

\subsubsection{由任意位置和速度确定轨道方向}

轨道的尺度和形状由 $a,e$ 确定，平面中的朝向还需由初速度方向确定。已知 $r_0,v_{r0},v_{\varphi0}$ 时，先算
\[
 h=r_0v_{\varphi0},\quad
 \varepsilon=\frac12(v_{r0}^2+v_{\varphi0}^2)-\frac{\mu}{r_0},
\]
再由
\begin{equation}\label{sup8:eq:d8-initial-angle}
 e\cos f_0=\frac{h^2}{\mu r_0}-1,\qquad
 e\sin f_0=\frac{hv_{r0}}{\mu}
\end{equation}
确定 $f_0$ 的象限。$v_{r0}>0$ 表示向外飞行，$v_{r0}<0$ 表示向内飞行；只求 $\arccos$ 会丢失这一分支信息。圆轨道没有唯一近心点，不能再要求唯一的 $f_0$。

也可定义偏心率矢量
\begin{equation}
 \vct e=\frac{\vct v\times\vct h}{\mu}-\vct e_r.
\end{equation}
该矢量在反平方引力下守恒，逐项求导的证明见第\ref{sec:complete8}节。它的模为$e$，非圆轨道中指向近心点，可用于求瞬时冲量前后轨道朝向的改变。


\subsection{例题与模型变式}

\begin{example}[带径向分量的初速度]\label{sup8:problem-example-2}\label{sup8:ex:d8-e2}
固定点质量产生引力参数为 $\mu$ 的场。质点位于图\ref{sup8:fig:d8-e2-q} 中的 $P_0=(0,r_0)$，初始径向速度向外，大小 $v_{r0}=\tfrac12\sqrt{\mu/r_0}$；横向速度沿 $-x$ 方向，大小 $v_{\varphi0}=\sqrt{\mu/r_0}$。求完整轨道方程、近远心距离及初始真近点角，并确定此后先经过哪个拱点。
\end{example}
\begin{figure}[H]\centering
\begin{tikzpicture}
 \draw[axis] (-2,0)--(2.3,0) node[right] {$x$};\draw[axis] (0,-.5)--(0,3.6) node[above] {$y$};
 \fill (0,0) circle (1.7pt) node[below right] {$O$};\fill (0,2) circle (1.7pt) node[right] {$P_0$};
 \draw[curve] (0,0)--(0,2) node[midway,right] {$r_0$};
 \draw[axis] (0,2)--(0,2.8) node[right] {$v_{r0}$};
 \draw[axis] (0,2)--(-1.6,2) node[left] {$v_{\varphi0}$};
 \draw[axis] (0,2)--(-1.6,2.8) node[above left] {$\vct v_0$};
\end{tikzpicture}
\caption{初速度同时含径向与横向分量，运动方向为逆时针。}\label{sup8:fig:d8-e2-q}
\end{figure}
\begin{solution}
以图\ref{sup8:fig:d8-e2-q} 的 $O$ 为力心。横向分量决定 $h$，两个分量共同决定能量：
\[
 h=\sqrt{\mu r_0},\qquad
 \varepsilon=\frac12\left(\frac14+1\right)\frac\mu{r_0}-\frac\mu{r_0}
 =-\frac{3\mu}{8r_0}.
\]
因此
\[
 p=r_0,\quad a=\frac{4r_0}3,\quad
 e=\sqrt{1+\frac{2\varepsilon h^2}{\mu^2}}=\frac12.
\]
式\eqref{sup8:eq:d8-initial-angle} 给出
\[
 e\cos f_0=0,\qquad e\sin f_0=\frac12.
\]
两式共同确定 $f_0=\pi/2$。既然初始半径指向 $+y$，近心点方向便指向 $+x$。用题图的极角 $\varphi$ 表示，$f=\varphi$，完整轨道为
\[
 \boxed{r=\frac{r_0}{1+\tfrac12\cos\varphi}}.
\]
\begin{figure}[H]\centering
\begin{tikzpicture}[scale=1.5]
 \draw[curve] (-.667,0) ellipse (1.333 and 1.155);
 \draw[axis] (-2.3,0)--(1,0) node[right] {$x$};\draw[guide] (0,-.2)--(0,1.4);
 \fill (0,0) circle (1.1pt) node[below right] {$O$};
 \fill (0,1) circle (1.1pt) node[right] {$P_0$};
 \fill (.667,0) circle (1.1pt) node[below right] {$P_p$};\fill (-2,0) circle (1.1pt) node[left] {$P_a$};
 \draw[axis] (0,1)--(-.65,1.325) node[above] {$\vct v_0$};
 \draw (0.35,0) arc (0:90:.35) node[midway,above right] {$f_0$};
 \draw[axis] plot[domain=104:135,samples=20] ({cos(\x)/(1+.5*cos(\x))},{sin(\x)/(1+.5*cos(\x))});
 \node at (-.65,-1.55) {$OP_p=2r_0/3,\quad OP_a=2r_0$};
\end{tikzpicture}
\caption{解出的轨道朝向。$P_0$ 处向外运动，沿上半椭圆先到远心点 $P_a$。}\label{sup8:fig:d8-e2-s}
\end{figure}
图\ref{sup8:fig:d8-e2-s} 中两个拱点的距离为
\[
 r_p=\frac{p}{1+e}=\frac{2r_0}3,\qquad
 r_a=\frac{p}{1-e}=2r_0.
\]
初始 $v_r>0$，所以 $f$ 从 $\pi/2$ 增大到 $\pi$，质点先到 $P_a$，再沿下半椭圆返回 $P_p$。在 $P_0$ 处，轨道的切向斜率为 $v_y/v_x=-1/2$，与速度箭头相符。初始速率 $v_0=\sqrt{5\mu/(4r_0)}$ 只确定能量；横向分量与径向分量还分别提供了确定 $p$、朝向和运动分支所需的信息。
\end{solution}



\begin{example}[切向发射的全部分支]\label{sup8:problem-example-3}\label{sup8:ex:d8-e3}
在理想点质量引力场中，质点于距力心 $r_0$ 处沿切向发射，速率为 $v_0=\beta\sqrt{\mu/r_0}$，其中 $\beta>0$。按 $\beta$ 的取值确定轨道种类、发射点是近心点还是远心点，并求有限的拱点距离。
\end{example}
\begin{figure}[H]\centering
\begin{tikzpicture}
 \fill (0,0) circle (2pt) node[below] {$O$};\fill (3,0) circle (1.7pt) node[below] {$P_0$};
 \draw[curve] (0,0)--(3,0) node[midway,below] {$r_0$};
 \draw[axis] (3,0)--(3,1.8) node[right] {$\beta\sqrt{\mu/r_0}$};
 \draw[guide] (0,0) circle (3);
 \node at (0,1.9) {参考圆轨道};
\end{tikzpicture}
\caption{同一发射位置，初速度纯切向；虚线圆仅作为 $\beta=1$ 的参照。}\label{sup8:fig:d8-e3-q}
\end{figure}
\begin{solution}
由图\ref{sup8:fig:d8-e3-q} 的垂直关系，$h=r_0v_0$，所以
\[
 \varepsilon=\frac{\mu}{2r_0}(\beta^2-2),\qquad
 p=\beta^2r_0,\qquad e=|\beta^2-1|.
\]
发射点的 $\dot r=0$，但它是哪一个拱点，须由
\[
 \ddot r(0)=\frac{\mu}{r_0^2}(\beta^2-1)
\]
判断。$\beta<1$ 时从发射点开始向内运动，$\beta>1$ 时开始向外运动。
\begin{figure}[H]\centering
\begin{tikzpicture}[scale=.75]
 \begin{scope}[xshift=-3.6cm]
  \draw[curve] (0,0) ellipse (1.5 and 1.2);\fill (-.9,0) circle (2pt) node[below] {$O$};
  \fill (1.5,0) circle (2pt) node[right] {$P_0=P_a$};\draw[axis] (1.5,0)--(1.5,1);
  \draw[axis] plot[domain=0:45,samples=20] ({1.5*cos(\x)},{1.2*sin(\x)});
  \node at (0,-1.8) {$0<\beta<1$};
 \end{scope}
 \begin{scope}[xshift=2.5cm]
  \draw[curve] (-.9,0) ellipse (1.5 and 1.2);\fill (0,0) circle (2pt) node[below] {$O$};
  \fill (.6,0) circle (2pt) node[right] {$P_0=P_p$};\draw[axis] (.6,0)--(.6,1);
  \draw[axis] plot[domain=0:45,samples=20] ({-.9+1.5*cos(\x)},{1.2*sin(\x)});
  \node at (-.7,-1.8) {$1<\beta<\sqrt2$};
 \end{scope}
 \node[align=center] at (0,-3) {$\beta=1$：圆；$\beta=\sqrt2$：抛物线；$\beta>\sqrt2$：双曲线};
\end{tikzpicture}
\caption{两种椭圆分支的发射点不同。抛物线和双曲线的切向发射点也是近心点。}\label{sup8:fig:d8-e3-s}
\end{figure}
对 $0<\beta<1$，$e=1-\beta^2$，故
\[
 r_a=r_0,\qquad r_p=\frac{\beta^2}{2-\beta^2}r_0,\qquad
 a=\frac{r_0}{2-\beta^2}.
\]
图\ref{sup8:fig:d8-e3-s} 左图的力心在发射点的远侧焦点，符合发射点为远心点的几何关系。对 $1<\beta<\sqrt2$，$e=\beta^2-1$，则
\[
 r_p=r_0,\qquad r_a=\frac{\beta^2}{2-\beta^2}r_0,
 \qquad a=\frac{r_0}{2-\beta^2}.
\]
$\beta=1$ 时 $e=0$，轨道为半径 $r_0$ 的圆，没有区分开的近、远心点。$\beta=\sqrt2$ 时 $\varepsilon=0,e=1$，轨道为近心距离 $r_0$ 的抛物线；$\beta>\sqrt2$ 时 $\varepsilon>0,e>1$，轨道为双曲线，最近距离仍为 $r_0$，均没有有限的远心点。

若中心天体有实体半径 $R<r_0$，第一种椭圆还须比较 $r_p$ 与 $R$。形式上闭合的椭圆可能被表面截断。题设排除了 $\beta=0$；该情形是径向直落，应另列径向方程。
\end{solution}



\begin{example}[垂直速度与二倍速率的最小偏心率]\label{ex:hw8-6}

彗星在同一条开普勒轨道上两点的速度相互垂直，且$|\vct v_1|=2|\vct v_2|$。求轨道偏心率的最小可能值。

\end{example}


\begin{figure}[H]\centering
\begin{tikzpicture}[>=Stealth,line cap=round,line join=round]
\draw[axis](0,0)--(3,0)node[right]{$v_x$};\draw[axis](0,0)--(0,3)node[above]{$v_y$};\coordinate(A)at(0,2);\coordinate(B)at(1,0);\fill(A)circle(2pt)node[above right]{$\vct v_1$};\fill(B)circle(2pt)node[below right]{$\vct v_2$};\draw[curve](-.5,.5)circle(1.581);\fill(-.5,.5)circle(1pt)node[left]{$H$};
\draw[axis](0,0)--(0,2);\draw[axis](0,0)--(1,0);\fill(0,0)circle(1pt)node[below right]{$O$};
\end{tikzpicture}
\caption{两个相互垂直且速率比为二的速度}\label{fig:hw8-6-q}
\end{figure}

\begin{solution}
逆平方力的速度轨迹是圆：由$\vct e=\vct v\times\vct h/\mu-\vct e_r$可知，速度圆半径为$R_v=\mu/h$，圆心距速度原点为$eR_v$。

\begin{figure}[H]\centering
\begin{tikzpicture}[>=Stealth,line cap=round,line join=round]
\draw[axis](0,0)--(3,0)node[right]{$v_x$};\draw[axis](0,0)--(0,3)node[above]{$v_y$};\coordinate(A)at(0,2);\coordinate(B)at(1,0);\fill(A)circle(2pt)node[above right]{$\vct v_1$};\fill(B)circle(2pt)node[below right]{$\vct v_2$};\draw[curve](-.5,.5)circle(1.581);\fill(-.5,.5)circle(1pt)node[left]{$H$};\draw[guide](-.5,.5)--(0,0);\draw[guide](-.5,.5)--(1,0);\node at(2.5,-1){$e=OH/R_v$};
\draw[axis](0,0)--(0,2);\draw[axis](0,0)--(1,0);\fill(0,0)circle(1pt)node[below right]{$O$};
\end{tikzpicture}
\caption{最小偏心率对应的速度圆心}\label{fig:hw8-6-s}
\end{figure}

取近心方向为位置$x$轴时，速度分量为$(-\mu\sin\theta/h,\mu(e+\cos\theta)/h)$，直接满足
\[
 v_x^2+\left(v_y-\frac{\mu e}{h}\right)^2=
 \left(\frac\mu h\right)^2.
\]
旋转速度坐标轴不改变圆心距与半径之比，故可把两已知速度放在互相垂直的坐标轴上求最小值。

以$u=|\vct v_2|$为单位，将图\ref{fig:hw8-6-s}的两个速度端点写成$A=(0,2)$、$B=(1,0)$。设速度圆心$H=(a,b)$，等距条件给出$a=2b-3/2$，因此
\[
 |OH|^2=5b^2-6b+9/4=:N,\quad
 R_v^2=5b^2-10b+25/4=:D.
\]
偏心率平方为$e^2=N/D$。直接计算
\[
 5N-D=5(2b-1)^2\ge0,
\]
故$e^2\ge1/5$。当$b=1/2,a=-1/2$时等号成立，圆半径为$\sqrt{5/2}$、圆心距为$1/\sqrt2$，确有一条椭圆满足两速度条件。
\[
 \boxed{e_{\min}=1/\sqrt5}.
\]
上述速度圆给出了达到下界的轨道。两个已知速度相互垂直，分别把它们认作近、远心点速度并不符合题设。


最后确认等号情形的两个速度端点确实可达。用$\vct w_i=\vct v_i/u$，圆心$H=(-1/2,1/2)$、半径$\sqrt{5/2}$，有
\[
 \frac{A-H}{R_v}=\frac{(1,3)}{\sqrt{10}},\qquad
 \frac{B-H}{R_v}=\frac{(3,-1)}{\sqrt{10}}.
\]
它们都是单位向量，且互相垂直，所以是速度圆上两个真实而不同的点。所求$e=1/\sqrt5<1$，对应的开普勒椭圆会遍历整个速度圆，故不存在因只取双曲线某一圆弧而使这两点不可达的问题。图\ref{fig:hw8-6-s}的构造给出的下界可以由实际椭圆轨道达到。

\end{solution}



\begin{example}[力心在椭圆中心的轨道]\label{sup8:problem-example-9}\label{sup8:ex:d8-e9}
质点受平面中心力 $\vct F=-k\vct r$，质量为 $m$。初始位置为 $(A,0)$，初速度为 $(0,v_0)$，其中 $0<v_0<A\sqrt{k/m}$。求轨道、角动量、能量及相邻拱点的角间隔，说明它与引力椭圆的区别。
\end{example}
\begin{figure}[H]\centering
\begin{tikzpicture}
 \draw[axis] (-1,0)--(3.6,0) node[right] {$x$};\draw[axis] (0,-.5)--(0,2.8) node[above] {$y$};
 \fill (0,0) circle (1.6pt) node[below left] {$O$};\fill (3,0) circle (1.6pt) node[below] {$P_0=(A,0)$};
 \draw[axis] (3,0)--(3,1.5) node[right] {$v_0$};
 \draw[axis] (3,0)--(1.6,0) node[midway,above] {$kA$};
\end{tikzpicture}
\caption{各向同性线性中心力，初始速度与半径垂直。}\label{sup8:fig:d8-e9-q}
\end{figure}
\begin{solution}
取 $\omega=\sqrt{k/m}$。图\ref{sup8:fig:d8-e9-q} 的两个坐标满足同频率的独立方程
\[
 \ddot x+\omega^2x=0,\qquad\ddot y+\omega^2y=0.
\]
各坐标的通解为 $x=C_x\cos\omega t+D_x\sin\omega t$、$y=C_y\cos\omega t+D_y\sin\omega t$。初位置给出 $C_x=A,C_y=0$；通解求导后，初速度给出 $\omega D_x=0,\omega D_y=v_0$。代入这些常数，有
\[
 x=A\cos\omega t,\qquad y=B\sin\omega t,\qquad B=\frac{v_0}{\omega}<A.
\]
消去时间得 $x^2/A^2+y^2/B^2=1$，力心 $O$ 位于椭圆几何中心。相对于力心的单位质量角动量、总角动量与机械能分别为
\[
 h=x\dot y-y\dot x=\omega AB=Av_0,\qquad L=mh=mAv_0,
\]
\[
 E=\frac12m\omega^2(A^2+B^2).
\]
这里势能零点取在 $O$，所有非静止轨道的 $E$ 均为正，仍然束缚。引力问题中的“负能量对应束缚”依赖其在无穷远取零的势能，不能搬到本题。
\begin{figure}[H]\centering
\begin{tikzpicture}[scale=.9]
 \draw[curve] (0,0) ellipse (3 and 1.7);\draw[guide] (-3.4,0)--(3.5,0);\draw[guide] (0,-2.2)--(0,2.3);
 \fill (0,0) circle (1.7pt) node[below left] {$O$};
 \fill (3,0) circle (1.5pt) node[right] {$P_a$};\fill (0,1.7) circle (1.5pt) node[above] {$P_p$};
 \draw[axis] plot[domain=15:55,samples=25] ({3*cos(\x)},{1.7*sin(\x)});
 \draw (.6,0) arc (0:90:.6) node[midway,above right] {$\pi/2$};
 \draw[<->,>=Stealth] (-3,-2.1)--(3,-2.1) node[midway,below] {$2A$};
 \node[right] at (3.7,1) {$r_{\max}=A$};\node[right] at (3.7,.3) {$r_{\min}=B$};
\end{tikzpicture}
\caption{线性中心力的椭圆以力心为中心，相邻远、近拱点夹角为 $\pi/2$。}\label{sup8:fig:d8-e9-s}
\end{figure}
由图\ref{sup8:fig:d8-e9-s}，从 $P_a$ 出发，经过 $\pi/(2\omega)$ 到 $P_p$。进一步写出
\[
 r^2=A^2\cos^2\omega t+B^2\sin^2\omega t
 =\frac{A^2+B^2}2+\frac{A^2-B^2}2\cos2\omega t.
\]
所以径向运动周期为 $\pi/\omega$，而质点回到同一位置和速度的周期为 $2\pi/\omega$。一圈中有两个最近点、两个最远点，相邻同类拱点的角间隔为 $\pi$。在圆极限 $A=B$，半径恒定，不再有区分开的拱点。

引力椭圆的力心在焦点，相邻远、近拱点相隔 $\pi$。两种中心力都能产生椭圆，力心的位置却不同；由轨道形状反推力律时，这个位置也是所需条件。
\end{solution}



\section{面积定律、飞行时间与会合相位}

轨道形状确定后，飞行时间由面积定律或径向积分求出。转移轨道的半周期给出飞行时长，会合还要求两物体在终点具有相同相位。



\subsection{面积定律与飞行时间}
\subsubsection{周期由面积求出}

半径矢量在 $\dd t$ 内扫过的面积为 $\dd S=\tfrac12r^2\dd\varphi$，所以
\begin{equation}
 \frac{\dd S}{\dd t}=\frac h2.
\end{equation}
椭圆面积 $\pi ab_{\rm ell}$ 除以面积速度，得到
\begin{equation}\label{sup8:eq:d8-period}
 T=\frac{2\pi ab_{\rm ell}}h=2\pi\sqrt{\frac{a^3}{\mu}}.
\end{equation}
这个周期用于相对轨道时，$a$ 是两物体间距轨道的半长轴，$\mu=G(M+m)$。若代入单个恒星绕质心的半长轴，必须同时改变相应参数，不能原样套用。

\subsubsection{偏近点角与开普勒方程}

求一段轨道的时间时，应先确定运动分支，再求扫过面积。对椭圆，引入偏近点角 $\mathcal E$。图\ref{sup8:fig:d8-anomaly} 以几何中心 $C$ 为原点画辅助圆，椭圆点的参数为
\[
 x_C=a\cos\mathcal E,\qquad y_C=b_{\rm ell}\sin\mathcal E.
\]
力心在 $x_C=ae$ 处，因此相对于力心
\begin{equation}
 x=a(\cos\mathcal E-e),\quad
 y=b_{\rm ell}\sin\mathcal E,\quad
 r=a(1-e\cos\mathcal E).
\end{equation}
又因
\[
 x\frac{\dd y}{\dd\mathcal E}-y\frac{\dd x}{\dd\mathcal E}
 =ab_{\rm ell}(1-e\cos\mathcal E),
\]
面积定律要求 $h\dd t=(x\dd y/\dd\mathcal E-y\dd x/\dd\mathcal E)\dd\mathcal E$。从近心点 $\mathcal E=0,t=t_p$ 积分，
\[
t-t_p=\frac{ab_{\rm ell}}h\int_0^{\mathcal E}(1-e\cos u)\dd u
=\frac{ab_{\rm ell}}h[u-e\sin u]_0^{\mathcal E}.
\]
又 $h/(ab_{\rm ell})=\sqrt{\mu/a^3}=n$，因此
\begin{equation}\label{sup8:eq:d8-kepler}
 n(t-t_p)=\mathcal E-e\sin\mathcal E,\qquad
 n=\sqrt{\frac{\mu}{a^3}}.
\end{equation}
$t_p$ 为一次经过近心点的时刻。跨过远心点后，$\mathcal E$ 应继续从 $\pi$ 增大到 $2\pi$，不能在反余弦的主值区间中折回。

\begin{figure}[H]\centering
\begin{tikzpicture}[scale=.9]
 \coordinate (C) at (0,0);\coordinate (O) at (1.8,0);
 \coordinate (Q) at (1.268,2.719);\coordinate (P) at (1.268,2.175);
 \draw[guide] (C) circle (3);
 \draw[curve] (C) ellipse (3 and 2.4);
 \draw[guide] (-3.4,0)--(3.5,0);\draw[guide] (Q)--(1.268,0);
 \draw (C)--(Q);\draw[axis] (O)--(P) node[midway,right] {$r$};
 \fill (C) circle (1.3pt) node[below] {$C$};\fill (O) circle (1.7pt) node[below] {$O$};
 \fill (Q) circle (1.3pt) node[above right] {$Q$};\fill (P) circle (1.3pt) node[below left] {$P$};
 \draw (.7,0) arc (0:65:.7) node[midway,right] {$\mathcal E$};
 \draw[<->,>=Stealth] (0,-.55)--(1.8,-.55) node[midway,below] {$ae$};
 \node[right,align=left] at (3.2,2) {辅助圆半径 $a$\\椭圆半短轴 $b_{\rm ell}$};
\end{tikzpicture}
\caption{偏近点角在几何中心 $C$ 处量取；真近点角在焦点 $O$ 处量取，两者一般不相等。}\label{sup8:fig:d8-anomaly}
\end{figure}

从真近点角换到偏近点角，可先把上述坐标除以$r$：
\[
 \cos f=\frac{\cos\mathcal E-e}{1-e\cos\mathcal E},\qquad
 \sin f=\frac{\sqrt{1-e^2}\sin\mathcal E}{1-e\cos\mathcal E}.
\]
第一式交叉相乘，得$(1+e\cos f)\cos\mathcal E=e+\cos f$；由此又有$1-e\cos\mathcal E=(1-e^2)/(1+e\cos f)$。将它代入第二式，便得到
\[
 \cos\mathcal E=\frac{e+\cos f}{1+e\cos f},\qquad
 \sin\mathcal E=\frac{\sqrt{1-e^2}\sin f}{1+e\cos f}.
\]
两个式子一起保留象限。若只求半个周期，近心点到远心点直接用 $T/2$，无需再积分。


\subsection{例题与模型变式}

\begin{example}[由速度方向确定椭圆与飞行时间]\label{ex:hw8-1}

行星椭圆轨道周期为$T$。真近点角为$\theta_1$时，向外径向与速度矢量的有向夹角为$\alpha_1$，行星按$\theta$增加方向运行。求真近点角$\theta_2$处的速度夹角$\alpha_2$，以及从$\theta_1$到$\theta_2$的运行时间。角度以同一近心方向为基准。

\end{example}


\begin{figure}[H]\centering
\begin{tikzpicture}[>=Stealth,line cap=round,line join=round]
\draw(0,0)ellipse(2.3 and 1.5);\coordinate(F)at(1.74356,0);\coordinate(P)at(1.47841,1.14907);\fill(F)circle(2pt)node[below]{$F$};\fill(P)circle(2pt)node[above right]{$P$};\draw[guide](F)--(P);\draw[axis](P)--(.59746,1.63116)node[above]{$\vct v$};\draw[guide](-2.5,0)--(2.7,0);\draw(2.14356,0)arc(0:103:.4)node[midway,above right]{$\theta$};
\draw[thin](P)--++(-.2,.87);\draw(1.3411,1.7446)arc(103:151.32:.611);\node at(1.03,1.78){$\alpha$};
\end{tikzpicture}
\caption{椭圆焦点、真近点角与切向速度}\label{fig:hw8-1-q}
\end{figure}

\begin{solution}
由 $r=p/(1+e\cos\theta)$，按链式法则对极角求导，$\dd r/\dd\theta=pe\sin\theta/(1+e\cos\theta)^2$。再乘 $\dot\theta=h/r^2=h(1+e\cos\theta)^2/p^2$，得 $v_r=(he/p)\sin\theta$；横向速度则是 $v_\theta=r\dot\theta=h/r=h(1+e\cos\theta)/p$。用 $p=h^2/\mu$，图\ref{fig:hw8-1-s}中的两个分量分别为
\[
 v_r=\frac{\mu}{h}e\sin\theta,\qquad
 v_\theta=\frac{\mu}{h}(1+e\cos\theta).
\]

\begin{figure}[H]\centering
\begin{tikzpicture}[>=Stealth,line cap=round,line join=round]
\draw(0,0)ellipse(2.3 and 1.5);\coordinate(F)at(1.74356,0);\coordinate(P)at(1.47841,1.14907);\fill(F)circle(2pt)node[below]{$F$};\fill(P)circle(2pt)node[above right]{$P$};\draw[guide](F)--(P);\draw[axis](P)--(.59746,1.63116)node[above]{$\vct v$};\draw[guide](-2.5,0)--(2.7,0);\draw(2.14356,0)arc(0:103:.4)node[midway,above right]{$\theta$};\draw[axis](1.47841,1.14907)--(1.3284,1.799)node[right]{$v_r$};\draw[axis](1.47841,1.14907)--(.7475,.9802)node[below]{$v_\theta$};\node at(-.5,-2){$M=E-e\sin E$};\draw[thin](P)--++(-.2,.87);\draw(1.3411,1.7446)arc(103:151.32:.611);\node at(1.03,1.78){$\alpha$};
\end{tikzpicture}
\caption{速度分解及径向与速度的夹角}\label{fig:hw8-1-s}
\end{figure}

因为 $\cot\alpha=v_r/v_\theta$，令 $q=\cot\alpha_1$，有 $q=e\sin\theta_1/(1+e\cos\theta_1)$。交叉相乘并将含 $e$ 的项合并，$q=e(\sin\theta_1-q\cos\theta_1)$，因此一般情况下
\[
 \boxed{e=\frac{q}{\sin\theta_1-q\cos\theta_1}},\qquad
 \boxed{\alpha_2=\operatorname{atan2}(1+e\cos\theta_2,e\sin\theta_2)}.
\]
应检验$0\le e<1$。若已知点恰是拱点且速度切向，两个角不能单独确定$e$；上述分式此时为$0/0$，题目数据不足，须再给一个形状条件。

用偏近点角 $E$ 处理时间，先由真近点角求
\[
E(\theta)=\operatorname{atan2}(\sqrt{1-e^2}\sin\theta,e+\cos\theta).
\]
沿真实运动将 $E$ 连续展开。为推导时间积分，采用焦点为原点、近心方向为 $x$ 轴的辅助圆参数 $x=a(\cos E-e)$、$y=b\sin E$，其中 $b=a\sqrt{1-e^2}$。对 $E$ 求导，$x_E=-a\sin E$、$y_E=b\cos E$，于是
\[
xy_E-yx_E=ab[(\cos E-e)\cos E+\sin^2E]=ab(1-e\cos E).
\]
面积定律为 $h=x\dot y-y\dot x=(xy_E-yx_E)\dot E$。用 $h=nab$、$n=2\pi/T$，得 $\dd t/\dd E=(1-e\cos E)/n$。因此
\[
\Delta t=\frac1n\int_{E_1}^{E_2}(1-e\cos E)\,\dd E
=\frac1n[E-e\sin E]_{E_1}^{E_2},
\]
即
\[
\boxed{\Delta t=\frac{T}{2\pi}[E_2-E_1-e(\sin E_2-\sin E_1)]}.
\]
记 $M=E-e\sin E$，则 $\dot M=n$。$\operatorname{atan2}$ 保留分子、分母的象限信息，因此比只用正切函数更适合跨象限的轨道计时。

若跨过近心方向或多运行整圈，展开角中的$2\pi$不可丢掉。圆轨道时$e=0$，结果退化为均匀角速度运动。


\end{solution}



\begin{example}[两次切向脉冲与会合相位]\label{deep8:problem-example-5}\label{deep8:ex:d8-e5}
卫星原在半径 $r_1$ 的圆轨道上逆时针运动，拟用与内外圆相切的半椭圆，转移至半径 $r_2=4r_1$ 的同向圆轨道。求两次速度增量及转移时间。若外圆上另有目标按圆周运动，求卫星第一次脉冲时目标应领先的相位角。
\end{example}
\begin{figure}[H]\centering
\begin{tikzpicture}[scale=.68]
 \draw[guide] (0,0) circle (1);\draw[curve] (0,0) circle (4);
 \draw[curve] plot[domain=0:180,samples=100] ({-1.5+2.5*cos(\x)},{2*sin(\x)});
 \fill (0,0) circle (2pt) node[below right] {$O$};
 \fill (1,0) circle (2pt) node[below right] {$P_1$};\fill (-4,0) circle (2pt) node[left] {$P_2$};
 \draw[axis] (1,0)--(1,1.4) node[right] {第一次脉冲};
 \draw[axis] (-4,0)--(-4,-1) node[left] {第二次脉冲};
 \fill ({4*cos(91.05)},{4*sin(91.05)}) circle (2pt) node[above] {目标初位置 $Q_0$};
 \draw (.6,0) arc (0:91.05:.6) node[midway,above right] {$\delta$};
\end{tikzpicture}
\caption{转移从右端 $P_1$ 开始，沿上半椭圆到左端 $P_2$；目标也沿逆时针方向运动。}\label{deep8:fig:d8-e5-q}
\end{figure}
\begin{solution}
图\ref{deep8:fig:d8-e5-q} 中 $OP_1=r_1,OP_2=4r_1$ 分别为转移椭圆的近、远心距离，所以
\[
 a_t=\frac52r_1,\qquad e_t=\frac35.
\]
令 $v_c=\sqrt{\mu/r_1}$。两条圆轨道上的速率分别为 $v_c,v_c/2$。由活力公式，转移椭圆两端速率为
\[
 v_{t1}=\sqrt{\mu\left(\frac2{r_1}-\frac2{5r_1}\right)}
       =\frac4{\sqrt{10}}v_c,\qquad
 v_{t2}=\sqrt{\mu\left(\frac2{4r_1}-\frac2{5r_1}\right)}
       =\frac1{\sqrt{10}}v_c.
\]
两端均须沿当时运动方向加速：
\[
 \Delta v_1=\left(\frac4{\sqrt{10}}-1\right)v_c,\qquad
 \Delta v_2=\left(\frac12-\frac1{\sqrt{10}}\right)v_c.
\]
发动机两次提供的速度增量按大小相加，得 $\Delta v_{\rm sum}=(3/\sqrt{10}-1/2)v_c$。两个脉冲发生在不同位置，方向也不同，燃料需求所对应的是这个标量和。
\begin{figure}[H]\centering
\begin{tikzpicture}
 \begin{scope}[xshift=-2.3cm]
  \fill (0,0) circle (1.6pt) node[below] {$P_1$};
  \draw[axis,line width=1pt] (0,0)--(0,2.53) node[above] {$v_{t1}$};
  \draw[axis] (.65,0)--(.65,2) node[right] {$v_c$};
  \draw[axis] (1.65,2)--(1.65,2.53) node[midway,right] {$\Delta v_1$};
  \draw[guide] (0,2)--(1.65,2);\draw[guide] (0,2.53)--(1.65,2.53);
 \end{scope}
 \begin{scope}[xshift=2.2cm]
  \fill (0,2.53) circle (1.6pt) node[above] {$P_2$};
  \draw[axis] (0,2.53)--(0,{2.53-2/sqrt(10)}) node[left] {$v_{t2}$};
  \draw[axis,line width=1pt] (.65,2.53)--(.65,1.53) node[below right=2pt] {$v_c/2$};
  \draw[axis] (1.75,{2.53-2/sqrt(10)})--(1.75,1.53) node[midway,right] {$\Delta v_2$};
  \draw[guide] (0,{2.53-2/sqrt(10)})--(1.75,{2.53-2/sqrt(10)});\draw[guide] (.65,1.53)--(1.75,1.53);
 \end{scope}
 \node at (0,-.55) {$\Delta\varphi_{\text{卫星}}=\pi,\quad\Delta\varphi_{\text{目标}}=n_2t_t$};
\end{tikzpicture}
\caption{两图的速度箭头采用相同标度。$P_1$ 处速度向上，$P_2$ 处速度向下，两个脉冲均沿各自的运动方向。}\label{deep8:fig:d8-e5-s}
\end{figure}
图\ref{deep8:fig:d8-e5-s} 的箭头方向对应题图的两个接触点。飞行经过半个转移椭圆，故
\[
 t_t=\pi\sqrt{\frac{(5r_1/2)^3}{\mu}}.
\]
目标在外圆上的角速度为 $n_2=\sqrt{\mu/(4r_1)^3}$。以卫星的初始半径 $OP_1$ 为零角，目标初始角度为 $\delta$，会合条件为
\[
 \delta+n_2t_t=\pi\pmod{2\pi}.
\]
取 $0\le\delta<2\pi$，得到
\[
 \boxed{\delta=\pi\left[1-\left(\frac58\right)^{3/2}\right]
 \simeq1.5893\ \mathrm{rad}\simeq91.1^\circ}.
\]
会合的初相位由转移时间和目标角速度共同确定。上述结果限于指定的两次切向转移；若允许增加脉冲或改变中间轨道，最优方案须另行求解。
\end{solution}



\section{瞬时变轨、碰撞截面与散射}

冲量把发生在同一位置的前后两条轨道连接起来。讨论碰撞或散射时，轨道最近距离还须与天体半径及相对入射速度对应。



\subsection{冲量、碰撞与散射的边界}
\subsubsection{瞬时变轨分前后两条轨道}

短时发动机作用可理想化为固定位置上的速度突变 $\Delta\vct v$。忽略这段时间内引力的冲量，位置连续，而
\begin{equation}\label{sup8:eq:d8-impulse}
 \varepsilon_+-\varepsilon_-=\vct v_-\cdot\Delta\vct v
       +\frac12|\Delta\vct v|^2,\qquad
 \vct h_+-\vct h_-=\vct r\times\Delta\vct v.
\end{equation}
冲量前后分别使用守恒量；跨过冲量的瞬间，机械能和角动量一般不守恒。纯径向冲量不改 $\vct h$，切向冲量通常同时改变能量和角动量。若中心天体的质量被外界突然改变，则即使速度连续，引力势也会突变，前后能量仍不能直接相等。

两圆轨道间的两次切向变轨可选用与两圆相切的半椭圆。内、外圆半径为 $r_1<r_2$ 时，转移椭圆半长轴 $a_t=(r_1+r_2)/2$，由活力公式分别计算两端速度差，转移时间为 $\pi\sqrt{a_t^3/\mu}$。这是指定共面、两次切向脉冲方案的结果；是否与目标会合，还须满足目标的相位条件。

\subsubsection{有限半径天体和引力聚焦}

物体从无穷远以速率 $v_\infty$ 飞来，瞄准距离 $b$ 是入射渐近直线到力心的垂距，而不是最近距离。守恒量为
\[
 \varepsilon=\frac12v_\infty^2,\qquad h=bv_\infty.
\]
最近点 $\dot r=0$，速率全为切向分量 $v_p=h/r_p=bv_\infty/r_p$。由能量守恒，
\[
\frac12v_\infty^2=\frac{b^2v_\infty^2}{2r_p^2}-\frac\mu{r_p}.
\]
乘以 $2r_p^2/v_\infty^2$，最近距离 $r_p$ 满足
\begin{equation}\label{sup8:eq:d8-focusing}
 b^2=r_p^2+\frac{2\mu r_p}{v_\infty^2}.
\end{equation}
对于半径 $R$ 的吸收天体，令 $r_p=R$ 得擦边条件。比这一临界 $b$ 更小的入射轨道先接触表面，点质量轨道向表面以内的延拓没有实际意义。计算碰撞截面时用 $\pi b_{\max}^2$，不能用 $\pi R^2$ 代替引力聚焦后的入射区域。

\subsubsection{散射角和参考系}

散射角是速度入射方向与出射方向之间的夹角，取$0\le\chi\le\pi$。考虑$b>0$的吸引性反平方散射。将$\varepsilon=v_\infty^2/2$、$h=bv_\infty$代入偏心率公式，得
\[
 e=\sqrt{1+\frac{b^2v_\infty^4}{\mu^2}}>1.
\]
轨道在$f=\pm f_\infty$处趋于无穷远，分母趋零给出$\cos f_\infty=-1/e$，故$\pi/2<f_\infty<\pi$。入射位置射线与入射速度反向，速度偏转角为$\chi=2f_\infty-\pi$，因此
\[
 \sin\frac\chi2=\sin\left(f_\infty-\frac\pi2\right)
 =-\cos f_\infty=\frac1e.
\]
在$0<\chi<\pi$内，$\cos(\chi/2)=\sqrt{e^2-1}/e$为正，再取两者之比，得到
\[
 \tan\frac\chi2=\frac1{\sqrt{e^2-1}}
 =\frac{\mu}{bv_\infty^2}.
\]
求这些角度时，入射渐近射线指向远处，入射速度却指向遭遇区，二者方向相反。速度偏转角须按入射、出射速度量取。排斥性反平方力的偏转角大小也有相同形式，弯曲方向则与吸引情形相反，最近距离公式中的符号随之改变。

若力心以速度 $\vct V$ 运动，应在近似惯性的力心系中先求散射。在这一参考系中，入射、出射相对速度的模相等；变回原参考系后，速率一般不等。引力辅助中的能量转移正是由这一步速度相加得到的。参考系、势能零点和中心天体是否可视为固定，均须在列式前确定。


\subsection{例题与模型变式}

\begin{example}[圆轨道卫星抛出载荷]\label{ex:hw8-4}

质量$m$的卫星绕质量$M$的中心天体作半径$r_0$的圆周运动，中心天体视为固定。瞬间抛出质量$m/10$的载荷，发射的相对速度与原轨道正向切线夹角为$60^\circ$，并有向外径向分量。发射释放机械能$\Delta E=9GMm/(2r_0)$，抛出后两部分之间的相互引力忽略不计。求余下卫星新轨道的近心距、远心距及长轴方向。

\end{example}


\begin{figure}[H]\centering
\begin{tikzpicture}[>=Stealth,line cap=round,line join=round]
\draw(0,0)circle(1.6);\fill(0,0)circle(2pt)node[left]{$M$};\coordinate(P)at(1.6,0);\fill(P)circle(2pt)node[below]{$P$};\draw[axis](P)--(1.6,1.3)node[left]{$v_c$};\draw[axis](P)--(2.6,.58)node[right]{抛出方向};\draw[guide](0,0)--(P);\node at(.8,-.3){$r_0$};
\end{tikzpicture}
\caption{载荷从圆轨道卫星向外侧抛出}\label{fig:hw8-4-q}
\end{figure}

\begin{solution}
置$v_c=\sqrt{GM/r_0}$。发射前后在同一点，外引力势能不变，故释放能量全部增加系统动能。在质心系中，剩余卫星与载荷的约化质量为$9m/100$，于是
\[
 \frac12\frac{9m}{100}u^2=\frac92mv_c^2,
 \qquad u=10v_c.
\]

\begin{figure}[H]\centering
\begin{tikzpicture}[>=Stealth,line cap=round,line join=round]
\draw[guide,gray](0,0)circle(1.8);
\begin{scope}[shift={(1.35,-.77942286)},rotate=150]\draw[curve](0,0)ellipse(1.8 and .9);\end{scope}
\draw[guide](-.9,.519615)--(3.12,-1.80133);
\fill(0,0)circle(2pt)node[below left]{$M$};\fill(1.8,0)circle(2pt)node[below right]{$P$};
\fill(-.20884573,.12057714)circle(1.5pt);\draw[thin](-.20884573,.12057714)--(-.55,.65);\node[fill=white,inner sep=1pt]at(-.55,.85){$Q_p$};
\fill(2.90884573,-1.67942286)circle(1.5pt)node[below right]{$Q_a$};
\draw[axis](0,0)--(-1.04,.6)node[above left]{$\vct e$};
\draw[axis](1.8,0)--(1.02,0)node[above]{$v_r'$};\draw[axis](1.8,0)--(1.8,.45)node[right]{$v_t'$};
\draw[guide](0,0)--(1.8,0);\node[below]at(.8,0){$r_0$};
\draw[guide,gray](-1.4,-2.4)--(-.7,-2.4)node[right]{原圆轨道};
\draw[curve](1.1,-2.4)--(1.8,-2.4)node[right]{反冲后轨道};
\end{tikzpicture}
\caption{反冲后的椭圆与原圆相交于发射点，长轴沿偏心率矢量}\label{fig:hw8-4-s}
\end{figure}

图\ref{fig:hw8-4-s}中载荷相对卫星的速度为$u(\sin60^\circ\vct e_r+\cos60^\circ\vct e_t)$。动量守恒使卫星反冲速度等于该矢量的$-1/10$，因此
\[
 (v_r',v_t')=(-\sqrt3v_c/2,v_c/2).
\]
余下卫星的比能量仍为$-GM/(2r_0)$，所以$a=r_0$；比角动量$h=r_0v_c/2$，故$p=r_0/4$、$e=\sqrt3/2$。于是
\[
 \boxed{r_p=r_0(1-\sqrt3/2),\quad r_a=r_0(1+\sqrt3/2)}.
\]
偏心率矢量为
\[
 \vct e=\left(\frac{v_t'^2}{v_c^2}-1\right)\vct e_r
 -\frac{v_r'v_t'}{v_c^2}\vct e_t
 =-\frac34\vct e_r+\frac{\sqrt3}{4}\vct e_t.
\]
图\ref{fig:hw8-4-s}将新轨道画成以中心天体为焦点的椭圆，$Q_p,Q_a$是新近、远心点，$P$为发射点。半短轴为$b=a\sqrt{1-e^2}=r_0/2$。发射时$v_r'<0$，所以卫星先朝$Q_p$运行；把$P$直接标作远心点会与这一径向速度矛盾。

故从发射点向外径向转向原运动切向量起，近心方向为$150^\circ$，无向长轴与原径向线成$30^\circ$。若载荷改向内侧抛出，轨道形状相同，长轴方向镜像反射。


\end{solution}



\begin{example}[减速后的轨道与撞击条件]\label{deep8:problem-example-6}\label{deep8:ex:d8-e6}
卫星原在半径 $r_0>R$ 的圆轨道上运行，中心天体的半径为 $R$，引力参数为 $\mu$。一次瞬时切向制动使速率变为原来的 $q$ 倍，$0<q<1$，方向不变。求新轨道参数、能量和角动量的变化，并求不会撞击天体的 $q$ 的范围。擦过表面单列为临界情形。
\end{example}
\begin{figure}[H]\centering
\begin{tikzpicture}[scale=.85]
 \draw[curve,fill=black!7] (0,0) circle (.65);\draw[guide] (0,0) circle (2);
 \fill (0,0) circle (1.6pt) node[below] {$O$};\fill (2,0) circle (1.6pt) node[below right] {$A$};
 \draw[axis] (2,0)--(2,2) node[above] {$v_c$};
 \draw[axis,line width=1pt] (2.8,0)--(2.8,1.2) node[right] {$qv_c$};
 \draw[axis] (3.8,2)--(3.8,1.2) node[midway,right] {制动};
 \draw[curve] (0,0)--(2,0) node[midway,below] {$r_0$};
 \draw (0,0)--(-.46,.46) node[midway,above left] {$R$};
\end{tikzpicture}
\caption{在 $A$ 处减小切向速度。$v_c=\sqrt{\mu/r_0}$，制动后仍沿原转向运行。}\label{deep8:fig:d8-e6-q}
\end{figure}
\begin{solution}
由图\ref{deep8:fig:d8-e6-q}，制动前后位置相同、径向速度均为零。制动后的守恒量为
\[
 \varepsilon_+=-\frac{\mu}{2r_0}(2-q^2),\qquad
 h_+=q\sqrt{\mu r_0}.
\]
因此
\[
 a=\frac{r_0}{2-q^2},\quad e=1-q^2,\quad
 r_a=r_0,\quad r_p=\frac{q^2r_0}{2-q^2}.
\]
由于 $q<1$，$A$ 为远心点。制动后 $\ddot r<0$，卫星先沿新的椭圆向内飞行。对质量为 $m$ 的卫星，制动造成的变化为
\[
 E_+-E_-=-\frac{m\mu}{2r_0}(1-q^2),\qquad
 L_+-L_-=m\sqrt{\mu r_0}(q-1).
\]
制动力所做的功和冲量矩分别改变能量与角动量。这两项量各自在制动前、后的自由飞行阶段守恒。
\begin{figure}[H]\centering
\begin{tikzpicture}[scale=1.25]
 \draw[curve,fill=black!7] (0,0) circle (.65);
 \pgfmathsetmacro{\periEight}{2*.6^2/(2-.6^2)}
 \pgfmathsetmacro{\centerEight}{(2-\periEight)/2}
 \pgfmathsetmacro{\aEight}{(2+\periEight)/2}
 \pgfmathsetmacro{\bEight}{sqrt(2*\periEight)}
 \draw[curve] (.675,0) ellipse [x radius=1.325,y radius={sqrt(1.3)}];
 \draw[guide] (\centerEight,0) ellipse [x radius=\aEight,y radius=\bEight];
 \fill (0,0) circle (1.2pt) node[below] {$O$};
 \fill (2,0) circle (1.3pt) node[right] {$A$};
 \fill (-.65,0) circle (1.3pt) node[left] {$P_p$};
 \draw[axis] plot[domain=40:75,samples=20] ({.675+1.325*cos(\x)},{sqrt(1.3)*sin(\x)});
 \node[align=left] at (3,1) {实线：$r_p=R$\\虚线：$r_p<R$};
 \draw[axis] (3,.1)--(3,-.9) node[below] {$q\text{ 减小}$};
\end{tikzpicture}
\caption{两图均取 $R/r_0=0.325$。实线为擦边椭圆；虚线取 $q=0.6$，在表面以内的部分只是点质量轨道的延拓。}\label{deep8:fig:d8-e6-s}
\end{figure}
图\ref{deep8:fig:d8-e6-s} 中的不撞击条件为 $r_p>R$。令 $\lambda=R/r_0\in(0,1)$，有
\[
 \frac{q^2}{2-q^2}>\lambda
 \quad\Longleftrightarrow\quad
 \boxed{\sqrt{\frac{2\lambda}{1+\lambda}}<q<1}.
\]
等号时擦过表面；小于临界值时撞击。若轨道未被表面截断，从 $A$ 到 $P_p$ 的时间为
\[
 t_{a\to p}=\pi\sqrt{\frac{r_0^3}{\mu(2-q^2)^3}}.
\]
撞击时，首次接触早于这个时间。若需接触时刻，可取 $\mathcal E=\pi$ 为制动点，令
\[
 C_R=\frac{1-R/a}{e},\qquad
 \mathcal E_R=2\pi-\arccos C_R,
\]
在撞击分支 $r_p<R<r_a$ 上有 $-1<C_R<1$。以 $n=\sqrt{\mu/a^3}$ 表示平均角速度，首次接触的时间积分为
\[
t_R=\frac1n\int_\pi^{\mathcal E_R}(1-e\cos\mathcal E)\,\dd\mathcal E
=\frac1n[\mathcal E-e\sin\mathcal E]_\pi^{\mathcal E_R}.
\]
利用 $\sin\pi=0$，得到
\[
 t_R=\frac{\mathcal E_R-e\sin\mathcal E_R-\pi}{\sqrt{\mu/a^3}}.
\]
这里取 $\mathcal E_R\in(\pi,2\pi)$，才对应图中由远心点向内运动的一段。
\end{solution}



\begin{example}[引力聚焦的碰撞截面]\label{deep8:problem-example-7}\label{deep8:ex:d8-e7}
半径 $R$、引力参数 $\mu$ 的球形天体静止。一束稀薄质点从无穷远平行入射，速率均为 $v_\infty$。入射直线到球心的垂距为瞄准距离 $b$。忽略质点之间的作用，接触表面即被吸收。求最近距离与 $b$ 的关系、最大碰撞瞄准距离和碰撞截面，并讨论高速与低速极限。
\end{example}
\begin{figure}[H]\centering
\begin{tikzpicture}
 \draw[curve,fill=black!7] (1,0) circle (1);\fill (1,0) circle (1.5pt) node[below] {$O$};
 \draw[guide] (-6,0)--(3,0);\draw[guide] (-5.8,2)--(1.8,2);
 \draw[axis] (-5.8,2)--(-4.5,2) node[midway,above] {$v_\infty$};
 \draw[<->,>=Stealth] (-5.8,0)--(-5.8,2) node[midway,left] {$b$};
 \draw[curve,domain=165:90,samples=140,variable=\s] plot ({1+2*cos(\s)/(1+cos(\s)+sin(\s))},{2*sin(\s)/(1+cos(\s)+sin(\s))});
 \draw[axis,domain=140:120,samples=30,variable=\s] plot ({1+2*cos(\s)/(1+cos(\s)+sin(\s))},{2*sin(\s)/(1+cos(\s)+sin(\s))});
 \fill (1,1) circle (1.5pt) node[above right] {$P$};
 \draw (1,0)--(1,1) node[midway,right] {$R$};
\end{tikzpicture}
\caption{入射双曲线在 $P$ 处接触球面并终止。虚线是入射渐近直线，$b$ 是它到球心的垂距。}\label{deep8:fig:d8-e7-q}
\end{figure}
\begin{solution}
图\ref{deep8:fig:d8-e7-q} 中的 $b$ 和入射速率给出
\[
 \varepsilon=\frac12v_\infty^2,\qquad h=bv_\infty.
\]
图\ref{deep8:fig:d8-e7-q} 中的接触点 $P$ 通常仍有径向速度。先按点质量场延拓轨道，并把最近点记为 $P_p$；在该处 $\dot r=0$。设切向速率为 $v_p$，由角动量和能量守恒分别得
\[
 r_pv_p=bv_\infty,\qquad
 \frac12v_p^2-\frac\mu{r_p}=\frac12v_\infty^2.
\]
消去 $v_p$，得到
\[
 r_p^2+\frac{2\mu}{v_\infty^2}r_p-b^2=0.
\]
只取非负根，故
\begin{equation}
 \boxed{r_p=\frac{\sqrt{\mu^2+b^2v_\infty^4}-\mu}{v_\infty^2}}
 =\frac{b^2v_\infty^2}{\sqrt{\mu^2+b^2v_\infty^4}+\mu}.
\end{equation}
当 $b$ 很小时，后一形式避免了两个接近的数相减，适合数值计算。$b=0$ 的径向入射可由连续极限或径向方程单独求得。
\begin{figure}[H]\centering
\begin{tikzpicture}
 \begin{scope}[xshift=-2.5cm]
  \draw[curve,fill=black!7] (0,0) circle (1);
  \fill (0,0) circle (1.5pt) node[below] {$O$};\fill (1,0) circle (1.5pt) node[right] {$P_p$};
  \draw[axis] (1,0)--(1,1.4) node[right] {$v_p$};
  \draw[curve] (0,0)--(1,0) node[midway,below] {$R$};
  \node at (0,-1.5) {临界最近点：$r_p=R$};
 \end{scope}
 \begin{scope}[xshift=2.3cm]
  \draw[fill=black!4] (0,0) circle ({sqrt(5)});\draw[curve] (0,0) circle (1);
  \fill (0,0) circle (1.5pt);\draw[axis] (0,0)--({sqrt(2.5)},{sqrt(2.5)}) node[above right] {$b_{\max}$};
  \draw (0,0)--(-1,0) node[midway,below] {$R$};
  \node at (0,-2.65) {垂直入射方向的瞄准平面};
 \end{scope}
\end{tikzpicture}
\caption{左图用擦边轨道求临界守恒量；右图的外圆是能导致碰撞的入射区域，内圆是天体几何投影。}\label{deep8:fig:d8-e7-s}
\end{figure}
由图\ref{deep8:fig:d8-e7-s} 左图，擦边满足 $r_p=R$。由于 $b^2=r_p^2+2\mu r_p/v_\infty^2$ 随 $r_p\ge0$ 严格增加，碰撞条件恰为 $0\le b\le b_{\max}$，其中
\[
 b_{\max}=R\sqrt{1+\frac{2\mu}{Rv_\infty^2}},\qquad
 \boxed{\sigma=\pi R^2\left(1+\frac{2\mu}{Rv_\infty^2}\right)}.
\]
图中右侧面积是入射平面上的面积，不是天体表面的面积。令 $v_{\rm esc}=\sqrt{2\mu/R}$，聚焦因子为 $1+v_{\rm esc}^2/v_\infty^2$。

若 $v_\infty\gg v_{\rm esc}$，$\sigma\simeq\pi R^2$，轨道接近直线；若 $v_\infty\ll v_{\rm esc}$，$\sigma\simeq2\pi\mu R/v_\infty^2$。对孤立天体，后式随 $v_\infty\to0$ 发散。它描述从无穷远入射的临界面积，实际可用范围仍受飞行时间、其他天体和外场限制。
\end{solution}



\begin{example}[排斥性反平方散射]\label{deep8:problem-example-8}\label{deep8:ex:d8-e8}
固定力心对质量为 $m$ 的质点施加排斥力 $\vct F=k\vct e_r/r^2$，$k>0$。质点从无穷远以速率 $v_\infty$、瞄准距离 $b>0$ 入射。求最近距离、散射角 $\chi$，并由入射环带面积推导微分散射截面 $\dd\sigma/\dd\Omega$。
\end{example}
\begin{figure}[H]\centering
\begin{tikzpicture}[scale=.8]
 \fill (0,0) circle (2pt) node[left] {$O$};
 \fill (1.5,0) circle (1.5pt) node[below right] {$P_p$};
 \draw[axis] (1.5,0)--(2.4,0) node[above] {$\vct F$};
 \draw[curve,domain=-40:40,samples=100,variable=\s] plot ({1.5*cos(\s)/(2*cos(\s)-1)},{1.5*sin(\s)/(2*cos(\s)-1)});
 \draw[axis,domain=38:28,samples=20,variable=\s] plot ({1.5*cos(\s)/(2*cos(\s)-1)},{1.5*sin(\s)/(2*cos(\s)-1)}) node[above left] {$\vct v_{\rm in}$};
 \draw[axis,domain=-28:-38,samples=20,variable=\s] plot ({1.5*cos(\s)/(2*cos(\s)-1)},{1.5*sin(\s)/(2*cos(\s)-1)}) node[below left] {$\vct v_{\rm out}$};
 \draw[guide] (.5,-.866)--(3,3.464);\draw[guide] (1,0)--(3,-3.464);
 \draw[<->,>=Stealth] (0,0)--(.75,-.433) node[below left] {$b$};
 \node[right,align=left] at (3.2,0) {排斥力\\指向外侧};
\end{tikzpicture}
\caption{排斥双曲线的一个示意分支。两条位置渐近射线的夹角与速度偏转角互为补角。}\label{deep8:fig:d8-e8-q}
\end{figure}
\begin{solution}
令 $\alpha=k/m$，并用 $h=bv_\infty>0$ 表示单位质量角动量的模。单位质量势能为 $+\alpha/r$，与引力势的符号相反。最近点处的能量式为
\[
 \frac12v_\infty^2=\frac{b^2v_\infty^2}{2r_p^2}+\frac\alpha{r_p},
\]
所以
\[
 \boxed{r_p=\frac{\alpha+\sqrt{\alpha^2+b^2v_\infty^4}}{v_\infty^2}}>b.
\]
这一不等式符合图\ref{deep8:fig:d8-e8-q} 中排斥力把质点推离力心的方向。排斥情形的 Binet 方程为 $u''+u=-\alpha/h^2$。常数特解是 $-\alpha/h^2$，齐次通解为 $C\cos f+D\sin f$。令最近点为 $f=0$，此处 $u'(0)=0$，而 $u'=-C\sin f+D\cos f$，所以 $D=0$；再由 $u(0)=1/r_p$ 得 $C=1/r_p+\alpha/h^2$。将 $C$ 写为 $\alpha e/h^2$，并代入已求出的 $r_p$，得到
\[
 r=\frac{h^2/\alpha}{e\cos f-1},\qquad
 e=\sqrt{1+\frac{b^2v_\infty^4}{\alpha^2}}.
\]
无穷远对应 $f=\pm f_\infty$，其中 $\cos f_\infty=1/e$。入射速度指向力心一侧，故速度偏转角为 $\chi=\pi-2f_\infty$，于是
\[
 \sin\frac\chi2=\frac1e,\qquad
 \boxed{\tan\frac\chi2=\frac\alpha{bv_\infty^2}},\qquad
 b=\frac\alpha{v_\infty^2}\cot\frac\chi2.
\]
\begin{figure}[H]\centering
\begin{tikzpicture}
 \begin{scope}[xshift=-2.7cm]
  \draw[curve] (0,0) circle (1.1);\draw[curve] (0,0) circle (1.4);
  \fill[black!12,even odd rule] (0,0) circle (1.4) (0,0) circle (1.1);
  \fill (0,0) circle (1.4pt);\draw[axis] (0,0)--(1.1,0) node[midway,below] {$b$};
  \draw[<->,>=Stealth] (.778,.778)--(.99,.99) node[right] {$\dd b$};
  \node at (0,-1.9) {$\dd\sigma=2\pi b\,|\dd b|$};
 \end{scope}
 \begin{scope}[xshift=2.4cm]
  \draw[axis] (0,0)--(3,0) node[right] {入射方向};
  \draw[axis] (0,0)--(1.3,2.252) node[above] {出射方向};
  \draw[guide] (0,0)--(1.67,1.99);
  \draw (.7,0) arc (0:60:.7) node[midway,right] {$\chi$};
  \draw[<->,>=Stealth] ({1.7*cos(50)},{1.7*sin(50)}) arc (50:60:1.7) node[midway,above] {$\dd\chi$};
  \node[align=center] at (1,-1.9) {$|\dd\Omega|=2\pi\sin\chi\,|\dd\chi|$};
 \end{scope}
\end{tikzpicture}
\caption{入射环带与出射角带对应。$b$ 随 $\chi$ 增大而减小，截面计算取导数绝对值。}\label{deep8:fig:d8-e8-s}
\end{figure}
图\ref{deep8:fig:d8-e8-s} 中入射环带的面积为 $\dd\sigma=2\pi b|\dd b|$，对应出射角带的立体角为 $\dd\Omega=2\pi\sin\chi\,\dd\chi$。由于本题 $b(\chi)$ 单调，这两个角带一一对应，相除得
\[
 \frac{\dd\sigma}{\dd\Omega}
 =\frac b{\sin\chi}\left|\frac{\dd b}{\dd\chi}\right|.
\]
令 $A=\alpha/v_\infty^2$。对 $b=A\cot(\chi/2)$ 求导，外层余切产生 $-\csc^2(\chi/2)$，内层半角产生 $1/2$，即 $\dd b/\dd\chi=-(A/2)\csc^2(\chi/2)$。因此
\[
 \left|\frac{\dd b}{\dd\chi}\right|=\frac A2\csc^2\frac\chi2,
 \qquad \sin\chi=2\sin\frac\chi2\cos\frac\chi2.
\]
代入后，$\cos(\chi/2)$ 约去，得到
\[
 \boxed{\frac{\dd\sigma}{\dd\Omega}=
 \left(\frac{k}{2mv_\infty^2}\right)^2\csc^4\frac\chi2},\qquad 0<\chi<\pi.
\]
量纲为面积。前向 $\chi\to0$ 的发散对应任意大的 $b$ 仍有微小偏转，是无限作用范围模型的结果；若只统计偏转超过固定 $\chi_0>0$ 的粒子，截面有限，为 $\pi A^2\cot^2(\chi_0/2)$。
\end{solution}



\section{轨道微扰、进动与潮汐}

偏离反平方力后，能量与角动量仍可能守恒，椭圆方向却可能发生进动。偏心率矢量和近圆展开用于分析这种变化；潮汐形变与耗散导致的长期演化另作区分。



\subsection{LRL矢量及其微扰变化}
\label{sec:complete8}
比角动量为$\vct h=\vct r\times\vct v$，引力参数为$\mu=GM$。令
\[
 \vct e=\frac{\vct v\times\vct h}{\mu}-\vct e_r,
 \qquad \vct B=m\mu\vct e.
\]
后者为本讲采用的LRL矢量。角动量非零时运动位于固定平面，取$\vct h=h\vct e_z$。因$\vct e_r\times\vct h=-h\vct e_\theta$，在反平方引力下
\[
 \frac{\dd\vct e}{\dd t}
 =-\frac{\vct e_r\times\vct h}{r^2}
 -\dot\theta\vct e_\theta
 =\left(\frac h{r^2}-\dot\theta\right)\vct e_\theta=0.
\]
与$\vct e_r$点乘，$\vct e\cdot\vct e_r=h^2/(\mu r)-1$；取$\vct e$方向为真近点角零点，便得到$r=(h^2/\mu)/(1+e\cos\theta)$。偏心率矢量的模确定形状，方向确定近心位置。圆轨道$\vct e=0$，近心方向没有唯一含义。

若再叠加小的径向加速度$\delta f(r)\vct e_r$，$\vct h$仍守恒，而
\[
 \dot{\vct e}=-\frac{h\,\delta f(r)}{\mu}\vct e_\theta.
\]
例\ref{ex:hw8-5}的均匀密度背景产生$\delta f=-\nu^2r$，其中$\nu^2=4\pi G\rho/3$。先沿未扰动的微偏心椭圆求一周期积累，使用$\dd t=r^2\dd\theta/h$：
\[
 \Delta\vct e=\frac{\nu^2}{\mu}\int_0^{2\pi}r^3\vct e_\theta\dd\theta.
\]
取近心方向为 $x$ 正向，$r=a(1-e\cos\theta)+O(e^2)$，故 $r^3=a^3(1-3e\cos\theta)+O(e^2)$；又 $\vct e_\theta=(-\sin\theta,\cos\theta)$。分别计算两个分量所需的积分，
\[
\begin{aligned}
\int_0^{2\pi}(-\sin\theta+3e\sin\theta\cos\theta)\dd\theta
&=[\cos\theta-\tfrac32e\cos^2\theta]_0^{2\pi}=0,\\
\int_0^{2\pi}(\cos\theta-3e\cos^2\theta)\dd\theta
&=\left[\sin\theta-3e\left(\frac\theta2+\frac{\sin2\theta}4\right)\right]_0^{2\pi}
=-3\pi e.
\end{aligned}
\]
乘回 $\nu^2a^3/\mu$，再以 $\Delta e_y/e$ 取一阶方向变化，得到
\[
 \Delta\vct e=-\frac{3\pi\nu^2a^3}{\mu}e\vct e_y+O(e^2),
 \quad \Delta\varpi=-\frac{3\pi\nu^2a^3}{\mu}.
\]
以$n=\sqrt{\mu/a^3}$记未扰动平均角速度，进动角速度为$-3\nu^2/(2n)$，负号表示逆行进动。例题用径向频率与角向频率之差求解，在密度扰动和偏心率所保留的阶数内，两种推导给出相同结果。以上积累只保留密度扰动的一阶；不能把一周期近似在长时间内重复使用而忽略轨道要素的缓慢变化。

\subsubsection{中心势中的近圆频率}
一般中心加速度$f(r)$下，圆轨道满足$h^2=-r_0^3f(r_0)$。沿用前文固定$h$的径向线性化，将式\eqref{sup8:eq:d8-stability}中的$\Phi'=-f$、$\Phi''=-f'$代入，并令$r=r_0+\xi$，得到
\[
 \ddot\xi+\kappa_r^2\xi=0,\qquad
 \kappa_r^2=-f'(r_0)-\frac{3f(r_0)}{r_0},\qquad
 \Omega^2=-\frac{f(r_0)}{r_0}.
\]
径向往返一周对应角位移$2\pi\Omega/\kappa_r$。若两频率之比有理，线性近圆近似中轨迹闭合；一般有限偏心率轨道是否闭合，还须检验完整轨道方程。二阶径向回复消失时，应回到有效势的高阶项，不能只凭$\kappa_r=0$断言稳定。



\LectureNeedspace{4\baselineskip}


\subsection{例题与模型变式}

\begin{example}[均匀暗物质导致的近圆轨道进动]\label{ex:hw8-5}

在质量$M$的中心天体附近有均匀密度$\rho$的球对称暗物质分布。某小天体的参考圆轨道半径为$r_0$，$\rho r_0^3\ll M$。对圆轨道施加无穷小偏心扰动，求近心方向的进动周期及方向。

\end{example}


\begin{figure}[H]\centering
\begin{tikzpicture}[>=Stealth,line cap=round,line join=round]
\fill[gray!10](0,0)circle(2);\draw[guide](0,0)circle(1.5);\draw[curve,domain=0:346.153846,samples=240]plot({1.5*(1-.12*cos(1.04*\x))*cos(\x)},{1.5*(1-.12*cos(1.04*\x))*sin(\x)});\draw[thin,domain=346.153846:692.307692,samples=240]plot({1.5*(1-.12*cos(1.04*\x))*cos(\x)},{1.5*(1-.12*cos(1.04*\x))*sin(\x)});\fill(0,0)circle(2pt)node[below]{$M$};\node at(3,1){$\rho r_0^3\ll M$};
\end{tikzpicture}
\caption{均匀背景引力中的近圆轨道}\label{fig:hw8-5-q}
\end{figure}

\begin{solution}
半径$r$内暗物质质量为$4\pi\rho r^3/3$，中心加速度为$f(r)=-\mu/r^2-\nu^2r$，其中$\mu=GM$、$\nu^2=4\pi G\rho/3$。圆轨道角频率满足
\[
 \Omega^2=\frac{\mu}{r_0^3}+\nu^2.
\]

\begin{figure}[H]\centering
\begin{tikzpicture}[>=Stealth,line cap=round,line join=round]
\fill[gray!10](0,0)circle(2);\draw[guide](0,0)circle(1.5);\draw[curve,domain=0:346.153846,samples=240]plot({1.5*(1-.12*cos(1.04*\x))*cos(\x)},{1.5*(1-.12*cos(1.04*\x))*sin(\x)});\draw[thin,domain=346.153846:692.307692,samples=240]plot({1.5*(1-.12*cos(1.04*\x))*cos(\x)},{1.5*(1-.12*cos(1.04*\x))*sin(\x)});\fill(0,0)circle(2pt)node[below]{$M$};\node at(3,1){$\rho r_0^3\ll M$};\draw[axis](0,0)--(1.32,0);\draw[axis](0,0)--(1.28165,-.3159);\node[right]at(1.65,.15){原近心};\node[right]at(1.65,-.5){后继近心};\node at(0,-2.5){$\kappa>\Omega$};
\draw[axis](.8,0)arc(0:-13.84615:.8);\draw[axis](-1.59748,-.58143)arc(200:230:1.7);\node at(-1.25,-1.7){公转正向};
\end{tikzpicture}
\caption{连续两次近心方向沿公转反向偏移}\label{fig:hw8-5-s}
\end{figure}

图\ref{fig:hw8-5-s}的微小径向偏离以频率 $\kappa$ 振动。求导时须保持比角动量 $h$ 不变，而不是在扰动过程中始终令角速度等于圆轨道角速度。有效势和参考圆轨道条件为
\[
 V_{\rm eff}=\frac{h^2}{2r^2}-\frac\mu r+\frac{\nu^2r^2}{2},\qquad
 h^2=\mu r_0+\nu^2r_0^4.
\]
先按幂函数逐项求导，再代入圆轨道对应的 $h$：
\[
V_{\rm eff}'=-\frac{h^2}{r^3}+\frac\mu{r^2}+\nu^2r,\qquad
V_{\rm eff}''=\frac{3h^2}{r^4}-\frac{2\mu}{r^3}+\nu^2.
\]
令 $r=r_0+\xi$，径向方程 $\ddot r=-V_{\rm eff}'(r)$ 在平衡点附近化为 $\ddot\xi+V_{\rm eff}''(r_0)\xi=0$，所以
\[
\kappa^2=\frac{3(\mu r_0+\nu^2r_0^4)}{r_0^4}-\frac{2\mu}{r_0^3}+\nu^2
=\frac\mu{r_0^3}+4\nu^2.
\]

每次从一个近心点到下一个，转过角 $2\pi\Omega/\kappa<2\pi$，因此进动逆于公转方向。令 $n=\sqrt{\mu/r_0^3}$、$\delta=\nu^2/n^2\ll1$，用 $(1+s)^{1/2}=1+s/2+O(s^2)$ 分别展开：
\[
\Omega=n\sqrt{1+\delta}=n(1+\delta/2+O(\delta^2)),\qquad
\kappa=n\sqrt{1+4\delta}=n(1+2\delta+O(\delta^2)).
\]
两式相减，给出
\[
 \Omega_{\text{进}}=\Omega-\kappa\simeq-\frac{3\nu^2}{2n}
 =-\frac{2\pi G\rho}{n},\quad
 \boxed{T_{\text{进}}\simeq\frac{n}{G\rho}}.
\]
严格圆轨道没有唯一近心方向，故这里取无穷小偏心轨道的极限。偏心率有限时，应另作摄动平均；$r_0$处的线性频率只描述近圆运动。


进动角也可按一个径向周期计算。利用上面两个展开，作比值时再用 $(1+2\delta)^{-1}=1-2\delta+O(\delta^2)$，有
\[
 \frac\Omega\kappa=\sqrt{\frac{1+\delta}{1+4\delta}}
 =1-\frac32\delta+O(\delta^2),
 \qquad
 \Delta\varpi=2\pi\left(\frac\Omega\kappa-1\right)
 \simeq-\frac{4\pi^2\rho r_0^3}{M}.
\]
图\ref{fig:hw8-5-s}中后继近心方向落在原近心方向的公转反侧，与负号相符。用径向周期$2\pi/\kappa$除以每次进动角的绝对值、再乘$2\pi$，得到$T_{\text{进}}=2\pi/|\Omega-\kappa|$，从而还原$n/(G\rho)$。这里周期指近心线在惯性空间转过整周所需时间。

\end{solution}



\begin{example}[带逆平方势修正的进动轨道（选做）]\label{sup8:problem-example-10}\label{sup8:ex:d8-e10}
某平面中心势的单位质量势能为
\[
 \Phi(r)=-\frac\mu r+\frac{\beta}{2r^2},\qquad \mu>0.
\]
其中 $\beta$ 为可正可负的常数。对给定的非零 $h$，求非圆轨道方程、圆轨道存在条件及近圆轨道的近心点进动量。讨论 $h^2+\beta\le0$ 时为何不能沿用稳定圆轨道结论。
\end{example}
\begin{figure}[H]\centering
\begin{tikzpicture}[x=1.6cm,y=1.3cm]
 \draw[axis] (0,0)--(4.6,0) node[right] {$r$};\draw[axis] (0,-1.1)--(0,1.5) node[above] {$\Phi_{\rm eff}$};
 \draw[curve,domain=.3:4.2,samples=140] plot (\x,{.4/(\x*\x)-1/\x});
 \draw[guide] (.8,0)--(.8,-.625);\fill (.8,-.625) circle (1pt) node[below] {$r_c$};
 \node[align=left] at (2.9,1) {$\Phi_{\rm eff}=\dfrac{h^2+\beta}{2r^2}-\dfrac\mu r$\\图示 $h^2+\beta>0$};
\end{tikzpicture}
\caption{修正项与横向动能共同决定小半径处的有效势屏障。}\label{sup8:fig:d8-e10-q}
\end{figure}
\begin{solution}
实际单位质量力由势能求导：$-\Phi'(r)=-[\mu r^{-2}-\beta r^{-3}]=-\mu/r^2+\beta/r^3$。代入 Binet 方程，
\[
 u''+\left(1+\frac\beta{h^2}\right)u=\frac\mu{h^2}.
\]
设 $h^2+\beta>0$，令 $s=\sqrt{1+\beta/h^2}$。取常数特解 $u_p=\mu/(h^2s^2)=\mu/(h^2+\beta)$；齐次方程的特征根为 $\pm is$，所以 $u=u_p+C\cos s\varphi+D\sin s\varphi$。把两个三角项合成 $u_pe_*\cos[s(\varphi-\varphi_p)]$，其中 $e_*=\sqrt{C^2+D^2}/u_p$，再取倒数，得到
\begin{equation}
 r=\frac{p_*}{1+e_*\cos[s(\varphi-\varphi_p)]},\qquad
 p_*=\frac{h^2+\beta}{\mu}.
\end{equation}
$e_*$ 是本题解的振幅参数，一般不能解释为某条固定圆锥曲线的几何偏心率。对 $0<e_*<1$，半径有界且为正；相邻两个近心点满足相位差 $2\pi$，所以
\[
 \Delta\varphi_{p\to p}=\frac{2\pi}s,\qquad
 \boxed{\delta\varphi=2\pi\left(\frac1s-1\right)}.
\]
进动量是相邻两次经过近心点之间的累计转角减去 $2\pi$，方位角按运动过程连续计数。$\beta>0$ 时进动为负，表现为退行；$-h^2<\beta<0$ 时为正，表现为顺行。若 $|\beta|\ll h^2$，将 $1/s=(1+\beta/h^2)^{-1/2}=1-\beta/(2h^2)+O(\beta^2/h^4)$ 代入，有
\[
 \delta\varphi\simeq-\frac{\pi\beta}{h^2}.
\]
\begin{figure}[H]\centering
\begin{tikzpicture}[scale=1.6]
 \draw[curve,domain=0:800,samples=550,variable=\s] plot ({cos(\s)/(1+.3*cos(.9*\s))},{sin(\s)/(1+.3*cos(.9*\s))});
 \fill (0,0) circle (1.1pt) node[below] {$O$};
 \draw[guide] (0,0)--(1.6,0) node[right] {$P_{p,0}$方向};
 \draw[guide] (0,0)--({1.6*cos(40)},{1.6*sin(40)}) node[above right] {$P_{p,1}$方向};
 \fill (.7692,0) circle (1.2pt);\fill ({.7692*cos(40)},{.7692*sin(40)}) circle (1.2pt);
 \draw (.5,0) arc (0:40:.5) node[midway,right] {$\delta\varphi$};
 \node[align=center] at (0,-1.95) {$s=0.9$，相邻近心点相隔 $400^\circ$\\扣除一周后顺行 $40^\circ$};
\end{tikzpicture}
\caption{有进动的轨道不再是一条固定椭圆。图示 $e_*=0.3,\beta/h^2=-0.19$，绘制两个径向周期。}\label{sup8:fig:d8-e10-s}
\end{figure}
图\ref{sup8:fig:d8-e10-s} 的两个近心点方向相差 $40^\circ$，对应完整转角 $400^\circ$。圆轨道由图\ref{sup8:fig:d8-e10-q} 中的有效势极小点确定：
\[
 r_c=\frac{h^2+\beta}{\mu}>0.
\]
具体地，$\Phi_{\rm eff}=(h^2+\beta)/(2r^2)-\mu/r$，逐次求导得
\[
\Phi_{\rm eff}'=-\frac{h^2+\beta}{r^3}+\frac\mu{r^2},\qquad
\Phi_{\rm eff}''=\frac{3(h^2+\beta)}{r^4}-\frac{2\mu}{r^3}.
\]
由 $\Phi_{\rm eff}'(r_c)=0$ 得 $h^2+\beta=\mu r_c$，代入二阶导数，$3\mu/r_c^3-2\mu/r_c^3=\mu/r_c^3$，所以在此处
\[
 \omega_r^2=\Phi_{\rm eff}''(r_c)=\frac\mu{r_c^3}>0,
 \qquad \omega_\varphi=\frac{|h|}{r_c^2},\qquad
 \frac{\omega_r}{\omega_\varphi}=s.
\]
近圆频率法给出相同的拱点角。此势的 Binet 方程为线性，所以上述拱点角实际上也适用于所有 $0<e_*<1$ 的有界非圆轨道。

其余两种情形须分别求解。取 $\theta=\varphi-\varphi_0$，其中 $\varphi_0$ 是任选的参考方向。当 $h^2+\beta=0$ 时，$u''=\mu/h^2$，先积分得 $u'=C_2+(\mu/h^2)\theta$，再积分为
\[
 u''=\frac\mu{h^2},\qquad
 \boxed{u=C_1+C_2\theta+\frac\mu{2h^2}\theta^2}.
\]
当 $h^2+\beta<0$ 时，令 $\sigma=\sqrt{-(1+\beta/h^2)}>0$。方程为 $u''-\sigma^2u=\mu/h^2$，取常数特解 $u_p=-\mu/(h^2\sigma^2)$，齐次根为 $\pm\sigma$，故
\[
 u''-\sigma^2u=\frac\mu{h^2},\qquad
 \boxed{u=C_1\cosh(\sigma\theta)+C_2\sinh(\sigma\theta)
       -\frac\mu{h^2\sigma^2}}.
\]
设 $u(0)=u_0$、$u'(0)=u_0'$。临界解有 $C_1=u_0,C_2=u_0'$；双曲函数解则由 $u_0=C_1-\mu/(h^2\sigma^2)$、$u_0'=\sigma C_2$ 确定 $C_1=u_0+\mu/(h^2\sigma^2),C_2=u_0'/\sigma$。只取 $u>0$ 的区间，再用 $r=1/u$ 还原轨道。临界情形的有效势为 $\Phi_{\rm eff}=-\mu/r$；负值情形的有效势在 $r\to0$ 时趋于负无穷。它们在 $r>0$ 上均单调增加，因此不存在圆轨道。向内运动可以落到中心，角动量不再形成落心屏障。若力心有实体表面，则在首次接触处终止运动。
\end{solution}



\begin{example}[潮汐势与潮汐力]\label{ex:hw8-7}

质量$M$的远方天体距行星中心$O$为$d$。取随$O$自由落体平移、坐标轴不旋转的参考系，以$OM$为极轴。求距$O$为$\varrho\ll d$、极角$\theta$处质量$\delta m$质元的潮汐势及潮汐力。只考虑远方天体造成的差异引力。

\end{example}


\begin{figure}[H]\centering
\begin{tikzpicture}[>=Stealth,line cap=round,line join=round]
\draw(0,0)circle(1.2);\fill(0,0)circle(2pt)node[below]{$O$};\fill(5,0)circle(3pt)node[right]{$M$};\coordinate(P)at(.7,.75);\fill(P)circle(2pt)node[above]{$P$};\draw[guide](0,0)--(5,0);\draw(0,0)--(P)node[midway,left]{$\varrho$};\node at(2.7,-.4){$d$};\node at(.5,.18){$\theta$};
\draw[axis](0,0)--(0,1.6)node[above]{$y$};\node[above]at(2.2,0){$x$};
\end{tikzpicture}
\caption{远方天体与行星内部质元的位置}\label{fig:hw8-7-q}
\end{figure}

\begin{solution}
图\ref{fig:hw8-7-s}中$P$到远方天体的距离为$\sqrt{d^2+\varrho^2-2d\varrho\cos\theta}$。参考系平移加速度为$GM/d^2$，相应惯性势抵消引力势的一阶项。取中心潮汐势为零，精确表达式为
\[
 U_t=GM\delta m\left[\frac1d+\frac{\varrho\cos\theta}{d^2}
 -\frac1{\sqrt{d^2+\varrho^2-2d\varrho\cos\theta}}\right].
\]

\begin{figure}[H]\centering
\begin{tikzpicture}[>=Stealth,line cap=round,line join=round]
\draw(0,0)circle(1.2);\fill(0,0)circle(2pt)node[below]{$O$};\fill(5,0)circle(3pt)node[right]{$M$};\coordinate(P)at(.7,.75);\fill(P)circle(2pt)node[above]{$P$};\draw[guide](0,0)--(5,0);\draw(0,0)--(P)node[midway,left]{$\varrho$};\node at(2.7,-.4){$d$};\node at(.5,.18){$\theta$};\draw[axis](.7,.75)--(1.5,.75)node[above]{$2Cx$};\draw[axis](.7,.75)--(.7,.2)node[right]{$-Cy$};\node at(3,1.5){$C=GM\delta m/d^3$};
\draw[axis](0,0)--(0,1.6)node[above]{$y$};\node[above]at(2.2,0){$x$};
\end{tikzpicture}
\caption{潮汐作用沿极轴拉伸、沿横向压缩}\label{fig:hw8-7-s}
\end{figure}

令 $s=\varrho/d$，距离倒数为 $d^{-1}(1-2s\cos\theta+s^2)^{-1/2}$。取 $z=-2s\cos\theta+s^2$，用二项式展开 $(1+z)^{-1/2}=1-z/2+3z^2/8+O(z^3)$。保留至 $s^2$ 阶，$-z/2=s\cos\theta-s^2/2$、$3z^2/8=(3/2)s^2\cos^2\theta+O(s^3)$，相加得
\[
 \frac1{|\vct d-\vct r|}=\frac1d
 \left[1+s\cos\theta+\frac{s^2}{2}(3\cos^2\theta-1)\right]+O(s^3/d).
\]
中心值由所选势能零点消去，一次项由随行星中心自由落体的惯性势消去，留下的二次项才描述差异引力。

展开到$\varrho^2$阶，得
\[
 \boxed{U_t=-\frac{GM\delta m}{2d^3}\varrho^2(3\cos^2\theta-1)}.
\]
记 $C=GM\delta m/d^3$。按极坐标取梯度时，径向是 $F_\varrho=-\partial U_t/\partial\varrho$，角向是 $F_\theta=-(1/\varrho)\partial U_t/\partial\theta$。分别求导：
\[
\frac{\partial U_t}{\partial\varrho}=-C\varrho(3\cos^2\theta-1),\qquad
\frac{\partial U_t}{\partial\theta}
=-\frac C2\varrho^2\cdot3\cdot2\cos\theta(-\sin\theta)
=3C\varrho^2\sin\theta\cos\theta.
\]
代入两个分量的定义，得到
\[
 \boxed{F_\varrho=\frac{GM\delta m}{d^3}\varrho(3\cos^2\theta-1),\quad
 F_\theta=-\frac{3GM\delta m}{d^3}\varrho\sin\theta\cos\theta}.
\]
沿$OM$方向的直角坐标为$x$时，同一结果为$\vct F_t=(GM\delta m/d^3)(2x,-y,-z)$：轴向拉伸，横向压缩。这里坐标轴不旋转，没有另加绕$O$的离心势；若换用随公转转动的坐标系，应连同离心、科里奥利项一起改变模型。


\end{solution}



\begin{example}[潮汐耗散与双星的稳定阈值]\label{ex:hw8-8}

孤立双星质量分别为$M,m$，其中$m$视为质点，$M$绕自身对称轴的惯量为$I$。两星绕质心作圆周公转，轨道角动量为$\ell$；$M$的自转轴与公转轴同向。总角动量$L$守恒，潮汐耗散使机械能下降，并假定演化中轨道始终近圆。求体系存在稳定同步状态所需的临界总角动量$L_0$。

\end{example}


\begin{figure}[H]\centering
\begin{tikzpicture}[scale=.9,>=Stealth,line cap=round,line join=round]
\fill(-1.5,0)circle(3pt)node[above]{$M,I$};\fill(2,0)circle(2pt)node[above]{$m$};\fill(0,0)circle(1pt)node[below]{$C$};\draw[guide](-1.5,0)--(2,0);\draw[<->](-1.5,-.6)--(2,-.6)node[midway,below]{$r$};\draw[axis](-1.8,.5)arc(160:-100:.4)node[below left]{$\Omega_s$};
\end{tikzpicture}
\caption{双星公转与其中一天体的自转}\label{fig:hw8-8-q}
\end{figure}

\begin{solution}
令约化质量$\mu_r=Mm/(M+m)$、$K=GMm$。图\ref{fig:hw8-8-s}中圆轨道有$\ell^2=\mu_rKr$，公转能量为$-K/(2r)$；自转角动量为$L-\ell$。因此在固定$L$下
\[
 E(\ell)=\frac{(L-\ell)^2}{2I}-\frac{A}{2\ell^2},\qquad A=\mu_rK^2.
\]

\begin{figure}[H]\centering
\begin{tikzpicture}[>=Stealth,line cap=round,line join=round]
\begin{scope}[shift={(0,1)}]
\fill(-1.5,0)circle(3pt)node[above]{$M,I$};\fill(1.8,0)circle(2pt)node[above]{$m$};\fill(0,0)circle(1pt)node[below]{$C$};
\draw[guide](-1.5,0)--(1.8,0);\draw[<->](-1.5,-.65)--(1.8,-.65)node[midway,below]{$r$};
\draw[axis](-1.8,.5)arc(160:-100:.4)node[below left]{$\Omega_s$};\end{scope}
\begin{scope}[shift={(4.3,0)},x=1.2cm,y=1.2cm]
\draw[axis](0,0)--(3,0)node[right]{$\ell$};\draw[axis](0,0)--(0,2.6)node[above]{$L(\ell)$};
\draw[curve,domain=.55:2.8,samples=100]plot(\x,{\x+1/(3*\x^3)});
\draw[guide](0,1.333333)--(2.8,1.333333);\node[left]at(0,1.333333){$L_0$};
\draw[thin](0,1.7)--(2.8,1.7);\node[left]at(0,1.7){$L$};
\fill(1,1.333333)circle(1.2pt);\draw[guide](1,0)--(1,1.333333);\node[below]at(1,0){$\ell_c$};
\fill(0.691380504,1.7)circle(1pt);\fill(1.621867308,1.7)circle(1pt);
\node[above left]at(0.691380504,1.7){内支};\node[above right]at(1.621867308,1.7){稳定外支};
\node at(1.6,-.65){$L(\ell)=\ell+IA/\ell^3$};\end{scope}
\end{tikzpicture}
\caption{总角动量超过最低值时，同步解分为不稳定内支与稳定外支}\label{fig:hw8-8-s}
\end{figure}

\LectureNeedspace{8\baselineskip}
驻值满足$E'=(\ell-L)/I+A/\ell^3=0$，即
\[
 L=\ell+\frac{IA}{\ell^3}.
\]
记右端为 $f(\ell)=\ell+IA\ell^{-3}$，求导得 $f'=1-3IA\ell^{-4}$、$f''=12IA\ell^{-5}>0$。由 $f'=0$ 得 $\ell_c=(3IA)^{1/4}$；又 $f$ 在 $\ell\downarrow0$ 和 $\ell\to\infty$ 时均趋于无穷，因此这是全局最小点。用 $IA/\ell_c^3=\ell_c/3$，最小值为 $L_0=4\ell_c/3$，即
\[
 \boxed{L_0=\frac43(3I\mu_rG^2M^2m^2)^{1/4}}.
\]
若$L>L_0$，有两个驻值。$E''=1/I-3A/\ell^4$说明$\ell>\ell_c$的一支是稳定局部极小，内侧一支是极大。$L=L_0$时两驻值合并，二阶回复消失，不是稳定极小；$L<L_0$则没有同步驻值。

在临界点$E'=E''=0$，而$\dd^3E/\dd\ell^3=12A/\ell_c^5>0$，能量的局部主导项是三次项，故该合并点不是局部极小。所以稳定同步状态要求严格的$L>L_0$。图\ref{fig:hw8-8-s}将两支及合并阈值标出。

稳定支满足同步条件。由 $\ell^2=\mu_rKr$，有
\[
\frac{L-\ell}{I}=\frac{A}{\ell^3}
=\frac{\mu_rK^2}{(\mu_rKr)^{3/2}}
=\sqrt{\frac{K}{\mu_rr^3}}=\sqrt{\frac{G(M+m)}{r^3}},
\]
恰为公转角速度。稳定支 $\ell^4>3IA$ 又给出 $IA/\ell^3<\ell/3$，所以 $L<4\ell/3$，即 $\ell>3L/4$。对应临界距离为 $r_c=\ell_c^2/(\mu_rK)=\sqrt{3I/\mu_r}$。实体天体还须满足不相交的几何条件；超过总角动量阈值只保证动力学平衡存在。


\end{solution}



\section{习题与解答}

\begin{exercise}[在半长轴以内停留多久]\label{sup8:problem-exercise-2}\label{sup8:ex:d8-h2}
在引力参数为 $\mu$ 的场中，质点沿半长轴为 $a$、偏心率为 $e=3/5$ 的椭圆运动。求从近心点向外飞到 $r=a$ 所需的时间，并求一周期中位于 $r<a$ 区域内的时间占比。
\end{exercise}
\begin{figure}[H]\centering
\begin{tikzpicture}[scale=.9]
 \draw[curve] (-1.8,0) ellipse (3 and 2.4);\draw[guide] (0,0) circle (3);
 \fill (0,0) circle (1.6pt) node[below right] {$O$};\fill (1.2,0) circle (1.5pt) node[right] {$P_p$};
 \fill (-1.8,2.4) circle (1.5pt) node[above] {$P_+$};\fill (-1.8,-2.4) circle (1.5pt) node[below] {$P_-$};
 \draw (0,0)--(-1.8,2.4) node[midway,right] {$a$};
 \node at (2.8,-2.8) {虚线：$r=a$};
\end{tikzpicture}
\caption{圆 $r=a$ 与椭圆交于 $P_+,P_-$，近心点位于圆内。}\label{sup8:fig:d8-h2-q}
\end{figure}
\begin{solution}
以近心点为 $t=0$，令 $n=\sqrt{\mu/a^3}$。面积定律对应的时间积分为
\[
t(\mathcal E)=\frac1n\int_0^{\mathcal E}(1-e\cos s)\,\dd s
=\frac1n[s-e\sin s]_0^{\mathcal E}
=\frac{\mathcal E-e\sin\mathcal E}{n}.
\]
图\ref{sup8:fig:d8-h2-q} 的 $P_+$ 在向外分支，$r=a(1-e\cos\mathcal E)=a$ 给出 $\mathcal E=\pi/2$，故
\[
 t_+=\frac1n\left(\frac\pi2-\frac35\right).
\]
向内分支的 $P_-$ 对应 $\mathcal E=3\pi/2$，相应时刻为
\[
 t_-=\frac1n\left(\frac{3\pi}2+\frac35\right).
\]
一周期中，$r<a$ 的时间区间为 $[0,t_+)$ 与 $(t_-,T]$，不是 $P_+$ 到 $P_-$ 的长弧。
\begin{figure}[H]\centering
\begin{tikzpicture}[x=1.05cm]
 \draw[axis] (0,0)--(6.8,0) node[right] {$\mathcal E$};
 \draw[line width=3pt] (0,0)--(1.5708,0);\draw[line width=3pt] (4.7124,0)--(6.2832,0);
 \foreach \x/\lab in {0/0,1.5708/{\pi/2},3.1416/\pi,4.7124/{3\pi/2},6.2832/{2\pi}} {\draw (\x,.12)--(\x,-.12) node[below=5pt] {$\lab$};}
 \node at (.78,.55) {$r<a$};\node at (3.14,.55) {$r>a$};\node at (5.5,.55) {$r<a$};
 \node[align=center] at (3.14,-1.25) {$\dd t=\dfrac{1-e\cos\mathcal E}{n}\dd\mathcal E$\\偏近点角等间隔不对应等时间};
\end{tikzpicture}
\caption{以偏近点角标记圆内、圆外的区间，时间需按 $1-e\cos\mathcal E$ 加权。}\label{sup8:fig:d8-h2-s}
\end{figure}
由图\ref{sup8:fig:d8-h2-s}，两个圆内区间关于近心点对称，总时间为 $2t_+$。因此
\[
 \boxed{\frac{\Delta t_{r<a}}T=\frac{2t_+}{2\pi/n}
 =\frac12-\frac e\pi=\frac12-\frac3{5\pi}\simeq0.3090}.
\]
这两个区间合占偏近点角的一半，却用去不足半个周期。在圆内 $1-e\cos\mathcal E<1$，同样的偏近点角增量对应较短时间；面积速度始终为常量 $h/2$。区间端点是否包含，不影响这一时间占比。
\end{solution}


\begin{exercise}[抛物线的一段飞行时间]\label{sup8:problem-exercise-9}\label{sup8:ex:d8-h9}
引力抛物线的近心距离为 $q$，引力参数为 $\mu$。求质点从近心点飞到 $r=2q$ 的时间及该处速度分量；再求它从向内运动的 $r=2q$ 处，经近心点到向外运动的 $r=2q$ 处的总时间。
\end{exercise}
\begin{figure}[H]\centering
\begin{tikzpicture}[scale=1.2]
 \draw[curve,domain=-110:110,samples=130,variable=\s] plot ({2*cos(\s)/(1+cos(\s))},{2*sin(\s)/(1+cos(\s))});
 \fill (0,0) circle (1.4pt) node[below right] {$O$};\fill (1,0) circle (1.4pt) node[right] {$P_p$};
 \fill (0,2) circle (1.4pt) node[right] {$P_+$};\fill (0,-2) circle (1.4pt) node[right] {$P_-$};
 \draw[guide] (0,-2)--(0,2);\draw[<->,>=Stealth] (0,-.5)--(1,-.5) node[midway,below] {$q$};
 \draw[axis] plot[domain=45:75,samples=20] ({2*cos(\x)/(1+cos(\x))},{2*sin(\x)/(1+cos(\x))});
\end{tikzpicture}
\caption{抛物线 $r=2q/(1+\cos f)$，$P_-,P_+$ 分别对应 $f=-\pi/2,\pi/2$。}\label{sup8:fig:d8-h9-q}
\end{figure}
\begin{solution}
抛物线 $e=1$，故 $p=2q,h=\sqrt{2\mu q}$。图\ref{sup8:fig:d8-h9-q} 中 $r=2q$ 给出 $\cos f=0$，向外分支取 $f=\pi/2$。

由面积定律 $\dd t=r^2\dd f/h$。为完成积分，令 $D=\tan(f/2)$，其导数为 $\dd D/\dd f=\tfrac12\sec^2(f/2)=\tfrac12(1+D^2)$；再用 $1+\cos f=2/(1+D^2)$，则
\[
 r=q(1+D^2),\qquad \dd f=\frac{2\dd D}{1+D^2}.
\]
以经过近心点的时刻为零，
\[
 t=\frac{2q^2}h\int_0^D(1+s^2)\dd s
 =\frac{2q^2}h\left[s+\frac{s^3}3\right]_0^D
 =\sqrt{\frac{2q^3}\mu}\left(D+\frac{D^3}3\right).
\]
在 $P_+$ 处 $D=1$，所以
\[
 \boxed{t_{p\to+}=\frac43\sqrt{\frac{2q^3}\mu}}.
\]
\begin{figure}[H]\centering
\begin{tikzpicture}
 \begin{scope}[xshift=-2.5cm]
  \draw[axis] (-.3,0)--(2,0) node[right] {$D$};\draw[axis] (0,-.2)--(0,2) node[above] {$t/\sqrt{2q^3/\mu}$};
  \draw[curve,domain=0:1.1,samples=60] plot (\x,{\x+\x*\x*\x/3});
  \draw[guide] (1,0)--(1,1.333)--(0,1.333);\node[below] at (1,0) {$1$};\node[left] at (0,1.333) {$4/3$};
 \end{scope}
 \begin{scope}[xshift=2.9cm]
  \fill (0,0) circle (1.5pt) node[below right] {$P_+$};
  \draw[axis] (0,0)--(0,1.2) node[above] {$v_r$};\draw[axis] (0,0)--(-1.2,0) node[left] {$v_\varphi$};
  \draw[axis,line width=.8pt] (0,0)--(-1.2,1.2) node[above left] {$\vct v$};
  \draw[guide] (-1.2,0)--(-1.2,1.2)--(0,1.2);
 \end{scope}
\end{tikzpicture}
\caption{左图为抛物飞行时间；右图在 $P_+$ 分解速度，两个分量大小相等。}\label{sup8:fig:d8-h9-s}
\end{figure}
在 $P_+$，由径向、横向速度公式，
\[
 v_r=\frac\mu h=\sqrt{\frac\mu{2q}},\qquad
 v_\varphi=\frac h{2q}=\sqrt{\frac\mu{2q}}.
\]
图\ref{sup8:fig:d8-h9-s} 中向上和向左两个分量合成的速率为 $\sqrt{\mu/q}$，也等于零能量式 $v^2=2\mu/r$ 在 $r=2q$ 处的结果。

$t(D)$ 是奇函数，$P_-$ 对应 $D=-1$。从 $P_-$ 经近心点到 $P_+$ 的总时间为
\[
 \boxed{t_{-\to+}=\frac83\sqrt{\frac{2q^3}\mu}}.
\]
抛物轨道没有有限周期，这段飞行时间却是有限的。椭圆周期在 $a\to\infty$ 时发散，有限轨道弧段的时间仍须由局部面积积分求出。
\end{solution}


\begin{exercise}[双星同时失去部分质量]\label{deep8:problem-exercise-4}\label{deep8:ex:d8-h4}
两颗质量均为 $m$ 的质点恒星，间距为 $d$，原来绕质心作圆周运动。某瞬间，两星各自以球对称方式失去相同质量比例，剩余质量均为 $\eta m$，$0<\eta<1$。假定抛出物迅速离开，此后不再施加引力，失质量过程无反冲，剩余恒星的瞬时位置、速度连续。求新相对轨道；确定双星仍束缚的条件。
\end{exercise}
\begin{figure}[H]\centering
\begin{tikzpicture}
 \fill (-1.6,0) circle (2pt);\fill (1.6,0) circle (2pt);
 \node at (-2.7,-.75) {$m\to\eta m$};\node at (2.7,-.75) {$m\to\eta m$};
 \fill (0,0) circle (1.3pt) node[below] {$C$};\draw[<->,>=Stealth] (-1.6,.6)--(1.6,.6) node[midway,above] {$d$};
 \draw[axis] (-1.6,0)--(-1.6,-1.2) node[left] {$v$};\draw[axis] (1.6,0)--(1.6,1.2) node[right] {$v$};
 \foreach \a in {35,145,215,325} {\draw[axis] (-1.6,0)--++(\a:.65);\draw[axis] (1.6,0)--++(\a:.65);}
\end{tikzpicture}
\caption{两星同时、对称失质量，粗略辐射箭头仅表示无净反冲的抛出过程。}\label{deep8:fig:d8-h4-q}
\end{figure}
\begin{solution}
图\ref{deep8:fig:d8-h4-q} 中每星的原圆周半径为 $d/2$。由 $mv^2/(d/2)=Gm^2/d^2$，
\[
 v=\sqrt{\frac{Gm}{2d}},\qquad v_{\rm rel}=2v=\sqrt{\frac{2Gm}d}.
\]
失质量后，相对速度不变，但引力参数从 $\mu_0=2Gm$ 变为 $\mu_1=2G\eta m$。新守恒量为
\[
 h=dv_{\rm rel}=\sqrt{2Gmd},\qquad
 \varepsilon_+=\frac{Gm}d-\frac{2G\eta m}d
 =\frac{Gm}d(1-2\eta).
\]
这里守恒的是失质量之后的 $h,\varepsilon_+$。剩余系统的约化质量也改变，故不能把总角动量与单位约化质量角动量混作同一不变量。
\begin{figure}[H]\centering
\begin{tikzpicture}[scale=.85]
 \draw[curve] (-1,0) ellipse (2 and 1.732);\draw[guide] (0,0) circle (1);
 \fill (0,0) circle (1.6pt) node[below] {$O$};\fill (1,0) circle (1.5pt) node[right] {$P_0=P_p$};
 \draw[axis] (1,0)--(1,1.2) node[right] {$v_{\rm rel}$};
 \node[align=left,anchor=west] at (2.8,1.4) {$\eta>1/2$：椭圆\\$\eta=1/2$：抛物线\\$\eta<1/2$：双曲线};
 \node at (-1,-2.2) {实线示例：$\eta=2/3$};
\end{tikzpicture}
\caption{失质量后引力减弱，相对运动从原圆上的点向外离开。虚线为原相对圆轨道。}\label{deep8:fig:d8-h4-s}
\end{figure}
由能量符号，双星仍束缚的条件是 $\eta>1/2$。在此范围内
\[
 a=\frac{\eta d}{2\eta-1},\qquad
 e=\frac{1-\eta}{\eta},\qquad
 r_p=d,\qquad r_a=\frac d{2\eta-1}.
\]
图\ref{deep8:fig:d8-h4-s} 中初始点为近心点，因为原相对圆周速率大于新引力参数对应的圆周速率。新周期为
\[
 T_1=2\pi\sqrt{\frac{a^3}{2G\eta m}},\qquad
 \frac{T_1}{T_0}=\frac{\eta}{(2\eta-1)^{3/2}},
 \quad T_0=2\pi\sqrt{\frac{d^3}{2Gm}}.
\]
当 $\eta=1/2$ 时 $e=1$，恰好抛物逃逸；当 $\eta<1/2$ 时 $e=(1-\eta)/\eta>1$，沿双曲线分离。三种情况下最近距离均为 $d$。这些结果依赖题设的瞬时失质量模型；若抛出物仍在附近或有反冲，需重新计算初态和引力场。
\end{solution}


\begin{exercise}[穿过均匀球体的直通隧道]\label{deep8:problem-exercise-5}\label{deep8:ex:d8-h5}
质量为 $M$、半径为 $R$ 的均匀球体中有一条穿过球心的极细直隧道，隧道对引力场的影响忽略。质点从球外轴线上距球心 $r_0>R$ 处静止释放，无摩擦地穿过球体。求一次完整往返的周期，分别写出球外和球内所用的运动规律。
\end{exercise}
\begin{figure}[H]\centering
\begin{tikzpicture}
 \draw[curve,fill=black!6] (0,0) circle (1.5);\draw[line width=3pt,white] (-1.5,0)--(1.5,0);\draw[guide] (-3.8,0)--(4,0);
 \fill (0,0) circle (1.3pt) node[below] {$O$};\fill (3.3,0) circle (1.6pt) node[above] {$P_0$};
 \node[above] at (-1.5,0) {$-R$};\node[above] at (1.5,0) {$R$};
 \draw[axis] (3.1,0)--(2.4,0);\draw[<->,>=Stealth] (0,-1.9)--(3.3,-1.9) node[midway,below] {$r_0$};
\end{tikzpicture}
\caption{质点沿固定直径运动，以带符号坐标 $x$ 区分球心两侧。}\label{deep8:fig:d8-h5-q}
\end{figure}
\begin{solution}
令 $\mu=GM$。在图\ref{deep8:fig:d8-h5-q} 的右侧球外，$x>R$，单位质量力为 $-\mu/x^2$。从 $r_0$ 落到 $R$ 的时间和入球速率由例\ref{sup8:ex:d8-e4} 的径向积分给出：
\[
 t_{\rm out}=\sqrt{\frac{r_0^3}{2\mu}}(\eta_R+\sin\eta_R\cos\eta_R),
 \quad\eta_R=\arccos\sqrt{\frac R{r_0}},
\]
\[
 v_R=\sqrt{2\mu\left(\frac1R-\frac1{r_0}\right)}.
\]
球内半径 $|x|$ 的包围质量为 $M|x|^3/R^3$，球壳部分的合力为零，故在整个 $-R<x<R$ 区间
\[
 \ddot x=-\frac\mu{R^3}x=-\omega^2x,\qquad
 \omega=\sqrt{\frac\mu{R^3}}.
\]
取刚进入右端表面的时刻为球内时间 $s=0$。方程特征根为 $\pm i\omega$，故通解 $x=C\cos\omega s+D\sin\omega s$，求导为 $\dot x=-C\omega\sin\omega s+D\omega\cos\omega s$。由 $x(0)=R$ 得 $C=R$，由 $\dot x(0)=-v_R$ 得 $D=-v_R/\omega$，于是
\[
 x(s)=R\cos\omega s-\frac{v_R}{\omega}\sin\omega s
     =A\cos(\omega s+\delta),
\]
其中
\[
 A=\sqrt{R^2+\frac{v_R^2}{\omega^2}}
 =R\sqrt{3-\frac{2R}{r_0}},\qquad
 \delta=\arctan\frac{v_R}{\omega R}\in(0,\pi/2).
\]
其中 $A>R$ 是球内简谐解的形式振幅；质点一到球面就转入球外运动，这段简谐解的使用区间只到球面为止。
\begin{figure}[H]\centering
\begin{tikzpicture}[x=1.5cm,y=1.5cm]
 \draw[axis] (-2.8,0)--(2.9,0) node[right] {$x/R$};\draw[axis] (0,-1.35)--(0,1.45) node[above] {$f/(\mu/R^2)$};
 \draw[curve,domain=-2.5:-1,samples=60] plot (\x,{1/(\x*\x)});
 \draw[curve] (-1,1)--(1,-1);
 \draw[curve,domain=1:2.5,samples=60] plot (\x,{-1/(\x*\x)});
 \draw[guide] (-1,0)--(-1,1);\draw[guide] (1,0)--(1,-1);
 \node[below] at (-1,0) {$-1$};\node[above] at (1,0) {$1$};
 \node at (-1.9,1.2) {球外};\node[above right] at (.25,.45) {球内线性力};\node at (1.9,-1.2) {球外};
\end{tikzpicture}
\caption{带符号的引力加速度在两个表面连续，球内与球外分别采用不同函数。}\label{deep8:fig:d8-h5-s}
\end{figure}
图\ref{deep8:fig:d8-h5-s} 说明球内不能延用反平方力。首次到达 $x=-R$ 时相位为 $\pi-\delta$，故穿越球体的时间
\[
 t_{\rm in}=\frac{\pi-2\delta}{\omega}
 =\frac2\omega\arctan\frac{\omega R}{v_R}.
\]
由对称性，从 $+r_0$ 到 $-r_0$ 需 $2t_{\rm out}+t_{\rm in}$；再返回原位置与原速度，周期为
\[
 \boxed{T=4t_{\rm out}+2t_{\rm in}}.
\]
当 $r_0\downarrow R$ 时，$t_{\rm out}\to0,v_R\to0$，从而 $t_{\rm in}\to\pi/\omega$、$T\to2\pi/\omega$，与从表面静止释放的全程简谐运动一致。一般 $r_0>R$ 时，只有球内的一段是简谐运动，不能把 $2\pi/\omega$ 当作完整周期。
\end{solution}


\begin{exercise}[引力轨道的速度端点圆]\label{deep8:problem-exercise-8}\label{deep8:ex:d8-h8}
在引力参数为 $\mu$ 的反平方引力场中，取近心点方向为 $+x$，运动方向为逆时针，$h>0$。求速度分量 $v_x,v_y$ 关于真近点角 $f$ 的表达式，证明速度矢量端点落在一个圆上。用此圆解释椭圆、抛物线和双曲线的区别，并求 $e=1/2$ 椭圆的最大、最小速率。
\end{exercise}
\begin{figure}[H]\centering
\begin{tikzpicture}[scale=1.2]
 \draw[axis] (-1.2,0)--(2.5,0) node[right] {近心点方向 $x$};\draw[axis] (0,-.4)--(0,3.1) node[above] {$y$};
 \coordinate (P) at (1,1.732);\draw[axis] (0,0)--(P) node[midway,right] {$r$};
 \fill (0,0) circle (1.3pt) node[below left] {$O$};\fill (P) circle (1.3pt) node[right] {$P$};
 \draw[axis] (P)--++(.4,.693) node[above] {$v_r$};\draw[axis] (P)--++(-1.039,.6) node[left] {$v_\varphi$};
 \draw (.55,0) arc (0:60:.55) node[midway,right] {$f$};
\end{tikzpicture}
\caption{速度沿径向与横向分解，基矢 $\vct e_r,\vct e_\varphi$ 随 $f$ 转动。}\label{deep8:fig:d8-h8-q}
\end{figure}
\begin{solution}
图\ref{deep8:fig:d8-h8-q} 中
\[
 \vct e_r=(\cos f,\sin f),\qquad
 \vct e_\varphi=(-\sin f,\cos f).
\]
将式\eqref{sup8:eq:d8-velocities} 的两个分量组合为 $\vct v=v_r\vct e_r+v_\varphi\vct e_\varphi$，得
\begin{align*}
 v_x&=\frac\mu h[e\sin f\cos f-(1+e\cos f)\sin f]
     =-\frac\mu h\sin f,\\
 v_y&=\frac\mu h[e\sin^2f+(1+e\cos f)\cos f]
     =\frac\mu h(e+\cos f).
\end{align*}
故速度端点满足
\begin{equation}
 \boxed{v_x^2+\left(v_y-\frac{\mu e}h\right)^2
 =\left(\frac\mu h\right)^2}.
\end{equation}
圆心为 $C_v=(0,\mu e/h)$，半径为 $\mu/h$。这是速度空间中的圆，不是位置空间中的轨道。
\begin{figure}[H]\centering
\begin{tikzpicture}[scale=1.6]
 \draw[axis] (-1.35,0)--(1.55,0) node[right] {$v_x/(\mu/h)$};\draw[axis] (0,-.8)--(0,1.9) node[above] {$v_y/(\mu/h)$};
 \draw[curve] (0,.5) circle (1);\fill (0,.5) circle (1.1pt) node[right] {$C_v$};
 \fill (0,0) circle (1.1pt) node[below left] {$O_v$};
 \fill (0,1.5) circle (1.1pt) node[right] {$P_p^v$};\fill (0,-.5) circle (1.1pt) node[right] {$P_a^v$};
 \draw[axis] (0,0)--(-.866,1) node[above left] {$\vct v(f=\pi/3)$};
 \draw[guide] (0,.5)--(-.866,1) node[midway,above] {$1$};
 \node at (0,-1.2) {图示 $e=1/2$};
\end{tikzpicture}
\caption{速度端点圆。速率是从速度原点 $O_v$ 到圆上端点的距离，而非圆半径。}\label{deep8:fig:d8-h8-s}
\end{figure}
图\ref{deep8:fig:d8-h8-s} 中 $e<1$，原点 $O_v$ 在圆内，质点在一周期中走完整圆。最大、最小速率分别为
\[
 v_{\max}=\frac\mu h(1+e),\qquad
 v_{\min}=\frac\mu h(1-e).
\]
对 $e=1/2$，二者为 $3\mu/(2h)$ 与 $\mu/(2h)$，对应近心点与远心点。

$e=1$ 时，速度圆经过原点。沿抛物线飞至无穷远，速度端点趋近原点，有限时刻则仍有非零速率。$e>1$ 时，原点在圆外，物理轨道只覆盖满足 $1+e\cos f>0$ 的圆弧；圆弧两端是入射、出射渐近速度，模均为 $v_\infty$。

双曲线的速率只能在这段物理圆弧上取值。整圆上离原点最近的点不在该圆弧内；沿实际轨道，速率的下确界是无穷远的 $v_\infty$，有限位置的速率都更大。
\end{solution}


\begin{exercise}[反幂次中心力的圆轨道稳定性（选做）]\label{deep8:problem-exercise-10}\label{deep8:ex:d8-h10}
质量为 $m$ 的质点受 $\vct F=-k\vct e_r/r^n$，$k>0$，$n$ 为实数。对于存在的圆轨道，求角频率、固定角动量下的小径向振动频率及相邻远、近拱点的角间隔。分别讨论 $n<3,n=3,n>3$，说明近圆轨道闭合所需的频率条件。
\end{exercise}
\begin{figure}[H]\centering
\begin{tikzpicture}
 \draw[guide] (0,0) circle (2);\draw[guide] (0,0)--(2.3,0);
 \fill (0,0) circle (1.6pt) node[below] {$O$};\fill (2.3,0) circle (1.5pt) node[right] {$P$};
 \draw[axis] (2.3,0)--(2.3,1.2) node[right] {$h/r$};\draw[axis] (2.3,0)--(1.2,0) node[midway,above] {$k/(mr^n)$};
 \draw[guide] (2,0)--(2,-.7);\draw[guide] (2.3,0)--(2.3,-.7);
 \draw[<->,>=Stealth] (2,-.5)--(2.3,-.5) node[below] {$\xi$};
 \node at (-2.5,2) {虚线：$r=r_c$};
\end{tikzpicture}
\caption{虚线圆是未扰动圆轨道，$P$ 是径向偏移后的瞬时位置。固定 $h$ 时，横向速率随半径改变。}\label{deep8:fig:d8-h10-q}
\end{figure}
\begin{solution}
令 $K=k/m$。圆轨道满足 $r_c\omega_\varphi^2=K/r_c^n$，故
\[
 \omega_\varphi^2=\frac K{r_c^{n+1}},\qquad
 h^2=r_c^4\omega_\varphi^2=Kr_c^{3-n}.
\]
当 $n\ne3$ 且 $h\ne0$ 时，圆轨道半径由上式唯一确定：$r_c=(h^2/K)^{1/(3-n)}$。保持 $h$ 不变，径向方程为
\[
 \ddot r=\frac{h^2}{r^3}-\frac K{r^n}.
\]
令 $F_r(r)=h^2r^{-3}-Kr^{-n}$，保持 $h$ 不变求导得 $F_r'=-3h^2r^{-4}+nKr^{-n-1}$。再令 $r=r_c+\xi$，按图\ref{deep8:fig:d8-h10-q} 展开 $F_r(r_c+\xi)=F_r(r_c)+F_r'(r_c)\xi+O(\xi^2)$。圆轨道条件使 $F_r(r_c)=0$，所以一阶方程为
\[
 \ddot\xi=\left(-\frac{3h^2}{r_c^4}+\frac{nK}{r_c^{n+1}}\right)\xi
 =-(3-n)\omega_\varphi^2\xi.
\]
因此 $n<3$ 时
\[
 \boxed{\omega_r=\sqrt{3-n}\,\omega_\varphi}.
\]
由一个远拱点到下一个近拱点需半个径向周期；在近圆极限中，角速度取零阶值，故
\[
 \boxed{\Delta\varphi_{a\to p}=\frac\pi{\sqrt{3-n}}},\qquad
 \Delta\varphi_{p\to p}=\frac{2\pi}{\sqrt{3-n}}.
\]
\begin{figure}[H]\centering
\begin{tikzpicture}[x=1.4cm,y=.85cm]
 \draw[axis] (-1.2,0)--(4.8,0) node[right] {$n$};\draw[axis] (0,-2.2)--(0,4.5) node[above] {$\omega_r^2/\omega_\varphi^2$};
 \draw[curve] (-1,4)--(4.5,-1.5);\fill (3,0) circle (1.4pt) node[above right] {$n=3$};
 \fill (2,1) circle (1.4pt) node[above right] {引力 $n=2$};
 \fill (-1,4) circle (1.4pt) node[above left] {线性力 $n=-1$};
 \node at (1.2,-1.4) {稳定：$n<3$};\node at (4,-2) {不稳定};
\end{tikzpicture}
\caption{频率平方比为 $3-n$。$n=3$ 的零值需用完整径向方程另行判断。}\label{deep8:fig:d8-h10-s}
\end{figure}
图\ref{deep8:fig:d8-h10-s} 中 $n>3$ 时，令 $\lambda=\sqrt{n-3}\,\omega_\varphi$，扰动方程为 $\ddot\xi-\lambda^2\xi=0$。特征根 $\pm\lambda$ 给出 $\xi=C\cosh\lambda t+D\sinh\lambda t$；由初值 $C=\xi(0)$、$D=\dot\xi(0)/\lambda$，一般初态含增长分支，圆轨道不稳定。
$n=3$ 时圆轨道存在须 $h^2=K$；在这个特殊 $h$ 下，完整径向方程对所有 $r$ 都给出 $\ddot r=0$。积分先得 $\dot r=\dot r(0)$，再得 $r=r(0)+\dot r(0)t$。纯位置扰动且 $\dot r=0$ 可停留在另一圆上，但任意微小径向速度会导致持续漂移。因此它没有恢复振动，也不是对一般小扰动稳定的圆轨道。若 $h^2\ne K$，则不存在圆轨道。

近圆轨道在一阶近似下闭合，要求 $\omega_r/\omega_\varphi=\sqrt{3-n}$ 为有理数，使径向相位和方位角在有限圈后同时复原。这个条件只来自近圆展开，不能推出任意振幅轨道均闭合。引力 $n=2$ 和线性力 $n=-1$ 的闭合性质可由各自的完整解确认，分别对应频率比 $1$ 和 $2$。
\end{solution}


\begin{exercise}[同一椭圆，不同力心（选做）]\label{deep8:problem-exercise-11}\label{deep8:ex:d8-h11}
已知位置轨道为半长轴 $A$、半短轴 $B$ 的椭圆，$A>B>0$。分别假设中心力的力心位于椭圆几何中心 $C$ 或一个焦点 $O$，且相应的单位质量角动量为非零常量 $h_C,h_O$。推导两种情况下沿该轨道所需的力与距离的关系。
\end{exercise}
\begin{figure}[H]\centering
\begin{tikzpicture}
 \draw[curve] (0,0) ellipse (3 and 2);\draw[guide] (-3.4,0)--(3.5,0);
 \fill (0,0) circle (1.6pt) node[below] {$C$};\fill (2.236,0) circle (1.6pt) node[below] {$O$};
 \coordinate (P) at (1.5,1.732);\fill (P) circle (1.6pt) node[above] {$P$};
 \draw[axis] (0,0)--(P) node[midway,left] {$r_C$};\draw[axis] (2.236,0)--(P) node[midway,right] {$r_O$};
 \node at (-1.2,-2.5) {$CO=\sqrt{A^2-B^2}$};
\end{tikzpicture}
\caption{两种假设分别以 $C$ 或 $O$ 为力心，不能在同一次运动中同时要求这两个角动量守恒。}\label{deep8:fig:d8-h11-q}
\end{figure}
\begin{solution}
先以图\ref{deep8:fig:d8-h11-q} 的 $C$ 为坐标原点。用几何参数 $\psi$ 表示椭圆：
\[
 x=A\cos\psi,\qquad y=B\sin\psi.
\]
这里尚未假设 $\psi$ 随时间匀速。计算角动量，有
\[
 h_C=x\dot y-y\dot x=AB\dot\psi.
\]
因为 $h_C$ 为常量，才得到 $\dot\psi=h_C/(AB)$ 为常量，因而 $\ddot\psi=0$。链式法则逐次给出 $\dot x=-A\sin\psi\,\dot\psi$、$\ddot x=-A\cos\psi\,\dot\psi^2-A\sin\psi\,\ddot\psi$，以及 $\dot y=B\cos\psi\,\dot\psi$、$\ddot y=-B\sin\psi\,\dot\psi^2+B\cos\psi\,\ddot\psi$。代入常角速度，得到
\[
 \ddot x=-\left(\frac{h_C}{AB}\right)^2x,\qquad
 \ddot y=-\left(\frac{h_C}{AB}\right)^2y.
\]
所以力心在中心时，沿轨道所需的力为
\begin{equation}
 \boxed{\vct F_C=-\frac{mh_C^2}{A^2B^2}\vct r_C}.
\end{equation}
\begin{figure}[H]\centering
\begin{tikzpicture}
 \begin{scope}[xshift=-2.4cm]
  \draw[curve] (0,0) ellipse (1.5 and 1);
  \fill (0,0) circle (1.5pt) node[below] {$C$};\fill (.75,.866) circle (1.5pt) node[above] {$P$};
  \draw[axis] (.75,.866)--(.225,.26) node[midway,left] {$\vct F_C$};
  \node at (0,-1.55) {$|F_C|\propto r_C$};
 \end{scope}
 \begin{scope}[xshift=2.6cm]
  \draw[curve] (0,0) ellipse (1.5 and 1);
  \fill (1.118,0) circle (1.5pt) node[below] {$O$};\fill (.75,.866) circle (1.5pt) node[above] {$P$};
  \draw[axis] (.75,.866)--(1.008,.26) node[midway,right] {$\vct F_O$};
  \node at (0,-1.55) {$|F_O|\propto r_O^{-2}$};
 \end{scope}
\end{tikzpicture}
\caption{在同一个几何点 $P$，两种假设下的力方向分别指向中心与焦点，时间参数化也不同。}\label{deep8:fig:d8-h11-s}
\end{figure}
再把力心移到焦点 $O$，并从与该焦点最近的顶点量真近点角 $f$。该椭圆的极坐标式为
\[
 u=\frac1r=\frac{1+e\cos f}{p},\qquad
 e=\sqrt{1-\frac{B^2}{A^2}},\quad p=\frac{B^2}A.
\]
对真近点角逐次求导，$u'=-e\sin f/p$、$u''=-e\cos f/p$，所以 $u''+u=[-e\cos f+(1+e\cos f)]/p=1/p$。由 Binet 方程，
\[
 \frac{F_O}m=-h_O^2u^2(u''+u)=-\frac{h_O^2}{pr_O^2},
\]
即
\begin{equation}
 \boxed{\vct F_O=-\frac{mh_O^2A}{B^2r_O^2}\vct e_{r_O}}.
\end{equation}
图\ref{deep8:fig:d8-h11-s} 的力方向和两种距离定义进入了不同的运动方程。所得关系只由这条轨道采样到的距离确定：中心情形 $r_C\in[B,A]$，焦点情形 $r_O\in[A(1-e),A(1+e)]$。若要断言轨道范围以外也服从同一力律，须另加模型假设，不能仅凭一条观测椭圆推出。
\end{solution}


\begin{exercise}[有限半径行星的引力辅助（选做）]\label{deep8:problem-exercise-12}\label{deep8:ex:d8-h12}
在一次短时遭遇中，行星在某惯性系中近似以恒速度 $\vct V=V\vct e_x$ 运动。行星引力参数为 $\mu_p$，吸收半径为 $R$；飞行器质量可忽略。在行星系远处，飞行器入射速度为 $\vct u_{\rm in}=-u\vct e_x$，$u>0$。飞行器由双曲线散射后离开，轨道不接触表面。忽略遭遇期间太阳等外场对相对运动的影响。

求飞行器在原惯性系中的单位质量动能变化与偏转角 $\chi$ 的关系；求不接触表面的最大偏转角上界及可获得的能量增量上确界，说明能量的来源。
\end{exercise}
\begin{figure}[H]\centering
\begin{tikzpicture}
 \begin{scope}[xshift=-2.7cm]
  \draw[curve,fill=black!7] (0,0) circle (.7);\fill (0,0) circle (1.5pt) node[below] {行星};
  \draw[axis] (0,0)--(1.6,0) node[above] {$\vct V$};
  \draw[axis] (2,1.5)--(.8,1.5) node[midway,above] {$\vct u_{\rm in}$};
  \node at (0,-1.35) {相对入射从行星前方而来};
 \end{scope}
 \begin{scope}[xshift=2.5cm]
  \draw[axis] (0,0)--(-1.7,0) node[left] {$\vct u_{\rm in}$};
  \draw[axis] (0,0)--(.85,1.472) node[above] {$\vct u_{\rm out}$};
  \draw (0,0)++(180:.6) arc (180:60:.6);
  \node at (-.1,.85) {$\chi$};\fill (0,0) circle (1.5pt);
 \end{scope}
\end{tikzpicture}
\caption{左图给出行星运动与相对入射的方向，右图单独画相对速度的偏转。}\label{deep8:fig:d8-h12-q}
\end{figure}
\begin{solution}
\begingroup
\setlength{\parskip}{2pt}
\setlength{\abovedisplayskip}{2pt}\setlength{\belowdisplayskip}{2pt}
\setlength{\belowdisplayshortskip}{2pt}
在图\ref{deep8:fig:d8-h12-q} 的行星系中，引力是静态保守中心力，故入射、出射的远处相对速率相同，均为 $u$。取偏转后的横向分量为正，
\[
 \vct u_{\rm out}=u(-\cos\chi,\sin\chi).
\]
原惯性系中的速度要加上 $\vct V$：
\[
 \vct v_{\rm in}=(V-u,0),\qquad
 \vct v_{\rm out}=(V-u\cos\chi,u\sin\chi).
\]
所以
\begin{align*}
 \Delta K/m
 &=\frac12\left(|\vct V+\vct u_{\rm out}|^2
                  -|\vct V+\vct u_{\rm in}|^2\right)\\
 &=\vct V\cdot(\vct u_{\rm out}-\vct u_{\rm in})
 =\boxed{Vu(1-\cos\chi)}.
\end{align*}
入射与出射的相对动能 $u^2/2$ 相同，原惯性系中的差别来自速度相加的交叉项。若向另一侧偏转，$y$ 分量变号，能量变化保持相同。
\begin{figure}[H]\centering
\begin{tikzpicture}[scale=1.4]
 \coordinate (O) at (0,0);\coordinate (C) at (1.7,0);\coordinate (I) at (.4,0);\coordinate (F) at (2.35,1.126);
 \draw[guide] (C) circle (1.3);
 \draw[axis] (O)--(C) node[pos=.75,below=12pt] {$\vct V$};
 \draw[axis] (C)--(I) node[midway,above] {$\vct u_{\rm in}$};
 \draw[axis] (C)--(F) node[midway,right] {$\vct u_{\rm out}$};
 \draw[axis,line width=.8pt] (O)--(F) node[midway,above left] {$\vct v_{\rm out}$};
 \fill (O) circle (1.1pt) node[below] {$O_v$};\fill (C) circle (1.1pt) node[below] {$C_v$};
 \fill (I) circle (1.1pt) node[below] {$I$};\fill (F) circle (1.1pt) node[above] {$F$};
 \draw (C)++(180:.4) arc (180:60:.4) node[midway,above] {$\chi$};
\end{tikzpicture}
\caption{原惯性系的速度端点位于以 $\vct V$ 为圆心、半径 $u$ 的圆上；距 $O_v$ 的距离随偏转而改变。}\label{deep8:fig:d8-h12-s}
\end{figure}
图\ref{deep8:fig:d8-h12-s} 中，偏转越大，出射速度端点离 $O_v$ 越远。但实体半径限制了可达到的偏转。行星系的吸引双曲线满足
\[
 e=1+\frac{r_pu^2}{\mu_p},\qquad
 \sin\frac\chi2=\frac1e.
\]
要求 $r_p>R$，得到
\[
 \chi<\chi_{\rm graze}=2\arcsin\frac1{1+Ru^2/\mu_p}.
\]
等号对应擦过表面，不属于严格不接触的轨道。因此最大值一般不取到，偏转角的上确界为 $\chi_{\rm graze}$。代入能量式，
\[
 \boxed{\sup\frac{\Delta K}m
 =\frac{2Vu}{(1+Ru^2/\mu_p)^2}}.
\]
与之对应的瞄准距离下界为
\[
 b_{\rm graze}=R\sqrt{1+\frac{2\mu_p}{Ru^2}},\qquad b>b_{\rm graze}.
\]
若把天体视作没有半径的点质量，形式上可令 $\chi\to\pi$，能量增量趋于 $2Vu$；有限半径排除了这一极限。

完整二体模型中，飞行器获得的能量来自行星的平动能，行星会有极小的反冲。本题把行星速度视作恒量，是飞行器质量远小于行星质量的近似；不能据此认为飞行器在原惯性系中的机械能守恒。在太阳系应用中，若遭遇过程足够短，太阳势在进、出遭遇区时的差别可忽略，上式才可近似解释为日心轨道比能的改变。
\par\endgroup
\end{solution}


\end{document}
